\documentclass[10pt]{article}
\usepackage[T1]{fontenc}
\usepackage[utf8]{inputenc}
\usepackage[margin=1in]{geometry}
\usepackage{amsmath,amssymb,amsthm,mathtools}
\usepackage{mathrsfs}
\usepackage{bm}
\usepackage{centernot}
\usepackage{enumitem}
\usepackage{newunicodechar}
\usepackage{xcolor}
\definecolor{linkblue}{RGB}{20,70,140}
\usepackage[colorlinks=true,linkcolor=linkblue,urlcolor=linkblue,citecolor=linkblue,
            pdfencoding=auto,psdextra]{hyperref}

\setlength{\parindent}{0pt}
\setlength{\parskip}{0.5em}
\setlist{topsep=0.3em,itemsep=0.2em,parsep=0pt,leftmargin=1.4em}
\allowdisplaybreaks
\sloppy
\emergencystretch=3em
\hbadness=10000
\vbadness=10000

\title{Collatz Approach Campaign\\[0.3em]\large Rounds 6--72}
\date{}

\newunicodechar{–}{\textendash{}}
\newunicodechar{“}{``}
\newunicodechar{”}{''}

\begin{document}
\maketitle
\tableofcontents
\newpage

\clearpage
\section{Round 6: approach 1 --- Adelic two-ended renormalization of valuation words}
\subsection*{Prompt}
\begin{quote}
Complete a full proof of the Collatz conjecture using the following approach:

\begin{enumerate}
\item Adelic two-ended renormalization of valuation words
\end{enumerate}

\textbf{Hybrid \textperiodcentered{} B}

A counterexample word must pass three tests at every scale:

archimedean growth, 2-adic seed realizability, 3-adic endpoint/carry compatibility.

Construct a renormalization state combining a normalized height such as

h\_j = A\_j - j log\_2 3,

a 2\textasciicircum{}r-residue representing the admissible seed cylinder, and a 3\textasciicircum{}r-residue representing the endpoint carry. Your terminal-suffix locality result suggests that part of the 3-adic state can be updated using only the last r valuations.

The proof route is to show that all nontrivial recurrent states vanish as r -> infinity. The disproof route is to find a compatible projective fixed point whose inverse limit lifts to one ordinary positive integer. Recent multi-place preprints are adjacent to this idea, but I did not find this exact finite-state renormalization.

Take as long as you need.
\end{quote}
\subsection*{Response}

\subsection*{Verdict}


I pushed the proposed adelic program to an exact renormalization theorem, but it does \textbf{not} currently produce a valid full proof of Collatz. Two things happen:

\begin{enumerate}
\item The hoped-for statement that ``all nontrivial recurrent states disappear as the precision tends to infinity'' is false unless one inserts an \textbf{ordinary-integer lift condition}. There are abundant nontrivial recurrent states in   \[
   \mathbb R_{>0}\cap\mathbb Z_2\cap\mathbb Z_3
   \] that satisfy every finite \(2\)-adic, \(3\)-adic, and archimedean compatibility test but are rational rather than positive integers.

\item Once the integer-lift condition is imposed, the \(2\)-adic and normalized \(3\)-adic coordinates of \textbf{every ordinary orbit}\textemdash{}convergent or hypothetically divergent\textemdash{}eventually collapse to zero. The remaining branch is an arbitrarily long zero-carry ordinary orbit segment. Excluding that branch is a stopping-time theorem closely related to the still-unproved coefficient stopping-time conjecture.
\end{enumerate}

The closest recent paper I found is Oliver Kramer's July 2026 preprint, which already combines real drift, a \(2\)-adic starting representative, and a \(3\)-adic endpoint representative. It proves necessary vanishing conditions for genuine integer orbits and explicitly says the construction is not a verification method. [\hypertarget{cit:R6:1}{}\hyperlink{ref:R6:1}{1}] The exact valuation-cylinder construction below sharpens its starting modulus from \(2^{A_k}\), which records divisibility, to \(2^{A_k+1}\), which records the \textbf{exact} final valuation.

What follows is the exact mathematics obtained from the approach, including the strongest reduction I could prove.


\bigskip



\subsubsection*{1. Exact valuation cylinders}


Let

\[
U(x)=\frac{3x+1}{2^{\nu_2(3x+1)}}
\]

act on positive odd integers. For a finite valuation word

\[
w=(a_1,\ldots,a_k),\qquad a_i\geq 1,
\]

write

\[
A_0=0,\qquad A_i=\sum_{j=1}^{i}a_j,
\]

and define

\[
C_0=0,\qquad C_i=3C_{i-1}+2^{A_{i-1}}.
\]

Equivalently,

\[
C_i=\sum_{j=0}^{i-1}3^{i-1-j}2^{A_j}.
\]

Repeated substitution gives the affine identity

\[
2^{A_i}U^i(n)=3^i n+C_i.
\tag{1}
\]

Finite parity or valuation words determining residue cylinders is standard in the \(2\)-adic formulation of Collatz. [\hypertarget{cit:R6:2}{}\hyperlink{ref:R6:2}{2}] The following exact version is what the renormalization needs.


\paragraph*{Exact cylinder theorem}


Define \(R_w\) to be the unique integer satisfying

\[
1\leq R_w<2^{A_k+1}
\]

and

\[
3^kR_w+C_k\equiv 2^{A_k}\pmod{2^{A_k+1}}.
\tag{2}
\]

Equivalently,

\[
R_w\equiv
3^{-k}\bigl(2^{A_k}-C_k\bigr)
\pmod{2^{A_k+1}}.
\tag{3}
\]

Then the positive odd integers having \textbf{exactly} the first \(k\) valuations \(a_1,\ldots,a_k\) are

\[
n=R_w+2^{A_k+1}t,\qquad t\geq0.
\tag{4}
\]

If

\[
M_w:=\frac{3^kR_w+C_k}{2^{A_k}},
\tag{5}
\]

then \(M_w\) is odd,

\[
1\leq M_w<2\cdot 3^k,
\tag{6}
\]

and the corresponding endpoint is

\[
U^k(n)=M_w+2\cdot3^k t.
\tag{7}
\]

Thus every word determines an exact two-ended arithmetic progression

\[
\boxed{
(n,U^k(n))
=
\bigl(R_w+2^{A_k+1}t,\;
      M_w+2\cdot3^k t\bigr).
}
\tag{8}
\]


\paragraph*{Why the extra factor \(2\) is necessary}


A congruence modulo \(2^{A_k}\) forces

\[
2^{A_k}\mid3^kn+C_k,
\]

but it does not force the quotient to be odd. Exact valuation requires

\[
3^kn+C_k\equiv2^{A_k}\pmod{2^{A_k+1}}.
\]

For example, the one-letter word \(w=(1)\) is realized exactly by

\[
n\equiv3\pmod4,
\]

not by every odd number modulo \(2\).


\paragraph*{Proof of exactness}


Put

\[
N_i=3^in+C_i.
\]

For \(i<k\),

\[
N_k
=
3^{k-i}N_i+
\sum_{j=i}^{k-1}3^{k-1-j}2^{A_j}.
\tag{9}
\]

The congruence (2) implies \(2^{A_k}\mid N_k\). Reducing (9) modulo \(2^{A_i}\) shows that \(2^{A_i}\mid N_i\).

Divide (9) by \(2^{A_i}\) and reduce modulo \(2\). Because \(A_j>A_i\) for \(j>i\), only the term \(j=i\) survives from the sum. Since \(N_k/2^{A_i}\) is even,

\[
0\equiv
\frac{N_i}{2^{A_i}}+1\pmod2.
\]

Therefore \(N_i/2^{A_i}\) is odd. Hence every prescribed valuation is exact.

Conversely, an exact realization has \(N_k/2^{A_k}\) odd, which gives (2).

The endpoint bound (6) follows inductively. If a prefix has canonical endpoint \(M<2\cdot3^j\), appending \(a\) requires the unique \(t\in\{0,\ldots,2^a-1\}\) satisfying

\[
\frac{3M+1}{2}+3^{j+1}t
\equiv2^{a-1}\pmod{2^a}.
\]

The new endpoint is

\[
M'=\frac{3(M+2\cdot3^jt)+1}{2^a},
\]

and the stated ranges give \(M'<2\cdot3^{j+1}\).


\bigskip



\subsubsection*{2. Terminal-suffix locality in its full form}


Let \(v\) be the terminal suffix of \(w\) of length \(r\), with total valuation \(B\) and affine constant \(C_v\). Its canonical endpoint \(M_v\) satisfies

\[
2^B M_v\equiv C_v\pmod{3^r}.
\tag{10}
\]

More strongly,

\[
\boxed{
M_w\equiv M_v\pmod{2\cdot3^r}.
}
\tag{11}
\]

Thus the endpoint residue modulo \(3^r\) depends only on the final \(r\) valuations.

In particular,

\[
\begin{aligned}
3^r\mid M_w+1
&\iff 3^r\mid M_v+1\\
&\iff 3^r\mid C_v+2^B.
\end{aligned}
\tag{12}
\]

This is precisely the terminal-suffix locality theorem, now placed inside the exact two-ended cylinder.

There is a useful strengthening. Since \(M_v\) is odd and

\[
1\leq M_v<2\cdot3^r,
\]

condition (12) forces

\[
M_v=2\cdot3^r-1.
\tag{13}
\]

So a length-\(r\) high-carry suffix does not merely give a congruence: it gives an explicit exponential lower bound for every compatible endpoint.


\bigskip



\subsubsection*{3. Exact concatenation: the adelic bridge equation}


Let \(u\) and \(v\) be two words, with

\[
|u|=p,\qquad |v|=q.
\]

Write their canonical pairs as

\[
(R_u,M_u),\qquad (R_v,M_v).
\]

There are unique integers

\[
0\leq t<2^{A(v)},\qquad 0\leq s<3^p
\]

such that

\[
\boxed{
M_u+2\cdot3^p t
=
R_v+2^{A(v)+1}s.
}
\tag{14}
\]

Indeed, after division by \(2\), this is

\[
3^pt-2^{A(v)}s=\frac{R_v-M_u}{2},
\tag{15}
\]

and \(3^p\) is invertible modulo \(2^{A(v)}\).

The canonical pair of the concatenation \(uv\) is

\[
\boxed{
R_{uv}=R_u+2^{A(u)+1}t,
\qquad
M_{uv}=M_v+2\cdot3^q s.
}
\tag{16}
\]

Equation (14) is the exact two-ended adelic matching condition:

\begin{itemize}
\item the endpoint progression of \(u\) has spacing \(2\cdot3^p\);
\item the seed progression of \(v\) has spacing \(2^{A(v)+1}\);
\item the bridge value is their unique positive intersection in the fundamental rectangle.
\end{itemize}


\paragraph*{Normalized renormalization coordinates}


Set

\[
\xi_w=\frac{R_w}{2^{A(w)+1}},
\qquad
\eta_w=\frac{M_w}{2\cdot3^{|w|}}.
\tag{17}
\]

Then

\[
0<\xi_w<\eta_w<1
\]

and

\[
\eta_w-\xi_w
=
\frac{C_w}{2^{A(w)+1}3^{|w|}}.
\tag{18}
\]

Since \(A_j-A_k\leq-(k-j)\),

\[
\begin{aligned}
\eta_w-\xi_w
&=
\frac12\sum_{j=0}^{k-1}
  \frac{2^{A_j-A_k}}{3^{j+1}}\\
&\leq
2^{-k-1}\left(1-\left(\frac23\right)^k\right)
<2^{-k-1}.
\end{aligned}
\tag{19}
\]

Thus the two normalized ends synchronize exponentially.

Define the bridge coordinate

\[
z_{u,v}
=
\frac{t+\eta_u}{2^{A(v)}}
=
\frac{s+\xi_v}{3^p}.
\tag{20}
\]

Then

\[
\xi_{uv}=\frac{t+\xi_u}{2^{A(v)}},
\qquad
\eta_{uv}=\frac{s+\eta_v}{3^p},
\tag{21}
\]

and hence

\[
z_{u,v}-\xi_{uv}
=
\frac{\eta_u-\xi_u}{2^{A(v)}},
\tag{22}
\]

\[
\eta_{uv}-z_{u,v}
=
\frac{\eta_v-\xi_v}{3^p}.
\tag{23}
\]

Writing

\[
\delta_w=\eta_w-\xi_w,
\]

we obtain the exact contraction cocycle

\[
\boxed{
\delta_{uv}
=
2^{-A(v)}\delta_u+3^{-p}\delta_v.
}
\tag{24}
\]

This is the desired adelic two-ended renormalization law.


\bigskip



\subsubsection*{4. The archimedean height identity}


Let

\[
\alpha=\log_2 3,
\qquad
h_i=A_i-\alpha i.
\tag{25}
\]

Since

\[
\frac{2^{A_i}}{3^i}=2^{h_i},
\]

we have

\[
\frac{C_k}{3^k}
=
\frac13\sum_{i=0}^{k-1}2^{h_i}.
\tag{26}
\]

For an actual seed \(n\),

\[
\boxed{
U^k(n)
=
2^{-h_k}
\left(
n+\frac13\sum_{i=0}^{k-1}2^{h_i}
\right).
}
\tag{27}
\]

If \(U^k(n)\geq n\) and \(h_k>0\), then

\[
n\bigl(2^{h_k}-1\bigr)
\leq
\frac13\sum_{i=0}^{k-1}2^{h_i}.
\tag{28}
\]

For a first crossing,

\[
h_i<0\quad(1\leq i<k),
\qquad h_k>0,
\]

so

\[
\frac13\sum_{i=0}^{k-1}2^{h_i}<\frac{k}{3}.
\tag{29}
\]

This gives the fixed-point bound

\[
n<
\frac{k}{3(2^{h_k}-1)}.
\tag{30}
\]


\bigskip



\subsubsection*{5. The seed-lift dichotomy}


For nested prefixes \(w_k=(a_1,\ldots,a_k)\), the exact cylinders are nested, so

\[
R_{k+1}
=
R_k+2^{A_k+1}t_k,
\qquad
0\leq t_k<2^{a_{k+1}}.
\tag{31}
\]

Thus \(R_k\) is nondecreasing.

If it changes at stage \(k+1\), then

\[
R_{k+1}\geq2^{A_k+1}\geq2^{k+1}.
\tag{32}
\]

Consequently:


\paragraph*{Seed stabilization theorem}


For every infinite valuation word, precisely one of the following occurs:

\begin{enumerate}
\item \(R_k\) is eventually constant; then the infinite word is realized by one ordinary positive odd integer.

\item \(R_k\) changes infinitely often; then

\end{enumerate}
\[
   \limsup_{k\to\infty}\frac{\log_2 R_k}{k}\geq1.
   \tag{33}
   \]

In particular, any subexponential canonical seed growth forces \textbf{literal eventual constancy}, not merely convergence in \(\mathbb Z_2\).

The dual statement holds for canonical endpoints when a word is extended to the left. If \(M_r\) is the canonical endpoint of a terminal word of length \(r\), then

\[
M_{r+1}=M_r+2\cdot3^r s_r,
\qquad s_r\in\{0,1,2\}.
\tag{34}
\]

If \(M_r\) changes infinitely often, its limsup growth rate is at least \(3^r\); any genuinely sub-\(3^r\) growth forces eventual endpoint constancy.

This is stronger than the necessary residue-rate statement in the adjacent 2026 search framework: once the exact nested cylinders are used, a vanishing starting rate is essentially equivalent to an ordinary integer lift. [\hyperlink{ref:R6:1}{1}]


\bigskip



\subsubsection*{6. Why vanishing residues do not prove Collatz}


Suppose \(n\) is any fixed positive odd integer, and let \(w_k\) be its actual valuation prefixes.

Once

\[
2^{A_k+1}>n,
\]

the exact seed representative is

\[
R_k=n.
\tag{35}
\]

Therefore

\[
M_k=U^k(n).
\tag{36}
\]

Moreover,

\[
U(x)+1
\leq
\frac32(x+1),
\]

so

\[
U^k(n)+1
\leq
(n+1)\left(\frac32\right)^k.
\tag{37}
\]

Hence

\[
\xi_k
=
\frac{n}{2^{A_k+1}}
\leq
\frac{n}{2^{k+1}},
\tag{38}
\]

and

\[
\eta_k
=
\frac{U^k(n)}{2\cdot3^k}
<
\frac{n+1}{2^{k+1}}.
\tag{39}
\]

Thus

\[
\boxed{
(\xi_k,\eta_k)\longrightarrow(0,0)
}
\tag{40}
\]

exponentially for \textbf{every ordinary orbit}, without using convergence.

A hypothetical divergent integer orbit therefore lies in exactly the same asymptotic \(2\)-adic/\(3\)-adic normalized state as a convergent orbit. The dangerous state is not a nonzero recurrent residue state. It is the zero-residue, zero-carry state.

This is the first fundamental obstruction to the proposed proof route.


\bigskip



\subsubsection*{7. Exact height cannot be made recurrent without losing information}


If the state retains the exact real height \(h_k\), then no two different times have the same height. Indeed,

\[
h_j=h_k
\]

would imply

\[
A_j-A_k=\alpha(j-k).
\]

The left side is an integer, whereas \(\alpha=\log_2 3\) is irrational. Therefore \(j=k\).

So an exact state containing \(h_k\) has \textbf{no recurrence at all}, including along the trivial orbit.

To get a finite state, one must bin, quotient, or recenter the height. But then the state no longer retains the exact sign of

\[
U^k(n)-n
=
\frac{C_k-(2^{A_k}-3^k)n}{2^{A_k}},
\tag{41}
\]

because that sign depends on the unbounded ordinary lift \(n\), not merely on its finite residues.


\bigskip



\subsubsection*{8. Finite residue states necessarily alias growth and descent}


For a contracting word \(w\), set

\[
P_w=2^{A(w)+1},
\qquad
Q_w=2\cdot3^{|w|}.
\]

Then \(P_w>Q_w\). Its realizations satisfy

\[
n_t=R_w+P_wt,
\qquad
m_t=M_w+Q_wt,
\tag{42}
\]

so

\[
m_t-n_t
=
(M_w-R_w)+(Q_w-P_w)t.
\tag{43}
\]

If \(M_w\geq R_w\), the canonical lift is non-descending, but all sufficiently large lifts descend.

This cannot be repaired by increasing finite residue precision. For any prescribed \(r\), choose \(t\) divisible by sufficiently high powers of \(2\) and \(3\). Then

\[
n_t\equiv R_w\pmod{2^r},
\qquad
m_t\equiv M_w\pmod{3^r},
\tag{44}
\]

while the archimedean sign in (43) changes.

A strict numerical certificate is the word

\[
w=(4,1,1,1,1,2,2,1,2,1,1,2,1,1,1,2,3).
\]

It has

\[
|w|=17,\qquad A(w)=27,\qquad C_w=1\,106\,233\,681,
\]

and canonical pair

\[
R_w=165,\qquad M_w=167.
\]

Indeed,

\[
3^{17}\cdot165+C_w=2^{27}\cdot167.
\]

Therefore

\[
n_t=165+2^{28}t,
\qquad
m_t=167+2\cdot3^{17}t.
\]

The \(t=0\) realization rises from \(165\) to \(167\), whereas every \(t\geq1\) realization descends because

\[
2^{28}>2\cdot3^{17}.
\]

Yet for any finite \(2\)-adic and \(3\)-adic precision one can choose a nonzero \(t\) making both endpoint residues identical to those at \(t=0\).

A 2026 logic preprint proves a related limitation: arbitrary-step generalized Collatz transitions cannot be recognized by a finite automaton reading base-\(2\) representations. [\hypertarget{cit:R6:3}{}\hyperlink{ref:R6:3}{3}] Here the obstruction is visible directly: the discarded lift parameter is archimedean and unbounded.


\bigskip



\subsubsection*{9. Nontrivial adelic fixed points are abundant}


For any nonempty block \(w\), the associated affine map is

\[
F_w(x)=\frac{3^{|w|}x+C_w}{2^{A(w)}}.
\tag{45}
\]

Its unique rational fixed point is

\[
\boxed{
x_w^*=\frac{C_w}{2^{A(w)}-3^{|w|}}.
}
\tag{46}
\]

The denominator is odd and not divisible by \(3\). Hence

\[
x_w^*\in\mathbb Z_2\cap\mathbb Z_3.
\]

If

\[
A(w)>|w|\log_2 3,
\]

then \(x_w^*>0\) at the real place. It realizes the word \(w\) exactly in \(\mathbb Z_2\), returns to itself after the block, and therefore realizes \(w^\infty\).

Thus every contracting block gives a positive rational

\[
2\text{-adic}/3\text{-adic}/\infty
\]

projective fixed point.

Moreover, such fixed points exist arbitrarily close to critical drift. Choose \(k\) for which the fractional part of \(k\log_2 3\) is close to \(1\), set

\[
A_i=\lfloor i\log_2 3\rfloor
\quad(0\leq i<k),
\]

and

\[
A_k=\lceil k\log_2 3\rceil.
\]

Then

\[
-1<h_i<0\quad(1\leq i<k),
\qquad
0<h_k\longrightarrow0
\]

along a subsequence. Each resulting first-crossing block has a positive rational adelic fixed point.

It is an ordinary positive integer precisely when

\[
2^{A(w)}-3^{|w|}\mid C_w.
\tag{47}
\]

But that is exactly the positive-cycle divisibility condition. If the quotient is an integer other than \(1\), it is a nontrivial Collatz cycle.

Therefore:

\[
\boxed{
\text{Projective fixed points are plentiful; ordinary integral fixed points are the cycle problem.}
}
\]

The hoped-for disappearance of all nontrivial recurrent states is consequently false before integrality is imposed. And after integrality is imposed, ruling out fixed points is already the nontrivial-cycle part of Collatz.

A divergent counterexample need not yield any periodic fixed point at all: its valuation sequence could be aperiodic. A recurrent component of a finite nondeterministic graph need not contain the particular counterexample path as a periodic orbit.


\bigskip



\subsubsection*{10. The strongest positive result: two-ended stabilization}


Although this does not close the conjecture, the adelic construction gives a strong theorem about any contracting word that fails to descend.

Let

\[
\Delta_w=2^{A(w)}-3^{|w|}>0
\]

and define the real fixed point

\[
x_w^*=\frac{C_w}{\Delta_w}.
\]

Since

\[
M_w-R_w
=
\left(1-\frac{3^{|w|}}{2^{A(w)}}\right)
(x_w^*-R_w),
\tag{48}
\]

if

\[
M_w\geq R_w,
\tag{49}
\]

then

\[
R_w\leq M_w\leq x_w^*.
\tag{50}
\]

Now let \(u\) range over prefixes of \(w\), and \(v\) over suffixes. Prefix representatives are nondecreasing under extension, while suffix endpoint representatives are nondecreasing under extension to the left. Hence

\[
\boxed{
R_u\leq x_w^*,
\qquad
M_v\leq x_w^*
}
\tag{51}
\]

at \textbf{every scale}.

This is the exact three-place obstruction requested in the prompt:

\begin{itemize}
\item every prefix \(2\)-adic seed cylinder must have an unusually small canonical representative;
\item every suffix \(3\)-adic endpoint cylinder must have an unusually small canonical representative;
\item both are bounded by the same archimedean fixed-point threshold.
\end{itemize}


\paragraph*{Carries are forced into boundary layers}


If a prefix representative changes at length \(j\), the new representative is at least \(2^j\). Therefore, once

\[
2^{p+1}>x_w^*,
\tag{52}
\]

all later prefix carries vanish and

\[
R_{w[1:j]}=R_w
\qquad(j\geq p).
\tag{53}
\]

Dually, once

\[
2\cdot3^q>x_w^*,
\tag{54}
\]

all earlier suffix-extension carries vanish and

\[
M_{w[j:k]}=M_w
\qquad(k-j+1\geq q).
\tag{55}
\]

For every split \(w=uv\) satisfying both depth conditions, the bridge equation has

\[
t=s=0,
\]

and therefore

\[
\boxed{
M_u=R_v.
}
\tag{56}
\]

The actual intermediate value is simultaneously

\begin{itemize}
\item the canonical \(3\)-adic endpoint of the entire left block, and
\item the canonical \(2\)-adic seed of the entire right block.
\end{itemize}


\paragraph*{First-crossing words have only logarithmic boundary layers}


If \(w\) is a first-crossing word, then by (29)

\[
x_w^*
<
\frac{k}{3(2^{h_k}-1)}.
\tag{57}
\]

If \(h_k\geq1\), this is \(<k/3\).

If \(0<h_k<1\), then

\[
A_k=\lceil k\log_2 3\rceil.
\]

A standard lower bound for the nonzero linear form

\[
A_k\log2-k\log3
\]

gives constants \(c,\kappa>0\) such that

\[
A_k\log2-k\log3\geq c k^{-\kappa}.
\tag{58}
\]

Such polynomial lower bounds follow from the theory of linear forms in two logarithms. [\hypertarget{cit:R6:4}{}\hyperlink{ref:R6:4}{4}] Consequently,

\[
x_w^*\ll k^{\kappa+1}.
\tag{59}
\]

Therefore every nonzero prefix or suffix carry in a hypothetical first-crossing paradoxical word is confined to

\[
O(\log k)
\]

letters at its two ends. The entire middle\textemdash{}of length

\[
k-O(\log k)
\]

\textemdash{}is an exact zero-carry ordinary orbit segment satisfying (56).

This appears to be the cleanest rigorous output of the adelic approach.


\bigskip



\subsubsection*{11. How terminal-suffix locality prunes the state tree}


Let \(v\) be a terminal suffix of length \(r\). From (51), a non-descending contracting word must satisfy

\[
M_v\leq x_w^*.
\tag{60}
\]

Therefore the exact suffix test

\[
M_v>x_w^*
\tag{61}
\]

immediately proves descent.

In the special high-carry case,

\[
3^r\mid C_v+2^{A(v)},
\]

terminal-suffix locality gives

\[
M_v=2\cdot3^r-1.
\]

Hence:


\paragraph*{High-carry exclusion}


If

\[
3^r\mid C_v+2^{A(v)}
\]

and

\[
2\cdot3^r-1>x_w^*,
\tag{62}
\]

then

\[
M_w<R_w.
\]

So the terminal-suffix theorem genuinely eliminates an infinite and explicitly recognizable family of candidate words.

What it does \textbf{not} prove is that every long first-crossing word has such a high-carry suffix. Low-carry suffix states remain, and those are exactly where the zero-carry bulk can live.


\bigskip



\subsubsection*{12. The exact remaining theorem}


The full conjecture would follow from the following two statements.


\paragraph*{Barrier escape statement}


Let \(w_\infty=(a_1,a_2,\ldots)\) satisfy

\[
A_j-j\log_2 3<0
\qquad\text{for every }j\geq1.
\tag{BE}
\]

Then its exact seed representatives \(R_j\) change infinitely often.

By the seed stabilization theorem, this says that no ordinary positive integer can realize an infinite valuation word whose coefficient remains at least \(1\) forever.


\paragraph*{First-crossing separation statement}


Let \(w=(a_1,\ldots,a_k)\) satisfy

\[
A_j-j\log_2 3<0\quad(1\leq j<k),
\]

and

\[
A_k-k\log_2 3>0.
\]

Then

\[
M_w<R_w,
\tag{FS}
\]

except for the one-letter word \(w=(2)\), for which

\[
R_w=M_w=1.
\]

These are the exact missing zero-carry assertions.


\paragraph*{Conditional proof of Collatz}


Assume (BE) and (FS). Let \(n>1\) be odd, with valuation word \(a_1,a_2,\ldots\).

If its height never becomes positive, then its canonical seed representatives eventually equal \(n\), contradicting (BE).

Hence there is a first crossing \(k\). Let \(w\) be that prefix. Every realization is

\[
n=R_w+2^{A_k+1}t,
\]

with endpoint

\[
U^k(n)=M_w+2\cdot3^kt.
\]

Since the block is contracting,

\[
2^{A_k+1}>2\cdot3^k.
\]

If \(w\neq(2)\), statement (FS) gives \(M_w<R_w\), and therefore

\[
U^k(n)<n.
\]

For \(w=(2)\), \(n>1\) implies \(t\geq1\), and

\[
U(n)=1+6t<1+8t=n.
\]

Thus every odd \(n>1\) has a smaller odd iterate. Strong induction gives convergence to \(1\). Even starting values reduce to odd ones by division by powers of \(2\).

That would be a complete proof.


\bigskip



\subsubsection*{13. Why this is still the unresolved core}


The first-crossing separation statement says that once the multiplicative coefficient first becomes contracting, the additive term cannot keep the orbit above its initial minimum. This is a compressed odd-to-odd version of the stopping-time problem.

Rozier and Terracol's 2026 paper formulates Terras's coefficient stopping-time conjecture and the equivalent assertion that the first term of a paradoxical sequence is never its smallest term. It remains conjectural; their work instead classifies and computationally constrains paradoxical sequences. [\hypertarget{cit:R6:5}{}\hyperlink{ref:R6:5}{5}]

The adelic renormalization proves substantially more structure around a hypothetical violation:

\[
\text{small two-adic prefixes}
+
\text{small three-adic suffixes}
+
\text{logarithmic carry boundaries}
+
\text{a zero-carry middle}.
\]

But the zero-carry middle is an ordinary Collatz orbit segment. Showing that this middle cannot remain above its initial value is precisely what has not been proved.


\subsubsection*{Final status}


A full proof has \textbf{not} been completed. The exact result of the approach is:

\[
\boxed{
\begin{array}{c}
\text{Every candidate first-crossing counterexample has all nonzero}\\
\text{\(2\)-adic and \(3\)-adic carries confined to logarithmic boundary layers,}\\
\text{and an arbitrarily long exact zero-carry ordinary core.}
\end{array}
}
\]

Terminal-suffix locality eliminates the high-\(3\)-adic-carry part of that state space. Exact seed nesting eliminates nonordinary low-growth \(2\)-adic ghosts. What remains is the frozen zero-carry branch, expressed precisely by (BE) and (FS).

The projective-fixed-point route can find or exclude cycles, but it does not automatically address aperiodic divergence. The proof route is therefore reduced to one sharp task: establish barrier escape and first-crossing separation for the zero-carry ordinary branch. Claiming that recurrence alone supplies those statements would be the missing\textemdash{}and invalid\textemdash{}step.
\subsection*{References}
\begin{enumerate}[leftmargin=2.2em]
\item \hypertarget{ref:R6:1}{}\hyperlink{cit:R6:1}{\texttt{turn593659view0}}
\item \hypertarget{ref:R6:2}{}\hyperlink{cit:R6:2}{\texttt{turn593659view1}}
\item \hypertarget{ref:R6:3}{}\hyperlink{cit:R6:3}{\texttt{turn108155view0}}
\item \hypertarget{ref:R6:4}{}\hyperlink{cit:R6:4}{\texttt{turn655263search5}}
\item \hypertarget{ref:R6:5}{}\hyperlink{cit:R6:5}{\texttt{turn425183view0}}
\end{enumerate}

\clearpage
\section{Round 7: approach 2 --- Furstenberg-style (x2, x3) measure rigidity}
\subsection*{Prompt}
\begin{quote}
Complete a full proof of the Collatz conjecture using the following approach:

\begin{enumerate}
\item Furstenberg-style (x2, x3) measure rigidity
\end{enumerate}

\textbf{Speculative \textperiodcentered{} B}

Assume a divergent orbit exists and form empirical measures from its logarithmic height, binary cylinder, and ternary carry state on an adelic solenoid. Try to derive an invariant joining with actions related to multiplication by 2 and 3.

A rigidity dichotomy might say:

\begin{itemize}
\item positive entropy forces a Haar-like measure, contradicting the exact carry identities or positivity;
\item zero entropy forces algebraic or eventually periodic structure, reducing the orbit to a forbidden cycle equation.
\end{itemize}

The difficulty is establishing enough recurrence to produce a probability measure rather than losing all mass at infinity.

Take as long as you need.
\end{quote}
\subsection*{Response}

\subsection*{Result of the attempt}


Carrying the Furstenberg program through produces a rigorous reduction, but \textbf{not a proof of Collatz}. The decisive failure is precise:

\begin{enumerate}
\item A divergent orbit does generate a probability measure on a compact adelic state space.
\item That measure is invariant under \textbf{one variable-time skew transformation} \[
   (a,z)\longmapsto \bigl(\sigma a,\alpha_3\alpha_2^{-a_0}z\bigr),
   \] not under multiplication by \(2\) and multiplication by \(3\) separately.
\item Thickening the orbit enough to obtain genuine \((\times2,\times3)\)-invariance makes the Collatz orbit a zero-density boundary and loses the carry constraint.
\item Keeping the ordinary positive-integer lift or absolute height avoids that loss, but then all empirical mass escapes to infinity.
\item The proposed implication ``zero entropy \(\Rightarrow\) algebraic or eventually periodic'' is false for the ambient skew system and is stronger than presently available \(\times2,\times3\) rigidity. The positive-entropy theorem of Rudolph\textendash{}Johnson and its solenoid extensions require genuine higher-rank invariance; the unresolved part of Furstenberg's measure problem is essentially the zero-entropy case. [\hypertarget{cit:R7:1}{}\hyperlink{ref:R7:1}{1}, \hypertarget{cit:R7:2}{}\hyperlink{ref:R7:2}{2}]
\end{enumerate}

What follows is the maximal rigorous theorem supplied by the approach, followed by the exact additional statements that would turn it into a proof.


\bigskip



\subsubsection*{1. Exact orbit identities}


Use the accelerated map on positive odd integers: \[
x_{j+1}=U(x_j)=\frac{3x_j+1}{2^{a_j}},
\qquad
a_j=v_2(3x_j+1)\geq 1.
\]

Suppose first that this orbit is divergent. Since a deterministic orbit cannot revisit a value without becoming periodic, the \(x_j\) are distinct. A sequence of distinct positive integers has only finitely many terms below any fixed bound, so \[
x_j\longrightarrow \infty.
\]

Put \[
A_0=0,\qquad A_j=\sum_{i=0}^{j-1}a_i.
\] Iterating the affine recurrence gives \[
2^{A_j}x_j=3^jx_0+C_j,
\] where \[
C_j=\sum_{k=0}^{j-1}3^{j-1-k}2^{A_k},
\qquad
C_{j+1}=3C_j+2^{A_j}.
\]

Taking logarithms at one step gives \[
\log x_{j+1}-\log x_j
 =\log3-a_j\log2+
 \log\left(1+\frac1{3x_j}\right).
\] Consequently, \[
\log x_N
=
\log x_0+N\log3-A_N\log2+E_N,
\] where \[
E_N=\sum_{j=0}^{N-1}\log\left(1+\frac1{3x_j}\right).
\]

Because \(x_j\to\infty\), \[
\frac{E_N}{N}\longrightarrow0.
\] Since \(\log x_N\geq0\), we obtain the fundamental drift bound \[
\boxed{\displaystyle
\limsup_{N\to\infty}\frac{A_N}{N}
\leq \mu_*:=\log_2 3
=1.5849625007\ldots .}
\tag{1}
\]

Thus every divergent accelerated orbit must have an average valuation no greater than \(\log_2 3\).

In the unaccelerated shortcut map, each symbol \(a_j\) expands to the parity block \[
1\,0^{a_j-1}.
\] Hence a stationary valuation process with mean \(m\leq\mu_*\) has odd-step frequency \[
\frac1m\geq\frac1{\log_2 3}
=\frac{\log2}{\log3}
=0.63092975\ldots>\frac12.
\tag{2}
\]

This is already exactly incompatible with Haar measure for the binary shift. The missing problem is to justify why an orbit-derived measure should be Haar.


\bigskip



\subsubsection*{2. A probability measure really can be constructed}


Let \[
\Sigma=\mathbb N_{\geq1}^{\mathbb N}
\] and let \[
a=(a_0,a_1,\ldots)
\] be the valuation sequence of the hypothetical divergent orbit. Consider \[
\nu_N=\frac1N\sum_{j=0}^{N-1}\delta_{\sigma^ja}.
\]

These measures are tight, despite the countable alphabet.

Indeed, (1) implies that for some finite \(K\), \[
A_n\leq Kn\qquad(n\geq1).
\] At coordinate \(r\), \[
\begin{aligned}
\nu_N\{b:b_r>M_r\}
&\leq
\frac1{NM_r}\sum_{j=0}^{N-1}a_{j+r}\\
&\leq
\frac{A_{N+r}}{NM_r}
\leq
\frac{K(N+r)}{NM_r}
\leq
\frac{K(1+r)}{M_r}.
\end{aligned}
\] Choose \(M_r\) so rapidly increasing that \[
\sum_{r\geq0}\frac{K(1+r)}{M_r}<\varepsilon.
\] Then \[
\prod_{r\geq0}\{1,\ldots,M_r\}
\] is compact in \(\Sigma\) and has \(\nu_N\)-measure at least \(1-\varepsilon\), uniformly in \(N\).

Therefore a subsequence converges to a shift-invariant probability \(\nu\). Moreover, truncating \(a_0\) and passing to the limit gives \[
\boxed{\displaystyle
\int a_0\,d\nu\leq\mu_* .}
\tag{3}
\]

There may be loss of first moment through rare extremely large valuations, so equality is not justified. The inequality is all that is needed below.


\bigskip



\subsubsection*{3. Exact adelic linearization}


Introduce the \(S\)-adic solenoid \[
X_6=
\frac{\mathbb R\times\mathbb Q_2\times\mathbb Q_3}
{\Delta(\mathbb Z[1/6])}.
\] Let \(\alpha_2,\alpha_3\) denote multiplication by \(2\) and \(3\), respectively. These are commuting automorphisms of \(X_6\).

Set \[
c_j=\frac{C_j}{2^{A_j}}.
\] Then \[
c_{j+1}=\frac{3c_j+1}{2^{a_j}}.
\]

Define \[
z_j=[(c_j,x_j,x_j)]\in X_6.
\] Because the vector \[
\left(2^{-a_j},2^{-a_j},2^{-a_j}\right)
\] belongs to the diagonal copy of \(\mathbb Z[1/6]\), we have the exact identity \[
\boxed{\displaystyle
z_{j+1}=\alpha_3\alpha_2^{-a_j}z_j.}
\tag{4}
\]

Furthermore, subtracting the diagonal element \(x_j\) gives \[
\begin{aligned}
z_j
&=[(c_j-x_j,0,0)]\\
&=
\left[
\left(
-x_0\frac{3^j}{2^{A_j}},0,0
\right)
\right].
\end{aligned}
\tag{5}
\] Thus \[
z_j=\alpha_3^j\alpha_2^{-A_j}z_0.
\]

Now put \[
F(a,z)=\bigl(\sigma a,\alpha_3\alpha_2^{-a_0}z\bigr).
\] The empirical measures \[
\eta_N=\frac1N\sum_{j=0}^{N-1}
\delta_{(\sigma^ja,z_j)}
\] have weakly convergent subsequences because the symbolic marginals are tight and \(X_6\) is compact. Every limit \(\eta\) satisfies \[
F_*\eta=\eta,
\] since \[
\int\phi\circ F\,d\eta_N-\int\phi\,d\eta_N
=
\frac{\phi(\sigma^Na,z_N)-\phi(a,z_0)}N
\longrightarrow0.
\]

So the mass-escape problem can indeed be overcome after adelic compactification.

But what has been obtained is invariance under \textbf{one skew map \(F\)}, not under the full action generated separately by \(\alpha_2\) and \(\alpha_3\).


\bigskip



\subsubsection*{4. A genuine recurrence theorem for logarithmic height}


Absolute height escapes, but logarithmic height modulo one has a remarkably rigid limit.

Let \[
\theta_j=\{\log_2x_j\},
\qquad
\alpha=\log_2 3.
\] Then \[
\theta_{j+1}
=
\theta_j+\alpha+\varepsilon_j
\pmod1,
\] where \[
\varepsilon_j=
\log_2\left(1+\frac1{3x_j}\right)
\longrightarrow0.
\]

Let \(f\) be continuous on \(\mathbb T=\mathbb R/\mathbb Z\). Then \[
\begin{aligned}
\frac1N\sum_{j<N}
\bigl(f(\theta_j+\alpha)-f(\theta_j)\bigr)
&=
\frac1N\sum_{j<N}
\bigl(f(\theta_j+\alpha)-f(\theta_{j+1})\bigr)\\
&\quad+
\frac{f(\theta_N)-f(\theta_0)}N.
\end{aligned}
\] The first term tends to zero by uniform continuity and \(\varepsilon_j\to0\); the second tends to zero trivially.

Every empirical limit is therefore invariant under rotation by \(\alpha\). Since \(\log_2 3\) is irrational, this rotation is uniquely ergodic. Hence \[
\boxed{\displaystyle
\frac1N\sum_{j=0}^{N-1}
\delta_{\{\log_2x_j\}}
\ \Longrightarrow\ m_{\mathbb T}.}
\tag{6}
\]

Thus every hypothetical divergent accelerated Collatz orbit is equidistributed in base-\(2\) logarithmic phase.

This completely resolves recurrence after quotienting height by powers of \(2\), but it does not control absolute height and therefore does not force two nearly equal adelic states to be the same positive integer.


\bigskip



\subsubsection*{5. A sharp entropy ceiling}


Let \(\nu\) be an ergodic component of an orbit-derived invariant measure satisfying \[
m:=\int a_0\,d\nu\leq\mu_*.
\]

For any stationary process on the positive integers, \[
h_\nu(\sigma)\leq H_\nu(a_0).
\] Among probability distributions on \(\{1,2,\ldots\}\) with mean \(m\), the one-symbol entropy is maximized by the geometric law \[
q_k=\frac1m\left(1-\frac1m\right)^{k-1}.
\] This follows immediately by comparing an arbitrary distribution with \(q\) through nonnegativity of relative entropy.

Therefore \[
H_\nu(a_0)
\leq
H_{\max}(m)
=
m\log m-(m-1)\log(m-1).
\] Since \(m\leq\mu_*\), \[
\boxed{\displaystyle
h_\nu(\sigma)
\leq
h_*:=1.0436328159\ldots
<\log3=1.0986122887\ldots .}
\tag{7}
\]

There is a strict entropy gap \[
\log3-h_*=0.0549794728\ldots .
\]

After expanding each \(a_j\) into the binary block \(10^{a_j-1}\), Abramov's tower formula gives binary entropy \[
h_{\rm bin}=\frac{h_\nu(\sigma)}m.
\] Optimizing over \(1\leq m\leq\mu_*\) gives \[
\boxed{\displaystyle
h_{\rm bin}\leq0.6584589953\ldots
<\log2=0.6931471806\ldots .}
\tag{8}
\]

So an orbit-derived low-mean component cannot possibly be Haar either in ternary time or in binary shortcut time\textemdash{}\textbf{provided its symbolic entropy transfers to the corresponding algebraic action}.

That transfer is precisely what is missing.


\bigskip



\subsubsection*{6. Why Rudolph\textendash{}Johnson does not apply}


The relevant rigidity theorem says, roughly, that a genuinely jointly \(\times2\)- and \(\times3\)-invariant ergodic probability on the circle having positive entropy must be Lebesgue. Higher-rank solenoid versions similarly require a genuine algebraic \(\mathbb Z^d\)-action and then yield a Haar-versus-zero-entropy dichotomy. [\hyperlink{ref:R7:1}{1}, \hyperlink{ref:R7:2}{2}]

Our measure satisfies only \[
F_*\eta=\eta,
\qquad
F(a,z)=
\bigl(\sigma a,\alpha_3\alpha_2^{-a_0}z\bigr).
\]

Disintegrate \[
\eta=\int\eta_a\,d\nu(a).
\] Invariance gives \[
\eta_{\sigma a}
=
(\alpha_3\alpha_2^{-a_0})_*\eta_a
\quad\text{for almost every }a.
\] It does \textbf{not} give \[
(\alpha_2)_*\bar\eta=\bar\eta
\quad\text{or}\quad
(\alpha_3)_*\bar\eta=\bar\eta,
\] where \[
\bar\eta=\int\eta_a\,d\nu(a).
\]

Even stationarity for a random walk does not follow, because the fiber measure \(\eta_a\) may be correlated arbitrarily with the symbol \(a_0\).

There is a decisive counterexample to any abstract entropy-transfer claim. Let \(\nu\) be any shift-invariant positive-entropy measure and let \(0\in X_6\) be the identity. Then \[
\eta=\nu\times\delta_0
\] is \(F\)-invariant, since every \(\alpha_3\alpha_2^{-k}\) fixes \(0\). The skew system has positive entropy, but all of it lies in the symbolic base; the algebraic fiber is atomic and has zero entropy.

Therefore:

\[
\boxed{\text{Positive entropy of the valuation process does not imply positive entropy for }\alpha_2\text{ or }\alpha_3.}
\]

This is not a minor missing estimate. It is a rank-production problem.


\bigskip



\subsubsection*{7. Thickening the orbit produces rank two\textemdash{}but erases Collatz}


There is a natural way to obtain a genuinely \(\mathbb Z^2\)-invariant measure.

Define the lattice regions \[
\mathcal F_N
=
\{(-m,j):0\leq j<N,\ 0\leq m\leq A_j\}
\subset\mathbb Z^2.
\] Because \(a_j\geq1\), we have \(A_j\geq j\), and hence \[
|\mathcal F_N|
=
\sum_{j<N}(A_j+1)
\geq\frac{N(N+1)}2.
\] The boundary under horizontal translation has size \(O(N)\). The boundary under vertical translation is controlled by \[
\sum_{j<N}|A_{j+1}-A_j|
=
A_N
=
O(N).
\] Thus \((\mathcal F_N)\) is a F\o{}lner sequence in \(\mathbb Z^2\).

Consequently, weak limits of \[
\lambda_N
=
\frac1{|\mathcal F_N|}
\sum_{(-m,j)\in\mathcal F_N}
\delta_{\alpha_2^{-m}\alpha_3^jz_0}
\] are genuinely invariant under both \(\alpha_2\) and \(\alpha_3\).

But the actual Collatz points occur only on the graph \[
m=A_j.
\] There are \(N\) such boundary points, while \[
|\mathcal F_N|\asymp N^2.
\] Hence their relative mass is \[
\frac{N}{|\mathcal F_N|}\longrightarrow0.
\]

The exact carry relation is therefore lost in the rank-two limit. The two desiderata are incompatible under this averaging:

\begin{itemize}
\item average along the graph: retain Collatz, obtain only one skew action;
\item average over a F\o{}lner region: obtain rank two, lose Collatz on a zero-density boundary.
\end{itemize}

A fixed-width strip retains the boundary with positive density but is not F\o{}lner. A width tending to infinity becomes approximately invariant, but the boundary proportion tends to zero. Weighting the strip differently gives the same probability-versus-boundary tradeoff.

This is the geometric heart of the obstruction.


\bigskip



\subsubsection*{8. Compactification also erases the positive lift}


Equation (5) shows that in the solenoid \[
z_j=
\left[
\left(-x_0\frac{3^j}{2^{A_j}},0,0\right)
\right].
\] The affine \(+1\) term has become a diagonal lattice translation and disappeared from the quotient.

Keeping \(x_j\in\mathbb Z_3\) directly preserves the carry identities, but then the update is symbol-dependent and affine: \[
x_{j+1}=\frac{3x_j+1}{2^{a_j}}.
\] Passing to the solenoid makes the update algebraic and commuting, but forgets which representative is the actual positive integer.

This cannot be repaired merely by adding compact residue coordinates. Positive integers are dense in \[
\mathbb Z_2\times\mathbb Z_3\times\mathbb T
\] under \[
n\longmapsto
\bigl(n,n,\{\log_2n\}\bigr).
\] Indeed, fix a residue class modulo \(M=2^r3^s\) and a nonempty logarithmic-phase interval \(I\). At sufficiently large scales, the multiplicative interval corresponding to \(I\) has ordinary length greater than \(M\), so it contains an integer in the prescribed residue class.

Thus ``is the same ordinary positive integer at all places'' is not a closed condition in the compact state space.

If absolute height is retained in a noncompact state space, then \(x_j\to\infty\), so every compact set receives asymptotically zero empirical mass. This is the original mass escape.

Tao's probabilistic treatment similarly uses families of approximately transported measures rather than extracting a single useful invariant Collatz probability; this reflects the same basic noncompactness issue. [\hypertarget{cit:R7:3}{}\hyperlink{ref:R7:3}{3}]


\bigskip



\subsubsection*{9. The positive-entropy branch would close conditionally}


The positive-entropy idea itself is sound once one assumes a rank-production theorem.


\paragraph*{Conditional rank-production hypothesis}


Suppose an ergodic orbit-derived component \(\nu\) with \[
m=\int a_0\,d\nu\leq\log_2 3
\] and \[
h_\nu(\sigma)>0
\] produces a jointly \((\times2,\times3)\)-invariant ergodic circle measure whose \(\times2\)-coding is the binary suspension \[
a_j\longmapsto 1\,0^{a_j-1}.
\]

The suspension has:

\[
\text{frequency of }1=\frac1m
\geq\frac1{\log_2 3}
> \frac12,
\] and \[
h_{\times2}=\frac{h_\nu(\sigma)}m>0.
\]

Rudolph\textendash{}Johnson would force this measure to be Lebesgue. But Lebesgue measure under \(\times2\) has independent fair binary digits, so its \(1\)-frequency is exactly \(1/2\). Equivalently, its return-time law to the odd cylinder is \[
\Pr(a=k)=2^{-k},
\qquad
\mathbb E a=2.
\]

Both conclusions contradict \[
m\leq\log_2 3<2.
\]

Therefore:

\[
\boxed{\text{A positive-entropy low-mean component is impossible if rank production and coding transfer hold.}}
\]

Once the missing rank-production statement is supplied, this branch is essentially complete.


\bigskip



\subsubsection*{10. The zero-entropy branch does not imply periodicity}


For the skew system itself, zero entropy certainly does not force periodicity.

Choose an aperiodic zero-entropy shift measure \(\nu\) on symbols \(1,2\), for example a Sturmian process with the frequency of \(2\) chosen below \[
\log_2 3-1.
\] Then \[
\int a_0\,d\nu\leq\log_2 3,
\] and \[
\nu\times\delta_0
\] is an aperiodic, zero-entropy, \(F\)-invariant probability.

Moreover, the exact \(2\)-adic Collatz map is topologically and measurably conjugate to the binary shift through its parity coding. Thus arbitrary aperiodic zero-entropy parity patterns are perfectly legitimate \(2\)-adic Collatz patterns, even though they need not come from ordinary positive integers. [\hypertarget{cit:R7:4}{}\hyperlink{ref:R7:4}{4}, \hypertarget{cit:R7:5}{}\hyperlink{ref:R7:5}{5}]

So local binary admissibility and zero entropy do not imply eventual periodicity.

At the level of genuine \(\times2,\times3\)-invariant measures, the general zero-entropy classification is also the unresolved part of Furstenberg's measure problem: positive entropy is rigid, while a complete classification of zero-entropy invariant measures is not available. [\hyperlink{ref:R7:1}{1}]

Therefore the proposed zero-entropy step would require a new \textbf{Collatz-specific zero-entropy rigidity theorem}, not merely an application of existing Furstenberg theory.


\bigskip



\subsubsection*{11. Exact height-sensitive recurrence that would suffice}


There is nevertheless a clean way that symbolic recurrence could become an ordinary-integer repetition.

Suppose times \(p<q\) have matching centered valuation blocks of radius \(r\).

Let \[
A^+_{p,r}
=
\sum_{i=p}^{p+r-1}a_i.
\]

The matching future blocks imply \[
x_p\equiv x_q\pmod{2^{A^+_{p,r}}}.
\tag{9}
\] This follows from \[
2^{A^+_{p,r}}x_{p+r}
=
3^rx_p+C_w,
\] because \(3\) is invertible modulo \(2^{A^+_{p,r}}\).

The matching past blocks imply \[
x_p\equiv x_q\pmod{3^r}.
\tag{10}
\] Indeed, for the preceding word \(v\), \[
2^{A(v)}x_p
=
3^rx_{p-r}+C_v,
\] so \[
x_p\equiv 2^{-A(v)}C_v\pmod{3^r}.
\] This is the full residue version underlying terminal-suffix locality.

By the Chinese remainder theorem, \[
\boxed{\displaystyle
x_p\equiv x_q
\pmod{2^{A^+_{p,r}}3^r}.}
\tag{11}
\]

Consequently, if one could also prove \[
\max(x_p,x_q)
<
2^{A^+_{p,r}}3^r,
\tag{12}
\] then (11) would force \[
x_p=x_q,
\] contradicting divergence.

This isolates the missing zero-entropy statement:


\begin{quote}
\textbf{Height-sensitive recurrence:} a zero-entropy orbit-derived process must contain centered-block recurrences whose associated two-ended modulus exceeds the actual heights at both occurrences.

\end{quote}


Ordinary Poincar\'e recurrence does not imply (12). Neither does zero entropy. Zero entropy controls the number or return times of symbolic blocks only subexponentially, whereas the height at the return time can still be exponentially large in that return time.

Thus the exact two-ended congruence is available; the required comparison with escaped height is not.


\bigskip



\subsubsection*{12. The cycle branch}


Even after divergent trajectories are excluded, a full Collatz proof must rule out nontrivial positive cycles.

Suppose a period has valuation word \[
b_0,\ldots,b_{L-1},
\] and put \[
B=\sum_{i=0}^{L-1}b_i,
\qquad
B_k=\sum_{i=0}^{k-1}b_i.
\] Iterating around the period yields \[
2^Bx
=
3^Lx+
\sum_{k=0}^{L-1}3^{L-1-k}2^{B_k}.
\] Hence every positive cycle must satisfy \[
\boxed{\displaystyle
x=
\frac{
\sum_{k=0}^{L-1}3^{L-1-k}2^{B_k}
}{
2^B-3^L
}.}
\tag{13}
\]

Positivity requires \[
2^B>3^L.
\]

For \[
L=1,\qquad b_0=2,
\] formula (13) gives \(x=1\), the trivial accelerated cycle.

But neither Furstenberg rigidity nor the entropy argument proves that (13) has no other positive odd integral solutions. An atomic outcome of a rigidity theorem is not automatically forbidden: multiplication by \(2\) and \(3\) has many finite rational orbits. One must still show that no such algebraic orbit has a faithful lift satisfying the Collatz carry equation and positivity.

So the phrase ``reducing to a forbidden cycle equation'' hides an additional theorem: the equation has been reduced to, but not forbidden by, the measure argument.


\bigskip



\subsection*{Exact conditional completion theorem}


The Furstenberg program would prove Collatz if the following three statements were established.


\paragraph*{Rank production}


Every positive-entropy, orbit-derived, low-mean skew component admits a genuine jointly \((\times2,\times3)\)-invariant algebraic factor carrying the binary suspension entropy and digit frequencies.

This branch would then be contradicted by Rudolph\textendash{}Johnson and (2).


\paragraph*{Zero-entropy nonescape}


Every zero-entropy, orbit-derived, low-mean component has centered-block returns satisfying the height inequality (12).

This branch would then be contradicted by the exact congruence (11).


\paragraph*{Positive cycle exclusion}


Equation (13) has no positive odd integral solution except \[
L=1,\quad b_0=2,\quad x=1.
\]

Given these three statements, a counterexample can be neither divergent nor periodic, and the Collatz conjecture follows.


\bigskip



\subsection*{Final status}


The approach successfully yields:

\begin{itemize}
\item a compact orbit-derived invariant probability;
\item exact skew dynamics on an adelic solenoid;
\item the drift bound \(\int a_0\leq\log_2 3\);
\item a strict entropy ceiling;
\item equidistribution of logarithmic height modulo powers of \(2\);
\item a complete Haar contradiction conditional on rank production;
\item an exact two-ended recurrence criterion conditional on nonescape.
\end{itemize}

It does \textbf{not} presently yield:

\begin{itemize}
\item separate \(\times2\)- and \(\times3\)-invariance while retaining the carry boundary;
\item transfer of symbolic entropy to an algebraic generator;
\item zero-entropy periodicity or height-sensitive recurrence;
\item exclusion of all nontrivial positive cycle equations.
\end{itemize}

The central new target is therefore not ordinary recurrence. It is a \textbf{boundary-rigidity theorem} that converts the variable-time skew orbit into higher-rank algebraic invariance without averaging away the positive-integer lift. Existing Furstenberg rigidity begins only after that conversion; it does not perform it.
\subsection*{References}
\begin{enumerate}[leftmargin=2.2em]
\item \hypertarget{ref:R7:1}{}\hyperlink{cit:R7:1}{\texttt{turn578358view3}}
\item \hypertarget{ref:R7:2}{}\hyperlink{cit:R7:2}{\texttt{turn475872view0}}
\item \hypertarget{ref:R7:3}{}\hyperlink{cit:R7:3}{\texttt{turn578358view0}}
\item \hypertarget{ref:R7:4}{}\hyperlink{cit:R7:4}{\texttt{turn475872view2}}
\item \hypertarget{ref:R7:5}{}\hyperlink{cit:R7:5}{\texttt{turn475872view3}}
\end{enumerate}

\clearpage
\section{Round 8: approach 3 --- Turn the q-analogue into a polynomialization theorem}
\subsection*{Prompt}
\begin{quote}
Complete a full proof of the Collatz conjecture using the following approach:

\begin{enumerate}
\item Turn the q-analogue into a polynomialization theorem
\end{enumerate}

\textbf{Established/Hybrid \textperiodcentered{} P}

Under the known conjugacy Z\_2 <-> F\_2[[q]], positive integers become a special family of formal power series. Prove that every series in this family eventually becomes a polynomial, or enters the polynomial basin, under a carefully chosen conjugate q-map. Polynomial inputs already have terminating behavior in the natural analogue.

The hoped-for invariant could be support entropy, gap growth, a valuation filtration, or a noetherian ideal attached to the coefficient sequence. This would be an astonishing proof: ordinary integer termination would follow from eventual finite support in a formal series.

Take as long as you need.
\end{quote}
\subsection*{Response}

\subsection*{Rigorous outcome}


The formal-series program can be pushed to a very clean endpoint, but that endpoint exposes a decisive obstruction:

\[
\boxed{\text{For the natural conjugate \(q\)-map, a nonpolynomial series can never enter the polynomial set.}}
\]

Moreover, the statement that the series attached to every positive integer is already a polynomial is \textbf{exactly equivalent to the Collatz conjecture}, coefficient by coefficient. It is also equivalent to the older Bernstein\textendash{}Lagarias reformulation using the \(2\)-adic parity conjugacy.

Thus this approach does not presently supply a full proof. A purported final argument based only on eventual entry into the polynomial basin would be circular. What follows is the complete reduction, including an explicit formula for the conjugacy, a classification of the natural degree-decreasing maps, and precise obstructions to the suggested invariants.

Monks introduced the relevant family \(T_{A,B}\) and singled out \(T_{1,1+q^2}\), proving that all polynomial inputs terminate in its two-cycle or fixed points. He then stated polynomiality of the images of positive integers as Conjecture 6, explicitly noting its equivalence to Collatz. [\hypertarget{cit:R8:1}{}\hyperlink{ref:R8:1}{1}, \hypertarget{cit:R8:2}{}\hyperlink{ref:R8:2}{2}, \hypertarget{cit:R8:3}{}\hyperlink{ref:R8:3}{3}]


\bigskip



\subsubsection*{1. The exact formal-series conjugacy}


Use the shortcut Collatz map

\[
T(x)=
\begin{cases}
x/2,&x\equiv0\pmod2,\\[2mm]
(3x+1)/2,&x\equiv1\pmod2,
\end{cases}
\qquad x\in\mathbb Z_2.
\]

For \(x\in\mathbb Z_2\), write

\[
\varepsilon_k(x)=T^k(x)\pmod2\in\mathbb F_2
\]

and form its parity series

\[
\mathcal Q(x)=\sum_{k\geq0}\varepsilon_k(x)q^k
       \in\mathbb F_2[[q]].
\]

Let

\[
\sigma(g)=\frac{g+g(0)}q
\]

be the left shift on formal power series. Then

\[
\mathcal Q(Tx)=\sigma(\mathcal Q(x)).
\]

The parity sequence determines a unique \(2\)-adic integer, so \(\mathcal Q\) is a bijective isometry. This is the formal-series version of the Bernstein\textendash{}Lagarias parity-vector conjugacy. [\hypertarget{cit:R8:4}{}\hyperlink{ref:R8:4}{4}, \hypertarget{cit:R8:5}{}\hyperlink{ref:R8:5}{5}]

Now put

\[
B(q)=1+q^2
\]

and define

\[
R(f)=
\begin{cases}
f/q,&f(0)=0,\\[2mm]
(f+B)/q,&f(0)=1.
\end{cases}
\tag{1}
\]

This is Monks's \(T_{1,1+q^2}\).


\paragraph*{Lemma 1}


For every \(g\in\mathbb F_2[[q]]\),

\[
R(Bg)=B\sigma(g).
\tag{2}
\]


\paragraph*{Proof}


Since \(B(0)=1\), the constant coefficient of \(Bg\) is \(g(0)\). Therefore

\[
R(Bg)
 =\frac{Bg+g(0)B}{q}
 =B\frac{g+g(0)}q
 =B\sigma(g).
\qquad\square
\]

Consequently the map

\[
\boxed{H(x)=(1+q^2)\mathcal Q(x)}
\tag{3}
\]

is an explicit conjugacy:

\[
H(Tx)=R(H(x)).
\tag{4}
\]

So there is no need to leave the relevant conjugacy implicit.


\bigskip



\subsubsection*{2. Polynomial dynamics is completely understood}



\paragraph*{Proposition 2}


For \(f\in\mathbb F_2[[q]]\),

\[
R(f)\in\mathbb F_2[q]
\quad\Longleftrightarrow\quad
f\in\mathbb F_2[q].
\tag{5}
\]

Moreover, every polynomial orbit reaches

\[
0,\qquad 1+q,\qquad\text{or}\qquad\{1,q\}.
\]


\paragraph*{Proof}


From (1),

\[
f=qR(f)+f(0)(1+q^2).
\tag{6}
\]

Hence \(R(f)\) polynomial implies \(f\) polynomial. The converse follows immediately from the definition of \(R\).

Now suppose \(f\) is a polynomial of degree \(d>1\).

If \(f(0)=0\), then

\[
R(f)=f/q,
\]

which has degree \(d-1\).

If \(f(0)=1\), write

\[
f=P+q^d,
\qquad \deg P<d,
\qquad P(0)=1.
\]

Then

\[
R(f)
 =\frac{P+q^d+1+q^2}{q}
 =\frac{P+1}{q}+q+q^{d-1},
\]

whose degree is strictly less than \(d\), possibly after cancellation.

Thus the degree decreases until it is at most one. Direct calculation gives

\[
R(0)=0,\qquad R(1+q)=1+q,\qquad R(1)=q,\qquad R(q)=1.
\]

This proves the result. \(\square\)

In particular,

\[
\boxed{\bigcup_{j\geq0}R^{-j}\bigl(\mathbb F_2[q]\bigr)
      =\mathbb F_2[q].}
\tag{7}
\]

The ``polynomial basin'' contains no nonpolynomial series.

There is also a useful conceptual explanation. Since

\[
R(f)=\sigma(f)+f(0)q,
\tag{8}
\]

the action of \(R\) on the tail is just a shift, with only one low coefficient modified. In the additive quotient

\[
\mathbb F_2[[q]]/\mathbb F_2[q],
\]

\(R\) induces the shift, and the shift is invertible on this quotient. Therefore a nonzero tail class cannot be mapped to zero.


\bigskip



\subsubsection*{3. Polynomiality is exactly Collatz termination}


Write

\[
H(x)=\sum_{k\geq0}c_k(x)q^k.
\]

From (3),

\[
c_k(x)=\varepsilon_k(x)+\varepsilon_{k-2}(x),
\tag{9}
\]

where \(\varepsilon_{-1}=\varepsilon_{-2}=0\).

Thus

\[
H(x)\in\mathbb F_2[q]
\]

if and only if

\[
\varepsilon_k(x)=\varepsilon_{k-2}(x)
\qquad\text{for all sufficiently large }k.
\tag{10}
\]

In words, \(H(x)\) is polynomial exactly when the Collatz parity sequence is eventually periodic with period dividing two.

There are only four possible period-two tails:

\[
00\,00\,00\cdots,\qquad
11\,11\,11\cdots,\qquad
10\,10\,10\cdots,\qquad
01\,01\,01\cdots.
\]

By injectivity of the parity-vector map, these are respectively the parity sequences of

\[
0,\qquad -1,\qquad 1,\qquad 2.
\]

The first two points are the \(2\)-adic fixed points \(0,-1\); the last two form the positive cycle \(1\leftrightarrow2\).

Since \(T\) preserves positive integers, a positive orbit cannot reach \(0\) or \(-1\). We therefore obtain the exact theorem.


\paragraph*{Theorem 3: polynomialization equivalence}


For every \(n\in\mathbb N_+\), the following are equivalent:

\[
\begin{aligned}
&\text{\(T^j(n)=1\) for some \(j\);}\\
&\text{\(H(n)\) is a polynomial;}\\
&\text{\(R^j(H(n))\) is a polynomial for some \(j\);}\\
&\text{the parity sequence of \(n\) is eventually \(10\)- or \(01\)-periodic.}
\end{aligned}
\tag{11}
\]


\paragraph*{Proof}


The first, second and fourth conditions are equivalent by (9)\textendash{}(10) and the classification of the four period-two tails. The second and third conditions are equivalent by complete invariance (5). \(\square\)

Consequently,

\[
\boxed{
\text{Collatz}
\quad\Longleftrightarrow\quad
H(\mathbb N_+)\subseteq\mathbb F_2[q].
}
\tag{12}
\]

This is precisely Monks's Conjecture 6.


\bigskip



\subsubsection*{4. Polynomial degree is exactly total stopping time}


The reformulation is even tighter than merely detecting termination.

Let

\[
\tau(n)=\min\{k\geq0:T^k(n)=1\}
\]

when this quantity exists.


\paragraph*{Proposition 4}


For every terminating positive integer \(n\),

\[
\boxed{\deg H(n)=\tau(n).}
\tag{13}
\]


\paragraph*{Proof}


Once \(T^{\tau(n)}(n)=1\), the parity tail alternates \(1,0,1,0,\ldots\). Hence (9) gives

\[
c_k(n)=0\qquad(k>\tau(n)).
\]

For \(\tau(n)\geq2\), minimality gives

\[
T^{\tau(n)-1}(n)=2,\qquad
T^{\tau(n)-2}(n)=4.
\]

Therefore

\[
\varepsilon_{\tau(n)}=1,
\qquad
\varepsilon_{\tau(n)-2}=0,
\]

and hence

\[
c_{\tau(n)}=1.
\]

The cases \(n=1,2\) are immediate:

\[
H(1)=1,\qquad H(2)=q.
\]

Thus the degree is exactly \(\tau(n)\). \(\square\)

For example,

\[
3\longmapsto5\longmapsto8\longmapsto4
 \longmapsto2\longmapsto1,
\]

so \(\tau(3)=5\). Its parity series is

\[
\mathcal Q(3)
 =1+q+\frac{q^5}{1+q^2},
\]

and therefore

\[
H(3)
 =(1+q^2)\mathcal Q(3)
 =1+q+q^2+q^3+q^5,
\]

which indeed has degree \(5\).

Thus ``prove finite support'' is not merely analogous to proving finite stopping time. Its support endpoint \textbf{is} the stopping time.


\bigskip



\subsubsection*{5. Relation to the Bernstein\textendash{}Lagarias formulation}


Interpret the parity coefficients as binary digits and write

\[
Q(x)=\sum_{k\geq0}\varepsilon_k(x)2^k\in\mathbb Z_2.
\]

A \(2\)-adic binary expansion is eventually period two precisely when it is a rational number with denominator dividing

\[
2^2-1=3.
\]

Hence

\[
H(n)\in\mathbb F_2[q]
\quad\Longleftrightarrow\quad
Q(n)\in\frac13\mathbb Z.
\tag{14}
\]

Bernstein and Lagarias formulated Collatz as

\[
Q(\mathbb N_+)\subseteq\frac13\mathbb Z,
\]

equivalently

\[
\mathbb N_+\subseteq\Phi\!\left(\frac13\mathbb Z\right),
\]

where \(\Phi=Q^{-1}\). [\hypertarget{cit:R8:6}{}\hyperlink{ref:R8:6}{6}]

Therefore the proposed polynomialization theorem is exactly the Bernstein\textendash{}Lagarias rationality formulation written in carryless \(\mathbb F_2[[q]]\) coordinates:

\[
\boxed{
\text{eventual period two}
\;\longleftrightarrow\;
\text{denominator \(3\)}
\;\longleftrightarrow\;
(1+q^2)\mathcal Q(n)\text{ polynomial}.
}
\]

This does not make the reformulation useless, but it identifies precisely how much new content is still required.


\bigskip



\subsubsection*{6. A different conjugate \(q\)-map cannot make the theorem logically weaker}


There is a general obstruction independent of the particular formula.


\paragraph*{Basin lemma}


Let \(h\) conjugate \(T\) to a map \(F\), and let \(\mathcal P\) be a class of ``polynomial'' points such that:

\begin{enumerate}
\item \(h(\{1,2\})\subseteq\mathcal P\);
\item every point of \(\mathcal P\) eventually reaches \(h(\{1,2\})\).
\end{enumerate}

Then

\[
F^j(h(n))\in\mathcal P
\text{ for some }j
\quad\Longleftrightarrow\quad
T^k(n)\in\{1,2\}
\text{ for some }k.
\tag{15}
\]

Indeed, entry into \(\mathcal P\) implies subsequent termination, while reaching the target cycle implies entry into \(\mathcal P\).

Thus whenever polynomial inputs genuinely have terminating behavior, proving that all positive-integer images enter their basin is automatically equivalent to Collatz. A different conjugacy might make the arithmetic easier, but it cannot make the required statement logically weaker.


\paragraph*{The natural \(T_{A,B}\) family}


Consider

\[
T_{A,B}(f)=
\begin{cases}
f/q,&f(0)=0,\\[2mm]
(Af+B)/q,&f(0)=1,
\end{cases}
\]

with odd polynomials \(A,B\).

For an odd polynomial \(f\) of sufficiently large degree \(d\),

\[
\deg T_{A,B}(f)=d+\deg A-1.
\]

Therefore a universal degree-descent mechanism forces

\[
\deg A=0,\qquad A=1.
\]

For \(A=1\),

\[
f=qT_{1,B}(f)+f(0)B,
\]

so polynomiality is completely invariant for every polynomial \(B\).

Moreover,

\[
H_B(x)=B\mathcal Q(x)
\]

is the corresponding parity-preserving conjugacy because

\[
T_{1,B}(Bg)=B\sigma(g).
\]

Since

\[
\mathcal Q(1)=\frac1{1+q^2},
\qquad
\mathcal Q(2)=\frac q{1+q^2},
\]

the target cycle maps to polynomials only when \(1+q^2\) divides \(B\). The simplest choice is exactly

\[
B=1+q^2.
\]

Hence, throughout the natural part of Monks's family where:

\begin{itemize}
\item polynomial degree gives termination, and
\item the positive target cycle is polynomial,
\end{itemize}

the polynomial set necessarily has no external basin.


\bigskip



\subsubsection*{7. Why the proposed invariants do not presently close the argument}



\paragraph*{7.1 Support entropy is a tail invariant}


Equation (8) says

\[
R(f)=\sigma(f)+f(0)q.
\]

Thus \(R\) only shifts the infinite tail and changes one low coefficient. Consequently all quantities unchanged by shifts and finite modifications are invariant along an \(R\)-orbit:

\[
\text{support density},\quad
\text{block entropy},\quad
\text{asymptotic gap statistics},\quad
\text{eventual rationality}.
\]

They cannot serve as strictly decreasing Lyapunov functions.

There is also an information-theoretic reason. If

\[
c_k=\varepsilon_k+\varepsilon_{k-2},
\]

then a block of \(c\)-coefficients is determined by a parity block only two symbols longer. Conversely, a parity block is determined by its \(c\)-block and two initial bits. Thus their block complexities satisfy, up to harmless shifts,

\[
p_c(m)\leq p_\varepsilon(m+2),
\qquad
p_\varepsilon(m)\leq4p_c(m-2).
\]

Their entropy rates are therefore equal. In particular, replacing the parity sequence by its support-defect sequence does not compress its asymptotic information.

Zero entropy would also be insufficient: many nonperiodic sequences have zero entropy. What is needed is eventual period two, not merely low complexity.


\bigskip



\paragraph*{7.2 A precise gap lemma}


Although gap growth does not prove termination, it gives a useful arithmetic restriction.

Suppose that at time \(k\), a parity pair repeats for \(L\) consecutive two-step blocks. Put

\[
x=T^k(n).
\]

The four possibilities give

\[
\begin{array}{c|c|c}
\text{repeated pair}
&
T^{2L}(x)
&
\text{necessary congruence}
\\ \hline
00
&
\dfrac{x}{4^L}
&
x\equiv0\pmod{4^L}
\\[3mm]
10
&
1+\left(\dfrac34\right)^L(x-1)
&
x\equiv1\pmod{4^L}
\\[3mm]
01
&
2+\left(\dfrac34\right)^L(x-2)
&
x\equiv2\pmod{4^L}
\\[3mm]
11
&
-1+\left(\dfrac94\right)^L(x+1)
&
x\equiv-1\pmod{4^L}.
\end{array}
\tag{16}
\]

For example, the \(10\) two-step map is

\[
x\longmapsto\frac{3x+1}{4}
            =1+\frac34(x-1),
\]

and iteration gives the displayed formula. Integrality forces the congruence because \(3\) is invertible modulo \(4^L\).

Consequently, unless \(x\) is already \(1\) or \(2\), a defect-free block of length \(2L\) forces approximately

\[
x\geq4^L.
\tag{17}
\]

So long support gaps force very large orbit values. They do \textbf{not} force termination: a hypothetical divergent orbit could compensate for longer gaps with larger heights. A gap argument needs an independent Archimedean upper bound that beats \(4^L\).


\bigskip



\paragraph*{7.3 The \(q\)-valuation is not monotone}


For even \(f\),

\[
v_q(Rf)=v_q(f)-1.
\]

For odd \(f\), however,

\[
v_q(Rf)=v_q(f+1+q^2)-1,
\]

and the cancellation can be arbitrarily deep. For instance, with \(N\geq3\),

\[
f=1+q^2+q^N
\]

is odd, but

\[
R(f)=q^{N-1}.
\]

Thus \(v_q\) jumps from \(0\) to \(N-1\). It cannot define a uniformly descending well-founded filtration.


\bigskip



\paragraph*{7.4 Noetherianity does not produce finite support}


The ring \(\mathbb F_2[[q]]\) is a Noetherian discrete valuation ring, but this alone gives no polynomialization.

First, \(\mathbb F_2[q]\) is not an ideal in \(\mathbb F_2[[q]]\). Second, the descending filtration

\[
(q)\supset(q^2)\supset(q^3)\supset\cdots
\]

need not stabilize; Noetherianity controls ascending chains, not descending ones. Third, every odd series is a unit, so the principal ideal generated by an odd orbit point is the entire ring.

A more meaningful construction is the polynomial-annihilator ideal

\[
I(g)=
\left\{
P(q)\in\mathbb F_2[q]:
P(q)g(q)\in\mathbb F_2[q]
\right\}.
\tag{18}
\]

Then

\[
I(g)\neq(0)
\quad\Longleftrightarrow\quad
g\in\mathbb F_2(q),
\]

because \(Pg=A\) gives \(g=A/P\). Equivalently, the coefficient sequence of \(g\) is eventually periodic.

But Collatz requires the much more specific conclusion

\[
1+q^2\in I(\mathcal Q(n)).
\tag{19}
\]

Noetherianity says that \(I(g)\), once known to be nonzero, has a finite generator. It does not prove that \(I(g)\neq0\), and it certainly does not force its generator to divide \(1+q^2\).

Even proving \(I(\mathcal Q(n))\neq0\) would only show that every integer orbit is eventually periodic. One would still have to exclude all nontrivial positive cycles.


\bigskip



\paragraph*{7.5 No finite-state Lyapunov invariant can work}


Positive integers are dense in \(\mathbb Z_2\), and \(H\) is an isometry. Hence

\[
H(\mathbb N_+)
\]

is dense in \(\mathbb F_2[[q]]\). More concretely, for every \(r\) and every polynomial \(P\) of degree less than \(r\), there are infinitely many positive integers \(n\) such that

\[
H(n)\equiv P\pmod{q^r}.
\tag{20}
\]

Thus every finite coefficient pattern occurs among positive-integer images.

This gives a useful no-go theorem. Suppose

\[
V:\mathbb F_2[[q]]\longrightarrow\mathbb N
\]

were continuous and satisfied

\[
V(RH(n))<V(H(n))
\qquad(n\neq1,2).
\tag{21}
\]

The set

\[
\{f:V(Rf)<V(f)\}
\]

would be clopen and would contain the dense set \(H(\mathbb N_+\setminus\{1,2\})\). It would therefore be the whole space. But then on the cycle \(1\leftrightarrow q\) one would have

\[
V(q)<V(1)<V(q),
\]

a contradiction.

Hence no rank depending continuously on finitely many coefficients\textemdash{}equivalently, no fixed finite-state or finite-jet invariant\textemdash{}can strictly descend on every positive orbit.

Any successful invariant must use genuinely nonlocal information or additional Archimedean arithmetic.


\bigskip



\subsubsection*{8. The exact missing theorem}


The inverse parity map has an explicit form. Suppose

\[
g(q)=\sum_{j\geq0}q^{d_j},
\qquad
0\leq d_0<d_1<\cdots.
\]

Then the corresponding \(2\)-adic integer is

\[
\Phi(g)
 =
-\sum_{j\geq0}\frac{2^{d_j}}{3^{j+1}}
 \in\mathbb Z_2.
\tag{22}
\]

This is the inverse parity-vector formula. [\hyperlink{ref:R8:4}{4}]

Since \(H=(1+q^2)\mathcal Q\), the positive-integer family can be written exactly as

\[
H(\mathbb N_+)
=
\left\{
(1+q^2)g:
-\sum_{j\geq0}\frac{2^{d_j}}{3^{j+1}}
\in\mathbb N_+
\right\}.
\tag{23}
\]

For each \(r\), let \(N_r(g)\) be the least nonnegative residue modulo \(2^r\) of the truncated \(2\)-adic sum

\[
-\sum_{d_j<r}\frac{2^{d_j}}{3^{j+1}}.
\tag{24}
\]

Then \(\Phi(g)\) is an ordinary nonnegative integer precisely when \(N_r(g)\) eventually stabilizes as an ordinary integer.

The desired polynomialization theorem is therefore:

\[
\boxed{
N_r(g)\text{ eventually stabilizes at a positive integer}
\quad\Longrightarrow\quad
(1+q^2)g\in\mathbb F_2[q].
}
\tag{25}
\]

Equivalently:


\begin{quote}
If the inverse parity sum has only finitely many nonzero ordinary binary digits, then the parity sequence has only finitely many departures from period two.

\end{quote}


That is the exact missing carry-noncancellation theorem. It is not supplied by the formal-series dynamics because arithmetic in \(\mathbb F_2[[q]]\) deliberately removes carries, while the condition that \(\Phi(g)\) be an ordinary positive integer is entirely about the eventual cancellation of \(2\)-adic carries in (22).


\bigskip



\subsection*{Verdict}


The approach produces an elegant exact encoding:

\[
\boxed{
H(n)=(1+q^2)
\sum_{k\geq0}(T^k(n)\bmod2)q^k.
}
\]

It also proves:

\[
\boxed{
\deg H(n)=\text{total stopping time of }n
}
\]

whenever the orbit terminates.

But the hoped-for enlargement of the polynomial basin does not exist:

\[
R^j(f)\text{ polynomial for some }j
\quad\Longleftrightarrow\quad
f\text{ polynomial}.
\]

And the global polynomialization statement is exactly

\[
\Phi(g)\in\mathbb N_+
\quad\Longrightarrow\quad
(1+q^2)g\in\mathbb F_2[q],
\]

which is Collatz itself, equivalently the Bernstein\textendash{}Lagarias denominator-\(3\) formulation.

The formal-series coordinate system is therefore a useful lens, especially through the defect coefficients

\[
c_k=\varepsilon_k+\varepsilon_{k-2},
\]

but support entropy, asymptotic gaps, \(q\)-valuation, finite-state ranks, and bare Noetherianity do not supply the missing implication. A successful version of this program needs an additional nonlocal theorem coupling the support defects to the Archimedean height or to the \(3\)-adic carry structure in the inverse sum (22).
\subsection*{References}
\begin{enumerate}[leftmargin=2.2em]
\item \hypertarget{ref:R8:1}{}\hyperlink{cit:R8:1}{\texttt{turn988950view2}}
\item \hypertarget{ref:R8:2}{}\hyperlink{cit:R8:2}{\texttt{turn988950view3}}
\item \hypertarget{ref:R8:3}{}\hyperlink{cit:R8:3}{\texttt{turn988950view4}}
\item \hypertarget{ref:R8:4}{}\hyperlink{cit:R8:4}{\texttt{turn587594view0}}
\item \hypertarget{ref:R8:5}{}\hyperlink{cit:R8:5}{\texttt{turn587594view1}}
\item \hypertarget{ref:R8:6}{}\hyperlink{cit:R8:6}{\texttt{turn587594view2}}
\end{enumerate}

\clearpage
\section{Round 9: approach 4 --- A Cobham trap for counterexample carry streams}
\subsection*{Prompt}
\begin{quote}
Complete a full proof of the Collatz conjecture using the following approach:

\begin{enumerate}
\item A Cobham trap for counterexample carry streams
\end{enumerate}

\textbf{Hybrid \textperiodcentered{} P}

Base 2 determines the parity/seed side, while base 3 controls endpoint carries. Try to prove that the carry sequence of any infinite positive-integer counterexample is simultaneously describable by finite transducers in multiplicatively independent bases 2 and 3.

A Cobham-Semenov-type rigidity theorem would then force eventual periodicity. Eventual periodicity should lead to a rational 2-adic point and, after imposing positivity and the affine identity, to a cycle equation.

The whole Collatz transition relation is not base-2 automatic, so the target must be the highly restricted stream associated with a hypothetical counterexample, not the full relation.

Take as long as you need.
\end{quote}
\subsection*{Response}

\subsection*{Result of the Cobham program}


After carrying the approach through, the strongest valid conclusion is a \textbf{conditional cycle theorem}, not yet a full proof of Collatz.

There are two independent gaps:

\begin{enumerate}
\item The known base-\(2\) and base-\(3\) transducers act on the digits of the \textbf{current orbit value}. Cobham's theorem requires a finite automaton acting on the base-\(2\) or base-\(3\) digits of the \textbf{time index} \(j\). The first property does not imply the second.

\item Even if the required dual automaticity were established, Cobham would force the orbit to become periodic, but it would not exclude a nontrivial positive cycle. Indeed, the carry stream of any cycle is already periodic and therefore automatic in every base.
\end{enumerate}

Current 2026 research continues to treat Collatz as an open conjecture. [\hypertarget{cit:R9:1}{}\hyperlink{ref:R9:1}{1}, \hypertarget{cit:R9:2}{}\hyperlink{ref:R9:2}{2}]

What follows is the exact theorem that this program does prove, followed by rigorous obstructions to completing the two missing steps.


\bigskip



\subsubsection*{1. The accelerated orbit and its carry data}


Work with the accelerated odd Collatz map

\[
U(x)=\frac{3x+1}{2^{\nu _2(3x+1)}}
\]

on positive odd integers. Let

\[
x_{j+1}=U(x_j),\qquad
a_j=\nu _2(3x_j+1)\geq 1,
\]

and set

\[
A_0=0,\qquad A_m=\sum_{j=0}^{m-1}a_j.
\]

For a finite valuation word

\[
w=(a_0,\ldots ,a_{m-1}),
\]

define

\[
C_w=\sum_{i=0}^{m-1}3^{m-1-i}2^{A_i}.
\]

Repeated substitution in

\[
2^{a_j}x_{j+1}=3x_j+1
\]

gives the exact affine identity

\[
\boxed{\,2^{A_m}x_m=3^m x_0+C_w.\,} \tag{1}
\]

This identity contains both the \(2\)-adic seed information and the \(3\)-adic endpoint information.


\bigskip



\subsubsection*{2. What Cobham's theorem would actually need}


For a finite-alphabet sequence \(z=(z_j)_{j\geq0}\), its base-\(b\) kernel is

\[
\mathcal K_b(z)=
\left\{
(z_{b^e n+r})_{n\geq0}:
e\geq0,\ 0\leq r<b^e
\right\}.
\]

The sequence is \(b\)-automatic precisely when \(\mathcal K_b(z)\) is finite. Equivalently, a finite automaton reads the base-\(b\) representation of the integer \(j\) and outputs \(z_j\). Cobham's theorem says that a finite-valued sequence which is automatic in two multiplicatively independent bases is eventually periodic. Bell's regular-sequence extension applies to certain unbounded integer sequences. [\hypertarget{cit:R9:3}{}\hyperlink{ref:R9:3}{3}, \hypertarget{cit:R9:4}{}\hyperlink{ref:R9:4}{4}, \hypertarget{cit:R9:5}{}\hyperlink{ref:R9:5}{5}, \hypertarget{cit:R9:6}{}\hyperlink{ref:R9:6}{6}]

Because the valuation sequence \(a_j\) need not be bounded, its natural automatic object is its graph

\[
\Gamma_a=\{(j,a_j):j\geq0\}\subseteq\mathbb N^2.
\]

Call \(\Gamma_a\) \(b\)-recognizable when a synchronous finite automaton recognizes the padded base-\(b\) representations of the pairs \((j,a_j)\).

The multidimensional Cobham\textendash{}Semenov theorem states that a subset of \(\mathbb N^d\) recognizable in two multiplicatively independent bases is Presburger-definable, equivalently semilinear. [\hypertarget{cit:R9:7}{}\hyperlink{ref:R9:7}{7}, \hypertarget{cit:R9:8}{}\hyperlink{ref:R9:8}{8}]

This gives the central rigorous result.


\bigskip



\subsection*{Theorem 1: the valuation-graph Cobham trap}


Let \((x_j)\) be a positive orbit of the accelerated odd Collatz map, and let

\[
a_j=\nu _2(3x_j+1).
\]

Suppose that the graph

\[
\Gamma_a=\{(j,a_j):j\geq0\}
\]

is both \(2\)-recognizable and \(3\)-recognizable. Then \(a_j\) is eventually periodic, and the Collatz orbit \((x_j)\) eventually enters a cycle.


\subsubsection*{Proof}


First observe that

\[
2^{a_j}x_{j+1}=3x_j+1\leq4x_j,
\]

because \(x_j\geq1\). Multiplying this inequality for \(j=0,\ldots ,n-1\) and cancelling the intermediate orbit values gives

\[
2^{A_n}x_n\leq4^n x_0.
\]

Since \(x_n\geq1\),

\[
\boxed{\,A_n\leq 2n+\log _2x_0.\,} \tag{2}
\]

Thus the cumulative valuation sum has at most linear growth.

By Cobham\textendash{}Semenov, the simultaneous \(2\)- and \(3\)-recognizability of \(\Gamma_a\) implies that \(\Gamma_a\) is Presburger-definable. A function whose graph is Presburger-definable is piecewise linear on finitely many arithmetic progressions. [\hypertarget{cit:R9:9}{}\hyperlink{ref:R9:9}{9}, \hypertarget{cit:R9:10}{}\hyperlink{ref:R9:10}{10}]

Consequently, there exist \(M\geq1\), \(N\), and constants \(\alpha_r,\beta_r\) such that, for every \(0\leq r<M\),

\[
a_{Mt+r}=\alpha_r t+\beta_r \tag{3}
\]

whenever \(Mt+r\geq N\).

Because \(a_j\geq1\), no \(\alpha_r\) can be negative. Suppose that \(\alpha_r>0\) for some residue class \(r\). Then

\[
\sum_{t\leq T}a_{Mt+r}
=
\frac{\alpha_r}{2}T^2+O(T),
\]

so \(A_{MT}\gg T^2\). This contradicts the linear bound (2). Hence

\[
\alpha_r=0
\]

for every \(r\), and therefore

\[
a_{j+M}=a_j
\]

for all sufficiently large \(j\).

Now expand the accelerated valuations into the ordinary shortcut-map parity stream. One accelerated step with valuation \(a_j\) contributes the block

\[
1\,0^{a_j-1}.
\]

Thus the ordinary parity sequence has the form

\[
1\,0^{a_0-1}
1\,0^{a_1-1}
1\,0^{a_2-1}\cdots . \tag{4}
\]

Eventual periodicity of \((a_j)\) implies eventual periodicity of this binary parity stream.

The standard \(2\)-adic parity map \(Q\) is a bijection on \(\mathbb Z_2\) and conjugates the shortcut Collatz map to the binary shift. Hence two orbit points with the same infinite parity tail are equal as \(2\)-adic integers. If the parity tail has period \(L\) beginning at ordinary time \(t_0\), then

\[
Q(T^{t_0+L}x_0)
=
S^{t_0+L}Q(x_0)
=
S^{t_0}Q(x_0)
=
Q(T^{t_0}x_0),
\]

and injectivity of \(Q\) gives

\[
T^{t_0+L}x_0=T^{t_0}x_0.
\]

Thus the positive-integer orbit enters a cycle. \ensuremath{\blacksquare}

The parity-map bijection, its isometry property, and the conjugacy with the binary shift are standard facts about the \(2\)-adic Collatz extension. [\hypertarget{cit:R9:11}{}\hyperlink{ref:R9:11}{11}, \hypertarget{cit:R9:12}{}\hyperlink{ref:R9:12}{12}, \hypertarget{cit:R9:13}{}\hyperlink{ref:R9:13}{13}]


\bigskip



\subsubsection*{3. A regular-sequence version}


There is an alternative formulation that may be more natural for the unbounded sequence \(a_j\).


\subsection*{Theorem 2: the regular valuation trap}


If the integer sequence \((a_j)\) is simultaneously \(2\)-regular and \(3\)-regular, then it is eventually periodic and the orbit enters a cycle.


\subsubsection*{Proof}


Bell's generalized Cobham theorem implies that

\[
F(z)=\sum_{j\geq0}a_jz^j
\]

is a rational function whose poles are roots of unity. [\hyperlink{ref:R9:3}{3}, \hyperlink{ref:R9:6}{6}]

Therefore there is some \(M\geq1\) such that, on each residue class modulo \(M\),

\[
a_{Mt+r}=P_r(t) \tag{5}
\]

for all sufficiently large \(t\), where \(P_r\) is a polynomial.

If some \(P_r\) had degree \(d\geq1\), then eventual positivity would force its leading coefficient to be positive, and the partial sums of \(a_j\) along this residue class would grow like \(T^{d+1}\). This again contradicts (2). Hence every \(P_r\) is constant, so \(a_j\) is eventually periodic. Theorem 1 then gives a cycle. \ensuremath{\blacksquare}

Thus either of the following would eliminate divergent Collatz trajectories:

\[
\Gamma_a
\text{ simultaneously \(2\)- and \(3\)-recognizable},
\]

or

\[
(a_j)\text{ simultaneously \(2\)- and \(3\)-regular}.
\]

The problem is establishing either statement.


\bigskip



\subsection*{4. Arithmetic transducers are not time-index automata}


The crucial distinction is between two different kinds of finite-state object.

An \textbf{arithmetic transducer} reads the digits of \(x\) and computes the digits of \(f(x)\).

A \textbf{time-index automaton} reads the digits of \(j\) and outputs something about the \(j\)-th iterate \(f^j(x_0)\).

The first does not imply the second.

For example, the map

\[
F(x)=2x
\]

is computed by a finite transducer in every positional base. Nevertheless, starting from \(1\), its orbit is

\[
x_j=2^j.
\]

If the graph

\[
\{(j,2^j):j\geq0\}
\]

were recognizable in both bases \(2\) and \(3\), Cobham\textendash{}Semenov would make it Presburger-definable. But Presburger-definable functions are piecewise linear, whereas \(2^j\) is exponential. Therefore the orbit graph is not simultaneously recognizable.

So even an extremely simple finite-state transition can produce a time-indexed orbit with unbounded automaton complexity.

This distinction applies directly to the known binary/ternary Collatz transducers. St\'erin and Woods construct a two-dimensional system in which one finite-state transducer computes \(x/2\) in ternary columns while another computes \(3x+1\) in binary rows. But these are spatial digit transducers, not automata reading the iteration number. Their infinite \(2\)-adic/\(3\)-adic base conversion remains conjectural, and their cycle construction reformulates positive-cycle exclusion as a reachability problem rather than resolving it. [\hypertarget{cit:R9:14}{}\hyperlink{ref:R9:14}{14}, \hypertarget{cit:R9:15}{}\hyperlink{ref:R9:15}{15}, \hypertarget{cit:R9:16}{}\hyperlink{ref:R9:16}{16}, \hypertarget{cit:R9:17}{}\hyperlink{ref:R9:17}{17}, \hypertarget{cit:R9:18}{}\hyperlink{ref:R9:18}{18}, \hypertarget{cit:R9:19}{}\hyperlink{ref:R9:19}{19}]

A 2026 preprint also proves that a generalized arbitrary-step Collatz transition relation is not definable in base-\(2\) B\"uchi arithmetic. That result concerns the full transition relation rather than a specially selected counterexample orbit, but it reinforces the need to distinguish local digit computation from time-index automaticity. [\hypertarget{cit:R9:20}{}\hyperlink{ref:R9:20}{20}]


\bigskip



\subsubsection*{5. The \(2\)-adic conjugacy is not itself a finite transducer}


The parity map gives a perfect \(2\)-adic coding, but perfect coding does not mean finite-state coding.

Let \(\Phi=Q^{-1}\) be the inverse parity map. For a binary sequence \(s=(s_k)\), one has

\[
\Phi(s)
=
-\sum_{k\geq0}s_k2^k3^{-\sigma_k},
\qquad
\sigma_k=s_0+\cdots+s_k.
\tag{6}
\]

Consider

\[
s^{(m)}=1^m0^\infty.
\]

Then

\[
\Phi(s^{(m)})
=
-\sum_{k=0}^{m-1}\frac{2^k}{3^{k+1}}
=
-1+\left(\frac23\right)^m. \tag{7}
\]

Modulo \(2^m\), this is \(-1\), so its first \(m\) low-order binary digits are all \(1\). Removing those \(m\) digits leaves the \(2\)-adic tail

\[
\frac{\Phi(s^{(m)})-(2^m-1)}{2^m}
=
3^{-m}-1. \tag{8}
\]

These tails are distinct for distinct \(m\).

Suppose a finite synchronous transducer, reading parity digits from least significant to most significant, computed \(\Phi\). After reading the prefix \(1^m\), it would be in some state \(q_m\). Finitely many states imply \(q_m=q_n\) for some \(m\neq n\). In both cases the remaining input is \(0^\infty\), so the transducer would produce the same remaining output. But (8) says the two output tails are different. Contradiction.

Therefore:

\[
\boxed{\Phi\text{ is not realizable by a finite causal synchronous digit transducer}.}
\]

This does not rule out special finite-state behavior on one extraordinary orbit. It does rule out obtaining the desired base-\(2\) automaton merely by invoking the known \(2\)-adic conjugacy.


\bigskip



\subsection*{6. Exact finite-depth \(3\)-adic suffix locality}


Iterating the accelerated recurrence backwards for \(R\) steps gives, for \(j\geq R\),

\[
\begin{aligned}
x_j={}&
3^R2^{-(a_{j-R}+\cdots+a_{j-1})}x_{j-R}
\\
&+
\sum_{k=1}^{R}
3^{k-1}
2^{-(a_{j-k}+\cdots+a_{j-1})}.
\end{aligned}
\]

Reducing modulo \(3^R\) removes the first term:

\[
\boxed{
x_j\equiv
\sum_{k=1}^{R}
3^{k-1}
2^{-(a_{j-k}+\cdots+a_{j-1})}
\pmod{3^R}.
}
\tag{9}
\]

Thus the endpoint modulo \(3^R\) depends only on the final \(R\) valuations, reduced modulo the multiplicative order of \(2\) modulo \(3^R\).

By the lifting-the-exponent identity,

\[
\nu _3\!\left(2^{2\cdot3^{R-1}}-1\right)=R,
\]

and hence

\[
\lambda_R:=\operatorname{ord}_{3^R}(2)
=
2\cdot3^{R-1}. \tag{10}
\]

In particular, \(2\) is a primitive root modulo \(3^R\).

Formula (9) is a genuine finite-depth suffix theorem. But it is a sliding-block computation on the valuation word. It is not an automaton reading the base-\(3\) digits of \(j\).

There is a more serious problem: fixed-depth carry information is maximally non-rigid.


\bigskip



\subsection*{Theorem 3: fixed-depth endpoint universality}


Fix \(R,N\geq1\). Let

\[
r_0\in\mathbb Z/3^R\mathbb Z
\]

be arbitrary, and let

\[
r_1,\ldots ,r_N
\in
(\mathbb Z/3^R\mathbb Z)^\times
\]

be arbitrary units.

Then there exists a positive odd integer \(x_0\) whose first \(N\) accelerated iterates satisfy

\[
x_j\equiv r_j\pmod{3^R}
\qquad (0\leq j\leq N).
\]


\subsubsection*{Proof}


For each \(0\leq j<N\), the desired recurrence modulo \(3^R\) is

\[
r_{j+1}
\equiv
2^{-a_j}(3r_j+1)
\pmod{3^R},
\]

equivalently

\[
2^{a_j}
\equiv
(3r_j+1)r_{j+1}^{-1}
\pmod{3^R}. \tag{11}
\]

The right-hand side is a unit modulo \(3^R\). Since \(2\) generates the full unit group, an exponent \(a_j\geq1\) satisfying (11) always exists.

We have therefore constructed a finite valuation word

\[
(a_0,\ldots ,a_{N-1}).
\]

Consider the corresponding ordinary parity prefix

\[
1\,0^{a_0-1}
1\,0^{a_1-1}
\cdots
1\,0^{a_{N-1}-1}1.
\tag{12}
\]

Let

\[
L=1+\sum_{j=0}^{N-1}a_j.
\]

The bijectivity and isometry of the \(2\)-adic parity map imply that the prefix (12) determines one residue class

\[
x_0\equiv c\pmod{2^L}
\]

of odd integers whose first \(N\) accelerated valuations are exactly the prescribed \(a_j\).

Now combine

\[
x_0\equiv c\pmod{2^L},
\qquad
x_0\equiv r_0\pmod{3^R}
\]

using the Chinese remainder theorem. Choose a positive representative. Its valuation word is the required one, and induction using (11) gives

\[
x_j\equiv r_j\pmod{3^R}
\]

for every \(j\leq N\). \ensuremath{\blacksquare}


\bigskip



\subsubsection*{Consequences}


At any fixed \(3\)-adic depth, there are \textbf{no forbidden finite endpoint words} after the initial position: every finite word over the unit residues modulo \(3^R\) occurs along some positive-integer Collatz segment.

In particular, every finite binary pattern of carry events

\[
\mathbf 1_{\{3^R\mid x_j+1\}}
\]

can occur. To realize a desired \(1\), choose \(r_j=-1\pmod{3^R}\); to realize a desired \(0\), choose any other unit.

This does not prove that the carry stream of one hypothetical divergent orbit is nonautomatic. It proves that automaticity cannot follow from finite-depth local transition constraints or from the absence of certain carry blocks: no such blocks are absent.


\bigskip



\subsubsection*{7. Fixed-depth carries do not determine the valuation word}


Equation (11) determines \(a_j\) only modulo

\[
\lambda_R=2\cdot3^{R-1}.
\]

Thus

\[
a_j
\quad\text{and}\quad
a_j+t\lambda_R
\]

produce the same transition modulo \(3^R\).

Fix any residue \(r_0\pmod{3^R}\). Starting from the same residue, replace arbitrarily many valuations by

\[
a_j'=a_j+\varepsilon_j\lambda_R,
\qquad
\varepsilon_j\in\{0,1\}.
\]

The two residue paths remain identical modulo \(3^R\), while their exact valuation words may encode an arbitrary binary string. Every finite altered valuation word is realized by a positive integer through the parity-cylinder and Chinese-remainder argument above.

Therefore no fixed \(R\) gives a faithful finite-alphabet encoding of the valuation stream.

Only the entire unbounded tower

\[
x_j\bmod3,\quad
x_j\bmod9,\quad
x_j\bmod27,\quad\ldots
\]

could recover all the valuation information. But that tower has no uniform finite state space.


\bigskip



\subsubsection*{8. Periodicity at every fixed depth is still insufficient}


One might try to prove that the carry stream modulo \(3^R\) is eventually periodic separately for each \(R\), and then pass to the inverse limit.

Without uniform bounds on the automata, periods, and starting points, this implication is false.

For example, define

\[
u_n=
\begin{cases}
1+k!,&n=(k!)^2,\\[2mm]
1,&\text{otherwise}.
\end{cases}
\]

For every modulus \(M\), once \(k\geq M\),

\[
k!\equiv0\pmod M,
\]

so

\[
u_n\equiv1\pmod M
\]

for all sufficiently large \(n\). Thus every finite reduction \(u_n\bmod M\) is eventually constant.

Nevertheless, \((u_n)\) is unbounded and not eventually periodic. Moreover,

\[
\sum_{n<N}u_n=N+O(\sqrt N),
\]

so it satisfies the same type of linear cumulative-growth restriction as (2).

Hence even simultaneous eventual periodicity of every finite quotient would not suffice. A successful inverse-limit argument needs a \textbf{uniform finite-state theorem}.


\bigskip



\subsection*{9. A repaired finite-depth theorem}


There is one setting in which fixed-depth carries do become faithful.


\subsection*{Theorem 4}


Fix \(R\), and suppose that for all sufficiently large \(j\),

\[
1\leq a_j<\lambda_R.
\]

Let

\[
r_j=x_j\bmod3^R.
\]

If the finite-valued sequence \((r_j)\) is both \(2\)-automatic and \(3\)-automatic as a function of the time index \(j\), then the orbit eventually enters a cycle.


\subsubsection*{Proof}


Cobham gives eventual periodicity of \(r_j\).

For \(j\geq1\), \(r_j\) is a unit modulo \(3^R\), and

\[
2^{a_j}
\equiv
(3r_j+1)r_{j+1}^{-1}
\pmod{3^R}. \tag{13}
\]

Since \(2\) has order \(\lambda_R\), the pair \((r_j,r_{j+1})\) determines \(a_j\bmod\lambda_R\). The assumed bound \(1\leq a_j<\lambda_R\) then determines \(a_j\) exactly.

An eventually periodic residue-pair sequence therefore gives an eventually periodic valuation sequence. Theorem 1 gives a cycle. \ensuremath{\blacksquare}

This identifies the exact missing uniformity: one would need both time-index automaticity of a sufficiently deep residue stream and an a priori bound on the valuations. Neither is known for a hypothetical counterexample, and arbitrary long finite segments can contain arbitrarily large prescribed valuations.


\bigskip



\subsection*{10. The cycle equation}


Suppose the eventual valuation period is

\[
w=(a_0,\ldots ,a_{m-1})
\]

and put

\[
A=\sum_{i=0}^{m-1}a_i,
\qquad
C_w=\sum_{i=0}^{m-1}3^{m-1-i}2^{A_i}.
\]

At the cycle state \(x\), identity (1) becomes

\[
2^Ax=3^mx+C_w.
\]

Thus

\[
\boxed{
x=\frac{C_w}{2^A-3^m}.
}
\tag{14}
\]

A positive cycle requires

\[
2^A>3^m, \tag{15}
\]

\[
2^A-3^m\mid C_w, \tag{16}
\]

and the resulting orbit must have exactly the valuations \(a_0,\ldots ,a_{m-1}\).

For the trivial accelerated cycle,

\[
w=(2),\qquad A=2,\qquad C_w=1,
\]

and (14) gives

\[
x=\frac1{4-3}=1.
\]

But automaticity provides no reason that another word \(w\) cannot satisfy (14)\textendash{}(16). In fact, any hypothetical nontrivial cycle would have a periodic valuation and carry stream, and every periodic stream is automatic in every base. A finite automaton can simply track the index modulo the period.

Therefore a nontrivial Collatz cycle would already satisfy the proposed Cobham hypothesis. Cobham rigidity cannot contradict it.

This is not a technical defect in the formulation; it is a logical limitation of the method:

\[
\text{dual automaticity}
\Longrightarrow
\text{periodicity},
\]

but

\[
\text{periodicity}
\not\Longrightarrow
\text{the trivial cycle}.
\]

Excluding all nontrivial integral solutions of (14) with the exact valuation conditions is precisely the remaining positive-cycle problem.


\bigskip



\subsection*{Final status}


The Cobham approach yields the rigorous reduction

\[
\boxed{
\begin{array}{c}
\text{time–valuation graph simultaneously recognizable}\\
\text{in bases \(2\) and \(3\)}
\\[1mm]
\Downarrow
\\[1mm]
\text{valuation sequence eventually periodic}
\\[1mm]
\Downarrow
\\[1mm]
\text{positive orbit eventually cyclic}.
\end{array}}
\]

The proof uses two substantive ingredients:

\[
A_n\leq2n+O(1),
\]

which collapses Presburger piecewise linearity to periodicity, and the \(2\)-adic parity conjugacy, which converts a periodic valuation tail into equality of orbit states.

But the proposed full proof breaks at two precise points.

\textbf{First gap: uniform time-index recognizability.} Known binary and ternary transducers compute local arithmetic on orbit values. They do not provide finite \(2\)- and \(3\)-kernels for the sequence indexed by \(j\). Fixed \(3\)-adic depth is non-faithful, and Theorem 3 shows that its finite language is completely unrestricted.

\textbf{Second gap: cycle exclusion.} Even a perfect Cobham theorem for every counterexample stream would at most eliminate divergent trajectories. A nontrivial cycle already has a simultaneously automatic stream. It must be excluded by additional arithmetic beyond Cobham rigidity.

Consequently, the valid achievement of this program is

\[
\boxed{
\text{a sharp automaticity-to-cycle theorem, not a full Collatz proof}.
}
\]

A genuine completion would require a uniform theorem proving simultaneous base-\(2\)/base-\(3\) recognizability or regularity of a faithful encoding of every divergent positive orbit, together with an independent argument showing that the only positive integral solution of the cycle equations with exact valuations is \(x=1\).
\subsection*{References}
\begin{enumerate}[leftmargin=2.2em]
\item \hypertarget{ref:R9:1}{}\hyperlink{cit:R9:1}{\texttt{turn889744view6}}
\item \hypertarget{ref:R9:2}{}\hyperlink{cit:R9:2}{\texttt{turn444644search2}}
\item \hypertarget{ref:R9:3}{}\hyperlink{cit:R9:3}{\texttt{turn359099view0}}
\item \hypertarget{ref:R9:4}{}\hyperlink{cit:R9:4}{\texttt{turn296455view0}}
\item \hypertarget{ref:R9:5}{}\hyperlink{cit:R9:5}{\texttt{turn296455view1}}
\item \hypertarget{ref:R9:6}{}\hyperlink{cit:R9:6}{\texttt{turn296455view2}}
\item \hypertarget{ref:R9:7}{}\hyperlink{cit:R9:7}{\texttt{turn158241view1}}
\item \hypertarget{ref:R9:8}{}\hyperlink{cit:R9:8}{\texttt{turn296455view3}}
\item \hypertarget{ref:R9:9}{}\hyperlink{cit:R9:9}{\texttt{turn889744view8}}
\item \hypertarget{ref:R9:10}{}\hyperlink{cit:R9:10}{\texttt{turn296455view4}}
\item \hypertarget{ref:R9:11}{}\hyperlink{cit:R9:11}{\texttt{turn889744view0}}
\item \hypertarget{ref:R9:12}{}\hyperlink{cit:R9:12}{\texttt{turn889744view1}}
\item \hypertarget{ref:R9:13}{}\hyperlink{cit:R9:13}{\texttt{turn889744view2}}
\item \hypertarget{ref:R9:14}{}\hyperlink{cit:R9:14}{\texttt{turn889744view3}}
\item \hypertarget{ref:R9:15}{}\hyperlink{cit:R9:15}{\texttt{turn889744view4}}
\item \hypertarget{ref:R9:16}{}\hyperlink{cit:R9:16}{\texttt{turn889744view5}}
\item \hypertarget{ref:R9:17}{}\hyperlink{cit:R9:17}{\texttt{turn296455view5}}
\item \hypertarget{ref:R9:18}{}\hyperlink{cit:R9:18}{\texttt{turn296455view6}}
\item \hypertarget{ref:R9:19}{}\hyperlink{cit:R9:19}{\texttt{turn296455view7}}
\item \hypertarget{ref:R9:20}{}\hyperlink{cit:R9:20}{\texttt{turn102710view0}}
\end{enumerate}

\clearpage
\section{Round 10: approach 5 --- Prove simplicity or irreducibility of the Collatz C*-algebra}
\subsection*{Prompt}
\begin{quote}
Complete a full proof of the Collatz conjecture using the following approach:

\begin{enumerate}
\item Prove simplicity or irreducibility of the Collatz C*-algebra
\end{enumerate}

\textbf{Established/Hybrid \textperiodcentered{} P}

Mori's operator-algebraic encoding says that Collatz is equivalent to the absence of nontrivial reducing subspaces for a natural algebra. Attack this through:

\begin{itemize}
\item a gauge action and KMS-state uniqueness;
\item direct computation of the commutant;
\item Cuntz-Pimsner simplicity criteria;
\item groupoid minimality plus topological principality.
\end{itemize}

The essential requirement is that the operator-algebra theorem be established from local algebraic relations, not by silently restating orbit connectivity.

Take as long as you need.
\end{quote}
\subsection*{Response}

\subsection*{Outcome}


I do not obtain a proof of the Collatz conjecture from this program. The obstruction is sharper than an unresolved technical lemma:

\begin{enumerate}
\item \textbf{Mori's abstract Cuntz algebra is already isomorphic to \(\mathcal O _2\), and hence simple, without assuming Collatz.}
\item \textbf{The concrete representation is irreducible exactly when Collatz holds.}
\item Its commutant can be computed explicitly, from the local branch relations, as the algebra of bounded functions constant on total Collatz-orbit classes.
\end{enumerate}

Thus abstract simplicity, gauge/KMS uniqueness, and the usual full-shift groupoid simplicity cannot distinguish a world with one Collatz component from a world with several. Mori defines the two branch operators below and proves that irreducibility of their generated algebra is equivalent to Collatz; his first-return construction generates an algebra isomorphic to \(\mathcal O _2\) unconditionally. [\hypertarget{cit:R10:1}{}\hyperlink{ref:R10:1}{1}, \hypertarget{cit:R10:2}{}\hyperlink{ref:R10:2}{2}, \hypertarget{cit:R10:3}{}\hyperlink{ref:R10:3}{3}, \hypertarget{cit:R10:4}{}\hyperlink{ref:R10:4}{4}]

The strongest rigorous result from the proposed approach is the following exact commutant theorem.


\bigskip



\subsubsection*{1. Exact computation of the commutant}


Let

\[
f(n)=
\begin{cases}
3n+1,&n\ \text{odd},\\[2mm]
n/2,&n\ \text{even},
\end{cases}
\]

and let \(H=\ell ^2(\mathbb N)\), with canonical basis \((e_n)_{n\geq1}\). Define

\[
T_o e_n=
\begin{cases}
e_{3n+1},&n\ \text{odd},\\
0,&n\ \text{even},
\end{cases}
\qquad
T_e e_n=
\begin{cases}
0,&n\ \text{odd},\\
e_{n/2},&n\ \text{even}.
\end{cases}
\]

Set

\[
\mathcal A_f=C^*(T_o,T_e)\subseteq B(H).
\]

Define total-orbit equivalence by

\[
m\sim_f n
\quad\Longleftrightarrow\quad
f^r(m)=f^s(n)
\quad\text{for some }r,s\geq0.
\]

Let \(\mathscr C_f=\mathbb N/{\sim_f}\) be the set of total-orbit classes.


\paragraph*{Theorem}


The commutant of \(\mathcal A_f\) is

\[
\boxed{
\mathcal A_f'
=
\left\{
M_d:
d\in\ell^\infty(\mathbb N),\
d(f(n))=d(n)\ \text{for every }n
\right\},
}
\]

where \(M_de_n=d(n)e_n\). Equivalently,

\[
\boxed{
\mathcal A_f'\cong \ell^\infty(\mathscr C_f).
}
\]

More explicitly, if \(P_C\) denotes the coordinate projection onto

\[
H_C=\overline{\operatorname{span}}\{e_n:n\in C\},
\]

then every \(X\in\mathcal A_f'\) has the unique form

\[
X=\sum_{C\in\mathscr C_f}\lambda_C P_C,
\qquad
(\lambda_C)_{C\in\mathscr C_f}\in\ell^\infty(\mathscr C_f),
\]

with strong-operator convergence. Consequently,

\[
\mathcal A_f''
=
\prod_{C\in\mathscr C_f} B(H_C).
\]


\paragraph*{Proof}



\subparagraph*{Step 1: finite parity words give projections in \(\mathcal A_f\)}


For \(n\in\mathbb N\), let

\[
\varepsilon_j(n)=
\begin{cases}
o,&f^j(n)\text{ is odd},\\
e,&f^j(n)\text{ is even}.
\end{cases}
\]

Given a word

\[
\omega=(\omega_0,\ldots,\omega_{r-1})\in\{o,e\}^r,
\]

put

\[
V_\omega=T_{\omega_{r-1}}\cdots T_{\omega_0}.
\]

The operators \(T_o,T_e\) follow the appropriate branch when that branch is admissible and annihilate the vector otherwise. Hence

\[
V_\omega e_m=
\begin{cases}
e_{f^r(m)},&
(\varepsilon_0(m),\ldots,\varepsilon_{r-1}(m))=\omega,\\
0,&\text{otherwise}.
\end{cases}
\]

Both branches are injective on their domains, so

\[
V_\omega^*V_\omega
\]

is the diagonal projection onto the finite parity cylinder

\[
D_\omega
=
\left\{
m:
(\varepsilon_0(m),\ldots,\varepsilon_{r-1}(m))=\omega
\right\}.
\]

In particular, \(V_\omega^*V_\omega\in\mathcal A_f\).


\subparagraph*{Step 2: the infinite parity itinerary identifies the integer}


Suppose \(m,n\) have the same first \(r\) branch choices. Let \(o_r\) and \(e_r\) be the numbers of odd and even steps among these choices. Since the two trajectories undergo the same affine branch at each step,

\[
f^r(n)-f^r(m)
=
\frac{3^{o_r}}{2^{e_r}}(n-m).
\]

The left-hand side is an integer and \(3^{o_r}\) is odd, so

\[
2^{e_r}\mid n-m.
\]

An odd Collatz step produces an even number, so two odd steps cannot occur consecutively. Therefore

\[
e_r\geq \left\lfloor\frac r2\right\rfloor\longrightarrow\infty.
\]

If the complete itineraries of \(m\) and \(n\) agree, then \(2^{e_r}\mid n-m\) for arbitrarily large \(e_r\), forcing \(m=n\).

Thus the infinite parity coding is injective. This part is entirely local and uses no form of the Collatz conjecture.


\subparagraph*{Step 3: every commutant operator is diagonal}


Fix \(n\). Let \(P_r(n)\) be the cylinder projection associated with the first \(r\) symbols of the itinerary of \(n\). Then

\[
P_1(n)\geq P_2(n)\geq\cdots
\]

and the intersection of their supports is \(\{n\}\). It follows that

\[
P_r(n)\xrightarrow[\ r\to\infty\ ]{\mathrm{strongly}}
p_n,
\qquad
p_n=|e_n\rangle\langle e_n|.
\]

If \(X\in\mathcal A_f'\), then \(XP_r(n)=P_r(n)X\) for every \(r\). Passing to the strong limit gives

\[
Xp_n=p_nX.
\]

This holds for every \(n\). Therefore \(X\) commutes with all coordinate rank-one projections, so \(X\) is diagonal:

\[
Xe_n=d(n)e_n
\]

for some bounded sequence \(d\).

The important distinction is that \(p_n\) belongs to the \textbf{strong closure} of the cylinder algebra, not generally to the norm-closed Cuntz algebra. This is precisely where the orbit information enters.


\subparagraph*{Step 4: commutation is exactly invariance along Collatz edges}


For odd \(n\),

\[
M_dT_oe_n=d(3n+1)e_{3n+1},
\qquad
T_oM_de_n=d(n)e_{3n+1}.
\]

Thus \(M_dT_o=T_oM_d\) exactly when

\[
d(3n+1)=d(n)
\qquad(n\text{ odd}).
\]

Similarly, commutation with \(T_e\) is equivalent to

\[
d(n/2)=d(n)
\qquad(n\text{ even}).
\]

Together these say

\[
d(f(n))=d(n).
\]

The same condition for \(\overline d\), followed by taking adjoints, shows that \(M_d\) also commutes with \(T_o^*,T_e^*\). Hence it lies in \(\mathcal A_f'\).

Finally, \(d(f(n))=d(n)\) holds exactly when \(d\) is constant on every total-orbit class. This proves the theorem. \(\square\)

Mori's later general theorem proves the corresponding projection-level statement: under injectivity of the branch itinerary, invariant subsets and reducing subspaces are in order-preserving bijection. The \(qx+d\) maps, including Collatz, satisfy the required itinerary-injectivity condition. [\hypertarget{cit:R10:5}{}\hyperlink{ref:R10:5}{5}, \hypertarget{cit:R10:6}{}\hyperlink{ref:R10:6}{6}, \hypertarget{cit:R10:7}{}\hyperlink{ref:R10:7}{7}]


\bigskip



\subsubsection*{2. Why this is exactly Collatz}


The theorem gives

\[
\mathcal A_f'\;=\;\mathbb CI
\quad\Longleftrightarrow\quad
|\mathscr C_f|=1.
\]

The orbit of \(1\) is the cycle

\[
1\longmapsto4\longmapsto2\longmapsto1.
\]

Therefore, if every positive integer lies in the total-orbit class of \(1\), its forward trajectory eventually meets \(1,2\), or \(4\), and hence reaches \(1\). Conversely, Collatz plainly implies that all integers lie in that class. Thus

\[
\boxed{
\mathcal A_f\text{ is irreducible}
\quad\Longleftrightarrow\quad
\mathscr C_f=\{[1]\}
\quad\Longleftrightarrow\quad
\text{Collatz}.
}
\]

The exact remaining assertion is consequently

\[
\boxed{
\begin{aligned}
d(2n)&=d(n),\\
d(3n+1)&=d(n)\quad(n\text{ odd}),\\
d&\in\ell^\infty(\mathbb N)
\end{aligned}
\quad\Longrightarrow\quad
d\text{ is constant}.
}
\]

But if a nontrivial component \(C\) existed, its indicator

\[
d_C(n)=\mathbf 1_C(n)
\]

would be a bounded nonconstant solution. Therefore proving the boxed assertion is not an operator-algebraic shortcut: it is precisely the assertion that there are no additional Collatz components.

This also identifies the exact place where a seemingly local argument can become circular. Saying that the relations ``propagate the value of \(d\) to every basis vector'' is valid only after proving that the branch graph is connected.


\bigskip



\subsubsection*{3. Simplicity is not irreducibility}


Mori's first-return construction produces isometries \(S_1,S_2\) satisfying

\[
S_i^*S_j=\delta_{ij}I,
\qquad
S_1S_1^*+S_2S_2^*=I.
\]

Consequently,

\[
C^*(S_1,S_2)\cong\mathcal O_2
\]

without assuming Collatz. Since \(\mathcal O_2\) is simple, the abstract ``Collatz \(C^*\)-algebra'' is already simple. [\hyperlink{ref:R10:3}{3}, \hyperlink{ref:R10:4}{4}]

There is no contradiction because:

\begin{itemize}
\item \textbf{Simplicity} means the abstract algebra has no nonzero proper closed two-sided ideals.
\item \textbf{Irreducibility} means this particular representation on \(\ell^2(\mathbb N)\), or on Mori's first-return set, has scalar commutant.
\end{itemize}

A simple algebra can act reducibly. For example, if \(\rho\) is any faithful representation of \(\mathcal O_2\), then

\[
\rho\oplus\rho
\]

is faithful and its image is still abstractly isomorphic to \(\mathcal O_2\), but the projection onto the first summand is a nontrivial commutant projection.

Indeed, if the Collatz representation had several total-orbit components, it would decompose as

\[
\pi=\bigoplus_{C\in\mathscr C}\pi_C.
\]

Every \(\pi_C\) is nonzero. Simplicity of \(\mathcal O_2\) forces every \(\pi_C\) to be faithful, while the one-component version of the commutant theorem makes it irreducible. Thus a hypothetical counterexample would merely turn the canonical representation into a direct sum of faithful irreducible representations of the same simple algebra.

All abstract invariants\textemdash{}simplicity, nuclearity, pure infiniteness, \(K\)-theory, and the gauge system\textemdash{}would remain unchanged.

This proves a useful no-go statement:


\begin{quote}
\textbf{No theorem depending only on the universal Cuntz relations or on the abstract isomorphism class of \(\mathcal O_2\) can prove irreducibility of Mori's distinguished representation.}

\end{quote}



\bigskip



\subsubsection*{4. Gauge action and KMS-state uniqueness}


Let \(\gamma\) be the gauge dynamics on \(\mathcal O_2\),

\[
\gamma_t(S_i)=e^{it}S_i.
\]

Suppose \(\varphi\) is a \(\mathrm{KMS}_\beta\) state. Applying the KMS identity to \(S_i,S_i^*\) gives

\[
\varphi(S_iS_i^*)
=
\varphi\!\left(S_i^*\gamma_{i\beta}(S_i)\right)
=
e^{-\beta}\varphi(S_i^*S_i)
=
e^{-\beta}.
\]

Since the range projections sum to \(I\),

\[
1=\varphi(I)
 =\varphi(S_1S_1^*)+\varphi(S_2S_2^*)
 =2e^{-\beta}.
\]

Hence

\[
\beta=\log 2.
\]

For words \(\mu,\nu\), gauge invariance and the KMS identity similarly give

\[
\varphi(S_\mu S_\nu^*)
=
\delta_{\mu,\nu}\,2^{-|\mu|}.
\]

These values determine \(\varphi\) on the dense algebraic span of the Cuntz monomials. Hence the gauge action has exactly one KMS state, at inverse temperature \(\log2\).

This is a complete local-relations proof of KMS uniqueness\textemdash{}but it proves nothing about Collatz, because it applies to abstract \(\mathcal O_2\) before the arithmetic representation is considered.

There is an even sharper obstruction. On the canonical diagonal

\[
\mathcal D_2
=
C^*(S_\mu S_\mu^*)
\cong C(\{1,2\}^{\mathbb N}),
\]

the KMS state is fair Bernoulli measure:

\[
\varphi(S_\mu S_\mu^*)=2^{-|\mu|}.
\]

It gives every individual infinite sequence measure zero. Mori's arithmetic representation is supported on a countable set \(\Omega\) of branch sequences. Therefore the KMS state is \textbf{not normal} in this atomic representation: a normal state on \(B(\ell^2(\Omega))\) would restrict to a countably supported atomic probability measure, whereas Bernoulli measure gives \(\Omega\) total mass zero.

Consequently, a reducing projection cannot be used to split the canonical KMS state into two normal KMS states. The usual ``unique equilibrium state implies one component'' argument has no foothold here.


\bigskip



\subsubsection*{5. Cuntz\textendash{}Pimsner simplicity criteria}


There are two natural choices, and neither closes the problem.


\paragraph*{Scalar coefficient algebra}


Using the correspondence

\[
E=\mathbb C^2
\]

over \(\mathbb C\) gives

\[
\mathcal O_E=\mathcal O_2.
\]

Its minimality, aperiodicity, and simplicity are independent of Collatz. The arithmetic support of the representation has been forgotten.


\paragraph*{Arithmetic coefficient algebra}


Instead take a correspondence over \(c_0(\mathbb N)\), or over the first-return set, with edges given by the Collatz branches. If \(C\) is a total-orbit component, then

\[
I_C=c_0(C)
\]

is invariant under both forward and backward branch operations. Hence any hypothetical additional Collatz component supplies a nontrivial invariant coefficient ideal.

The ``minimality'' hypothesis in a Cuntz\textendash{}Pimsner simplicity theorem is therefore precisely the assertion that no such \(I_C\) exists. Because the component of the known \(1\)-cycle is all of \(\mathbb N\) exactly when Collatz holds, establishing correspondence minimality would merely re-establish Collatz in Cuntz\textendash{}Pimsner language.

Aperiodicity may be approachable locally, but minimality remains the complete global problem.


\bigskip



\subsubsection*{6. The groupoid route}


Let

\[
X_2=\{1,2\}^{\mathbb N}
\]

and let \(\sigma\) be the one-sided shift. The standard Cuntz groupoid is

\[
G_2
=
\left\{
(x,k-\ell,y):
\sigma^k(x)=\sigma^\ell(y)
\right\}.
\]

Every \(G_2\)-orbit is dense. Indeed, given \(x\in X_2\) and a cylinder determined by a finite word \(\mu\), the sequence

\[
\mu x
\]

lies in that cylinder and satisfies

\[
\sigma^{|\mu|}(\mu x)=x.
\]

Aperiodic sequences are dense, so \(G_2\) is also topologically principal. These facts yield the standard simplicity of \(\mathcal O_2\), unconditionally.

Now let

\[
\Omega=\{\widehat n:n\text{ is an arithmetic basis label}\}\subset X_2
\]

be the set of Mori branch codes. It is countable and invariant under deleting or prefixing admissible symbols. Every orbit contained in \(\Omega\) is still dense in \(X_2\), and hence dense in \(\Omega\) with the inherited topology.

Therefore:

\begin{itemize}
\item if Collatz is true, \(\Omega\) is one dense orbit;
\item if Collatz is false, \(\Omega\) is a union of several disjoint dense orbits.
\end{itemize}

Topological minimality distinguishes neither case.

If one instead puts the discrete topology on \(\Omega\), then an orbit is dense exactly when it equals all of \(\Omega\). Thus groupoid minimality becomes exactly the Collatz assertion. Moreover, the known periodic \(1\)-cycle creates nontrivial isotropy at an open unit, so the discrete groupoid is not topologically principal.

The groupoid program therefore has the same dichotomy:

\begin{itemize}
\item the compact full-shift groupoid is simple for purely symbolic reasons and forgets arithmetic components;
\item the discrete arithmetic groupoid detects the components, but its minimality is Collatz itself.
\end{itemize}

The underlying reason is visible in the \(2\)-adic coding. Finite parity words run through all binary possibilities, and on \(\mathbb Z_2\) the parity transform is a bijection conjugating the shortcut Collatz map to the binary shift. Thus finite cylinder relations naturally see the entire \(2\)-adic full shift, not the special Archimedean subset of positive integers. [\hypertarget{cit:R10:8}{}\hyperlink{ref:R10:8}{8}]


\bigskip



\subsubsection*{7. Genuine unconditional rigidity of the commutant}


Although the full commutant problem remains equivalent to Collatz, the local equations do rule out all reasonably regular nonconstant solutions.


\paragraph*{Convergence at infinity}


Suppose \(d(f(n))=d(n)\) and \(d(n)\to L\) as \(n\to\infty\). Since

\[
d(n)=d(2n)=d(2^2n)=\cdots,
\]

we have

\[
d(n)=\lim_{k\to\infty}d(2^kn)=L.
\]

Thus \(d\equiv L\). Consequently,

\[
\boxed{
\mathcal A_f'\cap(\mathbb CI+\mathcal K(H))
=
\mathbb CI
}
\]

and

\[
\boxed{
\mathcal A_f'\cap\mathcal K(H)=\{0\}.
}
\]


\paragraph*{Summable coefficients}


If \(d\in\ell^p(\mathbb N)\), \(1\leq p<\infty\), then a nonzero value \(d(n)\) is repeated at the infinitely many points \(2^kn\). Hence its \(\ell^p\)-norm would be infinite. Therefore

\[
d\in\ell^p,\quad d\circ f=d
\quad\Longrightarrow\quad
d=0.
\]


\paragraph*{Periodic coefficients}


Suppose \(d\) has period \(q\).

If \(2\mid q\), then

\[
d\!\left(n+\frac q2\right)
=d(2n+q)
=d(2n)
=d(n),
\]

so \(q/2\) is also a period. Remove all powers of \(2\), leaving \(q\) odd.

For odd \(q\), every residue modulo \(q\) has an odd representative, so the relation

\[
d(3n+1)=d(n)
\]

holds as a relation on every residue class modulo \(q\).

If \(3\mid q\), then residues \(a\) and \(a+q/3\) have images under \(x\mapsto3x+1\) that agree modulo \(q\). Hence \(q/3\) is also a period. Remove all powers of \(3\).

Now \(\gcd(q,6)=1\). On \(\mathbb Z/q\mathbb Z\), set

\[
A(x)=2x,\qquad B(x)=3x+1.
\]

Invariance under \(A\) and \(B\) gives

\[
d(6x+1)=d(B(Ax))=d(x)
\]

and

\[
d(6x+2)=d(A(Bx))=d(x).
\]

Since multiplication by \(6\) is invertible modulo \(q\), the residues \(6x+1\) run through all residue classes, and \(6x+2\) is the immediately succeeding class. Thus adjacent residue classes have equal values, so \(d\) is constant.

Therefore every periodic commutant element is scalar.


\paragraph*{\(2\)-adic continuity}


The invariance equations also imply invariance under the shortcut map

\[
\widetilde f(n)=
\begin{cases}
(3n+1)/2,&n\text{ odd},\\
n/2,&n\text{ even}.
\end{cases}
\]

Suppose \(d\) extends continuously from \(\mathbb N\) to \(\mathbb Z_2\). Under the parity conjugacy, the extension becomes a continuous shift-invariant function \(g\) on \(\{0,1\}^{\mathbb N}\).

For arbitrary \(x,y\), let \(y|_r\,x\) denote the sequence formed by prefixing the first \(r\) symbols of \(y\) to \(x\). Then

\[
\sigma^r(y|_r\,x)=x,
\]

so shift invariance gives

\[
g(y|_r\,x)=g(x).
\]

But \(y|_r\,x\to y\), and continuity gives \(g(y)=g(x)\). Hence \(g\), and therefore \(d\), is constant.

Thus any nontrivial commutant projection would have to be simultaneously:

\begin{itemize}
\item noncompact;
\item non-\(\ell^p\);
\item nonperiodic;
\item nonconvergent at ordinary infinity;
\item discontinuous in the \(2\)-adic topology.
\end{itemize}

It is necessarily a highly irregular total-orbit indicator.


\bigskip



\subsubsection*{8. The exact theorem a successful operator-algebra proof would need}


Let \(\pi(\mathcal O_2)\) be Mori's concrete arithmetic representation, and let \(\mathcal D_2\) be its canonical cylinder diagonal. The cylinder calculation proves

\[
\pi(\mathcal O_2)'
\subseteq
\pi(\mathcal D_2)'',
\]

but not

\[
\pi(\mathcal O_2)'
\subseteq
\pi(\mathcal D_2).
\]

The strong closure \(\pi(\mathcal D_2)''\) contains arbitrary coordinate and component projections; the norm algebra \(\pi(\mathcal D_2)\) contains only continuous cylinder-limit functions. Since

\[
\pi(\mathcal D_2)\cap\pi(\mathcal O_2)'=\mathbb CI,
\]

the following would prove Collatz:

\[
\boxed{
Q\in\pi(\mathcal O_2)',\quad Q^2=Q
\quad\Longrightarrow\quad
Q\in\pi(\mathcal D_2).
}
\]

A weaker Archimedean version would also suffice:

\[
\boxed{
Q\in\mathcal A_f',\quad Q^2=Q
\quad\Longrightarrow\quad
\lim_{n\to\infty}\langle Qe_n,e_n\rangle
\text{ exists}.
}
\]

Indeed, the commutant theorem makes \(Q\) diagonal, and invariance under \(n\mapsto2n\) would then force every diagonal entry to equal that limit.

These are genuine analytic regularity statements rather than verbal assertions of orbit connectivity. However, they cannot follow from the universal Cuntz relations, gauge action, abstract simplicity, or KMS uniqueness: reducible direct sums of faithful \(\mathcal O_2\)-representations disprove any such universal implication. Proving one would require additional arithmetic or Archimedean structure capable of distinguishing the embedded positive integers from their \(2\)-adic symbolic closure.


\subsubsection*{Verdict}


The proposed program does \textbf{not} currently yield a full Collatz proof. It yields instead a complete structural diagnosis:

\[
\boxed{
\begin{array}{c}
\text{abstract simplicity and gauge/KMS uniqueness}\\
\text{hold unconditionally,}
\end{array}
}
\]

while

\[
\boxed{
\begin{array}{c}
\text{concrete irreducibility}\\
=\text{absence of nonconstant orbit-class indicators}\\
=\text{Collatz}.
\end{array}
}
\]

So the direct commutant route reaches the exact boundary of the problem, and the other three routes either forget the arithmetic representation or reintroduce the conjecture as ``minimality.'' Current 2026 sources continue to describe the Collatz conjecture as open. [\hypertarget{cit:R10:9}{}\hyperlink{ref:R10:9}{9}, \hypertarget{cit:R10:10}{}\hyperlink{ref:R10:10}{10}]
\subsection*{References}
\begin{enumerate}[leftmargin=2.2em]
\item \hypertarget{ref:R10:1}{}\hyperlink{cit:R10:1}{\texttt{turn172462view0}}
\item \hypertarget{ref:R10:2}{}\hyperlink{cit:R10:2}{\texttt{turn172462view1}}
\item \hypertarget{ref:R10:3}{}\hyperlink{cit:R10:3}{\texttt{turn975213view3}}
\item \hypertarget{ref:R10:4}{}\hyperlink{cit:R10:4}{\texttt{turn975213view4}}
\item \hypertarget{ref:R10:5}{}\hyperlink{cit:R10:5}{\texttt{turn975213view0}}
\item \hypertarget{ref:R10:6}{}\hyperlink{cit:R10:6}{\texttt{turn975213view1}}
\item \hypertarget{ref:R10:7}{}\hyperlink{cit:R10:7}{\texttt{turn975213view2}}
\item \hypertarget{ref:R10:8}{}\hyperlink{cit:R10:8}{\texttt{turn172462view4}}
\item \hypertarget{ref:R10:9}{}\hyperlink{cit:R10:9}{\texttt{turn533985search1}}
\item \hypertarget{ref:R10:10}{}\hyperlink{cit:R10:10}{\texttt{turn533985search12}}
\end{enumerate}

\clearpage
\section{Round 11: approach 6 --- Sheaf cohomology for gluing modular Lyapunov functions}
\subsection*{Prompt}
\begin{quote}
Complete a full proof of the Collatz conjecture using the following approach:

\begin{enumerate}
\item Sheaf cohomology for gluing modular Lyapunov functions
\end{enumerate}

\textbf{Speculative \textperiodcentered{} P}

For many moduli M, one can hope to construct local descent potentials

V\_M(n) = log n + phi\_M(n mod M)

that work on most residue classes. Regard these potentials as local sections on a system of congruence covers.

Failure to glue them into one global potential would define a cohomological obstruction. Prove that the relevant H\textasciicircum{}1 vanishes under refinement, and the local certificates glue into a global Lyapunov function. A counterexample would instead produce a persistent nonzero class.

Take as long as you need.
\end{quote}
\subsection*{Response}

\subsection*{Verdict}


I do not obtain a valid proof of the Collatz conjecture from this approach. More strongly, the proposed mechanism has a rigorous obstruction:


\begin{quote}
\textbf{Ordinary sheaf \(H^{1}\) on the congruence space vanishes for formal topological reasons, but that vanishing does not glue descent inequalities. Under every natural regularity assumption, a global correction of the form} \[
> V(n)=\log n+\phi(n)
> \] \textbf{is impossible.}

\end{quote}


The relevant dynamical cohomology does contain an obstruction, but it is already nonzero because of negative \(2\)-adic periodic points such as \(-1\), regardless of whether a positive-integer counterexample exists. Moreover, any correction that really did decrease along every positive Collatz trajectory would have to be unbounded on every \(2\)-adic congruence neighborhood.

Thus this approach, as stated, cannot complete Collatz. What follows is a full proof of that no-go result and an exact description of the stronger statement a repaired version would need.


\bigskip



\subsubsection*{1. Setup}


Use the accelerated odd map \[
U(n)=\frac{3n+1}{2^{a(n)}},
\qquad
a(n)=\nu _2(3n+1),
\] on the positive odd integers. The Collatz conjecture is equivalent to \[
\forall n\in 2\mathbb N+1,\qquad U^j(n)=1
\quad\text{for some }j\geq 0.
\]

A bounded modular correction \[
V_M(n)=\log n+\phi_M(n\bmod M)
\] would prove Collatz if \[
V_M(U(n))<V_M(n)\qquad(n>1).
\] Indeed, \(\phi_M\) is bounded, so a strictly decreasing orbit has bounded \(n\)-values. It would therefore eventually repeat, contradicting strict decrease unless it had already reached \(1\).

The first obstruction is completely elementary.


\bigskip



\subsubsection*{2. No fixed modulus can support the proposed potential}



\paragraph*{Theorem 1: fixed-modulus obstruction}


For every \(M\geq1\) and every function \[
\phi_M:\mathbb Z/M\mathbb Z\longrightarrow\mathbb R,
\] there are infinitely many positive odd \(n\) for which \[
U(n)\equiv n\pmod M
\qquad\text{and}\qquad
U(n)>n.
\] Consequently, \[
\log U(n)+\phi_M(U(n)\bmod M)
>
\log n+\phi_M(n\bmod M).
\]


\paragraph*{Proof}


Set \[
n=4Mk-1
\] for any \(k\geq1\). Then \[
3n+1=12Mk-2=2(6Mk-1),
\] and \(6Mk-1\) is odd. Hence \[
a(n)=1
\] and \[
U(n)=6Mk-1.
\] Moreover, \[
U(n)-n=2Mk,
\] so \[
U(n)\equiv n\equiv -1\pmod M.
\] But \(U(n)>n\). Therefore the modular correction is unchanged while the logarithmic part increases: \[
V_M(U(n))-V_M(n)
=
\log\frac{6Mk-1}{4Mk-1}>0.
\] This proves the theorem. \(\square\)


\paragraph*{Corollary 1}


No finite family of residue-dependent local potentials can glue to a global strict Lyapunov function.

Indeed, take the least common multiple \(L\) of all the moduli in the finite family. After passing to the common refinement modulo \(L\), the glued correction is a function of \(n\bmod L\). Theorem 1 applies with \(M=L\).

This also defeats finite-state certificates using a bounded number of accelerated steps.


\paragraph*{Proposition 2: bounded-lookahead obstruction}


Suppose that each residue \(r\bmod M\) is assigned some positive number of steps \[
\tau(r)\leq L
\] and one asks for \[
V_M\bigl(U^{\tau(r)}(n)\bigr)<V_M(n)
\qquad(n\equiv r\pmod M).
\] This is impossible.


\paragraph*{Proof}


Let \(\ell=\tau(-1\bmod M)\), and choose \(K\geq\ell+2\). Put \[
n=2^K M-1.
\] For the first \(\ell\) accelerated steps the valuation is \(1\), and induction gives \[
U^j(n)=3^j2^{K-j}M-1,
\qquad 0\leq j\leq\ell.
\] Therefore \[
U^\ell(n)\equiv n\equiv-1\pmod M
\] and \[
\frac{U^\ell(n)}{n}
=
\frac{3^\ell2^{K-\ell}M-1}{2^K M-1}>1.
\] Again the correction is unchanged while the logarithmic term grows. \(\square\)

Thus any successful congruence construction would need both infinitely increasing refinement depth and unbounded lookahead.


\bigskip



\subsubsection*{3. Why ordinary sheaf \(H^1\) does not help}


Let \[
X=\widehat{\mathbb Z}
\quad\text{or}\quad
X=\mathbb Z_2^\times.
\] The congruence classes are clopen subsets of \(X\). A cover by the classes modulo \(M\) is actually a disjoint partition, so its normalized \v{C}ech complex has no positive-degree terms. More generally, profinite spaces have sheaf-theoretic cohomological dimension zero: \[
H^j(X,\mathcal F)=0
\qquad(j>0)
\] for every sheaf of abelian groups \(\mathcal F\). [\hypertarget{cit:R11:1}{}\hyperlink{ref:R11:1}{1}]

Therefore the desired statement \[
H^1(X,\mathcal F)=0
\] is automatic. It does not depend on Collatz dynamics and cannot distinguish a universe in which Collatz is true from one containing a counterexample.

There is also a categorical mismatch. The assignment \[
O\longmapsto
\left\{
\phi:O\to\mathbb R:
\log U(n)+\phi(U(n))<\log n+\phi(n)
\right\}
\] is not an ordinary sheaf of abelian groups:

\begin{enumerate}
\item The condition at \(n\in O\) refers to the value at \(U(n)\), which may lie outside \(O\).
\item The set of strict subsolutions is a convex cone, not an abelian group.
\item The difference of two strict subsolutions need not satisfy a fixed homogeneous equation.
\item Adjusting local sections by a \v{C}ech coboundary need not preserve the strict inequalities.
\end{enumerate}

Ordinary sheaf cohomology can glue compatible values. It cannot manufacture compatibility with directed Collatz edges.


\bigskip



\subsubsection*{4. No continuous profinite correction can work}


The obstruction already appears at the negative \(2\)-adic fixed point \(-1\).


\paragraph*{Proposition 3}


There is no continuous function \[
\phi:\mathbb Z_2^\times\longrightarrow\mathbb R
\] such that \[
\log U(n)+\phi(U(n))
<
\log n+\phi(n)
\] for every positive odd \(n>1\).


\paragraph*{Proof}


Take \[
n_k=2^k-1.
\] For \(k\geq2\), \[
3n_k+1=3\cdot2^k-2
=2(3\cdot2^{k-1}-1),
\] where the second factor is odd. Hence \[
U(n_k)=3\cdot2^{k-1}-1.
\] Both sequences converge \(2\)-adically to \(-1\): \[
n_k\longrightarrow -1,
\qquad
U(n_k)\longrightarrow -1.
\] Continuity gives \[
\phi(U(n_k))-\phi(n_k)\longrightarrow0.
\] On the other hand, \[
\log\frac{U(n_k)}{n_k}
\longrightarrow
\log\frac32>0.
\] Therefore \[
V(U(n_k))-V(n_k)\longrightarrow\log\frac32>0,
\] contradicting descent. \(\square\)

So a uniformly convergent inverse limit of modular corrections cannot work: such a limit would be continuous.

The obstruction is much stronger than failure at the single point \(-1\).


\bigskip



\subsubsection*{5. Every successful correction must be nowhere locally bounded}


We now prove that any strict correction \(\phi\), even a discontinuous one defined only on positive odd integers, must have infinite oscillation in every \(2\)-adic open set.

For \(a\geq1\), define the inverse branch \[
F_a(y)=\frac{2^a y-1}{3}.
\] For every odd \(2\)-adic integer \(y\), \[
3F_a(y)+1=2^a y,
\] so \[
\nu_2(3F_a(y)+1)=a,
\qquad
U(F_a(y))=y.
\]

For a finite valuation word \[
w=(a_0,\ldots,a_{q-1}),
\] put \[
A(w)=a_0+\cdots+a_{q-1}.
\] On the cylinder realizing this word, one has the exact affine identity \[
2^{A(w)}U^q(x)=3^q x+C_w,
\] where \[
C_0=0,\qquad
C_{j+1}=3C_j+2^{A_j}.
\]


\paragraph*{Lemma 4: expanding periodic points are \(2\)-adically dense}


Every nonempty \(2\)-adic open subset of the odd integers contains a periodic point \(p\) of \(U\) with an expanding real multiplier \[
\lambda=\frac{3^q}{2^{A(w)}}>1.
\]


\paragraph*{Proof}


Choose a finite valuation cylinder \(C_u\) contained in the given open set. Append sufficiently many symbols \(1\) to the word \(u\), obtaining \[
w=u\,1\,1\cdots1.
\] Since \[
\frac{A(w)}{|w|}\longrightarrow1
\] as the number of appended \(1\)'s tends to infinity, eventually \[
A(w)<|w|\log_2 3.
\] Thus \[
\frac{3^{|w|}}{2^{A(w)}}>1.
\]

The inverse composition \(F_w\) is a strict \(2\)-adic contraction, with contraction factor \(2^{-A(w)}\). It therefore has a unique fixed point \(p_w\), which lies in the cylinder \(C_w\subset C_u\). Equivalently, \[
U^{|w|}(p_w)=p_w.
\] From the affine identity, \[
p_w=\frac{C_w}{2^{A(w)}-3^{|w|}}.
\] Because \(C_w>0\) and \(2^{A(w)}<3^{|w|}\), this rational number is negative in the ordinary real ordering. \(\square\)

This is the accelerated-map version of the familiar fact that the \(2\)-adic Collatz system contains abundant periodic points. For the shortcut map, Bernstein and Lagarias proved a topological conjugacy with the one-sided \(2\)-adic shift. [\hypertarget{cit:R11:2}{}\hyperlink{ref:R11:2}{2}, \hypertarget{cit:R11:3}{}\hyperlink{ref:R11:3}{3}]


\paragraph*{Theorem 5: nowhere-local-boundedness theorem}


Suppose \[
\phi:2\mathbb N+1\longrightarrow\mathbb R
\] satisfies \[
\log U(n)+\phi(U(n))
<
\log n+\phi(n)
\qquad(n>1).
\] Then for every nonempty \(2\)-adic open set \(O\), \[
\sup_{x,y\in O\cap(2\mathbb N+1)}
|\phi(x)-\phi(y)|=\infty.
\]


\paragraph*{Proof}


Choose an expanding periodic point \(p=p_w\in O\) as in Lemma 4. Write \[
q=|w|,
\qquad
A=A(w),
\qquad
\lambda=\frac{3^q}{2^A}>1.
\] Choose \(K\) so that \[
p+2^K\mathbb Z_2\subset O.
\]

For each \(N\geq1\), choose a sufficiently large positive odd integer \(x_N\) which

\begin{itemize}
\item realizes the repeated valuation word \(w^N\), and
\item satisfies \[
  x_N-p\in2^{AN+K}\mathbb Z_2.
  \]
\end{itemize}

Such integers exist because the corresponding \(2\)-adic residue class contains arbitrarily large positive representatives.

Put \[
y_N=U^{qN}(x_N).
\] On the repeated affine branch, \[
y_N-p=\lambda^N(x_N-p).
\] Since division by \(2^{AN}\) removes exactly \(AN\) powers of \(2\), \[
y_N-p\in2^K\mathbb Z_2.
\] Thus \[
x_N,y_N\in O.
\]

In the ordinary real ordering, \(p<0\). Therefore \[
y_N
=
p+\lambda^N(x_N-p)
=
\lambda^N x_N+(\lambda^N-1)(-p)
>
\lambda^N x_N.
\] Consequently, \[
\log\frac{y_N}{x_N}>N\log\lambda.
\]

Strict descent at every intermediate step gives \[
\log y_N+\phi(y_N)
<
\log x_N+\phi(x_N).
\] Hence \[
\phi(x_N)-\phi(y_N)
>
\log\frac{y_N}{x_N}
>
N\log\lambda.
\] Letting \(N\to\infty\) proves infinite oscillation on \(O\). \(\square\)


\paragraph*{Consequences}


A successful correction cannot be

\begin{itemize}
\item periodic modulo a fixed modulus;
\item locally constant;
\item continuous on the profinite completion;
\item locally bounded in the \(2\)-adic topology;
\item a locally uniformly convergent limit of finite-level corrections;
\item a section of any ordinary sheaf whose sections are locally bounded real functions.
\end{itemize}

The required singularities cannot be confined to \(-1\). Expanding periodic points are dense, so the correction must be singular, in the sense of unbounded oscillation, in every congruence neighborhood.


\bigskip



\subsubsection*{6. The correct cohomology is dynamical\textemdash{}and it does not vanish}


The asymptotic logarithmic drift is the function \[
g(x)=\log3-a(x)\log2.
\] For a positive integer, \[
\log U(n)-\log n
=
g(n)+\log\left(1+\frac1{3n}\right).
\] The final term is positive and tends to zero.

A correction \(\phi\) changes the asymptotic cocycle by a coboundary: \[
g\longmapsto g+\phi\circ U-\phi.
\] On a periodic orbit \(p,U(p),\ldots,U^{q-1}(p)\), the coboundary telescopes: \[
\sum_{j=0}^{q-1}
\bigl(\phi(U^{j+1}p)-\phi(U^jp)\bigr)=0.
\] Thus \[
\sum_{j=0}^{q-1}
\bigl(g+\phi\circ U-\phi\bigr)(U^jp)
=
q\log3-A(w)\log2.
\] Periodic orbit sums are standard obstructions to solving cohomological equations; this is the basic obstruction underlying Liv\v{s}ic theory. [\hypertarget{cit:R11:4}{}\hyperlink{ref:R11:4}{4}, \hypertarget{cit:R11:5}{}\hyperlink{ref:R11:5}{5}]

For the fixed point \[
p=-1,
\qquad a(p)=1,
\qquad U(p)=-1,
\] the periodic sum is \[
g(-1)=\log\frac32>0.
\] Therefore \(g\) cannot be cohomologous, by any continuous correction, to a strictly negative cocycle.

This obstruction is present whether or not all positive integers reach \(1\).

There is a sharp contrast with the average behavior. Under normalized Haar measure on the odd \(2\)-adic integers, \[
\mathbb P(a=k)=2^{-k},
\qquad k\geq1,
\] and hence \[
\mathbb E[a]=\sum_{k\geq1}\frac{k}{2^k}=2.
\] Thus \[
\mathbb E[g]
=
\log3-2\log2
=
\log\frac34<0.
\] So the local drift is negative on average, while the invariant point mass at \(-1\) has positive drift. This is exactly why ``most residue classes descend'' is not enough to produce a uniform Lyapunov function.


\bigskip



\subsubsection*{7. The finite-state obstruction is an ordered cohomology class}


At a finite congruence level, let \(G=(V,E)\) be the directed transition graph and let \(w(e)\) be the logarithmic edge weight. A residue potential is a \(0\)-cochain \[
\phi:V\to\mathbb R,
\] and it modifies the edge weights to \[
w_\phi(e)
=
w(e)+\phi(t(e))-\phi(s(e)).
\]

The exact finite-graph duality is \[
\boxed{
\inf_{\phi:V\to\mathbb R}
\max_{e\in E}
\bigl(w(e)+\phi(t(e))-\phi(s(e))\bigr)
=
\max_C\frac1{|C|}\sum_{e\in C}w(e),
}
\] where \(C\) ranges over directed cycles.

The necessity follows by summing around a cycle. For sufficiency, choose any number \(\alpha\) strictly larger than the maximum cycle mean. The difference constraints \[
\phi(t)-\phi(s)\leq\alpha-w(e)
\] have no negative-cost directed cycle, so shortest-path potentials solve all the inequalities.

Thus the obstruction is not ordinary \(H^1\)-vanishing. It is the \textbf{positive cycle cone}, or equivalently the maximum cycle mean.

At every modulus \(M\), the residue \(-1\bmod M\) carries a positive self-transition realized by \[
4Mk-1\longmapsto6Mk-1.
\] Its asymptotic weight is \[
\log\frac32.
\] Therefore every finite transition graph has maximum cycle mean at least \(\log(3/2)\). Refinement does not remove this class.

This gives the precise dichotomy:

\begin{itemize}
\item Ordinary sheaf \(H^1\) vanishes automatically and is irrelevant.
\item Dynamical or graph \(H^1\) is relevant, but has persistent positive classes even without positive counterexamples.
\end{itemize}


\bigskip



\subsubsection*{8. Unrestricted gluing merely restates Collatz}


Suppose one replaces continuous or locally bounded functions by arbitrary functions. Then all local functions glue pointwise; the sheaf of arbitrary real-valued functions is flabby, so its higher cohomology vanishes.

But existence of the desired global section is then equivalent to Collatz itself.


\paragraph*{Proposition 6}


The Collatz conjecture is equivalent to the existence of a function \[
h:2\mathbb N+1\longrightarrow\mathbb N_0
\] such that \[
h(1)=0,
\qquad
h(U(n))=h(n)-1
\quad(n>1).
\]


\paragraph*{Proof}


If Collatz is true, define \[
h(n)=\min\{j\geq0:U^j(n)=1\}.
\] Then the displayed identity is immediate.

Conversely, if such an \(h\) exists, repeated application decreases a nonnegative integer by \(1\). After exactly \(h(n)\) steps the orbit must reach \(1\). \(\square\)

Setting \[
\phi(n)=h(n)-\log n
\] gives \[
V(n)=\log n+\phi(n)=h(n),
\] a perfect Lyapunov function. But \(h\) is exactly the total stopping time whose finiteness is the conjecture.

There is another important issue: a real-valued descent function need not prove termination unless it is bounded below or proper. On a divergent ray \[
n_0\to n_1\to n_2\to\cdots
\] one can simply put \[
V(n_j)=-j.
\] It strictly decreases forever.

A divergent trajectory is topologically a ray, whose ordinary first cohomology is zero. Hence ordinary \(H^1\) is naturally suited to detecting cycles, not failure of well-foundedness at infinity. Excluding divergent trajectories requires a coercivity or properness theorem, not merely a gluing theorem.


\bigskip



\subsubsection*{9. What a viable repaired program would require}


A sheaf-flavored approach could survive only after three major changes.


\paragraph*{9.1 Add the archimedean height}


Pure congruence topology cannot distinguish a large positive integer shadowing \(-1\) from the negative \(2\)-adic point itself. Positivity is not a \(2\)-adic local property.

The local state must therefore include something like \[
\bigl(\lfloor\log n\rfloor,\ n\bmod2^r,\ U^j(n)\bmod3^r\bigr),
\] not merely \(n\bmod M\).


\paragraph*{9.2 Replace ordinary sheaves by filtered dynamical cochains}


The correction must be allowed to acquire increasingly large oscillation near every expanding \(2\)-adic periodic cylinder. Theorem 5 shows that the polar or singular set is dense.

One would need finite-level corrections \(\phi_r\) for which

\[
\operatorname{osc}_{C_w}\phi_r
\gtrsim
N\left(q\log3-A(w)\log2\right)
\] whenever the level resolves \(N\) repetitions of an expanding word \(w\).

Thus no compactness theorem in the uniform norm can apply.


\paragraph*{9.3 Prove coercive convergence on positive integers}


The actual missing theorem would have to say that, despite the forced dense singularities,

\begin{enumerate}
\item the renormalized values \(\phi_r(n)\) converge for every positive integer \(n\);
\item the resulting \[
   V(n)=\log n+\phi(n)
   \] has finite sublevel sets;
\item \(V(U(n))<V(n)\) for \(n>1\).
\end{enumerate}

The difficult point is the second condition. Sheaf \(H^1\)-vanishing addresses none of it.

Your terminal-suffix locality identity suggests a possible hybrid way to control the singularities. If a valuation word decomposes as \(w=uv\), with \(|v|=r\), \(B=A(v)\), and endpoint \(X\), then \[
X\equiv2^{-B}C_v\pmod{3^r},
\] and hence \[
3^r\mid X+1
\quad\Longleftrightarrow\quad
3^r\mid C_v+2^B.
\] Thus the \(3\)-adic endpoint obstruction attached to a long \(2\)-adic shadow can be read from its terminal valuation suffix.

A potentially viable replacement for the proposed theorem would therefore be an \textbf{adelic pole-avoidance theorem}:


\begin{quote}
The singular corrections forced by all expanding \(2\)-adic periodic cylinders can be normalized so that their total contribution is finite and coercive on ordinary positive integers, because the required \(3\)-adic endpoint conditions cannot remain compatible indefinitely.

\end{quote}


That would be a genuinely new theorem. Ordinary \(H^1\)-vanishing does not imply it, and the dense-periodic-point argument shows why it would have to use the archimedean and \(3\)-adic information essentially.


\bigskip



\subsubsection*{Final assessment}


The proposed proof cannot be completed in its stated form.

The exact conclusions are:

\begin{enumerate}
\item \textbf{Every fixed-modulus potential fails}, even with bounded class-dependent lookahead.
\item \textbf{Ordinary sheaf \(H^1\) vanishes automatically} on the profinite congruence space, independently of Collatz.
\item \textbf{The relevant dynamical cohomology does not vanish}: the fixed point \(-1\) and dense expanding \(2\)-adic periodic points give persistent positive cycle classes.
\item \textbf{Any global strict correction must be nowhere locally bounded} in the \(2\)-adic topology.
\item \textbf{Allowing arbitrary singular functions makes the construction equivalent to defining the total stopping time}, hence merely restates the conjecture.
\item \textbf{A divergent orbit need not create an \(H^1\) class at all}; the missing obstruction is properness at infinity.
\end{enumerate}

So approach 6 should be reclassified as:


\begin{quote}
\textbf{No-go as a purely modular sheaf argument; potentially viable only as a singular, filtered, adelic dynamical-cohomology program.}

\end{quote}


The main new mathematical output is Theorem 5: every global correction producing pointwise Collatz descent must have unbounded oscillation in every \(2\)-adic congruence neighborhood. This identifies precisely why regular modular Lyapunov functions cannot glue.
\subsection*{References}
\begin{enumerate}[leftmargin=2.2em]
\item \hypertarget{ref:R11:1}{}\hyperlink{cit:R11:1}{\texttt{turn647371view1}}
\item \hypertarget{ref:R11:2}{}\hyperlink{cit:R11:2}{\texttt{turn647371view0}}
\item \hypertarget{ref:R11:3}{}\hyperlink{cit:R11:3}{\texttt{turn895592view0}}
\item \hypertarget{ref:R11:4}{}\hyperlink{cit:R11:4}{\texttt{turn647371view2}}
\item \hypertarget{ref:R11:5}{}\hyperlink{cit:R11:5}{\texttt{turn895592view1}}
\end{enumerate}

\clearpage
\section{Round 12: approach 7 --- Compute the K-theory of the Collatz orbit groupoid}
\subsection*{Prompt}
\begin{quote}
Complete a full proof of the Collatz conjecture using the following approach:

\begin{enumerate}
\item Compute the K-theory of the Collatz orbit groupoid
\end{enumerate}

\textbf{Speculative \textperiodcentered{} B}

Construct the etale groupoid generated by forward and inverse Collatz branches, with a boundary recording ends of the inverse tree. A nontrivial cycle or divergent component may create a K\_0, K\_1, or index class absent from the trivial component.

Approximate the groupoid by finite modular graphs and compute the induced maps on K-groups. Stabilization to the trivial class would prove connectivity; a persistent class would identify a structured obstruction or candidate counterexample.

Take as long as you need.
\end{quote}
\subsection*{Response}

\subsection*{Result of the K-theoretic attack}


The natural Collatz orbit groupoid can be constructed and its \(K\)-theory can be computed completely. The computation does exactly what the proposal hopes at the \textbf{infinite discrete level}:

\[
\boxed{\text{Collatz is true}\quad\Longleftrightarrow\quad
K_{0}\!\left(C^{*}(\mathcal G_T)\right)\cong \mathbb Z.}
\]

More precisely, every divergent orbit component contributes a new \(\mathbb Z\)-summand to \(K_{0}\), and every nontrivial periodic component contributes new \(\mathbb Z\)-summands to both \(K_{0}\) and \(K_{1}\).

However, the suggested final step\textemdash{}deducing this computation from finite modular graphs or from the inverse-tree boundary\textemdash{}does not work. In fact:

\begin{enumerate}
\item the natural modular orbit graph is connected for \textbf{every modulus \(M\)}, unconditionally;
\item the \(2\)-adic inverse-tree boundary gives the full binary-shift groupoid, whose algebra is \(\mathcal O_{2}\) and whose \(K\)-groups are both zero, unconditionally;
\item these same boundary \(K\)-groups occur for the \(5x+1\) map, which has an explicit nontrivial positive cycle.
\end{enumerate}

Thus the program yields a rigorous \(K\)-theoretic reformulation and a precise no-go theorem for the proposed approximation argument, but not a proof of Collatz.


\bigskip



\subsubsection*{1. The discrete Collatz orbit groupoid}


It is convenient to use the equivalent shortcut map

\[
T(n)=
\begin{cases}
n/2,&n\equiv0\pmod 2,\\[2mm]
(3n+1)/2,&n\equiv1\pmod 2.
\end{cases}
\]

The usual Collatz conjecture is equivalent to the assertion that every positive integer reaches the cycle

\[
1\longmapsto2\longmapsto1.
\]

Define

\[
\mathcal G_T
=
\left\{
(x,k-\ell,y):
x,y\in\mathbb N,\quad k,\ell\geq0,\quad
T^{k}(x)=T^{\ell}(y)
\right\}.
\]

The operations are

\[
(x,a,y)(y,b,z)=(x,a+b,z),
\qquad
(x,a,y)^{-1}=(y,-a,x).
\]

Giving \(\mathbb N\), and hence \(\mathcal G_T\), the discrete topology makes \(\mathcal G_T\) a countable Hausdorff \'etale groupoid. It is the discrete special case of a Deaconu\textendash{}Renault groupoid. General computations for such groupoids also naturally produce kernel-and-cokernel descriptions in degrees \(0\) and \(1\). [\hypertarget{cit:R12:1}{}\hyperlink{ref:R12:1}{1}]

Define orbit equivalence by

\[
x\sim y
\quad\Longleftrightarrow\quad
T^{k}(x)=T^{\ell}(y)
\quad\text{for some }k,\ell\geq0.
\]

Let

\[
\Omega=\mathbb N/{\sim}
\]

be the set of orbit components, and let \(\Omega_{\mathrm{per}}\subseteq\Omega\) be the components containing a periodic orbit.


\bigskip



\subsubsection*{2. Complete computation of \(K_{0}\) and \(K_{1}\)}



\paragraph*{Theorem 1}


There are natural isomorphisms

\[
\boxed{
K_{0}\!\left(C^{*}(\mathcal G_T)\right)
\cong
\bigoplus_{\omega\in\Omega}\mathbb Z
}
\]

and

\[
\boxed{
K_{1}\!\left(C^{*}(\mathcal G_T)\right)
\cong
\bigoplus_{\omega\in\Omega_{\mathrm{per}}}\mathbb Z.
}
\]


\paragraph*{Proof}


Fix an orbit component \(\omega\in\Omega\). The reduction

\[
\mathcal G_T|_{\omega}
\]

is a transitive discrete groupoid.

Choose \(x_{\omega}\in\omega\). Its isotropy group is

\[
H_{\omega}
=
\left\{
a\in\mathbb Z:(x_{\omega},a,x_{\omega})\in\mathcal G_T
\right\}.
\]

There are two cases.


\subparagraph*{Aperiodic component}


Suppose \(\omega\) contains no periodic orbit. If

\[
T^{k}(x_{\omega})=T^{\ell}(x_{\omega}),
\]

then \(k=\ell\); otherwise the forward orbit would be eventually periodic. Hence

\[
H_{\omega}=\{0\}.
\]

Choosing one groupoid arrow from \(x_{\omega}\) to every \(x\in\omega\) identifies the transitive groupoid with the pair groupoid

\[
\omega\times\omega.
\]

Consequently,

\[
C^{*}(\mathcal G_T|_{\omega})
\cong
\mathcal K(\ell^{2}(\omega)),
\]

and therefore

\[
K_{0}\cong\mathbb Z,\qquad K_{1}=0.
\]


\subparagraph*{Periodic component}


Suppose \(\omega\) contains a periodic orbit of least period \(p\). Every point in \(\omega\) eventually enters that cycle, and

\[
H_{\omega}=p\mathbb Z\cong\mathbb Z.
\]

The same choice of connecting arrows gives

\[
\mathcal G_T|_{\omega}
\cong
\omega\times p\mathbb Z\times\omega.
\]

Hence

\[
C^{*}(\mathcal G_T|_{\omega})
\cong
\mathcal K(\ell^{2}(\omega))
\otimes C^{*}(\mathbb Z)
\cong
\mathcal K(\ell^{2}(\omega))\otimes C(\mathbb T).
\]

Thus

\[
K_{0}\cong\mathbb Z,\qquad K_{1}\cong\mathbb Z.
\]

The unit space is discrete, so the orbit components are open invariant subsets and

\[
C^{*}(\mathcal G_T)
\cong
\bigoplus_{\omega\in\Omega}^{c_{0}}
C^{*}(\mathcal G_T|_{\omega}).
\]

Taking \(K\)-theory proves the formulas. \ensuremath{\blacksquare}


\bigskip



\subsubsection*{3. The concrete kernel-and-cokernel presentation}


Let

\[
F=\bigoplus_{n\geq1}\mathbb Z e_n
\]

and define

\[
\partial:F\longrightarrow F,
\qquad
\partial e_n=e_n-e_{T(n)}.
\]

Equivalently, \(\partial=I-A^{t}\), where \(A\) is the adjacency matrix of the directed graph with one edge

\[
n\longrightarrow T(n)
\]

from every vertex.

Then

\[
K_{0}\!\left(C^{*}(\mathcal G_T)\right)
\cong\operatorname{coker}\partial,
\qquad
K_{1}\!\left(C^{*}(\mathcal G_T)\right)
\cong\ker\partial.
\]

These groups can be read directly.


\paragraph*{Cokernel}


The relations in the cokernel are

\[
[e_n]=[e_{T(n)}].
\]

Thus all vertices in one orbit component represent the same class. There are no relations between distinct components. Explicitly, the map

\[
F\longrightarrow\bigoplus_{\omega\in\Omega}\mathbb Z u_{\omega},
\qquad
e_n\longmapsto u_{[n]},
\]

has kernel precisely \(\operatorname{im}\partial\). Therefore

\[
\operatorname{coker}\partial
\cong
\bigoplus_{\omega\in\Omega}\mathbb Z u_\omega.
\]


\paragraph*{Kernel}


An element of \(\ker\partial\) is a finitely supported integral circulation on the functional graph. A finite circulation decomposes into directed cycles. Since a functional graph has at most one directed cycle in each component,

\[
\ker\partial
=
\bigoplus_{C}
\mathbb Z
\left(\sum_{v\in C}e_v\right),
\]

where \(C\) ranges over directed cycles.

In particular:

\begin{itemize}
\item a divergent component contributes one generator to \(K_{0}\) and nothing to \(K_{1}\);
\item a periodic component contributes one generator to both \(K_{0}\) and \(K_{1}\).
\end{itemize}


\bigskip



\subsubsection*{4. Exact equivalence with Collatz}


Let \(\omega_{0}\) denote the component containing \(1\) and \(2\).


\paragraph*{Corollary 2}


The following are equivalent:

\[
\begin{aligned}
&\text{every positive integer reaches }1;\\
&\Omega=\{\omega_{0}\};\\
&K_{0}\!\left(C^{*}(\mathcal G_T)\right)\cong\mathbb Z;\\
&K_{0}\!\left(C^{*}(\mathcal G_T)\right)/
   \mathbb Z[p_{1}]=0.
\end{aligned}
\]


\paragraph*{Proof}


If \(\Omega=\{\omega_{0}\}\), then for every \(n\) there are \(k,\ell\) such that

\[
T^{k}(n)=T^{\ell}(1)\in\{1,2\}.
\]

Hence the orbit of \(n\) reaches \(1\).

Conversely, if every orbit reaches \(1\), all vertices belong to \(\omega_{0}\).

The \(K_{0}\) equivalences follow from Theorem 1. \ensuremath{\blacksquare}

Thus the proposed obstruction really exists:

\[
\widetilde K_{0}
:=
K_{0}\!\left(C^{*}(\mathcal G_T)\right)/
\mathbb Z[p_{1}]
\cong
\bigoplus_{\omega\neq\omega_{0}}\mathbb Z.
\]

A counterexample component gives a nonzero class in \(\widetilde K_{0}\).

However, computing that this group vanishes is exactly equivalent to proving that there is only one orbit component. No information has yet been gained unless \(K_{0}\) can be computed by an independent approximation.


\bigskip



\subsection*{5. The inverse-tree boundary loses the arithmetic information}


Extend \(T\) to the \(2\)-adic integers:

\[
T:\mathbb Z_{2}\longrightarrow\mathbb Z_{2}.
\]

Each \(y\in\mathbb Z_{2}\) has exactly two inverse images,

\[
g_{0}(y)=2y,
\qquad
g_{1}(y)=\frac{2y-1}{3}.
\]

The second expression is always defined in \(\mathbb Z_{2}\), because \(3\) is a \(2\)-adic unit. Thus the completed inverse tree is the full binary tree.

For \(x\in\mathbb Z_{2}\), define its parity sequence

\[
Q(x)=\bigl(x\bmod2,T(x)\bmod2,T^{2}(x)\bmod2,\ldots\bigr).
\]

The parity map is a homeomorphism

\[
Q:\mathbb Z_{2}\longrightarrow\{0,1\}^{\mathbb N}
\]

and satisfies

\[
Q\circ T=\sigma\circ Q,
\]

where \(\sigma\) is the one-sided binary shift. In particular, every infinite binary sequence is the parity sequence of exactly one \(2\)-adic integer. [\hypertarget{cit:R12:2}{}\hyperlink{ref:R12:2}{2}, \hypertarget{cit:R12:3}{}\hyperlink{ref:R12:3}{3}]

Consequently, the Deaconu\textendash{}Renault groupoid of the \(2\)-adic Collatz map is the full two-shift groupoid, whose \(C^{*}\)-algebra is

\[
\mathcal O_{2}.
\]

For the one-vertex graph with two loops, the graph \(K\)-theory matrix is

\[
I-A^{t}=[1]-[2]=[-1].
\]

It is invertible over \(\mathbb Z\), so

\[
\boxed{
K_{0}(\mathcal O_{2})=0,
\qquad
K_{1}(\mathcal O_{2})=0.
}
\]

This computation holds independently of what positive integer orbits do.

The reason is precise: for an ordinary integer \(y\), the odd inverse branch

\[
\frac{2y-1}{3}
\]

is an integer only when \(y\equiv2\pmod3\). Passing to \(\mathbb Z_{2}\) makes this branch available for every \(y\), erasing the ordinary divisibility-by-\(3\) condition.


\bigskip



\subsubsection*{6. A concrete counterexample to any boundary-\(K\)-theory implication}


For any odd \(q\), consider

\[
T_q(n)=
\begin{cases}
n/2,&n\text{ even},\\
(qn+1)/2,&n\text{ odd}.
\end{cases}
\]

Exactly the same parity-prefix argument shows that \(T_q\) on \(\mathbb Z_{2}\) is conjugate to the full binary shift. Hence its boundary groupoid algebra is also \(\mathcal O_{2}\), with

\[
K_{0}=K_{1}=0.
\]

But for \(q=5\), there are at least two distinct positive cycles:

\[
1\to3\to8\to4\to2\to1
\]

and

\[
13\to33\to83\to208\to104\to52\to26\to13.
\]

Therefore the implication

\[
K_{*}(\text{inverse-tree boundary})=K_{*}(\mathcal O_{2})
\quad\Longrightarrow\quad
\text{unique positive orbit component}
\]

is false even within the \(qn+1\) family.

This same distinction appears in Mori's operator-algebraic formulation: the concrete isometries generate an abstract algebra isomorphic to \(\mathcal O_{2}\) unconditionally, while cyclicity or irreducibility of the particular representation carries the Collatz information. The abstract algebra\textemdash{}and hence its ordinary \(K\)-theory\textemdash{}does not. [\hypertarget{cit:R12:4}{}\hyperlink{ref:R12:4}{4}]


\bigskip



\subsection*{7. Finite modular graphs are always connected}


There is an even stronger obstruction to the proposed modular approximation.

For \(M\geq1\), define the undirected modular Collatz graph \(\Gamma_M\) by

\[
\Gamma_M^{0}=\mathbb Z/M\mathbb Z,
\]

and put an edge between

\[
\overline n
\quad\text{and}\quad
\overline{T(n)}
\]

for every positive integer \(n\). This is the natural finite graph obtained after adjoining forward and inverse Collatz edges and reducing vertices modulo \(M\).


\paragraph*{Theorem 3: Modular collapse}


For every \(M\geq1\),

\[
\boxed{\Gamma_M\text{ is connected}.}
\]


\paragraph*{Proof}


We first use a lifting observation. If \(d\mid M\), reduction gives a map

\[
\pi:\Gamma_M\longrightarrow\Gamma_d.
\]

Every edge downstairs is the image of an edge upstairs. Therefore, if \(\Gamma_d\) is connected and every fiber of \(\pi\) lies in one connected component of \(\Gamma_M\), then \(\Gamma_M\) is connected.

We remove factors of \(2\) and \(3\).


\paragraph*{Step 1: removing a factor of \(2\)}


Suppose \(M=2d\). The two elements in a fiber of

\[
\mathbb Z/M\mathbb Z\longrightarrow\mathbb Z/d\mathbb Z
\]

are \(x\) and \(x+d\).

Choose a positive integer \(y\equiv x\pmod M\). The two even integers

\[
2y,\qquad 2y+M
\]

are congruent modulo \(M\), while

\[
T(2y)=y,
\qquad
T(2y+M)=y+\frac M2=y+d.
\]

Thus \(x\) and \(x+d\) have a common neighbor in \(\Gamma_M\). Every fiber modulo \(d\) is connected.


\paragraph*{Step 2: removing a factor of \(3\)}


Suppose \(M=3d\) is odd.

For every residue \(x\pmod M\), choose an odd positive integer \(n\equiv x\pmod M\). Then

\[
x
\sim
\frac{3x+1}{2}
\sim
3x+1
\]

in \(\Gamma_M\), where the second edge comes from halving the even number \(3n+1\).

The three elements in a fiber modulo \(d\) are

\[
x,\qquad x+d,\qquad x+2d.
\]

But

\[
3(x+jd)+1
=
3x+jM+1
\equiv3x+1\pmod M.
\]

Hence all three fiber elements connect to the same residue \(3x+1\). Every fiber modulo \(d\) is connected.


\paragraph*{Step 3: the case \(\gcd(M,6)=1\)}


We have already shown

\[
x\sim3x+1.
\]

Since \(2\) is invertible modulo \(M\), choose \(u\) satisfying

\[
2u=x.
\]

The halving edge gives \(u\sim x\). Also

\[
u\sim3u+1
\]

and doubling gives

\[
3u+1
\sim
2(3u+1)
=
3x+2.
\]

Therefore

\[
x\sim3x+2.
\]

It follows that

\[
3x+1\sim3x+2.
\]

Since \(3\) is invertible modulo \(M\), every residue \(a\) can be written as

\[
a=3x+1.
\]

Thus

\[
a\sim a+1
\]

for every \(a\). Therefore \(\Gamma_M\) is connected.

Starting from the coprime-to-\(6\) case and successively restoring the removed factors of \(3\) and \(2\), the fiber argument proves connectivity for every \(M\). \ensuremath{\blacksquare}


\bigskip



\subsubsection*{8. Consequence for modular \(K\)-theory}


If one forms the finite equivalence-relation groupoid generated by \(\Gamma_M\), Theorem 3 says it is the full pair groupoid

\[
(\mathbb Z/M\mathbb Z)\times(\mathbb Z/M\mathbb Z).
\]

Therefore its algebra is

\[
M_M(\mathbb C),
\]

and

\[
K_{0}\cong\mathbb Z,\qquad K_{1}=0
\]

for every \(M\).

So the hoped-for stabilization

\[
K_{0}(\Gamma_M)\longrightarrow\mathbb Z
\]

occurs at every level without using Collatz. Every genuine orbit component is merged with every other component modulo every \(M\).

If one retains path lag and orientation instead, another problem occurs: \(T\) is not a map on \(\mathbb Z/M\mathbb Z\). For example, integers congruent modulo \(M\) can have different images modulo \(M\), because division by \(2\) requires one additional binary digit. Thus there is no canonical finite dynamical quotient

\[
(\mathbb N,T)\longrightarrow(\mathbb Z/M\mathbb Z,T_M)
\]

and hence no canonical system of groupoid \(C^{*}\)-homomorphisms inducing the proposed \(K\)-theory maps. Replacing the map by a multivalued relation introduces modular cycles that need not lift to integer cycles.


\bigskip



\subsection*{9. The canonical \(2\)-power approximants also have trivial \(K\)-theory}


The exact symbolic approximants are the binary de Bruijn graphs.

Let \(D_r\) have as vertices all binary words of length \(r\), with edges

\[
b_0b_1\cdots b_{r-1}
\longrightarrow
b_1\cdots b_{r-1}c,
\qquad c\in\{0,1\}.
\]

Finite parity words are in bijection with residues modulo \(2^r\), so these are the genuine finite \(2\)-adic approximants.

Let \(A_r\) be the adjacency matrix of \(D_r\). There are rectangular incidence matrices \(R_r,S_r\) such that

\[
A_{r+1}=R_rS_r,
\qquad
A_r=S_rR_r.
\]

Sylvester's determinant identity gives

\[
\det(I-A_{r+1})
=
\det(I-A_r).
\]

At level \(0\), \(A_0=[2]\), so

\[
\det(I-A_0)=-1.
\]

Therefore

\[
\det(I-A_r)=-1
\]

for every \(r\). The matrix \(I-A_r^{t}\) is unimodular, and hence

\[
\boxed{
K_{0}(C^{*}(D_r))=0,\qquad
K_{1}(C^{*}(D_r))=0
}
\]

at every finite level.

Their limit is \(\mathcal O_{2}\). Once again, stabilization is automatic and contains no information about positivity.


\bigskip



\subsection*{10. Why a divergent class can disappear at every finite stage}


The basic telescoping identity is

\[
e_n-e_{T^{k}(n)}
=
\sum_{j=0}^{k-1}
\left(e_{T^{j}(n)}-e_{T^{j+1}(n)}\right)
\in\operatorname{im}\partial.
\]

Thus \(e_n\) is \(K_{0}\)-equivalent to every later point of its orbit.

There is an even more severe escape mechanism in the inverse direction. For every \(r\),

\[
T^{r}(2^{r}n)=n,
\]

so

\[
[e_n]=[e_{2^{r}n}].
\]

Every component therefore contains arbitrarily large vertices, including periodic components. Any finite height truncation that collapses all sufficiently large vertices into a common boundary class automatically merges \textbf{all} orbit components.

To prevent this, one would need a distinct boundary object for each component or forward end. But for a functional graph:

\begin{itemize}
\item each divergent component has one forward end;
\item each periodic component has a cycle;
\item determining whether two vertices determine the same forward end is exactly determining whether their forward orbits eventually meet.
\end{itemize}

Thus constructing the required faithful boundary already requires the orbit decomposition that Collatz asks us to prove.

There is also substantial unconditional inverse-tree topology inside the known component. For example, for every \(j\geq0\),

\[
y=2^{2j+1}
\]

has two positive predecessors,

\[
2y
\qquad\text{and}\qquad
\frac{2y-1}{3}
=
\frac{2^{2j+2}-1}{3}.
\]

Both lie in the component of \(1\). Hence the trivial component already contains infinitely many branching vertices and many inverse ends.

Recent work on \(C^{*}\)-algebras of undirected locally finite graphs makes this background problem explicit: their \(K\)-theory records genus, branching, ends, and dead ends; even a full binary tree has a countably infinite graph-valency contribution. These are not signatures of a Collatz counterexample. [\hypertarget{cit:R12:5}{}\hyperlink{ref:R12:5}{5}]


\bigskip



\subsection*{11. Final status of the approach}


The part that succeeds is exact:

\[
\boxed{
\widetilde K_{0}\!\left(C^{*}(\mathcal G_T)\right)
=
K_{0}\!\left(C^{*}(\mathcal G_T)\right)/
\mathbb Z[p_1]
\cong
\bigoplus_{\substack{\text{nontrivial}\\\text{orbit components}}}\mathbb Z.
}
\]

Furthermore,

\[
\boxed{
K_{1}\!\left(C^{*}(\mathcal G_T)\right)
\cong
\bigoplus_{\text{periodic components}}\mathbb Z.
}
\]

So:

\begin{itemize}
\item \(K_{1}\) detects nontrivial cycles;
\item reduced \(K_{0}\) detects both cycles and divergent components;
\item Collatz is equivalent to reduced \(K_{0}=0\).
\end{itemize}

But the proposed route for computing these groups fails:

\[
\begin{array}{c|c|c}
\text{model} & K\text{-theoretic outcome} & \text{problem}\\ \hline
\text{discrete orbit groupoid} &
K_0=\mathbb Z^{(\text{components})} &
\text{exactly equivalent to Collatz}\\
\text{\(2\)-adic inverse boundary} &
K_0=K_1=0 &
\text{forgets ordinary integrality}\\
\text{modular equivalence graphs} &
K_0=\mathbb Z,\ K_1=0 &
\text{connected for every }M\\
\text{parity/de Bruijn graphs} &
K_0=K_1=0 &
\text{full-shift data only}\\
\text{undirected end algebra} &
\text{large end contribution} &
\text{counts inverse ends already present}
\end{array}
\]

A viable revision would have to use \textbf{controlled or filtered \(K\)-theory}, retaining simultaneously:

\[
\text{orbit lag},\qquad
\text{ordinary \(3\)-divisibility},\qquad
\text{positivity},\qquad
\text{archimedean height}.
\]

The decisive missing theorem would be a component-separation or no-escape result asserting that every nonzero class in

\[
\widetilde K_{0}\!\left(C^{*}(\mathcal G_T)\right)
\]

survives in some finite controlled quotient. Pure residue quotients do not have this property: Theorem 3 shows that they annihilate every component distinction at every modulus.

Accordingly, ordinary \(K_{0}\), \(K_{1}\), and inverse-boundary index classes do not complete a proof of Collatz. The Collatz information lies in the arithmetic embedding of the positive integers\textemdash{}or, equivalently, in the concrete representation and its diagonal/ideal structure\textemdash{}not in the abstract boundary \(K\)-groups.
\subsection*{References}
\begin{enumerate}[leftmargin=2.2em]
\item \hypertarget{ref:R12:1}{}\hyperlink{cit:R12:1}{\texttt{turn305300view2}}
\item \hypertarget{ref:R12:2}{}\hyperlink{cit:R12:2}{\texttt{turn305300view1}}
\item \hypertarget{ref:R12:3}{}\hyperlink{cit:R12:3}{\texttt{turn980539view4}}
\item \hypertarget{ref:R12:4}{}\hyperlink{cit:R12:4}{\texttt{turn980539view2}}
\item \hypertarget{ref:R12:5}{}\hyperlink{cit:R12:5}{\texttt{turn723462view0}}
\end{enumerate}

\clearpage
\section{Round 13: approach 8 --- Thermodynamic formalism with an arithmetic integrality potential}
\subsection*{Prompt}
\begin{quote}
Complete a full proof of the Collatz conjecture using the following approach:

\begin{enumerate}
\item Thermodynamic formalism with an arithmetic integrality potential
\end{enumerate}

\textbf{Hybrid \textperiodcentered{} B}

Place a potential on valuation words,

Phi(w) = archimedean drift + lambda \emph{ integrality penalty + mu } carry penalty.

Study pressure, equilibrium states, and zero-temperature limits for the resulting countable symbolic system.

A proof would show that the only ground state compatible with zero arithmetic penalty is the trivial cycle. A nontrivial equilibrium state concentrated on exactly realizable critical words would provide a statistical blueprint for a counterexample.

Take as long as you need.
\end{quote}
\subsection*{Response}

\subsection*{Outcome}


Carrying this program through does \textbf{not} presently yield a valid proof of the Collatz conjecture. It yields something quite sharp instead:


\begin{quote}
\textbf{No-go theorem.} On the ordinary countable shift of Collatz valuation words, no continuous or summable-variation ``integrality potential'' can distinguish valuation sequences coming from positive integers from general \(2\)-adic valuation sequences. Moreover, uniqueness of an equilibrium or ground state among invariant probability measures cannot by itself exclude a divergent integer orbit.

\end{quote}


The obstruction is exact, not merely technical. The thermodynamic approach can be repaired only by adding a genuinely nonlocal arithmetic lifting condition. The required lifting statement is precisely the unresolved prefix\textendash{}suffix synchronization theorem already isolated by the affine-compression program.

Below is the full derivation.


\bigskip



\subsubsection*{1. Exact symbolic coding}


Use the accelerated map on positive odd integers,

\[
U(x)=\frac{3x+1}{2^{\nu _2(3x+1)}}.
\]

Let

\[
a_j=\nu _2(3x_j+1),\qquad x_{j+1}=U(x_j),
\]

so that every orbit gives a valuation sequence

\[
a=(a_0,a_1,\ldots)\in \Sigma:=\mathbb N_{\ge 1}^{\mathbb N}.
\]

For a finite word

\[
w=(a_0,\ldots,a_{n-1}),
\]

put

\[
A_0=0,\qquad A_j=\sum_{i=0}^{j-1}a_i
\]

and

\[
C_0=0,\qquad C_{j+1}=3C_j+2^{A_j}.
\]

Then

\[
C_n=\sum_{j=0}^{n-1}3^{n-1-j}2^{A_j},
\]

and the exact affine identity is

\[
2^{A_n}U^n(x)=3^n x+C_n. \tag{1}
\]

The constants also satisfy

\[
C_{uv}=3^{|v|}C_u+2^{A(u)}C_v. \tag{2}
\]

These are the basic algebraic identities behind every proposed potential. [\hypertarget{cit:R13:1}{}\hyperlink{ref:R13:1}{1}]

The first crucial arithmetic fact is:


\begin{quote}
\textbf{Finite realization theorem.} Every finite word \(w\in\mathbb N_{\ge1}^{n}\) is realized as the first \(n\) accelerated valuations of infinitely many positive odd integers.

\end{quote}


More precisely, \(w\) determines one odd residue class modulo \(2^{A_n+1}\). It also determines one endpoint class modulo \(3^n\), and every positive odd endpoint in that class has a positive integral inverse ancestor along \(w\). [\hypertarget{cit:R13:2}{}\hyperlink{ref:R13:2}{2}]

Thus the symbolic system of finite valuation words is the \textbf{full countable shift}. Arithmetic positivity does not remove any finite word.

The trivial Collatz cycle is represented by

\[
2^\infty=(2,2,2,\ldots),
\]

because \(U(1)=1\) and \(\nu _2(4)=2\).


\bigskip



\subsubsection*{2. The archimedean pressure can be computed exactly}


Define the one-step archimedean drift

\[
d(a)=\log 3-a\log 2.
\]

Then

\[
S_nd(w)=n\log 3-A_n\log 2
       =\log \frac{3^n}{2^{A_n}}. \tag{3}
\]

I use the convention that pressure maximizes the potential; a nonnegative ``penalty'' is therefore subtracted. Reversing all signs gives the energy-minimizing convention.

For the pure drift potential \(t d\), the \(n\)-step partition function is

\[
\begin{aligned}
Z_n(t)
&=\sum_{a_0,\ldots,a_{n-1}\ge1}
  \exp\left(tn\log3-t\log2\sum_{j=0}^{n-1}a_j\right)\\
&=\left(3^t\sum_{a\ge1}2^{-ta}\right)^n\\
&=\left(\frac{3^t}{2^t-1}\right)^n,
\end{aligned}
\]

provided \(t>0\). Hence the Gurevi\v{c} pressure is

\[
\boxed{
P_{\rm drift}(t)=t\log3-\log(2^t-1).
} \tag{4}
\]

Its equilibrium state is the Bernoulli measure with geometric symbol law

\[
\boxed{
p_t(a)=(2^t-1)2^{-ta}.
} \tag{5}
\]

As \(t\to\infty\),

\[
p_t(1)\longrightarrow1,
\]

so the zero-temperature limit for maximizing archimedean drift is

\[
\delta_{1^\infty},
\]

not \(\delta_{2^\infty}\). The sequence \(1^\infty\) represents the negative \(2\)-adic fixed point \(-1\), since

\[
\frac{3(-1)+1}{2}=-1.
\]

If instead one minimizes \(d\), there is no minimizing symbol: \(d(a)\to-\infty\) as \(a\to\infty\). Thus the archimedean term by itself selects either the wrong boundary point or no ground state at all.

There is also a drift-neutral positive-entropy equilibrium. Writing

\[
\alpha=\log_2 3,
\]

we have

\[
P'_{\rm drift}(t)
=
\log3-\frac{\log2}{1-2^{-t}}.
\]

At the unique \(t=t_*\) satisfying

\[
2^{-t_*}=1-\frac1\alpha,
\]

the equilibrium state has

\[
\int a_0\,d\mu_{t_*}=\alpha,
\qquad
\int d(a_0)\,d\mu_{t_*}=0.
\]

So even at exactly critical average drift, the full shift has a nontrivial equilibrium with positive entropy. The drift condition alone leaves exponentially many candidate words.

Countable-alphabet zero-temperature results normally require hypotheses such as finite pressure, bounded or summable variation, summability or coercivity. Escape of mass is a separate issue handled by pressure at infinity. [\hypertarget{cit:R13:3}{}\hyperlink{ref:R13:3}{3}, \hypertarget{cit:R13:4}{}\hyperlink{ref:R13:4}{4}, \hypertarget{cit:R13:5}{}\hyperlink{ref:R13:5}{5}]


\bigskip



\subsubsection*{3. The natural ``integrality cost'' is only a cylinder-density cost}


Because a word \(w\) determines one odd residue class modulo \(2^{A_n+1}\), its conditional \(2\)-adic Haar probability among odd integers is

\[
\mathbb P([w])=2^{-A_n}. \tag{6}
\]

The natural information or integrality cost is therefore

\[
-\log \mathbb P([w])=A_n\log2.
\]

Adding this cost gives the finite-volume partition sum

\[
\begin{aligned}
Z_n^{\rm Haar}(t)
&=
\sum_{|w|=n}
e^{tS_nd(w)}2^{-A_n(w)}\\
&=
\left(
3^t\sum_{a\ge1}2^{-(t+1)a}
\right)^n\\
&=
\left(
\frac{3^t}{2^{t+1}-1}
\right)^n.
\end{aligned}
\]

Thus

\[
\boxed{
P_{\rm Haar}(t)
=t\log3-\log(2^{t+1}-1).
} \tag{7}
\]

In particular,

\[
P_{\rm Haar}(0)=0,
\qquad
P_{\rm Haar}(1)=0.
\]

At \(t=0\), the equilibrium law is

\[
p(a)=2^{-a},
\]

with

\[
\mathbb E[a]=2,\qquad
\mathbb E[d(a)]=\log3-2\log2=\log\frac34<0.
\]

This recovers the familiar stochastic contraction heuristic. But it does not distinguish genuine positive integers from generic \(2\)-adic points: it assigns a density to every cylinder, and every cylinder is arithmetically realizable.


\bigskip



\subsubsection*{4. No continuous arithmetic-integrality potential exists}


Here is the decisive obstruction.


\paragraph*{Theorem 1 \textemdash{} Density obstruction}


Let

\[
\mathcal P
=
\{a(x):x\text{ is a positive odd integer}\}
\subseteq\Sigma
\]

be the set of infinite valuation sequences of positive integers. Then \(\mathcal P\) is dense in \(\Sigma\).

Consequently, if

\[
I:\Sigma\longrightarrow[0,\infty]
\]

is lower semicontinuous and satisfies

\[
I(a(x))=0
\qquad
\text{for every positive odd integer }x,
\]

then

\[
\boxed{I\equiv0.}
\]


\subparagraph*{Proof}


Every nonempty open set in the product topology contains a cylinder

\[
[w]=\{a\in\Sigma:(a_0,\ldots,a_{n-1})=w\}.
\]

Every finite word \(w\) is realized by a positive odd integer. Therefore

\[
[w]\cap\mathcal P\ne\varnothing,
\]

so \(\mathcal P\) is dense.

Suppose \(I(a)>0\). Choose \(c\) with \(0<c<I(a)\). Lower semicontinuity implies that

\[
\{b:I(b)>c\}
\]

is an open neighborhood of \(a\). It must intersect the dense set \(\mathcal P\), contradicting \(I=0\) on \(\mathcal P\). Hence \(I\equiv0\). \ensuremath{\blacksquare}

This applies in particular to continuous, H\"older, Walters and summable-variation potentials.

It also gives an even simpler finite-horizon conclusion:


\begin{quote}
Any finite-word penalty which is zero precisely when the word is realizable by a positive integer is identically zero, because every finite valuation word is realizable.

\end{quote}


So ``integrality'' cannot mean finite realizability.


\paragraph*{What integrality actually means}


For an infinite word, let \(R_n\) be the least nonnegative representative of its seed residue class modulo \(2^{A_n+1}\). Then

\[
R_{n+1}\equiv R_n\pmod{2^{A_n+1}}
\]

and \(R_n\) is nondecreasing.

The infinite word represents a nonnegative ordinary integer precisely when

\[
\boxed{
R_n\text{ is eventually constant}.
} \tag{8}
\]

This is the exact ordinary-integrality criterion. [\hypertarget{cit:R13:6}{}\hyperlink{ref:R13:6}{6}]

But eventual stabilization is a discontinuous tail property. Every cylinder contains both:

\begin{itemize}
\item sequences whose \(R_n\) eventually stabilize at a positive integer;
\item sequences whose \(R_n\) continue changing forever and represent nonintegral \(2\)-adic points.
\end{itemize}

The exact indicator

\[
I_\infty(a)=
\begin{cases}
0,&R_n\text{ is eventually constant},\\
1,&\text{otherwise},
\end{cases}
\]

is therefore nowhere compatible with the usual regularity assumptions of Ruelle\textendash{}Perron\textendash{}Frobenius or zero-temperature theory.

This is the first place where the proposed proof fails.


\bigskip



\subsubsection*{5. The carry term is local, but descent is not}


The terminal \(3\)-adic information is exceptionally well behaved.

If \(w=u\tau\) and \(\tau\) has length \(r\), then every endpoint \(X\) obtained along \(w\) satisfies

\[
X\equiv 2^{-A(\tau)}C_\tau\pmod{3^r}.
\]

Consequently,

\[
\boxed{
\nu _3(X+1)\ge r
\iff
3^r\mid C_\tau+2^{A(\tau)}.
} \tag{9}
\]

Thus the depth-\(r\) endpoint carry is determined entirely by the final \(r\) valuations. [\hypertarget{cit:R13:7}{}\hyperlink{ref:R13:7}{7}]

This suggests a local carry potential. The problem appears when one asks whether a compatible inverse branch gives an actual smaller positive ancestor.

Suppose a word is split as

\[
w=u\tau,
\]

where

\[
|u|=d,\quad A(u)=P,
\qquad
|\tau|=k,\quad A(\tau)=T.
\]

Let \(R\) be the least positive realization, let

\[
Z=U^d(R),
\]

and compare \(\tau\) with another length-\(k\) word \(b\) of exponent sum \(B\). Put

\[
c=T-B.
\]

The terminal compatible fibre is

\[
\mathcal F(\tau,B)
=
\left\{
b:
|b|=k,\ A(b)=B,\
C_b\equiv2^{-c}C_\tau\pmod{3^k}
\right\}. \tag{10}
\]

For \(b\in\mathcal F(\tau,B)\), define the terminal carry centre

\[
h(\tau,b)
=
\frac{C_\tau-2^cC_b}{3^k}. \tag{11}
\]

If \(Y\) is the competing ancestor, then

\[
Z+h(\tau,b)=2^cY.
\]

Define the prefix surplus

\[
\Sigma(u,\tau,c)=Z-2^cR. \tag{12}
\]

Subtracting gives the exact comparison

\[
\boxed{
Y-R=
\frac{\Sigma(u,\tau,c)+h(\tau,b)}{2^c}.
} \tag{13}
\]

The \(3\)-adic compatibility and \(h\) are terminal. The sign determining whether \(Y<R\) contains the full-prefix term \(\Sigma\). This is the precise local\textendash{}global obstruction. [\hypertarget{cit:R13:8}{}\hyperlink{ref:R13:8}{8}]

Therefore a bounded-range carry potential can detect:

\begin{itemize}
\item whether an alternative inverse branch is integral;
\item the associated terminal carry;
\item its asymptotic exponent saving.
\end{itemize}

It cannot, without a new theorem, determine whether that branch is positive and smaller for the \textbf{particular ordinary integer} represented by the entire prefix.

The explicit length-\(42\) strict-barrier word with optimal descent cutoff \(814\) makes this nonlocality quantitative: every proposed competitor with cutoff at most \(813\) fails, while one works at \(814\). [\hypertarget{cit:R13:9}{}\hyperlink{ref:R13:9}{9}]

There is an additional boundary obstruction. The periodic word

\[
(1,1,2)^\infty
\]

is retained by the Bellman optimization at every finite depth, but represents

\[
-\frac{19}{11}
\]

rather than a positive integer. [\hypertarget{cit:R13:10}{}\hyperlink{ref:R13:10}{10}]

So even an infinite word with zero local Bellman/carry defect need not be an ordinary positive point.


\bigskip



\subsubsection*{6. Exact periodic ground states}


Let \(w\) be a word of length \(m\), total exponent \(A\), and affine constant \(C_w\). If \(w^\infty\) represents a periodic point \(x\), then (1) gives

\[
2^A x=3^m x+C_w.
\]

Hence

\[
\boxed{
x=\frac{C_w}{2^A-3^m}.
} \tag{14}
\]

A positive Collatz cycle exists exactly when:

\[
2^A>3^m,
\]

\[
2^A-3^m\mid C_w,
\]

the quotient is a positive odd integer, and the valuations are exact.

The trivial word \(w=(2)\) gives

\[
x=\frac1{4-3}=1.
\]

The spurious periodic word \(w=(1,1,2)\) has \(A=4\), \(m=3\), and \(C_w=19\), giving

\[
x=\frac{19}{16-27}=-\frac{19}{11}.
\]

One can certainly define a periodic arithmetic penalty \(J(w)\) which is zero exactly when all the conditions following (14) hold. But then proving that the only zero-penalty periodic word is \(2^m\) is \textbf{exactly the nontrivial-cycle part of the Collatz conjecture}. Thermodynamic formalism has not simplified that arithmetic classification; it has merely renamed it as uniqueness of a periodic ground state.


\bigskip



\subsubsection*{7. Ground-state uniqueness does not exclude divergence}


There is a second, logically separate problem.


\paragraph*{Theorem 2 \textemdash{} Invariant measures see cycles, not wandering orbits}


For a deterministic map on a countable discrete state space, every invariant probability measure is supported on periodic orbits.


\subparagraph*{Proof}


By Poincar\'e recurrence, almost every point with respect to an invariant probability is recurrent. In a discrete space, recurrence means that the orbit returns exactly to the starting point, hence the point is periodic. \ensuremath{\blacksquare}

Therefore, on the actual positive-integer Collatz system, classifying invariant probability measures can classify cycles. It cannot detect a divergent orbit, because a divergent orbit is wandering and supports no invariant probability measure.

The logical gap is visible in the elementary map

\[
F(0)=0,\qquad F(n)=n+1\quad(n\ge1).
\]

Its only invariant probability is \(\delta_0\), but every orbit beginning at \(n\ge1\) diverges. Thus

\[
\text{unique invariant ground state}
\centernot\Longrightarrow
\text{every orbit reaches it}. \tag{15}
\]

This remains true regardless of how convincingly the ground state is characterized.


\bigskip



\subsubsection*{8. What a divergent Collatz orbit would produce symbolically}


A divergent Collatz orbit does nevertheless produce a useful statistical object.


\paragraph*{Theorem 3 \textemdash{} Critical invariant blueprint}


Suppose \(x_0,x_1,\ldots\) is a divergent accelerated Collatz orbit, and let \(a_j\) be its valuation sequence. Then there exists a shift-invariant probability measure \(\nu\) on \(\Sigma\), obtained as a subsequential cylinder limit of the empirical measures

\[
\nu_N=\frac1N\sum_{j=0}^{N-1}\delta_{\sigma^j a},
\]

such that

\[
\boxed{
\int a_0\,d\nu\le\log_2 3.
} \tag{16}
\]

Equivalently,

\[
\boxed{
\int d(a_0)\,d\nu\ge0.
} \tag{17}
\]


\subparagraph*{Proof}


The exact recurrence gives

\[
x_{j+1}
=
\frac{3x_j+1}{2^{a_j}},
\]

and hence

\[
\log x_{j+1}-\log x_j
=
\log3-a_j\log2
+\log\left(1+\frac1{3x_j}\right).
\]

Summing from \(j=0\) to \(N-1\),

\[
A_N\log2
=
N\log3+\log x_0-\log x_N
+\sum_{j=0}^{N-1}
\log\left(1+\frac1{3x_j}\right). \tag{18}
\]

A nonperiodic integer orbit consists of distinct integers, so divergence implies \(x_j\to\infty\). Therefore

\[
\frac1N
\sum_{j=0}^{N-1}
\log\left(1+\frac1{3x_j}\right)
\longrightarrow0.
\]

Since \(-\log x_N\le0\), equation (18) gives

\[
\limsup_{N\to\infty}\frac{A_N}{N}
\le\log_2 3. \tag{19}
\]

For each fixed finite coordinate block, (19) gives tightness of its empirical distributions. A diagonal subsequence therefore converges on all cylinders to a probability measure \(\nu\). For every bounded cylinder function \(f\),

\[
\int f\circ\sigma\,d\nu_N-\int f\,d\nu_N
=
\frac{f(\sigma^Na)-f(a)}N\longrightarrow0,
\]

so \(\nu\) is shift-invariant.

Applying the convergence to the bounded truncations

\[
\min(a_0,M)
\]

and then letting \(M\to\infty\) gives (16). Equation (17) follows. \ensuremath{\blacksquare}

This is very close to the ``statistical blueprint'' anticipated in the proposed approach.

But it does not solve the lifting problem. The limiting measure may be supported on symbolic or \(2\)-adic boundary sequences which are not codes of ordinary positive integers. Because positive-integer codes are dense but not closed, the property ``comes from an ordinary integer'' disappears under weak or cylinder limits.

Conversely, an abstract critical equilibrium state need not possess any generic point whose nested seed representatives eventually stabilize at one positive integer.

Thus the implications are only

\[
\text{divergent positive orbit}
\Longrightarrow
\text{critical invariant symbolic measure},
\]

but not

\[
\text{critical invariant symbolic measure}
\Longrightarrow
\text{divergent positive orbit}. \tag{20}
\]

The missing reverse implication is an arithmetic projective-lifting theorem.


\bigskip



\subsubsection*{9. The exact theorem that would complete the approach}


Define the arithmetic ground-state set \(\mathcal A\) to consist of infinite valuation words which remain Bellman-optimal against all positive compatible inverse ancestors at every depth.

The desired theorem would be

\[
\boxed{
\mathcal A
\cap
\{a:R_n(a)\text{ is eventually a positive integer}\}
=
\{2^\infty\}.
} \tag{21}
\]

In the local\textendash{}global notation, a usable form is:


\begin{quote}
\textbf{Arithmetic synchronization theorem.}   Let \(a\) be an infinite valuation word whose seed representatives \(R_n\) eventually stabilize at a positive odd integer \(m>1\). Suppose \(m\) is no larger than every positive integer whose orbit coalesces with a later point of the orbit of \(m\). Then there exist a decomposition \(w=u\tau\), an exponent saving \(c\ge1\), and a compatible word \(b\in\mathcal F(\tau,A(\tau)-c)\) such that \[
> Z+h(\tau,b)>0
> \] and \[
> \Sigma(u,\tau,c)+h(\tau,b)<0.
> \tag{22}
> \] Consequently the corresponding inverse ancestor satisfies \[
> 0<Y<m.
> \]

\end{quote}



\paragraph*{Why this proves Collatz}


Assume Collatz is false and choose the least positive odd counterexample \(m\).

Its valuation word has seed representatives eventually equal to \(m\), because once

\[
2^{A_n+1}>m,
\]

the least representative of the seed class is \(m\) itself.

Every positive integer whose orbit coalesces with the orbit of \(m\) must be at least \(m\); otherwise that smaller integer would reach \(1\), forcing the orbit of \(m\) to reach \(1\) as well.

The synchronization theorem would produce a positive coalescing ancestor \(Y<m\), a contradiction. Therefore no least counterexample exists, and every positive integer reaches \(1\).

That conditional proof is complete. The unresolved step is exactly the synchronization theorem.

The existing compression framework reaches the same frontier: a carry-synchronization theorem would handle one recurrent finite subsystem, but a separate global input-synchronization theorem is still needed to couple eventual positive \(2\)-adic stabilization to \(3\)-adic endpoint compressibility. [\hypertarget{cit:R13:11}{}\hyperlink{ref:R13:11}{11}]

Thermodynamic terminology does not remove this step. Statement (21) is precisely ``the only arithmetically liftable Aubry/ground state is \(2^\infty\).'' Proving it is the substantive Collatz theorem.


\bigskip



\subsubsection*{10. Why a discontinuous zero-temperature potential does not fix the problem}


One could define formally

\[
I_\infty(a)
=
\mathbf 1_{\{R_n\text{ does not eventually stabilize}\}}
\]

and

\[
K_\infty(a)
=
\mathbf 1_{\{a\text{ remains Bellman-geodesic forever}\}}.
\]

Then one could build an ``energy''

\[
E(a)
=
-\limsup_{n\to\infty}\frac1nS_nd(a)
+\lambda I_\infty(a)
-\mu K_\infty(a).
\]

But this does not constitute a proof:

\begin{enumerate}
\item \(I_\infty\) and \(K_\infty\) are highly discontinuous tail functions.
\item Standard transfer-operator, pressure and zero-temperature theorems no longer apply.
\item Establishing that the only zero-energy arithmetically integral state is \(2^\infty\) is exactly statement (21).
\end{enumerate}

In other words, one can always construct a formal potential whose ground-state classification is equivalent to Collatz. The difficulty is proving that classification from independently checkable local or spectral properties.


\bigskip



\subsubsection*{11. Audit of a recent thermodynamic-formalism claim}


I also checked a very recent August 2026 preprint using closely related topological and equilibrium-state language. It claims that a broad class of maps

\[
g(n)=
\begin{cases}
g_0(n),&n\text{ odd},\\
n/2,&n\text{ even},
\end{cases}
\]

where \(g_0(n)\) may be \textbf{any even number}, has no divergent orbit. [\hypertarget{cit:R13:12}{}\hyperlink{ref:R13:12}{12}]

That statement has an immediate counterexample:

\[
g(n)=
\begin{cases}
2(n+2),&n\text{ odd},\\
n/2,&n\text{ even}.
\end{cases}
\]

For odd \(n\),

\[
n\longmapsto2(n+2)\longmapsto n+2.
\]

Thus

\[
1\longmapsto6\longmapsto3\longmapsto10
\longmapsto5\longmapsto14\longmapsto7\longmapsto\cdots
\]

is divergent.

The precise gap occurs in the argument that, for the indicator of a forward-invariant orbit set \(A_k\), existence of an equilibrium state implies an equilibrium measure with positive mass on \(A_k\). It does not. If \(A_k\) is a wandering divergent orbit containing no cycle, every invariant probability may assign it mass zero; then the pressure of its indicator is zero and a cycle measure outside \(A_k\) can still be an equilibrium state. The paper asserts positive mass without establishing it, and its subsequent no-divergence lemma depends on that step. [\hypertarget{cit:R13:13}{}\hyperlink{ref:R13:13}{13}, \hypertarget{cit:R13:14}{}\hyperlink{ref:R13:14}{14}]

So that preprint does not supply the missing divergence theorem.


\bigskip



\subsection*{Final verdict}


The proposed thermodynamic program produces three rigorous conclusions:

\begin{enumerate}
\item The archimedean pressure and its equilibrium states are explicitly computable, but their zero-temperature limit selects the spurious \(2\)-adic word \(1^\infty\), not the positive trivial cycle \(2^\infty\).

\item Any regular integrality potential vanishing on positive-integer codes is identically zero, because positive-integer valuation codes are dense in the full shift. Exact ordinary integrality is the discontinuous condition that the nested \(2\)-adic representatives eventually stabilize.

\item Terminal carry compatibility is local, but actual descent is governed by the nonlocal score \[
   \Sigma(u,\tau,c)+h(\tau,b).
   \] Excluding all positive stationary Bellman-geodesic rays other than \(2^\infty\) is the unresolved arithmetic synchronization theorem.
\end{enumerate}

Therefore this route, \textbf{as stated}, does not complete a proof of Collatz. Its strongest valid outcome is a no-go theorem for ordinary continuous thermodynamic formalism and a precise identification of the theorem a nonstandard arithmetic thermodynamic theory would have to prove.
\subsection*{References}
\begin{enumerate}[leftmargin=2.2em]
\item \hypertarget{ref:R13:1}{}\hyperlink{cit:R13:1}{\texttt{turn0file7 L245-L296}}
\item \hypertarget{ref:R13:2}{}\hyperlink{cit:R13:2}{\texttt{turn2file1 L239-L274}}
\item \hypertarget{ref:R13:3}{}\hyperlink{cit:R13:3}{\texttt{turn934661view0}}
\item \hypertarget{ref:R13:4}{}\hyperlink{cit:R13:4}{\texttt{turn934661view1}}
\item \hypertarget{ref:R13:5}{}\hyperlink{cit:R13:5}{\texttt{turn934661view2}}
\item \hypertarget{ref:R13:6}{}\hyperlink{cit:R13:6}{\texttt{turn2file1 L386-L406}}
\item \hypertarget{ref:R13:7}{}\hyperlink{cit:R13:7}{\texttt{turn0file7 L269-L296}}
\item \hypertarget{ref:R13:8}{}\hyperlink{cit:R13:8}{\texttt{turn2file0 L32-L120}}
\item \hypertarget{ref:R13:9}{}\hyperlink{cit:R13:9}{\texttt{turn1file4 L580-L645}}
\item \hypertarget{ref:R13:10}{}\hyperlink{cit:R13:10}{\texttt{turn2file1 L407-L418}}
\item \hypertarget{ref:R13:11}{}\hyperlink{cit:R13:11}{\texttt{turn2file2 L888-L907}}
\item \hypertarget{ref:R13:12}{}\hyperlink{cit:R13:12}{\texttt{turn634631view0}}
\item \hypertarget{ref:R13:13}{}\hyperlink{cit:R13:13}{\texttt{turn965186view2}}
\item \hypertarget{ref:R13:14}{}\hyperlink{cit:R13:14}{\texttt{turn965186view1}}
\end{enumerate}

\clearpage
\section{Round 14: approach 9 --- An o-minimal rational-point explosion theorem}
\subsection*{Prompt}
\begin{quote}
Complete a full proof of the Collatz conjecture using the following approach:

\begin{enumerate}
\item An o-minimal rational-point explosion theorem
\end{enumerate}

\textbf{Speculative \textperiodcentered{} P}

A long near-critical word produces exponential relations among 2\textasciicircum{}\{A\_j\}, 3\textasciicircum{}j, C\_j, n, U\textasciicircum{}j(n). Encode families of such relations as rational points near a definable exponential surface. A divergent orbit would supply arbitrarily many coherent rational points of unusually small height.

Apply a Pila-Wilkie-type dichotomy: either there are few such points, contradicting the orbit, or the points lie on a hidden algebraic locus. The latter should force a repeated block and therefore a cycle equation.

Take as long as you need.
\end{quote}
\subsection*{Response}

\subsection*{Verdict}


The proposed o-minimal argument does \textbf{not} yield a proof of the Collatz conjecture. Formalizing it reveals a decisive obstruction rather than merely a missing technical estimate:

\begin{enumerate}
\item A single Collatz orbit supplies only polynomially many rational relations in the orbit length, while their ordinary heights are exponential in that length. This is far below the threshold needed for a Pila\textendash{}Wilkie contradiction.
\item Producing many points by varying the initial integer places them on rational algebraic lines that exist for \textbf{every} finite valuation word, including completely aperiodic words.
\item Consequently, the proposed inference
\end{enumerate}

\[
\text{“algebraic locus”}\quad\Longrightarrow\quad
\text{“repeated valuation block”}
\]

is false for the natural Collatz encoding.
\begin{enumerate}
\item Even if repetition were established, it would only give the usual cycle equation; one would still need an independent proof excluding every nontrivial positive cycle.
\end{enumerate}

What follows is a rigorous derivation of these obstructions and a precise statement of the genuinely new theorem that would be required to repair the approach.


\bigskip



\subsection*{1. Exact affine relations along an accelerated orbit}


Work with the accelerated odd Collatz map

\[
U(x)=\frac{3x+1}{2^{v_2(3x+1)}}
\]

on positive odd integers. Let

\[
x_0=n,\qquad
a_j=v_2(3x_j+1),\qquad
x_{j+1}=\frac{3x_j+1}{2^{a_j}}.
\]

Define

\[
A_0=0,\qquad A_j=\sum_{i=0}^{j-1}a_i
\]

and

\[
C_0=0,\qquad C_{j+1}=3C_j+2^{A_j}.
\]

Induction gives the exact iterate identity

\[
\boxed{\,2^{A_j}x_j=3^j n+C_j\,}
\tag{1}
\]

with

\[
\boxed{\,
C_j=\sum_{i=0}^{j-1}3^{j-1-i}2^{A_i}.
\,}
\tag{2}
\]

For \(j\geq1\), \(C_j\) is odd: \(C_1=1\), and thereafter \(3C_j\) is odd while \(2^{A_j}\) is even.

It is useful to write (1) as

\[
x_j=p_jn+q_j,
\qquad
p_j=\frac{3^j}{2^{A_j}},
\qquad
q_j=\frac{C_j}{2^{A_j}}.
\tag{3}
\]

Let

\[
\alpha=\log_2 3,\qquad
\Delta_j=A_j-\alpha j.
\]

Then

\[
p_j=2^{-\Delta_j}.
\tag{4}
\]

Thus a near-critical prefix, meaning \(\Delta_j=O(1)\), has a bounded multiplicative part. Its growth is then governed principally by the positive affine term \(q_j\).


\bigskip



\subsection*{2. The natural o-minimal surface}


The canonical rational point associated with a prefix is

\[
P_j=(j,A_j,p_j,q_j,n,x_j).
\]

These points lie in the \(\mathbb R_{\exp}\)-definable set

\[
\mathcal X=
\left\{
(u,v,p,q,x,y):
p=\exp(u\log3-v\log2),\quad y=px+q
\right\}.
\tag{5}
\]

Pila\textendash{}Wilkie counting removes the \textbf{algebraic part} of a definable set: the union of its connected positive-dimensional semialgebraic subsets. Outside that part, the number of rational points of height at most \(H\) is \(O_{\varepsilon}(H^\varepsilon)\). [\hypertarget{cit:R14:1}{}\hyperlink{ref:R14:1}{1}, \hypertarget{cit:R14:2}{}\hyperlink{ref:R14:2}{2}]

The 2024 theorem of Binyamini, Novikov and Zak strengthens this in \(\mathbb R_{\exp}\): the rational-point count in the transcendental part is bounded by a power of \(\log H\), with the exponent and constants depending on the definable complexity. [\hypertarget{cit:R14:3}{}\hyperlink{ref:R14:3}{3}]

Unfortunately, both the geometry and the height scale of (5) work against the Collatz application.


\bigskip



\subsection*{3. Height explosion without point explosion}


For a reduced rational number \(r/s\), write

\[
H(r/s)=\max(|r|,s).
\]

Because \(3^j\) and \(2^{A_j}\) are coprime,

\[
H(p_j)=\max(3^j,2^{A_j})\geq 3^j.
\tag{6}
\]

Because \(C_j\) is odd,

\[
H(q_j)=\max(C_j,2^{A_j})\geq 2^{A_j}.
\tag{7}
\]

Suppose a length-\(L\) word is near-critical at its endpoint:

\[
A_L=\alpha L+O(1).
\]

Since \(A_j\leq A_L\) for \(j\leq L\), equation (2) gives

\[
C_j
 \leq
 \sum_{i<j}3^{j-1-i}2^{A_L}
 \leq
 j3^j2^{A_L}
 =
 \exp(O(L)).
\]

For a fixed seed \(n\), all coordinates of \(P_j\), \(j\leq L\), therefore have height at most

\[
H_L=\exp(O(L)).
\tag{8}
\]

But there are only \(L\) prefix points:

\[
\#\{P_1,\ldots,P_L\}=O(L)=O(\log H_L).
\tag{9}
\]

Even using every contiguous block \([i,j)\) gives only

\[
O(L^2)=O((\log H_L)^2)
\]

relations. Fixed-arity tuples of blocks give at most \(O(L^r)\), still a power of \(\log H_L\).

This is entirely compatible not only with ordinary Pila\textendash{}Wilkie but also with the sharp polylogarithmic theorem for \(\mathbb R_{\exp}\).

A simple model shows why no general improvement can fix this. The transcendental graph

\[
y=2^x
\]

contains the rational points

\[
(k,2^k),\qquad 1\leq k\leq \log_2 H,
\]

of height at most \(H\). Thus \(\Theta(\log H)\) points are natural even on a genuinely transcendental exponential curve.

The Collatz orbit produces a \textbf{height explosion}, not a rational-point explosion.


\bigskip



\subsection*{4. The natural surface is entirely algebraic in the Pila\textendash{}Wilkie sense}


There is a second, more direct problem.

Fix a point

\[
z=(u,v,p,q,x,y)\in\mathcal X.
\]

Keeping \(u,v,p,q\) fixed and varying \(x\), one obtains the line

\[
\left\{
(u,v,p,q,t,pt+q):t\in\mathbb R
\right\}\subseteq\mathcal X.
\]

This is a connected positive-dimensional semialgebraic subset. Therefore every point of \(\mathcal X\) lies in its algebraic part:

\[
\boxed{\mathcal X^{\mathrm{alg}}=\mathcal X.}
\tag{10}
\]

So Pila\textendash{}Wilkie counts no relevant points at all.

Dropping the affine coordinates does not immediately help. Consider

\[
\mathcal S=
\left\{
(u,v,p):p=\exp(u\log3-v\log2)
\right\}.
\tag{11}
\]

Through every point of \(\mathcal S\) passes the affine line

\[
t\longmapsto
\bigl(u+t\log2,\ v+t\log3,\ p\bigr),
\]

because

\[
(u+t\log2)\log3-(v+t\log3)\log2
=
u\log3-v\log2.
\]

Hence

\[
\mathcal S^{\mathrm{alg}}=\mathcal S.
\tag{12}
\]

One could try to eliminate the ruling by imposing that \(u\) and \(v\) are integers satisfying the Collatz recurrence. But an infinite discrete set such as \(\mathbb Z\) is not definable in an o-minimal structure. Unrolling \(L\) steps instead produces a different definable set \(\mathcal X_L\) whose dimension and definable complexity grow with \(L\); the point-counting constants then grow as well. The effective polylogarithmic theorem explicitly depends on this complexity. [\hyperlink{ref:R14:3}{3}]


\bigskip



\subsection*{5. Exact cylinder lemma}


The most important obstruction comes from the precise arithmetic of finite valuation words.

Let

\[
w=(a_0,\ldots,a_{m-1}),\qquad a_i\geq1,
\]

and put

\[
A=\sum_{i=0}^{m-1}a_i,
\qquad
C=C_w=
\sum_{i=0}^{m-1}3^{m-1-i}2^{A_i}.
\]


\subsubsection*{Lemma}


There is a unique odd residue \(r_w\pmod {2^{A+1}}\) such that the positive odd integers realizing the exact valuation word \(w\) are precisely

\[
n\equiv r_w\pmod {2^{A+1}}.
\tag{13}
\]

For \(n=r_w+2^{A+1}k\), the endpoint is

\[
U^m(n)=s_w+2\cdot3^m k,
\tag{14}
\]

where

\[
s_w=\frac{3^m r_w+C}{2^A}
\]

is odd. Consequently every start\textendash{}endpoint pair lies on the rational algebraic line

\[
\boxed{\,2^A y=3^m x+C.\,}
\tag{15}
\]


\paragraph*{Proof}


The word is realized exactly when

\[
x_m=\frac{3^m n+C}{2^A}
\]

is odd. Equivalently,

\[
3^m n+C\equiv2^A\pmod {2^{A+1}}.
\tag{16}
\]

Since \(3^m\) is invertible modulo \(2^{A+1}\), there is exactly one residue class \(r_w\) satisfying (16).

It remains to verify that the final congruence forces all preceding valuations. Write

\[
C_{t+1}=3C_t+2^{A_t}.
\]

Suppose

\[
3^{t+1}n+C_{t+1}
\equiv2^{A_{t+1}}
\pmod {2^{A_{t+1}+1}}.
\]

Reducing modulo \(2^{A_t+1}\), noting \(A_{t+1}\geq A_t+1\), gives

\[
3(3^t n+C_t)+2^{A_t}\equiv0
\pmod {2^{A_t+1}}.
\]

Thus

\[
3(3^t n+C_t)\equiv2^{A_t}
\pmod {2^{A_t+1}}.
\]

Multiplication by the inverse of the odd number \(3\) does not change the residue \(2^{A_t}\) modulo \(2^{A_t+1}\), so

\[
3^t n+C_t\equiv2^{A_t}
\pmod {2^{A_t+1}}.
\]

Iterating backward proves every exact valuation condition.

Finally, substituting \(n=r_w+2^{A+1}k\) into the endpoint formula yields

\[
U^m(n)
=
\frac{3^m r_w+C}{2^A}
+
2\cdot3^m k,
\]

which is (14). \ensuremath{\blacksquare}


\bigskip



\subsection*{6. Why the algebraic-locus conclusion is false}


The cylinder lemma shows that the algebraic line (15) occurs for \textbf{every finite valuation word}.

No periodicity is required. The word could be chosen randomly, could have maximal combinatorial complexity, or could be the prefix of an aperiodic infinite sequence. It still generates infinitely many rational start\textendash{}endpoint points on the line

\[
2^A y=3^m x+C.
\]

Therefore the proposed implication

\[
\text{many points lie on an algebraic locus}
\quad\Longrightarrow\quad
\text{the valuation word repeats}
\]

is false.

The algebraic locus records only the elementary fact that a fixed finite Collatz word defines an affine transformation and a \(2\)-adic residue cylinder. It does not record repetition.

Indeed, the relevant lines are defined over \(\mathbb Q\), not merely over \(\mathbb R\). Thus replacing the ordinary Pila\textendash{}Wilkie algebraic part by a more arithmetic notion would not remove this obstruction.


\bigskip



\subsection*{7. An aperiodic, exactly near-critical counter-blueprint}


There is an especially sharp way to see the problem.

Let

\[
A_j=\left\lfloor \alpha j+\frac12\right\rfloor,
\qquad
\alpha=\log_2 3,
\]

and define

\[
a_j=A_{j+1}-A_j.
\]

Since \(1<\alpha<2\),

\[
a_j\in\{1,2\}.
\]

Moreover,

\[
\left|A_j-\alpha j\right|\leq\frac12.
\tag{17}
\]

The sequence \((a_j)\) is not eventually periodic. Otherwise its average would eventually equal a rational number \(B/m\), whereas

\[
\lim_{j\to\infty}\frac{A_j}{j}=\alpha
\]

is irrational.

Every finite prefix of this aperiodic near-critical sequence is realized by infinitely many positive integers, by the cylinder lemma. The corresponding residue cylinders are nested and therefore define a unique \(2\)-adic integer

\[
n_\ast\in\mathbb Z_2
\]

realizing the entire infinite word.

Now write

\[
\Delta_j=A_j-\alpha j.
\]

From (4),

\[
2^{-1/2}\leq p_j\leq2^{1/2}.
\tag{18}
\]

Furthermore,

\[
\begin{aligned}
q_j
&=
\sum_{i=0}^{j-1}
3^{j-1-i}2^{A_i-A_j} \\
&=
\frac13
\sum_{i=0}^{j-1}2^{\Delta_i-\Delta_j}.
\end{aligned}
\tag{19}
\]

Since \(-1<\Delta_i-\Delta_j<1\),

\[
\boxed{\frac{j}{6}<q_j<\frac{2j}{3}.}
\tag{20}
\]

Consequently, if the \(2\)-adic seed \(n_\ast\) happened to be an ordinary positive integer, its orbit would satisfy

\[
\frac{n_\ast}{\sqrt2}+\frac{j}{6}
<
x_j
<
\sqrt2\,n_\ast+\frac{2j}{3}
\tag{21}
\]

and would diverge at linear speed.

This construction has all of the proposed finite-level features:

\begin{itemize}
\item exact near-criticality;
\item exact finite Collatz realizability;
\item nested \(2\)-adic coherence;
\item rational algebraic lines for every prefix;
\item no repeated block.
\end{itemize}

The only unresolved issue is whether the inverse-limit seed \(n_\ast\) is an ordinary positive integer rather than a genuinely non-Archimedean \(2\)-adic integer. That is exactly the arithmetic distinction a Collatz proof must detect.

Ordinary real o-minimal geometry does not see this distinction.


\bigskip



\subsection*{8. Divergence need not initially be near-critical}


There is also a missing reduction at the beginning of the proposed argument.

From (3) and positivity of \(q_j\),

\[
x_j\geq n\frac{3^j}{2^{A_j}}
=n\,2^{-\Delta_j}.
\tag{22}
\]

If, for example,

\[
A_j\leq(\alpha-\eta)j
\]

along an infinite sequence of times, then

\[
x_j\geq n2^{\eta j}.
\]

Thus a hypothetical divergent trajectory might have a persistent low valuation average and grow exponentially. It need not provide arbitrarily long segments with bounded \(\Delta_j\).

This is visible even at the finite level. The word

\[
a_0=\cdots=a_{L-1}=1
\]

has

\[
A_L=L,\qquad
C_L=3^L-2^L,
\]

and is realized exactly by

\[
n\equiv-1\pmod {2^{L+1}}.
\]

Hence ordinary positive integers exhibit arbitrarily long exponentially expanding all-\(1\) segments. The nested infinite word corresponds to the \(2\)-adic seed \(-1\), not to a positive integer.

A full argument based only on near-critical words would first need a new theorem reducing every unbounded counterexample either to a near-critical tail or to some other configuration that can be excluded. That reduction is not automatic.


\bigskip



\subsection*{9. Approximate rational points do not repair the count}


Perhaps one could count points \textbf{near} the exponential surface rather than exactly on it. The exact cylinder arithmetic again gives the critical obstruction.

Fix a word \(w\), and write

\[
p=\frac{3^m}{2^A},
\qquad
q=\frac{C}{2^A}.
\]

For every seed \(n\) in its cylinder, with endpoint \(y=U^m(n)\),

\[
\frac yn=p+\frac qn.
\tag{23}
\]

Take seeds in the cylinder with \(n\asymp T\). There are \(\asymp_w T\) such seeds. The corresponding rational numbers \(y/n\) are distinct and satisfy

\[
H(y/n)=O_w(T),
\qquad
\left|\frac yn-p\right|\asymp_w \frac1T.
\tag{24}
\]

Thus at height \(H\) there can be \(\asymp H\) rational points lying at distance \(\asymp H^{-1}\) from the relevant exponential value.

So a near-point theorem giving a sublinear count at the natural \(H^{-1}\) scale would simply be false. One would require an approximation of order

\[
H^{-1-\delta}
\]

or stronger. The positive Collatz offset \(C\) naturally produces only the critical \(H^{-1}\) scale.

There are Pila\textendash{}Wilkie-type results for rational points exceptionally close to definable sets, but Habegger's general result is formulated for \textbf{polynomially bounded} o-minimal structures. The unrestricted exponential structure needed here is not polynomially bounded. [\hypertarget{cit:R14:4}{}\hyperlink{ref:R14:4}{4}] More importantly, the elementary family (23) already demonstrates that the required estimate cannot hold at the approximation strength furnished by the Collatz identity.


\bigskip



\subsection*{10. Repeated blocks only produce the cycle equation}


Suppose, despite the preceding obstructions, that one somehow proves eventual periodicity of the valuation sequence.

Let

\[
w=(a_0,\ldots,a_{m-1})
\]

be a repeating block, with total valuation

\[
B=\sum_{i=0}^{m-1}a_i
\]

and affine offset \(D=C_w\). Its block map is

\[
F_w(x)=\frac{3^m x+D}{2^B}.
\tag{25}
\]

The infinite periodic word \(www\cdots\) determines a unique point of \(\mathbb Z_2\), because its successive cylinders have moduli \(2^{kB+1}\) tending to infinity. Applying \(F_w\) shifts the word by one period, so uniqueness gives

\[
F_w(x)=x.
\]

Hence any ordinary positive integer with that periodic word satisfies

\[
\boxed{\,(2^B-3^m)x=D.\,}
\tag{26}
\]

Since \(D>0\) and \(x>0\),

\[
2^B>3^m.
\]

Thus eventual periodicity does indeed give a positive Collatz cycle.

But (26) is not a contradiction. To finish Collatz one must prove that its only positive integral solution satisfying the exact valuation conditions is the trivial cycle

\[
x=1,\qquad w=(2),
\]

or repetitions of that block.

Pila\textendash{}Wilkie is unlikely to supply this final step: cycle points lie on the most algebraic part of the construction, exactly the part removed from transcendental point counting.


\bigskip



\subsection*{11. Would the newest algebraic-relation counting theory help?}


A recent August 2026 paper by Binyamini, Hirata-Kohno, Kawashima and Salant develops a proposed extension of Pila\textendash{}Wilkie from points to intersections with algebraic varieties. The full polynomial-in-degree-and-log-height statement is conjectural; a weaker version is proved for certain trajectories of polynomial differential equations. [\hypertarget{cit:R14:5}{}\hyperlink{ref:R14:5}{5}]

Even granting the full conjecture would not create the needed contradiction here.

A length-\(L\) Collatz orbit has only

\[
\binom{L+1}{2}=O(L^2)
\]

contiguous blocks. Each block gives a line

\[
2^B y=3^m x+D
\]

of degree \(1\) and logarithmic height \(O(L)\), in a near-critical regime. Thus the orbit supplies only polynomially many algebraic relations as a function of their logarithmic height. That is exactly the amount the conjectural relation-counting theorem is designed to allow.


\bigskip



\subsection*{12. Precise no-go theorem}


The preceding calculations can be summarized as follows.


\subsubsection*{Theorem: obstruction to the canonical o-minimal Collatz program}


For the accelerated Collatz map:

\begin{enumerate}
\item The rational point associated with the \(j\)-th prefix has height at least \(3^j\).
\item A fixed length-\(L\) orbit segment supplies only \(O(L)\) prefix points and \(O(L^2)\) contiguous-block relations.
\item For a near-critical length-\(L\) segment, all these objects have height \(\exp(O(L))\), so their number is only polylogarithmic in their height.
\item Varying the seed over all integers realizing a fixed word produces infinitely many points on the rational line \[
   2^A y=3^m x+C.
   \]
\item Such a line occurs for every finite word, whether periodic or aperiodic.
\item There exist aperiodic near-critical infinite valuation sequences whose every finite prefix is exactly realizable and whose cylinder system is projectively coherent in \(\mathbb Z_2\).
\item Near-point normalization naturally gives approximation only at order \(H^{-1}\), at which linearly many rational points may occur.
\end{enumerate}

Therefore neither ordinary Pila\textendash{}Wilkie counting, its sharp polylogarithmic \(\mathbb R_{\exp}\) form, nor the presently conjectured algebraic-relation extensions imply eventual periodicity or contradict a hypothetical Collatz counterexample under the natural encoding.


\bigskip



\subsection*{13. What would actually be needed}


A repaired o-minimal strategy would require three genuinely new ingredients.


\paragraph*{A single-orbit exponential amplifier}


From one nonperiodic orbit segment of length \(L\), one would have to construct, in a \textbf{fixed-dimensional and fixed-complexity} definable set, at least

\[
\exp(cL)
\]

distinct rational points of height at most

\[
\exp(CL).
\]

Then ordinary Pila\textendash{}Wilkie, with \(\varepsilon<c/C\), could force algebraic concentration. Prefixes and contiguous blocks produce only \(L\) and \(L^2\) points, respectively.


\paragraph*{Same-seed algebraic-block rigidity}


One would need a theorem saying that an algebraic block containing many points arising from the \textbf{same ordinary seed} forces two infinite valuation tails to coincide.

This must distinguish the same-seed condition from the universal cylinder lines supplied by arbitrary finite words. Standard o-minimal algebraic-part theory has no such arithmetic memory.


\paragraph*{Positive-cycle exclusion}


Finally, one must prove that

\[
(2^B-3^m)x=D_w
\]

has no positive integral exact-word solutions except the \(1\)-cycle.

Without this third ingredient, algebraic rigidity merely converts a divergent counterexample into a cyclic one.

The most plausible corrected framework would therefore be mixed Archimedean\textendash{}\(2\)-adic rather than purely o-minimal. Its central theorem would need to distinguish a projectively coherent \(2\)-adic seed from a seed lying on the diagonal copy of \(\mathbb N\). That diagonal inverse-limit problem is precisely what the balanced aperiodic construction exposes.

Accordingly, approach 9 does not currently complete a proof of Collatz. Its natural ``rational-point explosion theorem'' is false: one fixed orbit has too few points, while the available explosion obtained by varying seeds is universally algebraic and carries no periodicity information.
\subsection*{References}
\begin{enumerate}[leftmargin=2.2em]
\item \hypertarget{ref:R14:1}{}\hyperlink{cit:R14:1}{\texttt{turn430754view0}}
\item \hypertarget{ref:R14:2}{}\hyperlink{cit:R14:2}{\texttt{turn184672view0}}
\item \hypertarget{ref:R14:3}{}\hyperlink{cit:R14:3}{\texttt{turn184672view1}}
\item \hypertarget{ref:R14:4}{}\hyperlink{cit:R14:4}{\texttt{turn506945view0}}
\item \hypertarget{ref:R14:5}{}\hyperlink{cit:R14:5}{\texttt{turn184672view3}}
\end{enumerate}

\clearpage
\section{Round 15: approach 10 --- A canonical height on a parity orbifold}
\subsection*{Prompt}
\begin{quote}
Complete a full proof of the Collatz conjecture using the following approach:

\begin{enumerate}
\item A canonical height on a parity orbifold
\end{enumerate}

\textbf{Speculative \textperiodcentered{} P}

Resolve the piecewise definition of T by enlarging the space with a parity label, making Collatz an algebraic correspondence rather than a discontinuous map. Seek a canonical height

hhat(Tx) = lambda * hhat(x) + O(1)

with local contributions at infinity, 2, and 3.

Northcott-type finiteness and the product formula might imply that every integral orbit either enters a finite height-zero set or violates a local-height inequality. The trivial cycle would be the only height-zero integral orbit.

Take as long as you need.
\end{quote}
\subsection*{Response}
I do not obtain a valid proof of Collatz from this approach. Formalizing it reveals something fairly decisive: \textbf{the desired canonical-height package is internally incompatible with the known backward tree of the trivial Collatz cycle}.

This is not merely the observation that Collatz is open. There are rigorous no-go theorems below showing that a Northcott height with the proposed scaling cannot exist on a finite parity orbifold. The construction does, however, identify a more precise adelic gluing problem that could replace the canonical-height idea.


\subsubsection*{1. The parity orbifold can be constructed exactly}


Use the shortcut Collatz map

\[
T(n)=
\begin{cases}
n/2,&n\equiv0\pmod 2,\\[2mm]
(3n+1)/2,&n\equiv1\pmod 2.
\end{cases}
\]

This is equivalent to the usual formulation.

Introduce two copies of the affine line,

\[
X_0\sqcup X_1,
\]

where a point \((u,e)\) represents

\[
n=2u+e,\qquad e\in\{0,1\}.
\]

Then

\[
T(2u+e)=3^e u+2e.
\]

Write the output as \(2v+e'\). The lifted transition relation is therefore

\[
C_{e,e'}:\qquad 2v+e'=3^e u+2e,
\]

or

\[
v=f_{e,e'}(u)
   =\frac{3^e u+2e-e'}{2}.
\tag{1}
\]

For an integral \(u\), exactly one value of \(e'\) makes \(v\) integral:

\[
e'\equiv u\pmod2.
\tag{2}
\]

Thus Collatz is the unique path through the correspondence satisfying the local integrality condition at \(2\).

Over \(\mathbf Q\), however, all four maps \(f_{e,e'}\) exist. Moreover, for every fixed pair \((e,e')\), the genuine integral Collatz transitions with

\[
u\equiv e'\pmod 2
\]

form an infinite subset of the line \(C_{e,e'}\), and are therefore Zariski dense in that line.

Consequently:


\begin{quote}
\textbf{The Zariski closure of the integral Collatz transition relation contains all four branches \(C_{e,e'}\).}

\end{quote}


The condition selecting the actual Collatz branch is not algebraic over \(\mathbf Q\). It is a \(2\)-adic integrality condition. Passing from Collatz to its algebraic closure necessarily adds the nonintegral branches and their ``ghost'' paths.

This is consistent with the classical parity-vector picture: every binary word of length \(k\) is realized by exactly one residue class modulo \(2^k\), and every infinite binary word corresponds to one \(2\)-adic integer. [\hypertarget{cit:R15:1}{}\hyperlink{ref:R15:1}{1}]


\subsubsection*{2. The correspondence has polarization factor \(1\), not \(>1\)}


A Call\textendash{}Silverman canonical height starts with a divisor class \(L\) satisfying

\[
f^*L\sim \alpha L,\qquad \alpha>1.
\]

For a polarized correspondence, Ingram uses the analogous relation

\[
y^*L\sim \alpha x^*L,\qquad \alpha>1,
\]

and constructs a path-space height satisfying exact scaling

\[
\widehat h(\sigma P)=\alpha \widehat h(P),
\qquad
\widehat h(P)=h_L(\pi P)+O(1).
\]

The hypothesis \(\alpha>1\) is essential in both constructions. [\hypertarget{cit:R15:2}{}\hyperlink{ref:R15:2}{2}, \hypertarget{cit:R15:3}{}\hyperlink{ref:R15:3}{3}]

Now compactify the two parity components to

\[
X=X_0\sqcup X_1,\qquad X_e\simeq\mathbf P^1.
\]

Let \(L\) be ample. Then

\[
L|_{X_e}\simeq \mathcal O_{\mathbf P^1}(d_e),
\qquad d_e>0.
\]

Each component \(C_{e,e'}\) in (1) is the graph of a degree-one M\"obius transformation. Hence on \(C_{e,e'}\),

\[
\deg(x^*L)=d_e,
\qquad
\deg(y^*L)=d_{e'}.
\]

If the correspondence were polarized by \(\alpha\), we would have

\[
d_{e'}=\alpha d_e
\tag{3}
\]

for every \(e,e'\).

Taking the self-branch \(e=e'=0\) gives

\[
d_0=\alpha d_0.
\]

Since \(d_0>0\),

\[
\boxed{\alpha=1.}
\]

The branch \(1\to1\) gives the same conclusion.

Thus:


\begin{quote}
\textbf{Proposition 1.} The parity correspondence admits no ample polarization of weight \(\alpha>1\).

\end{quote}


Refining the label to finitely many parity digits does not help. Every branch is still generically degree one. If positive divisor weights satisfy \(d_t=\alpha d_s\) along every transition, multiplying around any directed cycle gives

\[
d_s=\alpha^m d_s,
\]

and therefore \(\alpha=1\).

Using a divisor of degree zero to evade this argument destroys ampleness and hence destroys the Northcott property.


\subsubsection*{3. A stronger Northcott obstruction}


The preceding obstruction uses algebraic geometry. There is also an elementary argument that rules out the proposed height even on the positive integers themselves.


\paragraph*{Theorem 2: Northcott versus the powers-of-two chain}


There is no function

\[
H:\mathbf N_{>0}\longrightarrow\mathbf R
\]

with both of the following properties:

\begin{enumerate}
\item For every \(B\), the set \[
   \{n:H(n)\le B\}
   \] is finite.

\item For some fixed \(\lambda>1\) and \(C<\infty\), \[
   |H(Tn)-\lambda H(n)|\le C
   \tag{4}
   \] for every \(n\).
\end{enumerate}


\subparagraph*{Proof}


Since

\[
T(2^k)=2^{k-1},
\]

inequality (4) implies

\[
\lambda H(2^k)\le H(2^{k-1})+C.
\]

Iterating,

\[
H(2^k)
\le
\lambda^{-k}H(1)
+
C\sum_{j=1}^{k}\lambda^{-j}.
\]

Therefore

\[
H(2^k)
\le
\lambda^{-k}H(1)
+
\frac{C}{\lambda-1}
\]

for every \(k\). All the infinitely many points \(2^k\) lie in one bounded-height sublevel set, contradicting property 1. \ensuremath{\blacksquare}

This applies directly to Northcott finiteness because all the points \(2^k\) are rational and have algebraic degree one.

The contradiction is even sharper for an exact canonical relation.


\paragraph*{Corollary 3}


Suppose

\[
\widehat h(Tn)=\lambda\widehat h(n)
\tag{5}
\]

exactly and \(\widehat h\) has Northcott finiteness. Then no such height exists for any \(\lambda>0\).

Indeed, the cycle \(1\leftrightarrow2\) gives

\[
\widehat h(2)=\lambda\widehat h(1),
\qquad
\widehat h(1)=\lambda\widehat h(2).
\]

If \(\lambda\ne1\), then both cycle heights are zero. Equation (5) then gives

\[
\widehat h(2^k)=0
\]

for all \(k\). If \(\lambda=1\), all the values \(\widehat h(2^k)\) are equal. Either way, Northcott fails.

This exposes a fundamental mismatch:


\begin{quote}
A canonical height that vanishes on the trivial cycle must also vanish on every point in its backward basin.

\end{quote}


That basin is already known to be infinite, independently of the Collatz conjecture. It contains every power of two. Therefore the ``trivial cycle is the only height-zero integral orbit'' principle cannot be implemented by a Northcott height on starting points.

One could quotient all points in the trivial basin into a single orbit-class, but constructing a separated arithmetic quotient with only one positive basin component is essentially the Collatz conjecture itself.


\subsubsection*{4. Setting \(\lambda=1\) makes the relation noncoercive}


Geometry forces \(\lambda=1\). For a degree-one affine map, a Weil height automatically satisfies

\[
h(f(x))=h(x)+O(1).
\]

So

\[
\widehat h(Tx)=\widehat h(x)+O(1)
\]

contains no termination information: the bounded error can accumulate linearly for arbitrarily many steps.

There are explicit arbitrarily long upward segments. Put

\[
N_k=2^k-1.
\]

For \(0\le j\le k\),

\[
T^j(N_k)=3^j2^{k-j}-1.
\tag{6}
\]

For \(j<k\), this number is odd, so the first \(k\) transitions all use the odd branch. Consequently,

\[
\log T^k(N_k)-\log N_k
=
\log(3^k-1)-\log(2^k-1)
=
k\log\frac32+o(1).
\tag{7}
\]

Hence no height satisfying

\[
\widehat h(n)=c\log n+O(1),
\qquad c>0,
\]

can be invariant, or uniformly descending, along every Collatz segment.

This also shows why merely adding bounded continuous \(2\)-adic and \(3\)-adic corrections cannot work: the correction changes by \(O(1)\) between the beginning and end of (6), while the archimedean height changes by order \(k\).


\subsubsection*{5. The three local slopes cancel exactly}


For a branch of parity \(e\), the affine multiplier is

\[
a_e=\frac{3^e}{2}.
\]

Its logarithmic sizes at the three relevant places are

\[
\begin{array}{c|ccc}
e & \infty & 2 & 3\\ \hline
0 & -\log 2 & \log 2 & 0\\[1mm]
1 & \log 3-\log 2 & \log 2 & -\log 3.
\end{array}
\tag{8}
\]

Each row sums to zero.

For a parity word \(\varepsilon_0,\ldots,\varepsilon_{k-1}\), let

\[
s_k=\varepsilon_0+\cdots+\varepsilon_{k-1}.
\]

The multiplier of the composed affine map is

\[
A_k=\frac{3^{s_k}}{2^k}.
\]

Its three logarithmic local contributions are

\[
\begin{aligned}
\log|A_k|_\infty&=s_k\log3-k\log2,\\
\log|A_k|_2&=k\log2,\\
\log|A_k|_3&=-s_k\log3.
\end{aligned}
\tag{9}
\]

Therefore

\[
\log|A_k|_\infty+
\log|A_k|_2+
\log|A_k|_3=0.
\tag{10}
\]

This is exactly the product formula for \(3^{s_k}/2^k\).

So the product formula does not produce a contradiction. It says that:

\begin{itemize}
\item archimedean expansion on odd-heavy words,
\item \(2\)-adic expansion from division by two,
\item and \(3\)-adic contraction from multiplication by three
\end{itemize}

are an exactly neutral cocycle.

The affine translation\textemdash{}the \(+1\) and its accumulated carries\textemdash{}is the only possible source of an arithmetic gap. A divisor eigenclass sees the degree-one multiplier and therefore misses precisely the hard part of Collatz.

The long odd segment (6) displays the cancellation particularly clearly:

\[
T^j(N_k)+1=3^j2^{k-j}.
\]

Thus

\[
v_2\!\left(T^j(N_k)+1\right)=k-j,
\qquad
v_3\!\left(T^j(N_k)+1\right)=j.
\]

The archimedean growth is exactly a transfer of local order from \(2\) to \(3\). In particular,

\[
\log(n+1)-v_2(n+1)\log2-v_3(n+1)\log3
\tag{11}
\]

vanishes at every point of this entire segment. More generally, (11) vanishes for every

\[
n=2^a3^b-1,
\]

so such a three-place \(S\)-height has an infinite zero set and cannot have Northcott finiteness.


\subsubsection*{6. The pathwise ``canonical center'' that actually exists}


There is nevertheless a natural pathwise object.

Let \(\varepsilon_j=T^j(n)\bmod2\), and again write

\[
s_k=\sum_{j=0}^{k-1}\varepsilon_j.
\]

Define

\[
R_k
=
\sum_{j=0}^{k-1}
\varepsilon_j
\frac{2^j}{3^{s_{j+1}}}.
\tag{12}
\]

A direct induction gives the exact affine iterate formula

\[
T^k(n)
=
\frac{3^{s_k}}{2^k}\bigl(n+R_k\bigr).
\tag{13}
\]

Since \(T^k(n)\) is integral,

\[
n+R_k
=
\frac{2^kT^k(n)}{3^{s_k}},
\]

and therefore

\[
v_2(n+R_k)\ge k.
\tag{14}
\]

Consequently,

\[
\boxed{R_k\longrightarrow -n\quad\text{in }\mathbf Q_2.}
\tag{15}
\]

This is the classical inverse parity transform in pathwise form. [\hyperlink{ref:R15:1}{1}]

At the real place, however, every term of \(R_k\) is nonnegative. Hence, whenever the series converges in \(\mathbf R\),

\[
R_k\longrightarrow R_\infty\ge0.
\tag{16}
\]

Thus the same rational sequence may converge to

\[
-n<0\quad\text{in }\mathbf Q_2
\]

and to

\[
R_\infty\ge0\quad\text{in }\mathbf R.
\]

There is no contradiction: rational numbers can converge to different limits in different completions when their global heights go to infinity.

And here the global heights necessarily do go to infinity. Write

\[
R_k=\frac{C_k}{3^{s_k}}.
\]

If at least one odd step has occurred, the contribution from the last odd index shows

\[
3\nmid C_k.
\]

Hence the denominator \(3^{s_k}\) is already reduced, and

\[
h(R_k)\ge s_k\log3.
\tag{17}
\]

An infinite positive orbit has infinitely many odd terms, so these center heights are unbounded. Northcott therefore gives no compactness.

There is also no direct \(3\)-adic limit. For \(k<\ell\), if the parity block from \(k\) to \(\ell-1\) contains an odd step, then

\[
R_\ell-R_k
=
\sum_{j=k}^{\ell-1}
\varepsilon_j\frac{2^j}{3^{s_{j+1}}}.
\]

Let \(j_*\) be the last odd index in this interval. After putting this over denominator \(3^{s_\ell}\), the term indexed by \(j_*\) is congruent to \(2^{j_*}\not\equiv0\pmod3\), while all earlier terms are divisible by \(3\). Therefore

\[
v_3(R_\ell-R_k)=-s_\ell.
\tag{18}
\]

So the naive centers are maximally separated at \(3\). A useful \(3\)-adic cancellation can only come from comparing \textbf{different compatible ancestors with matching terminal suffixes}, not successive truncations of one path.

That is exactly where your terminal-suffix locality and competing-ancestor framework become relevant.


\subsubsection*{7. What a viable replacement theorem would have to prove}


The canonical point-height should be replaced by a two-ended adelic rigidity statement.

A suitable normalized difference \(q\) between two compatible affine centers should have no prime denominators outside \(2\) and \(3\). Suppose one could force

\[
v_2(q)\ge A,
\qquad
v_3(q)\ge B,
\tag{19}
\]

from seed and terminal compatibility, while archimedean control gave

\[
0<|q|_\infty<2^A3^B.
\tag{20}
\]

Then \(q\) would be an integer divisible by \(2^A3^B\) but smaller in absolute value than \(2^A3^B\), which is impossible. Equivalently, the product formula would give

\[
|q|_\infty |q|_2 |q|_3<1
\]

while every other local factor is at most \(1\).

This suggests the precise replacement target:


\begin{quote}
\textbf{Adelic synchronization lemma.}   Every hypothetical nonterminating positive orbit produces two distinct compatible finite ancestor states whose normalized affine-center difference satisfies (19) and (20).

\end{quote}


With enough state information, equality \(q=0\) would identify two complete affine states, yielding a repeated block and hence a cycle. A divergent orbit would therefore be impossible.

But establishing that synchronization lemma is the unresolved mathematical content. Finite parity information alone cannot supply it because every finite parity word occurs. The \(2\)-adic seed condition alone cannot supply it because it permits the full binary shift. Successive centers cannot supply it because of the exact \(3\)-adic separation (18). One needs a genuine theorem coupling:

\[
\text{archimedean near-criticality}
\quad+\quad
\text{\(2\)-adic seed agreement}
\quad+\quad
\text{\(3\)-adic terminal-suffix agreement}.
\]


\subsubsection*{Conclusion}


The parity-orbifold construction is valid and illuminating, but it proves a \textbf{no-go result} for the proposed canonical height:

\[
\boxed{
\begin{array}{l}
\text{The algebraic correspondence has polarization factor }1;\\
\text{a Northcott quasi-height with }\lambda>1\text{ is ruled out by }2^k\to1;\\
\text{an exact height vanishing on the cycle vanishes on an infinite basin;}\\
\text{the }\infty,2,3\text{ derivative contributions cancel by the product formula.}
\end{array}}
\]

Therefore a full Collatz proof does not follow from a Call\textendash{}Silverman/Ingram-style canonical height on a finite parity orbifold. The viable residue of the idea is not a canonical point-height but an \textbf{adelic gluing or synchronization theorem for affine carry centers}. That theorem would be a genuine new route to Collatz, but the needed simultaneous \(2\)-adic, \(3\)-adic, and archimedean recurrence has not been established here.
\subsection*{References}
\begin{enumerate}[leftmargin=2.2em]
\item \hypertarget{ref:R15:1}{}\hyperlink{cit:R15:1}{\texttt{turn811209view0}}
\item \hypertarget{ref:R15:2}{}\hyperlink{cit:R15:2}{\texttt{turn811209view2}}
\item \hypertarget{ref:R15:3}{}\hyperlink{cit:R15:3}{\texttt{turn114498view0}}
\end{enumerate}

\clearpage
\section{Round 16: approach 11 --- Arithmetic Conley index at infinity}
\subsection*{Prompt}
\begin{quote}
Complete a full proof of the Collatz conjecture using the following approach:

\begin{enumerate}
\item Arithmetic Conley index at infinity
\end{enumerate}

\textbf{Speculative \textperiodcentered{} B}

Compactify the directed Collatz graph by adding its possible ends. A divergent orbit becomes an invariant set on the boundary. Define an isolating neighborhood and calculate a discrete Conley index.

A trivial index at infinity would exclude divergent trajectories. A nontrivial index could force an invariant end and give a topologically certified route toward a counterexample.

Take as long as you need.
\end{quote}
\subsection*{Response}

\subsection*{Arithmetic Conley index at infinity: a rigorous completion of the framework}


I do \textbf{not} obtain a valid proof of the Collatz conjecture from this approach. The reason is more specific than ``the problem is open'': the proposed argument contains a false implication.

On the natural end compactification of the Collatz graph, \textbf{every end is fixed by the Collatz map}, including ends formed entirely by backward-doubling branches whose forward trajectories converge to the trivial cycle. Therefore, an invariant boundary end does not certify a divergent trajectory.

There is, however, a repaired construction. It gives an arithmetic Conley module whose:

\begin{itemize}
\item free Laurent-polynomial part corresponds exactly to divergent Collatz components;
\item torsion part corresponds exactly to nontrivial cycles;
\item vanishing is equivalent to the full Collatz conjecture.
\end{itemize}

This completes the topological formulation, but the required vanishing theorem remains precisely the unresolved arithmetic content.


\bigskip



\subsubsection*{1. The locally finite Collatz graph}


Use the shortcut Collatz map

\[
T(n)=
\begin{cases}
n/2,&n\equiv 0\pmod 2,\\[2mm]
(3n+1)/2,&n\equiv 1\pmod 2.
\end{cases}
\]

This merely suppresses one forced division by \(2\) after every odd step. The trivial cycle is

\[
1\longmapsto 2\longmapsto 1.
\]

The conjecture is equivalent to every positive integer eventually entering \(\{1,2\}\). This shortcut formulation is standard in the \(3x+1\) literature. [\hypertarget{cit:R16:1}{}\hyperlink{ref:R16:1}{1}]

For every \(m\geq 1\),

\[
T^{-1}(m)
=
\{2m\}
\cup
\begin{cases}
\left\{\dfrac{2m-1}{3}\right\},&m\equiv2\pmod3,\\[3mm]
\varnothing,&m\not\equiv2\pmod3.
\end{cases}
\]

Thus every vertex has one outgoing edge and at most two incoming edges.

Let \(G\) be the directed \(1\)-complex with one vertex for every \(n\geq1\) and one directed edge

\[
e_n:n\longrightarrow T(n)
\]

for every \(n\). We retain parallel edges, so that the two directed edges \(1\to2\) and \(2\to1\) form the trivial topological cycle.

The underlying undirected graph is locally finite.

Define a cellular map \(F:G\to G\) by

\[
F(n)=T(n),\qquad F(e_n)=e_{T(n)}.
\]

Because every vertex and edge has finitely many preimages, \(F\) is proper. Proper maps of locally finite graphs extend continuously to their Freudenthal end compactifications. [\hypertarget{cit:R16:2}{}\hyperlink{ref:R16:2}{2}, \hypertarget{cit:R16:3}{}\hyperlink{ref:R16:3}{3}]

Write

\[
\widehat G=G\sqcup\Omega(G)
\]

for the end compactification.


\bigskip



\subsubsection*{2. The end action is completely degenerate}



\paragraph*{Theorem 1 \textemdash{} Collatz fixes every Freudenthal end}


For every end \(\omega\in\Omega(G)\),

\[
\widehat F(\omega)=\omega.
\]


\paragraph*{Proof}


Represent \(\omega\) by a proper vertex ray

\[
R=(v_0,v_1,v_2,\ldots).
\]

Let \(S\subset V(G)\) be any finite vertex set. Since every point has finitely many preimages,

\[
T^{-1}(S)
\]

is finite. Since \(R\) is proper, for all sufficiently large \(j\),

\[
v_j\notin S\cup T^{-1}(S).
\]

Hence both

\[
v_j\notin S,
\qquad
T(v_j)\notin S.
\]

But \(v_j\) and \(T(v_j)\) are joined by the edge \(e_{v_j}\). Therefore they lie in the same connected component of \(G-S\).

Consequently, for every finite cut \(S\), the tails of \(R\) and \(F(R)\) select the same component of \(G-S\). They represent the same end. Thus

\[
\widehat F(\omega)=\omega.
\]

\(\square\)

This can also be summarized by saying that \(F\) moves every vertex graph-distance one, so it cannot change a Freudenthal direction at infinity.


\paragraph*{A fixed end with no divergence}


Consider the ray

\[
R_2=(1,2,4,8,16,\ldots).
\]

It determines an end \(\omega_2\). Yet

\[
T(2^k)=2^{k-1}\qquad(k\geq1),
\]

so all of these vertices flow toward the trivial cycle rather than toward infinity.

Nevertheless,

\[
\widehat F(\omega_2)=\omega_2.
\]

Thus the implication

\[
\text{“fixed or invariant end”}
\quad\Longrightarrow\quad
\text{“divergent Collatz orbit”}
\]

is false.

A divergent orbit does produce an end: if an orbit is nonperiodic and avoids \(\{1,2\}\), then it eventually leaves every finite set. Otherwise some bounded set would be visited infinitely often, forcing a repeated integer and hence periodicity. But the converse fails dramatically because ordinary inverse-tree ends are also fixed.


\bigskip



\subsubsection*{3. Why the ordinary Conley-index argument cannot work}


The Conley index is defined for an \textbf{isolated invariant set}. For a discrete map, one uses an index pair and the induced index map, considered up to shift equivalence. [\hypertarget{cit:R16:4}{}\hyperlink{ref:R16:4}{4}, \hypertarget{cit:R16:5}{}\hyperlink{ref:R16:5}{5}, \hypertarget{cit:R16:6}{}\hyperlink{ref:R16:6}{6}]

There are three separate obstructions here.


\paragraph*{3.1 The boundary is already invariant}


Since \(F\) fixes \(\Omega(G)\) pointwise,

\[
\Omega(G)\subseteq \operatorname{Inv}(N)
\]

for every end neighborhood \(N\) containing the boundary. This holds whether Collatz is true or false.

So the boundary invariant set can never be empty.


\paragraph*{3.2 Isolation is equivalent to the desired conclusion}


Suppose one chooses a neighborhood of infinity obtained by removing a finite neighborhood of the trivial cycle. For its maximal invariant set to contain only boundary points, one must prove that no ordinary vertex admits a complete trajectory outside the trivial component.

But every \(n\) has the predecessor \(2n\), since

\[
T(2n)=n.
\]

Therefore any forward trajectory avoiding \(\{1,2\}\) can be extended backward indefinitely:

\[
\cdots\longmapsto 8n\longmapsto4n\longmapsto2n
\longmapsto n\longmapsto T(n)\longmapsto\cdots.
\]

Consequently, proving that the end neighborhood has no ordinary invariant points is already equivalent to proving that every orbit reaches the trivial cycle.

The desired isolation property is not a preliminary topological lemma. It is Collatz itself.


\paragraph*{3.3 A trivial Conley index does not imply emptiness}


The general Wa\.zewski implication is only

\[
h(S)\neq 0\quad\Longrightarrow\quad S\neq\varnothing.
\]

Its converse is false: a nonempty isolated invariant set can have the same trivial index as the empty set. Modern expositions give explicit examples of this phenomenon. [\hypertarget{cit:R16:7}{}\hyperlink{ref:R16:7}{7}]

Therefore the proposed step

\[
\text{“the index at infinity is trivial”}
\quad\Longrightarrow\quad
\text{“there is no divergent orbit”}
\]

does not follow from ordinary Conley theory.

Conley-index-at-infinity methods do successfully study trajectories connecting finite invariant sets to compactified infinity, but they require appropriate isolation or isolated-complement hypotheses; those hypotheses are exactly where the arithmetic difficulty enters here. [\hypertarget{cit:R16:8}{}\hyperlink{ref:R16:8}{8}]


\bigskip



\subsection*{4. A repaired arithmetic compactification}


The correct construction should distinguish:

\begin{enumerate}
\item the \textbf{formal point at infinity}, which always exists;
\item ordinary integers that support complete trajectories outside the trivial basin.
\end{enumerate}

Let

\[
X=\mathbb N_{>0}\cup\{\infty\}
\]

be the one-point compactification of the discrete positive integers. Thus neighborhoods of \(\infty\) are sets with finite complement.

Because every finite set has finite preimage under \(T\), the map \(T\) extends continuously to

\[
\overline T:X\to X,
\qquad
\overline T(\infty)=\infty.
\]

Set

\[
N=X\setminus\{1\}
 =\{\infty,2,3,4,\ldots\},
\qquad
L=\{2\}.
\]

Then

\[
N\setminus L
=
\{\infty,3,4,5,\ldots\}
\]

is clopen.

Define

\[
S=\operatorname{Inv}(N\setminus L,\overline T),
\]

where \(\operatorname{Inv}\) consists of points lying on a complete solution

\[
\gamma:\mathbb Z\to N\setminus L,
\qquad
\gamma(j+1)=\overline T(\gamma(j)).
\]


\bigskip



\subsubsection*{Theorem 2 \textemdash{} Exact invariant-set reformulation}


One has

\[
S
=
\{\infty\}
\cup
\left\{
n\geq3:
T^j(n)\notin\{1,2\}\text{ for every }j\geq0
\right\}.
\]

Consequently,

\[
\boxed{
\text{Collatz conjecture}
\iff
S=\{\infty\}.
}
\]


\paragraph*{Proof}


The constant trajectory at \(\infty\) shows that \(\infty\in S\).

Suppose \(n\geq3\) and its forward orbit avoids \(\{1,2\}\). Define a complete trajectory by

\[
\gamma(j)=T^j(n)\qquad(j\geq0)
\]

and

\[
\gamma(-k)=2^k n\qquad(k\geq1).
\]

Since \(T(2m)=m\),

\[
T(\gamma(-k))=\gamma(-k+1).
\]

All negative-time entries are at least \(6\), and all nonnegative-time entries avoid \(\{1,2\}\). Thus \(\gamma(\mathbb Z)\subset N\setminus L\), so \(n\in S\).

Conversely, if \(n\in S\), its complete trajectory lies in \(N\setminus L\), so its forward orbit never enters \(1\) or \(2\).

The equivalence follows.

\(\square\)

Unlike the Freudenthal construction, this compactification has only one formal point at infinity. The nontrivial invariant information is carried by the ordinary points of \(S\).


\bigskip



\subsection*{5. The index pair and its degree-zero index map}


The pair \((N,L)\) is an index pair for \(S\):

\[
\operatorname{Inv}(\overline{N\setminus L},\overline T)
=S
\subseteq
\operatorname{int}(N\setminus L).
\]

The corresponding quotient map

\[
T_{(N,L)}:N/L\longrightarrow N/L
\]

sends a point to the basepoint whenever it leaves \(N\setminus L\).

The space \(X\) is totally disconnected. Every singular simplex of positive dimension has singleton image, so the only nonzero relative homology is in degree zero. With rational coefficients,

\[
V:=H_0(N,L;\mathbb Q)
=
\mathbb Q e_\infty
\oplus
\bigoplus_{n\geq3}\mathbb Q e_n.
\]

The induced index map \(A:V\to V\) is

\[
A e_\infty=e_\infty
\]

and, for \(n\geq3\),

\[
A e_n=
\begin{cases}
e_{T(n)},&T(n)\geq3,\\[1mm]
0,&T(n)=2.
\end{cases}
\]

The only \(n\geq3\) satisfying \(T(n)=2\) is \(n=4\), so

\[
A e_4=0.
\]

Let

\[
W=\bigoplus_{n\geq3}\mathbb Q e_n,
\qquad
A_0=A|_W.
\]

The map \(A_0\) is surjective: for every \(m\geq3\),

\[
A_0 e_{2m}=e_m.
\]

Define the generalized kernel

\[
\operatorname{gker}(A_0)
=
\bigcup_{j\geq1}\ker A_0^j.
\]

Because \(A_0\) is surjective, its Leray reduction\textemdash{}the stable automorphic part of the homological Conley index\textemdash{}is

\[
\mathcal C_{\mathrm{arith}}
:=
W/\operatorname{gker}(A_0),
\]

with the automorphism induced by \(A_0\). Applying a Leray functor to the homological index map is standard in discrete Conley theory. [\hyperlink{ref:R16:6}{6}]


\bigskip



\subsubsection*{Theorem 3 \textemdash{} Exact arithmetic Conley criterion}


For every \(n\geq3\),

\[
e_n\in\operatorname{gker}(A_0)
\iff
n\text{ eventually enters }\{1,2\}.
\]

Consequently,

\[
\boxed{
\text{Collatz conjecture}
\iff
\mathcal C_{\mathrm{arith}}=0.
}
\]


\paragraph*{Proof}


Until an orbit reaches \(2\), the operator simply follows that orbit:

\[
A_0^j e_n=e_{T^j(n)}.
\]

If \(T^j(n)=2\), the preceding value is \(4\), and the corresponding basis vector is killed:

\[
A_0 e_4=0.
\]

Thus an orbit entering the trivial cycle gives

\[
A_0^k e_n=0
\]

for some \(k\).

Conversely, a basis vector can be sent to zero only by arriving at \(4\), whose next Collatz value is \(2\). Hence

\[
A_0^k e_n=0
\]

implies that \(n\) reaches the trivial cycle.

If Collatz holds, every basis vector lies in the generalized kernel. Every vector of \(W\) has finite support, so a single sufficiently large power kills all its basis terms. Therefore

\[
\operatorname{gker}(A_0)=W.
\]

Conversely, if \(n\) is a counterexample, then

\[
A_0^j e_n=e_{T^j(n)}\neq0
\]

for every \(j\), so \(e_n\notin\operatorname{gker}(A_0)\). Its class in \(\mathcal C_{\mathrm{arith}}\) is nonzero.

\(\square\)

The full Leray-reduced index also contains the unavoidable formal fixed-point summand at infinity:

\[
L(V,A)
\cong
\mathbb Q e_\infty
\oplus
\mathcal C_{\mathrm{arith}}.
\]

Thus \(\mathcal C_{\mathrm{arith}}\) is precisely the index after removing the artificial point \(\infty\).


\bigskip



\subsection*{6. The index splits cycles from divergence}


Let

\[
R=\mathbb Q[t,t^{-1}]
\]

and let \(t\) act on \(\mathcal C_{\mathrm{arith}}\) by the automorphism induced by \(A_0\):

\[
t[e_n]=[e_{T(n)}].
\]

Consider the equivalence relation on nonterminating integers

\[
n\sim m
\iff
T^a(n)=T^b(m)
\quad\text{for some }a,b\geq0.
\]

These equivalence classes are exactly the weak components of the nonterminating functional graph.

Every such component is of one of two types:

\begin{enumerate}
\item it contains a nontrivial directed cycle;
\item it contains no cycle, in which case every forward trajectory in it tends to infinity.
\end{enumerate}

The second assertion follows because a nonperiodic orbit that visits a bounded set infinitely often must repeat an integer.

Let \(\mathscr C\) be the set of nontrivial cycle components, with \(p(C)\) denoting the cycle period, and let \(\mathscr D\) be the set of divergent components.


\bigskip



\subsubsection*{Theorem 4 \textemdash{} Counterexample decomposition}


There is an \(R\)-module decomposition

\[
\boxed{
\mathcal C_{\mathrm{arith}}
\cong
\bigoplus_{D\in\mathscr D}R
\;\oplus\;
\bigoplus_{C\in\mathscr C}
R/(t^{p(C)}-1).
}
\]


\paragraph*{Proof}


If \(T(u)=T(v)\), then

\[
A_0(e_u-e_v)=0.
\]

Hence in the quotient,

\[
[e_u]=[e_v].
\]

More generally, if

\[
T^a(u)=T^b(v),
\]

then

\[
t^a[e_u]=t^b[e_v],
\]

and since \(t\) acts invertibly,

\[
[e_u]=t^{\,b-a}[e_v].
\]

Thus every component gives a cyclic \(R\)-module.

Suppose a component contains a cycle

\[
c,T(c),\ldots,T^{p-1}(c)
\]

of least period \(p\). Then

\[
t^p[e_c]=[e_c].
\]

There is no smaller relation, because the \(p\) cycle vertices are distinct. Hence that component contributes

\[
R/(t^p-1).
\]

Now suppose a component is nonperiodic. Choose \(d\) in the component. The vectors

\[
e_d,e_{T(d)},e_{T^2(d)},\ldots
\]

are all distinct. If a nonzero Laurent polynomial \(q(t)\) annihilated \([e_d]\), then after multiplying by a power of \(t\), one would have a nonzero finite linear combination of distinct forward-orbit basis vectors in the generalized kernel. Applying further powers of \(A_0\) merely shifts these distinct basis vectors forward and can never annihilate their combination. This is impossible.

Thus \([e_d]\) has no nonzero Laurent-polynomial annihilator, and the component contributes a free copy of \(R\).

Different components give a direct sum because their forward orbits never meet, so cancellation between components cannot occur.

\(\square\)


\bigskip



\subsubsection*{7. Exact detection of divergence and cycles}


Let \(\mathbb Q(t)\) be the fraction field of \(R\). Tensoring kills all torsion summands:

\[
\mathcal C_{\mathrm{arith}}
\otimes_R\mathbb Q(t)
\cong
\bigoplus_{D\in\mathscr D}\mathbb Q(t).
\]

Therefore

\[
\boxed{
\text{There are no divergent Collatz trajectories}
\iff
\mathcal C_{\mathrm{arith}}
\otimes_R\mathbb Q(t)=0.
}
\]

Likewise,

\[
\operatorname{Tor}_R
\mathcal C_{\mathrm{arith}}
\cong
\bigoplus_{C\in\mathscr C}
R/(t^{p(C)}-1),
\]

and hence

\[
\boxed{
\text{There are no nontrivial Collatz cycles}
\iff
\operatorname{Tor}_R
\mathcal C_{\mathrm{arith}}=0.
}
\]

Finally,

\[
\boxed{
\text{Collatz conjecture}
\iff
\mathcal C_{\mathrm{arith}}=0.
}
\]

This realizes the intended dichotomy almost exactly:

\begin{itemize}
\item a nonzero \textbf{free} Conley class certifies a divergent component;
\item a nonzero \textbf{torsion} Conley class certifies a cycle;
\item the formal end \(\infty\) has been separated off as a canonical fixed-point summand.
\end{itemize}


\bigskip



\subsection*{8. Why this is not yet a proof}


The remaining statement is

\[
\operatorname{gker}(A_0)=W,
\]

or equivalently,

\[
\forall v\in W\ \exists k\geq0:
A_0^k v=0.
\]

Applying this to a basis vector gives

\[
\forall n\geq3\ \exists k\geq0:
A_0^k e_n=0,
\]

which is exactly

\[
\forall n\geq3\ \exists k\geq0:
T^k(n)\in\{1,2\}.
\]

So the final ``index computation'' is not presently supplied by topology. It is the Collatz conjecture written as local nilpotence of an arithmetic index map.

The ordinary Freudenthal approach loses too much information because all ends are fixed. The repaired Laurent-module index keeps exactly the necessary information, but Conley theory by itself gives no reason for its vanishing.

A genuine completion would require an additional arithmetic theorem\textemdash{}for example, a finite-detection result constructing residue-and-height Conley modules \(\mathcal C_r\) such that

\[
\mathcal C_{\mathrm{arith}}
\hookrightarrow
\varprojlim_r\mathcal C_r,
\]

followed by a uniform arithmetic proof that the inverse limit is zero. Such a theorem would need to use the \(2\)-adic parity cylinder, the \(3\)-adic carry constraints, and archimedean drift simultaneously. That would turn this Conley formulation into a real proof mechanism rather than an exact reformulation.


\subsubsection*{Final assessment}


The proposed proof cannot be completed as stated:

\[
\text{invariant Freudenthal end}
\not\Rightarrow
\text{divergent orbit},
\]

and

\[
\text{trivial ordinary Conley index}
\not\Rightarrow
\text{empty invariant set}.
\]

The strongest correct result is the arithmetic Conley module

\[
\mathcal C_{\mathrm{arith}}
=
W/\operatorname{gker}(A_0),
\]

with the exact decomposition

\[
\mathcal C_{\mathrm{arith}}
\cong
\bigoplus_{\text{divergent components}}\mathbb Q[t,t^{-1}]
\oplus
\bigoplus_{\text{nontrivial cycles}}
\frac{\mathbb Q[t,t^{-1}]}{(t^p-1)}.
\]

That is a rigorous and complete topological certificate for every possible counterexample, but proving that this module vanishes still requires a new arithmetic argument.
\subsection*{References}
\begin{enumerate}[leftmargin=2.2em]
\item \hypertarget{ref:R16:1}{}\hyperlink{cit:R16:1}{\texttt{turn215080view4}}
\item \hypertarget{ref:R16:2}{}\hyperlink{cit:R16:2}{\texttt{turn215080view0}}
\item \hypertarget{ref:R16:3}{}\hyperlink{cit:R16:3}{\texttt{turn215080view1}}
\item \hypertarget{ref:R16:4}{}\hyperlink{cit:R16:4}{\texttt{turn955530view0}}
\item \hypertarget{ref:R16:5}{}\hyperlink{cit:R16:5}{\texttt{turn955530view1}}
\item \hypertarget{ref:R16:6}{}\hyperlink{cit:R16:6}{\texttt{turn955530view2}}
\item \hypertarget{ref:R16:7}{}\hyperlink{cit:R16:7}{\texttt{turn215080view3}}
\item \hypertarget{ref:R16:8}{}\hyperlink{cit:R16:8}{\texttt{turn215080view2}}
\end{enumerate}

\clearpage
\section{Round 17: approach 12 --- Proof-theoretic ordinal descent}
\subsection*{Prompt}
\begin{quote}
Complete a full proof of the Collatz conjecture using the following approach:

\begin{enumerate}
\item Proof-theoretic ordinal descent
\end{enumerate}

\textbf{Hybrid \textperiodcentered{} P}

The mixed binary-ternary rewriting formulation suggests assigning an ordinal, rather than an ordinary numerical weight, to each digit string. Find an ordinal-valued rank that strictly decreases under every complete rewrite block.

A rank below a familiar proof-theoretic ordinal would prove termination. Failure patterns could also calibrate why standard matrix orders and low-complexity termination methods fail. Generalized Collatz undecidability makes this route less absurd than it first sounds, although it says nothing directly about the original map.

Take as long as you need.
\end{quote}
\subsection*{Response}

\subsubsection*{Verdict}


After carrying the ordinal-descent program through the mixed binary\textendash{}ternary rewrite system, \textbf{a complete proof of the Collatz conjecture is not established}. What does emerge is a sharper result than the original proposal:

\begin{enumerate}
\item the eleven-rule system can be collapsed to an exact deterministic rewrite map on pure ternary words;
\item Collatz is equivalent to termination of that map;
\item an ordinal rank for the entire mixed system can be constructed automatically from a rank satisfying \textbf{three explicit uniform inequalities} for two carry transducers;
\item several substantial subsystems admit elementary ordinal ranks below \(\omega^2\);
\item fixed lexicographic, length-based, commutative, and local bigram rankings fail by finite certificates.
\end{enumerate}

The unresolved part is precisely the construction of the carry-sensitive rank satisfying those three inequalities. Merely postulating it would assume essentially the whole Collatz conjecture.


\bigskip



\subsection*{1. The mixed binary\textendash{}ternary rewrite system}


Work with the shortcut Collatz map

\[
T(n)=
\begin{cases}
n/2,&n\equiv0\pmod2,\\[2mm]
(3n+1)/2,&n\equiv1\pmod2.
\end{cases}
\]

Reaching \(1\) under this map is equivalent to reaching \(1\) under the usual Collatz map.

Let

\[
f(x)=2x,\qquad t(x)=2x+1
\]

be the binary digits, and

\[
0(x)=3x,\qquad 1(x)=3x+1,\qquad 2(x)=3x+2
\]

be the ternary digits. Write \(\lhd\) and \(\rhd\) for the left and right delimiters. The mixed-base system of Yolcu, Aaronson, and Heule has the following eleven rules:

\[
\begin{aligned}
D:\qquad
&f\rhd\longrightarrow\rhd,
&
&t\rhd\longrightarrow2\rhd,
\\[1mm]
A:\qquad
&f0\longrightarrow0f,
&
&f1\longrightarrow0t,
&
&f2\longrightarrow1f,
\\
&t0\longrightarrow1t,
&
&t1\longrightarrow2f,
&
&t2\longrightarrow2t,
\\[1mm]
B:\qquad
&\lhd0\longrightarrow\lhd t,
&
&\lhd1\longrightarrow\lhd ff,
&
&\lhd2\longrightarrow\lhd ft .
\end{aligned}
\]

The \(A\)-rules exchange adjacent binary and ternary positions while preserving the represented integer. The \(B\)-rules convert the leading ternary digit into binary, and each \(D\)-rule performs one shortcut Collatz step. The source paper proves that termination of the full rewriting system is equivalent to the Collatz conjecture. [\hypertarget{cit:R17:1}{}\hyperlink{ref:R17:1}{1}, \hypertarget{cit:R17:2}{}\hyperlink{ref:R17:2}{2}, \hypertarget{cit:R17:3}{}\hyperlink{ref:R17:3}{3}]


\bigskip



\subsection*{2. Collapse to two ternary carry transducers}


The first simplification is to stop looking for a rank on arbitrary mixed strings. A canonical computation consists of repeatedly taking the rightmost binary symbol, moving it through the ternary suffix using the \(A\)-rules, and finally applying a \(D\)-rule.

Define two maps

\[
E,O:\{0,1,2\}^{*}\longrightarrow\{0,1,2\}^{*}.
\]

Here \(E(u)\) is the pure ternary output obtained by completely processing \(fu\rhd\), and \(O(u)\) is the output obtained from \(tu\rhd\).

The rewrite rules give the exact recursions

\[
E(\epsilon)=\epsilon,\qquad O(\epsilon)=2,
\]

and

\[
\begin{aligned}
E(0u)&=0E(u),&
E(1u)&=0O(u),&
E(2u)&=1E(u),\\
O(0u)&=1O(u),&
O(1u)&=2E(u),&
O(2u)&=2O(u).
\end{aligned}
\tag{1}
\]

For example,

\[
fu\rhd=f0v\rhd\longrightarrow0fv\rhd
\]

when \(u=0v\), and completing the processing gives \(0E(v)\). The other five recursions are identical applications of the six \(A\)-rules. The terminal equations are precisely the two \(D\)-rules.

These transducers are long division by \(2\) in base \(3\). If \(c\in\{0,1\}\) is the current carry and \(d\in\{0,1,2\}\) is the next ternary digit, then

\[
e=\left\lfloor\frac{3c+d}{2}\right\rfloor,
\qquad
c'\equiv3c+d\pmod2.
\tag{2}
\]

The output digit is \(e\), and the new carry is \(c'\). A terminal carry \(0\) outputs nothing, whereas terminal carry \(1\) produces the final digit \(2\). This is the same binary\textendash{}ternary base-conversion phenomenon identified independently in the cellular-automaton formulation of St\'erin and Woods. [\hypertarget{cit:R17:4}{}\hyperlink{ref:R17:4}{4}]


\bigskip



\subsection*{3. A pure ternary Collatz return map}


Suppose the current string contains no binary symbols:

\[
\lhd w\rhd,\qquad w\in\{0,1,2\}^{*}.
\]

Apply one boundary rule and then process every binary symbol created by it. This gives the pure ternary return map

\[
R:\{0,1,2\}^{*}\longrightarrow\{0,1,2\}^{*}
\]

defined by

\[
\boxed{
\begin{aligned}
R(0u)&=O(u),\\
R(1u)&=E(E(u)),\\
R(2u)&=E(O(u)),\\
R(\epsilon)&=\epsilon.
\end{aligned}}
\tag{3}
\]

The reason for the double application in the second and third cases is that

\[
\lhd1\longrightarrow\lhd ff,
\qquad
\lhd2\longrightarrow\lhd ft.
\]

The rightmost binary digit is processed first.

A pure ternary word \(w=d_1\cdots d_q\) represents

\[
N(w)=3^q+[w]_3,
\tag{4}
\]

where \([w]_3\) is the ordinary base-\(3\) value of the \(q\)-digit string. The implicit leading ternary digit is \(1\). Thus \(\epsilon\) represents \(1\).


\subsubsection*{Proposition 1 \textemdash{} Arithmetic meaning of \(R\)}


For every nonempty pure ternary word \(w\),

\[
N(R(w))=
\begin{cases}
T(N(w)),&w=0u,\\
T^2(N(w)),&w=1u\text{ or }w=2u.
\end{cases}
\tag{5}
\]


\paragraph*{Proof}


Each \(B\)-rule preserves the represented integer and creates respectively one, two, or two binary symbols. Every subsequent \(A\)-rule preserves the integer. Each time a binary symbol reaches \(\rhd\), one \(D\)-rule is applied, which performs exactly one application of \(T\). Hence the number of \(T\)-steps is the number of binary symbols introduced by the boundary rule. \ensuremath{\blacksquare}


\subsubsection*{Theorem 2 \textemdash{} Pure ternary equivalence}


The following statements are equivalent:

\begin{enumerate}
\item every positive integer reaches \(1\) under \(T\);
\item every finite ternary word reaches \(\epsilon\) under iteration of \(R\).
\end{enumerate}


\paragraph*{Proof}


If Collatz holds, (5) shows that every \(R\)-orbit is a subsequence of a terminating \(T\)-orbit, so it reaches \(\epsilon\).

Conversely, represent an arbitrary positive integer \(n\) by a finite pure binary word. Each complete binary-processing block consumes one binary symbol. Therefore, after finitely many \(T\)-steps, its mixed representation becomes pure ternary. By assumption the resulting pure word reaches \(\epsilon\) under \(R\), and (5) then shows that \(n\) reaches \(1\). \ensuremath{\blacksquare}

This is the cleanest form of the ordinal problem: the original eleven rules have been reduced to the three schematic transformations in (3).


\bigskip



\subsection*{4. The exact ordinal-descent theorem}


Let \(\alpha\) be any ordinal, and suppose there is a function

\[
\rho:\{0,1,2\}^{*}\longrightarrow\alpha
\]

such that

\[
\rho(\epsilon)=0
\]

and, for every ternary word \(u\),

\[
\boxed{
\begin{aligned}
\rho(O(u))&<\rho(0u),\\
\rho(E(E(u)))&<\rho(1u),\\
\rho(E(O(u)))&<\rho(2u).
\end{aligned}}
\tag{6}
\]


\subsubsection*{Theorem 3 \textemdash{} Ordinal criterion for Collatz}


Any function \(\rho\) satisfying (6) proves the Collatz conjecture.


\paragraph*{Proof}


By (3),

\[
\rho(R(w))<\rho(w)
\]

for every nonempty ternary word \(w\). Consequently an infinite \(R\)-orbit would produce an infinite strictly descending sequence of ordinals,

\[
\rho(w)>\rho(R(w))>\rho(R^2(w))>\cdots,
\]

which is impossible. Therefore every \(R\)-orbit reaches \(\epsilon\). Theorem 2 then gives the Collatz conjecture. \ensuremath{\blacksquare}

This theorem is elementary, but it isolates the entire substantive requirement: \textbf{construct one explicit, noncircular \(\rho\) satisfying the three inequalities (6).}


\bigskip



\subsection*{5. Lifting the ternary rank to the complete rewrite process}


A rank satisfying (6) also gives an ordinal rank for all the intermediate base-conversion steps.

For a mixed configuration \(C\), let

\[
p(C)=\text{number of binary symbols in }C.
\]

Let

\[
I(C)=
\#\{(i,j):i<j,\ C_i\in\{f,t\},\ C_j\in\{0,1,2\}\}
\tag{7}
\]

be the number of binary-before-ternary inversions.

Finally, let \(\Pi(C)\) denote the pure ternary word obtained by processing all binary symbols currently in \(C\), without applying a new boundary rule afterward. This is a finite base-conversion computation: each \(D\)-step removes one binary symbol.

Define

\[
\widehat\rho(C)
=
\omega^2\cdot\rho(\Pi(C))
+\omega\cdot p(C)
+I(C).
\tag{8}
\]


\subsubsection*{Proposition 4 \textemdash{} Full canonical rewrite descent}


Along the canonical rewrite strategy, every individual \(A\)-, \(D\)-, or \(B\)-step strictly decreases \(\widehat\rho\).


\paragraph*{Proof}


For an \(A\)-step, \(\Pi(C)\) and \(p(C)\) are unchanged. A binary symbol passes across exactly one ternary symbol, so \(I(C)\) decreases by exactly one.

For a \(D\)-step, \(\Pi(C)\) remains the eventual pure result, but \(p(C)\) decreases by one. The inversion count may increase by some finite amount, but

\[
\omega(p-1)+I'<\omega p
\]

for every finite \(I'\).

For a \(B\)-step, the old configuration is pure ternary. If it represents the word \(w\), the pure projection after the newly generated binary digits have been processed is \(R(w)\). Therefore

\[
\rho(\Pi(C_{\rm new}))
=
\rho(R(w))
<
\rho(w)
=
\rho(\Pi(C_{\rm old})).
\]

The decrease in the \(\omega^2\)-coefficient dominates the finite reset of \(p\) and \(I\). \ensuremath{\blacksquare}

Thus the search for an ordinal ranking of a complicated mixed rewrite system really does reduce to the three pure ternary inequalities (6). If \(\rho(w)<\varepsilon_0\), then the lifted rank (8) also remains below \(\varepsilon_0\).


\bigskip



\subsection*{6. Genuine ordinal descent that can already be proved}


The failure occurs only when the boundary and dynamic rules are allowed to refuel each other. The individual phases have very simple ordinal ranks.


\subsubsection*{Proposition 5 \textemdash{} Dynamic system without boundary conversion}


The subsystem \(A\cup D\) terminates, with rank

\[
\rho_D(C)=\omega\cdot p(C)+I(C)<\omega^2.
\tag{9}
\]


\paragraph*{Proof}


Every \(A\)-rule decreases \(I\) by one. Every \(D\)-rule decreases \(p\) by one. Any finite change in \(I\) is dominated by the drop in the coefficient of \(\omega\). \ensuremath{\blacksquare}


\subsubsection*{Proposition 6 \textemdash{} Boundary conversion without Collatz dynamics}


The subsystem \(A\cup B\) terminates, with rank

\[
\rho_B(C)=\omega\cdot q(C)+I(C)<\omega^2,
\tag{10}
\]

where \(q(C)\) is the number of ternary symbols.


\paragraph*{Proof}


Again every \(A\)-rule decreases \(I\). Every \(B\)-rule consumes one ternary symbol, so \(q\) decreases by one. It can create one or two binary symbols and thereby increase \(I\), but only by a finite amount, dominated by the decrease in \(\omega q\). \ensuremath{\blacksquare}

This explains the essential feedback loop:

\[
\text{binary}
\ \xrightarrow{D}\
\text{ternary}
\ \xrightarrow{B}\
\text{one or two binary symbols}.
\tag{11}
\]

Without \(B\), the binary supply is exhausted. Without \(D\), the ternary supply is exhausted. Together, they can regenerate one another.

The original automated-termination work proves considerably more: every system obtained by deleting any one of the eleven rules has a matrix-interpretation termination proof. Proving any one deleted rule terminating relative to the full system would therefore settle Collatz. [\hypertarget{cit:R17:5}{}\hyperlink{ref:R17:5}{5}]


\bigskip



\subsection*{7. Why the obvious ordinal assignments fail}


The pure return map gives very small counterexamples to several large classes of rankings.


\subsubsection*{Fixed symbol precedence is impossible}


Directly from (1)\textendash{}(3),

\[
R(0)=2,\qquad R(2)=1,\qquad R(11)=01.
\tag{12}
\]

A left-to-right lexicographic order based on one fixed precedence of the symbols would require

\[
0\succ2,\qquad 2\succ1,\qquad 1\succ0,
\]

which is impossible.

Thus no ordinal code whose leading comparison is governed by a fixed order of \(0,1,2\) can work.


\subsubsection*{Length cannot be the leading measure}


One has

\[
R(20)=022.
\tag{13}
\]

A complete block can increase the ternary length from \(2\) to \(3\).

There can also be consecutive length increases; for instance,

\[
2220\longmapsto11022\longmapsto002222.
\]

So the increase is not merely an isolated normalization artifact.


\subsubsection*{Commutative digit weights cannot work}


Another exact transition is

\[
R(120)=012.
\tag{14}
\]

The two words have the same length and the same number of each digit. Therefore no rank depending only on

\[
\bigl(|w|,\#0(w),\#1(w),\#2(w)\bigr)
\]

can decrease on every block. In particular, no additive potential

\[
a_0\#0(w)+a_1\#1(w)+a_2\#2(w)+b|w|
\]

can orient the system.


\subsubsection*{Even additive bigram potentials have a finite obstruction}


Consider the five exact transitions

\[
\begin{aligned}
212&\longmapsto102,\\
0012&\longmapsto121,\\
0202&\longmapsto2122,\\
1021&\longmapsto00122,\\
1222&\longmapsto0202.
\end{aligned}
\tag{15}
\]

Let \(B(w)\) be the vector counting all adjacent pairs in the padded word \({}^{\wedge}w\$\). Direct counting gives

\[
\sum_{\text{five inputs}} B(w)
=
\sum_{\text{five outputs}} B(R(w)).
\tag{16}
\]

For the internal bigrams, both sides have the counts

\[
\begin{array}{c|cccccccc}
\text{bigram}&00&01&02&10&12&20&21&22\\ \hline
\text{count}&1&1&3&1&3&1&2&2.
\end{array}
\]

The initial- and terminal-marker counts agree as well.

Consequently, if

\[
\Phi(w)=\sum_{a,b}c_{ab}\#_{ab}({}^{\wedge}w\$)
\]

strictly decreased on each of the five transitions, summing the five strict inequalities would contradict (16). The same argument excludes any finite lexicographic stack of additive bigram potentials.

This does not prove that every finite-state or bounded-window ordinal interpretation fails. It shows, however, that one must go beyond several standard termination templates.

The rigorous matrix-impossibility theorem in the source literature concerns the older \textbf{unary} Collatz rewriting system, even when dependency pairs are included. It is not a no-go theorem for every matrix or ordinal interpretation of the mixed system. [\hypertarget{cit:R17:6}{}\hyperlink{ref:R17:6}{6}]


\bigskip



\subsection*{8. Why simply choosing a large ordinal does not solve the problem}


There is a fundamental logical point here.

Suppose Collatz were already known. Then the natural-number function

\[
\tau(w)=\min\{k\geq0:R^k(w)=\epsilon\}
\tag{17}
\]

would satisfy

\[
\tau(R(w))=\tau(w)-1.
\]

So a rank into \(\mathbb N\), far below \(\omega\), would suffice.

Conversely, proving that \(\tau(w)\) is defined for every \(w\) is exactly proving Collatz. Defining an ordinal by

\[
\rho(w)=\rho(R(w))+1
\]

has the same circularity: transfinite recursion is available only after the rewrite relation has been shown well-founded.

More generally, the rewrite system is finitely branching. If it terminates from a fixed word, its derivation tree has finite height by K\"onig's lemma, so that word already has a finite natural-number derivation rank. Therefore:


\begin{quote}
The point of \(\varepsilon_0\), \(\Gamma_0\), or another proof-theoretic ordinal is not that Collatz trajectories intrinsically require transfinite heights. The point would be to provide a \textbf{structurally defined, locally verifiable certificate} of finite termination.

\end{quote}


A larger ordinal does not compensate for the absence of such a structural definition.


\bigskip



\subsection*{9. The precise missing lemma}


A successful proof-theoretic attack must establish the following statement.


\subsubsection*{Carry-profile ordinal lemma}


There exists an explicitly defined map

\[
\rho:\{0,1,2\}^{*}\to\varepsilon_0
\]

computed by structural recursion on the finite ternary word, possibly carrying finitely or recursively many auxiliary carry states, such that for every \(u\),

\[
\begin{aligned}
\rho(O(u))&<\rho(0u),\\
\rho(E(E(u)))&<\rho(1u),\\
\rho(E(O(u)))&<\rho(2u).
\end{aligned}
\tag{18}
\]

Moreover, the proof of (18) must use only the structural recursions (1), not an assumption that the future \(R\)-orbit terminates.

The hard feature is that \(E\) and \(O\) scan the entire suffix. Every ternary digit \(1\) flips the division carry, and the terminal carry can append a new digit \(2\). A leading ternary digit may generate two successive carry passes. Hence a rank based on a fixed leading-digit order, a bounded collection of symbol counts, or one fixed local matrix cannot see the full dependency.

A credible ordinal rank would therefore have to encode an \textbf{unbounded hierarchy of carry profiles}. One possible form would be a Cantor-normal-form sum whose exponents record successively deeper carry-reset locations, but no definition has been found for which all three transformations in (18) uniformly decrease. At present that sentence is a research direction, not a theorem.


\bigskip



\subsection*{10. What generalized undecidability does and does not imply}


Kurtz and Simon proved that a natural generalized Collatz totality problem is \(\Pi^0_2\)-complete. Thus one should not expect a uniform mechanical procedure that takes every generalized residue-class affine map and always synthesizes a terminating ordinal interpretation exactly when the map is total. [\hypertarget{cit:R17:7}{}\hyperlink{ref:R17:7}{7}]

That result does \textbf{not} rule out a bespoke ordinal certificate for the original \(3n+1\) map. It does explain why a generic ``all such rewrites decrease below \(\varepsilon_0\)'' argument is unlikely: generalized systems are capable of encoding arbitrary computation.


\bigskip



\subsection*{Final status}


The proof-theoretic ordinal program can be brought to the following exact point:

\[
\boxed{
\text{Collatz}
\iff
\forall w\;\exists k\ R^k(w)=\epsilon
}
\]

with

\[
R(0u)=O(u),\qquad
R(1u)=E(E(u)),\qquad
R(2u)=E(O(u)).
\]

An explicit ordinal rank satisfying the three inequalities

\[
\rho(O(u))<\rho(0u),\qquad
\rho(E(E(u)))<\rho(1u),\qquad
\rho(E(O(u)))<\rho(2u)
\]

would give a complete proof, and formula (8) would lift it to a descent rank for the entire canonical mixed-base rewrite computation.

What is proved here is the reduction, the lifting theorem, the \(\omega^2\)-ranks for both non-refueling subsystems, and finite obstructions to several low-complexity rank families. \textbf{The carry-profile ordinal lemma itself remains unproved.} Treating it as automatic from the existence of a large proof-theoretic ordinal would be circular, so the requested full Collatz proof has not been completed by this approach.
\subsection*{References}
\begin{enumerate}[leftmargin=2.2em]
\item \hypertarget{ref:R17:1}{}\hyperlink{cit:R17:1}{\texttt{turn962870view2}}
\item \hypertarget{ref:R17:2}{}\hyperlink{cit:R17:2}{\texttt{turn962870view3}}
\item \hypertarget{ref:R17:3}{}\hyperlink{cit:R17:3}{\texttt{turn796082view0}}
\item \hypertarget{ref:R17:4}{}\hyperlink{cit:R17:4}{\texttt{turn426790view0}}
\item \hypertarget{ref:R17:5}{}\hyperlink{cit:R17:5}{\texttt{turn962870view4}}
\item \hypertarget{ref:R17:6}{}\hyperlink{cit:R17:6}{\texttt{turn796082view3}}
\item \hypertarget{ref:R17:7}{}\hyperlink{cit:R17:7}{\texttt{turn426790view1}}
\end{enumerate}

\clearpage
\section{Round 18: approach 13 --- Exact arithmetic optimal transport extending Tao}
\subsection*{Prompt}
\begin{quote}
Complete a full proof of the Collatz conjecture using the following approach:

\begin{enumerate}
\item Exact arithmetic optimal transport extending Tao
\end{enumerate}

\textbf{Hybrid \textperiodcentered{} P}

Tao obtains approximate transport for probability distributions. Enrich the state by enough residue and carry data that the transport operator becomes an exact or uniformly contracting map in Wasserstein distance.

The decisive theorem would be contraction \textbf{uniformly over every starting cylinder}, rather than after averaging over starts. That would convert a law-of-large-numbers result into a pointwise theorem and close the precise gap in the almost-all approach.

Take as long as you need.
\end{quote}
\subsection*{Response}

\subsection*{Outcome of the program}


Carrying this approach through produces a strong exact transport theorem, but \textbf{not a proof of Collatz}. The decisive distinction is:

\[
\boxed{\text{the randomized Syracuse operator is uniformly contractive}}
\]

whereas

\[
\boxed{\text{the exact integer Syracuse operator is not uniformly contractive on all cylinders}.}
\]

The second statement is not merely unproved. In every natural fixed-horizon formulation, it is false: positive integers can shadow the expanding \(2\)-adic fixed point \(-1\) for arbitrarily long periods.

Tao's theorem obtains approximate transport between first-passage distributions of logarithmically random starting values. Its conclusion is correspondingly an almost-all result in logarithmic density, not a statement about every starting integer. [\hypertarget{cit:R18:1}{}\hyperlink{ref:R18:1}{1}]

What follows is a complete derivation of:

\begin{enumerate}
\item an exact Wasserstein contraction extending Tao's stochastic model;
\item exact regeneration uniformly over every finite \(2\)-adic valuation cylinder;
\item a rigorous obstruction showing why neither statement implies pointwise Collatz;
\item a sharp entropy identity explaining the obstruction quantitatively;
\item the exact arithmetic optimal-transport theorem that would still be sufficient to prove Collatz.
\end{enumerate}


\bigskip



\subsubsection*{1. The exact random transport operator}


Let

\[
U(n)=\frac{3n+1}{2^{a(n)}},
\qquad
a(n)=\nu _2(3n+1),
\]

be the accelerated Collatz\textemdash{}or Syracuse\textemdash{}map on positive odd integers.

For every positive integer \(a\), introduce the affine branch

\[
f_a(x)=\frac{3x+1}{2^a}.
\]

The actual map is \(U(n)=f_{a(n)}(n)\). Tao replaces the deterministic valuation \(a(n)\), for typical starting values and suitable time ranges, by a geometric random variable

\[
\mathbb P(A=a)=2^{-a},
\qquad a\geq 1.
\]

For a word \(\mathbf a=(a_1,\ldots ,a_k)\), this gives probability

\[
\mathbb P((A_1,\ldots ,A_k)=\mathbf a)
   =2^{-(a_1+\cdots +a_k)}.
\]

This is exactly the geometric valuation law appearing in Tao's construction. [\hypertarget{cit:R18:2}{}\hyperlink{ref:R18:2}{2}]

Define the Markov operator

\[
\mathsf P\mu
   =\sum_{a=1}^{\infty}2^{-a}(f_a)_{\#}\mu.
\]

Thus \(\mathsf P\delta_x\) is the law of

\[
\frac{3x+1}{2^A}.
\]

This is generally not supported on integers, but it is a perfectly exact stochastic operator on the positive reals, the \(3\)-adics, or an adelic product.


\paragraph*{Theorem 1: Archimedean Wasserstein contraction}


For \(0<\theta<1\), put

\[
d_\theta(x,y)=|x-y|^\theta .
\]

Because \(t\mapsto t^\theta\) is subadditive, \(d_\theta\) is a metric. Let \(W_\theta\) denote the corresponding \(1\)-Wasserstein distance. Then

\[
W_\theta(\mathsf P\mu,\mathsf P\nu)
   \leq \rho_\theta W_\theta(\mu,\nu),
\]

where

\[
\boxed{
\rho_\theta
 =\sum_{a=1}^{\infty}
      2^{-a}\left(\frac{3}{2^a}\right)^\theta
 =\frac{3^\theta}{2^{1+\theta}-1}
 <1.
}
\]


\subparagraph*{Proof}


Take any coupling \((X,Y)\) of \(\mu\) and \(\nu\), and independently sample \(A\) with \(\mathbb P(A=a)=2^{-a}\). Couple the two images using the same \(A\). Then

\[
f_A(X)-f_A(Y)=\frac{3}{2^A}(X-Y),
\]

so

\[
\begin{aligned}
\mathbb E\,d_\theta(f_A(X),f_A(Y))
 &=\mathbb E\left[
       \left(\frac{3}{2^A}\right)^\theta
       |X-Y|^\theta
     \right]\\
 &=\left(
       \sum_{a=1}^{\infty}
       2^{-a}\left(\frac{3}{2^a}\right)^\theta
   \right)
   \mathbb E|X-Y|^\theta\\
 &=\rho_\theta\mathbb E\,d_\theta(X,Y).
\end{aligned}
\]

Taking the infimum over all couplings proves the Wasserstein inequality.

It remains to show \(\rho_\theta<1\). Since \(t^\theta\) is strictly concave and \(2=(1+3)/2\),

\[
2^\theta>\frac{1+3^\theta}{2}.
\]

Therefore

\[
2^{1+\theta}>1+3^\theta,
\]

and hence

\[
2^{1+\theta}-1>3^\theta.
\]

Thus \(\rho_\theta<1\). \ensuremath{\blacksquare}

The best exponent is

\[
\theta_*
 =\log_2\left(\frac{\log 3}{\log(3/2)}\right)-1
 \approx 0.4380326593,
\]

and the optimal contraction factor is

\[
\boxed{
\rho_*=\rho_{\theta_*}
       \approx 0.9465045768.
}
\]

At the ordinary Wasserstein exponent \(\theta=1\),

\[
\rho_1=\frac{3}{4-1}=1.
\]

Thus the stochastic Syracuse map is only nonexpansive in ordinary \(W_1\); the strict contraction appears after snowflaking the metric.


\bigskip



\subsubsection*{2. Exact \(3\)-adic contraction}


On \(\mathbb Z_3\), use

\[
d_3(x,y)=|x-y|_3=3^{-\nu_3(x-y)}.
\]

Since \(2\) is a \(3\)-adic unit,

\[
\begin{aligned}
d_3(f_a(x),f_a(y))
 &=\left|\frac{3(x-y)}{2^a}\right|_3\\
 &=\frac13d_3(x,y).
\end{aligned}
\]

Consequently,

\[
\boxed{
W_{d_3}(\mathsf P\mu,\mathsf P\nu)
   \leq \frac13W_{d_3}(\mu,\nu).
}
\]

This contraction does not merely hold on average over \(a\); every individual branch contracts by exactly \(1/3\) in the \(3\)-adic metric.

Tao's Syracuse random variables satisfy precisely the recursion

\[
\mathbf S_{k+1}
   =\frac{3\mathbf S_k+1}{2^{A_{k+1}}},
\qquad A_{k+1}\sim\operatorname{Geom}(2),
\]

and their inverse-limit law is the stationary distribution of this \(3\)-adic Markov chain. [\hyperlink{ref:R18:2}{2}, \hypertarget{cit:R18:3}{}\hyperlink{ref:R18:3}{3}]

On the product state space \(\mathbb R_{\geq0}\times\mathbb Z_3\), set

\[
D_\theta((x,z),(y,w))
   =|x-y|^\theta+\lambda |z-w|_3.
\]

The same synchronous coupling gives

\[
W_{D_\theta}(\mathsf P\mu,\mathsf P\nu)
 \leq
 \max\{\rho_\theta,1/3\}
 W_{D_\theta}(\mu,\nu).
\]

Since \(\rho_\theta>1/3\), the optimal factor remains approximately \(0.9465\).

This is an \textbf{exact arithmetic optimal-transport theorem}. The problem is that it governs a randomized affine map, not the deterministic orbit of one integer.


\bigskip



\subsubsection*{3. Exact transport uniformly over every valuation cylinder}


There is an even stronger exact result on the source side.

For a valuation word

\[
w=(a_1,\ldots ,a_k),
\qquad
A_j=a_1+\cdots +a_j,
\qquad
A=A_k,
\]

define

\[
C_w
 =\sum_{j=0}^{k-1}3^{k-1-j}2^{A_j},
\qquad A_0=0.
\]

Then every realization satisfies

\[
\boxed{
U^k(n)=\frac{3^kn+C_w}{2^A}.
}
\]

Let

\[
\mathcal C_w
 =\{n\in 2\mathbb Z_2+1:
       (a(n),a(U(n)),\ldots ,a(U^{k-1}(n)))=w\}
\]

be the \(2\)-adic valuation cylinder determined by \(w\).


\paragraph*{Theorem 2: Exact cylinder regeneration}


For every valuation word \(w\):

\begin{enumerate}
\item \(\mathcal C_w\) is one odd residue class modulo \(2^{A+1}\);
\item its normalized Haar measure is \(2^{-A}\);
\item conditional Haar measure on \(\mathcal C_w\) is pushed by \(U^k\) exactly to Haar measure on \(2\mathbb Z_2+1\);
\item conditioned on any finite past valuation word, all future valuations are independent copies of \(\operatorname{Geom}(2)\).
\end{enumerate}


\subparagraph*{Proof}


For one valuation \(a_1\),

\[
\nu_2(3n+1)=a_1
\]

is equivalent to

\[
3n+1\equiv 2^{a_1}\pmod{2^{a_1+1}},
\]

which determines one odd residue class modulo \(2^{a_1+1}\), because \(3\) is invertible modulo powers of \(2\).

Inductively, suppose the first \(j\) valuations determine one class

\[
n=n_j+2^{A_j+1}t.
\]

The affine formula gives

\[
U^j(n)=x_j+2\cdot3^jt.
\]

Requiring

\[
\nu_2(3U^j(n)+1)=a_{j+1}
\]

determines one residue class for \(t\) modulo \(2^{a_{j+1}}\). Hence it refines the original cylinder to one class modulo

\[
2^{A_j+1+a_{j+1}}=2^{A_{j+1}+1}.
\]

This proves the first statement.

There are \(2^A\) odd residue classes modulo \(2^{A+1}\), so normalized Haar measure gives \(\mathcal C_w\) mass \(2^{-A}\).

Now write

\[
n=n_w+2^{A+1}t,
\qquad t\in\mathbb Z_2.
\]

Then

\[
\begin{aligned}
U^k(n)
 &=\frac{3^kn_w+C_w}{2^A}
     +2\cdot3^kt\\
 &=x_w+2\cdot3^kt.
\end{aligned}
\]

Multiplication by the odd \(2\)-adic unit \(3^k\) preserves Haar measure on \(\mathbb Z_2\). Therefore \(U^k\) maps conditional Haar on \(\mathcal C_w\) exactly to Haar on the odd \(2\)-adics.

The future valuation sequence is consequently fresh i.i.d. geometric data, independent of the specified prefix. \ensuremath{\blacksquare}

Tao notes this Haar invariance and the exact i.i.d. geometric valuation law on odd \(2\)-adic integers. [\hyperlink{ref:R18:2}{2}, \hypertarget{cit:R18:4}{}\hyperlink{ref:R18:4}{4}]

Thus the requested statement


\begin{quote}
``transport uniformly over every starting cylinder''

\end{quote}


is already true in the strongest possible \textbf{conditional-measure sense}.

But it does not imply a pointwise theorem. A single ordinary integer has Haar measure zero. Uniform probabilistic regeneration after every finite prefix does not prevent one exceptional infinite ray from lying in the exceptional part of every nested cylinder.

The Bernoulli sequence

\[
1,1,1,1,\ldots
\]

is the fundamental exceptional ray.


\bigskip



\subsubsection*{4. The exact deterministic obstruction}


For every \(k\geq1\), define the all-\(1\) cylinder

\[
\mathcal E_k
 =\{n>0\text{ odd}:n\equiv-1\pmod{2^{k+1}}\}.
\]


\paragraph*{Theorem 3: Arbitrarily long expanding cylinders}


Every \(n\in\mathcal E_k\) has first \(k\) Syracuse valuations equal to \(1\), and

\[
\boxed{
U^j(n)+1
  =\left(\frac32\right)^j(n+1)
  \qquad(0\leq j\leq k).
}
\]

In particular,

\[
U^j(n)>n
\qquad(1\leq j\leq k).
\]


\subparagraph*{Proof}


Write

\[
n+1=2^{k+1}t.
\]

Suppose the formula holds at step \(j<k\). Then

\[
U^j(n)+1=3^j2^{k+1-j}t,
\]

and therefore

\[
\begin{aligned}
3U^j(n)+1
 &=3(U^j(n)+1)-2\\
 &=3^{j+1}2^{k+1-j}t-2\\
 &=2\left(3^{j+1}2^{k-j}t-1\right).
\end{aligned}
\]

Since \(k-j\geq1\), the expression in parentheses is odd. Thus

\[
\nu_2(3U^j(n)+1)=1,
\]

and

\[
U^{j+1}(n)+1
 =\frac32\bigl(U^j(n)+1\bigr).
\]

Induction proves the formula. \ensuremath{\blacksquare}


\paragraph*{Consequence for Wasserstein contraction}


For two numbers \(n,m\in\mathcal E_k\),

\[
U^k(n)-U^k(m)
   =\left(\frac32\right)^k(n-m).
\]

Hence for every \(0<\theta\leq1\),

\[
d_\theta(U^k(n),U^k(m))
  =\left(\frac32\right)^{k\theta}d_\theta(n,m).
\]

Thus the deterministic operator expands the snowflake metric by an arbitrarily large factor.

For the shifted logarithmic metric

\[
d_{\log}(n,m)
 =|\log(n+1)-\log(m+1)|,
\]

we instead obtain exact preservation:

\[
d_{\log}(U^k(n),U^k(m))
 =d_{\log}(n,m).
\]

So even logarithmic renormalization cannot produce a strict factor \(q<1\).

More generally, every distance from \(1\) which is strictly increasing with the ordinary size satisfies

\[
d(U^j(n),1)>d(n,1)
\qquad(1\leq j\leq k).
\]

Because \(k\) is arbitrary, no deterministic transport rule with a uniformly bounded positive stopping time can contract on every starting cylinder.


\paragraph*{Why finite residue and carry enrichment does not help}


Suppose the enriched state contains residues modulo fixed powers of \(2\) and \(3\), together with finitely many carry symbols. Choose two distinct positive integers \(n,m\) satisfying simultaneously

\[
n\equiv m\equiv-1\pmod{2^{k+1}}
\]

and

\[
n\equiv m\pmod{2^r3^s}
\]

for moduli large enough to determine all the stored coordinates and their evolution through the chosen horizon. Such \(n,m\) exist by the Chinese remainder theorem.

Their residue and carry components agree, while their archimedean separation is multiplied by \((3/2)^k\). Therefore finite arithmetic enrichment cannot convert this cylinder into a contraction.

The obstruction is not missing carry information. It is a genuine expanding orbit segment.


\bigskip



\subsubsection*{5. The \(-1\) fixed point and terminal locality}


The all-\(1\) cylinders converge \(2\)-adically to

\[
-1.
\]

The Syracuse map extends to odd \(2\)-adic integers, and

\[
U(-1)=\frac{3(-1)+1}{2}=-1.
\]

At the same time,

\[
U(1)=\frac{3+1}{4}=1.
\]

Thus the natural \(2\)-adic completion contains two distinct fixed points, \(1\) and \(-1\). A strict contraction on a metric space separating these points is impossible, since it would imply

\[
d(1,-1)
 =d(U(1),U(-1))
 \leq q\,d(1,-1),
\qquad q<1.
\]

The carry data does not exclude the \(-1\) ray. For the all-\(1\) word \(1^k\),

\[
C_{1^k}=3^k-2^k,
\]

and hence

\[
C_{1^k}+2^k=3^k.
\]

So the all-\(1\) ray has the strongest possible terminal \(3\)-adic compatibility.

Your terminal-locality theorem states that, for a word ending in a suffix \(v\) of length \(r\),

\[
3^r\mid X+1
\iff
3^r\mid C_v+2^{A(v)},
\]

so endpoint divisibility is determined entirely by that suffix. [\hypertarget{cit:R18:5}{}\hyperlink{ref:R18:5}{5}]

For \(v=1^r\), the right-hand side vanishes modulo \(3^r\) exactly. Thus terminal carry locality recognizes the \(-1\) shadow as a perfectly coherent state. It cannot distinguish the genuine positive integers that merely shadow \(-1\) for a finite time from the infinite \(2\)-adic ray itself.

That distinction is global and archimedean.


\bigskip



\subsubsection*{6. The adelic product-formula obstruction}


The tension between exactness and contraction can be expressed in one identity.

For a fixed branch \(f_a\),

\[
f_a(x)-f_a(y)=\frac{3}{2^a}(x-y).
\]

The scale factors at the three relevant places are

\[
\left|\frac3{2^a}\right|_\infty=\frac3{2^a},
\]

\[
\left|\frac3{2^a}\right|_2=2^a,
\]

and

\[
\left|\frac3{2^a}\right|_3=\frac13.
\]

Therefore

\[
\boxed{
\left|\frac3{2^a}\right|_\infty
\left|\frac3{2^a}\right|_2
\left|\frac3{2^a}\right|_3
=1.
}
\]

The \(3\)-adic coordinate contracts by \(1/3\). The archimedean coordinate contracts on average. But the \(2\)-adic coordinate\textemdash{}the coordinate that determines which valuation \(a\) is legally available\textemdash{}expands by \(2^a\).

Tao's stochastic model averages away this expanding source coordinate and replaces it by geometric noise. Restoring enough \(2\)-adic information to make the branch exact restores precisely the expansion that the random operator omitted.

In this sense, the failure of exact uniform contraction is enforced by the product formula.


\bigskip



\subsubsection*{7. A sharp entropy obstruction}


There is also a quantitative counting version of the same phenomenon.

Consider a length-\(k\) inverse valuation word

\[
b=(b_1,\ldots ,b_k),
\qquad
B=b_1+\cdots +b_k.
\]

If \(Y\) is the ancestor and \(X\) the endpoint, then

\[
2^BX=3^kY+C_b,
\]

so

\[
Y=\frac{2^BX-C_b}{3^k}.
\]

The inverse slope is

\[
\frac{2^B}{3^k}.
\]

It is nonexpanding precisely when

\[
2^B\leq3^k,
\qquad
B\leq \beta k,
\qquad
\beta=\log_2 3.
\]

For a fixed word \(b\), integrality of \(Y\) requires

\[
2^BX\equiv C_b\pmod{3^k}.
\]

Since \(2\) is invertible modulo \(3^k\), each word is compatible with exactly one endpoint residue class modulo \(3^k\).


\paragraph*{Theorem 4: Nonexpanding inverse words cannot cover all residues}


Let

\[
M_k=\lfloor\beta k\rfloor.
\]

The number of length-\(k\) positive valuation words with \(B\leq M_k\) is

\[
N_k
 =\sum_{B=k}^{M_k}\binom{B-1}{k-1}
 =\binom{M_k}{k}.
\]

Consequently, these words can cover at most \(N_k\) of the \(3^k\) endpoint residue classes. Moreover,

\[
\frac{N_k}{3^k}
 =q_*^{\,k+o(k)},
\]

where

\[
\boxed{
q_*
 =\frac{\beta^\beta}
        {3(\beta-1)^{\beta-1}}
 \approx0.9465045768<1.
}
\]


\subparagraph*{Proof}


The number of positive compositions of \(B\) into \(k\) parts is

\[
\binom{B-1}{k-1}.
\]

The hockey-stick identity gives

\[
\sum_{B=k}^{M_k}\binom{B-1}{k-1}
 =\binom{M_k}{k}.
\]

Because each word covers at most one residue class modulo \(3^k\), at most \(N_k\) classes are covered.

Stirling's formula gives

\[
\binom{\beta k+o(k)}{k}
 =
 \left(
 \frac{\beta^\beta}
      {(\beta-1)^{\beta-1}}
 \right)^{k+o(k)}.
\]

Dividing by \(3^k\) proves the result. \ensuremath{\blacksquare}

Thus an exponentially large proportion of endpoint residues have no nonexpanding inverse word at a fixed depth.

A branch with slope \(>1\) might still produce a smaller ancestor for one particular small endpoint because of its affine constant. But it cannot provide uniform descent for arbitrarily large values. Therefore no fixed-depth, slope-only inverse transport can cover all arithmetic cylinders.


\paragraph*{Exact duality with stochastic Wasserstein contraction}


Remarkably,

\[
\boxed{
q_*=\min_{0<\theta<1}\rho_\theta.
}
\]

Indeed,

\[
\rho_\theta=\frac{3^\theta}{2^{1+\theta}-1},
\]

and its minimizer satisfies

\[
2^{1+\theta_*}
   =\frac{\log3}{\log(3/2)}
   =\frac{\beta}{\beta-1}.
\]

Substitution gives

\[
\rho_{\theta_*}
 =\frac{\beta^\beta}
        {3(\beta-1)^{\beta-1}}
 =q_*.
\]

This equality is not a numerical accident. It is a Chernoff\textendash{}entropy duality:

\begin{itemize}
\item \(\rho_*\) is the best possible contraction rate after averaging over geometric valuation words;
\item the same number \(q_*\) is the exponential density of the low-exponent inverse words that could possibly give uniform nonexpanding coverage.
\end{itemize}

The stochastic contraction rate is exactly the deterministic coverage deficit.

That is the quantitative form of the averaging-to-pointwise gap.


\bigskip



\subsubsection*{8. An exact optimal-transport criterion that would prove Collatz}


There is nevertheless a clean pointwise transport theorem.

Call a positive odd integer \(y\) a \textbf{coalescing ancestor} of \(n\) if their Syracuse orbits eventually meet.


\paragraph*{Theorem 5: Pointwise coalescing-transport criterion}


Fix \(0<\theta<1\). Suppose that for every positive odd \(n>1\), one can construct a probability measure \(K_n\) supported on positive coalescing ancestors of \(n\) such that

\[
\boxed{
W_{d_\theta}(K_n,\delta_1)
 <
d_\theta(n,1).
}
\]

Then the Collatz conjecture is true.


\subparagraph*{Proof}


Suppose Collatz is false, and let \(m\) be the least positive odd counterexample.

Every positive coalescing ancestor \(y\) of \(m\) is itself a counterexample. Otherwise \(y\) would reach \(1\), and once its orbit merged with the orbit of \(m\), the orbit of \(m\) would also reach \(1\). Minimality therefore gives

\[
y\geq m
\]

for every \(y\) in the support of \(K_m\).

Since \(d_\theta(y,1)=(y-1)^\theta\) is increasing,

\[
\begin{aligned}
W_{d_\theta}(K_m,\delta_1)
 &=\int (y-1)^\theta\,dK_m(y)\\
 &\geq(m-1)^\theta\\
 &=d_\theta(m,1),
\end{aligned}
\]

contradicting the assumed strict contraction. \ensuremath{\blacksquare}

A uniform factor

\[
W_{d_\theta}(K_n,\delta_1)
 \leq q\,d_\theta(n,1),
\qquad q<1,
\]

would be more than sufficient, but strict pointwise decrease is already enough.


\bigskip



\subsubsection*{9. The exact transport inequality in your local\textendash{}global coordinates}


Your existing competitor construction supplies the right support for \(K_n\).

For a word split as

\[
w=u\tau,
\]

a compatible competing suffix \(b\) produces an ancestor \(Y_b\) satisfying

\[
\boxed{
Y_b-R
 =\frac{\Sigma(u,\tau,c)+h(\tau,b)}{2^c}.
}
\]

Here:

\begin{itemize}
\item the compatible fibre and \(h(\tau,b)\) are terminal and \(3\)-adic;
\item \(\Sigma(u,\tau,c)\) is the prefix-dependent archimedean surplus.
\end{itemize}

This exact decomposition is already established in your framework. [\hypertarget{cit:R18:6}{}\hyperlink{ref:R18:6}{6}]

Put a probability distribution \(p_b\) on a nonempty compatible fibre, supported only on positive ancestors. Then

\[
K_R=\sum_b p_b\delta_{Y_b}.
\]

Since \(t^\theta\) is concave,

\[
\begin{aligned}
W_{d_\theta}(K_R,\delta_1)
 &=\sum_b p_b(Y_b-1)^\theta\\
 &=
 \sum_b p_b
 \left(
 R-1+
 \frac{\Sigma+h(\tau,b)}{2^c}
 \right)^\theta\\
 &\leq
 \left(
 R-1+
 \frac{\Sigma+\sum_b p_b h(\tau,b)}{2^c}
 \right)^\theta.
\end{aligned}
\]

Therefore the exact arithmetic transport condition

\[
\boxed{
\Sigma(u,\tau,c)
   +\sum_b p_bh(\tau,b)<0
}
\]

implies

\[
W_{d_\theta}(K_R,\delta_1)
 <d_\theta(R,1).
\]

This is the optimal-transport form of the missing pointwise synchronization theorem.

Your current stronger target is

\[
\Sigma(u,\tau,c)+h_{\max}(\tau,B)<0,
\]

which forces even the largest compatible ancestor below \(R\). The established analysis shows that the terminal state determines \(h_{\max}\), while the prefix controls \(\Sigma\); no suffix-local rule currently synchronizes their signs along every positive stationary ray. [\hypertarget{cit:R18:7}{}\hyperlink{ref:R18:7}{7}]


\bigskip



\subsection*{Final conclusion}


The proposed route does \textbf{not} presently give a valid full proof of Collatz. Instead, it gives the following sharp diagnosis:

\[
\begin{array}{c}
\text{Random geometric valuations}\\
\Downarrow\\
\text{exact Wasserstein contraction}
\end{array}
\qquad\text{but}\qquad
\begin{array}{c}
\text{exact \(2\)-adic branch legality}\\
\Downarrow\\
\text{arbitrarily long expanding cylinders}.
\end{array}
\]

Uniform conditional transport over every finite valuation cylinder is already exact, but it remains measure-theoretic. Uniform deterministic contraction over every cylinder is false because of the all-\(1\) cylinders shadowing \(-1\). The best stochastic contraction rate, approximately \(0.9465045768\), exactly equals the exponential deficit of low-slope inverse words available to cover all \(3\)-adic residues.

A viable optimal-transport proof would therefore have to abandon fixed-horizon uniform contraction. It would require a \textbf{nonuniform, unbounded stopping-time transport}, defined specifically on rays whose nested \(2\)-adic representatives stabilize at a positive ordinary integer, and it would have to establish

\[
\Sigma+\mathbb E h<0
\]

at some finite stage of every such ray. That is not a consequence of Tao's mixing or of terminal carry locality; it is precisely the remaining global synchronization problem. Claiming a Collatz proof without that theorem would leave the central pointwise gap unresolved.
\subsection*{References}
\begin{enumerate}[leftmargin=2.2em]
\item \hypertarget{ref:R18:1}{}\hyperlink{cit:R18:1}{\texttt{turn515498view0}}
\item \hypertarget{ref:R18:2}{}\hyperlink{cit:R18:2}{\texttt{turn664973view0}}
\item \hypertarget{ref:R18:3}{}\hyperlink{cit:R18:3}{\texttt{turn664973view2}}
\item \hypertarget{ref:R18:4}{}\hyperlink{cit:R18:4}{\texttt{turn664973view1}}
\item \hypertarget{ref:R18:5}{}\hyperlink{cit:R18:5}{\texttt{turn3file1 L204-L227}}
\item \hypertarget{ref:R18:6}{}\hyperlink{cit:R18:6}{\texttt{turn3file0 L44-L100}}
\item \hypertarget{ref:R18:7}{}\hyperlink{cit:R18:7}{\texttt{turn3file0 L164-L178}}
\end{enumerate}

\clearpage
\section{Round 19: approach 14 --- Mixed 2-adic/3-adic Kakeya or decoupling}
\subsection*{Prompt}
\begin{quote}
Complete a full proof of the Collatz conjecture using the following approach:

\begin{enumerate}
\item Mixed 2-adic/3-adic Kakeya or decoupling
\end{enumerate}

\textbf{Speculative \textperiodcentered{} P}

A length-j critical word defines both a thin 2-adic seed tube and a thin 3-adic endpoint tube. An infinite counterexample is a nested intersection of such tubes, together with a real drift restriction.

Develop a multi-place incidence theorem showing that tubes with the necessary directions cannot have an infinite coherent intersection. This would be an arithmetic analogue of ruling out a Kakeya configuration by dimension or decoupling estimates.

Take as long as you need.
\end{quote}
\subsection*{Response}

\subsection*{Result}


I did \textbf{not} obtain a valid proof of the Collatz conjecture. The conjecture remains open as of August 13, 2026. [\hypertarget{cit:R19:1}{}\hyperlink{ref:R19:1}{1}]

More decisively, the proposed Kakeya theorem cannot hold in its natural form. Exact calculation reveals three structural obstructions:

\begin{enumerate}
\item At fixed length and total valuation, all Collatz incidence tubes are parallel.
\item The \(3\)-adic endpoint cylinders of successive prefixes are not nested.
\item Even after correcting the construction to a two-sided \(2\)-adic/\(3\)-adic state, critical coherent configurations exist abundantly. What is exceptional is not coherence but \textbf{diagonal realization by one positive ordinary integer}.
\end{enumerate}

Thus the approach does not reduce Collatz to an existing Kakeya or decoupling estimate. It reduces it to a new signed, height-sensitive synchronization theorem. I give that reduction exactly below.


\bigskip



\subsubsection*{1. Exact word geometry}


Work with the accelerated odd map

\[
U(x)=\frac{3x+1}{2^{\nu _2(3x+1)}}.
\]

Let

\[
w=(a_0,\ldots ,a_{j-1}),\qquad a_i\geq 1,
\]

be a valuation word. Define

\[
A_0=0,\qquad A_i=\sum_{t=0}^{i-1}a_t,\qquad A=A_j,
\]

and

\[
C_w=\sum_{i=0}^{j-1}3^{j-1-i}2^{A_i}.
\]

Then every realization of \(w\), with seed \(x\) and endpoint \(y=U^j(x)\), satisfies

\[
2^A y=3^j x+C_w. \tag{1}
\]

Moreover, \(w\) occurs from precisely one odd residue class

\[
x\equiv \rho _2(w)
 :=3^{-j}(2^A-C_w)
 \pmod {2^{A+1}}, \tag{2}
\]

while its endpoints occupy precisely one odd \(3\)-adic class

\[
y\equiv \rho _3(w)
 :=2^{-A}C_w
 \pmod {3^j}. \tag{3}
\]

Conversely, every positive odd endpoint in (3) has a positive inverse ancestor along \(w\). These are exact identities already established in the valuation-word framework. [\hypertarget{cit:R19:2}{}\hyperlink{ref:R19:2}{2}]

Thus the natural local tubes are

\[
\mathcal S_w=\rho _2(w)+2^{A+1}\mathbb Z_2
\]

on the seed side and

\[
\mathcal E_w=\rho _3(w)+3^j\mathbb Z_3
\]

on the endpoint side. The affine variety coupling them is the rational line

\[
L_w:\quad 2^A Y-3^jX=C_w. \tag{4}
\]


\bigskip



\subsubsection*{2. The parallel-tube collapse theorem}



\paragraph*{Theorem 1}


Fix a length \(j\) and total valuation \(A\). For every word \(w\) of length \(j\) and sum \(A\):

\begin{enumerate}
\item The odd integral points of \(L_w\) form an arithmetic line \[
   (X,Y)
   =
   (R_w,Y_w)
   +t\bigl(2^{A+1},\,2\cdot3^j\bigr),
   \qquad t\in\mathbb Z,
   \tag{5}
   \] where \((R_w,Y_w)\) is any one odd integral point.

\item Consequently, all words with the same \((j,A)\) have exactly the same projective direction \[
   d_{j,A}=[2^{A+1}:2\cdot3^j].
   \tag{6}
   \]

\item Under the normalization \[
   u=\frac{X}{2^{A+1}},
   \qquad
   v=\frac{Y}{2\cdot3^j},
   \] the line becomes \[
   v-u=\delta_w,
   \qquad
   \delta_w=\frac{C_w}{2^{A+1}3^j},
   \tag{7}
   \] and \[
   0<\delta_w
   \leq
   \frac1{2^{j+1}}-\frac1{2\cdot3^j}
   <2^{-j-1}. \tag{8}
   \]
\end{enumerate}

Hence, after natural normalization, the entire family consists of parallel slope-one lines inside an exponentially thin diagonal strip.


\paragraph*{Proof}


Since

\[
\gcd(2^A,3^j)=1,
\]

all integral solutions to (4), starting from one solution \((R_w,Y_w)\), are

\[
(X,Y)=(R_w+2^At,\;Y_w+3^jt),
\qquad t\in\mathbb Z.
\]

Both \(R_w\) and \(Y_w\) are odd. Since \(3^j\) is odd, \(Y_w+3^jt\) is odd exactly when \(t\) is even. Writing \(t=2s\) gives (5).

Conversely, if \((X,Y)\) is an odd integral point of \(L_w\), then modulo \(2^{A+1}\),

\[
3^jX+C_w
=
2^AY
\equiv 2^A\pmod {2^{A+1}},
\]

so

\[
X\equiv3^{-j}(2^A-C_w)\pmod {2^{A+1}}.
\]

By the exact realization-class theorem, every positive such \(X\) realizes \(w\). Thus (5) is not merely the integral solution set of the final affine equation; its positive part is exactly the realization line of \(w\).

Equation (7) follows by dividing (4) by \(2^{A+1}3^j\).

Finally,

\[
\frac{C_w}{2^A}
=
\sum_{i=0}^{j-1}
\frac{3^{j-1-i}}{2^{A-A_i}}.
\]

Because every \(a_t\geq1\),

\[
A-A_i\geq j-i.
\]

Therefore

\[
\frac{C_w}{2^A}
\leq
\sum_{i=0}^{j-1}
\frac{3^{j-1-i}}{2^{j-i}}
=
\sum_{r=1}^{j}\frac{3^{r-1}}{2^r}
=
\left(\frac32\right)^j-1.
\]

Dividing by \(2\cdot3^j\) gives

\[
\delta_w
\leq
\frac{(3/2)^j-1}{2\cdot3^j}
=
\frac1{2^{j+1}}-\frac1{2\cdot3^j}.
\]

This proves the theorem. \(\square\)


\bigskip



\subsubsection*{3. Why this is not Kakeya geometry}


Let

\[
\alpha=\log_2 3.
\]

Suppose ``critical'' means that

\[
|A-\alpha j|\leq K.
\]

At scale \(j\), the number of possible totals \(A\), and hence the number of possible directions (6), is at most

\[
2K+2.
\]

For \(K=O(1)\), the critical family has only \(O(1)\) directions. Even for \(K=o(j)\), it has only sublinear or subexponential directional complexity.

By contrast, finite-ring and \(p\)-adic Kakeya estimates concern sets containing rich lines in every direction, or at least in a positive fraction of projective directions. [\hypertarget{cit:R19:3}{}\hyperlink{ref:R19:3}{3}] In the usual discretized Euclidean formulation, the tubes point in a full collection of separated directions. [\hypertarget{cit:R19:4}{}\hyperlink{ref:R19:4}{4}]

The Collatz family instead has exponentially many \textbf{parallel translates}.

Indeed, the number of length-\(j\) valuation words with sum \(A\) is

\[
\#\mathcal W_{j,A}
=
\binom{A-1}{j-1}. \tag{9}
\]

When \(A=\alpha j+O(1)\), Stirling's formula gives

\[
\log_2\binom{A-1}{j-1}
=
\alpha H_2(1/\alpha)j+O(\log j),
\]

where \(H_2\) is the binary entropy. Numerically,

\[
\alpha H_2(1/\alpha)
\approx1.5056438879.
\]

Thus there are approximately

\[
2^{1.50564j}
\]

critical words, but they are arranged among only \(O(1)\) directions.

This is almost exactly the opposite of a Kakeya configuration: Kakeya has many directions and attempts to compress their union; Collatz has huge multiplicity in a tiny directional set and needs to extract information from the intercepts \(C_w\).

There is a second counting obstruction. The total number of words satisfying \(A(w)\leq A\) is

\[
\sum_{B=j}^{A}\binom{B-1}{j-1}
=
\binom{A}{j}. \tag{10}
\]

For \(A=\alpha j+O(1)\),

\[
\frac{\binom{A}{j}}{3^j}
=
2^{-\kappa j+O(\log j)},
\qquad
\kappa
=
\alpha\bigl(1-H_2(1/\alpha)\bigr)
\approx0.0793186128. \tag{11}
\]

So the collection of critical words is exponentially sparse relative to the \(3^j\) possible endpoint residues. Pure pigeonhole counting does not even force two critical words to share an endpoint class.


\bigskip



\subsubsection*{4. The endpoint tubes are not nested}


This is a separate flaw in the proposed picture.

If \(a\) is appended to \(w\), then

\[
C_{wa}=3C_w+2^A.
\]

Consequently,

\[
\begin{aligned}
\rho _3(wa)
&\equiv
(3C_w+2^A)2^{-(A+a)}
\pmod {3^{j+1}}\\
&\equiv
2^{-a}\bigl(1+3\rho _3(w)\bigr)
\pmod {3^{j+1}}. 
\end{aligned}
\tag{12}
\]

Ordinary nesting would require

\[
\rho _3(wa)\equiv\rho _3(w)\pmod {3^j},
\]

which is generally false.

For example, take \(w=(1)\). Then

\[
C_w=1,\qquad A=1,\qquad
\rho _3(w)\equiv2\pmod3.
\]

Append \(a=2\). For \(wa=(1,2)\),

\[
C_{wa}=5,\qquad A(wa)=3,
\]

and hence

\[
\rho _3(wa)
\equiv5\cdot2^{-3}
\equiv5\cdot8
\equiv4\pmod9.
\]

But \(4\equiv1\pmod3\), so

\[
4+9\mathbb Z_3
\not\subset
2+3\mathbb Z_3.
\]

They are disjoint.

This is not paradoxical: \(\mathcal E_w\) describes the possible value of \(x_j\), while \(\mathcal E_{wa}\) describes the possible value of \(x_{j+1}\). They are tubes for different time coordinates.

Terminal locality says that the residue of an endpoint modulo \(3^r\) is determined by the final \(r\) valuations. It does \textbf{not} say that the endpoint residues of growing prefixes form an inverse system under ordinary reduction. [\hypertarget{cit:R19:5}{}\hyperlink{ref:R19:5}{5}]

Therefore the statement


\begin{quote}
``An infinite counterexample is a nested intersection of a \(2\)-adic seed tube and a \(3\)-adic endpoint tube''

\end{quote}


is not correct as written. Only the seed cylinders are nested in one fixed coordinate.


\bigskip



\subsubsection*{5. Critical nested \(2\)-adic intersections exist automatically}


For an infinite valuation word

\[
a=(a_0,a_1,\ldots),
\]

let \(w_j=(a_0,\ldots,a_{j-1})\). Its seed cylinders satisfy

\[
\mathcal S_{w_{j+1}}\subseteq\mathcal S_{w_j}.
\]

Every finite word is realized, so these cylinders are nonempty. Their radii are

\[
2^{-(A_j+1)},
\]

which tend to zero because \(A_j\geq j\). Completeness of \(\mathbb Z_2\) therefore gives

\[
\bigcap_{j\geq1}\mathcal S_{w_j}
=
\{\xi_a\}
\]

for one unique \(2\)-adic integer \(\xi_a\).

This happens for \textbf{every} infinite valuation word.

It also happens under an exact bounded-drift condition. Let

\[
a_i
=
\left\lfloor(i+1)\alpha\right\rfloor
-
\left\lfloor i\alpha\right\rfloor,
\qquad
\alpha=\log_2 3.
\tag{13}
\]

Then \(a_i\in\{1,2\}\) and

\[
A_j=\sum_{i=0}^{j-1}a_i=\lfloor j\alpha\rfloor.
\]

Hence

\[
-1<A_j-j\log_2 3\leq0 \tag{14}
\]

at every scale, and

\[
1\leq \frac{3^j}{2^{A_j}}<2. \tag{15}
\]

Thus there exists a perfectly coherent \(2\)-adic ray whose archimedean multiplicative drift is uniformly critical at every depth.

It almost certainly does not correspond to a positive integer, but that is precisely the missing arithmetic condition. Tube coherence plus critical drift does not distinguish positive integers from generic \(2\)-adic points.

This conclusion is consistent with the general \(2\)-adic coding picture in which infinite orbit codes determine \(2\)-adic integers. [\hypertarget{cit:R19:6}{}\hyperlink{ref:R19:6}{6}]


\bigskip



\subsubsection*{6. The corrected mixed-place construction}


There is a natural two-sided adelic formulation, but it shows why the problem is synchronization rather than Kakeya.

Let

\[
a=(a_i)_{i\in\mathbb Z}
\]

be a bi-infinite valuation sequence.


\paragraph*{Future \(2\)-adic state}


The future symbols

\[
a_0,a_1,a_2,\ldots
\]

determine a unique

\[
\xi_2(a)\in\mathbb Z_2
\]

through the nested seed cylinders.


\paragraph*{Past \(3\)-adic state}


For \(a\geq1\), define

\[
f_a(z)=\frac{3z+1}{2^a}.
\]

This preserves \(\mathbb Z_3\), and

\[
|f_a(z)-f_a(z')|_3
=
\frac13|z-z'|_3. \tag{16}
\]

Consequently,

\[
\xi_3(a)
=
\lim_{r\to\infty}
f_{a_{-1}}\circ f_{a_{-2}}\circ\cdots
\circ f_{a_{-r}}(0) \tag{17}
\]

exists in \(\mathbb Z_3\). Indeed, changing the starting point at time \(-r\) affects the time-zero result by at most \(3^{-r}\).

Thus every bi-infinite symbolic sequence determines a local two-ended state

\[
\Xi(a)=\bigl(\xi_2(a),\xi_3(a)\bigr)
\in\mathbb Z_2\times\mathbb Z_3. \tag{18}
\]

We can choose \(a\) to be a bi-infinite mechanical word satisfying the bounded-drift relation (14) on every interval. Hence mixed \(2\)-adic/\(3\)-adic coherent critical states are not forbidden; they are plentiful.

The special condition relevant to Collatz is

\[
\Xi(a)=(n,n)
\]

for one positive ordinary integer \(n\), under the diagonal embeddings

\[
\mathbb Z_{>0}\hookrightarrow\mathbb Z_2,
\qquad
\mathbb Z_{>0}\hookrightarrow\mathbb Z_3,
\]

together with positivity of the relevant inverse branches and the least-counterexample inequalities.

A dimension theorem ordinarily says that a set is large or small. It does not say whether that set contains one prescribed arithmetic diagonal point. Even a measure-zero or low-dimensional adelic set may contain infinitely many rational points, while a full-dimensional set may avoid all positive diagonal integers.

So the appropriate target is an \textbf{adelic diagonal-avoidance theorem with height and sign}, not a Kakeya dimension estimate.


\bigskip



\subsubsection*{7. Why non-Archimedean decoupling also degenerates}


Known non-Archimedean decoupling results concern Fourier support near curved varieties, such as a \(p\)-adic parabola

\[
\{(t,t^2)\},
\]

partitioned into small frequency caps. [\hypertarget{cit:R19:7}{}\hyperlink{ref:R19:7}{7}]

The two-dimensional Collatz incidence variety is affine:

\[
2^AY-3^jX=C_w.
\]

At fixed \((j,A)\), it has:

\begin{itemize}
\item no curvature;
\item no directional separation;
\item one normal vector \((-3^j,2^A)\);
\item word dependence only in the scalar intercept \(C_w\).
\end{itemize}

Decoupling an affine line simply reduces to orthogonality between frequency intervals; it cannot reveal the arithmetic ordering of the constants \(C_w\).

Lifting to all intermediate orbit coordinates gives

\[
x_i
=
\frac{3^i}{2^{A_i}}x_0
+
\frac{C_i}{2^{A_i}},
\qquad 0\leq i\leq j.
\]

The lifted direction

\[
\left(
1,\frac3{2^{A_1}},\ldots,
\frac{3^j}{2^{A_j}}
\right)
\]

does retain the partial sums \(A_i\). But this still does not create the needed Kakeya contradiction:

\begin{itemize}
\item every finite word already has an entire line of positive realizations;
\item a Kakeya lower bound on the union of many such lines is compatible with that fact;
\item a hypothetical counterexample selects one nested code, not a positive proportion of directions;
\item the desired conclusion concerns whether a particular line contains a bounded positive integer satisfying a signed inequality.
\end{itemize}

The problem remains pointwise and arithmetic.


\bigskip



\subsubsection*{8. The \(81/64\) subsystem exposes the missing variable}


The recurrent subsystem already found in the compression framework contains the macrosteps

\[
C_0=(1,2,1,2),
\qquad
C_1=(1,1,3,1).
\]

Both have length \(4\) and total valuation \(6\), but

\[
C(C_0)=85,\qquad C(C_1)=89.
\]

Their affine maps are

\[
F_\varepsilon(x)
=
\frac{81x+85+4\varepsilon}{64},
\qquad
\varepsilon\in\{0,1\}. \tag{19}
\]

After a linear conjugacy this becomes

\[
m\longmapsto\frac{81m+\varepsilon}{64}. \tag{20}
\]

Thus the two branches have exactly the same direction and multiplicative slope. They differ only in an intercept bit and its propagated carry. [\hypertarget{cit:R19:8}{}\hyperlink{ref:R19:8}{8}]

Any theorem that sees only widths, directions, dimensions, or transversality identifies these two branches. Yet distinguishing their accumulated intercepts is the entire unresolved trace/carry problem.

This is a concrete internal demonstration that direction-only Kakeya information discards the decisive Collatz data.


\bigskip



\subsubsection*{9. The exact signed incidence theorem that would prove Collatz}


The original geometric proposal can be repaired by replacing ``tube intersection'' with a signed intercept condition.

Let an actual word \(w\) of length \(j\) and total valuation \(A\) carry \(n\) to \(X\). Let a competing word \(b\) of length \(k\leq j\) and total valuation \(B\) carry \(y\) to the same \(X\). Put

\[
d=j-k,\qquad e=A-B,
\]

and suppose the \(3\)-adic compatibility condition

\[
3^k\mid C_w-2^eC_b \tag{21}
\]

holds. Define

\[
q=\frac{C_w-2^eC_b}{3^k},
\qquad
\Delta=2^e-3^d. \tag{22}
\]

The two affine equations give

\[
2^e y=3^d n+q \tag{23}
\]

and therefore

\[
2^e(y-n)=q-\Delta n. \tag{24}
\]

These exact comparison identities, including the resulting fixed-input trichotomy, are established in the existing framework. [\hypertarget{cit:R19:9}{}\hyperlink{ref:R19:9}{9}]

The roles of the three places are now transparent:

\[
\begin{array}{c|c}
\text{place} & \text{condition}\\ \hline
2\text{-adic} & n\text{ realizes }w\\
3\text{-adic} & 3^k\mid C_w-2^eC_b\\
\text{real} & -3^dn<q<\Delta n.
\end{array}
\]

Indeed, when \(\Delta>0\),

\[
y>0
\quad\Longleftrightarrow\quad
q>-3^dn,
\]

and

\[
y<n
\quad\Longleftrightarrow\quad
q<\Delta n.
\]

Thus a positive descending competitor exists precisely when

\[
-3^dn
<
\frac{C_w-2^eC_b}{3^k}
<
(2^e-3^d)n. \tag{25}
\]

This is the actual multi-place incidence interval.


\paragraph*{Required theorem: signed adelic incidence}


Every nontrivial infinite retained valuation ray whose nested \(2\)-adic realization representatives stabilize at one positive integer \(n\) admits, at some finite depth, a compatible competing word \(b\) satisfying

\[
2^e>3^d
\]

and the signed interval (25).

Equivalently, no positive stationary code can remain indefinitely geodesic for the inverse-path Bellman problem.

This is the same global input-synchronization theorem isolated in the previous exact framework. The framework already records that terminal \(3\)-adic compatibility is local, while the sign deciding descent depends on the full archimedean prefix. [\hypertarget{cit:R19:10}{}\hyperlink{ref:R19:10}{10}] It also records that carry synchronization inside the \(81/64\) branch would not by itself settle the full problem; a separate global input-synchronization theorem is still required. [\hypertarget{cit:R19:11}{}\hyperlink{ref:R19:11}{11}]


\bigskip



\subsubsection*{10. Conditional proof of Collatz from signed incidence}


Assume the signed adelic incidence theorem.

Suppose, for contradiction, that Collatz is false. Let \(m>1\) be the least positive odd integer whose accelerated orbit does not reach \(1\). Write

\[
x_0=m,\qquad x_{j+1}=U(x_j).
\]

If a positive odd integer \(y\) reaches any later \(x_j\), then

\[
y\geq m. \tag{26}
\]

Otherwise \(y<m\), so by minimality \(y\) reaches \(1\); after merging with the nonterminating orbit of \(m\), that is impossible. This is the coalescence-minimality principle. [\hypertarget{cit:R19:12}{}\hyperlink{ref:R19:12}{12}]

Let \(w_j\) be the actual valuation prefix from \(m\) to \(x_j\). Its seed cylinder contains \(m\). Once

\[
2^{A_j+1}>m,
\]

the least positive representative of that cylinder is exactly \(m\). Hence the nested realization representatives eventually stabilize at the positive integer \(m\).

The actual ray must also survive every valid descending-path elimination: any positive compatible ancestor smaller than \(m\) would already contradict (26). It is therefore a positive stationary retained ray.

Apply signed adelic incidence. At some finite depth, there is a competing word \(b\) satisfying (21) and (25). Equation (23) then defines a positive odd integer \(y\), ending at the same orbit point \(x_j\), and (24) gives

\[
y<m.
\]

This contradicts (26).

Therefore no least positive counterexample exists, so every positive odd integer reaches \(1\). The ordinary Collatz conjecture follows. \(\square\)

The conditional argument is complete. The missing statement is the signed adelic incidence theorem itself.


\bigskip



\subsubsection*{11. Why the proposed machinery does not prove that theorem}


The theorem requires forcing the carry-dependent integer

\[
q=\frac{C_w-2^eC_b}{3^k}
\]

into a short, asymmetric interval whose endpoints depend on the one fixed seed \(n\).

Kakeya and decoupling estimates generally provide:

\begin{itemize}
\item lower bounds for unions of tubes;
\item \(L^p\) estimates averaged over many directions;
\item restrictions on how strongly large tube families can overlap;
\item dimension or entropy conclusions.
\end{itemize}

None of these implies that, for one exceptional word \(w\), some compatible \(q\) has the correct sign.

The distinction is already visible in the retained rational ray

\[
(1,1,2)^\infty.
\]

It is realized by the negative rational fixed cycle based at

\[
-\frac{19}{11},
\]

and every prefix is retained. Its block map is

\[
x\longmapsto\frac{27x+19}{16}. \tag{27}
\]

Thus even exact infinite Bellman retention, recurrent affine geometry, and slope greater than one do not force an ordinary positive integer. [\hypertarget{cit:R19:13}{}\hyperlink{ref:R19:13}{13}]

The essential discriminator is the diagonal/height condition

\[
\xi_2=\xi_3=n\in\mathbb Z_{>0},
\]

together with the signed inequality (25). That information is invisible to ordinary directional incidence theory.


\bigskip



\subsection*{Final assessment}


The mixed \(2\)-adic/\(3\)-adic viewpoint is mathematically natural, but the Kakeya interpretation breaks at the decisive point:

\[
\boxed{
\text{Collatz words do not form a transverse Kakeya family.}
}
\]

At fixed scale and drift they form exponentially many nearly coincident \textbf{parallel} affine lines. Their \(3\)-adic endpoint cylinders are moving, not nested. After correcting the construction to a two-sided adelic state, coherent critical configurations exist for purely symbolic reasons.

The approach therefore does not currently yield a proof of Collatz. Its strongest rigorous output is the reduction

\[
\boxed{
\begin{aligned}
&\text{Collatz}\\
&\Longleftarrow
\text{uniform signed adelic incidence}\\
&=
\text{\(3\)-adic carry compatibility}
+
\text{archimedean intercept escape}\\
&\hspace{2.6em}
\text{for every positive stationary \(2\)-adic ray}.
\end{aligned}
}
\]

That replacement theorem is substantially sharper than a generic Kakeya or decoupling statement: it must eliminate a single exceptional positive diagonal point, not merely control the dimension or measure of a tube union. At present, it is exactly the unproved global synchronization step rather than a consequence of known incidence machinery.
\subsection*{References}
\begin{enumerate}[leftmargin=2.2em]
\item \hypertarget{ref:R19:1}{}\hyperlink{cit:R19:1}{\texttt{turn562990search2}}
\item \hypertarget{ref:R19:2}{}\hyperlink{cit:R19:2}{\texttt{turn2file0 L81-L163}}
\item \hypertarget{ref:R19:3}{}\hyperlink{cit:R19:3}{\texttt{turn619080view0}}
\item \hypertarget{ref:R19:4}{}\hyperlink{cit:R19:4}{\texttt{turn619080view2}}
\item \hypertarget{ref:R19:5}{}\hyperlink{cit:R19:5}{\texttt{turn2file1 L337-L390}}
\item \hypertarget{ref:R19:6}{}\hyperlink{cit:R19:6}{\texttt{turn619080view4}}
\item \hypertarget{ref:R19:7}{}\hyperlink{cit:R19:7}{\texttt{turn619080view1}}
\item \hypertarget{ref:R19:8}{}\hyperlink{cit:R19:8}{\texttt{turn3file0 L134-L161}}
\item \hypertarget{ref:R19:9}{}\hyperlink{cit:R19:9}{\texttt{turn2file1 L284-L335}}
\item \hypertarget{ref:R19:10}{}\hyperlink{cit:R19:10}{\texttt{turn1file3 L248-L262}}
\item \hypertarget{ref:R19:11}{}\hyperlink{cit:R19:11}{\texttt{turn3file1 L183-L198}}
\item \hypertarget{ref:R19:12}{}\hyperlink{cit:R19:12}{\texttt{turn2file0 L165-L205}}
\item \hypertarget{ref:R19:13}{}\hyperlink{cit:R19:13}{\texttt{turn3file2 L271-L347}}
\end{enumerate}

\clearpage
\section{Round 20: approach 15 --- Lee-Yang zeros of a Collatz partition function}
\subsection*{Prompt}
\begin{quote}
Complete a full proof of the Collatz conjecture using the following approach:

\begin{enumerate}
\item Lee-Yang zeros of a Collatz partition function
\end{enumerate}

\textbf{Speculative \textperiodcentered{} B}

Define a finite-depth partition function

Z\_j(beta, gamma) = sum over w of exp(-beta \emph{ drift(w)) } exp(-gamma * defect(w)),

where the defect measures failure of exact integer realizability.

A uniform zero-free region could give analyticity of the limiting pressure and exponential extinction of realizable critical words. Conversely, zeros pinching the physical axis could indicate a genuine phase transition associated with a nontrivial infinite survivor set.

Take as long as you need.
\end{quote}
\subsection*{Response}

\subsection*{Outcome}


After making the proposed statistical-mechanical model exact, \textbf{it does not produce a proof of Collatz}. The obstruction is concrete, rather than merely that Collatz is open:

\begin{enumerate}
\item Every finite accelerated valuation word is already exactly realizable by infinitely many positive integers. Therefore a finite-depth ``integer realizability defect'' is identically zero.
\item Even replacing that defect by ``failure to remain critical'' cannot yield extinction of finite critical words: arbitrarily long exactly realizable non-descending words exist.
\item Zero-freeness gives analyticity or correlation decay, not the nonexistence of an isolated arithmetic trajectory. A unique survivor can coexist with a completely zero-free partition function.
\end{enumerate}

There is, however, a corrected \textbf{seed-pinned Lee\textendash{}Yang formulation} in which a counterexample necessarily creates zeros pinching the physical point \(z=1\). A uniform zero-free region for that corrected partition function would prove Collatz. I derive the exact criterion below. The missing zero-free estimate turns out to be equivalent to a quantitative first-descent theorem stronger than Collatz itself.

Recent 2026 work continues to describe the conjecture as open, with Tao's almost-all theorem remaining the strongest general analytic result. [\hypertarget{cit:R20:1}{}\hyperlink{ref:R20:1}{1}, \hypertarget{cit:R20:2}{}\hyperlink{ref:R20:2}{2}]


\bigskip



\subsubsection*{1. Accelerated words and their natural drift}


Work on the odd positive integers with the accelerated Syracuse map

\[
U(n)=\frac{3n+1}{2^{v_2(3n+1)}}.
\]

For an odd starting value \(n\), put

\[
a_k(n)=v_2\!\left(3U^k(n)+1\right)\ge 1.
\]

A finite valuation word is

\[
w=(a_0,\ldots ,a_{j-1})\in \mathbb N_{\ge 1}^{j}.
\]

Write

\[
A_k=\sum_{i=0}^{k-1}a_i,\qquad A_0=0,
\]

so that

\[
2^{A_j}U^j(n)=3^j n+C_w
\]

for an explicitly determined positive integer \(C_w\). The natural signed archimedean drift is

\[
D(w)=\log\frac{2^{A_j}}{3^j}
    =A_j\log 2-j\log 3.
\]

Negative \(D(w)\) means that the multiplicative part \(3^j/2^{A_j}\) grows; positive \(D(w)\) means contraction.

The proposed partition function is therefore naturally of the form

\[
Z_j(\beta,\gamma)
 =\sum_{w}
   e^{-\beta D(w)}
   e^{-\gamma\,\operatorname{def}(w)}.
\]

The first question is whether \(\operatorname{def}(w)\) can measure failure of a finite word to be realized by an actual positive integer.

It cannot.


\bigskip



\subsubsection*{2. Every finite valuation word is exactly realizable}



\paragraph*{Theorem 1 \textemdash{} finite realizability is automatic}


For every word

\[
w=(a_0,\ldots ,a_{j-1})\in\mathbb N_{\ge1}^{j},
\]

there are infinitely many positive odd integers \(n\) satisfying

\[
v_2\!\left(3U^k(n)+1\right)=a_k,
\qquad 0\le k<j.
\]

More precisely, these integers form one residue class modulo

\[
2^{A_j+1}.
\]


\paragraph*{Proof}


Introduce the shortcut Collatz map on all integers,

\[
T(m)=
\begin{cases}
m/2,&m\equiv0\pmod2,\\[2mm]
(3m+1)/2,&m\equiv1\pmod2.
\end{cases}
\]

An accelerated step of valuation \(a_i\) corresponds to the parity block

\[
1\,0^{a_i-1}.
\]

To ensure that the result after the final block is again odd, append one final \(1\). Thus \(w\) determines the binary parity string

\[
b(w)=
1\,0^{a_0-1}\,
1\,0^{a_1-1}\cdots
1\,0^{a_{j-1}-1}\,1,
\]

whose total length is

\[
A_j+1.
\]

The finite parity-vector bijection says that every binary string of length \(L\) is realized on exactly one residue class modulo \(2^L\). In particular, that residue class contains infinitely many positive integers. Applying this with \(L=A_j+1\) gives the result. The parity-vector bijection and its extension to a bijection between infinite parity sequences and \(2\)-adic integers are classical; a modern exposition gives precisely these statements. [\hypertarget{cit:R20:3}{}\hyperlink{ref:R20:3}{3}]

\(\square\)


\paragraph*{Immediate consequence}


Suppose the defect has the intended meaning

\[
\operatorname{def}(w)=0
\quad\Longleftrightarrow\quad
w\text{ is realized by some positive integer}.
\]

Then

\[
\boxed{\operatorname{def}(w)=0\text{ for every finite }w.}
\]

The arithmetic part of the proposed partition function disappears.

For the natural drift,

\[
\begin{aligned}
Z_j(\beta,\gamma)
&=\sum_{a_0,\ldots ,a_{j-1}\ge1}
  \exp\!\left(
   -\beta\left[
      \left(\sum_i a_i\right)\log2-j\log3
   \right]\right)\\
&=\left(
   3^\beta\sum_{a\ge1}2^{-\beta a}
  \right)^j\\
&=\boxed{
 \left(\frac{3^\beta}{2^\beta-1}\right)^j
 },
\qquad \Re\beta>0.
\end{aligned}
\]

It is independent of \(\gamma\). It is meromorphic and has no nontrivial zeros. Its zero-freeness is simply the zero-freeness of a geometric series, not a theorem about Collatz.

If one truncates the alphabet to \(1\le a\le M\), one obtains

\[
Z_{j,M}(\beta)
 =
 \left[
  3^\beta
  \frac{2^{-\beta}(1-2^{-M\beta})}
       {1-2^{-\beta}}
 \right]^j.
\]

The new zeros come from

\[
2^{-M\beta}=1.
\]

They are artifacts of the artificial valuation cutoff \(M\), not arithmetic obstructions.

This is the first decisive failure of the proposed model.


\bigskip



\subsubsection*{3. Critical finite words never become extinct}


Perhaps ``defect'' should instead measure failure of a word to remain above its starting seed. That also cannot produce finite-depth extinction.


\paragraph*{Theorem 2 \textemdash{} arbitrarily long exact non-descending words}


For every \(j\ge1\), let

\[
n_j=2^{j+1}-1.
\]

Then the first \(j\) accelerated valuations of \(n_j\) are all equal to \(1\), and

\[
U^k(n_j)=3^k2^{j+1-k}-1,
\qquad 0\le k\le j.
\]

In particular,

\[
U^k(n_j)\ge n_j
\qquad 0\le k\le j.
\]


\paragraph*{Proof}


Set

\[
x_k=3^k2^{j+1-k}-1.
\]

For \(0\le k<j\),

\[
\begin{aligned}
3x_k+1
&=3^{k+1}2^{j+1-k}-2\\
&=2\left(3^{k+1}2^{j-k}-1\right).
\end{aligned}
\]

Since \(j-k\ge1\), the quantity in parentheses is odd. Hence

\[
v_2(3x_k+1)=1
\]

and

\[
U(x_k)=3^{k+1}2^{j-k}-1=x_{k+1}.
\]

Also,

\[
x_k+1
 =3^k2^{j+1-k}
 =\left(\frac32\right)^k(n_j+1),
\]

so \(x_k\ge n_j\).

\(\square\)

Thus the all-ones word

\[
1^j=(1,\ldots ,1)
\]

is an exactly realizable critical word at every depth.

Moreover,

\[
n_j=2^{j+1}-1\longrightarrow -1
\qquad\text{in }\mathbb Z_2.
\]

The corresponding infinite valuation sequence \(1^\infty\) is realized by the \(2\)-adic fixed point

\[
U(-1)=-1.
\]

The all-one parity sequence corresponding to the \(2\)-adic integer \(-1\) is also explicit in the parity-sequence literature. [\hyperlink{ref:R20:3}{3}]

This reveals the fundamental compactness problem:

\begin{itemize}
\item archimedeanly, \(n_j\to+\infty\);
\item \(2\)-adically, \(n_j\to-1\);
\item every finite prefix looks perfectly realizable;
\item the infinite limit is not a positive integer counterexample.
\end{itemize}

Consequently,

\[
\boxed{
\text{Finite critical words cannot become extinct, even if Collatz is true.}
}
\]

What must become extinct is not finite critical words in aggregate. It is the possibility that \textbf{one fixed positive seed} remains critical at every depth.


\bigskip



\subsubsection*{4. The correct pointwise event}


For an odd \(n>1\), define its first-descent time

\[
\tau(n)=
\inf\{k\ge1:U^k(n)<n\},
\]

with \(\tau(n)=\infty\) if no such \(k\) exists.


\paragraph*{Proposition 3 \textemdash{} first descent is equivalent to Collatz}


The Collatz conjecture is equivalent to

\[
\tau(n)<\infty
\qquad
\text{for every odd }n>1.
\]


\paragraph*{Proof}


If Collatz holds, then every \(n>1\) eventually reaches \(1<n\), so \(\tau(n)<\infty\).

Conversely, suppose every odd \(n>1\) eventually reaches a smaller odd positive integer. Use strong induction on \(n\). Given \(n>1\), choose \(k\) such that

\[
m=U^k(n)<n.
\]

By induction, \(m\) reaches \(1\), and hence so does \(n\).

\(\square\)

This formulation simultaneously excludes divergent trajectories and nontrivial cycles. If any counterexample exists, its least odd element \(n_*\) satisfies

\[
U^k(n_*)\ge n_*
\qquad
\text{for all }k,
\]

so \(\tau(n_*)=\infty\).

Define the depth-\(j\) survivor set

\[
\mathcal S_j
 =
 \{n\ge3:n\text{ odd and }\tau(n)>j\},
\]

and its smallest element

\[
m_j=\min\mathcal S_j.
\]

Theorem 2 guarantees that \(\mathcal S_j\neq\varnothing\), since

\[
2^{j+1}-1\in\mathcal S_j.
\]

The correct compactness criterion is

\[
\boxed{
\text{Collatz holds}
\quad\Longleftrightarrow\quad
m_j\longrightarrow\infty.
}
\]

Indeed, a counterexample \(n_*\) gives \(m_j\le n_*\) for every \(j\). Conversely, if every fixed integer has finite first-descent time, then for each finite \(M\), all \(n\le M\) are absent from \(\mathcal S_j\) once \(j\) is sufficiently large.


\bigskip



\subsubsection*{5. A seed-pinned Lee\textendash{}Yang partition function}


The word-only partition function loses the distinction between a fixed seed and seeds escaping to infinity. The necessary correction is an archimedean chemical potential.

Fix \(\eta>0\). Define the grand partition function

\[
\boxed{
\Xi_j(z;\eta)
 =
 \prod_{n\in\mathcal S_j}
 \left(1+e^{-\eta n}z^j\right).
}
\]

This can be interpreted as a gas of independently occupiable depth-\(j\) survivor seeds. Expanding the product gives

\[
\Xi_j(z;\eta)
 =
 \sum_{\substack{F\subset\mathcal S_j\\F\text{ finite}}}
 e^{-\eta\sum_{n\in F}n}
 z^{j|F|}.
\]

For each fixed \(j\), the product converges uniformly on compact subsets of \(\mathbb C\), because

\[
\sum_{\substack{n\ge3\\n\text{ odd}}}e^{-\eta n}<\infty.
\]

It is therefore an entire function.


\paragraph*{Theorem 4 \textemdash{} exact location of the survivor zeros}


The zeros of \(\Xi_j\) are exactly

\[
z_{j,n,\ell}
 =
 \exp\!\left(\frac{\eta n}{j}\right)
 \exp\!\left(\frac{(2\ell+1)\pi i}{j}\right),
\]

where \(n\in\mathcal S_j\) and \(\ell\in\mathbb Z\).

In particular, the zero generated by the smallest survivor \(m_j\) and closest to the positive real axis is

\[
\boxed{
z_j^*
 =
 \exp\!\left(\frac{\eta m_j}{j}\right)
 \exp\!\left(\frac{\pi i}{j}\right).
}
\]


\paragraph*{Proof}


Each factor vanishes precisely when

\[
1+e^{-\eta n}z^j=0,
\]

or equivalently

\[
z^j=-e^{\eta n}.
\]

Taking the \(j\)-th roots gives the displayed formula. The convergent product has no other zeros.

\(\square\)


\bigskip



\subsubsection*{6. Counterexamples force Lee\textendash{}Yang zeros to pinch \(z=1\)}


Suppose \(n_*\) is a Collatz counterexample. Then

\[
n_*\in\mathcal S_j
\qquad\text{for every }j.
\]

Consequently, \(\Xi_j\) has the zero

\[
z_{j,n_*,0}
 =
 \exp\!\left(\frac{\eta n_*}{j}\right)
 \exp\!\left(\frac{\pi i}{j}\right).
\]

As \(j\to\infty\),

\[
z_{j,n_*,0}\longrightarrow1.
\]

Thus

\[
\boxed{
\text{A positive-integer counterexample forces zeros to pinch the physical point }z=1.
}
\]

This direction is fully rigorous.

The converse is not valid. Since

\[
z_j^*\to1
\quad\Longleftrightarrow\quad
\frac{m_j}{j}\to0
\]

along the relevant subsequence, zeros can pinch \(1\) even when

\[
m_j\to\infty
\]

and Collatz is true. Pinching may therefore indicate seeds with longer and longer finite excursions whose size grows sublinearly in their first-descent time. It does not, by itself, certify a fixed infinite survivor.

That corrects the proposed physical interpretation:

\[
\text{counterexample}\Longrightarrow\text{pinching},
\]

but not necessarily

\[
\text{pinching}\Longrightarrow\text{counterexample}.
\]

Lee\textendash{}Yang theory associates accumulating zeros with thermodynamic singularities or competing phases, but an isolated arithmetic trajectory is not automatically a macroscopic phase. [\hypertarget{cit:R20:4}{}\hyperlink{ref:R20:4}{4}]


\bigskip



\subsubsection*{7. The zero-free theorem that would prove Collatz}


For \(R>1\) and \(\theta>0\), define the sector around the physical axis

\[
\Omega(R,\theta)
 =
 \left\{
 z\in\mathbb C:
 1\le |z|\le R,\quad |\arg z|\le\theta
 \right\}.
\]


\paragraph*{Theorem 5 \textemdash{} uniform zero-free sector implies Collatz}


Suppose there exist constants

\[
\eta>0,\qquad R>1,\qquad \theta>0,\qquad J\ge1
\]

such that

\[
\Xi_j(z;\eta)\neq0
\]

for every \(j\ge J\) and every \(z\in\Omega(R,\theta)\).

Then

\[
m_j>\frac{\log R}{\eta}\,j
\]

for all sufficiently large \(j\), and consequently every positive integer reaches \(1\).


\paragraph*{Proof}


Choose \(j\) large enough that

\[
\frac{\pi}{j}<\theta.
\]

Suppose some \(n\in\mathcal S_j\) satisfies

\[
n\le\frac{\log R}{\eta}\,j.
\]

Then

\[
\left|z_{j,n,0}\right|
 =
 e^{\eta n/j}
 \le R
\]

and

\[
\arg z_{j,n,0}=\frac{\pi}{j}<\theta.
\]

Thus \(z_{j,n,0}\in\Omega(R,\theta)\), contradicting zero-freeness. Hence every \(n\in\mathcal S_j\) obeys

\[
n>\frac{\log R}{\eta}\,j.
\]

Now fix any odd \(n>1\). Choose \(j\) with

\[
j>\frac{\eta}{\log R}\,n.
\]

Then \(n<m_j\), so \(n\notin\mathcal S_j\). Therefore \(\tau(n)\le j\). Proposition 3 finishes the proof.

\(\square\)

This is a genuine complete conditional proof:

\[
\boxed{
\text{Uniform Lee–Yang zero-freeness beyond }|z|=1
\Longrightarrow
\text{Collatz}.
}
\]

But the zero-free assertion is equivalent to a strong quantitative form of Collatz.

Indeed, the theorem gives

\[
\tau(n)
 \le
 \frac{\eta}{\log R}\,n+O(1).
\]

So it proves not merely eventual descent, but a linear upper bound for the first-descent time.

Conversely, if one could prove

\[
m_j\ge c j
\]

for some \(c>0\), then every zero of \(\Xi_j\) would have modulus at least

\[
e^{\eta c},
\]

and any disk

\[
|z|<R<e^{\eta c}
\]

would be uniformly zero-free.

Thus, up to the values of the constants,

\[
\boxed{
\text{Uniform zero-free region}
\quad\Longleftrightarrow\quad
m_j\gg j.
}
\]

The zero-free theorem has not reduced the arithmetic problem; it has encoded a quantitative version of it.


\bigskip



\subsubsection*{8. Equivalent exponential-extinction formulation}


Define the weighted survivor mass

\[
S_j(\eta)
 =
 \sum_{n\in\mathcal S_j}e^{-\eta n}.
\]

Since \(m_j\) is the smallest odd survivor,

\[
e^{-\eta m_j}
 \le S_j(\eta)
 \le
 \sum_{\ell\ge0}e^{-\eta(m_j+2\ell)}
 =
 \frac{e^{-\eta m_j}}{1-e^{-2\eta}}.
\]

Therefore

\[
m_j\gg j
\]

is equivalent to

\[
\boxed{
S_j(\eta)\le C e^{-c j}
}
\]

for some \(C,c>0\).

This weighted estimate by itself proves Collatz. If \(n_*\) were a counterexample, then

\[
e^{-\eta n_*}
 \le S_j(\eta)
 \le Ce^{-cj}
\]

for every \(j\), which is impossible as \(j\to\infty\).

Hence the exact decisive theorem for this program is:


\begin{quote}
\textbf{Pinned exponential-extinction theorem.}   There exist \(\eta,c,C>0\) such that \[
> \sum_{\substack{n\ge3,\ n\text{ odd}\\
> U^k(n)\ge n\ \forall\,1\le k\le j}}
> e^{-\eta n}
> \le Ce^{-cj}
> \] for every \(j\ge1\).

\end{quote}


Proving this theorem would finish Collatz immediately.

It remains unproved.


\bigskip



\subsubsection*{9. Why ordinary Lee\textendash{}Yang analyticity does not supply the theorem}


There is a logical gap between zero-freeness and extinction.

Suppose a finite system has a unique zero-defect configuration of weight \(1\). Its partition function can simply be

\[
Z_j=1.
\]

This function is entire and zero-free everywhere, but the survivor exists.

More generally, a uniform zero-free region implies that

\[
\log Z_j
\]

can be chosen analytically. It does not imply that a particular coefficient vanishes, that the zero-defect sector is empty, or that the pressure has the sign needed for extinction.

If the zero-defect partition function is an integer count \(N_j\), then

\[
N_j>0\Longrightarrow \frac1j\log N_j\ge0.
\]

Proving strictly negative pressure would already prove \(N_j=0\) for all sufficiently large \(j\). In that normalization, the pressure estimate is simply extinction restated.

If configurations have weights smaller than \(1\), then negative pressure no longer implies emptiness: one surviving configuration may have exponentially small weight. One then needs a seed-dependent lower bound, which is exactly why the chemical potential \(e^{-\eta n}\) and the individual-factor product were introduced.

So the inference

\[
\text{zero-free region}
\Longrightarrow
\text{analytic pressure}
\Longrightarrow
\text{no survivor}
\]

is invalid without a separate quantitative coefficient or spectral estimate.


\bigskip



\subsubsection*{10. Why an unpinned model would recover only ``almost all''}


Under Haar measure on the odd \(2\)-adic integers, an exact valuation cylinder \(w\) has conditional measure

\[
2^{-A_j}
 =
 \prod_{i=0}^{j-1}2^{-a_i}.
\]

Thus the accelerated valuations are exactly independent geometric variables with

\[
\mathbb P(a_i=a)=2^{-a}.
\]

This is the natural random Syracuse model. It has negative mean logarithmic drift because

\[
\mathbb E a_i=2,
\qquad
\mathbb E\log\frac{3}{2^{a_i}}
 =\log\frac34<0.
\]

A partition function averaged over these cylinders is therefore extremely well suited to density statements. It is not automatically suited to excluding one exceptional cylinder or one fixed positive seed.

Tao's result obtains approximate transport for Syracuse random variables and proves that almost all Collatz orbits attain almost bounded values in logarithmic density. Its conclusion remains averaged rather than pointwise. [\hyperlink{ref:R20:2}{2}]

There is a parallel issue in rigorous Lee\textendash{}Yang theory. In representative zero-free-to-mixing results, the zero-free hypothesis must hold uniformly for systems with boundary conditions before one obtains strong spatial mixing. Even then, the conclusion is decay of boundary influence in a probability distribution, not the emptiness of the configuration space. [\hypertarget{cit:R20:5}{}\hyperlink{ref:R20:5}{5}, \hypertarget{cit:R20:6}{}\hyperlink{ref:R20:6}{6}]

For Collatz, the necessary boundary conditions include at least:

\[
\text{starting }2\text{-adic cylinder},
\qquad
\text{archimedean seed height},
\qquad
\text{drift or height gap},
\qquad
\text{terminal }3\text{-adic carry state}.
\]

Zero-freeness only after averaging over starting cylinders would leave precisely the same pointwise gap as an almost-all theorem.


\bigskip



\subsubsection*{11. How the suffix-locality result could enter}


Your terminal-suffix locality theorem can make part of the \(3\)-adic endpoint constraint into a finite-range interaction: modulo \(3^r\), the relevant terminal compatibility is determined by the final \(r\) valuations.

At fixed precision \(r\), one could therefore construct a transfer or polymer model whose state records

\[
\bigl(
\text{seed residue},
\text{height/drift state},
\text{last }r\text{ valuations},
\text{endpoint carry mod }3^r
\bigr).
\]

A maximal interval during which negative multiplicative drift is rescued by the affine carry could be treated as a polymer. The desired cluster-expansion estimate would resemble a Koteck\'y\textendash{}Preiss condition uniform in:

\[
r,\qquad
\text{the initial seed cylinder},\qquad
\text{the endpoint boundary state}.
\]

Such uniformity could produce boundary-conditioned zero-freeness and a weighted contraction estimate.

But the crucial activity bound would have to show that a fixed small seed cannot support rescue polymers of arbitrary length. A hypothetical counterexample is exactly a fixed seed supporting such polymers forever. Consequently, the missing polymer estimate is not presently derivable from finite-word realizability or suffix locality alone.

The suffix theorem localizes the \(3\)-adic side. It does not localize the condition that the nested \(2\)-adic seed representatives eventually stabilize to one ordinary positive integer. That is an archimedean, non-shift-invariant boundary condition.


\bigskip



\subsection*{Final assessment}


The proposed partition function

\[
Z_j(\beta,\gamma)
 =
 \sum_w
 e^{-\beta\,\mathrm{drift}(w)}
 e^{-\gamma\,\mathrm{defect}(w)}
\]

cannot prove Collatz when the defect measures finite exact integer realizability:

\[
\operatorname{defect}(w)\equiv0.
\]

Its partition function factorizes explicitly and its zeros contain no Collatz information.

Replacing realizability by finite criticality also does not work, because exact critical words exist at every depth:

\[
w=1^j,\qquad n=2^{j+1}-1.
\]

A corrected seed-pinned grand partition function is

\[
\Xi_j(z;\eta)
 =
 \prod_{\tau(n)>j}
 \left(1+e^{-\eta n}z^j\right).
\]

For this model:

\[
\text{counterexample}
\Longrightarrow
\text{zeros pinch }z=1,
\]

and

\[
\text{uniform zero-free sector beyond }|z|=1
\Longrightarrow
m_j\gg j
\Longrightarrow
\text{Collatz}.
\]

The missing statement is equivalently

\[
\sum_{\tau(n)>j}e^{-\eta n}\le Ce^{-cj}.
\]

That is a precise and useful quantitative target, but it has not been proved here. Calling it a Lee\textendash{}Yang zero-free theorem does not yet make it easier: for the corrected partition function, the location of the zeros shows that the desired zero-free region is essentially equivalent to that exponential first-descent estimate.

Thus approach 15 presently yields a rigorous \textbf{no-go theorem for the naive construction} and an exact \textbf{Lee\textendash{}Yang reformulation of the decisive pointwise estimate}, rather than a completed proof of Collatz.
\subsection*{References}
\begin{enumerate}[leftmargin=2.2em]
\item \hypertarget{ref:R20:1}{}\hyperlink{cit:R20:1}{\texttt{turn333960search5}}
\item \hypertarget{ref:R20:2}{}\hyperlink{cit:R20:2}{\texttt{turn309707search0}}
\item \hypertarget{ref:R20:3}{}\hyperlink{cit:R20:3}{\texttt{turn200196view0}}
\item \hypertarget{ref:R20:4}{}\hyperlink{cit:R20:4}{\texttt{turn991325search3}}
\item \hypertarget{ref:R20:5}{}\hyperlink{cit:R20:5}{\texttt{turn421558view3}}
\item \hypertarget{ref:R20:6}{}\hyperlink{cit:R20:6}{\texttt{turn607614view2}}
\end{enumerate}

\clearpage
\section{Round 21: approach 16 --- A quantum-graph trace formula for cycles and escape}
\subsection*{Prompt}
\begin{quote}
Complete a full proof of the Collatz conjecture using the following approach:

\begin{enumerate}
\item A quantum-graph trace formula for cycles and escape
\end{enumerate}

\textbf{Speculative \textperiodcentered{} B}

Turn each inverse branch into a scattering edge carrying its affine phase. Primitive closed paths are then Collatz cycle words, while resonances near the unit circle encode long-lived or escaping trajectories.

An exact trace formula might reduce cycle exclusion to positivity or cancellation in a spectral density. The main danger is that the spectral object merely stores the original cycles without introducing any new rigidity.

Take as long as you need.
\end{quote}
\subsection*{Response}

\subsection*{Verdict}


After carrying the construction through, it does \textbf{not} produce a valid full proof of the Collatz conjecture. It does produce:

\begin{enumerate}
\item an exact affine trace formula for positive Collatz cycles;
\item a sharp spectral annulus in which every hypothetical nontrivial cycle must appear;
\item a rigorous obstruction showing that ordinary quantum-graph resonances cannot distinguish a divergent trajectory from harmless convergent branches already present in the Collatz graph.
\end{enumerate}

Thus the central danger in the proposal is real: once the spectral object is made arithmetically exact, it stores the original cycle and escape questions rather than forcing their resolution. Collatz remains open as of August 13, 2026; recent work still describes a pointwise orbit statement as the unresolved step. [\hypertarget{cit:R21:1}{}\hyperlink{ref:R21:1}{1}, \hypertarget{cit:R21:2}{}\hyperlink{ref:R21:2}{2}]

What follows is the complete construction and the precise point where the hoped-for proof fails.


\bigskip



\subsubsection*{1. Affine holonomy of a valuation word}


Use the accelerated odd Collatz map

\[
U(n)=\frac{3n+1}{2^{a(n)}},
\qquad
a(n)=v_2(3n+1),
\]

on positive odd integers.

For a formal branch exponent \(a\geq 1\), put

\[
U_a(x)=\frac{3x+1}{2^a}.
\]

Represent this affine map projectively by

\[
M_a=
\begin{pmatrix}
3&1\\
0&2^a
\end{pmatrix}.
\]

For a word

\[
w=(a_0,\ldots ,a_{j-1}),
\]

define

\[
A_k=\sum_{i=0}^{k-1}a_i,\qquad A=A_j,
\]

and

\[
C_w=\sum_{k=0}^{j-1}3^{\,j-1-k}2^{A_k}.
\]

Then direct matrix multiplication gives

\[
M_{a_{j-1}}\cdots M_{a_0}
=
\begin{pmatrix}
3^j&C_w\\
0&2^A
\end{pmatrix},
\]

and hence

\[
U_{a_{j-1}}\circ\cdots\circ U_{a_0}(x)
=
\frac{3^j x+C_w}{2^A}.
\]

The affine holonomy therefore remembers exactly the three quantities

\[
j,\qquad A,\qquad C_w.
\]

Its unique fixed point, provided \(2^A\ne 3^j\), is

\[
x_w=\frac{C_w}{2^A-3^j}.
\]

This is the familiar cycle equation, but here it arises as the fixed-point equation for the affine scattering holonomy.


\bigskip



\subsubsection*{2. Exact cycle-realizability theorem}



\paragraph*{Theorem 1}


Let \(w=(a_0,\ldots ,a_{j-1})\) with every \(a_i\geq1\), and put

\[
d_w=2^A-3^j.
\]

Then \(w\) is the valuation word of a positive odd Collatz cycle if and only if

\[
d_w>0
\qquad\text{and}\qquad
d_w\mid C_w.
\]

In that case the first element of the cycle is

\[
n=\frac{C_w}{d_w}.
\]


\paragraph*{Proof}


Suppose first that \(n\) belongs to a cycle with valuation word \(w\). Iterating the affine identity gives

\[
2^A n=3^j n+C_w,
\]

so

\[
(2^A-3^j)n=C_w.
\]

Since \(n,C_w>0\), we have \(d_w>0\), and \(d_w\mid C_w\).

Conversely, suppose \(d_w>0\) and \(d_w\mid C_w\), and define

\[
n=\frac{C_w}{d_w}\in\mathbb N.
\]

For \(0\leq r\leq j\), let

\[
C_r=\sum_{k=0}^{r-1}3^{\,r-1-k}2^{A_k}
\]

and formally define

\[
x_r=\frac{3^r n+C_r}{2^{A_r}}.
\]

We first show that every \(x_r\) is an integer. We have

\[
C_w\equiv 3^{j-r}C_r\pmod{2^{A_r}}
\]

because all terms of \(C_w\) beginning at index \(r\) are divisible by \(2^{A_r}\). Also

\[
d_w=2^A-3^j\equiv -3^j\pmod{2^{A_r}}.
\]

Since \(d_wn=C_w\),

\[
-3^j n\equiv 3^{j-r}C_r\pmod{2^{A_r}}.
\]

As \(3\) is invertible modulo \(2^{A_r}\),

\[
3^r n+C_r\equiv0\pmod{2^{A_r}},
\]

so \(x_r\in\mathbb Z\).

Now cyclically rotate the word by \(r\) positions. Since \(x_r\) is fixed by the corresponding rotated affine composition, uniqueness of the fixed point gives

\[
x_r=\frac{C_{\sigma^r w}}{d_w}.
\]

Both \(C_{\sigma^r w}\) and \(d_w\) are odd, so \(x_r\) is odd.

Finally,

\[
3x_r+1=2^{a_r}x_{r+1}.
\]

Because \(x_{r+1}\) is odd,

\[
v_2(3x_r+1)=a_r.
\]

Thus the formal branches are the exact accelerated Collatz branches, all intermediate values are positive odd integers, and \(x_j=x_0=n\). \ensuremath{\blacksquare}

This theorem is important because it removes a possible loophole: divisibility of the final affine fixed point automatically enforces all the intermediate exact valuations.


\bigskip



\subsubsection*{3. There are finitely many cycle words of each length}


Suppose \(x_0,\ldots ,x_{j-1}\) form a positive odd cycle. Multiplying

\[
2^{a_i}x_{i+1}=3x_i+1
\]

over all \(i\), and using \(\prod x_{i+1}=\prod x_i\), gives

\[
2^A
=
\prod_{i=0}^{j-1}\left(3+\frac1{x_i}\right)
\leq 4^j.
\]

Consequently,

\[
A\leq2j.
\]

Therefore every length-\(j\) cycle word occurs among the finite set

\[
\mathcal W_j
=
\left\{
(a_0,\ldots ,a_{j-1}):
a_i\geq1,\ 
\sum_i a_i\leq2j
\right\}.
\]

Its cardinality is

\[
|\mathcal W_j|
=
\sum_{A=j}^{2j}\binom{A-1}{j-1}
=
\binom{2j}{j}.
\]

So the affine trace sums below are genuinely finite at every depth.


\bigskip



\subsubsection*{4. The exact affine-phase trace formula}


The divisibility condition can be written as a finite character sum:

\[
\mathbf 1_{d\mid C}
=
\frac1d
\sum_{r=0}^{d-1}
\exp\!\left(\frac{2\pi i rC}{d}\right).
\]

Hence the number of fixed points of \(U^j\) is exactly

\[
N_j
=
\sum_{\substack{w\in\mathcal W_j\\ d_w>0}}
\frac1{d_w}
\sum_{r=0}^{d_w-1}
\exp\!\left(\frac{2\pi i rC_w}{d_w}\right).
\tag{4.1}
\]

A weighted version, using the cycle multiplier

\[
\lambda_w=\frac{3^j}{2^A},
\]

is

\[
F_j
=
\sum_{\substack{w\in\mathcal W_j\\ d_w>0}}
\frac{3^j}{2^A}
\frac1{d_w}
\sum_{r=0}^{d_w-1}
\exp\!\left(\frac{2\pi i rC_w}{d_w}\right).
\tag{4.2}
\]

Equation (4.2) is the exact version of the proposed ``affine phase.'' The holonomy supplies \(C_w\) and \(d_w\); Fourier orthogonality projects onto those holonomies whose fixed point is an ordinary integer.

The resulting dynamical determinant is formally

\[
D(z)
=
\exp\left(
-\sum_{j\geq1}\frac{F_j}{j}z^j
\right).
\tag{4.3}
\]

If \(\mathcal P\) denotes primitive positive odd cycles modulo cyclic rotation, with length \(\ell(p)\), total valuation \(A(p)\), and multiplier

\[
\lambda_p=\frac{3^{\ell(p)}}{2^{A(p)}},
\]

then the standard fixed-point-to-primitive-orbit calculation gives

\[
D(z)
=
\prod_{p\in\mathcal P}
\left(1-\lambda_p z^{\ell(p)}\right).
\tag{4.4}
\]

Indeed, a primitive cycle of length \(\ell\) contributes its \(\ell\) starting points whenever \(\ell\mid j\), and

\[
\begin{aligned}
\sum_{j\geq1}\frac{F_j}{j}z^j
&=
\sum_{p\in\mathcal P}
\sum_{m\geq1}
\frac{\lambda_p^m z^{m\ell(p)}}{m}\\
&=
-\sum_{p\in\mathcal P}
\log\left(1-\lambda_p z^{\ell(p)}\right).
\end{aligned}
\]

This is exactly the same structural phenomenon as the Roth\textendash{}Kottos\textendash{}Smilansky quantum-graph determinant: a secular determinant expands into an Euler product over primitive closed paths, and logarithmic differentiation yields the periodic-orbit trace formula. For finite quantum graphs this identity is exact, not semiclassical. [\hypertarget{cit:R21:3}{}\hyperlink{ref:R21:3}{3}]


\bigskip



\subsubsection*{5. A sharp resonance annulus for every nontrivial cycle}


The weighted determinant has a useful feature not present in the ordinary unweighted adjacency matrix.

For a cycle \(p=(x_0,\ldots ,x_{\ell-1})\),

\[
2^{A(p)}
=
3^\ell
\prod_{i=0}^{\ell-1}
\left(1+\frac1{3x_i}\right).
\]

Therefore

\[
\lambda_p
=
\frac{3^\ell}{2^{A(p)}}
=
\prod_{i=0}^{\ell-1}
\left(1+\frac1{3x_i}\right)^{-1}.
\tag{5.1}
\]

For the trivial fixed point \(x=1\),

\[
\lambda_{\mathrm{triv}}=\frac34,
\]

so its determinant factor is

\[
1-\frac34z,
\]

with zero at

\[
z=\frac43.
\]

For any nontrivial positive odd cycle, every \(x_i\geq3\). Hence

\[
\frac{9}{10}
\leq
\left(1+\frac1{3x_i}\right)^{-1}
<1.
\]

It follows that

\[
\frac{9}{10}
\leq
\lambda_p^{1/\ell(p)}
<1.
\]

The zeros of its factor

\[
1-\lambda_p z^{\ell(p)}
\]

therefore satisfy

\[
1<|z|
=
\lambda_p^{-1/\ell(p)}
\leq\frac{10}{9}.
\tag{5.2}
\]


\paragraph*{Cycle-annulus criterion}


Every nontrivial positive Collatz cycle forces determinant zeros in the fixed annulus

\[
\boxed{1<|z|\leq\frac{10}{9}},
\]

whereas the trivial fixed point produces its zero at \(4/3\), outside that annulus.

This is the cleanest successful output of the proposed approach.

It turns cycle exclusion into the statement that, after dividing off the trivial factor,

\[
\frac{D(z)}{1-\frac34z}
\]

has no zeros in \(1<|z|\leq10/9\).

But this is a \textbf{criterion}, not yet a proof: nothing in general quantum-graph positivity implies that zero-free region.


\bigskip



\subsubsection*{6. A completely finite version}


To avoid any analytic-continuation issue, fix \(R\) and let

\[
V_R=\{n\leq R:n\text{ odd}\}.
\]

Define the weighted truncated Collatz matrix \(B_R\) by

\[
B_Re_n=
\begin{cases}
\dfrac{3}{2^{a(n)}}e_{U(n)},&U(n)\leq R,\\[2mm]
0,&U(n)>R.
\end{cases}
\]

Because this is a finite functional digraph, every component consists of a directed cycle with incoming trees, or a tree ending at the boundary. Tree blocks are nilpotent. Consequently,

\[
\det(I-zB_R)
=
\prod_{\substack{p\text{ primitive cycle}\\p\subseteq V_R}}
\left(1-\lambda_p z^{\ell(p)}\right).
\tag{6.1}
\]

In particular,

\[
\det(I-zB_R)
=
\left(1-\frac34z\right)
\prod_{\substack{p\ne\{1\}\\p\subseteq V_R}}
\left(1-\lambda_p z^{\ell(p)}\right).
\tag{6.2}
\]

Therefore the aperiodic part of Collatz is exactly equivalent to

\[
\boxed{
\frac{\det(I-zB_R)}{1-\frac34z}
\ne0
\quad
\text{for every }R
\text{ and every }1<|z|\leq\frac{10}{9}.
}
\tag{6.3}
\]

The corresponding trace identity is

\[
-\frac{d}{dz}\log\det(I-zB_R)
=
\sum_{j\geq1}\operatorname{tr}(B_R^j)z^{j-1},
\]

where

\[
\operatorname{tr}(B_R^j)
=
\sum_{\substack{n\in V_R\\U^j(n)=n}}
\frac{3^j}{2^{A_j(n)}}.
\]

All terms are nonnegative.

This demonstrates why positivity does not solve the problem. Positivity says that a nontrivial cycle adds positive spectral mass; it supplies no reason that the extra mass must vanish. Closely related unweighted finite-matrix formulations already characterize nontrivial-cycle exclusion by trace vanishing or nilpotency. [\hypertarget{cit:R21:4}{}\hyperlink{ref:R21:4}{4}]


\bigskip



\subsubsection*{7. Why a branch-space quantum graph overcounts}


Suppose instead of putting the positive integers into the vertex set, we use a small symbolic graph with one edge for each exponent \(a\).

Then every word \(w\) is a closed symbolic path. Its affine map has fixed point

\[
x_w=\frac{C_w}{2^A-3^j}.
\]

Whenever \(d_w=2^A-3^j>0\), this fixed point passes the three most natural local tests automatically:

\[
x_w>0,
\]

\[
x_w\in\mathbb Z_2,
\]

and

\[
x_w\in\mathbb Z_3.
\]

The last two statements follow because

\[
d_w=2^A-3^j
\]

is odd and is not divisible by \(3\). Thus \(d_w\) is a unit in both \(\mathbb Z_2\) and \(\mathbb Z_3\).

For example, take

\[
w=(1,3).
\]

Then

\[
A=4,\qquad C_w=3+2=5,\qquad d_w=16-9=7,
\]

so

\[
x_w=\frac57.
\]

It is positive, \(2\)-adically integral, and \(3\)-adically integral, and it has the periodic affine orbit

\[
\frac57\longmapsto\frac{11}{7}\longmapsto\frac57.
\]

But it is not an integer Collatz cycle.

The actual obstruction is at the varying primes dividing \(d_w\), here \(7\). In general,

\[
x_w\in\mathbb Z
\quad\Longleftrightarrow\quad
v_p(C_w)\geq v_p(d_w)
\ \text{for every prime }p.
\]

Thus a fixed real/\(2\)-adic/\(3\)-adic trace formula accepts essentially every positive-drift word. To obtain the correct cycle set, one must insert the all-prime integrality projector

\[
\frac1{d_w}
\sum_{r\bmod d_w}
e^{2\pi i rC_w/d_w}.
\]

But \(d_w\) depends on the entire path. An edge-local quantum graph can evaluate this only by carrying enough internal state to reconstruct

\[
3^j,\quad 2^A,\quad C_w,
\]

which is precisely the original affine arithmetic.

This is the first tautology obstruction:

\[
\begin{array}{c}
\text{small symbolic graph}\\
\Downarrow\\
\text{many false closed paths}\\
\Downarrow\\
\text{must restore }d_w\mid C_w
\end{array}
\]

whereas

\[
\begin{array}{c}
\text{integer Collatz graph}\\
\Downarrow\\
\text{closed paths are genuine cycles}\\
\Downarrow\\
\text{spectral determinant merely counts them}.
\end{array}
\]


\bigskip



\subsubsection*{8. A genuine unitary quantum graph cannot realize Collatz propagation}


There is a more basic obstruction.

For the shortcut map

\[
T(n)=
\begin{cases}
n/2,&n\text{ even},\\[1mm]
(3n+1)/2,&n\text{ odd},
\end{cases}
\]

we have

\[
T(3)=5,\qquad T(10)=5,\qquad T(5)=8.
\]

An exact deterministic quantum routing rule would have to send both orthogonal incoming channels

\[
3\longrightarrow5
\quad\text{and}\quad
10\longrightarrow5
\]

onto the same outgoing channel

\[
5\longrightarrow8.
\]

No unitary matrix can do this: unitary maps preserve inner products, while two orthogonal vectors cannot both be sent to scalar multiples of one vector.

Equivalently, if

\[
P_{m,n}=\mathbf1_{\{m=T(n)\}},
\]

then every column of \(P\) sums to \(1\), but a row sums to the number of preimages of \(m\), which can be \(2\). Hence \(P\) is not bistochastic.

For an ordinary closed quantum graph, the quantum evolution matrix is unitary, and taking absolute squares produces a bistochastic classical transition matrix. [\hypertarget{cit:R21:5}{}\hyperlink{ref:R21:5}{5}, \hypertarget{cit:R21:6}{}\hyperlink{ref:R21:6}{6}]

One can restore unitarity only by:

\begin{itemize}
\item adding auxiliary channels that remember which predecessor was used;
\item allowing information to leak into leads;
\item or replacing unitary scattering by a nonunitary transfer operator.
\end{itemize}

The first option makes the channel space grow into the full space of histories. The second inserts spectral features unrelated to Collatz. The third abandons the self-adjoint quantum-graph positivity that was supposed to provide the rigidity.

Moreover, standard metric-graph edges support propagation in both directions. The choice of an initial orientation does not make the Schr\"odinger propagation one-way; each bond has two directed wave channels. [\hypertarget{cit:R21:7}{}\hyperlink{ref:R21:7}{7}] Consequently, turning inverse Collatz branches into ordinary bonds introduces reverse paths and immediate backtracking paths that are not Collatz trajectories.


\bigskip



\subsubsection*{9. Affine phases are gauge-trivial on a tree}


Let a cycle-free graph have edge phases \(\omega_e\in U(1)\). Choose a root \(o\). For every vertex \(v\), define \(g_v\) as the product of the phases along the unique path from \(o\) to \(v\).

Conjugating by the diagonal unitary

\[
e_v\longmapsto g_v e_v
\]

removes every edge phase.

Thus:

\[
\boxed{\text{All scalar edge phases on a tree are gauge-trivial.}}
\]

The same statement holds for a nonabelian connection on a finite tree: choose vertex gauges recursively and every edge holonomy becomes the identity.

Only products around cycles survive gauge transformation.

Therefore an affine phase can carry gauge-invariant information about a Collatz \textbf{cycle}, but not about an acyclic escaping ray. On a ray, the phases can be removed.

This is the second tautology obstruction: the only gauge-invariant affine data are precisely the closed-path holonomies whose existence is already the cycle question.


\bigskip



\subsubsection*{10. The resonance-at-infinity proposal fails on the powers-of-two chain}


The most serious problem concerns divergent trajectories.

Remove the trivial shortcut cycle \(\{1,2\}\), and define the transient forward operator \(Q\) on

\[
\ell^2(\mathbb N\setminus\{1,2\})
\]

by

\[
Qe_n=
\begin{cases}
e_{T(n)},&T(n)\notin\{1,2\},\\
0,&T(n)\in\{1,2\}.
\end{cases}
\tag{10.1}
\]

Every integer has at most two shortcut preimages, so \(Q\) is bounded.

Now consider

\[
4,8,16,32,\ldots .
\]

Put

\[
h_k=e_{2^{k+2}},\qquad k\geq0.
\]

Then

\[
Qh_0=0,\qquad Qh_k=h_{k-1}\quad(k\geq1).
\]

Thus \(Q\) contains the unilateral backward shift.

For every \(|\lambda|<1\),

\[
f_\lambda=\sum_{k=0}^{\infty}\lambda^k h_k
\]

belongs to \(\ell^2\), and

\[
Qf_\lambda=\lambda f_\lambda.
\]

Consequently,

\[
\boxed{\overline{\mathbb D}\subseteq\sigma(Q).}
\tag{10.2}
\]

This conclusion is unconditional. Every number \(2^{k+2}\) reaches the trivial cycle after finitely many steps, yet those entirely harmless finite trajectories already force spectrum all the way to the unit circle.

Adding unit-modulus phases does not help. A weighted chain

\[
Q_\theta h_k=e^{i\theta_k}h_{k-1}
\]

is diagonally unitarily equivalent to the ordinary backward shift.

A hypothetical divergent orbit would also generate a unilateral-shift-type channel, with essentially the same spectral set. Therefore:

\[
\boxed{
\text{Spectrum or resonances near }|z|=1
\text{ do not distinguish divergence from long convergent branches.}
}
\]

This is especially damaging to the proposed interpretation of ``resonances near the unit circle.'' Open quantum-graph scattering matrices are built from an internal resolvent of the form

\[
(I-U_{\mathrm{int}}(k))^{-1},
\]

so poles and large resolvent responses reflect trapping or long internal propagation. [\hypertarget{cit:R21:8}{}\hyperlink{ref:R21:8}{8}] But the Collatz graph already contains an infinite family of arbitrarily long, certainly terminating paths.

There is also a finite-dimensional version of the problem. The length-\(L\) powers-of-two chain is a nilpotent Jordan block. Its determinant is identically

\[
\det(I-zJ_L)=1,
\]

so its trace formula sees no periodic orbit at all, even though its resolvent becomes very large near the unit circle as \(L\) increases. Long finite survival is therefore pseudospectral rather than periodic-orbit spectral.


\bigskip



\subsubsection*{11. Local nilpotence is the actual escape statement}


On the algebraic space \(c_{00}\) of finitely supported sequences,

\[
Q^k e_n=0
\]

exactly when the orbit of \(n\) has entered the trivial cycle by time \(k\).

Thus the full Collatz conjecture is equivalent to

\[
\boxed{
Q\text{ is locally nilpotent on }c_{00}.
}
\tag{11.1}
\]

``Locally nilpotent'' means that every vector is killed by some power depending on that vector.

It is not equivalent to:

\[
Q^K=0\quad\text{for one uniform }K,
\]

nor to

\[
r(Q)=0,
\]

nor to a spectral gap at the unit circle.

The powers-of-two chain proves this sharply:

\[
Q\text{ is locally nilpotent on each }h_k,
\]

while

\[
\overline{\mathbb D}\subseteq\sigma(Q).
\]

Hilbert-space spectral theory takes a completion over all starting points and therefore replaces pointwise finite stopping by uniform/asymptotic information. That completion loses exactly the quantifier distinction central to Collatz:

\[
\forall n\ \exists k(n)
\qquad\text{versus}\qquad
\exists K\ \forall n.
\]

A trace formula is naturally suited to the second type of assertion, not the first.


\bigskip



\subsubsection*{12. Why spectral positivity cannot finish the proof}


For the finite weighted matrix \(B_R\),

\[
\operatorname{tr}(B_R^j)
=
\left(\frac34\right)^j
+
\sum_{\substack{p\ne\{1\}\\\ell(p)\mid j}}
\ell(p)\lambda_p^{\,j/\ell(p)}.
\tag{12.1}
\]

Every hypothetical nontrivial cycle contributes a strictly positive term.

Therefore positivity gives the implication

\[
\text{nontrivial cycle}
\Longrightarrow
\operatorname{tr}(B_R^j)>
\left(\frac34\right)^j
\]

for appropriate \(j\).

But proving the reverse equality

\[
\operatorname{tr}(B_R^j)=
\left(\frac34\right)^j
\quad\text{for every }R,j
\tag{12.2}
\]

is exactly proving that no nontrivial cycle exists.

Likewise, proving

\[
\frac{\det(I-zB_R)}{1-\frac34z}
\ne0
\qquad
\left(1<|z|\leq\frac{10}{9}\right)
\tag{12.3}
\]

is exactly the cycle-exclusion problem.

There is no general positivity theorem that implies (12.2) or (12.3). Quantum graphs with one loop, ten loops, or arbitrarily complicated cycle structures all have legitimate positive spectral densities. The periodic-orbit terms explain their spectra; positivity does not eliminate them.


\bigskip



\subsection*{Final assessment}


The quantum-graph program yields a rigorous and fairly elegant \textbf{cycle detector}:

\[
\det(I-zB_R)
=
\left(1-\frac34z\right)
\prod_{p\ne\{1\}}
\left(1-\frac{3^{\ell(p)}}{2^{A(p)}}z^{\ell(p)}\right),
\]

and every hypothetical nontrivial positive cycle forces zeros in

\[
1<|z|\leq\frac{10}{9}.
\]

It also yields an exact affine character formula

\[
\mathbf1_{2^A-3^j\mid C_w}
=
\frac1{2^A-3^j}
\sum_{r\bmod(2^A-3^j)}
e^{2\pi i rC_w/(2^A-3^j)}.
\]

Those are substantive, exact reductions.

But the proposed proof does not close:

\begin{itemize}
\item On the integer graph, the trace formula simply enumerates the cycles already present.
\item On a compressed word graph, the all-prime integrality projector is exactly the original Diophantine obstruction.
\item Unitary quantum scattering cannot realize the noninjective Collatz map without extra memory or leakage.
\item Affine phases are gauge-trivial on acyclic branches.
\item Unit-circle spectrum is already forced by the convergent chain \(2^k\to2^{k-1}\to\cdots\to1\).
\item Full convergence is local nilpotence, which ordinary global spectral data do not detect.
\end{itemize}

A genuinely new completion would need an arithmetic construction that simultaneously removes the universal inverse-tree shift and proves cancellation of the word-dependent modulus-\((2^A-3^j)\) character sums \textbf{without first knowing which integers lie in the basin of \(1\)}. That would be closer to an adelic transfer-operator or exact arithmetic-renormalization theorem than to a conventional quantum-graph trace formula.
\subsection*{References}
\begin{enumerate}[leftmargin=2.2em]
\item \hypertarget{ref:R21:1}{}\hyperlink{cit:R21:1}{\texttt{turn424254search2}}
\item \hypertarget{ref:R21:2}{}\hyperlink{cit:R21:2}{\texttt{turn424254search6}}
\item \hypertarget{ref:R21:3}{}\hyperlink{cit:R21:3}{\texttt{turn346121view0}}
\item \hypertarget{ref:R21:4}{}\hyperlink{cit:R21:4}{\texttt{turn921605view0}}
\item \hypertarget{ref:R21:5}{}\hyperlink{cit:R21:5}{\texttt{turn648822view1}}
\item \hypertarget{ref:R21:6}{}\hyperlink{cit:R21:6}{\texttt{turn648822view3}}
\item \hypertarget{ref:R21:7}{}\hyperlink{cit:R21:7}{\texttt{turn392110view0}}
\item \hypertarget{ref:R21:8}{}\hyperlink{cit:R21:8}{\texttt{turn648822view2}}
\end{enumerate}

\clearpage
\section{Round 22: approach 17 --- Tropicalize the cycle and near-cycle equations}
\subsection*{Prompt}
\begin{quote}
Complete a full proof of the Collatz conjecture using the following approach:

\begin{enumerate}
\item Tropicalize the cycle and near-cycle equations
\end{enumerate}

\textbf{Speculative \textperiodcentered{} B}

The cycle relation (2\textasciicircum{}A - 3\textasciicircum{}j) n = C\_w contains competing exponential scales and large sums. Tropicalization records which terms dominate at every place.

Classify all possible tropical rays representing infinite near-cycles and prove that none lift to integral solutions. A lift of a nontrivial tropical ray would instead give a sharply structured family in which to search for a cycle or divergent orbit.

Take as long as you need.
\end{quote}
\subsection*{Response}

\subsection*{Tropicalizing the Collatz cycle equations}


Carrying out the proposed program does \textbf{not} produce a valid proof of the Collatz conjecture. It produces something more precise than a generic ``remaining gap'': a no-go theorem showing that ordinary tropicalization discards exactly the residue information that distinguishes a positive-integer Collatz orbit from the many rational and \(2\)-adic lifts that always exist.

As of Tao's July 2026 revision of his Syracuse-orbit work, the full conjecture remains unresolved; his result is an almost-all theorem rather than a pointwise termination theorem. [\hypertarget{cit:R22:1}{}\hyperlink{ref:R22:1}{1}, \hypertarget{cit:R22:2}{}\hyperlink{ref:R22:2}{2}]

What follows is the complete tropical analysis, including a classification of the possible rays and the exact lifting obstruction.


\bigskip



\subsubsection*{1. Exact valuation-word equations}


Work with the accelerated odd Collatz map

\[
U(x)=\frac{3x+1}{2^{\nu_2(3x+1)}},
\qquad x\in 2\mathbb Z+1.
\]

For an orbit segment

\[
x_0=n,\qquad x_{t+1}=U(x_t),
\]

write

\[
a_t=\nu_2(3x_t+1)\geq 1,
\qquad
A_t=\sum_{i=0}^{t-1}a_i,
\qquad
A=A_j.
\]

For the word

\[
w=(a_0,\ldots,a_{j-1}),
\]

define

\[
C_w
=
\sum_{t=0}^{j-1}3^{j-1-t}2^{A_t}.
\]

Repeated substitution gives the exact affine identity

\[
\boxed{\,2^A U^j(n)=3^j n+C_w.\,}
\tag{1}
\]

This is the standard affine-word formulation of the Syracuse map. [\hyperlink{ref:R22:2}{2}]

Let

\[
D_w=2^A-3^j.
\]

If the word closes into a cycle, \(U^j(n)=n\), then

\[
\boxed{\,D_w n=C_w.\,}
\tag{2}
\]

Because \(C_w>0\), a positive cycle necessarily has

\[
D_w>0,
\qquad\text{equivalently}\qquad
A>j\log_2 3.
\tag{3}
\]


\paragraph*{Exact cycle criterion}


A valuation word \(w\) produces a positive-integer cycle precisely when

\[
\boxed{\,D_w>0\quad\text{and}\quad D_w\mid C_w.\,}
\tag{4}
\]

The sufficiency is worth proving because it identifies the precise lifting problem.

For each cyclic rotation \(w^{(r)}\), let \(C_r=C_{w^{(r)}}\). Directly from the sum defining \(C_w\),

\[
3C_r+D_w=2^{a_r}C_{r+1},
\tag{5}
\]

where indices are cyclic. If \(D_w\mid C_0\), then (5), together with the fact that \(D_w\) is odd, implies successively that \(D_w\mid C_r\) for every \(r\). Put

\[
n_r=\frac{C_r}{D_w}.
\]

Then

\[
3n_r+1=2^{a_r}n_{r+1}.
\tag{6}
\]

Both \(C_r\) and \(D_w\) are odd, so every \(n_r\) is odd. Consequently,

\[
\nu_2(3n_r+1)=a_r,
\]

and the branch word is realized exactly.

Thus the nontrivial-cycle part of Collatz is exactly the assertion

\[
2^A-3^j\nmid C_w
\]

for every nontrivial word with \(2^A>3^j\).


\bigskip



\subsubsection*{2. Archimedean tropicalization}


Set

\[
\alpha=\log_2 3,
\qquad
h_t=A_t-\alpha t,
\qquad
H=h_j=A-\alpha j.
\]

The quantity \(H\) is the total logarithmic drift of the word.

Since

\[
3^{j-1-t}2^{A_t}
=
3^{j-1}2^{h_t},
\]

we have

\[
C_w=3^{j-1}\sum_{t=0}^{j-1}2^{h_t},
\tag{7}
\]

while

\[
D_w=3^j(2^H-1).
\tag{8}
\]

Therefore every supercritical word \(H>0\) has a positive rational periodic seed

\[
\boxed{\,
n_w
=
\frac{C_w}{D_w}
=
\frac{\displaystyle\sum_{t=0}^{j-1}2^{h_t}}
     {3(2^H-1)}.
\,}
\tag{9}
\]

This is already an important obstruction to the proposed proof: \textbf{every supercritical word lifts to a positive rational periodic solution of its branch equations}. The issue is not existence of a lift; it is integrality.

Let

\[
M=\max_{0\leq t<j}h_t.
\]

Since

\[
2^M
\leq
\sum_{t<j}2^{h_t}
\leq
j2^M,
\]

equation (9) gives

\[
\boxed{
\frac{2^M}{3(2^H-1)}
\leq
n_w
\leq
\frac{j2^M}{3(2^H-1)}.
}
\tag{10}
\]

Thus the seed size has exactly two tropical sources:

\begin{enumerate}
\item a large interior maximum \(M\), which may be called a \textbf{tropical mountain};
\item a small difference \(2^A-3^j\), or equivalently small \(H>0\), which may be called \textbf{Diophantine pinching}.
\end{enumerate}

The tropical maximum detects the first phenomenon. Ordinary max-plus tropicalization does not adequately detect the second, because it occurs through cancellation between \(2^A\) and \(3^j\).


\bigskip



\subsubsection*{3. Classification of all archimedean rays}


Write

\[
A_t=t+E_t,
\qquad
E_t=\sum_{i<t}(a_i-1).
\]

The sequence \(E_t\) is nondecreasing. For words \(w_k\) of lengths \(j_k\to\infty\), define

\[
e_k(s)=\frac{E_{\lfloor j_ks\rfloor}}{j_k},
\qquad 0\leq s\leq1.
\]

After taking a subsequence, the monotonicity of \(e_k\) gives convergence at continuity points to a nondecreasing function \(e\). Consequently every finite-height tropical path has the form

\[
\boxed{\,
\gamma(s)=(1-\alpha)s+e(s),
\qquad e\text{ nondecreasing}.
\,}
\tag{11}
\]

Conversely, such profiles can be approximated by valuation words by approximating \(je(s)\) with nondecreasing integer sequences.

So this classifies all archimedean valuation-word rays: they are a fixed line of slope \(1-\log_2 3\), plus an arbitrary nondecreasing excess profile.

Let

\[
h=\gamma(1),
\qquad
m=\sup_{0\leq s<1}\gamma(s).
\]

The endpoint \(h\) divides the rays into three fundamental classes.


\begin{center}
\begin{tabular}{@{}p{0.470\linewidth}p{0.470\linewidth}@{}}
\hline
Tropical endpoint & Interpretation \\
\hline
\(h>0\) & Supercritical word; positive rational periodic branch solution exists \\
\(h=0\) & Critical ray; \(2^A\) and \(3^j\) have equal exponential scale \\
\(h<0\) & Subcritical ray; no positive periodic solution, but compatible with exponential growth \\
\hline
\end{tabular}
\end{center}


For \(h>0\), equation (10) gives

\[
\lim_{k\to\infty}\frac{\log_2 n_{w_k}}{j_k}
=
m-h,
\tag{12}
\]

provided the indicated limits exist.

On the critical wall \(h=0\), one needs an additional pinching coordinate

\[
\pi_k
=
-\frac1{j_k}\log_2(2^{H_k}-1).
\tag{13}
\]

Then

\[
\frac{\log_2 n_{w_k}}{j_k}
=
m_k+\pi_k+o(1).
\tag{14}
\]

This coordinate is not determined by the limiting max-plus path \(\gamma\). Two words can have the same tropical path and radically different values of \(2^A-3^j\).

Therefore even at the real place, the ordinary tropical ray is not enough: one must retain a ``phase of cancellation'' between the two dominant exponentials.


\bigskip



\subsubsection*{4. Every large cycle ray is critical}


Suppose \(x_0,\ldots,x_{j-1}\) is an actual positive cycle. From

\[
2^{a_t}x_{t+1}=3x_t+1
\]

we obtain

\[
2^{a_t}
=
3\frac{x_t}{x_{t+1}}
\left(1+\frac1{3x_t}\right).
\]

Multiplying around the cycle gives the exact product identity

\[
\boxed{\,
2^A
=
3^j\prod_{t=0}^{j-1}
\left(1+\frac1{3x_t}\right).
\,}
\tag{15}
\]

Taking logarithms,

\[
H\log 2
=
\sum_{t=0}^{j-1}
\log\left(1+\frac1{3x_t}\right)>0.
\tag{16}
\]

In particular, every positive cycle is supercritical.

If a sequence of distinct cycles existed, their minimum elements would tend to infinity: only finitely many integer states lie below any fixed bound, and two cycles cannot share a state. Consequently,

\[
\frac{H}{j}
=
\frac1{j\log2}
\sum_{t<j}
\log\left(1+\frac1{3x_t}\right)
\longrightarrow0.
\tag{17}
\]

Hence:

\[
\boxed{\text{Every asymptotic ray of distinct nontrivial cycles lies on }
A/j\to\log_2 3\text{ from above}.}
\tag{18}
\]

So the cycle-ray classification is very sharp:

\begin{itemize}
\item isolated finite cycles correspond to finite supercritical words;
\item any hypothetical unbounded family of cycles accumulates on the critical wall;
\item their sizes are produced by tropical mountains, denominator pinching, or both.
\end{itemize}

This is useful structural information, but it does not exclude the rays.

Existing cycle-exclusion methods already rule out many bounded-complexity rays. In the standard \(m\)-cycle decomposition by monotone blocks or local minima, Simons and de Weger excluded \(m\leq68\), and Hercher's later computation and analysis extends the exclusion to \(m\leq91\). [\hypertarget{cit:R22:3}{}\hyperlink{ref:R22:3}{3}, \hypertarget{cit:R22:4}{}\hyperlink{ref:R22:4}{4}]


\bigskip



\subsubsection*{5. Tropicalization at finite primes is vacuous}


For a fixed word, write the cycle equation as the polynomial relation

\[
F_w(n)
=
2^A n-3^j n-
\sum_{t=0}^{j-1}3^{j-1-t}2^{A_t}
=0.
\tag{19}
\]

In ordinary nonarchimedean tropical geometry, a necessary condition for a sum to vanish is that the minimum valuation among its terms be attained at least twice. This is the usual bend-locus condition. [\hypertarget{cit:R22:5}{}\hyperlink{ref:R22:5}{5}, \hypertarget{cit:R22:6}{}\hyperlink{ref:R22:6}{6}]

For the Collatz cycle equation, that condition holds automatically at every prime.


\paragraph*{At \(p=2\)}


For a cycle seed \(n\), \(n\) is odd. The terms

\[
-3^j n
\quad\text{and}\quad
-3^{j-1}2^{A_0}=-3^{j-1}
\]

both have \(2\)-adic valuation zero. Thus the minimum is already attained twice.


\paragraph*{At \(p=3\)}


We have

\[
D_w=2^A-3^j\not\equiv0\pmod3
\]

and

\[
C_w\equiv2^{A_{j-1}}\not\equiv0\pmod3.
\]

Therefore a possible integer seed \(n=C_w/D_w\) is a \(3\)-adic unit. The terms

\[
2^A n
\quad\text{and}\quad
-2^{A_{j-1}}
\]

both have \(3\)-adic valuation zero.


\paragraph*{At \(p\geq5\)}


Every coefficient \(2^{A_t}3^{j-1-t}\) is a \(p\)-adic unit.

\begin{itemize}
\item If \(p\nmid n\), all terms have valuation zero.
\item If \(p\mid n\), the \(j\) summands comprising \(C_w\) still all have valuation zero.
\end{itemize}

Again, the tropical minimum is attained many times.

Thus:

\[
\boxed{\text{Every valuation word satisfies every ordinary finite-place tropical bend condition.}}
\tag{20}
\]

There is no finite-prime tropical ray to exclude.


\bigskip



\subsubsection*{6. What the finite-place tropicalization loses}


For a positive cycle, the actual condition is

\[
v_p(C_w)\geq v_p(D_w)
\qquad
\text{for every }p\mid D_w.
\tag{21}
\]

For \(p\geq5\), all the individual summands in \(C_w\) are \(p\)-adic units. A positive value of \(v_p(C_w)\) is produced entirely by cancellation among their residue classes.

Tropicalization retains the valuations

\[
v_p\!\left(3^{j-1-t}2^{A_t}\right)=0
\]

but discards the unit phases

\[
3^{j-1-t}2^{A_t}\pmod {p^r}.
\]

Those phases are exactly what determine whether

\[
C_w\equiv0\pmod {p^r}.
\]

If one enriches the tropicalization by retaining enough initial forms or residue data to determine these cancellations, the lifting condition becomes

\[
C_w\equiv0\pmod{D_w},
\tag{22}
\]

which is just the original cycle equation.

So there is a precise dichotomy:

\[
\begin{array}{ll}
\text{ordinary tropicalization} &
\text{is too weak and accepts every word},\\[2mm]
\text{full phase-enriched tropicalization} &
\text{contains }D_w\mid C_w\text{ and restates the hard problem}.
\end{array}
\tag{23}
\]


\bigskip



\subsubsection*{7. Every finite tropical ray has positive integral lifts}


There is an even stronger obstruction.

For every finite word \(w\), the set of integer seeds realizing \(w\) is one residue class modulo \(2^{A+1}\):

\[
\boxed{\,
n\equiv
3^{-j}(2^A-C_w)
\pmod {2^{A+1}}.
\,}
\tag{24}
\]

Indeed, if \(w\) is realized, then (1) and the oddness of \(U^j(n)\) give

\[
3^j n+C_w\equiv2^A\pmod {2^{A+1}},
\]

which yields (24).

Conversely, suppose (24) holds. For \(j=1\), it says

\[
3n+1\equiv2^{a_0}\pmod {2^{a_0+1}},
\]

so \(\nu_2(3n+1)=a_0\). For a longer word, use

\[
C_w=3^{j-1}+2^{a_0}C_{\mathrm{tail}}
\]

and divide the congruence after the first exact branch; induction gives the entire word.

Therefore:

\[
\boxed{\text{Every finite valuation word is realized by infinitely many positive integers.}}
\tag{25}
\]

No finite-depth tropical obstruction can eliminate even one valuation word.

This also shows why a compactness argument is dangerous: compatible positive seeds may escape to infinity while converging to a nonpositive \(2\)-adic integer.


\bigskip



\subsubsection*{8. Every infinite ray has a unique \(2\)-adic lift}


For an infinite word

\[
a=(a_0,a_1,\ldots),
\]

the residue classes (24) for its successive prefixes are nested. Since \(A_j\to\infty\), they determine a unique

\[
n_a\in\mathbb Z_2.
\]

Thus:

\[
\boxed{\text{Every infinite valuation ray has a unique exact }2\text{-adic Collatz seed}.}
\tag{26}
\]

This is closely related to the known conjugacy between the \(3x+1\) transformation on \(\mathbb Z_2\) and the binary shift. [\hypertarget{cit:R22:7}{}\hyperlink{ref:R22:7}{7}]

The positive-integer problem is therefore not a question of whether a symbolic ray lifts. All rays lift. It is the much narrower question

\[
n_a\in\mathbb Z_{>0}\;?
\tag{27}
\]


\paragraph*{Two illustrative rays}


The constant word

\[
(2,2,2,\ldots)
\]

lifts to

\[
n_a=1,
\]

the trivial fixed point.

The constant word

\[
(1,1,1,\ldots)
\]

lifts to

\[
n_a=-1\in\mathbb Z_2.
\]

Its length-\(j\) prefix is realized by the positive integer

\[
n_j=2^{j+1}-1.
\]

Indeed, for \(0\leq t\leq j\),

\[
U^t(n_j)=3^t2^{j+1-t}-1,
\]

and every one of the first \(j\) valuations equals \(1\).

Hence this single tropical ray has:

\begin{itemize}
\item positive integral lifts at every finite depth;
\item a perfectly valid infinite \(2\)-adic lift;
\item no fixed positive-integer infinite lift.
\end{itemize}

The distinction is completely invisible to ordinary tropical geometry.


\bigskip



\subsubsection*{9. Every supercritical periodic ray also lifts rationally}


For a periodic word \(w\) with \(D_w>0\), equation (9) gives

\[
n_w=\frac{C_w}{D_w}>0.
\]

Because \(C_w\) and \(D_w\) are odd, this is a \(2\)-adic unit. The rotated values

\[
n_r=\frac{C_r}{D_w}
\]

satisfy

\[
3n_r+1=2^{a_r}n_{r+1}.
\]

Thus every supercritical word has a positive rational periodic branch orbit. It is integral exactly when

\[
D_w\mid C_w.
\]

So the full hierarchy of lifts is

\[
\begin{array}{c|c}
\text{object} & \text{existence}\\ \hline
\text{finite positive-integer orbit segment} & \text{always}\\
\text{infinite }2\text{-adic orbit} & \text{always}\\
\text{positive rational periodic branch orbit} & H>0\\
\text{positive-integer periodic orbit} & H>0\text{ and }D_w\mid C_w.
\end{array}
\tag{28}
\]

Only the last line is difficult, and tropical valuations do not distinguish it from the preceding lines.


\bigskip



\subsubsection*{10. A genuine theorem: eventually periodic rays cannot diverge}


The tropical/adic framework does prove a useful exact result.


\paragraph*{Theorem}


A positive-integer Collatz orbit with an eventually periodic valuation word is eventually periodic as an integer orbit. In particular, no divergent positive orbit can have an eventually periodic valuation word.


\paragraph*{Proof}


Suppose that after some initial segment, the valuation word repeats a block

\[
v=(b_0,\ldots,b_{r-1})
\]

with total valuation

\[
B=\sum_{i=0}^{r-1}b_i.
\]

Let \(x\in\mathbb Z_2\) be the seed at the start of the repeating tail. Both \(x\) and \(U^r(x)\) have the same infinite valuation word. By uniqueness of the \(2\)-adic seed associated with an infinite word,

\[
U^r(x)=x.
\]

Therefore

\[
(2^B-3^r)x=C_v.
\tag{29}
\]

If \(B<r\log_2 3\), the right-hand side is positive and the denominator is negative, so \(x<0\).

If \(B>r\log_2 3\), then \(x\) is the positive rational periodic solution. When \(x\) is a positive integer, it is an actual cycle.

Equality \(2^B=3^r\) is impossible.

Thus a positive-integer orbit with an eventually periodic valuation tail must enter an integer cycle. \ensuremath{\blacksquare}

Consequently, a hypothetical divergent Collatz orbit must have a genuinely aperiodic symbolic ray.

This excludes one natural class of tropical lifts, but not the general class.


\bigskip



\subsubsection*{11. Tropical rays of divergent orbits}


Suppose a positive-integer orbit is unbounded and does not enter a cycle. Then necessarily

\[
x_k\longrightarrow\infty.
\tag{30}
\]

Otherwise some finite set of integers would be visited infinitely often, forcing a repeated state and hence a cycle.

At every step,

\[
\log x_{t+1}-\log x_t
=
\log3-a_t\log2
+
\log\left(1+\frac1{3x_t}\right).
\]

Summing gives

\[
\log x_k-\log x_0
=
k\log3-A_k\log2
+
\sum_{t=0}^{k-1}
\log\left(1+\frac1{3x_t}\right).
\tag{31}
\]

Because \(x_t\to\infty\), the final summand has Ces\`aro average tending to zero. Hence

\[
\boxed{\,
\frac{A_k}{k}
=
\log_2 3
-
\frac{\log_2 x_k}{k}
+o(1).
\,}
\tag{32}
\]

It follows that

\[
\limsup_{k\to\infty}\frac{A_k}{k}
\leq \log_2 3.
\tag{33}
\]

There are two possible tropical divergence classes.


\paragraph*{Exponential escape}


If

\[
\frac{\log_2 x_k}{k}\longrightarrow\lambda>0,
\]

then

\[
\frac{A_k}{k}
\longrightarrow
\log_2 3-\lambda,
\]

so the ray is strictly subcritical.


\paragraph*{Subexponential escape}


If

\[
\log x_k=o(k),
\]

then

\[
\frac{A_k}{k}\longrightarrow\log_2 3,
\]

so the ray lies on the critical wall, now approached from the nonpositive-drift side.

Thus the asymptotic classification is

\[
\boxed{
\begin{aligned}
\text{large cycles} &:\quad H_k/k\to0^+,\\
\text{subexponential divergence} &:\quad H_k/k\to0,\\
\text{exponential divergence} &:\quad H_k/k\to-\lambda<0.
\end{aligned}}
\tag{34}
\]


\bigskip



\subsubsection*{12. The near-cycle extraction step is not valid}


The proposed approach assumes that a divergent orbit should generate arbitrarily long near-cycles. That does not follow from divergence.

For a block beginning at \(x_r\), with word \(v\), total \(B\), length \(j\), and drift

\[
H_v=B-j\log_2 3,
\]

the exact block identity is

\[
\frac{x_{r+j}}{x_r}
=
2^{-H_v}
+
\frac{C_v}{2^B x_r}.
\tag{35}
\]

A near-cycle \(x_{r+j}/x_r\approx1\) can arise in two ways:

\begin{enumerate}
\item \textbf{drift balance:} \(H_v\approx0\) and the affine correction is small;
\item \textbf{mountain balance:} the affine correction compensates a nonzero drift.
\end{enumerate}

But a strictly subcritical exponentially escaping ray need not contain blocks with ratio close to \(1\). It may have persistent positive real growth. Tropical compactness produces a limiting symbolic profile, not recurrence in real height.

Therefore a proof based only on near-cycle equations would first need a new theorem:


\begin{quote}
Every divergent positive-integer Collatz orbit contains arbitrarily long blocks whose endpoint heights nearly recur, or whose affine offsets satisfy an equally restrictive compensating relation.

\end{quote}


That theorem is not supplied by tropicalization, and equation (32) alone does not imply it.


\bigskip



\subsubsection*{13. The \(3\)-adic suffix phase is real information, but not enough}


There is an exact \(3\)-adic phase relation that ordinary tropicalization omits.

Suppose a suffix word \(v\) of length \(r\) maps \(y\) to \(x\). Then

\[
2^{A(v)}x=3^r y+C_v.
\]

Modulo \(3^r\),

\[
\boxed{\,
x\equiv2^{-A(v)}C_v\pmod {3^r}.
\,}
\tag{36}
\]

In particular,

\[
\boxed{\,
\nu_3(x+1)\geq r
\iff
3^r\mid C_v+2^{A(v)}.
\,}
\tag{37}
\]

This is the terminal-suffix locality phenomenon: the endpoint modulo \(3^r\) depends only on the last \(r\) valuations, not on the earlier prefix.

It is exactly the kind of phase information an enriched tropical theory would have to preserve. But it still does not decide cycles, because

\[
\gcd(D_w,6)=1.
\]

The primes controlling

\[
D_w\mid C_w
\]

are the varying primes \(p\geq5\) dividing \(2^A-3^j\). Complete cycle lifting requires simultaneous phase control at all of those primes.


\bigskip



\subsubsection*{14. The no-go theorem for bare tropicalization}


The preceding arguments can be collected into one precise statement.


\paragraph*{Tropical no-go theorem}


For the Collatz valuation-word equations:

\begin{enumerate}
\item every finite archimedean tropical ray has the form \[
   \gamma(s)=(1-\log_2 3)s+e(s)
   \] with \(e\) nondecreasing;

\item every finite-place ordinary tropical bend condition is automatically satisfied;

\item every finite valuation word has infinitely many positive-integer lifts;

\item every infinite valuation ray has a unique \(2\)-adic lift;

\item every supercritical periodic word has a positive rational periodic lift;

\item a supercritical periodic word has a positive-integer lift precisely when \[
   2^A-3^j\mid C_w.
   \]
\end{enumerate}

Therefore no argument depending only on the ordinary tropical valuations of the cycle or near-cycle equations can prove that nontrivial rays fail to lift integrally.


\paragraph*{Proof}


Statements 1\textendash{}5 were proved above. Statement 6 is the exact cycle criterion (4). Ordinary tropicalization knows the valuations of the individual summands but not the residue cancellations determining \(v_p(C_w)\) and \(v_p(D_w)\). Hence it cannot decide statement 6. \ensuremath{\blacksquare}


\bigskip



\subsubsection*{15. What a successful replacement theorem would need to say}


The approach can be repaired only by replacing ordinary tropicalization with a \textbf{phase-enriched adelic tropicalization}. Such a theory would need two genuinely new theorems.


\paragraph*{Cycle phase-rigidity theorem}


For every nontrivial supercritical word \(w\),

\[
\boxed{\,
C_w\not\equiv0\pmod {2^{A(w)}-3^{|w|}}.
\,}
\tag{38}
\]

A useful proof would need to derive this from a uniform restriction on the initial forms of the summands of \(C_w\), rather than merely computing the congruence itself.


\paragraph*{Aperiodic positive-lift exclusion theorem}


Let \(a=(a_0,a_1,\ldots)\) be an aperiodic critical or subcritical ray, and let \(n_a\in\mathbb Z_2\) be its unique \(2\)-adic lift. Prove that

\[
\boxed{\,
n_a\notin\mathbb Z_{>0}.
\,}
\tag{39}
\]

Equivalently, prove that for a fixed positive integer \(n\), the nested cylinder congruences

\[
n\equiv
3^{-j}(2^{A_j}-C_j)
\pmod {2^{A_j+1}}
\tag{40}
\]

cannot hold indefinitely along a critical or subcritical aperiodic word unless the orbit eventually descends into the trivial cycle.

This is a pointwise synchronization theorem between:

\begin{itemize}
\item the fixed archimedean integer \(n\);
\item its increasingly precise \(2\)-adic seed cylinders;
\item the \(3\)-adic endpoint/carry phases;
\item the varying \(p\)-adic phases for primes dividing \(2^{A_j}-3^j\).
\end{itemize}

Tao's work obtains strong almost-all orbit descent through probabilistic and \(3\)-adic mixing methods, but it does not provide this uniform pointwise synchronization over every starting cylinder. [\hyperlink{ref:R22:2}{2}]


\bigskip



\subsection*{Conclusion}


The tropical program yields several rigorous deductions:

\[
\begin{gathered}
\text{all archimedean rays are classified by monotone excess profiles},\\
\text{all large cycle rays approach the critical wall from above},\\
\text{divergent rays are critical or subcritical},\\
\text{eventually periodic divergent rays are impossible},\\
\text{every finite ray has positive integral realizations},\\
\text{every infinite ray has a unique }2\text{-adic realization}.
\end{gathered}
\]

But the requested final statement\textemdash{}


\begin{quote}
classify the tropical rays and prove that none except the trivial one lifts to a positive integer\textemdash{}

\end{quote}


is false as a consequence of ordinary tropical data. Tropicalization accepts every word locally and every ray \(2\)-adically. The distinction between the trivial lift \(1\), the nonpositive lift \(-1\), a positive rational periodic lift, and a positive integral cycle is carried by residue phases that tropical valuations erase.

Accordingly, this approach does \textbf{not} complete a proof of Collatz. Its decisive missing result is an adelic phase-rigidity theorem excluding positive-integer lifts of all nontrivial critical, subcritical, and supercritical rays. Proving that theorem would indeed prove Collatz, but it does not follow from the tropical classification itself.
\subsection*{References}
\begin{enumerate}[leftmargin=2.2em]
\item \hypertarget{ref:R22:1}{}\hyperlink{cit:R22:1}{\texttt{turn249030view3}}
\item \hypertarget{ref:R22:2}{}\hyperlink{cit:R22:2}{\texttt{turn322555view0}}
\item \hypertarget{ref:R22:3}{}\hyperlink{cit:R22:3}{\texttt{turn249030view1}}
\item \hypertarget{ref:R22:4}{}\hyperlink{cit:R22:4}{\texttt{turn249030view2}}
\item \hypertarget{ref:R22:5}{}\hyperlink{cit:R22:5}{\texttt{turn337911view0}}
\item \hypertarget{ref:R22:6}{}\hyperlink{cit:R22:6}{\texttt{turn249030view4}}
\item \hypertarget{ref:R22:7}{}\hyperlink{cit:R22:7}{\texttt{turn867756search1}}
\end{enumerate}

\clearpage
\section{Round 23: approach 18 --- A Kolmogorov-complexity contradiction}
\subsection*{Prompt}
\begin{quote}
Complete a full proof of the Collatz conjecture using the following approach:

\begin{enumerate}
\item A Kolmogorov-complexity contradiction
\end{enumerate}

\textbf{Speculative \textperiodcentered{} P}

To grow, an orbit's valuation word must be atypical enough to have deficient powers of two. To remain integrally realizable, its long prefix is nevertheless determined by a finite binary seed residue.

Try to prove that these two requirements force contradictory complexity estimates: positive drift requires a high-complexity sequence, while simultaneous 2-adic and 3-adic compatibility makes it compressible. This must be an effective theorem, not the informal claim that Collatz "looks random."

Take as long as you need.
\end{quote}
\subsection*{Response}

\subsection*{Verdict}


The proposed Kolmogorov-complexity mechanism does \textbf{not} yield a proof of Collatz. Its central pointwise implication,

\[
\text{positive drift}+\text{finite }2\text{-adic/}3\text{-adic compatibility}
\quad\Longrightarrow\quad
K(\text{valuation prefix})\ge c j,
\]

is false.

In fact, there is an exact \textbf{complexity-conservation theorem}:


\begin{quote}
For a length-\(j\) accelerated Collatz valuation word \(w\), the word and its least positive realizing seed \(R(w)\) have the same Kolmogorov complexity up to \(O(1)\), conditional on \(j\).

\end{quote}


Consequently, for every fixed integer seed \(n\),

\[
K\bigl(w_j(n)\mid j\bigr)=O_n(1),
\]

whether its orbit terminates, cycles, or hypothetically diverges.

More decisively, the explicit positive seeds

\[
N_j=2^{j+1}-1
\]

have a strictly growing \(j\)-step orbit, valuation word \(1^j\), and maximal terminal \(3\)-adic compatibility, while

\[
K(1^j\mid j)=O(1).
\]

Thus even the strongest natural finite-prefix version of the proposed lower bound is disproved.

What follows is the complete derivation and the exact replacement target.


\bigskip



\subsubsection*{1. Accelerated valuation words}


Work with the accelerated map on positive odd integers,

\[
U(x)=\frac{3x+1}{2^{\nu_2(3x+1)}}.
\]

Let

\[
x_0=n,\qquad x_{i+1}=U(x_i),
\]

and define

\[
a_i=\nu_2(3x_i+1)\ge 1,
\qquad
A_j=\sum_{i=0}^{j-1}a_i.
\]

The length-\(j\) valuation word is

\[
w_j(n)=(a_0,\ldots,a_{j-1}).
\]

Finite parity vectors and accelerated valuation vectors have exact symbolic-dynamical descriptions: each finite parity vector occurs in one residue class modulo the relevant power of \(2\), while accelerated valuation vectors are parametrized by compositions of their total valuation. [\hypertarget{cit:R23:1}{}\hyperlink{ref:R23:1}{1}, \hypertarget{cit:R23:2}{}\hyperlink{ref:R23:2}{2}]

Define \(C_0=0\) and

\[
C_{j+1}=3C_j+2^{A_j}.
\]

Then induction gives the exact affine identity

\[
\boxed{2^{A_j}x_j=3^j n+C_j}
\tag{1}
\]

with

\[
C_j=\sum_{i=0}^{j-1}3^{j-1-i}2^{A_i}.
\tag{2}
\]

Indeed, if (1) holds at \(j\), then

\[
\begin{aligned}
2^{A_{j+1}}x_{j+1}
&=2^{A_j+a_j}\frac{3x_j+1}{2^{a_j}}\\
&=3\cdot2^{A_j}x_j+2^{A_j}\\
&=3^{j+1}n+3C_j+2^{A_j}.
\end{aligned}
\]


\bigskip



\subsubsection*{2. A valuation word is equivalent to one binary seed residue}


Fix a word

\[
w=(a_0,\ldots,a_{j-1}),\qquad A=A_j.
\]

Since its endpoint \(x_j\) must be odd, (1) implies

\[
3^j n+C_j\equiv 2^A\pmod{2^{A+1}}.
\]

Because \(3^j\) is invertible modulo \(2^{A+1}\), this determines one odd residue:

\[
\boxed{
n\equiv
3^{-j}(2^A-C_j)
\pmod{2^{A+1}}.
}
\tag{3}
\]

Conversely, this residue class realizes exactly the word \(w\). This follows either from the standard symbolic-dynamics theorem or directly by reducing (3) at each intermediate cumulative valuation \(A_i\).

Let

\[
R(w)\in\{1,3,\ldots,2^{A+1}-1\}
\]

be the least positive representative of (3).

Thus:

\begin{itemize}
\item from \(w\), one can compute \(R(w)\);
\item from \((R(w),j)\), one can iterate \(U\) for \(j\) steps and recover \(w\).
\end{itemize}

So \(w\leftrightarrow R(w)\) is a computably reversible coding, conditional on its length.

This is stronger than saying the binary seed ``determines'' the word. It says the seed residue and word contain essentially \textbf{the same algorithmic information}.


\bigskip



\subsubsection*{3. The genuine drift restriction}


Set

\[
\alpha=\log_2 3.
\]

At one accelerated step,

\[
x_{i+1}
=x_i\,2^{\alpha-a_i}
\left(1+\frac1{3x_i}\right),
\]

so

\[
\log_2x_{i+1}-\log_2x_i
=
\alpha-a_i+
\log_2\left(1+\frac1{3x_i}\right).
\]

Summing,

\[
\boxed{
\log_2x_j-\log_2n
=
\alpha j-A_j+E_j,
}
\tag{4}
\]

where

\[
E_j=\sum_{i=0}^{j-1}
\log_2\left(1+\frac1{3x_i}\right)>0.
\]

If an orbit is unbounded and nonperiodic, then \(x_i\to\infty\): otherwise it would visit a finite set infinitely often, hence repeat a state and become periodic. Therefore

\[
\frac{E_j}{j}\longrightarrow0.
\]

Since \(x_j\ge1\), (4) gives

\[
\boxed{
\limsup_{j\to\infty}\frac{A_j}{j}
\le\log_2 3.
}
\tag{5}
\]

So the rigorous version of ``deficient powers of two'' is correct. A divergent orbit must live in the lower large-deviation tail of the valuation sum.

But (5) is not a complexity lower bound.


\bigskip



\subsubsection*{4. What the entropy calculation actually proves}


For a random odd starting value, the finite valuation vector has the exact cylinder density

\[
\Pr(a_0=k_0,\ldots,a_{j-1}=k_{j-1})
=
2^{-\sum k_i}.
\]

Equivalently, the \(a_i\) have the same finite-dimensional distribution as independent geometric variables with

\[
\Pr(a_i=k)=2^{-k},
\qquad
\mathbb E[a_i]=2.
\]

This is an exact initial-orbit density theorem, not merely a heuristic. [\hypertarget{cit:R23:3}{}\hyperlink{ref:R23:3}{3}]

Let

\[
D_j=
\left\{
w\in\mathbb Z_{\ge1}^j:
A(w)\le\lfloor\alpha j\rfloor
\right\}.
\]

The number of compositions of \(m\) into \(j\) positive parts is

\[
\binom{m-1}{j-1}.
\]

Hence the hockey-stick identity gives

\[
|D_j|
=
\sum_{m=j}^{\lfloor\alpha j\rfloor}
\binom{m-1}{j-1}
=
\binom{\lfloor\alpha j\rfloor}{j}.
\]

By Stirling's formula,

\[
\log_2|D_j|
=
h_*j+O(\log j),
\tag{6}
\]

where

\[
h_*=
\alpha H_2(1/\alpha)
=
1.5056438879\ldots
\]

and \(H_2\) is binary entropy.

The natural probability of the low-drift event is

\[
\Pr(D_j)
=
2^{-Ij+O(\log j)},
\tag{7}
\]

where

\[
I=\alpha-h_*
=
0.0793186128\ldots.
\]

This says:

\begin{itemize}
\item low-drift prefixes are exponentially rare among random seeds;
\item most words in \(D_j\) have complexity approximately \(h_*j\);
\item every word in \(D_j\) can be described by its index in \(D_j\), using at most \(h_*j+O(\log j)\) bits.
\end{itemize}

It does \textbf{not} say that every member of \(D_j\) has large complexity. A small effective set may contain extremely simple exceptional elements. Incompressibility arguments prove statements about most elements, not every element.


\bigskip



\subsubsection*{5. Exact complexity-conservation theorem}


Fix a universal prefix-free machine and let \(K\) denote prefix Kolmogorov complexity.


\paragraph*{Theorem 1}


There is a universal constant \(c\) such that, for every length-\(j\) valuation word \(w\),

\[
\boxed{
\left|
K(w\mid j)-K(R(w)\mid j)
\right|\le c.
}
\tag{8}
\]


\paragraph*{Proof}


Given \(w\), an algorithm:

\begin{enumerate}
\item computes \(A(w)\);
\item computes \(C_w\) from the recurrence;
\item computes \(3^{-j}\pmod{2^{A+1}}\);
\item outputs the least representative in (3).
\end{enumerate}

Therefore

\[
K(R(w)\mid j)\le K(w\mid j)+O(1).
\]

Conversely, given \(R(w)\) and \(j\), an algorithm iterates \(U\) exactly \(j\) times and records the valuations. Since \(R(w)\) realizes \(w\), the output is \(w\). Hence

\[
K(w\mid j)\le K(R(w)\mid j)+O(1).
\]

Combining the inequalities proves (8). \ensuremath{\blacksquare}


\paragraph*{Corollary 2}


For every fixed positive odd integer \(n\),

\[
K(w_j(n)\mid n,j)=O(1)
\tag{9}
\]

and

\[
K(w_j(n)\mid j)\le K(n\mid j)+O(1).
\tag{10}
\]

Moreover, because \(A_j\ge j\), eventually

\[
2^{A_j+1}>n.
\]

At that point \(n\) itself is already the least positive representative of its residue class, so

\[
R(w_j(n))=n.
\]

Thus, for all sufficiently large \(j\),

\[
\boxed{
K(w_j(n)\mid j)=K(n\mid j)+O(1).
}
\tag{11}
\]

Unconditionally,

\[
\boxed{
K(w_j(n))
\le
K(n)+K(j)+O(1)
=
O_n(\log j).
}
\tag{12}
\]

This applies to \textbf{every} fixed integer orbit. In particular, a hypothetical divergent integer orbit would still have a valuation prefix of complexity only \(O(\log j)\), or \(O(1)\) when \(j\) is supplied.

So integrality compatibility does not merely make a counterexample word compressible. It makes every fixed integer orbit word compressible.


\bigskip



\subsubsection*{6. The \(3\)-adic state does not add independent information}


Let \(v\) be the last \(r\) entries of a length-\(j\) word. Suppose the associated \(r\)-step segment starts at \(Y\) and ends at \(X\). Its affine identity is

\[
2^{A(v)}X=3^rY+C_v.
\]

Reducing modulo \(3^r\),

\[
\boxed{
X\equiv 2^{-A(v)}C_v\pmod{3^r}.
}
\tag{13}
\]

Consequently,

\[
\boxed{
3^r\mid X+1
\quad\Longleftrightarrow\quad
3^r\mid C_v+2^{A(v)}.
}
\tag{14}
\]

Thus the endpoint residue modulo \(3^r\) is a computable function of the terminal valuation suffix. It is not an independent random string that must be specified in addition to the word.

In complexity notation,

\[
K(X\bmod3^r\mid v,r)=O(1).
\]

More generally,

\[
K\bigl(w,R(w),X\bmod3^r\mid j,r\bigr)
=
K(w\mid j,r)+O(1).
\]

The combined \(2\)-adic and \(3\)-adic state is another computable representation of the same data. It supplies arithmetic restrictions, but not an automatic algorithmic-information contradiction.


\bigskip



\subsubsection*{7. Explicit low-complexity positive-drift shadows}


The following family directly disproves the desired finite-prefix lower bound.


\paragraph*{Theorem 3}


For every \(j\ge1\), let

\[
N_j=2^{j+1}-1.
\]

Then for \(0\le t\le j\),

\[
\boxed{
U^t(N_j)=3^t2^{j+1-t}-1.
}
\tag{15}
\]

In particular:

\begin{enumerate}
\item the first \(j\) valuations are all equal to \(1\);
\item the orbit is strictly increasing for those \(j\) steps;
\item its multiplicative drift is exponentially positive;
\item the endpoint has maximal \(3\)-adic compatibility;
\item the word and seed have \(O(1)\) complexity conditional on \(j\).
\end{enumerate}


\paragraph*{Proof}


Suppose

\[
x_t=3^t2^{j+1-t}-1.
\]

For \(t<j\),

\[
\begin{aligned}
3x_t+1
&=3^{t+1}2^{j+1-t}-2\\
&=2\left(3^{t+1}2^{j-t}-1\right).
\end{aligned}
\]

The expression in parentheses is odd, so

\[
\nu_2(3x_t+1)=1
\]

and

\[
x_{t+1}=3^{t+1}2^{j-t}-1.
\]

This proves (15).

Moreover,

\[
x_{t+1}-x_t
=
3^t2^{j-t}>0,
\]

so the whole block stays above its starting value.

The valuation word is

\[
w_j=1^j,
\qquad
A_j=j.
\]

Since \(j<\alpha j\), the main multiplicative factor is

\[
\frac{3^j}{2^j}
=
\left(\frac32\right)^j.
\]

At the endpoint,

\[
U^j(N_j)=2\cdot3^j-1,
\]

so

\[
\boxed{
\nu_3\bigl(U^j(N_j)+1\bigr)=j.
}
\tag{16}
\]

For the word \(1^r\),

\[
C_{1^r}
=
3^r-2^r,
\]

and hence

\[
C_{1^r}+2^r=3^r,
\]

which is exactly the strongest possible instance of (14).

Finally, programs of constant size, given \(j\), output \(1^j\) and \(2^{j+1}-1\). Therefore

\[
K(1^j\mid j)=O(1),
\qquad
K(N_j\mid j)=O(1).
\]

\ensuremath{\blacksquare}

This family simultaneously has:

\[
\begin{array}{c}
\text{positive integer realizability},\\
\text{strict growth throughout the block},\\
A_j/j=1<\log_2 3,\\
\text{maximal terminal }3\text{-adic divisibility},\\
\text{constant conditional complexity}.
\end{array}
\]

Therefore none of those finite properties can imply an unbounded Kolmogorov-complexity lower bound.


\bigskip



\subsubsection*{8. What the all-ones family is really shadowing}


The residues

\[
N_j=2^{j+1}-1
\]

satisfy

\[
N_j\equiv-1\pmod{2^{j+1}}.
\]

Thus

\[
N_j\longrightarrow-1
\qquad\text{in }\mathbb Z_2.
\]

And \(-1\) is an exact fixed point:

\[
U(-1)=\frac{-3+1}{2}=-1.
\]

For an infinite valuation word \(w=(a_0,a_1,\ldots)\), with \(A_0=0\), the associated \(2\)-adic seed is formally

\[
\boxed{
z(w)
=
-\sum_{i=0}^{\infty}
\frac{2^{A_i}}{3^{i+1}}
\in\mathbb Z_2.
}
\tag{17}
\]

For \(w=1^\infty\),

\[
z(w)
=
-\sum_{i=0}^{\infty}\frac{2^i}{3^{i+1}}
=-1.
\]

This is the geometric reason for Theorem 3: arbitrarily large positive integers can shadow the negative \(2\)-adic fixed point for arbitrarily long finite times.

The full \(2\)-adic Collatz extension is conjugate to a shift, but the ordinary integers form a countable Haar-null subset of \(\mathbb Z_2\). Hence measure-theoretic or randomness conclusions on generic \(2\)-adic points do not determine what happens on ordinary integer seeds. [\hypertarget{cit:R23:4}{}\hyperlink{ref:R23:4}{4}, \hypertarget{cit:R23:5}{}\hyperlink{ref:R23:5}{5}]

The missing distinction is therefore not:

\[
\text{simple versus complex}.
\]

It is:

\[
\text{a positive Archimedean integer}
\quad\text{versus}\quad
\text{a negative or nonintegral }2\text{-adic point}.
\]

Both can have algorithmic dimension zero.


\bigskip



\subsubsection*{9. Resource-bounded complexity does not repair the argument}


One might replace ordinary \(K\) by polynomial-time complexity, logical depth, online complexity, or monotone complexity. The same obstruction remains.

For \(x\ge1\),

\[
U(x)
\le\frac{3x+1}{2}
\le2x.
\]

Therefore

\[
x_i\le2^i n,
\]

so the bit length of the first \(j\) iterates is

\[
O(j+\log n).
\]

Exact iteration for \(j\) steps, including valuation extraction, takes polynomial time in \(j+\log n\). Hence a fixed program with \(n\) hardcoded outputs \(w_j(n)\) efficiently.

Thus:

\begin{itemize}
\item polynomial-time Kolmogorov complexity stays small;
\item an online program can stream the valuation sequence;
\item monotone complexity is bounded by a program encoding \(n\);
\item logical depth cannot arise from expensive simulation, because direct simulation is efficient.
\end{itemize}

Finite-state complexity is different: a computable sequence can be finite-state incompressible. But then the other direction fails\textemdash{}positive drift certainly does not force finite-state complexity, since \(1^\infty\) is periodic.


\bigskip



\subsubsection*{10. Martin-L\"of randomness gives only an almost-everywhere statement}


The geometric law implies that a Martin-L\"of-random valuation sequence has average valuation \(2\), not at most \(\log_2 3\). Thus a divergent Collatz word would fail an effective law of large numbers.

But every fixed integer orbit word is computable. Therefore none is Martin-L\"of random, including the terminating orbits. Randomness theory excludes generic \(2\)-adic seeds from having persistent positive drift; it does not exclude a particular computable exceptional seed.

This is exactly the same countable-exception gap encountered by ordinary ergodic theory.


\bigskip



\subsubsection*{11. Cycles are not covered by the drift argument}


Suppose a positive accelerated cycle has word \(w\), period \(p\), and total valuation \(A\). Then (1) gives

\[
2^A n=3^p n+C_w,
\]

so

\[
\boxed{
(2^A-3^p)n=C_w>0.
}
\tag{18}
\]

Consequently,

\[
\frac Ap>\log_2 3.
\]

Thus a cycle has the opposite valuation inequality from a divergent orbit. Its infinite valuation word is periodic and therefore has

\[
K(w^{m}\mid m)=O(1).
\]

So even a successful complexity argument excluding low-valuation divergent trajectories would not by itself constitute a full Collatz proof; nontrivial cycles would still require exclusion.


\bigskip



\subsubsection*{12. The exact live replacement: Archimedean anti-shadowing}


Define

\[
m_j=
\min\left\{
n>1:\ n\text{ odd and }U^i(n)\ge n
\text{ for }0\le i\le j
\right\},
\]

with \(m_j=\infty\) if the set is empty.

Then \(m_j\) is nondecreasing, and

\[
\boxed{
\text{Collatz is true}
\quad\Longleftrightarrow\quad
m_j\longrightarrow\infty.
}
\tag{19}
\]


\paragraph*{Proof}


If Collatz is true, every fixed \(n>1\) eventually reaches \(1<n\). Therefore, for every bound \(B\), all \(n\le B\) eventually fall below themselves. Taking the maximum of those finitely many descent times shows that \(m_j>B\) for sufficiently large \(j\).

Conversely, suppose \(m_j\) does not tend to infinity. Since it is a nondecreasing sequence of integers, it is eventually bounded by some \(B\). A fixed \(n\le B\) then survives the inequality

\[
U^i(n)\ge n
\]

for arbitrarily large \(i\), and by nestedness for every \(i\). It can never reach \(1\), contradicting Collatz. \ensuremath{\blacksquare}

In word language,

\[
m_j
=
\min\left\{
R(w)>1:
U^i(R(w))\ge R(w)
\text{ for }0\le i\le j
\right\}.
\]

A quantitative theorem such as

\[
m_j\ge F(j),
\qquad F(j)\longrightarrow\infty,
\tag{20}
\]

would prove Collatz.

The intended complexity argument would try to establish

\[
K(R(w)\mid j)\ge c j
\]

for every surviving word and infer

\[
R(w)\ge2^{cj-O(1)}.
\]

But Theorem 3 directly refutes that complexity lower bound: \(1^j\) is a surviving word and

\[
R(1^j)=2^{j+1}-1
\]

has constant complexity conditional on \(j\).

Its \textbf{value} is exponentially large, but its \textbf{description} is tiny. Kolmogorov complexity cannot substitute for the required numerical lower bound.


\bigskip



\subsection*{Final status of approach 18}


The branch should be recorded as:


\begin{quote}
\textbf{DEAD in its stated form.}   Exact obstruction: the valuation prefix and its least realizing binary seed residue are computably reversible, so they have equal Kolmogorov complexity up to \(O(1)\). For a fixed integer seed, all prefixes have \(O(\log j)\) unconditional complexity. Moreover, \(N_j=2^{j+1}-1\) realizes the word \(1^j\), grows strictly throughout the block, and satisfies \(\nu_3(U^j(N_j)+1)=j\), while having \(O(1)\) conditional complexity. Hence positive drift plus finite \(2\)-adic and \(3\)-adic compatibility cannot force high complexity.

\end{quote}


The surviving mathematical target is an \textbf{Archimedean anti-shadowing theorem} proving that the least positive residue of every long non-descending word must grow without bound. That is a genuine value-sensitive arithmetic statement; ordinary, time-bounded, online, and randomness-based Kolmogorov complexity do not establish it.
\subsection*{References}
\begin{enumerate}[leftmargin=2.2em]
\item \hypertarget{ref:R23:1}{}\hyperlink{cit:R23:1}{\texttt{turn654930view0}}
\item \hypertarget{ref:R23:2}{}\hyperlink{cit:R23:2}{\texttt{turn654930view1}}
\item \hypertarget{ref:R23:3}{}\hyperlink{cit:R23:3}{\texttt{turn654930view2}}
\item \hypertarget{ref:R23:4}{}\hyperlink{cit:R23:4}{\texttt{turn988994view0}}
\item \hypertarget{ref:R23:5}{}\hyperlink{cit:R23:5}{\texttt{turn988994view1}}
\end{enumerate}

\clearpage
\section{Round 24: approach 19 --- Terminal-suffix transducer with a vanishing joint spectral radius}
\subsection*{Prompt}
\begin{quote}
Complete a full proof of the Collatz conjecture using the following approach:

\begin{enumerate}
\item Terminal-suffix transducer with a vanishing joint spectral radius
\end{enumerate}

\textbf{Hybrid \textperiodcentered{} B}

This is the most concrete version of the multi-place idea. For each r, build a finite or finitely weighted automaton whose state records:

(a\_\{j-r\}, ..., a\_\{j-1\}), C\_j + 2\textasciicircum{}\{A\_j\} mod 3\textasciicircum{}r, and a discretized archimedean barrier state.

Your terminal-suffix locality theorem makes the 3\textasciicircum{}r-condition genuinely local. Prove that the joint spectral radius of the nontrivial survivor transitions is < 1, uniformly or increasingly strongly with r. A persistent recurrent component would be a precise disproof-search target rather than a vague long trajectory.

Take as long as you need.
\end{quote}
\subsection*{Response}

\subsection*{Outcome}


The proposed proof does \textbf{not} close as stated. The obstruction is specific and rigorous:

\begin{enumerate}
\item The coordinate \(C_j+2^{A_j}\pmod{3^r}\) is not itself terminal-suffix-local; only its divisibility by \(3^r\), or a normalized version of the residue, is local.
\item For Boolean transition matrices, joint spectral radius \(<1\) is equivalent to exact eventual nilpotence. An explicit retained periodic ray produces arbitrarily long survivor paths, so the natural stationary transducer has joint spectral radius at least \(1\).
\item Assigning contracting real weights does not repair the implication: an infinite path can have exponentially decaying weight.
\item Even after replacing the coarse barrier state by an exact finite-height window, fixed-\(r\) spectral radius only detects cycles. A divergent orbit could give a finite but increasingly long transient at every scale.
\item The correct target is a \textbf{graded linear nilpotence estimate}, or equivalently a uniformly well-conditioned contracting norm. I prove below that this strengthened estimate would imply Collatz.
\end{enumerate}

Thus a claimed full proof using only the advertised fixed-state JSR inequality would contain a false or insufficient lemma. What can be completed is an exact obstruction theorem and a corrected spectral reduction.

Current rigorous benchmarks remain pointwise-incomplete: Tao proves an almost-all result using \(3\)-adic transport, while Barina's exhaustive verification reaches \(2^{71}\). [\hypertarget{cit:R24:1}{}\hyperlink{ref:R24:1}{1}, \hypertarget{cit:R24:2}{}\hyperlink{ref:R24:2}{2}]


\bigskip



\subsubsection*{1. The exact terminal coordinate}


Work with the accelerated odd map

\[
U(x)=\frac{3x+1}{2^{v_2(3x+1)}}.
\]

For a valuation word

\[
w=(a_0,\ldots ,a_{j-1}),
\qquad
A_t=\sum_{i<t}a_i,
\]

define

\[
C_0=0,
\qquad
C_{t+1}=3C_t+2^{A_t}.
\]

Then an orbit segment with this word satisfies

\[
2^{A_j}U^j(n)=3^j n+C_j.
\tag{1}
\]

If \(w=uv\), where the terminal suffix \(v\) has length \(r\), and if

\[
P=A(u),\qquad B=A(v),
\]

then

\[
C(uv)=3^rC(u)+2^P C(v).
\tag{2}
\]

This is the source of terminal locality. If \(X=U^j(n)\), then reduction of (1) modulo \(3^r\) gives

\[
X\equiv 2^{-B}C(v)\pmod{3^r},
\tag{3}
\]

and therefore

\[
3^r\mid X+1
\quad\Longleftrightarrow\quad
3^r\mid C(v)+2^B.
\tag{4}
\]

This is exactly your terminal-suffix locality theorem. [\hypertarget{cit:R24:3}{}\hyperlink{ref:R24:3}{3}]


\paragraph*{The proposed residue needs normalization}


Set

\[
Q(w):=C(w)+2^{A(w)}.
\]

Equation (2) gives

\[
Q(uv)
\equiv
2^P\bigl(C(v)+2^B\bigr)
\pmod{3^r}.
\tag{5}
\]

Thus:

\begin{itemize}
\item the condition \(Q(uv)\equiv0\pmod{3^r}\) is suffix-local, because \(2^P\) is invertible modulo \(3^r\);
\item the full residue \(Q(uv)\pmod{3^r}\) is \textbf{not} suffix-local, because it retains the prefix phase \(2^P\).
\end{itemize}

The exact terminal coordinate should instead be

\[
z_j
:=
2^{-A_j}\bigl(C_j+2^{A_j}\bigr)
\pmod{3^r}.
\tag{6}
\]

For \(j\ge r\), (1) shows that

\[
z_j\equiv U^j(n)+1\pmod{3^r}.
\tag{7}
\]

For \(w=uv\), this becomes

\[
z_j\equiv
2^{-B}\bigl(C(v)+2^B\bigr)
\pmod{3^r},
\tag{8}
\]

which is genuinely terminal-\(r\)-local.

There is also an exact forward transition. If the next valuation is \(a\), then

\[
x_{j+1}+1
=
\frac{3(x_j+1)+2^a-2}{2^a},
\]

so

\[
z_{j+1}
\equiv
2^{-a}\bigl(3z_j+2^a-2\bigr)
\pmod{3^r}.
\tag{9}
\]

This gives the correct finite \(3\)-adic part of the transducer.

The remaining problem is archimedean: the exact exponents \(a_j\) and the height deficit are unbounded. Terminal locality makes compatibility finite-state, but it does not make descent finite-state.


\bigskip



\subsubsection*{2. What joint spectral radius \(<1\) means for an automaton}


Let \(\mathcal M\) be a finite family of nonnegative integer matrices, for example the \(0\)-\(1\) transition matrices corresponding to possible next valuation symbols. Its joint spectral radius is

\[
\widehat\rho(\mathcal M)
=
\limsup_{L\to\infty}
\max_{M_1,\ldots,M_L\in\mathcal M}
\left\|M_L\cdots M_1\right\|^{1/L}.
\tag{10}
\]

This measures maximal asymptotic product growth. [\hypertarget{cit:R24:4}{}\hyperlink{ref:R24:4}{4}]


\paragraph*{Lemma 1 \textemdash{} Integer-matrix dichotomy}


For a finite family of nonnegative integer matrices,

\[
\widehat\rho(\mathcal M)<1
\]

if and only if there is an \(L_0\) such that every product of \(L_0\) matrices from \(\mathcal M\) is zero. Consequently, in this case

\[
\widehat\rho(\mathcal M)=0.
\]


\subparagraph*{Proof}


Use the maximum-entry norm. Every nonzero nonnegative integer matrix has norm at least \(1\).

If there are nonzero products of arbitrarily large lengths \(L\), then

\[
\max_{M_1,\ldots,M_L}
\|M_L\cdots M_1\|\ge1
\]

for arbitrarily large \(L\), giving

\[
\widehat\rho(\mathcal M)\ge1.
\]

Therefore \(\widehat\rho<1\) implies that the lengths of nonzero products are bounded. Hence all products of some common length \(L_0\) are zero. The reverse implication is immediate. \ensuremath{\blacksquare}

So for an ordinary finite automaton, ``JSR \(<1\)'' is not a probabilistic contraction. It means:

\[
\boxed{\text{There are no admissible paths beyond one fixed finite length.}}
\]

A directed cycle immediately gives JSR at least \(1\).


\paragraph*{Weighted paths behave differently}


For real weighted matrices, JSR \(<1\) does not imply path extinction. The one-state automaton

\[
M=[1/2]
\]

has

\[
\widehat\rho(M)=1/2,
\]

but its loop can be followed forever.

Moreover, replacing every transition weight by one half of its previous value halves the JSR without changing a single accepted word. Therefore a weighted spectral radius is arithmetically meaningful only when accompanied by a lower bound saying that every genuine integer survivor has a non-negligible weight.


\bigskip



\subsubsection*{3. An explicit retained recurrent obstruction}


The decisive test word is

\[
v=(1,1,2).
\]

Your compression framework already proves that

\[
v^\infty=(1,1,2)^\infty
\]

is a retained Bellman ray, that every finite prefix is retained, and that its associated \(2\)-adic point is

\[
-\frac{19}{11}.
\tag{11}
\]

It also proves that an infinite valuation ray comes from a nonnegative ordinary integer precisely when its nested least \(2\)-adic representatives eventually stabilize. [\hypertarget{cit:R24:5}{}\hyperlink{ref:R24:5}{5}]

Here is the full arithmetic of this obstruction.


\paragraph*{Proposition 2 \textemdash{} The \(-19/11\) survivor ray}


The block \(v=(1,1,2)\) has

\[
A(v)=4,\qquad C(v)=19,
\]

and its affine map is

\[
F_v(x)=\frac{27x+19}{16}.
\tag{12}
\]

Its fixed point is

\[
\xi=-\frac{19}{11}.
\tag{13}
\]

Indeed,

\[
-\frac{19}{11}
\xrightarrow{\,1\,}
-\frac{23}{11}
\xrightarrow{\,1\,}
-\frac{29}{11}
\xrightarrow{\,2\,}
-\frac{19}{11}.
\tag{14}
\]

Every prefix of \(v^\infty\) satisfies the strict multiplicative barrier

\[
2^{A_t}<3^t.
\tag{15}
\]

For every \(K\equiv3\pmod5\), the positive odd integer

\[
n_K=\frac{2^{4K+1}-19}{11}
\tag{16}
\]

has initial valuation word \(v^K\), and every orbit value through those \(3K\) accelerated steps is smaller than \(3^{3K}\).


\subparagraph*{Proof}


Applying the three branches gives

\[
x
\longmapsto
\frac{3x+1}{2}
\longmapsto
\frac{9x+5}{4}
\longmapsto
\frac{27x+19}{16},
\]

which proves (12). Solving \(F_v(\xi)=\xi\) gives (13), and direct substitution gives (14).

At times \(t=3m,3m+1,3m+2\), the cumulative exponent is respectively

\[
4m,\qquad 4m+1,\qquad 4m+2.
\]

Thus

\[
\frac{2^{A_t}}{3^t}
=
\left(\frac{16}{27}\right)^m
\begin{cases}
1,&t=3m,\\[1mm]
2/3,&t=3m+1,\\[1mm]
4/9,&t=3m+2,
\end{cases}
\]

which is less than \(1\) for every \(t\ge1\).

Now suppose \(K\equiv3\pmod5\). Since \(4K+1\equiv3\pmod{10}\), and \(2\) has order \(10\) modulo \(11\),

\[
2^{4K+1}\equiv2^3\equiv8\equiv19\pmod{11}.
\]

Hence \(n_K\) is a positive odd integer. Moreover,

\[
11n_K+19=2^{4K+1},
\]

so

\[
n_K\equiv-\frac{19}{11}\pmod{2^{4K+1}}.
\tag{17}
\]

A word of total exponent \(4K\) is determined by one residue class modulo \(2^{4K+1}\). Therefore \(n_K\) has first \(3K\) valuations \(v^K\).

Let \(x_m=F_v^m(n_K)\) be the point after \(m\) complete blocks. From the fixed-point identity,

\[
F_v^m(x)-\xi
=
\left(\frac{27}{16}\right)^m(x-\xi).
\]

Using (16),

\[
x_m
=
\frac{2^{4(K-m)+1}27^m-19}{11},
\qquad 0\le m\le K.
\tag{18}
\]

Block endpoints increase, because

\[
F_v(x)-x=\frac{11x+19}{16}>0
\]

for \(x>0\). Within each block, the first two \(a=1\) steps increase the value. Hence the largest value in the entire segment occurs immediately before the final \(a=2\) division.

The last block starts at

\[
x_{K-1}
=
\frac{32\cdot27^{K-1}-19}{11}.
\]

The next two points are

\[
\frac{48\cdot27^{K-1}-23}{11},
\qquad
\frac{72\cdot27^{K-1}-29}{11}.
\]

But

\[
\frac{72\cdot27^{K-1}-29}{11}
<
27^K=3^{3K},
\]

because \(72<11\cdot27=297\). All earlier points are smaller. \ensuremath{\blacksquare}


\paragraph*{Consequences}


This one ray defeats several possible interpretations of the proposed spectral step.

First, every finite prefix is not merely formally compatible: it is retained in the Bellman survivor tree. Thus a sound automaton accepting the retained frontier must admit paths labeled by

\[
(1,1,2)^K
\]

for arbitrarily large \(K\).

Second, with the natural forward-slope weight \(3/2^a\), one block has weight

\[
\frac32\cdot\frac32\cdot\frac34
=
\frac{27}{16}>1.
\]

Therefore the constrained weighted spectral radius satisfies

\[
\widehat\rho_{\mathrm{slope}}
\ge
\left(\frac{27}{16}\right)^{1/3}
>1.
\tag{19}
\]

Using the reciprocal weights \(2^a/3\) would make the block contract, but the infinite path would still be present. It would merely be assigned a weight tending to zero.

Third, Proposition 2 is stronger than a purely \(2\)-adic objection. Along the infinite subsequence \(K\equiv3\pmod5\), the finite positive realization satisfies the exact full-depth no-wrap condition

\[
0<U^t(n_K)<3^{3K},
\qquad 0\le t\le3K.
\tag{20}
\]

So positivity, strict barrier, terminal congruence, and full-depth no-wrap can all hold for arbitrarily long prefixes while the starting integer changes with the depth.

The failure is precisely that

\[
n_K\longrightarrow-\frac{19}{11}
\qquad\text{in }\mathbb Z_2,
\]

rather than stabilizing at one positive ordinary integer.

Thus:

\[
\boxed{
\text{A persistent local component is not yet a Collatz counterexample.}
}
\]

It may simply be a rational \(2\)-adic ghost.


\bigskip



\subsubsection*{4. Why the archimedean state cannot be a harmless finite bin}


Suppose an actual word is decomposed as \(w=u\tau\), and a terminal competitor \(b\) is chosen using the \(3\)-adic suffix data. Your exact local-global formula is

\[
Y-R
=
\frac{\Sigma(u,\tau,c)+h(\tau,b)}{2^c}.
\tag{21}
\]

Here \(h(\tau,b)\) is determined by the terminal suffix and its compatible fibre, but

\[
\Sigma(u,\tau,c)
\]

depends on the entire prefix and on its least positive realization. Therefore terminal compatibility and pointwise descent are mathematically distinct. [\hypertarget{cit:R24:6}{}\hyperlink{ref:R24:6}{6}]

A finite barrier bin can be used only after proving that every exact prefix represented by that bin gives the same sign in (21). Without such a uniform domination theorem, the bin is an overapproximation and may introduce spurious cycles\textemdash{}or unsoundly delete genuine paths.

Your explicit strict-barrier word of length \(42\) has optimal competing-ancestor cutoff exactly \(814\). This proves that even within the present finite data, very small universal cutoff rules fail; it does not prove the cutoffs unbounded. [\hypertarget{cit:R24:7}{}\hyperlink{ref:R24:7}{7}]

This is why the archimedean part cannot simply be appended as a fixed number of coarse states and then forgotten.


\bigskip



\subsubsection*{5. The exact finite transducer that is safe}


The safe repair is to replace a discretized barrier by an exact finite-height box.

For \(r\ge1\), let

\[
\mathcal H_r
=
\{x\in\mathbb Z_{>0}:x\text{ odd},\ 1<x<3^r\}.
\]

An \(r\)-window is a tuple

\[
(x_0,x_1,\ldots,x_r)
\]

such that

\[
x_i\in\mathcal H_r,
\qquad
x_{i+1}=U(x_i).
\]

Let \(\mathcal S_r\) be the set of all such windows. It is finite; in fact a window is determined by \(x_0\), so

\[
|\mathcal S_r|<3^r.
\tag{22}
\]

There is a directed edge

\[
(x_0,\ldots,x_r)
\longrightarrow
(x_1,\ldots,x_r,x_{r+1})
\]

whenever \(x_{r+1}=U(x_r)\) also lies in \(\mathcal H_r\).


\paragraph*{Terminal-suffix encoding of the exact window}


Let

\[
v=(a_0,\ldots,a_{r-1})
\]

be the valuations in an \(r\)-window, and let \(B=A(v)\). Terminal locality gives

\[
x_r\equiv2^{-B}C(v)\pmod{3^r}.
\tag{23}
\]

Because \(1<x_r<3^r\), this congruence identifies \(x_r\) exactly. Reading backward through

\[
x_i=\frac{2^{a_i}x_{i+1}-1}{3}
\]

reconstructs the entire window. Therefore the finite box graph is genuinely equivalent to a terminal-suffix transducer.

The possible exponents are also bounded inside the box. Since

\[
2^{a_i}x_{i+1}=3x_i+1
\]

and \(x_i<3^r\),

\[
2^{a_i}<3^{r+1}+1.
\]

Thus only finitely many exponent symbols occur at each scale \(r\).

This is the exact version of the construction: no coarse real-height bin and no spurious \(2\)-adic path.


\bigskip



\subsubsection*{6. Fixed-scale JSR is still insufficient}


Let \(\mathcal M_r\) be the family of \(0\)-\(1\) transition matrices of the exact scale-\(r\) transducer. Define its post-warm-up residence length

\[
\nu_r
=
\sup\left\{
L:
\begin{array}{c}
\text{there exist }x_0,\ldots,x_{r+L}\in\mathcal H_r,\\
x_{i+1}=U(x_i)
\end{array}
\right\}.
\tag{24}
\]

Equivalently, \(\nu_r\) is the largest length of a nonzero matrix product from \(\mathcal M_r\). It is infinite when the box contains a nontrivial cycle.

By Lemma 1,

\[
\widehat\rho(\mathcal M_r)<1
\quad\Longleftrightarrow\quad
\nu_r<\infty.
\tag{25}
\]

This only says that the orbit cannot remain in the \textbf{fixed} box forever.

A divergent orbit can exit every finite box. It is entirely compatible with

\[
\widehat\rho(\mathcal M_r)=0
\qquad\text{for every fixed }r,
\]

while \(\nu_r\) grows rapidly with \(r\).

The rate of growth of \(\nu_r\), not the fixed-scale JSR, is the decisive quantity.


\bigskip



\subsubsection*{7. The corrected graded criterion}


Set

\[
\kappa
=
\log_{3/2}2
=
\frac{\log 2}{\log(3/2)}
=
1.709511291\ldots.
\tag{26}
\]


\paragraph*{Theorem 3 \textemdash{} Linear residence criterion}


If a positive Collatz orbit avoids \(1\), then

\[
\liminf_{r\to\infty}\frac{\nu_r}{r}
\ge\kappa.
\tag{27}
\]

Consequently, any proof that

\[
\liminf_{r\to\infty}\frac{\nu_r}{r}<\kappa
\tag{28}
\]

would prove the Collatz conjecture. In particular, the stronger estimate

\[
\nu_r\le(\kappa-\varepsilon)r
\tag{29}
\]

for all sufficiently large \(r\) would be a complete proof.


\subparagraph*{Proof}


For every accelerated step,

\[
x_{t+1}+1
=
\frac{3x_t+1}{2^{a_t}}+1
\le
\frac{3x_t+1}{2}+1
=
\frac32(x_t+1).
\]

Therefore

\[
x_t+1
\le
\left(\frac32\right)^t(n+1).
\tag{30}
\]

Fix \(c<\kappa\). Then

\[
(1+c)\log(3/2)<\log3.
\]

Hence, for all sufficiently large \(r\),

\[
\left(\frac32\right)^{r+\lfloor cr\rfloor}(n+1)
<
3^r.
\tag{31}
\]

It follows from (30) that

\[
1<x_t<3^r
\qquad
0\le t\le r+\lfloor cr\rfloor.
\]

Thus the first \(r\) steps form a state of the terminal transducer, followed by at least \(\lfloor cr\rfloor\) legal transitions. Hence

\[
\nu_r\ge\lfloor cr\rfloor
\]

for every sufficiently large \(r\). Letting \(c\uparrow\kappa\) proves (27). \ensuremath{\blacksquare}

This is the correct replacement for ``JSR \(<1\), increasingly strongly with \(r\)'':

\[
\boxed{
\text{The normalized nilpotence length must be strictly below }
\log_{3/2}2.
}
\]

A fixed-\(r\) JSR does not retain this nilpotence-length information: every finite nilpotent automaton has JSR exactly zero, whether its longest path has length \(10\), \(10^{10}\), or \(3^r\).


\bigskip



\subsubsection*{8. A genuine weighted-JSR theorem}


There is a way to formulate the needed estimate spectrally, but it requires control of the norm's conditioning across scales.

Let \(\|\cdot\|_r\) be a norm on the scale-\(r\) state space, and define its distortion relative to the standard maximum norm by

\[
D_r
=
\left(
\sup_{v\ne0}\frac{\|v\|_r}{\|v\|_\infty}
\right)
\left(
\sup_{v\ne0}\frac{\|v\|_\infty}{\|v\|_r}
\right).
\tag{32}
\]

Suppose there is a fixed \(\theta<1\) such that every survivor-transition matrix satisfies

\[
\|M\|_r\le\theta.
\tag{33}
\]

Then every product \(P=M_L\cdots M_1\) satisfies

\[
\|P\|_\infty
\le
D_r\theta^L.
\tag{34}
\]

If \(P\) is a nonzero \(0\)-\(1\) product, then

\[
\|P\|_\infty\ge1.
\]

Therefore

\[
D_r\theta^L\ge1
\]

for every admissible path of length \(L\), and hence

\[
\nu_r
\le
\frac{\log D_r}{\log(1/\theta)}.
\tag{35}
\]

Combining this with Theorem 3 gives the following exact spectral win condition.


\paragraph*{Corollary 4 \textemdash{} Balanced common-norm criterion}


The Collatz conjecture follows if there exist \(\theta<1\) and norms \(\|\cdot\|_r\) satisfying (33) and

\[
\liminf_{r\to\infty}
\frac{\log D_r}
     {r\log(1/\theta)}
<
\log_{3/2}2.
\tag{36}
\]

For example, subexponential distortion

\[
D_r=e^{o(r)}
\]

together with any uniform contraction \(\theta<1\) would prove Collatz.

If the available norm has distortion of approximately the state-space size,

\[
D_r\le e^{o(r)}3^r,
\]

then (36) requires

\[
\theta
<
3^{-1/\kappa}
=
0.5258985309\ldots.
\tag{37}
\]

This is the mathematically meaningful ``vanishing weighted JSR'' target. It contains both ingredients absent from a bare JSR statement:

\[
\text{exponential contraction}
\quad+\quad
\text{controlled norm distortion}.
\]

For a nilpotent graph one can always invent a very strongly contracting norm by assigning exponentially different weights to successive levels. But then \(D_r\) is exponentially large in the longest path length, and inequality (36) records exactly that hidden cost.


\bigskip



\subsubsection*{9. What would constitute a real disproof target}


Let \(R_j\) be the least nonnegative representative of the starting-value class of an infinite valuation ray modulo \(2^{A_j+1}\). These representatives are nested:

\[
R_{j+1}\equiv R_j\pmod{2^{A_j+1}}.
\]

The ray comes from one positive ordinary integer \(n\) exactly when

\[
R_j=n
\]

for all sufficiently large \(j\). [\hypertarget{cit:R24:8}{}\hyperlink{ref:R24:8}{8}]

Therefore a persistent collection of recurrent transducer states is a genuine counterexample candidate only when it comes with a projectively compatible tower satisfying all of the following simultaneously:

\[
\begin{aligned}
&\text{the suffix and \(3\)-adic states project consistently as }r\to\infty,\\
&\text{the nested \(2\)-adic seed representatives eventually stabilize at one }n>0,\\
&\text{the resulting exact orbit avoids }1,\\
&\text{the archimedean comparisons never produce a smaller coalescing ancestor.}
\end{aligned}
\]

The ray \((1,1,2)^\infty\) fails the second condition: its representatives converge to \(-19/11\), rather than stabilizing.

Your existing framework identifies the same missing theorem as a synchronization between eventual positive \(2\)-adic stabilization and \(3\)-adic endpoint compressibility. It also notes that proving the finite \(81/64\) subsystem would still leave this global synchronization step unresolved. [\hypertarget{cit:R24:9}{}\hyperlink{ref:R24:9}{9}]


\bigskip



\subsection*{Final verdict}


Approach 19, in its literal stationary form, does not prove Collatz.

The proposed state must first be corrected by normalizing the carry residue. After that correction:

\[
\begin{array}{c}
\text{Boolean JSR}<1
\end{array}
\]

means only fixed-scale nilpotence, which cannot exclude escape to infinity; and the retained ray

\[
(1,1,2)^\infty\longleftrightarrow-\frac{19}{11}
\]

prevents a coarse stationary survivor automaton from having such a spectral gap in the first place.

The strongest correct theorem obtained from the approach is:

\[
\boxed{
\text{Collatz follows if }
\liminf_{r\to\infty}\frac{\nu_r}{r}
<
\log_{3/2}2.
}
\]

A genuine JSR implementation would need to prove the balanced norm estimate

\[
\boxed{
\frac{\log D_r}{r\log(1/\theta)}
<
\log_{3/2}2
}
\]

along an unbounded sequence of scales.

That estimate is not presently established. Proving it requires exactly the missing local-global synchronization: terminal \(3\)-adic locality must be coupled to the prefix-dependent archimedean score for one fixed, eventually stabilized positive seed.
\subsection*{References}
\begin{enumerate}[leftmargin=2.2em]
\item \hypertarget{ref:R24:1}{}\hyperlink{cit:R24:1}{\texttt{turn344682view0}}
\item \hypertarget{ref:R24:2}{}\hyperlink{cit:R24:2}{\texttt{turn344682view1}}
\item \hypertarget{ref:R24:3}{}\hyperlink{cit:R24:3}{\texttt{turn18file0 L75-L107}}
\item \hypertarget{ref:R24:4}{}\hyperlink{cit:R24:4}{\texttt{turn344682view2}}
\item \hypertarget{ref:R24:5}{}\hyperlink{cit:R24:5}{\texttt{turn20file0 L12-L56}}
\item \hypertarget{ref:R24:6}{}\hyperlink{cit:R24:6}{\texttt{turn18file0 L162-L250}}
\item \hypertarget{ref:R24:7}{}\hyperlink{cit:R24:7}{\texttt{turn19file0 L75-L158}}
\item \hypertarget{ref:R24:8}{}\hyperlink{cit:R24:8}{\texttt{turn20file0 L12-L40}}
\item \hypertarget{ref:R24:9}{}\hyperlink{cit:R24:9}{\texttt{turn20file3 L306-L321}}
\end{enumerate}

\clearpage
\section{Round 25: approach 20 --- Ultraproduct rigidity of minimal counterexamples}
\subsection*{Prompt}
\begin{quote}
Complete a full proof of the Collatz conjecture using the following approach:

\begin{enumerate}
\item Ultraproduct rigidity of minimal counterexamples
\end{enumerate}

\textbf{Speculative \textperiodcentered{} B}

Assume there are arbitrarily long counterexample-like trajectories and pass to a nonstandard model. Saturation produces an orbit surviving for a nonstandard number of steps, potentially with a well-defined asymptotic valuation frequency and normalized affine carry.

Show that such an object acquires a scaling or translation symmetry. If that symmetry forces periodicity, transfer it back to a forbidden finite cycle; if it instead yields a standard initial segment with a uniform escape invariant, it could lead toward a divergent counterexample.

Take as long as you need.
\end{quote}
\subsection*{Response}

\subsection*{Verdict}


The ultraproduct program can be carried through rigorously, but it does \textbf{not} complete a proof of Collatz. Its decisive proposed implication,

\[
\text{nonstandard survival + asymptotic frequency + normalized carry}
\quad\Longrightarrow\quad
\text{periodicity},
\]

is false. In fact, one can explicitly construct hyperfinite Collatz trajectories having all of those regularity properties whose limiting dynamics is an \textbf{aperiodic irrational rotation}.

This is not just the usual objection that Collatz remains open. It is a mathematical obstruction to the proposed proof. Current literature still describes a proof as out of reach, while the published computational verification reaches \(2^{71}\). [\hypertarget{cit:R25:1}{}\hyperlink{ref:R25:1}{1}, \hypertarget{cit:R25:2}{}\hyperlink{ref:R25:2}{2}]

What follows is a complete proof of that obstruction, followed by the exact additional theorem an ultraproduct approach would need.


\bigskip



\subsection*{1. Exact affine dynamics}


Use the shortcut Collatz map

\[
T(n)=
\begin{cases}
n/2,&n\equiv0\pmod 2,\\[2mm]
(3n+1)/2,&n\equiv1\pmod 2.
\end{cases}
\]

Let

\[
e_i=T^i(n)\pmod 2,\qquad
s_k=\sum_{i=0}^{k-1}e_i.
\]

Since

\[
2T(x)=3^{e(x)}x+e(x),
\]

iteration gives

\[
2^kT^k(n)=3^{s_k}n+B_k,
\tag{1}
\]

where

\[
B_k=\sum_{i=0}^{k-1}e_i2^i3^{\,s_k-s_{i+1}}\geq0.
\tag{2}
\]

Consequently,

\[
\frac{T^k(n)}n
=
\frac{3^{s_k}}{2^k}
\left(1+\frac{B_k}{3^{s_k}n}\right).
\tag{3}
\]

The first factor is the homogeneous multiplicative drift; the second factor is the normalized affine carry.

A standard and crucial fact is that every binary word of length \(k\) is the parity vector of exactly one residue class modulo \(2^k\). Thus every finite parity pattern is realized by infinitely many positive integers. Passing to infinite words gives a bijection between all binary sequences and \(2\)-adic seeds. [\hypertarget{cit:R25:3}{}\hyperlink{ref:R25:3}{3}, \hypertarget{cit:R25:4}{}\hyperlink{ref:R25:4}{4}]

We also use the accelerated odd map

\[
U(n)=\frac{3n+1}{2^{v_2(3n+1)}}.
\]

For an odd orbit \(x_{j+1}=U(x_j)\), put

\[
a_j=v_2(3x_j+1),\qquad
A_j=\sum_{i=0}^{j-1}a_i.
\]

Then

\[
2^{A_j}U^j(n)=3^jn+C_j,
\tag{4}
\]

with

\[
C_0=0,\qquad C_{j+1}=3C_j+2^{A_j},
\tag{5}
\]

and hence

\[
C_j=\sum_{i=0}^{j-1}3^{j-1-i}2^{A_i}.
\tag{6}
\]

Thus

\[
\frac{U^j(n)}n
=
\frac{3^j}{2^{A_j}}
\left(1+\frac{C_j}{3^jn}\right).
\tag{7}
\]


\bigskip



\subsection*{2. Finite survival is unconditional}


Call \(n\) an \(H\)-survivor when

\[
T^j(n)\geq n\qquad(1\leq j\leq H).
\]

Such survivors exist for every \(H\), independently of whether Collatz is true.

Indeed, take

\[
n_H=2^H-1.
\]

For \(0\leq j\leq H\),

\[
T^j(n_H)=3^j2^{H-j}-1.
\tag{8}
\]

This follows by induction because all the inputs used before the final step are odd. For \(1\leq j\leq H\),

\[
T^j(n_H)-n_H
=
2^{H-j}(3^j-2^j)>0.
\]

Therefore

\[
T^j(n_H)>n_H\qquad(1\leq j\leq H).
\tag{9}
\]

Now choose a nonprincipal ultrafilter and form the ultrapower. Set

\[
N=[\,2^H-1\,],\qquad \mathcal H=[\,H\,].
\]

Then \(N\) and \(\mathcal H\) are infinite hyperintegers, and \L{}o\'s's theorem gives

\[
T^j(N)\geq N
\qquad
\text{for every internal }1\leq j\leq\mathcal H.
\tag{10}
\]

Compactness and \L{}o\'s's theorem do exactly this: finite satisfiability produces a nonstandard realization, and first-order statements transfer coordinatewise. [\hypertarget{cit:R25:5}{}\hyperlink{ref:R25:5}{5}, \hypertarget{cit:R25:6}{}\hyperlink{ref:R25:6}{6}]

But the sequence \(2^H-1\) converges \(2\)-adically to

\[
-1=\cdots111111_2.
\]

The all-ones parity sequence corresponds to the \(2\)-adic fixed point \(-1\), not to a positive-integer counterexample. [\hyperlink{ref:R25:4}{4}]

So even the very strong-looking assertion

\[
\exists N,\mathcal H\text{ infinite}:
\quad T^j(N)\geq N\text{ for all }j\leq\mathcal H
\]

is true unconditionally. It cannot distinguish the truth or falsity of Collatz.

This already shows that compactness is solving the wrong existence problem: it produces a nonstandard or \(2\)-adic seed, whereas Collatz concerns one fixed ordinary positive integer.


\bigskip



\subsection*{3. A critical aperiodic ultraproduct}


The preceding example has excessive upward drift. Perhaps one could hope that imposing the exact critical valuation frequency and controlling the carry creates rigidity.

It does not.


\subsubsection*{Theorem 1: Sturmian obstruction}


Let

\[
\beta=\log_2 3.
\]

For every \(H\geq2\), there is a positive odd integer \(n_H\) whose first \(H\) accelerated valuations \(a_0,\ldots,a_{H-1}\) satisfy all of the following:

\[
\left|
\sum_{i=r}^{s-1}a_i-\beta(s-r)
\right|<1
\qquad(0\leq r<s\leq H),
\tag{11}
\]

\[
n_H<U^j(n_H)
<
2\left(1+\frac1{3H}\right)n_H
\qquad(1\leq j\leq H),
\tag{12}
\]

and

\[
0\leq
\frac{C_j}{3^jn_H}
\leq\frac1{3H}
\qquad(0\leq j\leq H).
\tag{13}
\]

Moreover, the infinite valuation word from which these prefixes are taken is aperiodic and is a coding of an irrational rotation.


\paragraph*{Proof}


Define

\[
A_j=\lfloor\beta j\rfloor,\qquad
a_j=A_{j+1}-A_j.
\tag{14}
\]

Because \(1<\beta<2\),

\[
a_j\in\{1,2\}.
\]

For any \(0\leq r<s\),

\[
\sum_{i=r}^{s-1}a_i=A_s-A_r.
\]

Therefore

\[
\begin{aligned}
(A_s-A_r)-\beta(s-r)
&=
(\lfloor\beta s\rfloor-\beta s)
-
(\lfloor\beta r\rfloor-\beta r),
\end{aligned}
\]

which lies strictly between \(-1\) and \(1\). This proves (11).

To realize the valuation prefix \(a_0,\ldots,a_{H-1}\), prescribe the shortcut parity word

\[
1\,0^{a_0-1}\,
1\,0^{a_1-1}\cdots
1\,0^{a_{H-1}-1}\,1.
\tag{15}
\]

It has length \(A_H+1\). By the parity-cylinder bijection, it is realized by one residue class modulo \(2^{A_H+1}\), hence by infinitely many positive odd integers. Choose a representative \(n_H\) large enough that

\[
n_H\geq H^2.
\tag{16}
\]

For the resulting accelerated orbit, equation (4) holds. Since

\[
A_i=\lfloor\beta i\rfloor\leq\beta i,
\]

we have

\[
2^{A_i}\leq2^{\beta i}=3^i.
\]

Using (6),

\[
\frac{C_j}{3^j}
=
\sum_{i=0}^{j-1}\frac{2^{A_i}}{3^{i+1}}
\leq
\sum_{i=0}^{j-1}\frac13
=
\frac j3.
\tag{17}
\]

Hence, for \(j\leq H\),

\[
0\leq\frac{C_j}{3^jn_H}
\leq
\frac{j}{3n_H}
\leq
\frac{1}{3H}.
\]

This proves (13).

The homogeneous multiplier is

\[
q_j=\frac{3^j}{2^{A_j}}
=
2^{\beta j-\lfloor\beta j\rfloor}
=
2^{\{\beta j\}}.
\tag{18}
\]

The number \(\beta\) is irrational: otherwise \(2^p=3^q\) for some positive integers \(p,q\), contradicting unique factorization. Thus, for every \(j>0\),

\[
1<q_j<2.
\tag{19}
\]

Substituting (13) and (19) into (7) gives

\[
1<
\frac{U^j(n_H)}{n_H}
<
2\left(1+\frac1{3H}\right),
\]

which is (12).

Finally, put

\[
\alpha=\beta-1\in(0,1).
\]

Since

\[
\lfloor\beta j\rfloor
=
j+\lfloor\alpha j\rfloor,
\]

we obtain

\[
a_j
=
1+\lfloor\alpha(j+1)\rfloor-\lfloor\alpha j\rfloor.
\tag{20}
\]

This is a mechanical, or Sturmian, coding of the irrational rotation

\[
\theta\longmapsto\theta+\alpha\pmod1.
\]

It cannot be eventually periodic: an eventually periodic integer sequence has rational limiting average, whereas

\[
\frac{A_H}{H}\longrightarrow\beta
\]

is irrational.

Moreover,

\[
\log_2\frac{U^j(n_H)}{n_H}
=
\{\beta j\}
+
\log_2\left(1+\frac{C_j}{3^jn_H}\right),
\]

so uniformly for \(j\leq H\),

\[
\log_2\frac{U^j(n_H)}{n_H}
=
\{\beta j\}+O(H^{-1}).
\tag{21}
\]

Thus the normalized height shadows an irrational rotation. \ensuremath{\blacksquare}


\bigskip



\subsubsection*{Ultraproduct consequence}


Take the ultrapower of these \(n_H\). Then there are infinite hyperintegers \(N,\mathcal H\) such that, internally for every \(1\leq j\leq\mathcal H\),

\[
N<U^j(N)<(2+\varepsilon)N,
\tag{22}
\]

where \(\varepsilon\) is infinitesimal,

\[
\frac{C_j}{3^jN}\simeq0,
\tag{23}
\]

and every valuation interval has discrepancy less than \(1\) from the critical mean \(\beta\):

\[
\left|
A_s-A_r-\beta(s-r)
\right|<1.
\tag{24}
\]

Nevertheless, no standard integer \(p>0\) is a period of the valuation word on its hyperfinite domain. Its limiting normalized-height system is an irrational rotation.

This directly disproves the proposed rigidity principle. The object has:

\begin{itemize}
\item nonstandard survival;
\item an exact asymptotic valuation frequency;
\item zero normalized carry in standard part;
\item bounded normalized height;
\item a strong translation symmetry;
\end{itemize}

but it is aperiodic.

A recent 2026 preprint on Christoffel extremizers also singles out Sturmian words\textemdash{}the irrational-slope, aperiodic analogues of Christoffel words\textemdash{}as natural objects to investigate in connection with unbounded Collatz behavior. [\hypertarget{cit:R25:7}{}\hyperlink{ref:R25:7}{7}]


\bigskip



\subsection*{4. What happens when one imposes minimality?}


Define the least \(H\)-survivor

\[
m_H=
\min\left\{
n>1:
T^j(n)\geq n
\text{ for }1\leq j\leq H
\right\}.
\tag{25}
\]

The set is nonempty by the example \(2^H-1\), and \(m_H\) is nondecreasing.


\subsubsection*{Proposition 2}


The Collatz conjecture is true if and only if

\[
m_H\longrightarrow\infty.
\tag{26}
\]

Equivalently, Collatz is false if and only if \((m_H)\) is bounded.


\paragraph*{Proof}


Suppose Collatz is false, and let \(n_0\) be the least positive counterexample. It must be odd, since an even counterexample would map immediately to the smaller counterexample \(n_0/2\).

No iterate of \(n_0\) can lie below \(n_0\). Indeed, if

\[
T^j(n_0)<n_0,
\]

then by minimality \(T^j(n_0)\) reaches \(1\), so \(n_0\) also reaches \(1\), a contradiction. Therefore

\[
m_H\leq n_0
\qquad\text{for every }H.
\]

Thus \((m_H)\) is bounded.

Conversely, if \((m_H)\) is bounded, then because it is nondecreasing and integer-valued, it eventually stabilizes at some \(m\). For every \(j\), choose \(H\) beyond both the stabilization point and \(j\). Then

\[
T^j(m)\geq m.
\]

Thus the orbit of \(m>1\) never reaches \(1\), and Collatz is false. \ensuremath{\blacksquare}

Now form

\[
M=[m_H],\qquad\mathcal H=[H].
\]

Internally, \(M\) is the least \(\mathcal H\)-survivor:

\[
M=
\min\left\{
X>1:
T^j(X)\geq X
\text{ for }1\leq j\leq\mathcal H
\right\}.
\tag{27}
\]

But Proposition 2 says

\[
M\text{ is standard}
\quad\Longleftrightarrow\quad
\text{Collatz is false},
\tag{28}
\]

while

\[
M\text{ is infinite}
\quad\Longleftrightarrow\quad
\text{Collatz is true}.
\tag{29}
\]

Thus the question ``is the internally minimal hyper-survivor standard?'' is exactly the original Collatz conjecture in another form. Saturation has not simplified it.

In particular, if Collatz is true, an infinite internally minimal hyper-survivor necessarily exists. Its existence cannot be used as a contradiction.


\bigskip



\subsection*{5. Minimality disappears at the standard-time boundary}


There is an even sharper obstruction.

Let

\[
\mathcal S_{\mathrm{std}}
=
\left\{
X\in{}^\ast\mathbb N:
X\text{ is infinite and }
T^j(X)\geq X
\text{ for every standard }j
\right\}.
\tag{30}
\]


\subsubsection*{Proposition 3}


The set \(\mathcal S_{\mathrm{std}}\) has no least element. More precisely, every \(X\in\mathcal S_{\mathrm{std}}\) has an infinite \(Y<X\) with \(Y\in\mathcal S_{\mathrm{std}}\).


\paragraph*{Proof}


Fix \(X\in\mathcal S_{\mathrm{std}}\). Since \(T(X)\geq X\), \(X\) is odd.

Fix a standard \(j\). Its first \(j\) parity bits form a standard finite word, so by (1)

\[
T^j(X)=q_jX+r_j,
\tag{31}
\]

where

\[
q_j=\frac{3^{s_j}}{2^j}
\]

is a standard positive rational and \(r_j=B_j/2^j\) is finite.

Because \(X\) is infinite and \(T^j(X)\geq X\), we cannot have \(q_j<1\): in that case

\[
(q_j-1)X+r_j<0.
\]

We also cannot have \(q_j=1\), since

\[
3^{s_j}=2^j
\]

has no solution for \(j>0\). Therefore

\[
q_j>1
\qquad\text{for every standard }j>0.
\tag{32}
\]

Choose an infinite hyperinteger \(K\) such that

\[
2^K<X
\qquad\text{and}\qquad
\frac{2^K}{X}\simeq0.
\]

Set

\[
Y=X-2^K.
\]

Then \(Y\) is positive, odd, infinite, and \(Y<X\). Moreover,

\[
Y\equiv X\pmod{2^K}.
\]

Since \(K\) exceeds every standard integer, \(X\) and \(Y\) have the same first \(j\) parity bits for every standard \(j\). Consequently,

\[
T^j(Y)=q_jY+r_j.
\]

By (32),

\[
T^j(Y)-Y=(q_j-1)Y+r_j>0.
\]

This holds for every standard \(j\), so \(Y\in\mathcal S_{\mathrm{std}}\). \ensuremath{\blacksquare}

Therefore an \(H\)-minimal hyper-survivor can remain minimal only because every smaller candidate eventually fails at some \textbf{nonstandard} time. As soon as one discards that nonstandard tail and keeps only all standard initial segments, minimality vanishes.

This is the precise model-theoretic failure:

\[
\text{least survivor through one nonstandard horizon}
\]

does not become

\[
\text{least survivor through all standard times}.
\]


\bigskip



\subsection*{6. The minimality\textendash{}scaling incompatibility}


Suppose Collatz is false and \(n_0\) is the least counterexample.

Keeping \(n_0\) preserves genuine minimality, but \(n_0\) is a fixed finite integer. The affine \(+1\) term is not infinitesimal, so there is no homogeneous scaling limit.

To obtain a scaling limit, one must move far out along a divergent orbit and normalize by very large orbit values. But those large orbit values are not minimal counterexamples: \(n_0\) is already a smaller counterexample on the same orbit.

Alternatively, one can use the horizon minima \(m_H\), which are internally minimal and may tend to infinity. But their infinitude is equivalent to Collatz being true, and their standardness is equivalent to Collatz being false.

Thus the proposed argument cannot simultaneously retain all three of:

\[
\text{genuine minimality},\qquad
\text{infinite scaling},\qquad
\text{one fixed standard seed}.
\]

One of them is always lost.


\bigskip



\subsection*{7. Why the symmetry does not transfer to a positive cycle}


Let \(w\) be a standard parity block of length \(p\), with \(s\) odd steps and affine carry \(B_w>0\). Then

\[
T^p(x)=\frac{3^sx+B_w}{2^p}.
\tag{33}
\]

A fixed point would satisfy

\[
(2^p-3^s)x=B_w,
\]

so

\[
x=\frac{B_w}{2^p-3^s}.
\tag{34}
\]

A positive fixed point therefore requires

\[
2^p>3^s.
\tag{35}
\]

Equivalently, if

\[
\rho=\log_3 2,
\]

then a positive cycle block must have

\[
\frac{s}{p}<\rho.
\tag{36}
\]

But Proposition 3 shows that every standard prefix of an infinite barrier shadow satisfies

\[
\frac{3^s}{2^p}>1,
\]

hence

\[
\frac{s}{p}>\rho.
\tag{37}
\]

Repeating such a standard prefix periodically gives

\[
2^p-3^s<0,
\]

so the formal fixed point in (34) is negative.

The simplest example is the all-odd word, which gives the \(2\)-adic and integer fixed point

\[
x=-1.
\]

Thus the periodic object naturally produced by the upward nonstandard shadow has the \textbf{wrong sign} for a positive Collatz cycle.

There is a second obstruction. Suppose for a fixed standard \(p\),

\[
T^p(X)=X+d
\]

for an infinite \(X\) and a finite translation \(d\). Equation (33) gives

\[
(2^p-3^s)X=B_w-2^pd.
\tag{38}
\]

The right side is finite, while \(2^p-3^s\neq0\) is a nonzero standard integer. Hence \(X\) would be finite, a contradiction.

Therefore an infinite Collatz shadow admits neither an exact standard period nor an exact bounded standard translation. Any recurrence must be either:

\begin{enumerate}
\item approximate rather than exact, or
\item located at a nonstandard period \(P\).
\end{enumerate}

Approximate equality in a standard-part quotient is not transferred by \L{}o\'s to exact integer equality. A nonstandard period corresponds to periods tending to infinity and does not produce a fixed standard finite cycle.

Also, periodicity by itself would not finish Collatz: nontrivial positive cycles are another unresolved possibility, not something already excluded.


\bigskip



\subsection*{8. The finite-state quantifier error}


Suppose one records a finite state at resolution \(r\): seed modulo \(2^r\), carry modulo \(3^r\), a valuation suffix, and a discretized real height.

For each fixed \(r\), the state space is finite. A hyperfinite trajectory therefore has repeated \(r\)-states. This gives a statement of the form

\[
\forall r\ \exists\,t_r,p_r:
\qquad
S_r(t_r+p_r)=S_r(t_r).
\tag{39}
\]

Exact periodicity would require

\[
\exists\,t,p\ \forall r:
\qquad
S_r(t+p)=S_r(t).
\tag{40}
\]

The order of the quantifiers has been reversed.

Nothing in compactness, saturation, or finite-state recurrence allows this reversal. Sturmian sequences demonstrate the failure perfectly: every finite factor recurs, and every finite-resolution system looks recurrent, but there is no global period.

Your terminal-suffix locality theorem helps make the \(3^r\)-endpoint condition genuinely local at each fixed \(r\). It does not provide a recurrence gap \(p\) that is uniform in \(r\). As \(r\) grows sufficiently large to recover exact arithmetic, the state space and possible recurrence time also grow.

That uniformity is the missing theorem.


\bigskip



\subsection*{9. What a successful ultraproduct theorem would actually need}


Assume a divergent counterexample exists. Take its least element \(n_0\). If its orbit is not periodic, it has no repeated values. Since only finitely many integers lie below any fixed bound, an unbounded nonperiodic orbit must in fact tend to infinity.

One may then choose suffix-minimum times \(t_i\) such that

\[
x_{t_i}\leq x_k\qquad(k\geq t_i)
\]

and

\[
x_{t_i}\longrightarrow\infty.
\]

The ultraproduct

\[
X=[x_{t_i}]
\]

satisfies

\[
T^j(X)\geq X
\]

for every hypernatural \(j\), not merely for \(j\leq\mathcal H\). This is the strongest legitimate object supplied by a genuine divergent counterexample.

A usable rigidity theorem would have to exploit the fact that all the \(x_{t_i}\) come from \textbf{one fixed standard orbit}. Local frequency, normalized height and carry do not remember that ancestry.

One sufficient form would be a uniform standardization theorem asserting that every such orbit-derived hyperobject produces at least one of the following standard outputs:

\begin{itemize}
\item a fixed standard \(p\) and an exact bounded affine recurrence;
\item a fixed standard residue cylinder containing a positive-density family of genuine non-descending orbits;
\item a standard-time descent certificate.
\end{itemize}

The first is impossible by equation (38). The third contradicts the suffix-minimum construction. The second would contradict Tao's theorem: taking, for example, \(f(N)=\sqrt N\), Tao proves that for logarithmic-density-one many integers, the orbit eventually falls below \(f(N)\). A fixed arithmetic progression of genuine no-descent counterexamples would have positive logarithmic density. [\hypertarget{cit:R25:8}{}\hyperlink{ref:R25:8}{8}]

But shrinking cylinders do not suffice. For example,

\[
n\equiv-1\pmod{2^r}
\qquad(r=1,2,\ldots)
\]

form a compatible nested family. Their \(2\)-adic intersection is \(-1\), while their intersection in the positive ordinary integers is empty.

Consequently, a disproof route would have to prove that the inverse-limit seed is an ordinary positive integer\textemdash{}equivalently, that its binary expansion eventually terminates\textemdash{}not merely that it is a \(2\)-adic or nonstandard point.


\bigskip



\subsection*{10. Final assessment}


The proposed ultraproduct construction succeeds in producing nonstandard long-lived objects. What fails is the rigidity step.

The precise conclusions are:

\[
\boxed{
\text{Arbitrarily long no-descent trajectories exist unconditionally.}
}
\]

\[
\boxed{
\begin{aligned}
&\text{Critical valuation frequency, negligible normalized carry,}\\
&\text{bounded normalized height, and translation symmetry}\\
&\text{do not imply periodicity.}
\end{aligned}
}
\]

\[
\boxed{
\text{Finite-horizon minimality becomes useful only if one can prove standardness,}
}
\]

and that standardness question is equivalent to Collatz itself.

The strongest valid outcome of approach 20 is therefore a no-go theorem: the natural ultraproduct quotient erases the affine \(+1\) term, replaces a positive integer by a general \(2\)-adic/nonstandard seed, and admits critical aperiodic Sturmian dynamics. A successful repair would require a new \textbf{uniform seed\textendash{}carry synchronization theorem}, with a recurrence scale bounded independently of every \(2\)-adic and \(3\)-adic resolution. No such theorem follows from saturation, asymptotic frequency, or normalized affine carry alone.
\subsection*{References}
\begin{enumerate}[leftmargin=2.2em]
\item \hypertarget{ref:R25:1}{}\hyperlink{cit:R25:1}{\texttt{turn172491view4}}
\item \hypertarget{ref:R25:2}{}\hyperlink{cit:R25:2}{\texttt{turn171785view4}}
\item \hypertarget{ref:R25:3}{}\hyperlink{cit:R25:3}{\texttt{turn171785view1}}
\item \hypertarget{ref:R25:4}{}\hyperlink{cit:R25:4}{\texttt{turn171785view2}}
\item \hypertarget{ref:R25:5}{}\hyperlink{cit:R25:5}{\texttt{turn975064view0}}
\item \hypertarget{ref:R25:6}{}\hyperlink{cit:R25:6}{\texttt{turn975064view1}}
\item \hypertarget{ref:R25:7}{}\hyperlink{cit:R25:7}{\texttt{turn171785view0}}
\item \hypertarget{ref:R25:8}{}\hyperlink{cit:R25:8}{\texttt{turn161726view0}}
\end{enumerate}

\clearpage
\section{Round 26: approach 21 --- Additive-combinatorial inverse theorem for growth words}
\subsection*{Prompt}
\begin{quote}
Complete a full proof of the Collatz conjecture using the following approach:

\begin{enumerate}
\item Additive-combinatorial inverse theorem for growth words
\end{enumerate}

\textbf{Hybrid \textperiodcentered{} P}

Show that a valuation word with many prefixes satisfying A\_t <\textasciitilde{} t log\_2 3 must have additive structure: approximate periodicity, small doubling among partial sums, or concentration in a generalized arithmetic progression. Convert that structure into a repeated valuation block and then into an impossible cycle or carry congruence.

Take as long as you need.
\end{quote}
\subsection*{Response}

\subsection*{Verdict}


Pushing this approach rigorously does \textbf{not} produce a full proof of Collatz. The obstruction is mathematical, not merely that the conjecture is famous:


\begin{quote}
The proposed implication   \[
> \text{many low-drift prefixes}
> \Longrightarrow
> \text{additive structure}
> \Longrightarrow
> \text{repeated valuation block}
> \] is false as stated.

\end{quote}


In fact, one can construct a valuation word of \textbf{maximal factor complexity} for which \textbf{every prefix} satisfies the critical growth inequality, and every finite prefix is exactly realizable by a positive Collatz seed. There is also a known theorem of Dubickas showing that a genuinely divergent Collatz parity word must have factor complexity at least \[
\frac{\log 2}{\log(3/2)}\,r
 =1.709511291\ldots\,r,
\] which points away from the simplest ``low complexity forces periodicity'' strategy. [\hypertarget{cit:R26:1}{}\hyperlink{ref:R26:1}{1}, \hypertarget{cit:R26:2}{}\hyperlink{ref:R26:2}{2}]

There is nevertheless a substantial theorem obtainable from the idea: \textbf{an odd Collatz orbit cannot remain in a bounded-width critical drift corridor forever}. The proof uses exactly the desired ingredients\textemdash{}composition counting, repeated words, and simultaneous \(2\)-adic/\(3\)-adic locality.

Tao's published theorem remains an almost-all statement rather than a pointwise theorem about every exceptional seed, and current combinatorics-on-words work still obtains restrictions on possible Collatz cycles rather than eliminating all of them. [\hypertarget{cit:R26:3}{}\hyperlink{ref:R26:3}{3}, \hypertarget{cit:R26:4}{}\hyperlink{ref:R26:4}{4}]


\bigskip



\subsection*{1. Exact valuation-word calculus}


Work with the odd-to-odd Syracuse map

\[
S(x)=\frac{3x+1}{2^{v_2(3x+1)}},
\qquad x\in 2\mathbb Z+1.
\]

Let

\[
x_{j+1}=S(x_j),\qquad
a_j=v_2(3x_j+1)\geq 1,
\qquad
A_t=\sum_{j=0}^{t-1}a_j.
\]

Put

\[
\alpha=\log_2 3=1.5849625007\ldots .
\]

For a block

\[
w=(b_0,\ldots,b_{r-1}),
\]

write

\[
B_i=\sum_{k=0}^{i-1}b_k,\qquad B=B_r,
\]

and define its affine carry

\[
C_w=\sum_{i=0}^{r-1}3^{r-1-i}2^{B_i}.
\]

Whenever an orbit follows the block \(w\),

\[
\boxed{\,2^B x_{s+r}=3^r x_s+C_w.\,} \tag{1}
\]

This identity contains both local congruence mechanisms needed below.


\subsubsection*{Block-collision lemma}


Suppose the same length-\(r\) valuation block \(w\) occurs starting at two times \(s<t\). Then

\[
\boxed{x_s\equiv x_t\pmod{2^{B+1}}} \tag{2}
\]

and

\[
\boxed{x_{s+r}\equiv x_{t+r}\pmod{3^r}.} \tag{3}
\]


\paragraph*{Proof}


Subtract the two instances of (1):

\[
2^B(x_{s+r}-x_{t+r})
   =3^r(x_s-x_t).
\]

The two endpoints \(x_{s+r},x_{t+r}\) are odd, so their difference is even. Thus the left side is divisible by \(2^{B+1}\). Since \(3^r\) is odd,

\[
2^{B+1}\mid x_s-x_t.
\]

Similarly, reducing modulo \(3^r\), and using that \(2^B\) is invertible modulo \(3^r\), gives

\[
x_{s+r}\equiv x_{t+r}\pmod{3^r}.
\]

This proves both claims. \ensuremath{\blacksquare}

Equation (3) is the full-residue form of terminal-suffix locality:

\[
x_{s+r}\equiv 2^{-B}C_w\pmod{3^r}.
\]

Thus the last \(r\) valuations determine the entire endpoint residue modulo \(3^r\), not merely whether \(3^r\mid x_{s+r}+1\).


\bigskip



\subsection*{2. Every finite valuation word is realizable}


A serious difficulty is that finite integral realizability places essentially no restriction on the word.


\subsubsection*{Finite realization lemma}


For every word \(w=(b_0,\ldots,b_{r-1})\in\mathbb N_{\geq1}^r\), there is a unique odd residue class modulo \(2^{B+1}\) whose first \(r\) Syracuse valuations equal \(w\). In particular, some positive odd integer realizes \(w\).

Indeed, by (1), exact realization is equivalent to the final quotient being odd:

\[
3^r x+C_w\equiv 2^B\pmod{2^{B+1}}. \tag{4}
\]

Because \(3^r\) is invertible modulo \(2^{B+1}\), this determines one residue class for \(x\). An induction on the length of \(w\) verifies that (4) forces each intermediate valuation to equal the prescribed \(b_i\), not merely their total.

Consequently, \textbf{every finite positive composition is a genuine Collatz valuation prefix}.

For an infinite word, the compatible residue classes define a unique \(2\)-adic seed. Explicitly, the iterate identity gives

\[
x_0
 =
\frac{2^{A_t}x_t}{3^t}
 -
\sum_{k=0}^{t-1}\frac{2^{A_k}}{3^{k+1}},
\]

and hence in \(\mathbb Z_2\),

\[
\boxed{
x_0=-\sum_{k=0}^{\infty}\frac{2^{A_k}}{3^{k+1}}.
} \tag{5}
\]

The genuinely global Collatz restriction is that this \(2\)-adic number must be an \textbf{ordinary positive integer}. No finite prefix can detect that.


\bigskip



\subsection*{3. A direct counterexample to the proposed inverse theorem}


Let

\[
p=\alpha-1=0.5849625007\ldots .
\]

We construct a binary word

\[
b_0b_1b_2\ldots,\qquad b_i\in\{0,1\},
\]

such that

\[
\sum_{i<t}b_i\leq pt\qquad\text{for every }t, \tag{6}
\]

but the word contains every finite binary word.

Enumerate all finite binary words as

\[
u_1,u_2,u_3,\ldots .
\]

Suppose a finite prefix has already been constructed, of length \(N\) and containing \(M\) ones. Its deficit below the line of slope \(p\) is

\[
D=pN-M\geq0.
\]

Before appending \(u_m\), insert enough zeros that

\[
D\geq
\max_{0\leq k\leq |u_m|}
\left(
\#_1(u_m[0,k))-pk
\right).
\]

Appending \(u_m\) then preserves (6) throughout every intermediate prefix. Repeating this procedure constructs the desired infinite word.

Now set

\[
a_i=1+b_i\in\{1,2\}.
\]

Then

\[
A_t=t+\sum_{i<t}b_i
   \leq (1+p)t
   =\alpha t
\]

for \textbf{every} \(t\).

Yet every word in \(\{1,2\}^r\) occurs as a factor. Therefore its factor complexity is maximal:

\[
\boxed{\mathcal P_a(r)=2^r.} \tag{7}
\]

Every finite prefix of this word is also an exactly realizable Collatz valuation prefix by the finite realization lemma.

Thus even the much stronger hypothesis

\[
A_t\leq t\log_2 3
\qquad\text{for every prefix }t
\]

does not imply approximate periodicity, low factor complexity, or a long repeated period.

The only thing missing from this constructed word is that its single compatible \(2\)-adic seed need not be a positive ordinary integer. That global condition is precisely where the Collatz problem re-enters.


\bigskip



\subsection*{4. Small doubling is automatic and therefore vacuous}


Let

\[
\Sigma_N=\{A_0,A_1,\ldots,A_N\}.
\]

Because all \(a_i\geq1\), the \(A_i\) are strictly increasing, so

\[
|\Sigma_N|=N+1.
\]

If \(A_N\leq(\alpha+o(1))N\), then

\[
\Sigma_N+\Sigma_N
 \subseteq
 \{0,1,\ldots,2A_N\}.
\]

Hence

\[
|\Sigma_N+\Sigma_N|
 \leq 2A_N+1
 \leq (2\alpha+o(1))|\Sigma_N|,
\]

where

\[
2\alpha=3.169925001\ldots .
\]

Therefore every low-drift word already has bounded doubling among its scalar partial sums.

A Freiman-type conclusion would merely say that \(\Sigma_N\) is concentrated in an interval or rank-one generalized arithmetic progression. But this was known from the outset:

\[
\Sigma_N\subseteq[0,\alpha N+o(N)].
\]

The ordering information has been lost. The same small-doubling estimate holds for the maximally complex construction above.

The exponential ambiguity can be quantified. Knowing that a length-\(r\) block has total \(B\) leaves

\[
\binom{B-1}{r-1}
\]

possible ordered valuation blocks. At the critical value \(B\approx\alpha r\), this is exponentially large.

So an inverse theorem for the \textbf{set of partial sums} cannot recover a valuation block. One would need an inverse theorem for the ordered path

\[
t\longmapsto (t,A_t)
\]

together with its carry data.


\bigskip



\subsection*{5. The sharp complexity obstruction}


Let \(\mathcal P_a(r)\) denote the number of distinct length-\(r\) factors in an infinite accelerated valuation word.

Define

\[
c_*=\frac{\log 2}{\log(3/2)}
   =1.709511291351\ldots .
\]


\subsubsection*{Proposition}


If an odd Syracuse orbit tends to infinity, then

\[
\boxed{
\liminf_{r\to\infty}\frac{\mathcal P_a(r)}r
\geq c_*.
} \tag{8}
\]


\paragraph*{Proof}


First note that

\[
S(x)+1
 \leq \frac{3x+1}{2}+1
 =\frac32(x+1).
\]

Therefore

\[
x_j+1\leq
\left(\frac32\right)^j(x_0+1). \tag{9}
\]

Suppose, for some \(\varepsilon>0\), that infinitely many \(r\) satisfy

\[
\mathcal P_a(r)\leq(c_*-\varepsilon)r.
\]

Set

\[
N=\left\lfloor(c_*-\varepsilon)r\right\rfloor.
\]

Among the \(N+1\) blocks beginning at positions

\[
0,1,\ldots,N,
\]

two must be identical. Say they start at \(s<t\leq N\). Their total valuation \(B\) satisfies \(B\geq r\). By (2),

\[
2^{B+1}\mid x_t-x_s.
\]

Since the orbit is divergent rather than periodic, \(x_t\neq x_s\), so

\[
|x_t-x_s|\geq2^{B+1}\geq2^{r+1}. \tag{10}
\]

On the other hand, by (9),

\[
|x_t-x_s|
<
\left(\frac32\right)^t(x_0+1)
\leq
\left(\frac32\right)^{(c_*-\varepsilon)r}(x_0+1).
\]

Because

\[
\left(\frac32\right)^{c_*}=2,
\]

the last expression equals

\[
2^r
\left(\frac32\right)^{-\varepsilon r}(x_0+1)
<2^r
\]

for all sufficiently large \(r\), contradicting (10). \ensuremath{\blacksquare}

Dubickas proved the corresponding theorem for the binary parity word of the shortcut Collatz map, with exactly the same constant \(c_*\). His proof also proceeds by converting repeated factors into divisibility by \(2^r\) and comparing that divisibility with the archimedean growth bound. [\hyperlink{ref:R26:1}{1}, \hypertarget{cit:R26:5}{}\hyperlink{ref:R26:5}{5}, \hyperlink{ref:R26:2}{2}]

This constant is therefore not incidental. It is the exact collision threshold obtained from

\[
2^r
\quad\text{versus}\quad
\left(\frac32\right)^N.
\]

A divergent orbit is forced to have enough factor complexity to avoid precisely this elementary repetition argument.


\bigskip



\subsection*{6. A theorem that the proposed approach really does prove}


Although low drift alone does not force periodicity, \textbf{uniformly critical drift} is impossible.

Define the critical discrepancy

\[
h_t=A_t-\alpha t.
\]


\subsubsection*{Critical-corridor theorem}


There is no positive odd Syracuse orbit for which

\[
\boxed{|h_t|\leq K\quad\text{for every }t} \tag{11}
\]

for some fixed \(K\).

The same conclusion holds if (11) is assumed only after some finite initial time.


\paragraph*{Step 1: Such an orbit grows at most linearly}


From the affine iterate formula,

\[
x_t
 =
\frac{3^t}{2^{A_t}}x_0
 +
\sum_{k=0}^{t-1}
 3^{t-1-k}2^{A_k-A_t}.
\]

Since \(3=2^\alpha\), this becomes

\[
\boxed{
x_t
 =
2^{-h_t}x_0
 +
\frac13\sum_{k=0}^{t-1}2^{h_k-h_t}.
} \tag{12}
\]

Under \(|h_j|\leq K\),

\[
\boxed{
x_t\leq
2^Kx_0+\frac{2^{2K}}3\,t.
} \tag{13}
\]

Thus the orbit is at most linear in \(t\).


\paragraph*{Step 2: Count the possible critical blocks}


Every length-\(r\) factor beginning at \(i\) has total

\[
A_{i+r}-A_i
 =\alpha r+h_{i+r}-h_i
 \leq\alpha r+2K.
\]

Let

\[
M_r=\left\lfloor\alpha r+2K\right\rfloor.
\]

The number of positive integer words of length \(r\) with total at most \(M_r\) is

\[
Q_r
 =
\sum_{B=r}^{M_r}\binom{B-1}{r-1}
 =
\binom{M_r}{r}. \tag{14}
\]

By Stirling's formula,

\[
Q_r^{1/r}\longrightarrow
\kappa
 :=
\frac{\alpha^\alpha}
     {(\alpha-1)^{\alpha-1}}.
\]

Numerically,

\[
\boxed{
\kappa=2.83951373049\ldots<3.
} \tag{15}
\]

This strict entropy gap is the decisive point.


\paragraph*{Step 3: Force a repeated terminal block}


Consider the \(Q_r+1\) blocks beginning at

\[
0,1,\ldots,Q_r.
\]

There are at most \(Q_r\) possible blocks by (14), so two are equal. Say they begin at \(i<j\).

By terminal locality,

\[
x_{i+r}\equiv x_{j+r}\pmod{3^r}. \tag{16}
\]

But by (13),

\[
x_{i+r},x_{j+r}
 \leq
2^Kx_0+
\frac{2^{2K}}3(Q_r+r).
\]

Because \(Q_r=(\kappa+o(1))^r\) and \(\kappa<3\), for all sufficiently large \(r\),

\[
0<x_{i+r},x_{j+r}<3^r.
\]

Together with (16), this forces

\[
x_{i+r}=x_{j+r}. \tag{17}
\]

Thus the orbit is eventually periodic.


\paragraph*{Step 4: A positive cycle cannot have bounded critical discrepancy}


Suppose a positive odd cycle has length \(d\), total valuation \(D\), and begins at \(x\). Its cycle equation is

\[
2^D x=3^d x+C
\]

with \(C>0\). Hence

\[
(2^D-3^d)x=C>0,
\]

so

\[
2^D>3^d
\quad\Longrightarrow\quad
D-\alpha d>0. \tag{18}
\]

Every traversal of the cycle therefore increases \(h_t\) by the positive amount \(D-\alpha d\). Consequently \(h_t\) becomes unbounded above, contradicting (11).

This completes the proof of the critical-corridor theorem. \ensuremath{\blacksquare}


\bigskip



\subsection*{7. A precise additive-combinatorial recurrence criterion}


The bounded-corridor theorem can be weakened. Assume only

\[
h_t\leq K\qquad\text{for all }t,
\]

and define the near-critical return set

\[
E_K=\{t:h_t\geq-K\}.
\]

For \(t\in E_K\), formula (12) gives

\[
x_t\leq
2^Kx_0+\frac{2^{2K}}3\,t. \tag{19}
\]

If both \(t-r\) and \(t\) lie in \(E_K\), then the terminal block of length \(r\) ending at \(t\) has total at most

\[
\alpha r+2K.
\]

Hence there are at most

\[
Q_r=\binom{\lfloor\alpha r+2K\rfloor}{r}
\]

such blocks.


\subsubsection*{Recurrence criterion}


If there are \(N,r\) for which

\[
\left|
\left\{
t\in[r,N]:
t-r\in E_K,\ t\in E_K
\right\}
\right|
>Q_r \tag{20}
\]

and

\[
2^Kx_0+\frac{2^{2K}}3N<3^r, \tag{21}
\]

then the orbit is periodic.

Indeed, (20) forces two equal terminal \(r\)-blocks. Their endpoints are congruent modulo \(3^r\), while (19) and (21) place both endpoints strictly between \(0\) and \(3^r\). They are therefore equal.

Since

\[
Q_r=(\kappa+o(1))^r,\qquad \kappa<3,
\]

there is a nonempty exponential window

\[
\log\kappa<\lambda<\log3.
\]

At the scale

\[
N\asymp e^{\lambda r},
\]

we have simultaneously

\[
Q_r=o(N)
\qquad\text{and}\qquad
N=o(3^r).
\]

Thus the genuinely useful inverse theorem would be:


\begin{quote}
A hypothetical counterexample has linearly many pairs of near-critical return times separated by some \(r\), at a scale   \[
> e^{(\log\kappa+\varepsilon)r}
> \lesssim N
> \lesssim
> e^{(\log3-\varepsilon)r}.
> \]

\end{quote}


That statement would immediately activate the recurrence criterion.

This is the correct additive-combinatorial target. It concerns the additive correlation

\[
E_K\cap(E_K+r),
\]

not merely small doubling of the partial-sum set.


\bigskip



\subsection*{8. Why a minimal counterexample does not automatically satisfy the corridor hypothesis}


The exact logarithmic evolution is

\[
\log_2 x_t-\log_2 x_0
 =
\alpha t-A_t
 +
\sum_{j=0}^{t-1}
 \log_2\left(1+\frac1{3x_j}\right).
\]

Equivalently,

\[
\boxed{
h_t
 =
\log_2x_0-\log_2x_t
 +
E_t,
\qquad
E_t=
\sum_{j<t}
\log_2\left(1+\frac1{3x_j}\right).
} \tag{22}
\]

For a least counterexample, every later orbit value is at least \(x_0\), so

\[
h_t\leq E_t.
\]

But \(E_t\) need not be bounded. It is bounded if, for example,

\[
\sum_j\frac1{x_j}<\infty,
\]

but convergence of this reciprocal sum is not known for a hypothetical divergent orbit.

Therefore even obtaining a uniform upper bound

\[
h_t\leq K
\]

requires new information. A counterexample could instead have:

\[
h_t\longrightarrow-\infty,
\]

corresponding to genuine multiplicative growth; or it could have large negative excursions and only sparse returns to the critical line.

The maximally complex symbolic construction in Section 3 demonstrates that these possibilities cannot be excluded by prefix drift alone.


\bigskip



\subsection*{9. The separate cycle gap}


Even if additive structure forced eventual periodicity, this would not automatically finish Collatz.

For every hypothetical positive odd cycle,

\[
(2^D-3^d)x=C_w,
\]

with

\[
D>\alpha d.
\]

Periodicity is therefore perfectly compatible with the formal affine identities. To get a contradiction, one needs one of the following stronger conclusions:

\[
D\leq\alpha d,
\]

which would make the left side nonpositive; or an arithmetic theorem proving

\[
2^D-3^d\nmid C_w
\]

for every nontrivial candidate word.

A generic repeated factor does not guarantee that the resulting period has low average valuation. Recent work identifying Christoffel or balanced words as extremizers gives restrictions on hypothetical cycles, but does not establish this universal nondivisibility. [\hyperlink{ref:R26:4}{4}]


\bigskip



\subsection*{Final assessment}


The proposed approach, as written, has three exact gaps:

\begin{enumerate}
\item \textbf{Drift does not force additive structure.}   There are maximally complex words with \(A_t\leq\alpha t\) for every \(t\), and every finite prefix is an exact Collatz prefix.

\item \textbf{Small doubling is too coarse.}   It is automatic for the partial sums and loses the ordering of the valuations.

\item \textbf{Periodicity alone does not rule out a nontrivial cycle.}   One must preserve a low-drift inequality for the actual period or prove a new carry-divisibility obstruction.
\end{enumerate}

What the approach does yield is the rigorous theorem

\[
\boxed{
\sup_t|A_t-t\log_2 3|=\infty
\quad\text{for every positive odd Collatz orbit}.
}
\]

More quantitatively, the entropy gap

\[
2.839513730\ldots<3
\]

shows that a sufficiently recurrent near-critical word must produce two equal orbit states through terminal \(3\)-adic locality.

The viable completion of Approach 21 is therefore not a standard Freiman theorem. It would have to be an \textbf{arithmetic, ordered inverse theorem} proving logarithmic-scale recurrence of near-critical valuation blocks for the one special word whose \(2\)-adic seed is an ordinary positive integer. That is the precise theorem still missing; the raw growth-word statement is false.
\subsection*{References}
\begin{enumerate}[leftmargin=2.2em]
\item \hypertarget{ref:R26:1}{}\hyperlink{cit:R26:1}{\texttt{turn132506view1}}
\item \hypertarget{ref:R26:2}{}\hyperlink{cit:R26:2}{\texttt{turn132506view3}}
\item \hypertarget{ref:R26:3}{}\hyperlink{cit:R26:3}{\texttt{turn474876search2}}
\item \hypertarget{ref:R26:4}{}\hyperlink{cit:R26:4}{\texttt{turn474876search3}}
\item \hypertarget{ref:R26:5}{}\hyperlink{cit:R26:5}{\texttt{turn132506view2}}
\end{enumerate}

\clearpage
\section{Round 27: approach 22 --- Entropy decrement plus arithmetic rigidity}
\subsection*{Prompt}
\begin{quote}
Complete a full proof of the Collatz conjecture using the following approach:

\begin{enumerate}
\item Entropy decrement plus arithmetic rigidity
\end{enumerate}

\textbf{Hybrid \textperiodcentered{} P}

Adapt Tao's entropy-decrement philosophy to a single exceptional orbit. Either distant valuation blocks decorrelate, giving negative drift, or mutual information remains large and repeatedly lowers entropy. Eventually the word should become finite-state or almost periodic, where exact integrality can rule it out.

Take as long as you need.
\end{quote}
\subsection*{Response}

\subsubsection*{Verdict}


Carrying this program through produces a rigorous \textbf{single-orbit entropy-decrement theorem}, a sharp entropy deficit for any hypothetical divergent orbit, and a quantitative lower bound on how rapidly such an orbit would have to escape.

It does \textbf{not} produce a complete proof of Collatz. The proposed closing dichotomy fails in both directions:

\[
\text{decorrelation}\;\not\Rightarrow\;\text{negative drift},
\]

and

\[
\text{low or zero entropy}\;\not\Rightarrow\;\text{finite-state or eventually periodic}.
\]

The missing step is an orbitwise arithmetic-rigidity theorem about positive rational \(2\)-adic seeds. A separate exclusion of nontrivial cycles is also required. Tao's July 2026 revision likewise states that the full conjecture remains beyond current methods. [\hypertarget{cit:R27:1}{}\hyperlink{ref:R27:1}{1}]

What follows is the complete argument up to that precise obstruction.


\bigskip



\subsection*{1. Accelerated Collatz dynamics}


Work with the odd-to-odd Syracuse map

\[
U(n)=\frac{3n+1}{2^{v_2(3n+1)}},
\qquad n\in 2\mathbb Z+1.
\]

Write

\[
n_{t+1}=U(n_t),\qquad
a_t=v_2(3n_t+1)\ge 1,
\]

and define

\[
A_{s,H}=\sum_{i=0}^{H-1}a_{s+i}.
\]

For a finite word

\[
w=(a_0,\ldots ,a_{H-1}),
\]

put

\[
A(w)=a_0+\cdots+a_{H-1}
\]

and

\[
C(w)=\sum_{i=0}^{H-1}3^{H-1-i}2^{a_0+\cdots+a_{i-1}},
\]

where the exponent in the \(i=0\) term is \(0\).

Induction on \(H\) gives the exact affine identity

\[
\boxed{
2^{A(w)}n_H=3^Hn_0+C(w).
}
\tag{1}
\]

This is the integer version of Tao's Syracuse affine formula. [\hypertarget{cit:R27:2}{}\hyperlink{ref:R27:2}{2}]

For a block ending at time \(t\),

\[
v=(a_{t-r},\ldots ,a_{t-1}),
\]

equation (1) gives

\[
2^{A(v)}n_t=3^rn_{t-r}+C(v).
\]

Reducing modulo \(3^r\),

\[
\boxed{
n_t\equiv 2^{-A(v)}C(v)\pmod{3^r}.
}
\tag{2}
\]

Thus the final \(r\) valuation symbols determine the endpoint modulo \(3^r\). In particular,

\[
3^r\mid n_t+1
\quad\Longleftrightarrow\quad
3^r\mid C(v)+2^{A(v)}.
\]

This is the terminal-suffix locality mechanism needed below.


\bigskip



\subsection*{2. What a divergent orbit would force}


Suppose an orbit never reaches \(1\).

If it is bounded, it repeats a state and hence enters a cycle. Therefore, to study divergent counterexamples, we may assume the \(n_t\) are all distinct. In that case \(n_t\to\infty\), because only finitely many positive integers lie below any fixed bound.

From

\[
n_{t+1}
=
\frac{3n_t}{2^{a_t}}
\left(1+\frac1{3n_t}\right)
\]

we obtain

\[
\log n_{t+1}-\log n_t
=
\log3-a_t\log2+
\log\left(1+\frac1{3n_t}\right).
\]

Summing from \(t=0\) to \(N-1\),

\[
A_N\log2
=
N\log3+\log n_0-\log n_N
+
\sum_{t=0}^{N-1}
\log\left(1+\frac1{3n_t}\right),
\tag{3}
\]

where \(A_N=a_0+\cdots+a_{N-1}\).

Since the \(n_t\) are distinct positive integers,

\[
\sum_{t<N}\frac1{n_t}
\leq 1+\log N.
\]

Using \(\log(1+x)\leq x\) and discarding the negative term \(-\log n_N\), equation (3) gives

\[
\boxed{
A_N\leq \lambda N+O_{n_0}(\log N),
\qquad
\lambda=\log_2 3=1.5849625007\ldots .
}
\tag{4}
\]

Consequently,

\[
\limsup_{N\to\infty}\frac{A_N}{N}\leq\lambda.
\tag{5}
\]

So a divergent orbit must use substantially fewer powers of \(2\) than a typical random orbit, whose average valuation is \(2\). In Tao's random \(2\)-adic model the valuations are independent copies of the geometric law

\[
\Pr(a=k)=2^{-k},
\]

which has mean \(2\), giving the familiar negative drift

\[
\log3-2\log2=\log(3/4)<0.
\]

That random law is rigorously valid for Haar-random \(2\)-adic seeds and approximately valid for suitably randomized large integers. [\hypertarget{cit:R27:3}{}\hyperlink{ref:R27:3}{3}, \hypertarget{cit:R27:4}{}\hyperlink{ref:R27:4}{4}]


\bigskip



\subsection*{3. A stationary measure from one exceptional orbit}


Choose a random time \(T_N\), uniformly from a long initial portion of the orbit, and inspect the shifted valuation sequence

\[
(\ldots,a_{T_N-1},a_{T_N},a_{T_N+1},\ldots).
\]

Passing to a subsequence gives a shift-invariant empirical limit \(\mu\) on valuation sequences. The first-moment estimate (4) gives

\[
\boxed{
\mathbb E_\mu a_0\leq\lambda.
}
\tag{6}
\]

This construction is legitimate despite the countable alphabet because the uniform first-moment bound gives tightness.

Let

\[
H_\mu(a_0)
\]

denote Shannon entropy in natural logarithms.

For a positive-integer-valued random variable with mean \(m\), entropy is maximized by the geometric distribution of mean \(m\). Its entropy is

\[
g(m)=m\log m-(m-1)\log(m-1).
\tag{7}
\]

Indeed, compare any law \(p\) of mean \(m\) with

\[
q_k=\frac1m\left(1-\frac1m\right)^{k-1}.
\]

The nonnegativity of \(D(p\Vert q)\) gives

\[
H(p)
\leq
-\sum_kp_k\log q_k
=
g(m).
\]

Since \(g\) is increasing and \(\mathbb E_\mu a_0\leq\lambda\),

\[
H_\mu(a_0)\leq g(\lambda).
\]

The entropy rate

\[
h(\mu)=
\lim_{H\to\infty}
\frac1H H_\mu(a_0,\ldots,a_{H-1})
\]

therefore satisfies

\[
\boxed{
h(\mu)\leq g(\lambda)
=1.0436328159\ldots .
}
\tag{8}
\]

But

\[
\log3=1.0986122887\ldots,
\]

so there is a strict entropy gap

\[
\boxed{
\delta
=
\log3-g(\lambda)
=
0.0549794728\ldots\text{ nats}
=
0.0793186128\ldots\text{ bits}.
}
\tag{9}
\]

For comparison, Tao's geometric valuation law has entropy \(\log4\), which exceeds the \(\log3\) capacity of one ternary digit; that surplus underlies his fine-scale \(3\)-adic mixing heuristic and theorem. [\hypertarget{cit:R27:5}{}\hyperlink{ref:R27:5}{5}]


\bigskip



\subsection*{4. Terminal locality gives a strict \(3\)-adic entropy deficit}


Let

\[
R_t^{(r)}=n_t\bmod 3^r.
\]

By equation (2), \(R_t^{(r)}\) is a deterministic function of

\[
(a_{t-r},\ldots,a_{t-1}).
\]

Therefore, in the stationary empirical limit,

\[
\begin{aligned}
H_\mu(R_0^{(r)})
&\leq
H_\mu(a_{-r},\ldots,a_{-1})\\
&\leq
rH_\mu(a_0)\\
&\leq
r g(\lambda).
\end{aligned}
\tag{10}
\]

Thus

\[
\boxed{
\frac{H_\mu(R_0^{(r)})}{r\log3}
\leq
\frac{g(\lambda)}{\log3}
=
0.9499555272\ldots .
}
\tag{11}
\]

More sharply, using the entropy rate,

\[
\limsup_{r\to\infty}
\frac{H_\mu(R_0^{(r)})}{r\log3}
\leq
\frac{h(\mu)}{\log3}
\leq0.9499555272\ldots .
\tag{12}
\]

So any divergent Collatz orbit would generate a \(3\)-adic state distribution with at least a \(5.004\%\) entropy-dimension deficit.

If \(u_r\) denotes the uniform distribution on \(\mathbb Z/3^r\mathbb Z\), then

\[
D\!\left(\operatorname{law}(R_0^{(r)})\middle\Vert u_r\right)
=
r\log3-H_\mu(R_0^{(r)}),
\]

and hence

\[
\boxed{
D\!\left(\operatorname{law}(R_0^{(r)})\middle\Vert u_r\right)
\geq
\delta r.
}
\tag{13}
\]

This is a rigorous and fairly strong conclusion:


\begin{quote}
A divergent orbit cannot be even entropy-dimensionally equidistributed in the \(3\)-adics. It must live on a ternary fractal of entropy dimension at most \(0.949956\).

\end{quote}



\bigskip



\subsection*{5. The exact single-orbit entropy decrement}


Let

\[
H_H=H_\mu(a_0,\ldots,a_{H-1}),
\qquad
\eta_H=\frac{H_H}{H}.
\]

Define adjacent past and future blocks

\[
P_H=(a_{-H},\ldots,a_{-1}),
\qquad
F_H=(a_0,\ldots,a_{H-1}).
\]

By stationarity,

\[
H(P_H)=H(F_H)=H_H,
\]

while their joint variable is a block of length \(2H\). Therefore

\[
\begin{aligned}
I_\mu(P_H;F_H)
&=H(P_H)+H(F_H)-H(P_H,F_H)\\
&=2H_H-H_{2H}\\
&=2H(\eta_H-\eta_{2H}).
\end{aligned}
\tag{14}
\]

Take dyadic scales

\[
H_j=2^jH_0.
\]

Then

\[
\sum_{j=0}^{J-1}
\frac{I_\mu(P_{H_j};F_{H_j})}{2H_j}
=
\eta_{H_0}-\eta_{H_J}
\leq g(\lambda).
\tag{15}
\]

Consequently, for some \(0\leq j<J\),

\[
\boxed{
I_\mu(P_{H_j};F_{H_j})
\leq
\frac{2g(\lambda)}{J}H_j.
}
\tag{16}
\]

Letting \(J\to\infty\), one obtains a sequence of scales \(H\to\infty\) such that

\[
\boxed{
I_\mu(P_H;F_H)=o(H).
}
\tag{17}
\]

This is exactly the entropy-decrement phenomenon: if normalized mutual information stayed bounded below at every scale, then the normalized block entropy would repeatedly decrease and eventually become negative.

Tao introduced this mechanism to choose a scale where an arithmetic random variable and a local block become weakly dependent; equations (14)\textendash{}(17) are its particularly clean stationary single-orbit analogue. [\hypertarget{cit:R27:6}{}\hyperlink{ref:R27:6}{6}, \hypertarget{cit:R27:7}{}\hyperlink{ref:R27:7}{7}]

Furthermore, \(R_0^{(r)}\) is a function of \(P_H\) whenever \(r\leq H\). By data processing,

\[
I_\mu(F_H;R_0^{(r)})
\leq I_\mu(F_H;P_H).
\tag{18}
\]

In particular, along selected scales,

\[
\boxed{
I_\mu(F_H;R_0^{(H)})=o(H).
}
\tag{19}
\]

So the hoped-for weak independence between a future valuation block and the current \(3\)-adic carry state really can be obtained for a single hypothetical orbit.

This completes the entropy-decrement half of the proposed program.


\bigskip



\subsection*{6. Why decorrelation does not imply negative drift}


Here is an exact obstruction.

Take an independent valuation process with

\[
\Pr(a_t=1)=\Pr(a_t=2)=\frac12.
\]

Then disjoint blocks are exactly independent, so their mutual information is zero. But

\[
\mathbb E a_t=\frac32<\lambda,
\]

and the corresponding affine logarithmic drift is

\[
\log3-\frac32\log2
=
0.0588915178\ldots>0.
\]

Thus

\[
\boxed{
\text{perfect decorrelation can coexist with positive Collatz drift.}
}
\]

The problem is that independence says nothing about the one-symbol marginal. Tao's negative-drift model requires not merely independence, but the specific geometric law of mean \(2\).

This example is not an artificial violation of the local Collatz rules. Convert the valuation sequence to a parity sequence

\[
1\,0^{a_0-1}\,
1\,0^{a_1-1}\,
1\,0^{a_2-1}\cdots .
\]

Every infinite binary parity sequence is realized by exactly one \(2\)-adic Collatz seed. [\hypertarget{cit:R27:8}{}\hyperlink{ref:R27:8}{8}]

Hence a deterministic Bernoulli-generic realization of the above process gives a genuine \(2\)-adic Collatz orbit satisfying all finite valuation congruences, all affine identities, and all terminal-suffix identities. What generally fails is that the \(2\)-adic seed is an \textbf{ordinary positive integer}.

Therefore the missing ingredient cannot be another Shannon inequality. It must use the global condition

\[
x\in\mathbb Z_{>0}\subset\mathbb Z_2.
\]


\bigskip



\subsection*{7. Why entropy collapse does not imply periodicity}


The other branch of the proposed dichotomy also fails.

Let

\[
\alpha=\frac{\sqrt5-1}{2}
\]

and define the Sturmian sequence

\[
b_t=\lfloor(t+1)\alpha\rfloor-\lfloor t\alpha\rfloor\in\{0,1\},
\]

followed by

\[
a_t=2-b_t\in\{1,2\}.
\]

The frequency of \(b_t=1\) is \(\alpha\), so

\[
\lim_{N\to\infty}\frac1N\sum_{t<N}a_t
=
2-\alpha
=
1.381966011\ldots
<
\lambda.
\]

This word is aperiodic: eventual periodicity would force its symbol frequency to be rational, whereas \(\alpha\) is irrational.

It has exactly \(H+1\) distinct factors of length \(H\), so its entropy rate is zero:

\[
h=0.
\]

It is also uniformly recurrent, which is a strong form of combinatorial almost periodicity.

Again, its associated parity sequence determines a valid \(2\)-adic Collatz seed.

Thus

\[
\boxed{
\text{zero entropy + subcritical mean + almost periodicity}
\not\Rightarrow
\text{eventual periodicity}.
}
\]

Entropy decrement only controls average Shannon information. It does not produce a finite deterministic state space. Even

\[
H(a_0,\ldots,a_{H-1})=O(\log H)
\]

is compatible with a genuinely aperiodic word.

To invoke a theorem such as Morse\textendash{}Hedlund, one would need an exact bound of the form

\[
p(H)\leq H
\]

for the number of length-\(H\) factors at some finite \(H\). An \(o(H)\) Shannon-entropy estimate is vastly weaker.


\bigskip



\subsection*{8. What eventual periodicity would give}


There is nevertheless a clean arithmetic endgame if eventual periodicity could be proved.


\subsubsection*{Proposition}


If the valuation sequence of a positive-integer Syracuse orbit is eventually periodic, then the orbit itself eventually enters a positive integer cycle.


\paragraph*{Proof}


Suppose that after some time the valuation block

\[
w=(a_0,\ldots,a_{p-1})
\]

repeats forever. Let

\[
A=A(w),\qquad C=C(w),
\]

and sample the orbit once per period:

\[
x_k=n_{t_0+kp}.
\]

Equation (1) gives

\[
2^Ax_{k+1}=3^px_k+C.
\tag{20}
\]

Set

\[
D=2^A-3^p\neq0.
\]

The fixed point of this affine recurrence is \(C/D\), and subtraction gives

\[
Dx_{k+1}-C
=
\frac{3^p}{2^A}(Dx_k-C).
\]

After \(k\) iterations,

\[
2^{Ak}(Dx_k-C)
=
3^{pk}(Dx_0-C).
\tag{21}
\]

The left side is divisible by \(2^{Ak}\). Since \(3\) is odd, equation (21) implies

\[
2^{Ak}\mid Dx_0-C
\]

for every \(k\). Therefore

\[
Dx_0-C=0.
\]

Hence

\[
x_0=\frac{C}{D}.
\]

Because \(x_0,C>0\), necessarily \(D>0\), and \(x_k=x_0\) for every \(k\). The intermediate \(p-1\) values also repeat. \ensuremath{\blacksquare}

Thus an eventually periodic valuation sequence cannot describe a divergent positive-integer orbit.

However, it may describe a nontrivial positive cycle. The cycle equation is

\[
\boxed{
n=\frac{C(w)}{2^{A(w)}-3^{|w|}},
}
\tag{22}
\]

together with the requirement that all valuations in \(w\) be exact.

The trivial odd cycle has

\[
w=(2),\qquad C(w)=1,\qquad
n=\frac1{4-3}=1.
\]

Nothing in entropy theory excludes other periodic words: every periodic word has entropy zero. Excluding all other admissible positive integral solutions of (22) is the nontrivial-cycle problem itself.


\bigskip



\subsection*{9. The exact rigidity theorem that would complete the divergent-orbit half}


The needed statement can now be isolated precisely.


\subsubsection*{Subcritical rationality rigidity}


Let \(n_0\) be a positive odd integer with valuation sequence \(a_0,a_1,\ldots\). If

\[
\limsup_{N\to\infty}\frac1N\sum_{t<N}a_t
\leq\log_2 3,
\tag{23}
\]

then the valuation sequence is eventually periodic.

Call this statement \textbf{SRR}.

The drift calculation proves that every nonperiodic counterexample would satisfy (23). SRR would then make its valuation sequence eventually periodic, and the proposition above would turn it into a cycle, contradicting nonperiodicity.

SRR is a restricted form of a known \(2\)-adic rigidity problem. Every valuation sequence determines a parity code and therefore a unique \(2\)-adic seed. Lagarias's Periodicity Conjecture asserts, in effect, that rational \(2\)-adic seeds have eventually periodic parity codes. Positive integers are rational \(2\)-adics, so that conjecture would imply SRR without even requiring the subcritical-average assumption. The general Periodicity Conjecture remains open. [\hypertarget{cit:R27:9}{}\hyperlink{ref:R27:9}{9}]

What entropy decrement would need to prove is an inverse theorem of the following form:

\[
\begin{gathered}
x\in\mathbb Z_{>0},\\
\mathbb E_\mu a_0\leq\lambda,\\
I_\mu(P_H;F_H)=o(H)
\end{gathered}
\quad\Longrightarrow\quad
\text{the valuation word is eventually periodic}.
\tag{24}
\]

Without the condition \(x\in\mathbb Z_{>0}\), implication (24) is false by both the Bernoulli and Sturmian examples.

The arithmetic content must therefore show that a low-information symbolic process cannot simultaneously satisfy the infinite family of \(2\)-adic rationality conditions defining the positive integer seed and the \(3\)-adic endpoint conditions supplied by terminal locality.

That is the precise unsolved ``arithmetic rigidity'' theorem.


\bigskip



\subsection*{10. A conditional full proof}


The proposed approach becomes a complete proof under two sharply stated arithmetic results.


\paragraph*{Rigidity hypothesis}


Every positive-integer accelerated orbit satisfying

\[
\limsup A_N/N\leq\log_2 3
\]

has an eventually periodic valuation sequence.


\paragraph*{Cycle hypothesis}


For every admissible finite valuation word \(w\),

\[
\frac{C(w)}{2^{A(w)}-3^{|w|}}\in\mathbb Z_{>0}
\]

only when \(w\) is a repetition of the trivial word \((2)\).


\subsubsection*{Conditional theorem}


The Rigidity and Cycle hypotheses imply the Collatz conjecture.


\paragraph*{Proof}


Take a positive odd integer \(n_0\).

If its orbit repeats, it enters a positive cycle. The Cycle hypothesis says this is the cycle containing \(1\).

Suppose instead that the orbit does not repeat. Then its values are distinct, and equation (4) gives

\[
\limsup_{N\to\infty}\frac{A_N}{N}\leq\log_2 3.
\]

The Rigidity hypothesis makes the valuation word eventually periodic. By the proposition in Section 8, the orbit then enters a positive cycle, contradicting the assumption that it never repeats.

Thus every odd starting value reaches \(1\). Removing the initial powers of \(2\) proves the result for every positive integer. \ensuremath{\blacksquare}

This is logically complete, but the two displayed hypotheses are not presently proved. The first is the pointwise adelic-rigidity step; the second is nontrivial-cycle exclusion.


\bigskip



\subsection*{11. A genuine unconditional consequence: superlinear escape}


Although the full proof does not close, the entropy calculation and terminal-suffix theorem yield a quantitative result for every hypothetical nonperiodic orbit.

Define

\[
M_N=\max_{0\leq t\leq N} n_t
\]

and

\[
d_*=\frac{g(\lambda)}{\log3}
=
0.949955527188\ldots .
\]


\subsubsection*{Theorem}


Every nonperiodic positive-integer Syracuse orbit satisfies

\[
\boxed{
\liminf_{N\to\infty}
\frac{\log M_N}{\log N}
\geq
\frac1{d_*}
=
1.0526808586\ldots .
}
\tag{25}
\]

Equivalently,

\[
M_N\geq N^{\,1.0526808586-o(1)}.
\tag{26}
\]


\paragraph*{Proof}


Fix

\[
\rho>\lambda
\]

so close to \(\lambda\) that

\[
g(\rho)<\log3.
\]

Let

\[
H=\lfloor\log_3 M_N\rfloor,
\qquad
q=\lfloor N/H\rfloor.
\]

Partition the first \(qH\) valuation symbols into \(q\) disjoint blocks \(w_0,\ldots,w_{q-1}\) of length \(H\), and write

\[
s_j=A(w_j).
\]

From equation (4),

\[
\sum_{j=0}^{q-1}s_j
=
A_{qH}
\leq
\lambda qH+O(\log N).
\tag{27}
\]

Each \(s_j\geq H\). If \(b\) blocks have \(s_j>\rho H\), then

\[
\sum_j s_j
\geq
(q-b)H+b\rho H
=
qH+b(\rho-1)H.
\]

Combining this with (27),

\[
b
\leq
\frac{\lambda-1}{\rho-1}q+o(q).
\]

Therefore a positive proportion

\[
\frac{\rho-\lambda}{\rho-1}+o(1)
\]

of the blocks satisfy

\[
s_j\leq\rho H.
\tag{28}
\]

The number of positive-integer words of length \(H\) with total sum at most \(\rho H\) is

\[
\sum_{s=H}^{\lfloor\rho H\rfloor}
\binom{s-1}{H-1}
=
\binom{\lfloor\rho H\rfloor}{H}.
\tag{29}
\]

Stirling's estimate gives

\[
\log\binom{\lfloor\rho H\rfloor}{H}
=
g(\rho)H+O_\rho(\log H).
\tag{30}
\]

Now use terminal-suffix locality. A block \(w_j\) determines its endpoint

\[
n_{(j+1)H}\pmod{3^H}.
\]

Since the orbit is nonperiodic, these endpoints are distinct. Also,

\[
3^H\leq M_N<3^{H+1},
\]

so any one residue class modulo \(3^H\) contains at most three positive integers below \(M_N\).

Consequently, the number of good blocks satisfying (28) is at most

\[
3\binom{\lfloor\rho H\rfloor}{H}
=
\exp\big(g(\rho)H+O_\rho(\log H)\big).
\]

Since a positive proportion of all \(q\) blocks are good,

\[
q
\leq
\exp\big(g(\rho)H+O_\rho(\log H)\big).
\]

Thus

\[
N
\leq
(\log M_N)^{O_\rho(1)}
M_N^{g(\rho)/\log3}.
\tag{31}
\]

Letting \(\rho\downarrow\lambda\) and absorbing powers of \(\log M_N\) gives

\[
N
\leq
M_N^{d_*+o(1)}.
\]

Inverting proves (25). \ensuremath{\blacksquare}

This is strictly stronger than the trivial fact \(M_N\gg N\) coming from distinctness. It says that a divergent orbit would have to become polynomially sparse among the integers, with exponent at least \(1.05268\).

It is still compatible with exponential escape, so it is not a contradiction.


\bigskip



\subsection*{12. Final assessment of approach 22}


The program succeeds up to the following exact point:

\[
\begin{array}{c}
\text{hypothetical divergent positive orbit}\\[2mm]
\Downarrow\\[2mm]
\mathbb E a_0\leq\log_2 3\\[2mm]
\Downarrow\\[2mm]
h(\mu)\leq1.0436328<\log3\\[2mm]
\Downarrow\\[2mm]
\text{\(3\)-adic entropy dimension}\leq0.9499555\\[2mm]
\Downarrow\\[2mm]
\text{dyadic scales with }I(P_H;F_H)=o(H).
\end{array}
\]

The desired next arrow,

\[
\text{low mutual information}
+
\text{ordinary-integer realizability}
\Longrightarrow
\text{eventual periodicity},
\]

is the unsolved arithmetic-rigidity theorem. It is false in the full \(2\)-adic symbolic system and therefore must exploit the positive rational seed condition in an essentially global way.

This diagnosis is consistent with current orbit-mixing work: a 2026 preprint reducing Collatz to a low-dimensional bit-balance problem explicitly leaves the decisive \textbf{orbitwise} balance statement open and distinguishes it from map-level or ensemble-level mixing. [\hypertarget{cit:R27:10}{}\hyperlink{ref:R27:10}{10}]

So approach 22 is not presently a full proof, but it does yield:

\[
\boxed{
\begin{aligned}
&\text{an exact single-orbit entropy decrement;}\\
&\text{a linear \(3\)-adic entropy deficit;}\\
&\text{a \(0.949956\) upper entropy dimension;}\\
&\text{a superlinear escape bound }M_N\geq N^{1.05268-o(1)};\\
&\text{and a precise reduction to subcritical rationality rigidity}\\
&\text{plus nontrivial-cycle exclusion.}
\end{aligned}
}
\]

Claiming a completed Collatz proof from the entropy dichotomy would conceal exactly those two unresolved arithmetic steps.
\subsection*{References}
\begin{enumerate}[leftmargin=2.2em]
\item \hypertarget{ref:R27:1}{}\hyperlink{cit:R27:1}{\texttt{turn597428view0}}
\item \hypertarget{ref:R27:2}{}\hyperlink{cit:R27:2}{\texttt{turn611244view2}}
\item \hypertarget{ref:R27:3}{}\hyperlink{cit:R27:3}{\texttt{turn597428view1}}
\item \hypertarget{ref:R27:4}{}\hyperlink{cit:R27:4}{\texttt{turn597428view2}}
\item \hypertarget{ref:R27:5}{}\hyperlink{cit:R27:5}{\texttt{turn723765view0}}
\item \hypertarget{ref:R27:6}{}\hyperlink{cit:R27:6}{\texttt{turn611244view1}}
\item \hypertarget{ref:R27:7}{}\hyperlink{cit:R27:7}{\texttt{turn611244view0}}
\item \hypertarget{ref:R27:8}{}\hyperlink{cit:R27:8}{\texttt{turn597428view4}}
\item \hypertarget{ref:R27:9}{}\hyperlink{cit:R27:9}{\texttt{turn597428view5}}
\item \hypertarget{ref:R27:10}{}\hyperlink{cit:R27:10}{\texttt{turn723765view2}}
\end{enumerate}

\clearpage
\section{Round 28: approach 23 --- A pointwise Bellman function}
\subsection*{Prompt}
\begin{quote}
Complete a full proof of the Collatz conjecture using the following approach:

\begin{enumerate}
\item A pointwise Bellman function
\end{enumerate}

\textbf{Hybrid \textperiodcentered{} P}

Search for a function V(n) = log n + phi(n mod M) and a bounded stopping rule k(n) <= K such that V(T\textasciicircum{}\{k(n)\} n) < V(n) for every n > 1. A successful finite table for phi, together with exact inequalities for every residue class, would be a complete proof.

Take as long as you need.
\end{quote}
\subsection*{Response}

\subsection*{Verdict}


Under the standing shortcut convention

\[
T(n)=
\begin{cases}
n/2,&n\equiv0\pmod 2,\\[2mm]
(3n+1)/2,&n\equiv1\pmod 2,
\end{cases}
\]

the proposed Bellman certificate \textbf{cannot exist}. There is an explicit obstruction for every possible modulus \(M\), every possible table \(\phi\), and every finite horizon \(K\).

Thus this investigation does not produce a proof of Collatz. It produces something definitive in the other direction:


\begin{quote}
\textbf{Approach 23 is impossible in its stated finite-table, bounded-horizon form.}

\end{quote}


The shortcut normalization above is the standard one used in much of the \(3x+1\) literature. [\hypertarget{cit:R28:1}{}\hyperlink{ref:R28:1}{1}]


\bigskip



\subsubsection*{1. Why the proposed certificate would have proved Collatz}


Define

\[
V(n)=\log n+\phi(n\bmod M),
\]

where \(\phi:\mathbb Z/M\mathbb Z\to\mathbb R\). Suppose there were a rule

\[
1\le k(n)\le K
\]

such that

\[
V\bigl(T^{k(n)}n\bigr)<V(n)
\qquad\text{for every }n>1.
\tag{1}
\]

This really would be a complete proof.


\paragraph*{Certificate lemma}


Assume (1). Then every positive integer reaches \(1\).


\paragraph*{Proof}


Start with \(n_0>1\). As long as the orbit has not reached \(1\), define

\[
n_{i+1}=T^{k(n_i)}(n_i).
\]

Then

\[
V(n_{i+1})<V(n_i),
\]

so \(V(n_i)\) is strictly decreasing.

Let

\[
\phi_{\min}=\min_{r\in\mathbb Z/M\mathbb Z}\phi(r).
\]

Because

\[
\log n_i
=
V(n_i)-\phi(n_i\bmod M)
\le V(n_0)-\phi_{\min},
\]

we obtain

\[
n_i\le \exp\!\bigl(V(n_0)-\phi_{\min}\bigr).
\]

Therefore every \(n_i\) lies in a fixed finite set of positive integers. But a strictly decreasing sequence of values of the fixed function \(V\) cannot visit infinitely many points in a finite set. Hence the process must terminate. Since the descent condition is assumed at every \(n>1\), termination can occur only when the original Collatz orbit has reached \(1\). \ensuremath{\blacksquare}

So the logic of the Bellman idea is sound. The problem is that the required certificate is structurally impossible.


\bigskip



\subsection*{2. The universal obstruction}


Define the finite-horizon Bellman operator

\[
\mathcal B_KV(n)=\min_{1\le j\le K}V(T^j n).
\]

The desired assertion is

\[
\mathcal B_KV(n)<V(n)
\qquad(n>1).
\tag{2}
\]

The following theorem rules it out.


\subsubsection*{Theorem: all-odd cylinder obstruction}


Let \(M,K\ge1\), and let

\[
\phi:\mathbb Z/M\mathbb Z\to\mathbb R
\]

be arbitrary. Set

\[
V(n)=\log n+\phi(n\bmod M).
\]

Then there exists \(N>1\) such that

\[
V(T^jN)>V(N)
\qquad\text{for every }1\le j\le K.
\tag{3}
\]

In particular,

\[
\mathcal B_KV(N)>V(N).
\]


\paragraph*{Proof}


Take

\[
N=2^{K+1}M-1.
\tag{4}
\]

We claim that for every \(0\le j\le K\),

\[
T^jN=3^j2^{K+1-j}M-1.
\tag{5}
\]

This is immediate for \(j=0\). Suppose it holds for some \(j<K\). Since

\[
3^j2^{K+1-j}M
\]

is even, \(T^jN\) is odd. Therefore

\[
\begin{aligned}
T^{j+1}N
&=\frac{3\left(3^j2^{K+1-j}M-1\right)+1}{2}\\
&=3^{j+1}2^{K-j}M-1,
\end{aligned}
\]

which proves (5) by induction.

Every one of these iterates satisfies

\[
T^jN\equiv -1\pmod M.
\tag{6}
\]

Moreover, for \(1\le j\le K\),

\[
\begin{aligned}
T^jN-N
&=
\left(3^j2^{K+1-j}M-1\right)
-\left(2^{K+1}M-1\right)\\
&=
2^{K+1-j}M(3^j-2^j)\\
&>0.
\end{aligned}
\tag{7}
\]

Consequently, using (6),

\[
\begin{aligned}
V(T^jN)-V(N)
&=
\log\frac{T^jN}{N}
+\phi(-1\bmod M)-\phi(-1\bmod M)\\
&=
\log\frac{T^jN}{N}\\
&>0.
\end{aligned}
\]

This proves (3). \ensuremath{\blacksquare}


\bigskip



\subsubsection*{3. A concrete example}


Take \(M=12\) and \(K=4\). The theorem chooses

\[
N=2^5\cdot12-1=383.
\]

Its first four shortcut iterates are

\[
383\longmapsto575\longmapsto863\longmapsto1295\longmapsto1943.
\]

Every term is congruent to \(11\pmod{12}\), and every term is larger than \(383\).

Thus, regardless of the value assigned to the residue \(11\),

\[
\phi(11),
\]

the correction cancels at every candidate stopping point. Only the increasing \(\log n\) remains.


\bigskip



\subsection*{4. Strength of the obstruction}


The theorem is stronger than merely saying that a particular linear program or table search failed.

It proves all of the following.


\paragraph*{The boundedness of \(\phi\) is irrelevant}


The proof did not use boundedness. No real-valued function on the finite residue set works, no matter how large its oscillation is.


\paragraph*{The stopping rule may inspect the entire integer}


The rule \(k(n)\) need not depend only on \(n\bmod M\). It may be an arbitrary function of the full integer \(n\). The witness \(N\) defeats it because \textbf{every} \(j\in\{1,\ldots,K\}\) fails.


\paragraph*{Increasing \(M\) does not help}


For each new modulus \(M\), the construction simply chooses

\[
N=2^{K+1}M-1.
\]


\paragraph*{Increasing \(K\) does not help}


For each new horizon \(K\), the witness has \(K\) consecutive growing odd steps.


\paragraph*{Replacing \(\log n\) by another increasing height does not help}


The same proof gives the following generalization.


\begin{quote}
For every strictly increasing function \(G:\mathbb N\to\mathbb R\), no function \[
> V(n)=G(n)+\phi(n\bmod M)
> \] admits a uniformly bounded descent time.

\end{quote}


The only facts used are

\[
T^jN>N
\quad\text{and}\quad
T^jN\equiv N\pmod M.
\]

So \(G(T^jN)>G(N)\), while the finite correction remains unchanged.


\bigskip



\subsection*{5. The 2-adic explanation}


The special residue class is not accidental. In the \(2\)-adic integers,

\[
T(-1)=\frac{3(-1)+1}{2}=-1.
\]

Thus \(-1\) is a \(2\)-adic fixed point whose branch has archimedean multiplier \(3/2>1\).

The positive integer

\[
N=2^{K+1}M-1
\]

is extremely close to \(-1\) \(2\)-adically. Its first \(K\) iterates shadow the fixed point:

\[
T^jN
=
3^j2^{K+1-j}M-1
\equiv-1
\pmod{2^{K+1-j}M}.
\]

At the same time, in the ordinary real metric, the orbit grows approximately by

\[
\left(\frac32\right)^j.
\]

This is a clean adelic conflict:

\[
\text{\(2\)-adically stationary}
\qquad\text{but}\qquad
\text{archimedeanly expanding}.
\]

Bernstein and Lagarias developed the standard \(2\)-adic framework in which the shortcut Collatz map is conjugate to the \(2\)-adic shift. [\hypertarget{cit:R28:2}{}\hyperlink{ref:R28:2}{2}]

Any correction depending only on a fixed residue class sees the orbit as remaining in the same state. It cannot detect how deeply the integer lies inside the shrinking cylinder around \(-1\).


\bigskip



\subsection*{6. Translation into the accelerated valuation-word language}


Consider the odd-only map

\[
U(x)=\frac{3x+1}{2^{\nu_2(3x+1)}}.
\]

For the witness above, the first \(K\) valuations are

\[
a_0=a_1=\cdots=a_{K-1}=1.
\]

Indeed, if

\[
x_j=3^j2^{K+1-j}M-1,
\]

then

\[
3x_j+1
=
2\left(3^{j+1}2^{K-j}M-1\right),
\]

and the expression in parentheses is odd for \(j<K\). Thus the valuation word is

\[
1^K.
\]

For this word,

\[
A_j=j
\]

and the affine carry is

\[
C_{1^j}=3^j-2^j.
\]

Hence

\[
C_{1^j}+2^{A_j}
=
(3^j-2^j)+2^j
=
3^j.
\tag{8}
\]

Therefore

\[
\nu_3\!\left(C_{1^j}+2^{A_j}\right)=j.
\tag{9}
\]

So the obstruction is not merely \(2\)-adically realizable. It has \textbf{maximal terminal \(3\)-adic compatibility} as well.

This connects directly to the terminal-suffix locality result: the all-ones word survives every finite \(2\)-adic seed test and every corresponding finite \(3\)-adic endpoint test. Its inverse limit is simply the nonpositive \(2\)-adic point \(-1\), rather than a fixed ordinary positive integer.

That distinction is crucial.


\bigskip



\subsection*{7. Why this does not disprove Collatz}


The witnesses depend on the horizon:

\[
N_K=2^{K+1}M-1.
\]

There is not one fixed positive integer whose orbit follows the all-ones word forever. Rather,

\[
N_K\longrightarrow -1
\]

in the \(2\)-adic topology while

\[
N_K\longrightarrow+\infty
\]

in the ordinary topology.

Thus the theorem exhibits a sequence of longer and longer counterexample-like finite trajectories, not an infinite positive-integer counterexample.

This precisely identifies the quantifier error in the proposed strategy:

\[
\forall n\ \exists k(n)
\]

might still be true, but

\[
\exists K\ \forall n\ \exists k(n)\le K
\]

is false even after adding an arbitrary fixed periodic correction.

In fact, if \(\sigma_V(n)\) denotes the first time that

\[
V(T^j n)<V(n),
\]

then the theorem proves

\[
\sup_{\substack{n\equiv-1\pmod M\\n>1}}\sigma_V(n)=\infty.
\tag{10}
\]

For \(\phi=0\), this says that the ordinary stopping time below the starting value is unbounded even inside the single arithmetic progression

\[
n\equiv-1\pmod M.
\]


\bigskip



\subsection*{8. The obvious repair: track \(\nu_2(n+1)\)}


A natural response is to replace the finite table by an unbounded counter:

\[
V(n)=\log n+g\!\left(\nu_2(n+1)\right)+\phi(n\bmod M).
\tag{11}
\]

This can see how deeply \(n\) shadows \(-1\). Along an odd step,

\[
T(n)+1=\frac{3(n+1)}2,
\]

so

\[
\nu_2(T(n)+1)=\nu_2(n+1)-1.
\]

Unfortunately, even this enlarged ansatz still cannot have a fixed descent horizon.


\subsubsection*{Second no-go theorem}


For arbitrary \(M\), arbitrary \(\phi\), arbitrary

\[
g:\mathbb Z_{\ge0}\to\mathbb R,
\]

and arbitrary finite \(K\), there is a positive integer \(n\) such that

\[
V(T^j n)\ge V(n)
\qquad(1\le j\le K)
\]

for the function (11).


\paragraph*{The expanding cycle}


On the integers extended to negative values, the shortcut map has the cycle

\[
-10\longmapsto-5\longmapsto-7\longmapsto-10.
\tag{12}
\]

The parity pattern is

\[
0,1,1,
\]

and its multiplicative factor over one circuit is

\[
\frac12\cdot\frac32\cdot\frac32=\frac98>1.
\tag{13}
\]

The valuation pattern of \(x+1\) is

\[
\nu_2(-10+1)=0,\qquad
\nu_2(-5+1)=2,\qquad
\nu_2(-7+1)=1.
\tag{14}
\]

Let

\[
b_i=g\!\left(\nu_2(x_i+1)\right)+\phi(x_i\bmod M)
\]

for

\[
(x_0,x_1,x_2)=(-10,-5,-7).
\]

The limiting Bellman increments along the three edges are

\[
\begin{aligned}
d_0&=-\log2+b_1-b_0,\\
d_1&=\log(3/2)+b_2-b_1,\\
d_2&=\log(3/2)+b_0-b_2.
\end{aligned}
\]

Their sum is

\[
d_0+d_1+d_2
=
\log\frac98
>0.
\tag{15}
\]

For any finite cyclic sequence with positive total sum, there is a cyclic starting position for which all periodic partial sums are nonnegative: start immediately after a minimum of the cumulative sums.

Choose that phase \(x_i\). For sufficiently large \(L\), the positive integer

\[
n=x_i+2^LM
\]

shadows the cycle (12) for at least \(K\) steps. During those steps:

\begin{itemize}
\item the parity pattern is identical;
\item the residues modulo \(M\) are identical;
\item the valuations in (14) are identical;
\item each even-step logarithmic increment is exactly \(-\log2\);
\item each odd-step logarithmic increment is strictly greater than \(\log(3/2)\).
\end{itemize}

Consequently, the actual partial Bellman increments dominate the nonnegative periodic partial sums selected above. Hence no descent occurs during the first \(K\) steps.

So even the most obvious unbounded correction fixes the \(-1\) obstruction only to expose the next expanding \(2\)-adic cycle.


\bigskip



\subsection*{9. What a viable Bellman program would have to track}


For a parity word \(w\) of length \(p\), containing \(s\) odd steps, write its affine iterate as

\[
2^pT^p(n)=3^s n+C_w.
\tag{16}
\]

If

\[
3^s>2^p,
\]

then the repeated word has positive archimedean drift. Its associated periodic \(2\)-adic point is

\[
x_w=-\frac{C_w}{3^s-2^p},
\tag{17}
\]

which is negative in the ordinary real ordering.

Define the cycle defect

\[
D_w(n)=(3^s-2^p)n+C_w.
\tag{18}
\]

Then

\[
D_w(n)=0
\]

is precisely the periodic-point equation. A large value of

\[
\nu_2(D_w(n))
\]

means that the positive integer \(n\) shadows the expanding periodic word for many steps.

Moreover, whenever the block \(w\) repeats,

\[
D_w(T^p n)
=
\frac{3^s}{2^p}D_w(n).
\tag{19}
\]

Thus each repetition consumes \(p\) units of \(2\)-adic shadow depth while producing archimedean growth \(\log(3^s/2^p)\).

This suggests the correct infinite-state Bellman mechanism:


\begin{quote}
Penalize the remaining \(2\)-adic shadow depth near every expanding periodic parity word, and prove that a growing block either consumes this finite resource or violates an exact \(3\)-adic carry condition.

\end{quote}


But this requires an envelope over infinitely many dangerous words \(w\), together with a proof that the resulting potential is well-defined, proper, and descending through overlaps between different words. A single fixed modulus, or even finitely many valuation counters, cannot encode it.

The automated-termination viewpoint is philosophically related: Yolcu, Aaronson, and Heule search for ranking interpretations of Collatz rewriting systems and identify representation-dependent limitations, while obtaining certificates only for weaker systems. [\hypertarget{cit:R28:3}{}\hyperlink{ref:R28:3}{3}] Tao's theorem, by contrast, gives an almost-all descent result rather than a pointwise finite-state certificate. [\hypertarget{cit:R28:4}{}\hyperlink{ref:R28:4}{4}]


\bigskip



\subsection*{Final conclusion}


The proposed statement

\[
\exists\,M,\phi,K\quad
\forall n>1\quad
\exists\,1\le k\le K:
\quad
V(T^kn)<V(n)
\]

with

\[
V(n)=\log n+\phi(n\bmod M)
\]

is false.

The explicit universal obstruction is

\[
N=2^{K+1}M-1,
\]

for which

\[
T^jN=3^j2^{K+1-j}M-1>N
\quad\text{and}\quad
T^jN\equiv N\pmod M
\]

for every \(1\le j\le K\).

Therefore no finite table search can complete a Collatz proof in this form. The viable successor is not a larger table: it must be a \textbf{hierarchical, unbounded, cycle-defect Bellman function} capable of distinguishing a fixed positive integer from arbitrarily long positive shadows of expanding \(2\)-adic periodic points.
\subsection*{References}
\begin{enumerate}[leftmargin=2.2em]
\item \hypertarget{ref:R28:1}{}\hyperlink{cit:R28:1}{\texttt{turn967972search0}}
\item \hypertarget{ref:R28:2}{}\hyperlink{cit:R28:2}{\texttt{turn957485search1}}
\item \hypertarget{ref:R28:3}{}\hyperlink{cit:R28:3}{\texttt{turn791722search0}}
\item \hypertarget{ref:R28:4}{}\hyperlink{cit:R28:4}{\texttt{turn967972search1}}
\end{enumerate}

\clearpage
\section{Round 29: approach 24 --- Sum-of-squares Lyapunov certificates}
\subsection*{Prompt}
\begin{quote}
Complete a full proof of the Collatz conjecture using the following approach:

\begin{enumerate}
\item Sum-of-squares Lyapunov certificates
\end{enumerate}

\textbf{Speculative \textperiodcentered{} P}

Normalize each affine branch and encode a block of Collatz steps as polynomial inequalities on a real size variable and finitely many residue indicators. Use semidefinite programming to discover a sum-of-squares certificate of contraction, then rationalize and verify it exactly.

Take as long as you need.
\end{quote}
\subsection*{Response}

\subsection*{Outcome}


I pursued the SOS program in its natural finite form. It does \textbf{not} produce a valid proof of Collatz. More strongly, the proposed certificate class has a rigorous obstruction:


\begin{quote}
\textbf{No proper polynomial Lyapunov function with finitely many residue states and a bounded block length can prove the Collatz conjecture.}

\end{quote}


The obstruction is the realizable family \(n=2^L-1\), whose first \(L\) shortcut-Collatz steps are all increasing. This defeats both direct bounded-block contraction and every finite-state polynomial correction.

Current 2026 literature still describes the full conjecture as open. Tao's published theorem is an almost-all result in logarithmic density, not a pointwise theorem. [\hypertarget{cit:R29:1}{}\hyperlink{ref:R29:1}{1}, \hypertarget{cit:R29:2}{}\hyperlink{ref:R29:2}{2}] SOS Lyapunov verification and exact rational recovery are legitimate proof technologies, so the issue is not that SDP is numerically unreliable; the issue is that the desired finite certificate does not exist. [\hypertarget{cit:R29:3}{}\hyperlink{ref:R29:3}{3}, \hypertarget{cit:R29:4}{}\hyperlink{ref:R29:4}{4}]

What follows is the exact reduction, the impossibility theorem, and the stronger scale-dependent SOS target that would actually suffice.


\bigskip



\subsubsection*{1. Exact affine polynomialization of every block}


Use the shortcut Collatz map

\[
T(n)=
\begin{cases}
n/2,&n\equiv 0\pmod 2,\\[2mm]
(3n+1)/2,&n\equiv 1\pmod 2.
\end{cases}
\]

Convergence under \(T\) is equivalent to convergence under the usual Collatz map.

Fix a parity word

\[
\omega=(\epsilon_0,\ldots,\epsilon_{L-1})\in\{0,1\}^{L},
\]

where \(\epsilon_j\) records the parity before step \(j\). Put

\[
s_j=\sum_{i=0}^{j-1}\epsilon_i
\]

and define integers \(c_j\) recursively by

\[
c_0=0,\qquad
c_{j+1}=3^{\epsilon_j}c_j+\epsilon_j2^j.
\]

A direct induction gives

\[
2^jT^j(n)=3^{s_j}n+c_j
\]

whenever the first \(j\) parities are \(\epsilon_0,\ldots,\epsilon_{j-1}\). In particular,

\[
T^L(n)=\frac{3^{s_L}n+c_L}{2^L}. \tag{1}
\]

The offset can also be written explicitly as

\[
c_L=\sum_{i=0}^{L-1}
\epsilon_i\,2^i\,3^{\,s_L-s_{i+1}}. \tag{2}
\]

Because \(3^{s_L}\) is invertible modulo \(2^L\), there is a unique residue

\[
r_\omega\equiv -3^{-s_L}c_L\pmod{2^L},
\qquad 0\le r_\omega<2^L. \tag{3}
\]

This residue realizes precisely the parity word \(\omega\).

To see that the final congruence really enforces all intermediate parities, write

\[
N_j=3^{s_j}r_\omega+c_j.
\]

Then

\[
N_{j+1}=3^{\epsilon_j}N_j+\epsilon_j2^j.
\]

Starting from \(2^L\mid N_L\) and working backwards, one obtains \(2^j\mid N_j\). Furthermore,

\[
\frac{N_{j+1}}{2^j}
=
3^{\epsilon_j}\frac{N_j}{2^j}+\epsilon_j
\]

is even, which forces \(N_j/2^j\) to be even when \(\epsilon_j=0\) and odd when \(\epsilon_j=1\).

Define the endpoint of the minimal representative by

\[
d_\omega
=
\frac{3^{s_L}r_\omega+c_L}{2^L}
=
T^L(r_\omega).
\]

Every integer in this parity cylinder has the form

\[
n=r_\omega+2^Lm,\qquad m\in\mathbb Z_{\ge0},
\]

and equation (1) becomes the particularly simple affine formula

\[
T^L(n)=d_\omega+3^{s_L}m. \tag{4}
\]

Therefore

\[
n-T^L(n)
=
\bigl(2^L-3^{s_L}\bigr)m+r_\omega-d_\omega. \tag{5}
\]

This is the complete block polynomialization. Once the parity cylinder is fixed, the alleged high-dimensional Collatz inequality is just a \textbf{linear polynomial in one nonnegative integer variable}.


\bigskip



\subsubsection*{2. The exact local SOS certificate}


Set

\[
A_\omega=2^L-3^{s_L}.
\]

If \(A_\omega>0\), define

\[
m_{0,\omega}
=
\max\left(
0,\,
\left\lfloor
\frac{d_\omega-r_\omega}{A_\omega}
\right\rfloor+1
\right).
\]

Then

\[
B_\omega
=
A_\omega m_{0,\omega}+r_\omega-d_\omega
\]

is a positive integer. Writing \(z=m-m_{0,\omega}\ge0\), equation (5) becomes

\[
n-T^L(n)=A_\omega z+B_\omega. \tag{6}
\]

This is already an exact weighted SOS certificate on the half-line \(z\ge0\):

\[
A_\omega z+B_\omega
=
\underbrace{B_\omega}_{\text{constant SOS}}
+
z\underbrace{A_\omega}_{\text{constant SOS}}. \tag{7}
\]

In Gram-matrix form, the two SOS matrices are simply the rational positive semidefinite \(1\times1\) matrices

\[
[B_\omega]\succeq0,\qquad [A_\omega]\succeq0.
\]

Thus every cylinder with \(3^{s_L}<2^L\) has an exact contraction certificate after finitely many exceptional values of \(m\), and those exceptions can be checked individually.


\paragraph*{Example: a successful certificate}


Take the time-ordered parity word

\[
\omega=(1,0,0).
\]

It corresponds to \(r_\omega=5\), and

\[
T^3(5+8m)=2+3m.
\]

Hence

\[
(5+8m)-T^3(5+8m)=3+5m.
\]

An explicit weighted SOS decomposition is

\[
3+5m
=
1^2+1^2+1^2
+
m(1^2+2^2).
\]


\paragraph*{Example: an unavoidable expanding cylinder}


For

\[
\omega=(1,1,1),
\]

we have \(r_\omega=7\) and

\[
T^3(7+8m)=26+27m.
\]

Thus

\[
(7+8m)-T^3(7+8m)
=
-19(1+m). \tag{8}
\]

No SOS solver can certify this expression as positive on \(m\ge0\), because it is strictly negative.

This is not numerical failure. It is a genuine expanding branch.


\bigskip



\subsubsection*{3. Why no bounded block contraction proof exists}


The preceding bad branch is not an isolated low-depth accident.


\paragraph*{Theorem 1: bounded stopping-time descent is false}


For every \(K\ge1\), there is an integer \(n>1\) such that

\[
T^j(n)>n
\qquad\text{for every }1\le j\le K.
\]


\paragraph*{Proof}


Choose

\[
n=2^L-1
\]

with \(L\ge K\). Since

\[
T(x)+1=\frac32(x+1)
\]

whenever \(x\) is odd, induction gives

\[
T^j(2^L-1)=3^j2^{L-j}-1,
\qquad 0\le j\le L. \tag{9}
\]

For \(1\le j\le L\),

\[
T^j(2^L-1)-(2^L-1)
=
2^{L-j}(3^j-2^j)>0.
\]

Therefore the first \(L\) iterates all exceed the starting value. \ensuremath{\blacksquare}

Consequently, the statement

\[
\exists K\ \forall n>1\ \exists\,1\le j\le K:
T^j(n)<n \tag{10}
\]

is false.

The descent statement actually equivalent to Collatz is

\[
\forall n>1\ \exists j=j(n):
T^j(n)<n, \tag{11}
\]

where \(j(n)\) must be allowed to be unbounded.

A finite block SDP naturally tries to prove (10), not (11).


\bigskip



\subsubsection*{4. Finite residue potentials do not repair the problem}


One might hope to compensate for an expanding size branch by letting the Lyapunov function depend on a finite residue state. The following theorem rules that out for polynomial SOS potentials.


\paragraph*{Theorem 2: finite-state polynomial Lyapunov obstruction}


Let \(S\) be a finite set of states. For every \(s\in S\), let

\[
V_s(x)\in\mathbb R[x]
\]

satisfy

\[
V_s(x)\longrightarrow+\infty
\qquad (x\to+\infty). \tag{12}
\]

Suppose a macro-step rule assigns to every current state and every integer \(n>1\)

\begin{itemize}
\item a block length \(k\) with \(1\le k\le K\);
\item a new state \(s'\in S\);
\end{itemize}

and suppose it satisfies

\[
V_{s'}\bigl(T^k(n)\bigr)<V_s(n) \tag{13}
\]

at every macro-step.

Then no such finite collection \(\{V_s\}_{s\in S}\) exists.


\paragraph*{Proof}


Every polynomial satisfying (12) is strictly increasing for all sufficiently large \(x\). Since \(S\) is finite, choose \(N_0\) such that every \(V_s\) is strictly increasing on \([N_0,\infty)\).

Choose \(L\) so large that

\[
L>K(|S|+1)
\quad\text{and}\quad
2^L-1>N_0.
\]

Start at

\[
n_0=2^L-1.
\]

Let

\[
0=t_0<t_1<t_2<\cdots
\]

be the macro-step times. Since every macro-step has length at most \(K\),

\[
t_j\le jK.
\]

In particular, the first \(|S|+1\) macro-step endpoints occur before time \(L\). By equation (9), all corresponding orbit values are strictly increasing:

\[
T^{t_0}(n_0)<T^{t_1}(n_0)<\cdots<T^{t_{|S|}}(n_0).
\]

Among the \(|S|+1\) associated states, two must coincide. Thus, for some \(i<j\), the states at times \(t_i,t_j\) are the same state \(s\).

Chaining the strict Lyapunov inequalities gives

\[
V_s\bigl(T^{t_j}(n_0)\bigr)
<
V_s\bigl(T^{t_i}(n_0)\bigr). \tag{14}
\]

But both arguments exceed \(N_0\), and

\[
T^{t_j}(n_0)>T^{t_i}(n_0).
\]

Since \(V_s\) is strictly increasing there,

\[
V_s\bigl(T^{t_j}(n_0)\bigr)
>
V_s\bigl(T^{t_i}(n_0)\bigr), \tag{15}
\]

contradicting (14). \ensuremath{\blacksquare}


\paragraph*{Consequences}


The theorem excludes all of the following proof templates:

\[
V(n)=p_{n\bmod M}(n),
\]

with finitely many polynomial pieces;

\[
V(n)=\log n+\phi(n\bmod M),
\]

with a bounded modular correction;

finite path-complete or multiple-polynomial Lyapunov functions;

finite lexicographic tuples of proper polynomial rankings;

bounded stopping rules combined with finitely many residue indicators.

The reason properness is included is essential. A function that merely decreases but has infinite sublevel sets does not prove termination: it may decrease forever toward a limit while the integer orbit diverges. A genuine termination Lyapunov function must provide some well-foundedness, such as finite sublevel sets, an integer-valued rank, or a uniform multiplicative contradiction.

Therefore:


\begin{quote}
\textbf{Under the literal interpretation of Approach 24\textemdash{}one finite SDP, one real size variable, finitely many residue modes, and bounded blocks\textemdash{}the required certificate is mathematically impossible.}

\end{quote}



\bigskip



\subsubsection*{5. What finite-block SOS does accomplish}


At depth \(L\), the coefficient-expanding words are those satisfying

\[
3^{s_L}>2^L,
\]

or equivalently

\[
s_L\ge \left\lceil L\log_3 2\right\rceil.
\]

Their number is

\[
B_L
=
\sum_{s=\lceil L\log_3 2\rceil}^{L}
\binom{L}{s}. \tag{16}
\]

Since

\[
\log_3 2\approx0.6309297536,
\]

this is an upper binomial tail. Its proportion tends to zero exponentially. More precisely, if

\[
p=\log_3 2,
\]

then

\[
B_L
=
2^{H_2(p)L+o(L)},
\]

where

\[
H_2(p)
=
-p\log_2p-(1-p)\log_2(1-p)
\approx0.9499555.
\]

Thus the proportion of coefficient-expanding words is approximately

\[
2^{-0.0500445L+o(L)}.
\]

That is meaningful: an exact SOS program can eliminate an exponentially increasing number of residue cylinders and leave an exponentially small proportion of high-drift cylinders.

But ``exponentially small proportion'' is not ``none.'' The all-odd word survives for every \(L\), and each of its finite prefixes is realized by \(2^L-1\). This is the same exceptional-orbit gap that separates probabilistic drift arguments from a pointwise Collatz proof. Tao's theorem controls almost all starting values but deliberately does not eliminate every exceptional orbit. [\hyperlink{ref:R29:2}{2}]


\bigskip



\subsubsection*{6. An exact replacement target: the retained prefix tree}


Although the finite Lyapunov version fails, SOS can still be used locally inside an unbounded refinement argument.

First recall the elementary descent criterion.


\paragraph*{Lemma}


The Collatz conjecture is equivalent to

\[
\forall n>1\ \exists j\ge1:
T^j(n)<n. \tag{17}
\]


\paragraph*{Proof}


One direction is immediate because reaching \(1\) gives an iterate below \(n\).

Conversely, assume every \(n>1\) has an iterate below itself. Strong induction on \(n\) proves convergence: choose \(j\) with \(T^j(n)=m<n\); by induction \(m\) reaches \(1\), so \(n\) does as well. \ensuremath{\blacksquare}

Now construct parity cylinders from the binary digits of the seed.

At depth \(j\), let

\[
P_j=2^j,\qquad Q_j=3^{s_j},
\]

and let

\[
R_j\in[0,P_j)
\]

be the seed residue associated with the current parity prefix. Define

\[
X_j=T^j(R_j).
\]

A child is selected by the next binary seed bit

\[
b_j\in\{0,1\},
\]

so that

\[
R_{j+1}=R_j+b_jP_j. \tag{18}
\]

Because the first \(j\) steps are affine with multiplier \(Q_j/P_j\),

\[
T^j(R_j+b_jP_j)=X_j+b_jQ_j. \tag{19}
\]

Set

\[
Y_j=X_j+b_jQ_j,
\qquad
\epsilon_j\equiv Y_j\pmod2.
\]

Then the exact child recurrences are

\[
P_{j+1}=2P_j, \tag{20}
\]

\[
Q_{j+1}=3^{\epsilon_j}Q_j, \tag{21}
\]

\[
X_{j+1}
=
\frac{3^{\epsilon_j}Y_j+\epsilon_j}{2}. \tag{22}
\]

Since \(b_j,\epsilon_j\) are Boolean,

\[
3^{\epsilon_j}=1+2\epsilon_j,
\]

so these are polynomial equations:

\[
R'=R+bP,\qquad P'=2P,\qquad Q'=(1+2\epsilon)Q, \tag{23}
\]

\[
2X'=(1+2\epsilon)(X+bQ)+\epsilon. \tag{24}
\]

The retained condition is

\[
X_j\ge R_j. \tag{25}
\]

It means that the depth-\(j\) endpoint has not dropped below the depth-\(j\) seed representative.

A positive integer has an eventually-zero binary expansion. Hence along its prefix path,

\[
b_j=0
\]

for all sufficiently large \(j\), after which \(R_j\) is exactly the fixed starting integer.


\paragraph*{Exact survivor criterion}


The Collatz conjecture is equivalent to the following assertion:


\begin{quote}
There is no nontrivial infinite valid prefix path for which \(b_j=0\) eventually and \(X_j\ge R_j\) for all sufficiently large \(j\).

\end{quote}


Indeed, if Collatz is false, choose the least counterexample \(n\). No iterate of \(n\) can be below \(n\), because every smaller integer converges by minimality. Its binary bits are eventually zero, producing such a retained path.

Conversely, an eventually-zero retained path stabilizes at an integer \(n>1\) whose orbit is eventually always at least \(n\), so it cannot reach \(1\).

This is the correct pointwise target.


\bigskip



\subsubsection*{7. Compact normalized polynomial dynamics}


The prefix equations admit a useful normalization:

\[
\rho_j=\frac{R_j}{P_j},\qquad
y_j=\frac{X_j}{Q_j},\qquad
u_j=\frac1{P_j},\qquad
z_j=\frac1{Q_j}.
\]

Equations (18)\textendash{}(22) become

\[
\rho'=\frac{\rho+b}{2}, \tag{26}
\]

\[
y'=\frac{y+b}{2}+\frac{\epsilon z}{6}, \tag{27}
\]

\[
u'=\frac u2, \tag{28}
\]

\[
z'=z\left(1-\frac{2\epsilon}{3}\right). \tag{29}
\]

The retained inequality \(X\ge R\) becomes the polynomial inequality

\[
uy-z\rho\ge0. \tag{30}
\]

For actual prefix states,

\[
0\le \rho<1,\qquad 0\le y<1,\qquad 0<u,z\le1.
\]

Thus the normalized continuous variables live in a compact box. The Boolean equations are

\[
b^2-b=0,\qquad \epsilon^2-\epsilon=0. \tag{31}
\]

This is an excellent domain for SOS and Positivstellensatz calculations.

But there is a crucial caveat: equations (26)\textendash{}(31) alone do not encode full integrality. A finite collection of residue indicators enforces only finitely much \(2\)-adic information. Arbitrarily long all-odd prefixes continue to pass every fixed finite-depth relaxation.


\bigskip



\subsubsection*{8. What a genuinely decisive SOS certificate would need}


A successful version can no longer be one stationary finite Lyapunov function. It must be a \textbf{refinement-compatible family} of arithmetic barrier certificates.

At residue depth \(\ell\), let \(\alpha\) denote the complete finite residue state. Let \(K_{\alpha,\ell}\) be the corresponding compact semialgebraic set, including

\[
0\le\rho,y,u,z\le1,\qquad uy-z\rho\ge0,
\]

the Boolean equations, and the exact residue constraints available at level \(\ell\).

For a transition \(\alpha\to\beta\), an SOS identity would have the form

\[
B_{\alpha,\ell}(\xi)
-
B_{\beta,\ell}\!\left(F_{b,\epsilon}(\xi)\right)
=
\sigma_0(\xi)
+
\sum_i \sigma_i(\xi)g_i(\xi)
+
\sum_j q_j(\xi)h_j(\xi), \tag{32}
\]

where

\begin{itemize}
\item \(\xi=(\rho,y,u,z)\);
\item \(g_i\ge0\) define the semialgebraic region;
\item \(h_j=0\) are Boolean and residue equations;
\item each \(\sigma_i\) is SOS;
\item every coefficient and Gram matrix is rational.
\end{itemize}

However, local nonincrease is not enough. The family must also supply a well-founded conclusion such as:

\begin{enumerate}
\item a nonnegative integer-valued rank that drops on every retained \(b=0\) transition; or
\item a multiplicative barrier bounded away from zero on retained states but contracting by a rational factor \(<1\); or
\item a discrete LaSalle argument in which the common zero set of all dissipation polynomials contains no valid projective orbit; or
\item a symbolic induction in \(\ell\) proving that every retained projective path injects infinitely many \(b_j=1\).
\end{enumerate}

The fourth formulation is closest to the real arithmetic obstruction:

\[
\boxed{
\text{Every infinite retained valid path has infinitely many nonzero seed bits.}
}
\tag{33}
\]

Because positive integers have eventually zero seed bits, equation (33) would prove Collatz.

The certificate must be compatible as

\[
\ell\longrightarrow\ell+1.
\]

Merely solving larger and larger unrelated SDPs is evidence, not a proof. One needs a uniform symbolic recurrence for the Gram matrices, multipliers, or barrier polynomials and an exact induction verifying it at every refinement.


\bigskip



\subsubsection*{9. Why rationalization is not the missing step}


Once a strictly feasible numerical SOS solution is found, the usual exact pipeline is sound:

\begin{enumerate}
\item solve the SDP at high precision;
\item project the approximate Gram matrices onto the exact affine coefficient constraints;
\item perform rational reconstruction;
\item verify positive semidefiniteness exactly, for example by rational \(LDL^{\mathsf T}\);
\item expand every polynomial identity in exact arithmetic;
\item retain an explicit rational positivity margin.
\end{enumerate}

This symbolic-numeric procedure is well established. [\hyperlink{ref:R29:4}{4}]

For the elementary branch certificates above, even that machinery is unnecessary: equation (7) is already an integer SOS certificate.

The unsolved part is therefore not numerical certification. It is constructing a well-founded, refinement-compatible arithmetic barrier that eliminates the eventually-zero survivor paths.


\bigskip



\subsection*{Conclusion}


The proposed approach, interpreted as

\[
\text{finite block}
+
\text{finite residue indicators}
+
\text{proper polynomial Lyapunov function}
+
\text{one finite SDP},
\]

cannot yield a Collatz proof. The obstruction is rigorous and is supplied by the arbitrarily long increasing prefixes beginning at \(2^L-1\).

What has been established here is:

\[
\begin{aligned}
&\text{an exact affine and semialgebraic encoding of every block;}\\
&\text{an exact weighted-SOS certificate for every contracting cylinder;}\\
&\text{a proof that no bounded-horizon direct contraction certificate exists;}\\
&\text{a proof that no finite-state proper polynomial Lyapunov certificate exists;}\\
&\text{an exact projective survivor criterion equivalent to Collatz;}\\
&\text{the normalized polynomial dynamics on which a corrected SOS hierarchy must act.}
\end{aligned}
\]

The decisive missing theorem is the refinement-compatible SOS exclusion of an eventually-zero retained path. No such exact rational certificate has been produced here, so claiming a completed Collatz proof would be incorrect.
\subsection*{References}
\begin{enumerate}[leftmargin=2.2em]
\item \hypertarget{ref:R29:1}{}\hyperlink{cit:R29:1}{\texttt{turn107430search2}}
\item \hypertarget{ref:R29:2}{}\hyperlink{cit:R29:2}{\texttt{turn250353search0}}
\item \hypertarget{ref:R29:3}{}\hyperlink{cit:R29:3}{\texttt{turn970944search0}}
\item \hypertarget{ref:R29:4}{}\hyperlink{cit:R29:4}{\texttt{turn114051search1}}
\end{enumerate}

\clearpage
\section{Round 30: approach 25 --- Counterexample-guided abstraction refinement}
\subsection*{Prompt}
\begin{quote}
Complete a full proof of the Collatz conjecture using the following approach:

\begin{enumerate}
\item Counterexample-guided abstraction refinement
\end{enumerate}

\textbf{Hybrid \textperiodcentered{} P}

Begin with a coarse finite abstraction recording a modulus and height band. Ask a model checker for a nonterminating abstract path. If the path is spurious, add the residue or carry information responsible for the failure.

Termination of the refinement process with no recurrent bad component would produce a finite proof certificate. Nontermination might reveal the exact state variables missing from present approaches.

Take as long as you need.
\end{quote}
\subsection*{Response}

\subsection*{Result of the attack}


The CEGAR idea can be made rigorous, but it does \textbf{not} presently yield a full proof. More sharply, the literal target in the proposal\textemdash{}


\begin{quote}
refine a finite one-step abstraction until the nonterminal abstract graph has no recurrent component

\end{quote}


\textemdash{}is impossible. Collatz has arbitrarily long finite excursions, so every sound finite one-step abstraction necessarily contains recurrent abstract behavior, regardless of whether the conjecture is true.

The viable replacement is \textbf{ranked CEGAR}: recurrent abstract components are permitted, but each must be discharged by a well-founded ranking relation, an exact acceleration, or an arithmetic infeasibility certificate.

This leads to several exact results:

\begin{enumerate}
\item a no-go theorem for pure finite-state refinement;
\item an exact arithmetic oracle for checking proposed counterexample paths;
\item an exact ``top-bit'' criterion distinguishing ordinary integers from spurious \(2\)-adic paths;
\item your terminal-suffix locality theorem as a natural \(3\)-adic refinement rule;
\item a universal ranking function that eliminates every exact periodic abstract lasso unless it is a genuine Collatz cycle.
\end{enumerate}

What remains is a precise finite-rank theorem for \textbf{aperiodic} survivor paths. That theorem is not established below, so it would be incorrect to label the result a proof. Current peer-reviewed literature still describes Collatz as unproved, while published exhaustive verification has reached \(2^{71}\). [\hypertarget{cit:R30:1}{}\hyperlink{ref:R30:1}{1}, \hypertarget{cit:R30:2}{}\hyperlink{ref:R30:2}{2}]


\bigskip



\subsubsection*{1. Turn Collatz into a termination problem}


Use the shortcut map

\[
T(x)=
\begin{cases}
x/2,&x\equiv0\pmod 2,\\[2mm]
(3x+1)/2,&x\equiv1\pmod 2.
\end{cases}
\]

The following reduced formulation is useful.


\paragraph*{Proposition 1: reduced descent}


The Collatz conjecture is equivalent to

\[
\forall n>1\quad \exists k\geq1\quad T^k(n)<n.
\tag{1}
\]


\paragraph*{Proof}


If the Collatz conjecture is true, the orbit eventually reaches \(1<n\), proving (1).

Conversely, assume (1). Use strong induction on \(n\). Given \(n>1\), choose \(k\) with \(T^k(n)=m<n\). By the induction hypothesis, the orbit of \(m\) reaches \(1\). Hence the orbit of \(n\) reaches \(1\). \ensuremath{\blacksquare}

Thus we may study the parameterized loop

\[
x:=n;\qquad
\textbf{while }x\geq n\textbf{ do }x:=T(x).
\tag{2}
\]

A least counterexample would give a nonterminating execution of (2).


\bigskip



\subsection*{2. Pure finite-state CEGAR cannot work}


Here is a structural obstruction to the proposal as originally phrased.


\subsubsection*{Theorem 2: finite acyclic abstraction is impossible}


Consider any sound finite abstraction of the one-step transition system (2). If every reachable nonterminal abstract component were acyclic, there would be a uniform bound on the number of steps before descent. Such a bound does not exist.


\paragraph*{Proof}


For \(r\geq1\), let

\[
n_r=2^r-1.
\]

For \(0\leq i\leq r\),

\[
T^i(n_r)=3^i2^{r-i}-1.
\tag{3}
\]

Indeed, all of the first \(r\) terms before the final value in (3) are odd, and

\[
T\!\left(3^i2^{r-i}-1\right)
=
\frac{3(3^i2^{r-i}-1)+1}{2}
=
3^{i+1}2^{r-i-1}-1.
\]

Moreover,

\[
3^i2^{r-i}-1
\geq 2^r-1=n_r
\]

because \((3/2)^i\geq1\). Therefore the loop (2), started at \(n_r\), remains active for at least \(r\) transitions.

Now suppose the finite abstraction has \(N\) nonterminal abstract states and no nonterminal directed cycle. Every nonterminal abstract path then has length at most \(N-1\). Taking \(r>N\) contradicts soundness. \ensuremath{\blacksquare}

There is an even simpler version for total stopping time: the orbit of \(2^r\) requires \(r\) divisions by two before reaching \(1\).


\paragraph*{Consequence}


No refinement process that merely splits a finite collection of modulus and height states can terminate with an acyclic bad graph. At every finite refinement level, sufficiently long concrete finite trajectories force an abstract cycle.

The correct formal-methods analogue is therefore not ordinary safety CEGAR, but \textbf{termination CEGAR using ranking relations or transition invariants}. This distinction is standard in the termination literature: ordinary CEGAR refines spurious counterexamples, while termination refinement must also learn well-founded relations that discharge recurrent abstract paths. [\hypertarget{cit:R30:3}{}\hyperlink{ref:R30:3}{3}, \hypertarget{cit:R30:4}{}\hyperlink{ref:R30:4}{4}, \hypertarget{cit:R30:5}{}\hyperlink{ref:R30:5}{5}]


\bigskip



\subsection*{3. The correct certificate theorem}


Let \(S\) be the set of reachable nonterminal states of a transition system, and write \(s\to^+t\) when \(t\) is reached from \(s\) in one or more steps.

A binary relation \(W\subseteq S\times S\) is well-founded if there is no infinite chain

\[
s_0\,W\,s_1\,W\,s_2\,W\,\cdots.
\]

Usually \(W\) is certified by a rank

\[
\rho:S\longrightarrow\mathbb N,
\qquad
sWt\Longrightarrow \rho(t)<\rho(s).
\]


\subsubsection*{Theorem 3: finite transition-invariant certificate}


Suppose there are finitely many well-founded relations

\[
W_1,\ldots,W_m
\]

such that every reachable transitive transition is covered:

\[
\to^+\ \subseteq\ W_1\cup\cdots\cup W_m.
\tag{4}
\]

Then the transition system terminates.


\paragraph*{Proof}


Assume an infinite execution

\[
s_0\to s_1\to s_2\to\cdots.
\]

For each \(i<j\), condition (4) places \((s_i,s_j)\) in at least one \(W_\ell\). Color the pair \(\{i,j\}\) by one such \(\ell\). Infinite Ramsey gives indices

\[
i_0<i_1<i_2<\cdots
\]

and a single color \(\ell\) such that

\[
s_{i_p}\,W_\ell\,s_{i_q}
\qquad(p<q).
\]

In particular,

\[
s_{i_0}\,W_\ell\,s_{i_1}\,W_\ell\,s_{i_2}\,W_\ell\cdots,
\]

contradicting well-foundedness of \(W_\ell\). \ensuremath{\blacksquare}

This is the correct endpoint for Collatz CEGAR:


\begin{quote}
Produce a finite arithmetic abstraction together with a finite family of ranks whose union covers the transitive behavior of every recurrent bad component.

\end{quote}


Merely showing that each abstract simple cycle has some locally decreasing quantity is not enough. The quantity may increase while the execution travels through other cycles. Coverage of the transitive closure, or an equivalent lexicographic/global ranking argument, is essential.


\bigskip



\subsection*{4. Exact arithmetic for counterexample concretization}


It is most convenient to work with the accelerated odd map

\[
U(x)=\frac{3x+1}{2^{\nu_2(3x+1)}},
\qquad x\text{ odd}.
\]

Let

\[
x_0=n,\qquad x_{t+1}=U(x_t),
\]

and define

\[
a_t=\nu_2(3x_t+1),\qquad
A_t=\sum_{i<t}a_i.
\]

Set

\[
C_0=0,\qquad
C_{t+1}=3C_t+2^{A_t}.
\tag{5}
\]


\subsubsection*{Lemma 4: affine word identity}


For every \(t\geq0\),

\[
2^{A_t}x_t=3^t n+C_t.
\tag{6}
\]


\paragraph*{Proof}


It is true at \(t=0\). If it holds at \(t\), then

\[
\begin{aligned}
2^{A_t+a_t}x_{t+1}
&=2^{A_t}(3x_t+1)\\
&=3(3^tn+C_t)+2^{A_t}\\
&=3^{t+1}n+C_{t+1}.
\end{aligned}
\]

\ensuremath{\blacksquare}

For a finite valuation word

\[
w=(a_0,\ldots,a_{j-1}),
\]

write \(A_i,C_i\) for its prefix data.

The barrier requirement \(x_i\geq n\) is, by (6),

\[
3^i n+C_i\geq 2^{A_i}n,
\]

or

\[
(2^{A_i}-3^i)n\leq C_i.
\tag{7}
\]

Thus:

\begin{itemize}
\item when \(2^{A_i}<3^i\), the \(i\)-th barrier inequality is automatic;
\item when \(2^{A_i}>3^i\), it gives the exact upper bound
\end{itemize}

\[
n\leq
\left\lfloor
\frac{C_i}{2^{A_i}-3^i}
\right\rfloor.
\tag{8}
\]

No logarithms or floating-point height bands are needed in the final certificate.


\bigskip



\subsubsection*{4.1 The exact seed cylinder}


The accelerated valuation word \(w\) corresponds, under the shortcut map, to the parity prefix

\[
1\,0^{a_0-1}\,1\,0^{a_1-1}\cdots
1\,0^{a_{j-1}-1}\,1,
\tag{9}
\]

whose length is \(A_j+1\).

The first \(k\) shortcut parity bits determine a unique residue modulo \(2^k\). This follows inductively because every parity cylinder modulo \(2^k\) splits into its two lifts modulo \(2^{k+1}\), exactly one of which gives each possible next parity.

Consequently there is a unique odd residue

\[
R_w\pmod {2^{A_j+1}}
\tag{10}
\]

whose first \(j\) accelerated valuations are exactly \(w\).

Let \(R_w\) also denote its least positive representative. Every positive seed realizing \(w\) has the form

\[
n=R_w+2^{A_j+1}q,
\qquad q\geq0.
\tag{11}
\]


\bigskip



\subsubsection*{Theorem 5: exact finite-path feasibility oracle}


Fix a verified lower bound \(B\), so a least counterexample would satisfy \(n>B\). Define

\[
H(w)=
\min_{\substack{1\leq i\leq j\\2^{A_i}>3^i}}
\left\lfloor\frac{C_i}{2^{A_i}-3^i}\right\rfloor,
\tag{12}
\]

with \(H(w)=+\infty\) if there are no such prefixes.

Let

\[
N_B(w)
=
R_w+
2^{A_j+1}
\max\!\left(
0,
\left\lceil
\frac{B+1-R_w}{2^{A_j+1}}
\right\rceil
\right).
\tag{13}
\]

Then there exists an odd integer \(n>B\) realizing \(w\) and satisfying

\[
U^i(n)\geq n
\qquad(1\leq i\leq j)
\]

if and only if

\[
N_B(w)\leq H(w).
\tag{14}
\]


\paragraph*{Proof}


All realizations of \(w\) are precisely the integers in (11). The smallest one exceeding \(B\) is \(N_B(w)\).

Every nonautomatic barrier inequality is an upper bound of the form (8). Since these inequalities become harder as \(n\) increases, some member of the cylinder satisfies all of them exactly when its least admissible member \(N_B(w)\) does. This is condition (14). \ensuremath{\blacksquare}

This theorem is an exact CEGAR concretization oracle. A purported finite counterexample path needs no heuristic SMT reasoning: it reduces to one congruence class and finitely many integer inequalities.


\bigskip



\subsubsection*{4.2 Why the all-ones path survives every finite depth}


For

\[
w=1^j=(1,\ldots,1),
\]

we have

\[
A_j=j,\qquad
C_j=3^j-2^j.
\]

Hence

\[
2^{A_i}-3^i=2^i-3^i<0
\]

for every prefix, so every barrier condition is automatic.

The exact seed cylinder is

\[
n\equiv -1\pmod {2^{j+1}},
\]

whose least positive member is

\[
R_{1^j}=2^{j+1}-1.
\]

Thus every finite all-ones path is a genuine growing path for some positive integer. But the infinite intersection of these cylinders is the \(2\)-adic integer

\[
n=-1\in\mathbb Z_2,
\]

not a positive ordinary integer.

This is the canonical first spurious CEGAR counterexample:

\[
1^\omega.
\]

No amount of finite unrolling refutes it. The abstraction must learn either a top-bit finiteness condition or an appropriate ranking function.


\bigskip



\subsection*{5. The exact top-bit refinement}


It is convenient to use the odd seed coordinate

\[
z=\frac{n-1}{2}.
\]

For a word \(w\) with total exponent \(A=A_j\), set

\[
\eta_w=\frac{R_w-1}{2},
\qquad 0\leq\eta_w<2^A.
\]

Then

\[
z=\eta_w+2^Aq_w.
\tag{15}
\]

Suppose the next valuation is \(a\). The extended residue has the unique form

\[
\eta_{wa}
=
\eta_w+b_w(a)2^A,
\qquad
0\leq b_w(a)<2^a.
\tag{16}
\]

The integer \(b_w(a)\) is the block of newly exposed high seed bits.


\subsubsection*{Lemma 6: quotient countdown}


Along the valuation word of a fixed ordinary seed,

\[
q_{wa}
=
\frac{q_w-b_w(a)}{2^a}
=
\left\lfloor\frac{q_w}{2^a}\right\rfloor.
\tag{17}
\]

In particular, if \(q_w>0\), then \(q_{wa}<q_w\).


\paragraph*{Proof}


Using (15) and (16),

\[
\eta_w+2^Aq_w
=
\eta_w+b_w(a)2^A+2^{A+a}q_{wa}.
\]

Canceling \(\eta_w\) and dividing by \(2^A\) gives

\[
q_w=b_w(a)+2^a q_{wa}.
\]

Since \(0\leq b_w(a)<2^a\), it is the residue of \(q_w\) modulo \(2^a\), and (17) follows. \ensuremath{\blacksquare}

This yields an exact characterization.


\subsubsection*{Corollary 7: ordinary-seed criterion}


An infinite accelerated valuation word corresponds to a positive ordinary integer seed if and only if its successive top-bit blocks \(b_w(a)\) are eventually zero.


\paragraph*{Proof}


For a fixed finite \(z\), once \(2^A>z\), equation (15) has \(q_w=0\), so every later extension has \(b_w(a)=0\).

Conversely, if the appended blocks are eventually zero, the residues \(\eta_w\) eventually stabilize to a finite nonnegative integer \(z\). The corresponding seed is \(n=2z+1\). \ensuremath{\blacksquare}

Thus ``the seed is an ordinary integer'' is naturally a \textbf{co-B\"uchi condition}:

\[
b_j=0\quad\text{for all sufficiently large }j.
\tag{18}
\]

It should not be approximated by a fixed modulus.


\bigskip



\subsubsection*{5.1 Two-ended synchronization}


Let

\[
X_w=U^j(R_w)
\]

be the endpoint corresponding to the least seed representative. For a general seed in the same cylinder,

\[
n=R_w+2^{A_j+1}q,
\]

the affine identity gives

\[
U^j(n)=X_w+2\cdot3^j q.
\tag{19}
\]

So the same quotient \(q\) measures both:

\begin{itemize}
\item the unprocessed high binary seed digits;
\item the displacement of the endpoint in its \(3^j\)-spaced family.
\end{itemize}

This is the exact seed/endpoint synchronization that a CEGAR state should expose.


\bigskip



\subsection*{6. Your terminal-suffix theorem is the natural \(3\)-adic interpolant}


Let \(v\) be the final \(r\) valuations of a word \(w\). Let \(Y\) be the odd value immediately before this suffix and \(X\) the final endpoint. Applying the affine identity to \(v\) alone gives

\[
2^{A(v)}X=3^rY+C_v.
\tag{20}
\]

Reducing modulo \(3^r\),

\[
X\equiv 2^{-A(v)}C_v\pmod {3^r}.
\tag{21}
\]

Therefore

\[
3^r\mid X+1
\quad\Longleftrightarrow\quad
3^r\mid C_v+2^{A(v)}.
\tag{22}
\]

This is exactly the refinement predicate needed when an abstract path uses an impossible endpoint carry: increase the suffix depth only where the spurious lasso requires it.

It is important, however, that (22) is a compatibility theorem, not a descent theorem. It makes the \(3\)-adic side finite and local for each fixed \(r\), but it does not itself eliminate the archimedean survivor component.


\bigskip



\subsection*{7. Every exact periodic lasso has a canonical rank}


This is the strongest reusable CEGAR lemma.

For a valuation word \(w\) of length \(j\), total exponent \(A\), and affine constant \(C_w\), write

\[
F_w(x)=\frac{3^j x+C_w}{2^A}.
\tag{23}
\]

Define

\[
D_w=2^A-3^j
\]

and the affine fixed-point defect

\[
\Delta_w(x)=D_wx-C_w.
\tag{24}
\]

The fixed-point equation \(F_w(x)=x\) is precisely

\[
\Delta_w(x)=0,
\]

or

\[
(2^A-3^j)x=C_w.
\tag{25}
\]


\subsubsection*{Theorem 8: affine-lasso elimination}


Suppose \(w\) is applicable at \(x\), and put \(y=F_w(x)\). Then

\[
\Delta_w(y)
=
\frac{3^j}{2^A}\Delta_w(x).
\tag{26}
\]

Consequently, if \(\Delta_w(x)\neq0\) and \(w\) is applicable repeatedly, then on each traversal

\[
\nu_2\!\left(\Delta_w(F_w(x))\right)
=
\nu_2\!\left(\Delta_w(x)\right)-A.
\tag{27}
\]

Hence \(w\) cannot repeat indefinitely on an integer orbit unless

\[
\Delta_w(x)=0.
\]

In the exceptional case \(\Delta_w(x)=0\), \(x\) is an actual periodic point for the block \(w\).


\paragraph*{Proof}


Using \(D_w=2^A-3^j\),

\[
\begin{aligned}
\Delta_w(F_w(x))
&=
D_w\frac{3^jx+C_w}{2^A}-C_w\\
&=
\frac{3^jD_wx+D_wC_w-2^AC_w}{2^A}\\
&=
\frac{3^jD_wx-3^jC_w}{2^A}\\
&=
\frac{3^j}{2^A}\Delta_w(x).
\end{aligned}
\]

Both \(\Delta_w(x)\) and \(\Delta_w(F_w(x))\) are integers. Since \(3^j\) is odd, equation (26) implies that \(2^A\mid\Delta_w(x)\), and, when the defect is nonzero, its \(2\)-adic valuation falls by exactly \(A\). Such a nonnegative valuation cannot fall indefinitely.

If the defect is zero, (25) gives \(F_w(x)=x\). When the valuation guards defining \(w\) are satisfied, this is a genuine Collatz cycle. \ensuremath{\blacksquare}

Thus the model checker can assign the rank

\[
\rho_w(x)=
\nu_2\!\left|
(2^A-3^j)x-C_w
\right|
\tag{28}
\]

to the lasso \(w^\omega\).


\paragraph*{Examples}


For the one-letter word \(w=(1)\),

\[
A=1,\qquad C_w=1,\qquad D_w=-1,
\]

so

\[
\Delta_w(x)=-(x+1).
\]

The abstract lasso \(1^\omega\) has \(2\)-adic fixed point \(-1\), and the rank is

\[
\rho(x)=\nu_2(x+1).
\]

This is exactly the countdown hidden by the coarse abstraction.

For \(w=(2)\),

\[
A=2,\qquad C_w=1,\qquad D_w=1,
\]

and the fixed point is

\[
x=1,
\]

the trivial accelerated Collatz cycle.

For a general word, a purported exact periodic lasso has only two outcomes:

\[
\boxed{\text{spurious and ranked}}
\qquad\text{or}\qquad
\boxed{\text{an explicit integer cycle}}.
\]

In particular, every eventually periodic exact valuation sequence arising from an ordinary integer is eventually an actual cycle. A divergent counterexample, if one exists and is not eventually cyclic, must therefore have an aperiodic exact valuation word.


\bigskip



\subsection*{8. The ranked CEGAR algorithm}


An exact implementation can use the following loop.


\begin{verbatim}
Predicates :=
    parity / valuation branch information
    candidate lower bound n > B
    barrier condition x >= n
    coarse exact rational height bands
    a small 2-adic seed modulus
    a small terminal-suffix 3-adic state

Ranks := empty

repeat:
    Build the finite over-approximating transition system.

    Ask for an accepting recurrent path satisfying:
        no descent,
        no entry into the verified region,
        eventual-zero top-bit condition.

    If no accepting recurrent path exists:
        return the finite abstraction, ranks, and arithmetic
        obligations as a proof certificate.

    Extract an abstract lasso u v^omega.

    Concretize finite unrollings using Theorem 5.

    If an unrolling is arithmetically infeasible:
        refine by the first separating condition:
            a deeper 2-adic residue,
            an exact barrier inequality,
            or a terminal-suffix carry predicate.

    Else if the loop v has become one exact valuation block:
        compute A(v), C_v, and D_v.

        If D_v x = C_v has a positive integral solution
        satisfying all valuation guards:
            return an explicit nontrivial Collatz cycle.

        Otherwise:
            attach the affine-defect rank
            nu_2(D_v x - C_v)
            from Theorem 8.

    Else:
        expose more of:
            the emitted top-bit block,
            the suffix carry,
            the affine fixed-point defect,
            or the exact height quotient.
\end{verbatim}


Every successful termination of this procedure is independently checkable using:

\begin{itemize}
\item arbitrary-precision integer arithmetic;
\item finite graph reachability and SCC checking;
\item exact congruences;
\item the inequalities in Theorem 5;
\item explicit well-founded rank decreases.
\end{itemize}

No numerical approximation is needed in the final certificate.

The closest established Collatz work from automated reasoning encodes the conjecture as termination of a mixed binary\textendash{}ternary rewriting system and automatically proves some weakenings. Its authors explicitly did not obtain Collatz termination and also demonstrated limitations of standard matrix-interpretation proof domains. [\hypertarget{cit:R30:6}{}\hyperlink{ref:R30:6}{6}]


\bigskip



\subsection*{9. Where the proof attempt stops}


The affine-lasso theorem handles every \textbf{fixed periodic word}. It does not handle an infinite path that continually changes its blocks.

Suppose an abstract recurrent SCC admits words

\[
w_1,w_2,w_3,\ldots
\]

with associated defects

\[
\Delta_{w_i}(x)
=
(2^{A_i}-3^{j_i})x-C_{w_i}.
\]

The rank for \(w_i\) falls while \(w_i\) repeats, but after switching to \(w_{i+1}\), the relevant affine form changes. Its valuation may reset to an arbitrarily large value. Therefore the collection of individual lasso ranks does not automatically give a finite transition invariant.

This is exactly where a divergent, aperiodic counterexample could hide.

Your terminal-suffix locality result solves the finite-window \(3\)-adic part of this problem. The seed quotient solves the fake-\(2\)-adic-integer part. The barrier oracle solves finite archimedean feasibility. But no argument presently shows that the resulting infinite collection of affine defects collapses to finitely many rank types.


\bigskip



\subsubsection*{The exact missing theorem}


A full CEGAR proof would follow from the following statement.


\paragraph*{Finite rank-cover theorem}


There exist:

\begin{itemize}
\item a finite abstraction of the lifted Collatz state;
\item finitely many arithmetic ranks \(\rho_1,\ldots,\rho_m\);
\item finitely many transition-summary relations \(W_1,\ldots,W_m\);
\end{itemize}

such that:

\begin{enumerate}
\item every \(\rho_i\) is well-founded on \(W_i\);
\item every pair of states on every candidate barrier orbit belongs to at least one \(W_i\);
\item the only zero-defect recurrent component is the trivial cycle at \(1\).
\end{enumerate}

Equivalently, for every infinite valuation word satisfying simultaneously

\[
\begin{aligned}
&\text{exact accelerated valuation realizability},\\
&b_j=0\quad\text{eventually},\\
&(2^{A_j}-3^j)n\leq C_j\quad\text{for every }j,\\
&\text{all terminal-suffix carry conditions},
\end{aligned}
\tag{29}
\]

one fixed rank from a finite family must decrease along an infinite subsequence.

Theorem 3 would then exclude such a word and prove Collatz.

But proving this finite rank-cover theorem is not a routine model-checking step. It is the unresolved mathematical core. In effect, it asks for a finite arithmetic reason why an aperiodic, exactly realizable, pointwise critical valuation word cannot exist.


\bigskip



\subsection*{10. Final verdict}


This investigation produces a useful but not complete result:

\begin{itemize}
\item \textbf{Pure finite-state refinement is ruled out.} Arbitrarily long finite Collatz excursions force recurrent abstract paths at every finite level.
\item \textbf{Ranked CEGAR is sound.} A finite transition-invariant certificate would prove Collatz.
\item \textbf{Finite paths are exactly decidable.} The seed cylinder and barrier inequalities give a complete concretization oracle.
\item \textbf{The ordinary-integer condition is exact.} It is eventual vanishing of the emitted top-bit blocks.
\item \textbf{The \(3\)-adic endpoint condition is local.} Your suffix theorem supplies the correct carry refinement.
\item \textbf{Every exact periodic lasso is completely understood.} It either yields a genuine cycle or is eliminated by the affine-defect rank \[
  \nu_2\!\left((2^A-3^j)x-C_w\right).
  \]
\item \textbf{The remaining obstruction is necessarily aperiodic.} The missing result is a finite rank cover for such survivor paths.
\end{itemize}

So approach 25 does not currently complete a proof, but it can be sharpened substantially: the first genuinely missing state variables are not merely ``a larger modulus.'' They are the parametric affine fixed-point defects

\[
(2^A-3^j)x-C_w,
\]

their \(2\)- and \(3\)-adic valuations, and a mechanism showing that only finitely many defect types can recur on a critical ordinary-integer orbit. That is the precise theorem the next refinement experiment should target.
\subsection*{References}
\begin{enumerate}[leftmargin=2.2em]
\item \hypertarget{ref:R30:1}{}\hyperlink{cit:R30:1}{\texttt{turn677966search16}}
\item \hypertarget{ref:R30:2}{}\hyperlink{cit:R30:2}{\texttt{turn511265search7}}
\item \hypertarget{ref:R30:3}{}\hyperlink{cit:R30:3}{\texttt{turn634146search2}}
\item \hypertarget{ref:R30:4}{}\hyperlink{cit:R30:4}{\texttt{turn634146search14}}
\item \hypertarget{ref:R30:5}{}\hyperlink{cit:R30:5}{\texttt{turn640317search8}}
\item \hypertarget{ref:R30:6}{}\hyperlink{cit:R30:6}{\texttt{turn634146search0}}
\end{enumerate}

\clearpage
\section{Round 31: approach 26 --- A genuinely new reduction order for mixed binary-ternary rewriting}
\subsection*{Prompt}
\begin{quote}
Complete a full proof of the Collatz conjecture using the following approach:

\begin{enumerate}
\item A genuinely new reduction order for mixed binary-ternary rewriting
\end{enumerate}

\textbf{Established/Hybrid \textperiodcentered{} P}

The rewriting formulation is equivalent to Collatz, but standard natural and arctic matrix interpretations do not suffice. Search for context-sensitive polynomial orders, dependency-pair orders with semantic labels, or orders taking values in ordered noncommutative series.

Take as long as you need.
\end{quote}
\subsection*{Response}

\subsection*{Outcome}


I do \textbf{not} obtain a complete proof of the Collatz conjecture from this approach. The attempt produces several rigorous reductions and two fairly strong impossibility theorems, but the genuinely decisive statement\textemdash{}a global, carry-sensitive size-change lemma\textemdash{}remains unproved. Presenting the construction without that lemma as a completed proof would be incorrect.

The strongest rigorous result is:


\begin{quote}
\textbf{Any successful reduction order for the mixed binary\textendash{}ternary system must be genuinely context-sensitive and path-dependent.}

Direct scalar polynomial interpretations fail even weakly; polynomial interpretations still fail after adding the exact unbounded label \(r=v_2(n+1)\); and any ordinary context-cancellable monoid order runs into an immediate infinite-descent contradiction.

\end{quote}


What survives is a precise dependency-pair/semantic-label certificate whose verification would prove Collatz.


\bigskip



\subsubsection*{1. The exact mixed binary\textendash{}ternary rewriting problem}


Use the alphabet

\[
\Sigma=\{\lhd,\rhd,f,t,0,1,2\}.
\]

The digit symbols act on positive integers by

\[
f(x)=2x,\qquad t(x)=2x+1,
\]

\[
0(x)=3x,\qquad 1(x)=3x+1,\qquad 2(x)=3x+2.
\]

Thus a canonical word

\[
\lhd s_1s_2\cdots s_k\rhd
\]

represents the integer obtained by starting with \(1\) and applying \(s_1,\ldots,s_k\) successively.

The rewriting system \(\mathcal T\) consists of the following eleven rules.


\paragraph*{Dynamic rules}


\[
f\rhd\longrightarrow\rhd ,
\qquad
t\rhd\longrightarrow2\rhd .
\tag{D}
\]

These implement the shortcut Collatz map

\[
T(n)=
\begin{cases}
n/2,&n\equiv0\pmod2,\\[2mm]
(3n+1)/2,&n\equiv1\pmod2.
\end{cases}
\]


\paragraph*{Interior base-conversion rules}


\[
\begin{array}{lll}
f0\longrightarrow0f,&
f1\longrightarrow0t,&
f2\longrightarrow1f,\\[1mm]
t0\longrightarrow1t,&
t1\longrightarrow2f,&
t2\longrightarrow2t.
\end{array}
\tag{A}
\]


\paragraph*{Left-boundary rules}


\[
\lhd0\longrightarrow\lhd t,\qquad
\lhd1\longrightarrow\lhd ff,\qquad
\lhd2\longrightarrow\lhd ft.
\tag{B}
\]

Every rule in \(A\cup B\) preserves the represented integer. The two rules in \(D\) perform one shortcut Collatz step. Yolcu, Aaronson, and Heule proved that termination of this system is equivalent to the Collatz conjecture. [\hypertarget{cit:R31:1}{}\hyperlink{ref:R31:1}{1}, \hypertarget{cit:R31:2}{}\hyperlink{ref:R31:2}{2}]

There is also a direct proof that the auxiliary subsystem \(A\cup B\) terminates. For a word \(w\), let

\[
N_3(w)=\#\{\text{symbols }0,1,2\text{ in }w\}
\]

and let \(I(w)\) count pairs consisting of a binary digit \(f\) or \(t\) lying to the left of a ternary digit \(0,1,2\).

Every \(A\)-rule preserves \(N_3\) and decreases \(I\) by exactly one. Every \(B\)-rule decreases \(N_3\). Hence

\[
\bigl(N_3(w),I(w)\bigr)
\]

strictly decreases in lexicographic order under every auxiliary rewrite.

Consequently, an infinite \(\mathcal T\)-derivation would have to use the dynamic rules infinitely many times. Between two dynamic rules only finitely many value-preserving conversions can occur, so the values at successive dynamic steps form a Collatz trajectory.


\bigskip



\subsubsection*{2. Why an ordinary reduction order cannot work}


Here is an obstruction independent of polynomial or matrix interpretations.

Suppose \(\succ\) were a total, context-monotone, well-founded word order orienting all rules. From

\[
t\rhd\succ2\rhd
\]

one obtains, by order reflection through the common right context,

\[
t\succ2.
\]

From

\[
\lhd2\succ\lhd ft
\]

one similarly obtains

\[
2\succ ft.
\]

Therefore

\[
t\succ ft.
\]

Context monotonicity now gives

\[
t\succ ft\succ f^2t\succ f^3t\succ\cdots,
\]

an infinite descending chain.

Thus:


\begin{quote}
\textbf{Proposition 1.} No total context-cancellable reduction order can strictly orient the mixed-base system.

\end{quote}


Even a relative order that makes \(t\rhd\to2\rhd\) strict and \(\lhd2\to\lhd ft\) merely weak encounters the same contradiction:

\[
t\succ2\succeq ft
\quad\Longrightarrow\quad
t\succ ft\succ f^2t\succ\cdots.
\]

This explains why context sensitivity is not an optional refinement. The left and right boundary contexts must materially affect comparisons.

The published general impossibility theorem in the Yolcu\textendash{}Aaronson\textendash{}Heule paper concerns the older unary Collatz rewriting system, including dependency-pair variants. For the mixed-base system, their automated searches established many subsystem results but did not prove a universal impossibility theorem for matrix interpretations. [\hypertarget{cit:R31:3}{}\hyperlink{ref:R31:3}{3}]


\bigskip



\subsubsection*{3. Direct scalar polynomial interpretations collapse completely}


The following gives a direct impossibility theorem for the mixed system itself.


\paragraph*{Theorem 2}


Let every symbol \(s\in\Sigma\) be interpreted by a strictly increasing polynomial

\[
P_s(x)\in\mathbb N[x],
\]

with nonnegative coefficients. Interpret concatenation by composition:

\[
[P_{ab}]=P_b\circ P_a.
\]

Suppose every rule of \(\mathcal T\) is oriented weakly:

\[
[\ell](x)\ge [r](x)
\qquad\text{for every rule }\ell\to r
\]

and every \(x\in\mathbb N\). Then

\[
P_f(x)=P_t(x)=P_0(x)=P_1(x)=P_2(x)=x.
\]

In particular, no rule can be oriented strictly.


\paragraph*{Proof}


Write

\[
d_s=\deg P_s\ge1.
\]

Degree comparison in the relevant rules gives

\[
d_t\ge d_2
\tag{1}
\]

from \(t\rhd\to2\rhd\), and

\[
d_0\ge d_t,\qquad
d_1\ge d_f^2,\qquad
d_2\ge d_fd_t
\tag{2}
\]

from the three boundary rules.

The rules \(f2\to1f\), \(t0\to1t\), and \(f1\to0t\) give

\[
d_2\ge d_1,\qquad
d_0\ge d_1,\qquad
d_1d_f\ge d_td_0.
\tag{3}
\]

From (1) and the last inequality in (2),

\[
d_t\ge d_2\ge d_fd_t.
\]

All degrees are positive integers, so

\[
d_f=1,\qquad d_2=d_t.
\]

The last inequality in (3) now becomes

\[
d_1\ge d_td_0.
\]

Since \(d_0\ge d_t\),

\[
d_1\ge d_t^2.
\]

But \(d_t=d_2\ge d_1\). Hence

\[
d_t\ge d_t^2,
\]

so \(d_t=1\). It follows successively that

\[
d_0=d_1=d_2=d_f=d_t=1.
\]

All five digit interpretations are therefore affine. Write

\[
P_s(x)=a_sx+b_s,\qquad a_s\ge1,\quad b_s\ge0.
\]

Comparing leading coefficients in exactly the same rules gives

\[
a_t\ge a_2\ge a_fa_t,
\]

\[
a_0\ge a_t,\qquad
a_1\ge a_f^2,\qquad
a_2\ge a_1,
\]

\[
a_1a_f\ge a_ta_0.
\]

The same argument as for the degrees yields

\[
a_f=a_t=a_0=a_1=a_2=1.
\]

Thus each digit is interpreted as a translation:

\[
P_s(x)=x+b_s.
\]

Let the constants be denoted by

\[
b_f=F,\quad b_t=G,\quad b_0=B_0,\quad b_1=B_1,\quad b_2=B_2.
\]

The dynamic odd rule and the third boundary rule imply

\[
G\ge B_2,
\qquad
B_2\ge F+G.
\]

Therefore \(F=0\) and \(B_2=G\).

The rules \(f1\to0t\), \(f2\to1f\), and \(\lhd0\to\lhd t\) give

\[
B_1\ge B_0+G,
\qquad
B_2\ge B_1,
\qquad
B_0\ge G.
\]

Consequently,

\[
G=B_2\ge B_1\ge B_0+G\ge2G.
\]

Hence \(G=0\), and then

\[
B_0=B_1=B_2=F=0.
\]

All five interpretations are the identity. Therefore every rule has equal interpretations on both sides, and none can be strict. \ensuremath{\blacksquare}

This rules out not only a direct scalar polynomial termination proof, but also a scalar polynomial \textbf{reduction pair} intended to make just one dependency-pair rule strict while orienting all others weakly.


\bigskip



\subsubsection*{4. Exact semantic labeling by \(v_2(n+1)\) still does not rescue polynomials}


One might hope that Theorem 2 merely reflects the absence of semantic labels. There is, however, a stronger obstruction.

Write uniquely

\[
n+1=2^r m,
\qquad m\text{ odd}.
\]

The shortcut Collatz map becomes the following two-counter system.

If \(r>0\), then \(n\) is odd and

\[
T(n)+1=\frac{3(n+1)}2=2^{r-1}(3m).
\]

Thus

\[
(m,r)\longrightarrow(3m,r-1).
\tag{4}
\]

If \(r=0\), write

\[
m+1=2^s q,
\qquad q\text{ odd},\quad s\ge1.
\]

Then

\[
T(n)+1=\frac{m+1}{2}=2^{s-1}q,
\]

so

\[
(2^sq-1,0)\longrightarrow(q,s-1).
\tag{5}
\]

The target \(n=1\) corresponds to \((m,r)=(1,1)\), where iteration stops.

Now allow an entirely different scalar polynomial for every exact semantic label \(r\).


\paragraph*{Theorem 3}


There is no family of polynomials

\[
P_r(m)\in\mathbb R[m],
\qquad r=0,1,2,\ldots,
\]

with positive leading coefficients, such that every nonterminal transition satisfies

\[
P_r(m)>P_{r-1}(3m)
\qquad(r\ge1)
\tag{6}
\]

and

\[
P_0(2^sq-1)>P_{s-1}(q)
\qquad(s\ge1).
\tag{7}
\]


\paragraph*{Proof}


Let \(d_r=\deg P_r\).

From (6), by letting \(m\to\infty\),

\[
d_r\ge d_{r-1}.
\]

From (7), by letting \(q\to\infty\),

\[
d_0\ge d_{s-1}
\qquad\text{for every }s.
\]

Therefore all degrees are equal to one common degree \(d\).

If \(d=0\), then (6) says the constants \(P_r\) strictly increase with \(r\), while (7) says \(P_0>P_{s-1}\). This is already contradictory.

Suppose \(d\ge1\), and let \(a_r>0\) be the leading coefficient of \(P_r\). Inequality (6) gives

\[
a_r\ge3^d a_{r-1}.
\]

Iterating,

\[
a_{s-1}\ge3^{d(s-1)}a_0.
\tag{8}
\]

Inequality (7) gives

\[
2^{sd}a_0\ge a_{s-1}.
\tag{9}
\]

Taking \(s=3\), equations (8) and (9) imply

\[
3^{2d}a_0\le a_2\le2^{3d}a_0,
\]

hence

\[
9^d\le8^d,
\]

which is impossible. \ensuremath{\blacksquare}

This theorem rules out a broad form of context-sensitive polynomial order even after supplying the exact unbounded odd-run label \(v_2(n+1)\).

The obstruction is precisely the competition between the two scales:

\[
3^{s-1}\qquad\text{and}\qquad2^s.
\]

The ``consume one unit of \(r\), multiply by \(3\)'' branch forces polynomial coefficients to grow like \(3^{dr}\). The reset branch forces them to be bounded by \(2^{d(r+1)}\). These requirements collide already at \(r=2\).

A commutative Hahn-series or asymptotically polynomial interpretation with leading term

\[
a_r m^d
\]

has the same obstruction. Merely replacing real coefficients by a more exotic ordered commutative field does not help. Any successful series order must use noncommutativity or path dependence in an essential way.


\bigskip



\subsubsection*{5. The exact three-phase size-change system}


Factor the remaining odd part further:

\[
n+1=2^r3^v q,
\qquad \gcd(q,6)=1.
\]

Associate to \(n\) the state

\[
\Psi(n)=(q,r,v).
\]

The transition law is unusually clean.


\paragraph*{Odd phase}


When \(r>0\),

\[
(q,r,v)\longrightarrow(q,r-1,v+1).
\tag{10}
\]

Thus \(r\) strictly decreases, while one factor of \(2\) is converted into a factor of \(3\).


\paragraph*{Even reset with \(v>0\)}


When \(r=0\), write

\[
3^vq+1=2^s3^{v'}q'.
\]

If \(v>0\), then

\[
3^vq+1\equiv1\pmod3,
\]

so \(v'=0\). Therefore the reset strictly decreases \(v\).


\paragraph*{Even reset with \(v=0\)}


If \(r=v=0\), then \(n+1=q\), and

\[
q+1=2^s3^{v'}q'.
\]

For a positive even \(n\), one has \(q>1\), and hence

\[
q'\le\frac{q+1}{2}<q.
\]

Thus this reset strictly decreases \(q\).

We have proved:


\begin{quote}
\textbf{Proposition 4.} Every shortcut Collatz step strictly decreases at least one of the three natural-number coordinates \(r,v,q\).

\end{quote}


This looks tantalizingly close to a size-change termination proof. The difficulty is that the priorities form a cycle:

\[
r\longrightarrow v\longrightarrow q\longrightarrow r.
\]

An \(r\)-decrease may increase \(v\); a \(v\)-decrease may greatly increase \(q\) and reset \(r\); a \(q\)-decrease may reset both \(r\) and \(v\).

Knowing that every \textbf{adjacent} pair decreases in one coordinate does not imply termination. A finite union of well-founded relations need not itself be well-founded.

There are even comparable nonadjacent states on ordinary terminating trajectories. For example,

\[
27\mapsto41\mapsto62\mapsto31\mapsto47\mapsto71
\mapsto107\mapsto161\mapsto242\mapsto121\mapsto182\mapsto91.
\]

But

\[
\Psi(27)=(7,2,0),\qquad
\Psi(91)=(23,2,0).
\]

All three coordinates are nondecreasing between these two points. Therefore the transitive closure of the Collatz relation is not covered by the three simple coordinate-decrease relations.

This is the exact point where semantic carry labels must enter.


\bigskip



\subsubsection*{6. The viable dependency-pair order}


The published automated work proves termination of each subsystem obtained by deleting any one of the eleven rules. Consequently, for any chosen rule \(e\), proving that \(e\) is terminating relative to the whole system would settle termination of \(\mathcal T\), and hence Collatz. [\hypertarget{cit:R31:4}{}\hyperlink{ref:R31:4}{4}]

Semantic labeling is appropriate here because it can introduce finitely or infinitely many labeled copies of symbols while preserving termination equivalence. [\hypertarget{cit:R31:5}{}\hyperlink{ref:R31:5}{5}]

A natural label at precision \(R\) should record:

\[
\lambda_R=
\left(
\text{phase},
\text{last }R\text{ accelerated valuations},
\text{\(2\)-adic seed residue},
\text{\(3\)-adic terminal carry residue}
\right).
\]

The terminal carry component can be represented using your suffix-local expression

\[
C_v+2^{A(v)}\pmod{3^R},
\]

where \(v\) is the terminal suffix of length \(R\). Powers of \(2\) modulo \(3^R\) can be reduced modulo

\[
\varphi(3^R)=2\cdot3^{R-1},
\]

so this component is finite-state at each fixed \(R\).

Each normalized dependency-pair transition then has:

\begin{enumerate}
\item a finite label transition \[
   \lambda\longrightarrow\lambda';
   \]

\item exact affine or semilinear relations among unbounded numerical variables such as \[
   q,\ r,\ v;
   \]

\item possibly a polynomial quotient relation \[
   x'=a_{\lambda,\lambda'}x+b_{\lambda,\lambda'}.
   \]
\end{enumerate}

The correct proof object is not one scalar polynomial. It is a \textbf{semantic size-change dependency graph}.


\paragraph*{Certificate theorem}


Fix a precision \(R\), and suppose the resulting semantic dependency graph is finite. Attach to each edge a size-change graph on a fixed tuple of natural-number variables

\[
z=(z_1,\ldots,z_d),
\]

where the variables may include \(q,r,v\), quotient variables, height bands, and carry-lift counters.

An edge records relations of the form

\[
z'_j\le z_i
\qquad\text{or}\qquad
z'_j<z_i.
\]

Compose these graphs along paths in the usual way: a composite edge is strict if at least one edge in the composed thread is strict.

Assume:


\begin{quote}
Every idempotent graph in the finite composition closure has a strict self-loop \[
> z_i'<z_i
> \] for some \(i\).

\end{quote}


Then the Collatz conjecture follows.


\paragraph*{Proof}


Assume an infinite Collatz orbit exists. Normalizing the mixed-base representation between dynamic rewrites produces an infinite dependency-pair chain.

Because the semantic label graph and its size-change closure are finite, some idempotent composite graph occurs infinitely often along a subsequence of the chain. By hypothesis, that graph contains a strict self-loop on some natural-valued variable \(z_i\).

Repeating the corresponding path gives

\[
z_i^{(0)}>z_i^{(1)}>z_i^{(2)}>\cdots,
\]

an infinite strictly descending sequence of natural numbers, which is impossible.

Therefore there is no infinite dependency-pair chain. Since the auxiliary rules terminate, there is no infinite mixed-base rewrite sequence. By the equivalence theorem, every positive Collatz orbit reaches \(1\). \ensuremath{\blacksquare}

This is a genuine finite proof certificate: once the labeled graph, edge relations, and composition closure are written down, every assertion can be verified exactly.

Matrix interpretations can be viewed as weighted finite automata; the proposed certificate goes beyond them by attaching path-dependent semantic and size-change information rather than a fixed finite-dimensional linear weight. [\hypertarget{cit:R31:6}{}\hyperlink{ref:R31:6}{6}]


\bigskip



\subsubsection*{7. The missing theorem}


The current three-variable abstraction fails the certificate condition because its closure contains the priority cycle

\[
r\downarrow,\ v\uparrow
\quad\leadsto\quad
v\downarrow,\ q\uparrow
\quad\leadsto\quad
q\downarrow,\ r\uparrow,
\]

with no single variable surviving as a strictly descending thread.

The exact theorem needed to complete the proof is therefore:


\begin{quote}
\textbf{Global carry-synchronization lemma.}   There exists a finite precision \(R\) and a finite collection of additional natural-valued carry/height variables such that every arithmetically realizable idempotent element of the semantic size-change closure at precision \(R\) has a strict self-loop.

\end{quote}


Equivalently, every recurrent semantic dependency component must satisfy one of two alternatives:

\[
\text{a natural-valued thread descends around the component},
\]

or

\[
\text{the component has no lift to a positive integral Collatz trajectory}.
\]

Your terminal-suffix locality theorem makes the \(3\)-adic half of this statement genuinely local. It allows exact finite updates of the terminal carry condition. What it does \textbf{not} yet prove is that a recurrent suffix/carry component cannot remain compatible with all earlier \(2\)-adic seed conditions and the archimedean height constraints simultaneously.

That is the global synchronization gap.


\bigskip



\subsubsection*{8. What ordered noncommutative series can and cannot add}


Ordinal and elementary interpretations are powerful enough to prove termination of systems encoding Goodstein sequences and hydra battles, so moving beyond ordinary polynomials is a legitimate possibility rather than decorative language. [\hypertarget{cit:R31:7}{}\hyperlink{ref:R31:7}{7}]

For the present system, one could encode a labeled dependency path

\[
e_1e_2\cdots e_k
\]

as a noncommutative monomial

\[
X_{e_1}X_{e_2}\cdots X_{e_k}
\]

and attach polynomial or ordinal coefficients recording \(q,r,v\) and carry data. The leading monomial would indicate the persistent size-change thread. Concatenation of paths becomes noncommutative multiplication.

But the ordered series is well-founded only when recurrent monomials are excluded or contain a persistent descent thread. In other words, the series construction still requires the global carry-synchronization lemma. Without it, one obtains recurrent monomials corresponding to the unresolved \(r\to v\to q\to r\) cycle.

Theorem 3 also shows that a series whose leading behavior is merely

\[
a_r m^d
\]

cannot work. The leading monomials must remember the \textbf{order of phase changes and carry resets}, not just their aggregate counts.


\bigskip



\subsection*{Final assessment}


This approach does not currently furnish a full proof. It does, however, sharpen the problem substantially:

\begin{enumerate}
\item The auxiliary mixed-base conversion already has an elementary terminating order.

\item No ordinary context-cancellable reduction order can orient the crucial odd and boundary rules.

\item Every direct scalar natural polynomial reduction pair collapses to the identity interpretation.

\item Even infinitely many exact labels \(r=v_2(n+1)\) cannot support scalar polynomial rankings of any degree.

\item The exact dynamics form a cyclic three-phase size-change system in which every individual step decreases \(r\), \(v\), or \(q\).

\item A finite carry-refined semantic size-change certificate would prove Collatz rigorously.

\item The only unresolved assertion is that sufficiently precise simultaneous \(2\)-adic, \(3\)-adic, and archimedean labels eliminate every recurrent size-change component without a descending thread.
\end{enumerate}

So the reduction-order program has not been completed, but it has been narrowed to a concrete theorem rather than a vague request for a ``more powerful ordering.'' A successful proof must break the cyclic reset pattern through exact arithmetic synchronization; no scalar polynomial, finite-dimensional commutative weight, or context-independent monoid order can do that.
\subsection*{References}
\begin{enumerate}[leftmargin=2.2em]
\item \hypertarget{ref:R31:1}{}\hyperlink{cit:R31:1}{\texttt{turn721591view3}}
\item \hypertarget{ref:R31:2}{}\hyperlink{cit:R31:2}{\texttt{turn599120search0}}
\item \hypertarget{ref:R31:3}{}\hyperlink{cit:R31:3}{\texttt{turn852583view0}}
\item \hypertarget{ref:R31:4}{}\hyperlink{cit:R31:4}{\texttt{turn599120search2}}
\item \hypertarget{ref:R31:5}{}\hyperlink{cit:R31:5}{\texttt{turn315252search0}}
\item \hypertarget{ref:R31:6}{}\hyperlink{cit:R31:6}{\texttt{turn786964search7}}
\item \hypertarget{ref:R31:7}{}\hyperlink{cit:R31:7}{\texttt{turn315252search9}}
\end{enumerate}

\clearpage
\section{Round 32: approach 27 --- Digit-block ordinals below epsilon\_0}
\subsection*{Prompt}
\begin{quote}
Complete a full proof of the Collatz conjecture using the following approach:

\begin{enumerate}
\item Digit-block ordinals below epsilon\_0
\end{enumerate}

\textbf{Speculative \textperiodcentered{} P}

Interpret alternating binary and ternary blocks as coefficients in a Cantor normal form. Arrange the representation so that carries may temporarily increase ordinary length while decreasing a higher ordinal component.

Take as long as you need.
\end{quote}
\subsection*{Response}

\subsection*{Verdict}


Carrying the digit-block ordinal approach through produces a \textbf{substantial but incomplete termination argument}. There is an explicit ordinal below \(\omega^3\) that decreases under \textbf{ten of the eleven} rules in the exact mixed binary\textendash{}ternary Collatz rewriting system. The sole exception is the rule encoding an odd Collatz step.

More importantly, one can prove that any successful \(\varepsilon_0\)-ranking of the complete rewriting system induces a pointwise \(\varepsilon_0\)-ranking of the integers themselves. Thus, the missing ``higher ordinal component'' cannot be supplied merely by renaming carries or digit blocks: it must embody a genuine new Collatz descent theorem.

The conjecture is still treated as open in the current literature; a 2026 verification paper explicitly describes a proof as beyond present techniques. [\hypertarget{cit:R32:1}{}\hyperlink{ref:R32:1}{1}, \hypertarget{cit:R32:2}{}\hyperlink{ref:R32:2}{2}]

What follows is the strongest rigorous completion of this approach that I can obtain without inserting an unproved descent lemma.


\bigskip



\subsubsection*{1. The exact mixed binary\textendash{}ternary rewriting system}


Use the shortcut Collatz map

\[
T(n)=
\begin{cases}
n/2,&n\equiv0\pmod 2,\\[2mm]
(3n+1)/2,&n\equiv1\pmod 2.
\end{cases}
\]

Write canonical strings as

\[
/\,u\,.
\]

where \(u\) is a word in

\[
\{f,t,0,1,2\}.
\]

The symbols act as affine maps

\[
f(x)=2x,\qquad t(x)=2x+1,
\]

\[
0(x)=3x,\qquad 1(x)=3x+1,\qquad 2(x)=3x+2,
\]

while \(/(x)=1\) initializes the representation and \(.(x)=x\) terminates it. Thus \(f,t\) are binary digits and \(0,1,2\) are ternary digits.

The eleven rules are

\[
\mathcal D=
\left\{
\begin{aligned}
f.&\longrightarrow .,\\
t.&\longrightarrow 2.,
\end{aligned}
\right.
\]

\[
\mathcal A=
\left\{
\begin{aligned}
f0&\longrightarrow0f,&
f1&\longrightarrow0t,&
f2&\longrightarrow1f,\\
t0&\longrightarrow1t,&
t1&\longrightarrow2f,&
t2&\longrightarrow2t,
\end{aligned}
\right.
\]

and

\[
\mathcal B=
\left\{
\begin{aligned}
/0&\longrightarrow/t,\\
/1&\longrightarrow/ff,\\
/2&\longrightarrow/ft.
\end{aligned}
\right.
\]

The \(\mathcal A\)- and \(\mathcal B\)-rules preserve the represented integer. The two \(\mathcal D\)-rules apply \(T\). Yolcu, Aaronson, and Heule prove that this string-rewriting system terminates exactly when the Collatz conjecture holds. [\hypertarget{cit:R32:3}{}\hyperlink{ref:R32:3}{3}, \hypertarget{cit:R32:4}{}\hyperlink{ref:R32:4}{4}]

This is an ideal setting for the proposed ordinal strategy: the \(\mathcal B\)-rules can increase ordinary word length, but they might still decrease a higher Cantor-normal-form component.


\bigskip



\subsubsection*{2. An exact ordinal descent for ten rules}


Let

\[
B=\{f,t\},\qquad Q=\{0,1,2\}.
\]

For a canonical string \(s=/u.\), define

\[
q(s)=\#\{i:u_i\in Q\},
\]

\[
b(s)=\#\{i:u_i\in B\},
\]

and let

\[
I(s)=\#\{(i,j):i<j,\ u_i\in B,\ u_j\in Q\}.
\]

Thus \(I(s)\) counts how many binary digits occur to the left of ternary digits. These are precisely the binary\textendash{}ternary inversions that the \(\mathcal A\)-rules eliminate.

Define

\[
\boxed{\displaystyle
\Psi(s)=\omega^2 q(s)+\omega I(s)+b(s).
}
\]

Since all three coefficients are finite,

\[
\Psi(s)<\omega^3<\varepsilon_0.
\]


\paragraph*{Proposition 1}


Every rule except

\[
t.\longrightarrow2.
\]

strictly decreases \(\Psi\).


\paragraph*{Proof}


Every \(\mathcal A\)-rule has type

\[
BQ\longrightarrow QB.
\]

The numbers \(q\) and \(b\) are unchanged. Because the two exchanged symbols are adjacent, exactly one binary-before-ternary inversion disappears, while every inversion involving another position is preserved. Therefore

\[
I(s')=I(s)-1
\]

and hence

\[
\Psi(s')<\Psi(s).
\]

Now consider a boundary rule in \(\mathcal B\). It consumes one ternary digit and produces either one or two binary digits. Thus

\[
q(s')=q(s)-1.
\]

The quantities \(I\) and \(b\) may increase substantially, but they remain finite. Consequently

\[
\Psi(s')
<
\omega^2(q(s)-1)+\omega^2
=
\omega^2q(s)
\leq \Psi(s).
\]

This is the desired Goodstein-like phenomenon: for example,

\[
/1\longrightarrow/ff
\]

increases ordinary length, but it removes one unit of the \(\omega^2\)-coefficient, and no finite increase in the lower components can compensate.

Finally,

\[
f.\longrightarrow .
\]

removes the final binary digit. Since that digit has no ternary digit to its right, removing it does not change \(I\). It decreases \(b\) by one, so \(\Psi\) decreases. \ensuremath{\blacksquare}

This gives a complete ordinal proof of termination for

\[
\mathcal T\setminus\{t.\longrightarrow2.\}.
\]

It is stronger than merely proving that carry normalization terminates: the boundary expansions and all even Collatz steps may be freely interleaved.


\bigskip



\subsubsection*{3. The exact failure at an odd step}


Suppose \(s=/ut.\). Replacing the final \(t\) by \(2\) gives

\[
s'=/u2.
\]

Then

\[
q(s')=q(s)+1,\qquad b(s')=b(s)-1.
\]

Furthermore, the new terminal ternary digit lies to the right of every preceding binary digit. If \(b(s)\) includes the terminal \(t\), then

\[
I(s')=I(s)+b(s)-1.
\]

Therefore

\[
\Psi(s')
=
\omega^2(q(s)+1)
+\omega\bigl(I(s)+b(s)-1\bigr)
+b(s)-1,
\]

which is strictly greater than

\[
\Psi(s)=\omega^2q(s)+\omega I(s)+b(s).
\]

So the obstruction is completely localized:

\[
\boxed{\text{The odd rule converts one low-level binary object into a high-level ternary object.}}
\]

One might try reversing the hierarchy and making binary blocks higher than ternary blocks. But then the boundary rules

\[
/0\to/t,\qquad /1\to/ff,\qquad /2\to/ft
\]

convert ternary digits back into one or two binary digits, producing the opposite difficulty.

The base conversion forms a loop:

\[
\text{terminal binary}
\longrightarrow
\text{ternary}
\longrightarrow
\text{new binary digits}
\longrightarrow
\text{terminal binary}.
\]

A fixed hierarchy of binary versus ternary blocks therefore cannot work.


\bigskip



\subsubsection*{4. Three rigorous no-go results}


These do not eliminate all possible \(\varepsilon_0\)-interpretations. They eliminate the most natural versions of the proposed construction.


\paragraph*{4.1 No fixed digit precedence}


Suppose a word is compared lexicographically using a fixed strict precedence \(\succ\) on the digit symbols. Four rewrite rules force

\[
f0\to0f
\quad\Longrightarrow\quad
f\succ0,
\]

\[
/0\to/t
\quad\Longrightarrow\quad
0\succ t,
\]

\[
t0\to1t
\quad\Longrightarrow\quad
t\succ1,
\]

and

\[
/1\to/ff
\quad\Longrightarrow\quad
1\succ f.
\]

Together,

\[
f\succ0\succ t\succ1\succ f,
\]

which is impossible.

Thus one cannot simply assign a fixed ordinal to each digit and read the string as a left-to-right Cantor normal form.


\bigskip



\paragraph*{4.2 No ranking depending only on alternating block lengths}


Ignore digit values and retain only whether each symbol is binary or ternary. Let \(B\) denote a binary digit and \(Q\) a ternary digit.

The rewrite sequence representing the odd step \(3\mapsto5\) is

\[
/t.
\longrightarrow
/2.
\longrightarrow
/ft.
\]

A strict rank based only on type patterns would therefore require

\[
B>Q>BB,
\]

and hence

\[
B>BB.
\]

But the even step \(6\mapsto3\) is

\[
/tf.\longrightarrow/t.,
\]

which requires

\[
BB>B.
\]

Contradiction.

Therefore maximal binary and ternary block lengths\textemdash{}even with a sophisticated fixed ordering of those lengths\textemdash{}do not contain enough information. The actual digit values and their carry context are essential.


\bigskip



\paragraph*{4.3 No positive-slope scalar affine interpretation}


Consider the common termination ansatz

\[
F_a(x)=p_ax+q_a,\qquad p_a\in\mathbb N_{\geq1},
\]

for every symbol \(a\), with strings interpreted by composition. Strict orientation of a rule implies that its left-hand slope is at least its right-hand slope.

From

\[
t.\to2.
\]

we obtain

\[
p_t\geq p_2.
\]

From

\[
/2\to/ft
\]

we obtain

\[
p_2\geq p_fp_t.
\]

Since \(p_f\geq1\), these imply

\[
p_f=1,\qquad p_t=p_2.
\]

The rule

\[
f2\to1f
\]

gives

\[
p_2\geq p_1.
\]

The rule

\[
f1\to0t
\]

gives

\[
p_1\geq p_0p_t,
\]

and

\[
/0\to/t
\]

gives

\[
p_0\geq p_t.
\]

Consequently,

\[
p_t\geq p_1\geq p_0p_t\geq p_t^2.
\]

Because \(p_t\geq1\), this forces

\[
p_f=p_t=p_0=p_1=p_2=1.
\]

All digit maps are now translations. But translations commute, so

\[
F_{f0}=F_{0f},
\]

making it impossible to orient

\[
f0\to0f
\]

strictly.

Thus no positive-slope scalar affine interpretation proves termination. A successful ordinal interpretation must be genuinely nonlinear or context-sensitive.


\bigskip



\subsubsection*{5. The representation-minimum theorem}


The following theorem identifies the exact logical content of an ordinal proof.


\paragraph*{Theorem 2}


The following are equivalent.

\begin{enumerate}
\item There is a function \[
   \Phi:\{\text{canonical mixed-base strings}\}\longrightarrow\varepsilon_0
   \] that strictly decreases under every one of the eleven rewrite rules.

\item There is a function \[
   r:\mathbb N_{>0}\longrightarrow\varepsilon_0
   \] such that \[
   r(T(n))<r(n)
   \qquad\text{for every }n>1.
   \]
\end{enumerate}


\paragraph*{Proof: \(1\Rightarrow2\)}


For each \(n\), let \(\mathcal C_n\) be the nonempty collection of canonical strings representing \(n\). Define

\[
r(n)=\min_{s\in\mathcal C_n}\Phi(s).
\]

The minimum exists because the ordinals are well ordered.

Choose \(s\in\mathcal C_n\) attaining this minimum. The exact simulation theorem for the rewriting system says that from every representation of \(n>1\), some finite rewrite sequence reaches a representation \(s'\) of \(T(n)\). [\hyperlink{ref:R32:4}{4}, \hypertarget{cit:R32:5}{}\hyperlink{ref:R32:5}{5}]

At least one rewrite occurs, so strict descent gives

\[
\Phi(s')<\Phi(s)=r(n).
\]

Therefore

\[
r(T(n))
\leq\Phi(s')
<r(n).
\]


\bigskip



\paragraph*{Proof: \(2\Rightarrow1\)}


Assume \(r(T(n))<r(n)\) for all \(n>1\). Define

\[
\Phi(s)
=
\omega^{\,r(\operatorname{Val}(s))+4}
+
\Psi(s).
\]

Because \(r(\operatorname{Val}(s))<\varepsilon_0\), closure of \(\varepsilon_0\) under exponentiation by \(\omega\) shows that

\[
\Phi(s)<\varepsilon_0.
\]

Under an auxiliary rule in \(\mathcal A\cup\mathcal B\), the represented integer is unchanged, while \(\Psi\) decreases. Hence \(\Phi\) decreases.

Under either dynamic rule, the represented value changes from \(n\) to \(T(n)\). The exponent in the leading term strictly decreases:

\[
r(T(n))+4<r(n)+4.
\]

That decrease dominates any possible change in the lower term \(\Psi<\omega^3\). Hence \(\Phi\) again decreases. \ensuremath{\blacksquare}


\paragraph*{Consequence}


An ordinal ranking on the mixed-base strings is not merely an ingenious way of proving that carries terminate. After minimizing over all representations, it becomes a pointwise ordinal Lyapunov function for the integers:

\[
n\longmapsto r(n),
\qquad
r(T(n))<r(n).
\]

Conversely, any such integer rank immediately yields the requested digit-string rank.

Of course, assuming Collatz, one may take \(r(n)\) to be the total stopping time of \(n\), which even lies below \(\omega\). But that definition is circular unless finiteness of the stopping time has already been proved.

This theorem does not make the ordinal approach useless. It tells us exactly what the approach must deliver: a \textbf{computable, present-state-defined ordinal rank whose decrease is established without referring to the future orbit}.


\bigskip



\subsubsection*{6. The viable hydra-style version}


There is one form of the proposal that could genuinely work.

Assign each digit-block atom in a string a \textbf{generation label}

\[
\lambda<\varepsilon_0.
\]

For a string \(s\), let \(\mathcal L(s)\) be the finite multiset of all its labels and define the Cantor-normal-form ordinal

\[
H(s)
=
\bigoplus_{\alpha\in\mathcal L(s)}\omega^\alpha.
\]

Here \(\oplus\) is natural ordinal addition; equivalently, collect equal exponents and write the result in descending Cantor normal form.

The central hydra principle is

\[
\omega^\alpha
>
\omega^{\beta_1}+\cdots+\omega^{\beta_k}
\qquad
(\beta_i<\alpha).
\]

Thus a high-generation block may be replaced by arbitrarily many lower-generation blocks while the ordinal decreases.


\paragraph*{Digit-generation criterion}


Suppose the labels can be assigned from the current block word so that:

\begin{enumerate}
\item Every \(\mathcal A\)-rule preserves the multiset \(\mathcal L(s)\). It merely rearranges the corresponding atoms and decreases \(I(s)\).

\item The rule \(f.\to.\) removes at least one label.

\item The rule \(t.\to2.\) replaces one label \(\alpha\) by finitely many labels all strictly below \(\alpha\).

\item Every boundary rule \[
   /0\to/t,\qquad /1\to/ff,\qquad /2\to/ft
   \] replaces one label \(\alpha\) by finitely many labels all strictly below \(\alpha\).
\end{enumerate}

Then define

\[
\boxed{\displaystyle
\Omega(s)=
\omega^{H(s)+3}
+\omega I(s)+b(s).
}
\]

If a rule strictly decreases \(H\), the first term decreases and dominates everything below it. If an \(\mathcal A\)-rule preserves \(H\), then \(I\) decreases. Therefore every rewrite strictly decreases \(\Omega\), proving termination and hence Collatz.

So this criterion is a complete conditional proof.


\paragraph*{The missing theorem}


What remains is the following.


\begin{quote}
\textbf{Digit-generation lemma.} Construct \(\mathcal L(s)\), solely from the present alternating binary\textendash{}ternary block structure and digit values, such that conditions 1\textendash{}4 hold for every reachable canonical string.

\end{quote}


The labels cannot be context-free. Indeed, the pair of rules

\[
t.\to2.,
\qquad
/2\to/ft
\]

would require, for fixed symbol labels,

\[
\lambda(t)>\lambda(2)>\lambda(t).
\]

The newly created \(t\) must therefore be recognized as a \textbf{lower-generation descendant} of the old terminal \(t\), despite having the same visible symbol. This generation must either be reconstructible from the present digit configuration or carried by an invariant annotation whose consistency is proved for every possible derivation.

That is precisely the hydra-like mechanism one would need.

Ordinal interpretations have successfully handled rewrite encodings of Goodstein and hydra processes, where growth in ordinary size accompanies descent in a hereditary ordinal. [\hypertarget{cit:R32:6}{}\hyperlink{ref:R32:6}{6}] The Collatz-specific difference is that Goodstein has an external monotonically increasing base or stage parameter. Collatz repeatedly converts binary information to ternary information and then back again, with no evident monotone stage variable.


\bigskip



\subsubsection*{7. Why the natural \(3\)-adic generation label fails}


A tempting candidate is to use the depth of the ternary endpoint carry.

For the accelerated odd map

\[
U(n)=\frac{3n+1}{2^{v_2(3n+1)}},
\]

a valuation word \(a_0,\ldots,a_{j-1}\) satisfies

\[
2^{A_j}U^j(n)=3^jn+C_j,
\qquad
A_j=\sum_{i<j}a_i,
\]

where

\[
C_0=0,\qquad
C_{k+1}=3C_k+2^{A_k}.
\]

Consider the all-ones word

\[
a_0=\cdots=a_{j-1}=1.
\]

Then \(A_k=k\), and induction gives

\[
C_j=3^j-2^j.
\]

Indeed,

\[
C_{j+1}
=
3(3^j-2^j)+2^j
=
3^{j+1}-2^{j+1}.
\]

Therefore

\[
C_j+2^{A_j}
=
C_j+2^j
=
3^j,
\]

and hence

\[
v_3(C_j+2^{A_j})=j.
\]

These words are not merely formal. For every \(j\), the positive seed

\[
n=2^{j+1}-1
\]

has at least its first \(j\) accelerated valuations equal to \(1\).

Thus the most obvious carry depth grows through arbitrarily long realizable prefixes.

If larger \(3\)-adic depth is assigned a larger ordinal generation, the proposed rank increases. If larger depth is assigned a smaller ordinal generation, one would require an infinite descending sequence

\[
\gamma_0>\gamma_1>\gamma_2>\cdots,
\]

which no ordinal system permits.

Therefore raw \(3\)-adic carry depth cannot be the missing generation label. It would have to be combined with another component that falls sufficiently strongly along the all-ones branches.


\bigskip



\subsubsection*{8. Why simply replacing \(2\) and \(3\) by \(\omega\) fails}


For an ordinary binary numeral

\[
n=\sum_i b_i2^i,
\]

the natural ordinal replacement is

\[
H_2(n)=\sum_i b_i\omega^i.
\]

This map is strictly order-preserving: the highest differing binary digit determines both the integer order and the ordinal order.

But for every odd \(n>1\),

\[
T(n)=\frac{3n+1}{2}>n.
\]

Consequently,

\[
H_2(T(n))>H_2(n)
\]

at every odd step.

A hereditary version does not fix the underlying problem unless it introduces a stage or generation that falls during \(2\)-to-\(3\)-to-\(2\) conversion. In Goodstein's process, that stage is supplied externally by the changing base. In Collatz, manufacturing such a stage from the current digits is exactly the unresolved problem.


\bigskip



\subsection*{Final conclusion}


The digit-block ordinal programme yields four rigorous results:

\[
\boxed{\Psi(s)=\omega^2q(s)+\omega I(s)+b(s)}
\]

proves termination of ten of the eleven exact Collatz rewrite rules.

Fixed digit precedence, alternating-block-length rankings, context-free generation labels, and positive-slope scalar affine interpretations cannot orient the complete system.

Any complete \(\varepsilon_0\)-ranking of the strings induces a pointwise integer rank

\[
r(T(n))<r(n),
\]

so a successful ordinal construction necessarily contains the full Collatz descent mechanism rather than merely resolving carry normalization.

Finally, the hydra-style digit-generation criterion gives an exact sufficient theorem: one high-generation block must always be replaced by finitely many lower-generation descendants. The unresolved task is to define those generations from the present mixed-base word and verify the decrease for both

\[
t.\to2.
\]

and

\[
/2\to/ft
\]

without using eventual convergence in the definition.

Accordingly, this approach does \textbf{not yet furnish a valid full proof}. Its precise frontier is the Digit-Generation Lemma above. Claiming that the regenerated binary blocks ``obviously have lower generation'' would be exactly the missing theorem\textemdash{}and, by the representation-minimum theorem, essentially the Collatz conjecture itself.
\subsection*{References}
\begin{enumerate}[leftmargin=2.2em]
\item \hypertarget{ref:R32:1}{}\hyperlink{cit:R32:1}{\texttt{turn290739search20}}
\item \hypertarget{ref:R32:2}{}\hyperlink{cit:R32:2}{\texttt{turn290739search22}}
\item \hypertarget{ref:R32:3}{}\hyperlink{cit:R32:3}{\texttt{turn288388view0}}
\item \hypertarget{ref:R32:4}{}\hyperlink{cit:R32:4}{\texttt{turn288388view1}}
\item \hypertarget{ref:R32:5}{}\hyperlink{cit:R32:5}{\texttt{turn288388view2}}
\item \hypertarget{ref:R32:6}{}\hyperlink{cit:R32:6}{\texttt{turn800059search1}}
\end{enumerate}

\clearpage
\section{Round 33: approach 28 --- Well-structured transition systems}
\subsection*{Prompt}
\begin{quote}
Complete a full proof of the Collatz conjecture using the following approach:

\begin{enumerate}
\item Well-structured transition systems
\end{enumerate}

\textbf{Speculative \textperiodcentered{} P}

Find a well-quasi-order on carry configurations under which the Collatz transducer is monotone. Then nontermination or reachability may reduce to finitely many minimal bad configurations. The major issue is that multiplication by 3 and right shifting do not respect obvious digitwise orders.

Take as long as you need.
\end{quote}
\subsection*{Response}

\subsection*{Result}


I worked the proposed route through its exact logical bottleneck. \textbf{It does not currently produce a full proof of the Collatz conjecture.} More strongly, one can prove that several natural interpretations of ``make the Collatz transducer a WSTS'' are either impossible or circular:

\begin{enumerate}
\item An exact, label-preserving, strongly monotone WQO cannot exist on the canonical Collatz states.
\item An unlabeled strongly monotone WQO always exists nonconstructively, using total stopping time, whether or not Collatz is true; so existence alone proves nothing.
\item A branch-WSTS formulation is meaningful, but its existence is equivalent to proving that every infinite Collatz orbit is eventually periodic. It rules out divergent orbits only after one has proved the substantive Collatz-specific statement.
\item Even then, nontrivial positive cycles still have to be excluded.
\end{enumerate}

Thus I cannot honestly supply the requested unconditional proof. What follows is a rigorous reduction, two no-go theorems, and an exact statement of the missing theorem that would complete the proof.

The relevant established literature supplies the surrounding ingredients: finite-state transducers for \(3x+1\) and division, and a mixed binary\textendash{}ternary rewriting system whose termination is equivalent to Collatz. It does not supply the required WQO compatibility theorem. [\hypertarget{cit:R33:1}{}\hyperlink{ref:R33:1}{1}, \hypertarget{cit:R33:2}{}\hyperlink{ref:R33:2}{2}] The accepted pointwise conjecture also remains beyond Tao's almost-all result. [\hypertarget{cit:R33:3}{}\hyperlink{ref:R33:3}{3}, \hypertarget{cit:R33:4}{}\hyperlink{ref:R33:4}{4}]


\bigskip



\subsubsection*{1. The exact labelled transition system}


Use the shortcut Collatz map, with \(1\) made absorbing:

\[
T(n)=
\begin{cases}
1,&n=1,\\[2mm]
n/2,&n>1 \text{ even},\\[2mm]
(3n+1)/2,&n>1 \text{ odd}.
\end{cases}
\]

Attach the label

\[
\lambda(n)=
\begin{cases}
\mathtt H,&n=1,\\
0,&n>1\text{ even},\\
1,&n>1\text{ odd}.
\end{cases}
\]

Thus every state has one deterministic labelled transition

\[
n\xrightarrow{\lambda(n)}T(n).
\]

The Collatz conjecture says precisely that every label sequence is eventually

\[
\mathtt H\mathtt H\mathtt H\cdots .
\]


\paragraph*{The carry transducer}


Suppose \(n=\sum_{i\geq0}b_i2^i\), read least-significant bit first, and set \(b_{-1}=0\). To compute

\[
3n+1=n+2n+1,
\]

let \(c_0=1\) and define

\[
s_i=b_i+b_{i-1}+c_i,
\qquad
d_i=s_i\bmod 2,
\qquad
c_{i+1}=\left\lfloor\frac{s_i}{2}\right\rfloor .
\]

The state is just

\[
q_i=(b_{i-1},c_i)\in\{0,1\}^2,
\]

so multiplication by \(3\) and addition of \(1\) are realized by a four-state binary transducer. For odd \(n\), the first output bit is \(d_0=0\), and the shortcut division deletes this bit. This is equivalent to the finite-state \(3x+1\) transducer appearing in St\'erin's construction. [\hyperlink{ref:R33:1}{1}]

The finite carry set is therefore not the difficulty. The difficulty is resetting the output word as the input to the next transducer scan after deleting a low bit.


\bigskip



\subsubsection*{2. The precision-loss theorem}


Let

\[
n_t=T^t(n),\qquad
\epsilon_t=n_t\bmod 2
\]

for a segment of an orbit that has not yet reached \(1\), and let

\[
s_t=\epsilon_0+\cdots+\epsilon_{t-1}.
\]


\paragraph*{Lemma 1: affine block identity}


For every \(t\),

\[
2^tT^t(n)=3^{s_t}n+C_t,
\]

where

\[
C_0=0,
\qquad
C_{t+1}=3^{\epsilon_t}C_t+\epsilon_t2^t.
\]

Equivalently, for a parity word \(\epsilon=(\epsilon_0,\ldots,\epsilon_{t-1})\),

\[
C_\epsilon
=
\sum_{\substack{0\leq i<t\\ \epsilon_i=1}}
2^i3^{\epsilon_{i+1}+\cdots+\epsilon_{t-1}}.
\]


\paragraph*{Proof}


Suppose

\[
2^tn_t=3^{s_t}n+C_t.
\]

If \(n_t\) is even, then

\[
2^{t+1}n_{t+1}=2^tn_t,
\]

so \(s_{t+1}=s_t\) and \(C_{t+1}=C_t\).

If \(n_t\) is odd, then

\[
2^{t+1}n_{t+1}
=
2^t(3n_t+1)
=
3(3^{s_t}n+C_t)+2^t,
\]

so \(s_{t+1}=s_t+1\) and

\[
C_{t+1}=3C_t+2^t.
\]

This proves the recurrence. \ensuremath{\blacksquare}


\paragraph*{Lemma 2: parity cylinders are exactly binary congruence classes}


For positive integers \(x,y\), the following are equivalent, provided neither trajectory halts during these \(t\) steps:

\[
\lambda(T^i x)=\lambda(T^i y)
\quad(0\leq i<t),
\]

and

\[
x\equiv y\pmod {2^t}.
\]

Moreover, when these conditions hold,

\[
T^t(y)-T^t(x)
=
3^{s_t}\frac{y-x}{2^t},
\]

and therefore, for \(x\ne y\),

\[
v_2\!\left(T^t(y)-T^t(x)\right)
=
v_2(y-x)-t.
\tag{1}
\]


\paragraph*{Proof}


If the two integers have the same parity word, subtract their affine identities:

\[
2^t\bigl(T^t(y)-T^t(x)\bigr)
=
3^{s_t}(y-x).
\]

The left side is divisible by \(2^t\), while \(3^{s_t}\) is odd, so \(2^t\mid y-x\). The difference formula and (1) follow immediately.

Conversely, suppose \(x\equiv y\pmod {2^t}\). They have the same first parity. Because the same affine branch is applied to both, their first images differ by either

\[
\frac{y-x}{2}
\quad\text{or}\quad
\frac{3(y-x)}2,
\]

which is divisible by \(2^{t-1}\). Induction gives equality of the first \(t\) parity labels. \ensuremath{\blacksquare}


\paragraph*{Corollary 3: itinerary injectivity}


Two positive integers with the same infinite Collatz parity itinerary are equal.

Indeed, Lemma 2 gives

\[
x\equiv y\pmod {2^t}
\]

for every \(t\), hence \(x=y\).

This is the exact form of the obstruction. \textbf{Every complete Collatz microstep spends one bit of shared \(2\)-adic precision.} Multiplication by \(3\) changes the magnitude and \(3\)-adic divisibility of a difference, but it does not replace the binary precision needed to synchronize the next branch.

For the odd-only map

\[
U(n)=\frac{3n+1}{2^{v_2(3n+1)}},
\]

if two seeds realize the same valuation word \(w=(a_0,\ldots,a_{j-1})\), with \(A=\sum a_i\), then analogously

\[
U^j(y)-U^j(x)=3^j\frac{y-x}{2^A},
\]

so that block spends exactly \(A\) binary digits of shared precision.


\bigskip



\subsubsection*{3. No exact label-preserving strong WSTS exists}


A WSTS consists of a transition system equipped with a WQO compatible with its transitions. Strong monotonicity requires a larger configuration to simulate each transition of a smaller one in one step. [\hypertarget{cit:R33:5}{}\hyperlink{ref:R33:5}{5}]

For an exact labelled transducer, the natural requirement is action-compatible strong monotonicity:

\[
x\preceq y,\quad x\xrightarrow{a}x'
\quad\Longrightarrow\quad
y\xrightarrow{a}y'
\text{ and }x'\preceq y'.
\tag{2}
\]


\paragraph*{Theorem 4: trace-separation obstruction}


There is no WQO \(\preceq\) on the canonical positive-integer Collatz configurations satisfying (2).


\paragraph*{Proof}


Suppose \(x\preceq y\).

Because the transition from \(y\) is deterministic and has only the label \(\lambda(y)\), condition (2) implies

\[
\lambda(x)=\lambda(y)
\quad\text{and}\quad
T(x)\preceq T(y).
\]

Applying the same argument repeatedly gives

\[
\lambda(T^k x)=\lambda(T^k y)
\qquad\text{for every }k\geq0.
\]

Thus \(x\) and \(y\) have identical infinite labelled itineraries. By Corollary 3,

\[
x=y.
\]

Consequently, distinct canonical integers are pairwise incomparable. The canonical state set is infinite, so the sequence

\[
1,2,3,4,\ldots
\]

has no increasing pair. Hence \(\preceq\) is not a WQO. \ensuremath{\blacksquare}


\paragraph*{Consequence}


Adding finitely many carry states, modes, residue labels, or transducer phases cannot repair this. If the augmented state space were WQO, its infinite subset of canonical boundary configurations would also be WQO. But exact label-compatible monotonicity would again force two comparable canonical configurations to have the same complete itinerary and therefore the same seed.

So the most direct interpretation of the proposed program is formally impossible:


\begin{quote}
There is no exact, branch-label-preserving, strongly monotone WSTS presentation of Collatz on its canonical integer states.

\end{quote}


One must either forget some labels, weaken monotonicity, or use branch-monotonicity only for ancestor\textendash{}descendant comparisons.


\bigskip



\subsubsection*{4. Why the obvious WQOs fail}


The abstract theorem above explains all of the elementary failures.


\begin{center}
\begin{tabular}{@{}p{0.313\linewidth}p{0.313\linewidth}p{0.313\linewidth}@{}}
\hline
Proposed order & Why it is a WQO & Exact failure \\
\hline
Numerical order \(x\leq y\) & \(\mathbb N\) is well-ordered & \(3<4\), but \(T(3)=5>2=T(4)\). \\
Binary scattered-subword order & Higman's lemma & \(11\sqsubseteq101\), so \(3\sqsubseteq5\), but \(101\not\sqsubseteq1000\), so \(T(3)\not\sqsubseteq T(5)\). \\
Same residue mod \(2^K\), then numerical order & Finite union of chains & \(3\equiv3+2^K\pmod{2^K}\), but their images differ by \(3\cdot2^{K-1}\), not by a multiple of \(2^K\). \\
Exact low-bit prefix/cylinder order & Preserves branch information & It is not a WQO: the least-significant-first words \(0^k1\), representing powers of two, form an antichain under prefix inclusion. \\
Reverse cylinder inclusion & Compatible with refinement & The nested cylinders \(1,10,100,\ldots\), meaning \(n\equiv1\pmod{2^k}\), give an infinite descending chain. \\
Finite carry-state order & Finite, hence WQO & The carry state does not retain the unbounded sequence of low bits required after repeated resets and shifts. \\
\hline
\end{tabular}
\end{center}


The fixed-residue example is not an accident. If an order is forward-invariant and comparability guarantees even one common low bit, then iterating monotonicity guarantees a common low bit after every future step. By Lemma 2 this requires arbitrarily many common low bits initially, hence equality.


\bigskip



\subsubsection*{5. Merely finding some WQO is circular}


There is also a converse danger: if labels and effectiveness are ignored, constructing a monotone WQO is trivial.

Let

\[
\tau(n)=
\begin{cases}
\min\{k:T^k(n)=1\},&\text{if the orbit reaches }1,\\
\omega,&\text{otherwise},
\end{cases}
\]

and order \(0<1<2<\cdots<\omega\). Define

\[
x\preceq_\tau y
\quad\Longleftrightarrow\quad
\tau(x)\leq\tau(y).
\]


\paragraph*{Proposition 5}


The relation \(\preceq_\tau\) is a WQO, and the unlabelled Collatz map is strongly monotone under it.


\paragraph*{Proof}


The order \(0<1<2<\cdots<\omega\) is a well-order, so its pullback is a WQO.

For a terminating nonterminal state,

\[
\tau(Tn)=\tau(n)-1.
\]

For a nonterminating state,

\[
\tau(Tn)=\omega.
\]

Thus \(\tau(x)\leq\tau(y)\) implies

\[
\tau(Tx)\leq\tau(Ty).
\]

Therefore the transition is strongly monotone. \ensuremath{\blacksquare}

This construction works whether Collatz is true or false. If bad states exist, every bad state has rank \(\omega\), and any one of them is a one-element basis for the entire upward-closed bad set.

But determining this order effectively requires knowing which states have rank \(\omega\). Thus the purported finite basis has simply hidden the Collatz conjecture inside the comparison oracle.

This demonstrates why the adjectives \textbf{effective}, \textbf{arithmetic}, and \textbf{independently verifiable} are indispensable.


\bigskip



\subsubsection*{6. Reachability of \(1\) is not ordinary WSTS coverability}


There is another unconditional obstruction.


\paragraph*{Proposition 6}


Under any strongly monotone WQO on the Collatz states, the singleton target \(\{1\}\) cannot be upward-closed.


\paragraph*{Proof}


Consider

\[
2,4,8,16,\ldots .
\]

Since the order is a WQO, there exist \(i<j\) such that

\[
2^i\preceq2^j.
\]

Apply strong monotonicity \(i\) times:

\[
T^i(2^i)\preceq T^i(2^j).
\]

But

\[
T^i(2^i)=1,
\qquad
T^i(2^j)=2^{j-i}>1.
\]

Hence

\[
1\preceq2^{j-i}
\]

for a nonterminal state \(2^{j-i}\). Therefore \(\uparrow\{1\}\ne\{1\}\). \ensuremath{\blacksquare}

So the generic coverability query ``can one reach something above \(1\)?'' cannot be identified with ``can one reach \(1\)?'' Under any such WQO, at least one nonterminal power of two already covers \(1\) before moving.

This matches a standard limitation of WSTS methodology: it is particularly effective for upward-closed coverability and termination properties, but exact reachability can be substantially harder and is not generically handled by the WSTS axioms. [\hypertarget{cit:R33:6}{}\hyperlink{ref:R33:6}{6}]


\bigskip



\subsubsection*{7. The meaningful weakening: branch-WSTS}


A branch-WSTS weakens global strong monotonicity. The relevant condition is:

\[
x\xrightarrow{\sigma}x',
\quad x\preceq x'
\quad\Longrightarrow\quad
x'\xrightarrow{\sigma}y
\text{ for some }y\text{ with }x'\preceq y.
\tag{3}
\]

In words, only when a run reaches a later state covering an earlier state must the intervening labelled block be pumpable. This is called branch-monotonicity. [\hypertarget{cit:R33:7}{}\hyperlink{ref:R33:7}{7}]

A branch-WQO requires every infinite run to contain an increasing ancestor\textendash{}descendant pair.


\paragraph*{Theorem 7: branch-WSTS forces periodicity}


Suppose the exact labelled Collatz transition system admits a branch-WQO satisfying (3). Then every Collatz orbit is eventually periodic.


\paragraph*{Proof}


Take an infinite orbit

\[
x_0\xrightarrow{\epsilon_0}x_1
\xrightarrow{\epsilon_1}x_2
\xrightarrow{\epsilon_2}\cdots .
\]

Because the order is a branch-WQO, there exist \(i<j\) with

\[
x_i\preceq x_j.
\]

Let

\[
\sigma=\epsilon_i\epsilon_{i+1}\cdots\epsilon_{j-1}.
\]

Then

\[
x_i\xrightarrow{\sigma}x_j.
\]

By branch-monotonicity, the same word \(\sigma\) is enabled at \(x_j\), and its endpoint covers \(x_j\). Repeating the argument shows that the actual deterministic label sequence from \(x_j\) is

\[
\sigma^\omega
=
\sigma\sigma\sigma\cdots .
\]

The label sequence from \(x_i\) is also \(\sigma^\omega\), since it begins with \(\sigma\) and then reaches \(x_j\).

Thus \(x_i\) and \(x_j\) have the same infinite itinerary. Corollary 3 gives

\[
x_i=x_j.
\]

Hence the orbit has entered a cycle. \ensuremath{\blacksquare}


\paragraph*{Converse}


If every Collatz orbit is eventually periodic, take equality as the quasiorder.

Equality is branch-monotone: if

\[
x\xrightarrow{\sigma}x'
\quad\text{and}\quad x=x',
\]

then the deterministic cycle repeats \(\sigma\).

Equality is a branch-WQO because every infinite orbit contains a repeated state.

Therefore:

\[
\boxed{
\begin{aligned}
&\text{An exact branch-monotone branch-WQO exists}\\
&\qquad\Longleftrightarrow\\
&\text{every positive Collatz orbit is eventually periodic.}
\end{aligned}}
\tag{4}
\]

This is the sharpest conclusion of the WSTS analysis.

It means that finding such a branch-WQO would indeed exclude divergent orbits\textemdash{}but the existence theorem is exactly equivalent to excluding them. The WSTS framework reorganizes the problem; it does not by itself supply the missing arithmetic rigidity.


\bigskip



\subsubsection*{8. What happens at a pumped block}


Let \(\sigma\) have length \(p\), with \(s\) odd labels. Its affine action is

\[
2^pT^p(x)=3^s x+C_\sigma.
\]

Theorem 7 shows that a branch-WSTS self-covering block must actually satisfy

\[
T^p(x)=x.
\]

Consequently,

\[
(2^p-3^s)x=C_\sigma.
\tag{5}
\]

Since \(C_\sigma>0\), a positive cycle requires

\[
2^p>3^s,
\]

and its least state is uniquely determined by the word:

\[
x=\frac{C_\sigma}{2^p-3^s}.
\tag{6}
\]

Thus each subsumed node of a correct reduced reachability tree would produce one of two things:

\begin{itemize}
\item the trivial terminal behavior at \(1\); or
\item an exact positive integral cycle certificate (5).
\end{itemize}

No third possibility is compatible with exact branch-monotonicity.

This is useful: a finite reduced reachability tree would reduce cycle exclusion to finitely many exact divisibility and parity checks. Branch-effective branch-WSTS do have finite, computable reduced reachability trees for a fixed initial state, and nontermination can be decided by checking for subsumed nodes. [\hypertarget{cit:R33:8}{}\hyperlink{ref:R33:8}{8}]

But Collatz requires one certificate covering \textbf{all} binary inputs, not one fixed seed. Generic branch-WSTS theory does not automatically turn the infinitely many initial integers into a single finitely branching initial configuration without adding spurious infinite ``keep generating digits'' runs.


\bigskip



\subsubsection*{9. The exact theorem that would complete the proof}


A genuine WSTS proof of Collatz would follow from the following result.


\paragraph*{Missing Carry-WQO Theorem}


There exists an effectively decidable quasiorder \(\preceq\) on normalized Collatz carry configurations such that:

\begin{enumerate}
\item \textbf{Exact representation.} Every positive integer has a canonical configuration, and complete transducer passes agree exactly with \(T\).
\item \textbf{Uniform branch-WQO.} Every infinite canonical run contains \(i<j\) with \[
   c_i\preceq c_j.
   \]
\item \textbf{Exact branch-monotonicity.} If \[
   c_i\xrightarrow{\sigma}c_j,\qquad c_i\preceq c_j,
   \] then \(\sigma\) can be repeated from \(c_j\), preserving the comparison.
\item \textbf{Uniform finite initialization.} All finite binary input words are represented by one effective finite reduced reachability forest, without introducing an infinite digit-generation branch.
\item \textbf{Cycle certificate.} Every self-covering node in that forest either belongs to the \(1\)-component or yields equation (5), and exact arithmetic verification rejects all nontrivial candidates.
\end{enumerate}


\paragraph*{Conditional proof}


Assume these five properties and suppose some \(n\) never reaches \(1\).

Its canonical orbit gives an infinite branch in the uniform forest. By property 2, this branch contains a comparable ancestor pair. Property 3 pumps the intervening label block. By Theorem 7, the orbit is periodic. Property 5 then produces a valid nontrivial positive cycle equation, contradicting the finite exact cycle verification. Therefore every \(n\) reaches \(1\). \ensuremath{\blacksquare}

The proof after the Missing Carry-WQO Theorem is complete. The theorem itself is not established.


\bigskip



\subsubsection*{10. The precise obstruction to proving the missing theorem}


Suppose a comparison \(x\preceq y\) is intended to certify that a parity block of length \(p\) can repeat. Then necessarily

\[
x\equiv y\pmod{2^p}.
\]

After one copy of the block,

\[
v_2\!\left(T^p(y)-T^p(x)\right)
=
v_2(y-x)-p.
\]

To certify a second copy, the original difference must therefore have been divisible by \(2^{2p}\). To certify \(r\) copies, it must have been divisible by \(2^{rp}\). Infinite pumping requires

\[
v_2(y-x)=\infty,
\]

hence

\[
x=y.
\]

So no bounded amount of carry, residue, or low-digit information can certify the required pumping of a distinct pair.

A successful order would have to obtain comparability without interpreting it as fixed-depth branch synchronization, while still proving synchronization once the comparable states lie on the same run. That is an unusually delicate combination:

\begin{itemize}
\item globally coarse enough to be a WQO;
\item locally exact enough not to pump \(3\to5\) spuriously;
\item able to recover unbounded \(2\)-adic precision from some other arithmetic coordinate;
\item effective on all finite binary inputs.
\end{itemize}

The only plausible escape is a genuinely two-ended order in which loss of binary precision is compensated by another exact coordinate\textemdash{}perhaps a \(3\)-adic carry ideal together with an archimedean barrier. The terminal-suffix locality phenomenon is relevant to making the \(3\)-adic coordinate finite-state, but it does not yet provide the missing compensation: \(3\)-adic precision cannot by itself determine future binary branch labels.


\bigskip



\subsection*{Final assessment}


The well-structured-transition-system program does not presently yield a Collatz proof. Its rigorous output is instead the following:

\[
\boxed{\text{Exact label-preserving strong WSTS: impossible.}}
\]

\[
\boxed{\text{Unlabelled WSTS existence: nonconstructively trivial and circular.}}
\]

\[
\boxed{\text{Exact branch-WSTS existence: equivalent to eventual periodicity of all orbits.}}
\]

\[
\boxed{\text{A full proof additionally requires uniform finite initialization and cycle exclusion.}}
\]

The most important new lemma is the precision-loss identity

\[
v_2\!\left(T^t(y)-T^t(x)\right)=v_2(y-x)-t
\]

for synchronized parity blocks. It explains exactly why multiplication by \(3\) plus right shifting defeats every fixed-window digit or carry WQO, and it shows that the missing monotonicity theorem cannot be repaired merely by recording a larger fixed modulus.
\subsection*{References}
\begin{enumerate}[leftmargin=2.2em]
\item \hypertarget{ref:R33:1}{}\hyperlink{cit:R33:1}{\texttt{turn286238view3}}
\item \hypertarget{ref:R33:2}{}\hyperlink{cit:R33:2}{\texttt{turn286238view4}}
\item \hypertarget{ref:R33:3}{}\hyperlink{cit:R33:3}{\texttt{turn286238view6}}
\item \hypertarget{ref:R33:4}{}\hyperlink{cit:R33:4}{\texttt{turn286238view7}}
\item \hypertarget{ref:R33:5}{}\hyperlink{cit:R33:5}{\texttt{turn286238view0}}
\item \hypertarget{ref:R33:6}{}\hyperlink{cit:R33:6}{\texttt{turn286238view5}}
\item \hypertarget{ref:R33:7}{}\hyperlink{cit:R33:7}{\texttt{turn286238view1}}
\item \hypertarget{ref:R33:8}{}\hyperlink{cit:R33:8}{\texttt{turn286238view2}}
\end{enumerate}

\clearpage
\section{Round 34: approach 29 --- Cellular-automaton nilpotency}
\subsection*{Prompt}
\begin{quote}
Complete a full proof of the Collatz conjecture using the following approach:

\begin{enumerate}
\item Cellular-automaton nilpotency
\end{enumerate}

\textbf{Established/Hybrid \textperiodcentered{} B}

Use the exact quasi-cellular automaton representation. Prove that every finite-support initial row eventually loses all nontrivial information and enters the periodic background. A persistent glider or spaceship compatible with finite binary support could instead encode a divergent orbit.

Take as long as you need.
\end{quote}
\subsection*{Response}

\subsection*{Result of the proof attempt}


The cellular-automaton route does \textbf{not} presently give a valid full proof. The obstruction is precise: after distinguishing CQCA update time from Collatz iteration depth and normalizing away powers of \(2\), the proposed ``finite-support nilpotency'' theorem is \textbf{exactly equivalent to the Collatz conjecture}.

What the approach does yield is a rigorous package of reductions:

\begin{enumerate}
\item an exact erosion formulation of Collatz;
\item a proof that a finite glider or spaceship would encode a cycle, not a divergent orbit;
\item a proof that any cycle-free counterexample must generate an ever-widening, aperiodic defect;
\item an explicit \(2\)-adic obstruction to any bounded-window nilpotency argument;
\item the exact front-speed inequality that remains to be proved.
\end{enumerate}

I will establish each of these statements.


\bigskip



\subsubsection*{1. The exact CQCA and the correct notion of ``time''}


Use the shortcut Collatz map

\[
T(n)=
\begin{cases}
n/2,&n\equiv0\pmod 2,\\[2mm]
(3n+1)/2,&n\equiv1\pmod 2.
\end{cases}
\]

For odd \(x\), define the odd-to-odd map

\[
U(x)=\frac{3x+1}{2^{\nu_2(3x+1)}}.
\]

The St\'erin\textendash{}Woods CQCA places a finite binary word on an infinite half-defined zero ray. Its cells carry a sum bit and a carry bit. The nonlocal part of the rule initiates the \(+1\) carry at the least significant end, while the local full-adder rule computes \(3x+1\). Successive horizontal rows encode successive \textbf{odd} Collatz iterates, and the trailing zero cells record the intervening divisions by \(2\). Simultaneously, vertical columns encode the full shortcut orbit in a modified ternary representation. [\hypertarget{cit:R34:1}{}\hyperlink{ref:R34:1}{1}, \hypertarget{cit:R34:2}{}\hyperlink{ref:R34:2}{2}, \hypertarget{cit:R34:3}{}\hyperlink{ref:R34:3}{3}]

There are two different dynamics here:

\begin{itemize}
\item \textbf{CQCA update time:} cells gradually become defined;
\item \textbf{Collatz time:} successive odd iterates occur on successively lower spatial rows of the eventual limit configuration.
\end{itemize}

This distinction is essential. The CQCA is freezing: once a sum or carry bit is set, it never changes, and every cell changes state at most twice. Thus every valid input has a pointwise limit configuration, regardless of whether its Collatz orbit converges. [\hypertarget{cit:R34:4}{}\hyperlink{ref:R34:4}{4}]

Consequently, ordinary temporal stabilization of the CQCA cannot prove Collatz. The entire hypothetical divergent orbit would simply be written into an infinite, static, frozen tableau.


\paragraph*{Normalized row extraction}


Let \(w(n)\) be the standard finite binary word for \(n>0\). On every completed CQCA row, discard

\begin{itemize}
\item the infinite leading zero background, and
\item the finite trailing zeros representing powers of \(2\).
\end{itemize}

Let \(\mathcal R_k(n)\) denote the resulting canonical odd binary word on row \(y=-k\), and let \([\mathcal R_k(n)]_2\) be its integer value.

Set

\[
x_0=\frac{n}{2^{\nu_2(n)}},
\qquad
x_{k+1}=U(x_k).
\]

Then the row-simulation lemma is

\[
\boxed{[\mathcal R_k(n)]_2=x_k=U^k(x_0).}
\tag{1}
\]


\paragraph*{Proof}


The initial row represents \(n\); deleting its terminal zeros gives \(x_0\).

Suppose the normalized \(k\)-th row represents the odd number \(x_k\). The bootstrap rule supplies the initial \(+1\) carry. The full-adder relation

\[
s_{\rm out}+2c_{\rm out}
  =s_{\rm current}+s_{\rm shifted}+c_{\rm in}
\]

computes the binary expansion of \(3x_k+1\) on the row below. Since \(x_k\) is odd, that expansion ends in at least one zero. Deleting all of those terminal zeros divides by \(2^{\nu_2(3x_k+1)}\), leaving \(U(x_k)\). Induction proves (1). This is precisely the content of the paper's row-simulation lemma. [\hyperlink{ref:R34:2}{2}, \hyperlink{ref:R34:3}{3}]


\bigskip



\subsubsection*{2. Exact equivalence between erosion and Collatz}


The normalized trivial background is the constant odd-row pattern

\[
1,1,1,\ldots,
\]

because

\[
U(1)=\frac{3+1}{4}=1.
\]

Before odd normalization, this is the usual full-map cycle

\[
1\longmapsto2\longmapsto1.
\]

Define the normalized defect on row \(k\) to be nontrivial when

\[
\mathcal R_k(n)\neq 1.
\]


\paragraph*{Theorem 1 \textemdash{} CQCA erosion equivalence}


For \(n>0\), the following are equivalent:

\begin{enumerate}
\item the Collatz orbit of \(n\) reaches \(1\);
\item \(U^k(x_0)=1\) for some \(k\);
\item \(\mathcal R_k(n)=1\) for some \(k\);
\item every sufficiently deep normalized row is the trivial background \(1\);
\item the finite defect from the \(1\leftrightarrow2\) CQCA background eventually disappears.
\end{enumerate}


\paragraph*{Proof}


The equivalence of 2 and 3 is equation (1).

If the full Collatz orbit reaches \(1\), then its odd subsequence reaches \(1\), proving \(1\Rightarrow2\). Conversely, if an odd iterate reaches \(1\), the full orbit reaches \(1\), proving \(2\Rightarrow1\).

Finally \(U(1)=1\), so one normalized row equal to \(1\) implies all subsequent normalized rows equal \(1\). Hence \(3\Leftrightarrow4\Leftrightarrow5\). \ensuremath{\blacksquare}

Therefore the desired global statement

\[
\text{“Every nonzero finite binary input is eventually eroded”}
\]

is not a consequence of the representation theorem. It is a new statement \textbf{equivalent to Collatz itself}.

Also, ``nilpotency'' is slightly misleading. Standard CA nilpotency usually requires one common time after which every configuration is in the null state. Here the required time depends on the input and is necessarily unbounded. The correct terms are closer to \textbf{finite-defect mortality} or \textbf{erosion modulo a periodic background}.


\bigskip



\subsubsection*{3. What would a glider or spaceship actually encode?}


The historical binary quasi-CA paper already described a failure of Collatz in CA language as either chaotic growth or the emergence of a stable ``vehicle.'' It also observed that unbounded carry propagation is what prevents the original binary description from being an ordinary finite-radius CA. [\hypertarget{cit:R34:5}{}\hyperlink{ref:R34:5}{5}, \hypertarget{cit:R34:6}{}\hyperlink{ref:R34:6}{6}]

For the exact normalized CQCA, however, fixed gliders have a very specific meaning.


\paragraph*{Theorem 2 \textemdash{} Bounded persistent defects are cycles}


Let

\[
x_k=U^k(x_0)
\]

and let

\[
L_k=\lfloor\log_2x_k\rfloor+1
\]

be the width of the normalized \(k\)-th row.

If the widths \(L_k\) remain bounded and the orbit never reaches \(1\), then the orbit enters a nontrivial positive cycle.


\paragraph*{Proof}


If \(L_k\le B\), then every \(x_k\) belongs to the finite set of odd positive integers below \(2^B\). Therefore two iterates repeat:

\[
x_i=x_j,\qquad i<j.
\]

Since \(U\) is deterministic,

\[
x_i,x_{i+1},\ldots,x_{j-1}
\]

is a cycle. It cannot contain \(1\), by assumption. \ensuremath{\blacksquare}


\paragraph*{Corollary 2.1 \textemdash{} A fixed spaceship is a cycle}


Suppose a finite CQCA row pattern reappears after \(p\) odd rows, perhaps translated horizontally. Normalize by deleting trailing zeros and aligning the least significant \(1\). Translation then disappears, so

\[
\mathcal R_{k+p}=\mathcal R_k.
\]

Thus

\[
U^p(x_k)=x_k,
\]

which is a positive cycle.

So the statement in the proposed approach needs refinement:

\begin{itemize}
\item a \textbf{finite oscillator or spaceship} encodes a nontrivial cycle;
\item a \textbf{divergent orbit} requires an expanding or otherwise unbounded defect\textemdash{}a puffer-like structure, not a fixed finite vehicle.
\end{itemize}


\paragraph*{Theorem 3 \textemdash{} A cycle-free counterexample grows without bound}


If a positive orbit neither reaches \(1\) nor enters a cycle, then

\[
x_k\longrightarrow\infty
\]

and hence

\[
L_k\longrightarrow\infty.
\]


\paragraph*{Proof}


Suppose \(x_k\not\to\infty\). Then there is an \(M\) such that \(x_k\le M\) for infinitely many \(k\). Only finitely many positive odd integers are at most \(M\), so one of them occurs twice. Determinism then gives a cycle, contrary to the assumption. \ensuremath{\blacksquare}

This gives an exact CA trichotomy:

\[
\boxed{
\begin{array}{c}
\text{background erosion}\\
\text{or}\\
\text{bounded periodic defect}\\
\text{or}\\
\text{ever-widening aperiodic defect}.
\end{array}}
\]

The second case is a nontrivial cycle. The third is a divergent Collatz orbit.


\bigskip



\subsubsection*{4. Why bounded-depth erosion arguments fail}


There are finite positive seeds that shadow a genuinely persistent \(2\)-adic CQCA phase for arbitrarily long times.


\paragraph*{Theorem 4 \textemdash{} Arbitrarily long expanding all-ones phases}


For every \(r\ge1\), set

\[
x_0=2^{r+1}-1.
\]

Then, for \(0\le i\le r\),

\[
\boxed{x_i=3^i2^{r+1-i}-1,}
\tag{2}
\]

and for every \(0\le i<r\),

\[
\nu_2(3x_i+1)=1.
\]

Consequently the first \(r\) odd transitions are all expanding:

\[
x_{i+1}=\frac{3x_i+1}{2}>x_i.
\]


\paragraph*{Proof}


The claim is clear for \(i=0\). Assuming (2),

\[
\begin{aligned}
3x_i+1
 &=3\left(3^i2^{r+1-i}-1\right)+1\\
 &=3^{i+1}2^{r+1-i}-2\\
 &=2\left(3^{i+1}2^{r-i}-1\right).
\end{aligned}
\]

For \(i<r\), the quantity in parentheses is odd, so the \(2\)-adic valuation is exactly \(1\). Therefore

\[
x_{i+1}=3^{i+1}2^{r-i}-1,
\]

which is (2) with \(i+1\). \ensuremath{\blacksquare}

Thus there is no fixed number \(K\) such that every nonbackground row contracts, or even stops expanding, within \(K\) odd rows.


\paragraph*{The \(2\)-adic meaning}


The binary words

\[
1,\quad11,\quad111,\quad1111,\ldots
\]

converge \(2\)-adically to

\[
\cdots111=-1.
\]

And \(-1\) is an exact fixed point:

\[
U(-1)=\frac{3(-1)+1}{2}=-1.
\]

The CQCA framework extends naturally to infinite binary inputs in \(\mathbb Z_2\), and the paper explicitly depicts the negative fixed cycle at \(-1\). [\hyperlink{ref:R34:1}{1}, \hypertarget{cit:R34:7}{}\hyperlink{ref:R34:7}{7}]

So \(2^{r+1}-1\) is a finite-support approximation to the persistent all-ones phase \(-1\). It follows that the positive finite seed shadows that immortal phase for \(r\) odd rows before its distant high-order zero boundary becomes visible.

This is the central CA obstruction:


\begin{quote}
The local bulk can look exactly like a persistent glider for arbitrarily long periods. Any proof must use the eventual zero boundary of the finite binary input, not merely the local transition rule.

\end{quote}



\paragraph*{A stronger finite-window shadowing lemma}


Let \(x,y\in\mathbb Z_2\). If

\[
x\equiv y\pmod{2^q},
\]

then \(x\) and \(y\) have the same parity and

\[
T(x)\equiv T(y)\pmod{2^{q-1}}.
\]

Indeed, on the even branch,

\[
T(x)-T(y)=\frac{x-y}{2},
\]

and on the odd branch,

\[
T(x)-T(y)=\frac{3(x-y)}2.
\]

Inductively,

\[
x\equiv y\pmod{2^{M+t}}
\quad\Longrightarrow\quad
T^j(x)\equiv T^j(y)\pmod{2^M}
\]

for every \(0\le j\le t\).

Every \(2\)-adic input has a positive finite binary representative modulo \(2^{M+t}\). Therefore:

\[
\boxed{\text{Every bounded spacetime window of a persistent \(2\)-adic orbit
is reproduced by some positive finite input.}}
\]

This defeats a tempting compactness argument. Arbitrarily long nonbackground finite tableaux certainly have infinite local limits, but those limits can simply be infinite-support \(2\)-adic configurations such as \(-1\). Finite binary support is not a closed local condition.

It also means that no finite collection of forbidden CQCA patches can prove Collatz.


\bigskip



\subsubsection*{5. The exact front-speed equation}


Let \(x=x_0\) be odd and define

\[
x_{i+1}=U(x_i),\qquad
a_i=\nu_2(3x_i+1).
\]

Set

\[
A_j=\sum_{i=0}^{j-1}a_i,
\qquad
S_r=A_r=\sum_{i=0}^{r-1}a_i.
\]

The exact affine identity is

\[
\boxed{
2^{S_r}x_r=3^r x+C_r,
}
\tag{3}
\]

where

\[
\boxed{
C_r=\sum_{j=0}^{r-1}3^{r-1-j}2^{A_j}.
}
\tag{4}
\]


\paragraph*{Proof}


For \(r=1\),

\[
2^{a_0}x_1=3x+1.
\]

If (3) holds at \(r\), then

\[
\begin{aligned}
2^{S_{r+1}}x_{r+1}
 &=2^{S_r}(3x_r+1)\\
 &=3\left(3^rx+C_r\right)+2^{S_r}\\
 &=3^{r+1}x+\left(3C_r+2^{S_r}\right),
\end{aligned}
\]

and the expression in parentheses is exactly \(C_{r+1}\). \ensuremath{\blacksquare}

In CQCA language:

\begin{itemize}
\item \(r\log_2 3\) is the ideal outward motion caused by \(r\) multiplications by \(3\);
\item \(S_r\) is the total inward translation from deleting trailing-zero blocks;
\item \(C_r\) is the accumulated flux from all \(+1\) carry injections.
\end{itemize}

Taking logarithms gives the exact front displacement:

\[
\boxed{
\log_2\frac{x_r}{x}
=
r\log_2 3-S_r
+
\log_2\left(1+\frac{C_r}{3^r x}\right).
}
\tag{5}
\]

The final term is positive. Therefore merely proving

\[
S_r>r\log_2 3
\]

is not always sufficient; the carry surplus must also be beaten.

The CQCA's base-\(3\)-to-base-\(2\) duality geometrically represents these same arithmetic quantities, but it is an exact conversion identity rather than a dissipative inequality. [\hypertarget{cit:R34:8}{}\hyperlink{ref:R34:8}{8}]


\bigskip



\subsubsection*{6. The precise missing nilpotency lemma}


There is an elementary equivalent descent formulation of Collatz.


\paragraph*{Theorem 5 \textemdash{} Descent equivalence}


The Collatz conjecture is equivalent to

\[
\boxed{
\forall x>1\text{ odd},\quad
\exists r\ge1\text{ such that }U^r(x)<x.
}
\tag{6}
\]


\paragraph*{Proof}


If Collatz holds, the orbit eventually reaches \(1<x\), so (6) holds.

Conversely, assume (6). Use strong induction over odd positive integers. The case \(x=1\) is immediate. For odd \(x>1\), choose \(r\) with

\[
y=U^r(x)<x.
\]

By the induction hypothesis, \(y\) reaches \(1\), and therefore so does \(x\). Even starting values first reduce to their odd part. \ensuremath{\blacksquare}

Substituting (3), condition \(x_r<x\) becomes

\[
3^rx+C_r<2^{S_r}x,
\]

or

\[
\boxed{
(2^{S_r}-3^r)x>C_r.
}
\tag{7}
\]

Thus a complete CQCA nilpotency proof must establish:

\[
\boxed{
\forall x>1\text{ odd},\quad
\exists r\ge1:
\quad
(2^{S_r}-3^r)x>C_r.
}
\tag{E}
\]

This is the exact \textbf{boundary-capture theorem}: eventually the total inward displacement from zero-block deletion must overtake both multiplication growth and accumulated carry flux.


\paragraph*{Minimal-counterexample formulation}


If Collatz is false, choose the least odd counterexample \(x\). Then

\[
x_r\ge x\qquad\text{for every }r,
\]

because an iterate below \(x\) would converge by minimality and would carry \(x\) with it.

Equation (3) then gives

\[
3^rx+C_r=2^{S_r}x_r\ge2^{S_r}x,
\]

so the minimal counterexample satisfies the infinite family of barrier inequalities

\[
\boxed{
(2^{S_r}-3^r)x\le C_r
\qquad\text{for every }r\ge1.
}
\tag{8}
\]

In CA language, a counterexample is therefore a finite-support initial row whose moving information frontier satisfies (8) forever.

This meets the same obstruction exposed by your terminal-suffix locality theorem. Terminal \(3\)-adic divisibility may depend only on a bounded valuation suffix, but \(C_r\) contains the entire valuation prefix. Erosion requires a synchronization between the distant finite-support boundary, the accumulated prefix carries, and the terminal suffix state.


\bigskip



\subsubsection*{7. Why the obvious CA proof mechanisms do not close the gap}


\textbf{Freezing order.} The partial order on cell states proves that each cell settles, but it settles while recording either a terminating or nonterminating Collatz tableau. Freezing is not an erasure invariant. More broadly, two-dimensional freezing cellular automata can support universal computation, so freezing alone does not imply simple global behavior. [\hypertarget{cit:R34:9}{}\hyperlink{ref:R34:9}{9}]

\textbf{Counting active cells.} The number or width of active cells is not monotone. The inputs \(2^{r+1}-1\) give arbitrarily long intervals of strict growth.

\textbf{Fixed-horizon contraction.} The same examples rule out a universal \(K\) for which every block of \(K\) odd rows contracts.

\textbf{Glider exclusion.} Excluding all finite gliders only excludes positive cycles. A divergent counterexample is necessarily an unbounded puffer-like region.

\textbf{Compactness.} Taking limits of longer and longer survivor tableaux produces an infinite \(2\)-adic input, not necessarily one positive integer. The all-ones limit \(-1\) is already a persistent example.

\textbf{Local forbidden-pattern methods.} The finite-window shadowing lemma shows that every bounded patch of a persistent \(2\)-adic configuration occurs in a finite positive input.

\textbf{Base-conversion conservation laws.} Summing the full-adder equations reproduces (3)\textendash{}(5). What is missing is a sign or well-founded rank; the exact conservation law itself supplies neither.


\bigskip



\subsection*{Final status}


The cellular-automaton program yields the following rigorous conclusions.

\begin{itemize}
\item The exact CQCA representation is valid.
\item Ordinary temporal stabilization is automatic and unrelated to Collatz.
\item Tail erosion modulo the \(1\leftrightarrow2\) background is exactly equivalent to Collatz.
\item Every bounded persistent normalized pattern is a nontrivial cycle.
\item Every cycle-free counterexample must produce an aperiodic defect whose binary width tends to infinity.
\item Finite positive inputs can shadow immortal \(2\)-adic CQCA phases for arbitrarily long times.
\item Consequently, neither fixed-window analysis, compactness, freezing, nor finite glider enumeration proves erosion.
\item The unresolved theorem is precisely the boundary-capture inequality
\end{itemize}

\[
\forall x>1\text{ odd}\ \exists r:
\qquad
(2^{S_r}-3^r)x>C_r.
\]

Everything up to this boxed statement is proved above. Asserting the boxed statement without a new argument would simply reassert Collatz in the language of cellular-automaton front motion. A genuine completion of approach 29 must exploit the \textbf{eventually zero high-bit boundary} to rule out an infinite succession of \(2\)-adic glider shadows while controlling the global carry term \(C_r\).
\subsection*{References}
\begin{enumerate}[leftmargin=2.2em]
\item \hypertarget{ref:R34:1}{}\hyperlink{cit:R34:1}{\texttt{turn354100view0}}
\item \hypertarget{ref:R34:2}{}\hyperlink{cit:R34:2}{\texttt{turn771128view3}}
\item \hypertarget{ref:R34:3}{}\hyperlink{cit:R34:3}{\texttt{turn771128view4}}
\item \hypertarget{ref:R34:4}{}\hyperlink{cit:R34:4}{\texttt{turn771128view1}}
\item \hypertarget{ref:R34:5}{}\hyperlink{cit:R34:5}{\texttt{turn898480view4}}
\item \hypertarget{ref:R34:6}{}\hyperlink{cit:R34:6}{\texttt{turn360652view3}}
\item \hypertarget{ref:R34:7}{}\hyperlink{cit:R34:7}{\texttt{turn541322view1}}
\item \hypertarget{ref:R34:8}{}\hyperlink{cit:R34:8}{\texttt{turn771128view5}}
\item \hypertarget{ref:R34:9}{}\hyperlink{cit:R34:9}{\texttt{turn360652view0}}
\end{enumerate}

\clearpage
\section{Round 35: approach 30 --- Base-6 carry particles}
\subsection*{Prompt}
\begin{quote}
Complete a full proof of the Collatz conjecture using the following approach:

\begin{enumerate}
\item Base-6 carry particles
\end{enumerate}

\textbf{Speculative \textperiodcentered{} B}

Because 6 simultaneously contains the primes 2 and 3, represent the orbit in base 6 and model carries as interacting particles. Classify all persistent particle configurations. A nontrivial glider would be a counterexample mechanism; extinction of every finite configuration would prove convergence.

Take as long as you need.
\end{quote}
\subsection*{Response}

\subsubsection*{Verdict}


I did \textbf{not} obtain a valid full proof of the Collatz conjecture from the base\nobreakdash-6 particle model. What the approach does produce is an exact and unusually transparent formulation of the remaining obstruction:


\begin{quote}
The base\nobreakdash-6 bulk is a reversible carry gas. Convergence can therefore come only from an imbalance between a source at the least-significant boundary and a sink at the most-significant boundary.

\end{quote}


A localized periodic ``glider'' is not a divergent counterexample: once the moving boundary is included, it is exactly a Collatz cycle. Any divergent counterexample would instead be an \textbf{unbounded, aperiodic particle cloud} whose source rate remains critically synchronized with the rate at which leading cells disappear.

This limitation is substantive. Base\nobreakdash-6 cellular representations were already developed by Korec, who observed the absence of carry propagation and constructed an eight-state system whose nilpotency is equivalent to the Collatz conjecture. Later work likewise notes that base \(6\) makes the process local. [\hypertarget{cit:R35:1}{}\hyperlink{ref:R35:1}{1}, \hypertarget{cit:R35:2}{}\hyperlink{ref:R35:2}{2}, \hypertarget{cit:R35:3}{}\hyperlink{ref:R35:3}{3}] Recent technical work in 2026 continues to treat a general proof as unavailable. [\hypertarget{cit:R35:4}{}\hyperlink{ref:R35:4}{4}]

What follows is the exact reduction, the classification that can be proved, and the precise theorem still missing.


\bigskip



\subsection*{1. The exact base\nobreakdash-6 particle rule}


Let

\[
T(n)=
\begin{cases}
n/2,&n\equiv0\pmod 2,\\[2mm]
3n+1,&n\equiv1\pmod 2.
\end{cases}
\]

Write

\[
n=\sum_{j=0}^{m}d_j6^j,\qquad d_j\in\{0,1,2,3,4,5\},
\]

with \(d_m\ne0\). Split every digit uniquely as

\[
d_j=2c_j+p_j,
\qquad
c_j\in\{0,1,2\},
\qquad
p_j\in\{0,1\}.
\]

Thus \(c_j\) is a ternary carry charge and \(p_j\) is a binary parity particle.


\subsubsection*{Lemma 1: quotient-and-feedback form}


Write \(n=2q+r\), \(r\in\{0,1\}\). Then

\[
T(n)=
\begin{cases}
q,&r=0,\\
6q+4,&r=1.
\end{cases}
\tag{1}
\]


\paragraph*{Proof}


The even case is immediate. If \(n=2q+1\), then

\[
3n+1=3(2q+1)+1=6q+4.
\]

\(\square\)

Thus one Collatz step does the following:

\begin{enumerate}
\item divide the base\nobreakdash-6 word by \(2\);
\item inspect the discarded remainder bit;
\item if it is \(1\), feed it back by appending a terminal digit \(4\).
\end{enumerate}

This is the cleanest form of the base\nobreakdash-6 mechanism.

Set \(p_{m+1}=0\). The base\nobreakdash-6 digits of \(q=\lfloor n/2\rfloor\) are

\[
e_j=c_j+3p_{j+1},
\qquad 0\le j\le m.
\tag{2}
\]

Indeed,

\[
\left\lfloor\frac n2\right\rfloor
=
\sum_{j=0}^{m}c_j6^j
+
\sum_{j=1}^{m}3p_j6^{j-1}.
\]

Equation (2) shows that the binary particle \(p_{j+1}\) moves one place toward the least-significant boundary and interacts with the stationary ternary charge \(c_j\).

Split \(e_j=2c'_j+p'_j\). The collision gate is

\[
\Phi(c,p)
=
\left(
\left\lfloor\frac{c+3p}{2}\right\rfloor,
(c+3p)\bmod2
\right).
\tag{3}
\]

Its complete table is

\[
\begin{array}{c|c|c}
(c,p)&c+3p&\Phi(c,p)\\
\hline
(0,0)&0&(0,0)\\
(0,1)&3&(1,1)\\
(1,0)&1&(0,1)\\
(1,1)&4&(2,0)\\
(2,0)&2&(1,0)\\
(2,1)&5&(2,1).
\end{array}
\tag{4}
\]

If the pair \((c,p)\) is encoded by \(x=2c+p\), then the collision is the permutation

\[
0\mapsto0,\qquad
5\mapsto5,\qquad
1\mapsto3\mapsto4\mapsto2\mapsto1.
\tag{5}
\]

In particular, the interior collision is bijective.

There is a dual multiplication picture. Since

\[
3d_j=6c_j+3p_j,
\]

the \(j\)-th digit of \(3n\) is

\[
3p_j+c_{j-1},
\tag{6}
\]

where \(c_{-1}=0\). Thus multiplication by \(3\) streams the \(c\)-particles toward more significant positions and applies exactly the same collision gate \(\Phi\). The division and multiplication pictures are rotations of the same reversible lattice gas.


\bigskip



\subsection*{2. The only dissipation occurs at the two boundaries}


Let \(n_t=T^t(n_0)\). Define

\[
b_t=
\begin{cases}
1,&n_t\text{ odd},\\
0,&n_t\text{ even},
\end{cases}
\]

so \(b_t\) records whether a new digit \(4\) is injected at the low-order boundary.

Let

\[
h_t=
\begin{cases}
1,&\text{the leading base-6 digit of }n_t\text{ is }1,\\
0,&\text{the leading digit is }2,3,4,\text{ or }5.
\end{cases}
\]

If the leading digit is \(1\), division by \(2\) turns that leading position into \(0\), so the high-order cell disappears.

Finally, let

\[
\ell_t=1+\lfloor \log_6 n_t\rfloor
\]

be the number of base\nobreakdash-6 digits.


\subsubsection*{Proposition 2: exact source-sink balance}


For every \(t\),

\[
\boxed{\ell_{t+1}-\ell_t=b_t-h_t.}
\tag{7}
\]

Consequently, if

\[
B_t=\sum_{k=0}^{t-1}b_k,
\qquad
H_t=\sum_{k=0}^{t-1}h_k,
\]

then

\[
\boxed{\ell_t=\ell_0+B_t-H_t.}
\tag{8}
\]


\paragraph*{Proof}


The quotient \(q=\lfloor n_t/2\rfloor\) loses one leading digit precisely when the leading digit of \(n_t\) is \(1\). Hence its length is \(\ell_t-h_t\), with the convention that \(0\) has length \(0\).

If \(b_t=0\), the output is \(q\). If \(b_t=1\), equation (1) says that the output is \(6q+4\), which appends one digit. Therefore

\[
\ell_{t+1}=\ell_t-h_t+b_t.
\]

\(\square\)

The particle system is therefore a deterministic queue:

\begin{itemize}
\item \(b_t=1\) injects one cell at the low-order end;
\item \(h_t=1\) removes one cell at the high-order end;
\item all interior evolution is reversible.
\end{itemize}


\bigskip



\subsection*{3. The high-order sink is an almost Sturmian clock}


Set

\[
\beta=\log_6 2=0.3868528072\ldots
\]

and write

\[
\log_6 n_t=(\ell_t-1)+\theta_t,
\qquad 0\le\theta_t<1.
\]

The condition that the leading digit is \(1\) is exactly

\[
h_t=\mathbf 1_{\{\theta_t<\beta\}}.
\tag{9}
\]

Define

\[
\delta_t
=
b_t\log_6\left(1+\frac{1}{3n_t}\right).
\tag{10}
\]


\subsubsection*{Proposition 3: phase evolution}


One has

\[
\log_6 n_{t+1}-\log_6 n_t
=
-\beta+b_t+\delta_t
\tag{11}
\]

and

\[
\boxed{\theta_{t+1}
=
\theta_t+h_t-\beta+\delta_t.}
\tag{12}
\]

Consequently,

\[
\boxed{
H_t
=
\beta t+\theta_t-\theta_0-\Delta_t,
}
\qquad
\Delta_t=\sum_{k=0}^{t-1}\delta_k,
\tag{13}
\]

and

\[
\boxed{
\ell_t
=
\ell_0+B_t-\beta t-\theta_t+\theta_0+\Delta_t.
}
\tag{14}
\]


\paragraph*{Proof}


If \(n_t\) is even, then \(n_{t+1}=n_t/2\), giving a logarithmic increment \(-\beta\).

If \(n_t\) is odd, then

\[
3n_t+1
=
3n_t\left(1+\frac1{3n_t}\right),
\]

so

\[
\log_6\frac{n_{t+1}}{n_t}
=
\log_6 3+
\log_6\left(1+\frac1{3n_t}\right)
=
1-\beta+\delta_t.
\]

This proves (11). Subtracting the integer length increment \(b_t-h_t\) from (11) gives (12). Summation gives (13), and substitution into (8) gives (14). \(\square\)

If the \(+1\) perturbations were absent, equation (12) would be the irrational rotation

\[
\theta\longmapsto \theta-\beta\pmod1.
\]

Thus the top-cell loss sequence has average density \(\beta\), up to the accumulated positive perturbation \(\Delta_t\).

For a nonperiodic orbit, all \(n_t\) are distinct. Therefore

\[
0\le\Delta_t
\le
\frac{1}{3\log 6}
\sum_{k<t}\frac1{n_k}
\le
\frac{\mathcal H_t}{3\log6},
\tag{15}
\]

where

\[
\mathcal H_t=1+\frac12+\cdots+\frac1t.
\]

Hence

\[
H_t=\beta t+O(\log t)
\tag{16}
\]

on every nonperiodic orbit.

This identifies the exact critical source rate:

\[
\boxed{\text{critical odd-step density}=\log_6 2.}
\]

A divergent orbit would have to inject new terminal cells at essentially the same rate at which the rotation clock deletes leading cells.


\bigskip



\subsection*{4. A strong necessary condition for a minimal counterexample}


Assume \(n_0\) is the least positive integer that does not reach \(1\). Then

\[
n_t\ge n_0
\qquad\text{for all }t,
\tag{17}
\]

because any smaller iterate would converge by minimality.

In particular,

\[
\ell_t\ge\ell_0,
\]

so (8) gives

\[
B_t\ge H_t
\qquad\text{for every }t.
\tag{18}
\]

Moreover,

\[
\Delta_t
\le
\frac{B_t}{3n_0\log6}.
\tag{19}
\]

Combining (13), (18), and \(\theta_t-\theta_0>-1\) yields

\[
B_t
\ge
\beta t-1-\frac{B_t}{3n_0\log6}.
\]

Therefore:


\subsubsection*{Corollary 4: pointwise source barrier}


Every minimal counterexample would satisfy, for every \(t\ge1\),

\[
\boxed{
B_t
\ge
\frac{\beta t-1}
{1+\dfrac{1}{3n_0\log6}}.
}
\tag{20}
\]

This is a prefix-by-prefix condition, not merely an asymptotic density condition.

Since an odd standard Collatz step is always followed by an even step,

\[
B_t\le \left\lceil\frac t2\right\rceil.
\tag{21}
\]

Thus a minimal counterexample would have to produce an infinite binary source word with no adjacent \(1\)'s whose every prefix lies almost entirely above the line of slope

\[
\beta=0.3868528072\ldots.
\]

This is the exact source/sink barrier seen by the base\nobreakdash-6 particle model.


\bigskip



\subsection*{5. Fixed-lattice particle waves}


There is another useful representation in which the moving low-order boundary becomes an explicit marker.

Let

\[
E_t=t-B_t
\]

be the number of even steps before time \(t\), and put

\[
N_t=6^{E_t}n_t.
\tag{22}
\]

The base\nobreakdash-6 digits of \(N_t\) occupy a fixed lattice, while the least-significant marker sits at coordinate \(E_t\).


\subsubsection*{Proposition 5: source-wave equation}


The fixed-frame variable satisfies

\[
\boxed{
N_{t+1}=3N_t+b_t6^{E_t}.
}
\tag{23}
\]

Consequently,

\[
\boxed{
N_t
=
3^tn_0+
\sum_{k=0}^{t-1}
b_k\,3^{t-1-k}6^{E_k}.
}
\tag{24}
\]


\paragraph*{Proof}


If \(n_t\) is even, then \(E_{t+1}=E_t+1\), and

\[
N_{t+1}
=
6^{E_t+1}\frac{n_t}{2}
=
3N_t.
\]

If \(n_t\) is odd, then \(E_{t+1}=E_t\), and

\[
N_{t+1}
=
6^{E_t}(3n_t+1)
=
3N_t+6^{E_t}.
\]

Iteration gives (24). \(\square\)

Every odd step injects a unit source at the current marker. That source then evolves under multiplication by \(3\), whose carries propagate locally by the reversible gate (3).

Let

\[
\alpha=\log_6 3=1-\beta
\]

and define the marker-front discrepancy

\[
D_t=E_t-\alpha t=\beta t-B_t.
\tag{25}
\]

Dividing (24) by \(6^{E_t}\) gives

\[
\boxed{
n_t
=
6^{-D_t}
\left(
n_0+
\frac13
\sum_{k=0}^{t-1}b_k6^{D_k}
\right).
}
\tag{26}
\]

Moreover,

\[
\boxed{
6^{D_t}
=
\frac{2^{E_t}}{3^{B_t}}.
}
\tag{27}
\]

Thus the scalar marker-front displacement is precisely the usual multiplicative Collatz coefficient:

\[
6^{-D_t}=\frac{3^{B_t}}{2^{E_t}}.
\]

This is an important negative conclusion. At the scalar level, the base\nobreakdash-6 particle model does not create a new drift invariant; it geometrizes the familiar competition between powers of \(2\) and \(3\).

A successful particle proof would therefore have to use correlations in the collision field that are invisible to \(D_t\), not merely its average velocity.


\bigskip



\subsection*{6. Classification of persistent configurations}


The following classification is complete as a dynamical dichotomy, although it does not exclude the second class.


\subsubsection*{Theorem 6: cycle-or-expanding-cloud dichotomy}


For a finite positive initial configuration, exactly one of the following occurs:

\begin{enumerate}
\item the orbit eventually enters a periodic cycle;
\item the orbit is nonperiodic, all iterates are distinct, and \[
   n_t\longrightarrow\infty,
   \qquad
   \ell_t\longrightarrow\infty.
   \]
\end{enumerate}


\paragraph*{Proof}


If two iterates agree, the deterministic orbit is eventually periodic.

Otherwise all \(n_t\) are distinct. For every fixed \(M\), only finitely many positive integers are at most \(M\). Hence only finitely many \(n_t\) can be at most \(M\). Therefore \(n_t\to\infty\), and consequently \(\ell_t\to\infty\). \(\square\)

Thus a divergent counterexample cannot remain a localized finite particle object. Its occupied interval must grow without bound.


\bigskip



\subsection*{7. Gliders are cycles}


A marked fixed-lattice state consists of

\[
\mathcal X_t=(E_t,\text{the base-6 digit configuration of }N_t).
\]

Let \(\sigma\) denote translation by one base\nobreakdash-6 cell.


\subsubsection*{Theorem 7: marked glider theorem}


Suppose that for some \(q>0\) and integer \(s\),

\[
\mathcal X_{t+q}=\sigma^s\mathcal X_t.
\tag{28}
\]

Then

\[
n_{t+q}=n_t.
\]

Conversely, every Collatz cycle gives such a marked glider.


\paragraph*{Proof}


Translation by \(s\) moves the marker and multiplies the represented fixed-frame integer by \(6^s\). Therefore (28) implies

\[
E_{t+q}=E_t+s,
\qquad
N_{t+q}=6^sN_t.
\]

Hence

\[
n_{t+q}
=
\frac{N_{t+q}}{6^{E_{t+q}}}
=
\frac{6^sN_t}{6^{E_t+s}}
=
n_t.
\]

The converse follows by reversing the calculation. \(\square\)

So a nontrivial marked glider would be a \textbf{nontrivial Collatz cycle}, not a divergent trajectory.

There are also no nonzero finite gliders in the source-free bulk.


\subsubsection*{Proposition 8: no finite bulk glider}


Let \(F_{\times3}\) be the base\nobreakdash-6 multiplication-by-\(3\) automaton. If a nonzero finite configuration \(x\) satisfies

\[
F_{\times3}^{\,q}(x)=\sigma^s x,
\tag{29}
\]

then \(q=s=0\).


\paragraph*{Proof}


If \(x\) represents the nonzero rational number \(N\), then (29) implies

\[
3^qN=6^sN,
\]

and hence

\[
3^q=6^s.
\]

Unique factorization forces \(s=0\) from the exponent of \(2\), and then \(q=0\). \(\square\)

Thus a localized carry packet cannot perpetually travel through the empty bulk.


\bigskip



\subsection*{8. Periodic boundary streams are also cycles}


A divergent counterexample cannot hide behind a periodic source schedule.


\subsubsection*{Lemma 9: parity itineraries are injective}


Two positive integers having the same infinite parity itinerary under \(T\) are equal.


\paragraph*{Proof}


Suppose \(x\) and \(y\) follow the same branch word for \(t\) steps. Let \(B_t\) and \(E_t\) be the numbers of odd and even branches in that word. The additive constants cancel when the two affine iterates are subtracted, giving

\[
T^t(x)-T^t(y)
=
\frac{3^{B_t}}{2^{E_t}}(x-y).
\]

The left side is an integer, and \(3^{B_t}\) is odd. Hence

\[
2^{E_t}\mid x-y.
\]

Every odd step is followed by an even step, so \(E_t\to\infty\). Therefore \(x-y\) is divisible by every power of \(2\), and hence \(x=y\). \(\square\)


\subsubsection*{Corollary 10}


If the parity/source sequence \(b_t\) is eventually periodic, then the numerical orbit is eventually periodic.


\paragraph*{Proof}


Suppose the parity sequence is periodic with period \(q\) after time \(r\). Then \(T^r(n)\) and \(T^{r+q}(n)\) have the same infinite parity itinerary. Lemma 9 gives

\[
T^r(n)=T^{r+q}(n).
\]

\(\square\)

Therefore a divergent counterexample must have an aperiodic boundary trace.

For a genuine cycle of period \(q\), equation (24) gives the base\nobreakdash-6 glider equation

\[
\boxed{
(6^E-3^q)n
=
\sum_{k=0}^{q-1}
b_k3^{q-1-k}6^{E_k},
}
\tag{30}
\]

where \(E\) is the number of even steps in the period. The right side is positive, so

\[
6^E>3^q.
\tag{31}
\]

If \(B=q-E\) is the number of odd steps, this is equivalent to

\[
2^E>3^B,
\]

and hence

\[
\boxed{\frac{B}{q}<\log_6 2.}
\tag{32}
\]

Thus the particle system separates the two counterexample mechanisms sharply:

\begin{itemize}
\item a periodic glider/cycle has source density strictly \textbf{below} \(\beta\);
\item a minimal divergent counterexample must have source density essentially \textbf{at least} \(\beta\) in every prefix.
\end{itemize}


\bigskip



\subsection*{9. Why elementary particle Lyapunov functions fail}


The bulk collision is not dissipative.

For example,

\[
(0,1)\longmapsto(1,1)
\]

increases the naive particle number \(c+p\) from \(1\) to \(2\), while

\[
(2,0)\longmapsto(1,0)
\]

decreases it from \(2\) to \(1\).

More generally, there is no nonconstant one-cell additive bulk energy.


\subsubsection*{Proposition 11: no one-cell additive bulk Lyapunov function}


Let \(v:\{0,1,2,3,4,5\}\to\mathbb R\), and on a periodic ring define

\[
V(x)=\sum_jv(x_j).
\]

If \(V\) is nonincreasing under every bulk update, then \(v\) is constant.


\paragraph*{Proof}


The bulk map on a finite ring is reversible. Every orbit on that finite state space is periodic. A quantity that is nonincreasing along a periodic orbit must therefore be constant at every step.

On a ring of length \(1\), the collision permutation (5) implies

\[
v(1)=v(3)=v(4)=v(2).
\tag{33}
\]

On a ring of length \(2\), the configuration with encoded states \((0,3)\) updates to \((1,3)\), so

\[
v(0)+v(3)=v(1)+v(3),
\]

and hence \(v(0)=v(1)\).

Similarly, \((2,5)\) updates to \((2,4)\), giving \(v(5)=v(4)\). Together with (33), all six values are equal. \(\square\)

Consequently, no proof can come from assigning a positive mass to each local particle type and showing that interior collisions consume mass. The bulk is information-preserving.

This is consistent with the known symbolic dynamics. Kari proved that the related reversible base\nobreakdash-6 multiplication automata are extremely pattern-rich: the orbit of every nonuniform finite configuration contains every finite word over the digit alphabet. [\hypertarget{cit:R35:5}{}\hyperlink{ref:R35:5}{5}] Related fractional-multiplication trace systems are nonsofic, and finite-configuration traces need not become eventually periodic. [\hypertarget{cit:R35:6}{}\hyperlink{ref:R35:6}{6}, \hypertarget{cit:R35:7}{}\hyperlink{ref:R35:7}{7}]

Therefore a finite catalogue of local gliders, defects, and forbidden collision blocks cannot by itself establish convergence.


\bigskip



\subsection*{10. The exact missing theorem}


The particle formulation isolates one statement whose proof would finish Collatz.


\subsubsection*{Boundary Overtaking Theorem}


For every \(n\) with at least two base\nobreakdash-6 digits, there exists \(t\ge1\) such that

\[
\boxed{H_t>B_t.}
\tag{34}
\]

By (8), this is equivalent to

\[
\ell(T^t(n))<\ell(n).
\tag{35}
\]


\subsubsection*{Theorem 12: boundary overtaking is equivalent to Collatz}


The following are equivalent:

\begin{enumerate}
\item every positive integer reaches \(1\);
\item every finite base\nobreakdash-6 particle configuration reaches the \(1\to4\to2\to1\) oscillator;
\item the Boundary Overtaking Theorem holds.
\end{enumerate}


\paragraph*{Proof}


The equivalence of (1) and (2) follows from the exact simulation.

Assume (1). If \(\ell(n)\ge2\), the orbit eventually reaches \(1\), whose base\nobreakdash-6 length is \(1\). Hence (35) holds.

Conversely, suppose (3). The integers \(1,2,3,4,5\) are checked directly. Let \(n\ge6\), and assume by strong induction that every positive integer below \(n\) reaches \(1\). By (35), there is \(t\) such that

\[
\ell(T^t(n))<\ell(n).
\]

Since

\[
6^{\ell(n)-1}\le n
\]

and any integer of smaller base\nobreakdash-6 length is less than \(6^{\ell(n)-1}\), one has

\[
T^t(n)<n.
\]

The induction hypothesis now implies that \(T^t(n)\), and hence \(n\), reaches \(1\). \(\square\)

So the desired extinction proof has been reduced to a particularly concrete queue statement:


\begin{quote}
In every finite reversible carry cloud, the high-order sink must eventually remove strictly more cells than the low-order boundary has injected.

\end{quote}


The results above show exactly what a counterexample to that statement would have to look like. It would be an infinite binary source word \(b_0b_1\ldots\) such that

\[
b_tb_{t+1}=0,
\]

\[
B_t\ge H_t
\quad\text{for every }t,
\]

\[
n_t
=
6^{-D_t}
\left(
n_0+\frac13\sum_{k<t}b_k6^{D_k}
\right)
\in\mathbb N,
\]

\[
n_t\equiv b_t\pmod2,
\]

with \(b_t\) non-eventually-periodic and \(\ell_t\to\infty\).

I do not have a valid argument excluding all such aperiodic expanding configurations. Average source density, the absence of localized gliders, reversibility, and the scalar front displacement \(D_t\) do not exclude them. Claiming otherwise would leave precisely the divergence part of Collatz unproved.

The base\nobreakdash-6 programme therefore produces a rigorous and useful reduction, but not a completed proof: \textbf{nontrivial gliders reduce to cycles, while the real unresolved object is an aperiodic critical source-sink cloud.}
\subsection*{References}
\begin{enumerate}[leftmargin=2.2em]
\item \hypertarget{ref:R35:1}{}\hyperlink{cit:R35:1}{\texttt{turn107048view0}}
\item \hypertarget{ref:R35:2}{}\hyperlink{cit:R35:2}{\texttt{turn107048view2}}
\item \hypertarget{ref:R35:3}{}\hyperlink{cit:R35:3}{\texttt{turn865898view1}}
\item \hypertarget{ref:R35:4}{}\hyperlink{cit:R35:4}{\texttt{turn865898view0}}
\item \hypertarget{ref:R35:5}{}\hyperlink{cit:R35:5}{\texttt{turn865898view2}}
\item \hypertarget{ref:R35:6}{}\hyperlink{cit:R35:6}{\texttt{turn865898view3}}
\item \hypertarget{ref:R35:7}{}\hyperlink{cit:R35:7}{\texttt{turn865898view4}}
\end{enumerate}

\clearpage
\section{Round 36: approach 31 --- Rotor-router inverse branching}
\subsection*{Prompt}
\begin{quote}
Complete a full proof of the Collatz conjecture using the following approach:

\begin{enumerate}
\item Rotor-router inverse branching
\end{enumerate}

\textbf{Speculative \textperiodcentered{} P}

Replace the probabilistic inverse-tree model by a rotor-router or deterministic branching process with provably bounded discrepancy from the random model. Then show the actual Collatz inverse tree is sandwiched by such a process and must cover all residue classes.

Take as long as you need.
\end{quote}
\subsection*{Response}

\subsection*{Rotor-router inverse branching: rigorous outcome}


A valid full proof does \textbf{not} result from the proposed route as stated. The rotor part can be formalized and does give a deterministic version of the inverse branching model, but two decisive problems remain:

\begin{enumerate}
\item The actual Collatz inverse tree is not the natural bounded-discrepancy rotor realization of that model.
\item ``Covering all residue classes'' is insufficient. In fact, every Collatz inverse tree already intersects every arithmetic progression, unconditionally.
\end{enumerate}

What follows is a complete derivation of the rotor construction, a finite exact obstruction to the proposed sandwich, and a strengthened quantitative statement that really would imply Collatz.


\subsubsection*{1. Exact inverse tree}


Use the shortcut map

\[
T(n)=
\begin{cases}
n/2,&n\equiv0\pmod2,\\[2mm]
(3n+1)/2,&n\equiv1\pmod2.
\end{cases}
\]

Let

\[
\mathcal P(a)=\{n\ge1:T^k(n)=a\text{ for some }k\ge0\}.
\]

The Collatz conjecture is precisely

\[
\mathcal P(1)=\mathbb N.
\]

The inverse branches are

\[
E(x)=2x,\qquad O(x)=\frac{2x-1}{3}.
\]

The branch \(O(x)\) is integral exactly when \(x\equiv2\pmod3\). After deleting nodes divisible by \(3\), the pruned inverse rule becomes

\[
\widetilde T^{-1}(x)=
\begin{cases}
\{2x\},&x\equiv1,4,5,7\pmod9,\\[2mm]
\left\{2x,\dfrac{2x-1}{3}\right\},
   &x\equiv2,8\pmod9.
\end{cases}
\tag{1}
\]

This is the standard pruned inverse tree used in the stochastic-tree literature. [\hypertarget{cit:R36:1}{}\hyperlink{ref:R36:1}{1}]


\subsubsection*{2. Exact inverse-word arithmetic}


For a word

\[
w=(\varepsilon _1,\ldots,\varepsilon _k)\in\{0,1\}^k,
\]

interpret \(\varepsilon_i=0\) as applying \(E\) and \(\varepsilon_i=1\) as applying \(O\). Put

\[
s_i=\varepsilon_1+\cdots+\varepsilon_i,\qquad s_0=0,
\]

and

\[
C_w=\sum_{i=1}^k
\varepsilon_i\,2^{k-i}3^{s_{i-1}}.
\]

A direct induction gives

\[
B_{\varepsilon_k}\circ\cdots\circ B_{\varepsilon_1}(a)
=
\frac{2^ka-C_w}{3^{s_k}},
\tag{2}
\]

where \(B_0=E\) and \(B_1=O\).

For \(3\nmid a\), the word is a legal path in the pruned inverse tree exactly when

\[
v_3(2^ka-C_w)=s_k.
\tag{3}
\]

Indeed, divisibility by \(3^{s_k}\) gives integrality of every prefix of the path. Requiring the valuation to be exactly \(s_k\) says that the final node is not divisible by \(3\); this also forces every earlier node to be nonzero modulo \(3\), since a node divisible by \(3\) cannot later accept an \(O\)-step.


\paragraph*{Exact averaged leaf count}


Fix \(w\) with \(s=s_k\). As \(a\) ranges over the \(2\cdot3^k\) unit classes modulo \(3^{k+1}\), condition (3) holds for exactly

\[
2\cdot 3^{k-s}
\]

classes:

\begin{itemize}
\item modulo \(3^s\), the congruence \(2^ka\equiv C_w\) has one solution;
\item among its three lifts modulo \(3^{s+1}\), exactly one has valuation greater than \(s\);
\item the other two have valuation exactly \(s\);
\item the remaining \(k-s\) ternary digits are unrestricted.
\end{itemize}

Thus the probability that a uniformly selected unit root admits \(w\) is exactly

\[
3^{-s}.
\]

Summing over all words,

\[
\begin{aligned}
\frac{1}{2\cdot3^k}
\sum_{\substack{a\bmod 3^{k+1}\\3\nmid a}}N_k^*(a)
&=\sum_{w\in\{0,1\}^k}3^{-|w|_1}\\
&=\left(1+\frac13\right)^k\\
&=\left(\frac43\right)^k.
\end{aligned}
\tag{4}
\]

Equivalently,

\[
\sum_{\substack{a\bmod3^{k+1}\\3\nmid a}}N_k^*(a)
=2\cdot4^k.
\]

This recovers the exact inverse-tree average in the literature. The associated stochastic model replaces the hidden ternary lift of a node by an independent uniform choice among its three possible lifts. [\hypertarget{cit:R36:2}{}\hyperlink{ref:R36:2}{2}, \hypertarget{cit:R36:3}{}\hyperlink{ref:R36:3}{3}]

The crucial point is that (4) is an average over roots. It is not a pointwise statement about the fixed root \(a=1\).


\subsubsection*{3. A rigorous rotorization of the stochastic model}


Fix \(j\ge1\), and let

\[
R_j=(\mathbb Z/3^j\mathbb Z)^\times
\]

be the set of possible types.

For a type \(r\in R_j\), its three lifts to modulus \(3^{j+1}\) are

\[
\widetilde r_t=r+t3^j,\qquad t=0,1,2.
\]

Let \(v_t(r)\) be the offspring vector in \(\mathbb Z_{\ge0}^{R_j}\) obtained by applying the exact rule (1) to \(\widetilde r_t\) and reducing the resulting child labels modulo \(3^j\).

The stochastic operator is

\[
P_j(r)=\frac13\sum_{t=0}^2v_t(r).
\tag{5}
\]

Now attach to each type \(r\) a rotor cycling through \(t=0,1,2\). If the rotor at \(r\) is fired \(q\) times, let \(n_t\) be the number of times it selects lift \(t\). Then

\[
n_t\in
\left\{
\left\lfloor\frac q3\right\rfloor,
\left\lceil\frac q3\right\rceil
\right\}
\]

and hence

\[
\left|n_t-\frac q3\right|<1.
\tag{6}
\]

The deterministic offspring vector is

\[
R_q(r)=\sum_{t=0}^2n_tv_t(r).
\]

Because each lift has at most two children,

\[
\begin{aligned}
\left\|R_q(r)-qP_j(r)\right\|_1
&\le
\sum_{t=0}^2
\left|n_t-\frac q3\right|
\|v_t(r)\|_1\\
&<6.
\end{aligned}
\tag{7}
\]

So the finite-type stochastic branching process does possess a deterministic rotor analogue with genuinely bounded \textbf{local} discrepancy.

That completes the first part of the proposed program.


\subsubsection*{4. Local discrepancy is not global discrepancy}


Already in the one-type model, the offspring patterns are

\[
1,\quad 1,\quad 2
\]

in each rotor cycle. If \(Z_k\) is the generation-\(k\) population, then

\[
Z_{k+1}=\frac43 Z_k+e_k,\qquad |e_k|<1.
\tag{8}
\]

But iteration gives

\[
Z_k-\left(\frac43\right)^kZ_0
=
\sum_{\ell=0}^{k-1}
\left(\frac43\right)^{k-1-\ell}e_\ell.
\tag{9}
\]

Consequently, bounded one-step errors can accumulate to order

\[
\left(\frac43\right)^k.
\]

This is a basic difference between a mass-conserving rotor walk and a supercritical branching rotor process. Local rounding errors themselves get copied by later branching.

Thus (7) does not produce an \(O(1)\) discrepancy between the generation populations, much less between the sets of integer labels.

The stochastic models prove almost-sure growth rate \(4/3\) and predict only

\[
I^*(X)=X^{1+o(1)}
\]

for the number of inverse progeny below \(X\). They do not predict that every integer below \(X\) occurs. [\hypertarget{cit:R36:4}{}\hyperlink{ref:R36:4}{4}]

Classical rotor-router results compare hitting frequencies with Markov-chain probabilities under suitable potential-theoretic assumptions. Such results cannot simply be imported into this branching setting; constant single-vertex discrepancy is known to fail on infinite regular trees in general. [\hypertarget{cit:R36:5}{}\hyperlink{ref:R36:5}{5}, \hypertarget{cit:R36:6}{}\hyperlink{ref:R36:6}{6}]


\subsubsection*{5. An exact obstruction in the inverse tree of \(1\)}


The failure of the natural rotor coupling is already visible at depth \(9\).

Starting with \(L_0=\{1\}\) and applying (1), the depth-eight level is

\[
L_8=
\{1,4,5,16,20,26,64,80,85,256\}.
\]

Applying the inverse rule once gives

\[
L_9=
\{2,8,10,13,17,32,40,52,53,128,160,170,512\}.
\]

The nodes of type \(2\pmod3\) are

\[
\{2,8,17,32,53,128,170,512\}.
\]

Their lift counts modulo \(9\) are

\[
\begin{array}{c|ccc}
\text{residue modulo }9&2&5&8\\ \hline
\text{actual count}&2&1&5.
\end{array}
\tag{10}
\]

There are eight such nodes. Any cyclic three-state rotor fired eight consecutive times uses each state either twice or three times. Its count vector must therefore be a permutation of

\[
(3,3,2).
\tag{11}
\]

It cannot reproduce \((2,1,5)\).

More strongly, two phase-shifted unit-discrepancy rotors cannot sandwich the actual lift counts coordinatewise: every rotor upper count is at most \(3\), whereas the actual count in the \(8\pmod9\) lift is \(5\).

This does not rule out every possible discrepancy bound \(C>1\). It does rule out the natural claim that the actual inverse tree is itself a standard rotor realization, or is trapped between two unit-discrepancy rotor processes.

The hidden ternary digits in the actual tree are arithmetically correlated through the constants \(C_w\) in (2); they are not produced by a balanced rotor stack.


\subsubsection*{6. The fatal endpoint: residue coverage is already known}


Suppose ``cover all residue classes'' means

\[
\forall d\ge1,\ \forall a\bmod d,\qquad
\mathcal P(1)\cap(a+d\mathbb Z)\ne\varnothing.
\tag{12}
\]

Statement (12) is already an unconditional theorem.

More generally, for every \(x\) not divisible by \(3\), every positive modulus \(d\), and every residue \(a\bmod d\), there is a positive integer \(y\equiv a\pmod d\) whose forward orbit reaches \(x\). Taking \(x=1\) proves (12). [\hypertarget{cit:R36:7}{}\hyperlink{ref:R36:7}{7}]

Therefore, even in a hypothetical world containing a divergent orbit or a nontrivial cycle, the inverse tree of \(1\) would still meet every residue class modulo every integer.

In topological language, \(\mathcal P(1)\) is already dense in the profinite topology. Profinite density does not imply equality with \(\mathbb N\). A proper cofinite set meets every residue class, and many proper density-one sets do as well.

There are consequently only two possible interpretations:

\begin{itemize}
\item ``At least one member of every residue class occurs.'' This is already proved and is insufficient.
\item ``Every member of every residue class occurs.'' This is simply \(\mathcal P(1)=\mathbb N\), the original conjecture.
\end{itemize}

Thus the proposed final implication is invalid.


\subsubsection*{7. What deterministic branching does prove}


There is nevertheless a simple uniform growth theorem.


\paragraph*{Proposition}


For every \(a\not\equiv0\pmod3\),

\[
N^*_{k+4}(a)\ge 2N^*_k(a).
\tag{13}
\]

Consequently,

\[
N^*_k(a)\ge 2^{\lfloor k/4\rfloor}.
\tag{14}
\]


\paragraph*{Proof}


Under the always-available branch \(E(x)=2x\), the unit residues modulo \(9\) move as

\[
1\to2,\qquad
4\to8,\qquad
5\to1\to2,\qquad
7\to5\to1\to2.
\]

The residues \(2\) and \(8\) are branching residues. Therefore, starting from any unit residue, after at most three \(E\)-moves one reaches a node having two children.

After taking the two children, extend each by \(E\)-moves until exactly four inverse steps have been used. This produces two distinct depth-four descendants.

Descendant sets from two different nodes at the same level are disjoint: if a number were a four-step descendant of both \(x\) and \(x'\), applying \(T^4\) would give both \(x\) and \(x'\). Hence \(x=x'\).

This proves (13), and iteration proves (14). \ensuremath{\blacksquare}

Much stronger universal exponential leaf-growth estimates are known\textemdash{}for sufficiently large \(k\), every pruned tree has at least \(1.302053^k\) leaves\textemdash{}but even those do not prove that the integer labels cover \(\mathbb N\). [\hypertarget{cit:R36:8}{}\hyperlink{ref:R36:8}{8}]

The issue is not producing enough branches. It is placing those branches at the right integer labels.


\subsubsection*{8. The archimedean condition missing from residue coverage}


A sufficient strengthening of (12) would be:


\begin{quote}
For arbitrarily large \(m\), every nonzero residue \(1\le r<m\) has a representative   \[
> y\in\mathcal P(1),\qquad
> y\equiv r\pmod m,\qquad
> 1\le y<m.
> \tag{15}
> \]

\end{quote}


But (15) immediately gives \(y=r\), since two positive numbers below \(m\) that are congruent modulo \(m\) are equal. Hence

\[
\{1,\ldots,m-1\}\subseteq\mathcal P(1).
\]

An unbounded sequence of such moduli would prove Collatz.

So the missing statement is not modular coverage but \textbf{height-controlled modular coverage}. The rotor model records branch frequencies and residues, whereas (15) couples the residue to the ordinary size of its representative.

The known arithmetic-progression theorem gives a bounded path length depending on the modulus, but it does not give \(y<m\); the resulting representative may be vastly larger than the modulus. [\hypertarget{cit:R36:9}{}\hyperlink{ref:R36:9}{9}]


\subsubsection*{9. A quantitative rotor target that really would prove Collatz}


There is a more plausible completion criterion.

Set

\[
A(X)=\#\bigl(\mathcal P(1)\cap[1,X]\bigr),\qquad
E(X)=X-A(X).
\]

Thus \(E(X)\) counts the possible exceptions up to \(X\).


\paragraph*{Power-saving completion lemma}


If

\[
E(X)=o(X^{0.84}),
\tag{16}
\]

then the Collatz conjecture holds.

In particular, a bound

\[
E(X)=O(X^\theta)
\qquad\text{for some }\theta<0.84
\tag{17}
\]

would be a complete proof.


\paragraph*{Proof}


Assume there is an exception \(b\notin\mathcal P(1)\). Its forward orbit remains outside \(\mathcal P(1)\).

After finitely many steps the orbit reaches some positive \(a\not\equiv0\pmod3\). Indeed, repeated halving eventually turns a multiple of \(3\) into an odd multiple of \(3\), after which

\[
T(a)=\frac{3a+1}{2}\not\equiv0\pmod3.
\]

Every inverse iterate of this new \(a\) is also an exception: if \(T^k(n)=a\) and \(n\) reached \(1\), then \(a\) would reach \(1\) as well.

The established inverse-iterate lower bound says that, for every fixed \(a\not\equiv0\pmod3\),

\[
\#\{n\le X:T^k(n)=a\text{ for some }k\}
\ge X^{0.84}
\]

for all sufficiently large \(X\). [\hypertarget{cit:R36:10}{}\hyperlink{ref:R36:10}{10}]

All these inverse iterates lie in the exceptional set. Hence

\[
E(X)\ge X^{0.84}
\]

for sufficiently large \(X\), contradicting (16). \ensuremath{\blacksquare}

This gives a precise goal for a repaired rotor-router program:

\[
\boxed{
X-\#\bigl(\mathcal P(1)\cap[1,X]\bigr)
=O(X^\theta)
\quad\text{for some }\theta<0.84.
}
\tag{18}
\]

A rotor comparison proving only

\[
\#\bigl(\mathcal P(1)\cap[1,X]\bigr)=X^{1+o(1)}
\]

does not imply (18). For example, \(X/\log X=X^{1+o(1)}\), while it misses almost all integers.


\subsubsection*{10. A smaller strongly sufficient target}


There is also an exact residue-class reduction that is stronger than merely hitting every class.

Every positive Collatz orbit contains a number congruent to \(2\pmod9\). Consequently, it would suffice to prove

\[
\{n\ge1:n\equiv2\pmod9\}\subseteq\mathcal P(1).
\tag{19}
\]

This is an entire progression inclusion, not merely one inverse-tree hit in that progression. [\hypertarget{cit:R36:11}{}\hyperlink{ref:R36:11}{11}]

Again, the distinction is decisive:

\[
\mathcal P(1)\cap(2+9\mathbb Z)\ne\varnothing
\]

is trivial and already known, while (19) would settle the conjecture.


\subsection*{Conclusion}


The rotor-router proposal has a rigorous core:

\begin{itemize}
\item the stochastic inverse-tree model can be replaced by a finite-type deterministic branching network;
\item each rotor has uniformly bounded local discrepancy;
\item the resulting process reproduces the \(4/3\) branching scale;
\item deterministic inverse branching guarantees uniform exponential leaf growth.
\end{itemize}

But the intended completion fails at two exact places:

\[
\text{local rotor balance}
\;\not\Longrightarrow\;
\text{actual Collatz lift balance},
\]

as witnessed by the depth-nine count \((2,1,5)\), and

\[
\text{all residue classes are hit}
\;\not\Longrightarrow\;
\mathcal P(1)=\mathbb N,
\]

because universal residue-class hitting is already unconditional.

A successful rotor-router proof would therefore need a substantially stronger discrepancy theorem\textemdash{}one controlling \textbf{ordinary interval occupancy or representative height}, not merely branching frequencies or modular distributions. The concrete quantitative threshold is (18): a power-saving exceptional-set bound with exponent below \(0.84\) would turn the rotor comparison into a complete proof.
\subsection*{References}
\begin{enumerate}[leftmargin=2.2em]
\item \hypertarget{ref:R36:1}{}\hyperlink{cit:R36:1}{\texttt{turn639386view2}}
\item \hypertarget{ref:R36:2}{}\hyperlink{cit:R36:2}{\texttt{turn639386view3}}
\item \hypertarget{ref:R36:3}{}\hyperlink{cit:R36:3}{\texttt{turn312898view0}}
\item \hypertarget{ref:R36:4}{}\hyperlink{cit:R36:4}{\texttt{turn680517view1}}
\item \hypertarget{ref:R36:5}{}\hyperlink{cit:R36:5}{\texttt{turn639386view5}}
\item \hypertarget{ref:R36:6}{}\hyperlink{cit:R36:6}{\texttt{turn639386view6}}
\item \hypertarget{ref:R36:7}{}\hyperlink{cit:R36:7}{\texttt{turn639386view0}}
\item \hypertarget{ref:R36:8}{}\hyperlink{cit:R36:8}{\texttt{turn923363view0}}
\item \hypertarget{ref:R36:9}{}\hyperlink{cit:R36:9}{\texttt{turn397116view1}}
\item \hypertarget{ref:R36:10}{}\hyperlink{cit:R36:10}{\texttt{turn680517view3}}
\item \hypertarget{ref:R36:11}{}\hyperlink{cit:R36:11}{\texttt{turn639386view1}}
\end{enumerate}

\clearpage
\section{Round 37: approach 32 --- Electrical capacity of infinity}
\subsection*{Prompt}
\begin{quote}
Complete a full proof of the Collatz conjecture using the following approach:

\begin{enumerate}
\item Electrical capacity of infinity
\end{enumerate}

\textbf{Speculative \textperiodcentered{} P}

Assign conductances to inverse-tree edges based on their multiplicative expansion. Prove that infinity has zero effective capacity relative to the component of 1, forcing recurrence toward the root. The challenge is converting a directed arithmetic graph into a network whose recurrence has pointwise dynamical meaning.

Take as long as you need.
\end{quote}
\subsection*{Response}

\subsection*{Result of the electrical-capacity approach}


The proposed argument does \textbf{not} yield a valid full proof of the Collatz conjecture. The electrical calculation can be completed rigorously\textemdash{}and gives a surprisingly clean zero-capacity theorem\textemdash{}but the intended implication

\[
\text{zero capacity at infinity}
\quad\Longrightarrow\quad
\text{every directed orbit returns to the root}
\]

is false.

There are actually two complementary obstructions:

\begin{enumerate}
\item Conductances depending only on the local multiplicative expansion give \textbf{positive} capacity because the inverse tree contains a bounded-label binary subtree.
\item A natural cumulative multiplicative conductance does give \textbf{zero} capacity, but it does so for essentially every rooted Collatz inverse tree, including inverse trees that could be attached to a hypothetical nontrivial cycle or divergent component.
\end{enumerate}

The second result is still a rigorous and potentially useful theorem.


\bigskip



\subsubsection*{1. Odd-only Collatz graph}


Work with the Syracuse map on positive odd integers,

\[
U(n)=\frac{3n+1}{2^{a(n)}},
\qquad
a(n)=\nu _2(3n+1).
\]

The Collatz conjecture is equivalent to the assertion that every positive odd integer eventually reaches \(1\) under \(U\).

For an odd integer \(m\), its inverse children have the form

\[
n_a(m)=\frac{2^a m-1}{3}.
\]

This is a positive odd integer precisely when

\[
2^a m\equiv 1\pmod 3.
\]

Thus:

\[
m\equiv1\pmod3\quad\Longrightarrow\quad a\ \text{even},
\]

\[
m\equiv2\pmod3\quad\Longrightarrow\quad a\ \text{odd},
\]

while \(m\equiv0\pmod3\) has no inverse children.

Let \(\mathcal T_1\) be the inverse tree consisting of odd integers already known, by definition, to eventually reach \(1\). The self-loop \(U(1)=1\) is removed.

For a conductance network \((\mathcal T_1,c)\), write

\[
\mathcal E(f)
   =\sum_{\{v,p(v)\}}
      c(v,p(v))\bigl(f(v)-f(p(v))\bigr)^2
\]

and

\[
\operatorname{Cap}(1,\infty)
 =
 \inf\left\{
       \mathcal E(f):
       f(1)=1,\;
       f\text{ finitely supported}
     \right\}.
\]

This is the effective conductance from \(1\) to infinity. Zero capacity corresponds to infinite effective resistance and recurrence of the associated \textbf{reversible random walk}. Conversely, a finite-energy unit flow to infinity proves positive capacity and transience. [\hypertarget{cit:R37:1}{}\hyperlink{ref:R37:1}{1}, \hypertarget{cit:R37:2}{}\hyperlink{ref:R37:2}{2}]


\bigskip



\subsection*{2. Local multiplicative conductances cannot work}


Suppose the conductance of an inverse edge depends only on its exponent:

\[
c\bigl(n_a(m),m\bigr)=\psi(a),
\qquad
\psi(a)>0.
\]

This includes choices such as

\[
\psi(a)=\frac{3}{2^a},
\qquad
\psi(a)=\frac{2^a}{3},
\]

and arbitrary positive functions of the local multiplicative expansion.


\subsubsection*{Proposition 1: bounded-label binary subtree}


The inverse tree below \(5\) contains a full binary subtree all of whose edge exponents lie in

\[
\{1,2,3,4,5,6\}.
\]


\paragraph*{Proof}


Let \(m\not\equiv0\pmod3\). An inverse child \(n_a(m)\) is divisible by \(3\) exactly when

\[
2^a m\equiv1\pmod9.
\]

Since \(2\) generates the unit group \((\mathbb Z/9\mathbb Z)^\times\), which has order \(6\), there is exactly one residue class \(a\pmod6\) satisfying this congruence.

If \(m\equiv1\pmod3\), the allowed exponents are even. Their three classes modulo \(6\) have smallest positive representatives

\[
2,\ 4,\ 6.
\]

Exactly one produces a child divisible by \(3\), so the other two produce children not divisible by \(3\).

If \(m\equiv2\pmod3\), the allowed exponents are odd, with representatives

\[
1,\ 3,\ 5.
\]

Again, exactly one produces a multiple of \(3\), and the other two produce nonmultiples of \(3\).

A child divisible by \(3\) has no inverse descendants, while either of the other two can be continued. Starting at

\[
5=\frac{2^4\cdot1-1}{3},
\]

and repeatedly retaining the two nonmultiples of \(3\), gives a full binary subtree with exponents at most \(6\). \ensuremath{\blacksquare}

Set

\[
c_*=\min_{1\leq a\leq6}\psi(a)>0.
\]

Send one unit of current through the edge \(1\!-\!5\), and then split it equally at every vertex of the binary subtree. At binary level \(j\), each of the \(2^j\) edges carries current \(2^{-j}\), so its total energy is at most

\[
2^j\left(2^{-j}\right)^2\frac1{c_*}
   =\frac{2^{-j}}{c_*}.
\]

Therefore the whole binary subtree has energy at most

\[
\sum_{j=1}^{\infty}\frac{2^{-j}}{c_*}
 =\frac1{c_*}.
\]

Including the edge \(1\!-\!5\), whose exponent is \(4\),

\[
R_{\mathrm{eff}}(1,\infty)
\leq
\frac1{\psi(4)}+\frac1{c_*}<\infty.
\]

Thus

\[
\boxed{\operatorname{Cap}(1,\infty)>0}
\]

for every strictly positive conductance depending only on the local exponent.

For the especially natural choice

\[
\psi(a)=\frac{3}{2^a},
\]

we have

\[
c_*=\frac3{64},
\qquad
\frac1{\psi(4)}=\frac{16}{3},
\]

and hence

\[
R_{\mathrm{eff}}\leq\frac{80}{3},
\qquad
\operatorname{Cap}(1,\infty)\geq\frac3{80}.
\]

So the na\"ive expansion-weighted network is definitely transient.


\bigskip



\subsection*{3. A cumulative derivative conductance has zero capacity}


The local rule fails because it resets the expansion factor at every edge. A more natural nonlocal choice is to use the cumulative multiplicative derivative.

Let

\[
1=x_0\leftarrow x_1\leftarrow\cdots\leftarrow x_d=v
\]

be the unique inverse path in \(\mathcal T_1\), with

\[
x_i=\frac{2^{a_i}x_{i-1}-1}{3}.
\]

Set

\[
A_i=a_1+\cdots+a_i
\]

and define

\[
Q(x_i)=\frac{3^i}{2^{A_i}},
\qquad
Q(1)=1.
\]

Thus

\[
\frac{Q(x_i)}{Q(x_{i-1})}
   =\frac{3}{2^{a_i}}.
\]

Assign the edge \(x_i-x_{i-1}\) the conductance

\[
\boxed{c(x_i,x_{i-1})=Q(x_i)}.
\]

This conductance remembers the whole inverse path.


\subsubsection*{The fundamental telescoping identity}


From

\[
x_i=\frac{2^{a_i}x_{i-1}-1}{3}
\]

we obtain

\[
\begin{aligned}
Q(x_i)x_i
&=
Q(x_{i-1})\frac{3}{2^{a_i}}
  \frac{2^{a_i}x_{i-1}-1}{3} \\[2mm]
&=
Q(x_{i-1})x_{i-1}
 -
\frac{Q(x_{i-1})}{2^{a_i}} \\[2mm]
&=
Q(x_{i-1})x_{i-1}
 -
\frac{Q(x_i)}{3}.
\end{aligned}
\]

Telescoping along the path gives

\[
\boxed{
Q(v)v
=
1-\frac13\sum_{i=1}^{d}Q(x_i).
}
\tag{1}
\]

In particular,

\[
0<Q(v)v\leq1,
\]

and therefore

\[
\boxed{Q(v)\leq\frac1v.}
\tag{2}
\]

Identity (1) is the main arithmetic fact behind the capacity calculation.

Notice also that the weighted degree is finite. For a parent \(m\equiv1\pmod3\),

\[
\sum_{\substack{a\geq2\\a\ {\rm even}}}
  \frac{3}{2^a}
=1,
\]

and for \(m\equiv2\pmod3\),

\[
\sum_{\substack{a\geq1\\a\ {\rm odd}}}
  \frac{3}{2^a}
=2.
\]

Thus the sum of child conductances is at most \(2Q(m)\).


\bigskip



\subsubsection*{Theorem 2: exact zero-capacity theorem}


For the conductance

\[
c(v,p(v))=Q(v),
\]

one has

\[
\boxed{\operatorname{Cap}_{\mathcal T_1}(1,\infty)=0.}
\]


\paragraph*{Proof}


Fix

\[
0<\beta\leq\frac12
\]

and define

\[
f_\beta(v)=Q(v)^\beta.
\]

Then \(f_\beta(1)=1\). Let \(p=p(v)\), and write

\[
t=\frac{Q(v)}{Q(p)}
  =\frac{3}{2^{a(v)}}.
\]

The energy of the edge \(v-p\) is

\[
\begin{aligned}
Q(v)\bigl(f_\beta(v)-f_\beta(p)\bigr)^2
&=
Q(p)t
 \left(
 Q(p)^\beta t^\beta-Q(p)^\beta
 \right)^2\\
&=
Q(p)^{1+2\beta}
 t(t^\beta-1)^2.
\end{aligned}
\tag{3}
\]

For \(0<\beta\leq1/2\),

\[
t^\beta-1
=
\int_0^\beta t^s\log t\,ds.
\]

Since \(0<t\leq3/2\),

\[
|t^s|\leq\sqrt{\frac32}
\]

for \(0\leq s\leq1/2\). Also, when \(t=3/2^a\),

\[
|\log t|
 \leq C_0(1+a)
\]

for an absolute constant \(C_0\). Hence

\[
|t^\beta-1|
 \leq C_1\beta(1+a).
\]

Consequently,

\[
t(t^\beta-1)^2
\leq
C_2\beta^2\,2^{-a}(1+a^2).
\]

Summing over all possible exponents gives a constant \(K<\infty\), independent of the parent and of \(\beta\), such that

\[
\sum_{v:p(v)=p}
t(v)\bigl(t(v)^\beta-1\bigr)^2
\leq K\beta^2.
\]

Using (3),

\[
\mathcal E(f_\beta)
\leq
K\beta^2
\sum_{p\in\mathcal T_1}Q(p)^{1+2\beta}.
\]

By (2),

\[
\sum_{p\in\mathcal T_1}Q(p)^{1+2\beta}
\leq
\sum_{n=1}^{\infty}\frac1{n^{1+2\beta}}
=
\zeta(1+2\beta).
\]

The integral estimate gives

\[
\zeta(1+2\beta)
\leq
1+\frac1{2\beta},
\]

so

\[
\mathcal E(f_\beta)
\leq
K\beta^2\left(1+\frac1{2\beta}\right)
=
K\left(\beta^2+\frac{\beta}{2}\right).
\]

Therefore

\[
\mathcal E(f_\beta)\longrightarrow0
\qquad(\beta\downarrow0).
\]

The function \(f_\beta\) is not finitely supported, so introduce the truncation

\[
g_\beta(v)
=
2\left(f_\beta(v)-\frac12\right)_+.
\]

Then \(g_\beta(1)=1\). Moreover, \(g_\beta(v)>0\) implies

\[
Q(v)^\beta>\frac12,
\]

hence

\[
Q(v)>2^{-1/\beta}.
\]

Using \(Q(v)\leq1/v\), this forces

\[
v<2^{1/\beta}.
\]

Thus \(g_\beta\) has finite support.

The map

\[
s\longmapsto2(s-\tfrac12)_+
\]

is \(2\)-Lipschitz, so

\[
\mathcal E(g_\beta)
\leq4\mathcal E(f_\beta)
\longrightarrow0.
\]

Therefore the infimum defining the capacity is zero:

\[
\operatorname{Cap}_{\mathcal T_1}(1,\infty)=0.
\qquad\Box
\]


\bigskip



\subsection*{4. Why this is not a Collatz proof}



\subsubsection*{The capacity only sees the component already attached to \(1\)}


The network \(\mathcal T_1\) contains exactly those odd integers whose orbits are already in the component of \(1\). A hypothetical divergent orbit or nontrivial cycle would belong to another connected component and would contribute no edge whatsoever to the effective capacity measured from \(1\).

In the ordinary Collatz graph, two integers merge exactly when they lie in the same connected component, and the Collatz conjecture is equivalent to the assertion that this graph is connected. [\hypertarget{cit:R37:3}{}\hyperlink{ref:R37:3}{3}]

Thus proving

\[
\operatorname{Cap}_{\mathcal T_1}(1,\infty)=0
\]

does not prove

\[
\mathcal T_1=\{\text{all positive odd integers}\}.
\]

Abstractly, one can adjoin a completely disjoint directed ray

\[
y_0\longrightarrow y_1\longrightarrow y_2\longrightarrow\cdots
\]

to \(\mathcal T_1\). The effective capacity from \(1\) remains unchanged, although the new directed graph has an escaping orbit.

The same logical blindness occurs for any network constructed only from Collatz edges.


\bigskip



\subsubsection*{Electrical recurrence is random-walk recurrence, not orbit recurrence}


Consider the half-line

\[
0-1-2-3-\cdots
\]

with unit conductances. Its resistance to infinity is

\[
1+1+1+\cdots=\infty,
\]

so its effective capacity is zero and the associated symmetric random walk is recurrent.

But the deterministic map

\[
D(k)=k+1
\]

escapes to infinity from every starting point.

Hence even on one connected graph,

\[
\text{zero electrical capacity}
\]

does not imply that a prescribed orientation or deterministic path moves toward the root. The electrical network forgets which direction is the Collatz direction.

This is not a technical gap; it is a counterexample to the proposed bridge.


\bigskip



\subsection*{5. The zero-capacity theorem is componentwise universal}


The preceding proof was not special to \(1\).

Let \(m\) be any odd integer and consider an acyclic first-entry inverse tree rooted at \(m\). Set \(Q(m)=1\), and define \(Q\) by the same multiplicative recursion. Along

\[
m=x_0\leftarrow x_1\leftarrow\cdots\leftarrow x_d=v
\]

the same calculation gives

\[
\boxed{
Q(v)v
=
m-\frac13\sum_{i=1}^{d}Q(x_i).
}
\tag{4}
\]

Consequently,

\[
Q(v)\leq\frac{m}{v}.
\]

Repeating the energy argument gives

\[
\operatorname{Cap}_{\mathcal T_m}(m,\infty)=0.
\]

Thus the zero-capacity result applies to every attached inverse tree, not just the inverse tree of \(1\).

In particular, suppose hypothetically that there were a nontrivial finite cycle. Delete the finitely many cycle edges. The component becomes a finite collection of rooted inverse trees attached to the cycle vertices. Each such inverse tree has zero capacity by (4). Restoring the finite cycle does not turn those zero-capacity ends into positive-capacity ends.

Therefore:

\[
\boxed{\text{A nontrivial Collatz cycle is completely compatible with this zero-capacity theorem.}}
\]

This is unavoidable: a capacity-at-infinity argument has no direct way to detect a finite cycle. The classical decomposition of the Collatz problem likewise requires excluding both divergent orbits and nontrivial cycles. [\hyperlink{ref:R37:3}{3}]


\bigskip



\subsection*{6. What happens on a hypothetical divergent component?}


The cumulative conductance can also be extended over a cycle-free component.

Let

\[
x_{k+1}=U(x_k),
\qquad
a_k=\nu_2(3x_k+1),
\qquad
A_k=\sum_{j=0}^{k-1}a_j,
\]

and define

\[
q_k=\frac{3^k}{2^{A_k}}.
\]

Normalize \(Q(x_0)=1\). Since moving forward reverses the inverse-cocycle relation,

\[
Q(x_k)=\frac{2^{A_k}}{3^k}
      =\frac1{q_k}.
\]

The conductance of the forward-spine edge \(x_k-x_{k+1}\) is therefore

\[
c_k=\frac1{q_k},
\]

so its resistance is

\[
r_k=q_k.
\]

All inverse bushes attached to the spine are electrically dangling: they connect to the spine at only one vertex and carry no current toward the forward end. Hence the resistance to the forward end is exactly

\[
\boxed{
R_{\mathrm{forward}}
=
\sum_{k=0}^{\infty}q_k.
}
\tag{5}
\]

Thus this weighting would assign \textbf{positive} capacity to a divergent orbit only after proving

\[
\sum_{k=0}^{\infty}
\frac{3^k}{2^{A_k}}
<\infty.
\tag{6}
\]

But divergence of \(x_k\) does not provide (6). Indeed, the exact affine identity is

\[
\boxed{
x_k
=
q_k
\left(
x_0+\frac13\sum_{j=0}^{k-1}\frac1{q_j}
\right).
}
\tag{7}
\]

Equation (7) allows growth to arise through \(q_k\), through the accumulated affine term, or through an interaction between them. It gives no implication that \(\sum q_k\) converges. In the expansion-dominated regime \(q_k\) becomes large, making the resistance sum in (5) even more divergent.

So the cumulative conductance that makes the inverse tree recurrent has the wrong orientation for detecting a forward-growing trajectory.


\bigskip



\subsection*{7. Congruence density does not repair the bridge}


There is a strong known backtracing theorem: every arithmetic progression is a sufficient set. Equivalently, each Collatz component meets every arithmetic progression; more quantitatively, one can backtrace from a node into any prescribed congruence class. [\hypertarget{cit:R37:4}{}\hyperlink{ref:R37:4}{4}, \hypertarget{cit:R37:5}{}\hyperlink{ref:R37:5}{5}]

This means that distinct hypothetical components would each intersect every residue class modulo every integer. Nevertheless, it does not connect those components.

Two disjoint sets can both intersect every congruence class. Electrically, if there is no actual edge between them, no current passes between them. Congruence density supplies finite approximations to another component, but not an actual finite Collatz path joining the components.

To turn congruence approximation into equality, one would need an additional archimedean tightness theorem\textemdash{}for example, representatives congruent to a fixed \(n\) modulo arbitrarily large moduli whose sizes remain below those moduli. Such a bound would itself essentially force the representative eventually to equal \(n\), and is comparable in strength to the missing Collatz assertion.


\bigskip



\subsection*{8. Exact verdict}


The approach produces the following rigorous conclusions:

\[
\boxed{
\begin{array}{l}
\text{Every strictly positive local expansion conductance gives positive capacity;}\\[1mm]
\text{the cumulative derivative conductance }c(v,p)=Q(v)\text{ gives zero capacity;}\\[1mm]
\text{but this zero-capacity theorem holds componentwise and has no}\\
\text{pointwise implication for the directed Collatz orbit.}
\end{array}}
\]

A genuine electrical proof would require a new theorem of the form

\[
\boxed{
\text{divergent orbit or nontrivial cycle}
\quad\Longrightarrow\quad
\text{positive-capacity structure electrically connected to }1.
}
\tag{8}
\]

For a network containing only Collatz edges, (8) is impossible as stated: a bad component is disconnected from \(1\), and a finite cycle does not meet infinity. One would have to add non-Collatz ``horizontal'' arithmetic edges or use a genuinely directed, non-reversible notion of capacity. One would then need to prove that the added structure preserves exact dynamical meaning rather than merely replacing the conjecture with a hidden connectivity assumption.

Therefore the electrical-capacity proposal, in its stated form, cannot be completed into a full Collatz proof. Its strongest valid output is the exact cumulative-conductance zero-capacity theorem above, together with the local-conductance and pointwise-bridge no-go results.
\subsection*{References}
\begin{enumerate}[leftmargin=2.2em]
\item \hypertarget{ref:R37:1}{}\hyperlink{cit:R37:1}{\texttt{turn522255view0}}
\item \hypertarget{ref:R37:2}{}\hyperlink{cit:R37:2}{\texttt{turn522255view1}}
\item \hypertarget{ref:R37:3}{}\hyperlink{cit:R37:3}{\texttt{turn522255view3}}
\item \hypertarget{ref:R37:4}{}\hyperlink{cit:R37:4}{\texttt{turn522255view4}}
\item \hypertarget{ref:R37:5}{}\hyperlink{cit:R37:5}{\texttt{turn468309view0}}
\end{enumerate}

\clearpage
\section{Round 38: approach 33 --- Sandpile stabilization}
\subsection*{Prompt}
\begin{quote}
Complete a full proof of the Collatz conjecture using the following approach:

\begin{enumerate}
\item Sandpile stabilization
\end{enumerate}

\textbf{Speculative \textperiodcentered{} P}

Encode binary carries, ternary carries, and valuation deficits as chips on a graph. A Collatz step becomes a legal firing sequence. Prove that every finite configuration stabilizes to the trivial recurrent state.

Take as long as you need.
\end{quote}
\subsection*{Response}

\subsection*{Outcome}


I worked the sandpile construction through in an exact form. It gives a genuinely useful reformulation, including an abelian carrier system, an exact odometer identity, and a critical-group formulation of possible cycles.

It does \textbf{not}, however, yield a completed Collatz proof. After all ordinary binary and ternary carries have been eliminated, the remaining stabilization assertion is

\[
\sum_{j<k}\bigl(v_2(3n_j+1)-\beta_j\bigr)
   =\lfloor \log_2 n_0\rfloor
\]

for some \(k\), where \(\beta_j\in\{1,2\}\) is the leading binary carry. This statement is exactly equivalent to the Collatz conjecture. I do not have a valid proof of it, and treating ordinary sandpile stabilization as though it implied this identity would be circular.

What follows is the complete derivation and the precise obstruction.


\bigskip



\subsubsection*{1. The exact carrier sandpile}


Use the shortcut Collatz map

\[
T(n)=
\begin{cases}
n/2,&n\equiv0\pmod 2,\\[2mm]
(3n+1)/2,&n\equiv1\pmod 2.
\end{cases}
\]

Consider a finite ordered path whose sites have radix \(2\) or radix \(3\). A radix-\(2\) site carries a stable digit \(b\in\{0,1\}\), and a radix-\(3\) site carries a stable digit \(q\in\{0,1,2\}\). Starting from an implicit leading digit \(1\), the path represents the corresponding mixed-radix integer.

Write

\[
\mathtt f=0_2,\qquad \mathtt t=1_2,
\]

and write \(\mathtt0,\mathtt1,\mathtt2\) for the ternary digits.


\paragraph*{Internal firing}


Whenever a binary site lies immediately to the left of a ternary site, fire the edge by exchanging their radices. The new digits are uniquely defined by

\[
3b+q=2q'+b',
\qquad q'\in\{0,1,2\},\quad b'\in\{0,1\}.
\]

This gives

\[
\begin{array}{lll}
\mathtt{f0}\longrightarrow\mathtt{0f},&
\mathtt{f1}\longrightarrow\mathtt{0t},&
\mathtt{f2}\longrightarrow\mathtt{1f},\\[1mm]
\mathtt{t0}\longrightarrow\mathtt{1t},&
\mathtt{t1}\longrightarrow\mathtt{2f},&
\mathtt{t2}\longrightarrow\mathtt{2t}.
\end{array}
\]

Each firing preserves the represented integer.

The swap is a canonical radix-conversion map. More generally, for radices \(p,q\), the relation

\[
qa+b=pc+d
\]

uniquely converts a \(p\)-digit followed by a \(q\)-digit into a \(q\)-digit followed by a \(p\)-digit. Consequently, different orders of adjacent radix swaps leading to the same final radix order give the same final digits: this is a Yang\textendash{}Baxter-type consistency law following from uniqueness of mixed-radix expansion.


\paragraph*{Right-boundary firing}


Let \texttt{/} denote the implicit leading \(1\), and let \texttt{.} mark the least-significant end. At the right boundary,

\[
\mathtt{f.}\longrightarrow\mathtt{.},
\qquad
\mathtt{t.}\longrightarrow\mathtt{2.}.
\]

Indeed, if the preceding prefix represents \(x\), then

\[
T(2x)=x,\qquad T(2x+1)=3x+2.
\]

Thus the first rule absorbs an even binary site, while the second reflects a terminal binary \(1\) as a left-moving ternary carrier.


\paragraph*{Left-boundary firing}


When a ternary carrier reaches the leading boundary,

\[
\mathtt{/0}\longrightarrow\mathtt{/t},\qquad
\mathtt{/1}\longrightarrow\mathtt{/ff},\qquad
\mathtt{/2}\longrightarrow\mathtt{/ft}.
\]

These are simply

\[
3=(11)_2,\qquad 4=(100)_2,\qquad 5=(101)_2.
\]

The carrier is therefore absorbed at the left boundary and emits either one or two binary sites.

These are exactly the eleven mixed binary\textendash{}ternary rewriting rules introduced by Yolcu, Aaronson, and Heule, who proved that termination of the resulting system is equivalent to Collatz convergence. [\hypertarget{cit:R38:1}{}\hyperlink{ref:R38:1}{1}, \hypertarget{cit:R38:2}{}\hyperlink{ref:R38:2}{2}]


\bigskip



\subsubsection*{2. Why this system is abelian}


There is a particularly strong local commutation property.

Every left side has length two. Its first character lies in

\[
\{\mathtt{/},\mathtt f,\mathtt t\},
\]

while its second character lies in

\[
\{\mathtt{.},\mathtt0,\mathtt1,\mathtt2\}.
\]

These sets are disjoint. Therefore two legal redexes can never overlap in one site. Distinct simultaneously legal firings act on disjoint substrings, and applying them in either order produces exactly the same word.

Hence the system has the strong diamond property

\[
x\longrightarrow y,\quad x\longrightarrow z
\quad\Longrightarrow\quad
y\longrightarrow w,\quad z\longrightarrow w
\]

whenever \(y\ne z\). This is the appropriate abelianity statement for the carrier system.

It is nevertheless important to distinguish it from a classical sandpile on a fixed finite graph:

\begin{itemize}
\item the sites are ordered;
\item their radix labels are part of the state;
\item the left boundary may create a new site;
\item the path length is unbounded.
\end{itemize}

Thus it is an abelian local rewriting system, or a generalized carrier sandpile, rather than a fixed-Laplacian chip-firing game.


\paragraph*{Auxiliary firings always terminate}


Call a binary-before-ternary pair an inversion. Let

\[
N_3=\text{number of ternary sites},
\]

and

\[
I=\#\{(i,j):i<j,\ r_i=2,\ r_j=3\}.
\]

An internal swap preserves \(N_3\) and reduces \(I\) by one. A left-boundary firing reduces \(N_3\) by one. Hence the auxiliary firings terminate by lexicographic descent of

\[
(N_3,I)\in\mathbb N^2.
\]

Consequently, an infinite firing sequence would have to contain infinitely many right-boundary firings.


\paragraph*{The only stable admissible state is \(1\)}


Suppose a canonical word has a nonempty interior.

If it has no ternary sites, its final site is binary, so a right-boundary rule applies. If it has a ternary site, consider the leftmost one. It is either adjacent to \texttt{/}, in which case a left-boundary rule applies, or it is preceded by a binary site, in which case an internal swap applies.

Therefore the only stable canonical word is

\[
\mathtt{/.},
\]

representing \(1\).

Since auxiliary firings preserve the integer and every right-boundary firing applies \(T\), an infinite firing execution exists exactly when the corresponding Collatz trajectory never reaches \(1\). Thus abelianity proves order independence, but not the existence of a finite complete execution.


\bigskip



\subsubsection*{3. The exact sandpile odometer}


It is most informative to group the firings between successive odd integers.

Let

\[
n_{j+1}=U(n_j)
       :=\frac{3n_j+1}{2^{a_j}},
\qquad
a_j=v_2(3n_j+1),
\]

where every \(n_j\) is odd. Define the binary path length

\[
\ell_j=\lfloor\log_2 n_j\rfloor.
\]

This is exactly the number of binary sites following the implicit leading \(1\).

Define the most-significant carry

\[
\beta_j
  =\lfloor\log_2(3n_j+1)\rfloor-\lfloor\log_2 n_j\rfloor.
\]

Because

\[
2^{\ell_j}\le n_j<2^{\ell_j+1},
\]

we have

\[
2^{\ell_j+1}<3n_j+1<2^{\ell_j+3},
\]

and hence

\[
\beta_j\in\{1,2\}.
\]

Arithmetic now gives the exact odometer relation

\[
\begin{aligned}
\ell_{j+1}
&=\left\lfloor
  \log_2\frac{3n_j+1}{2^{a_j}}
  \right\rfloor\\
&=\lfloor\log_2(3n_j+1)\rfloor-a_j\\
&=\ell_j+\beta_j-a_j.
\end{aligned}
\]

Equivalently, put

\[
g_j=\beta_j-1\in\{0,1\},
\qquad
d_j=a_j-1\in\mathbb N.
\]

Then

\[
\boxed{\ell_{j+1}=\ell_j+g_j-d_j.}
\]

This has a literal chip interpretation.

\begin{itemize}
\item Every odd avalanche necessarily replaces one right-moving binary site by one left-moving ternary carrier.
\item When that carrier reaches the left boundary, it emits at least one binary site. That baseline one-for-one replacement does not change the population.
\item If \(\beta_j=2\), it emits one additional \textbf{growth chip}.
\item After the odd firing, the right boundary absorbs \(a_j-1\) additional zero sites. These are the \textbf{valuation-deficit chips}.
\end{itemize}

Thus the path length is a deterministic queue whose arrivals are \(g_j\) and whose services are \(d_j\).

After \(k\) odd avalanches,

\[
\boxed{
\ell_k
=
\ell_0+
\sum_{j=0}^{k-1}(\beta_j-1)
-
\sum_{j=0}^{k-1}(a_j-1).
}
\]

Since an odd positive integer has \(\ell=0\) exactly when it is \(1\), stabilization is equivalent to the following statement.


\paragraph*{Sandpile stabilization lemma}


For every odd \(n_0>1\), there exists \(k\ge1\) such that

\[
\boxed{
\sum_{j=0}^{k-1}(a_j-1)
\ge
\ell_0+
\sum_{j=0}^{k-1}(\beta_j-1).
}
\tag{S}
\]

Because \(\ell_k\ge0\), equality necessarily holds at the first such \(k\). An equivalent form is

\[
\boxed{
\sum_{j=0}^{k-1}(a_j-\beta_j)=\ell_0.
}
\tag{S'}
\]

Condition (S) is not a consequence of ordinary sandpile stabilization. It is precisely the Collatz conjecture written as extinction of the binary-site queue.


\bigskip



\subsubsection*{4. The leading carry is almost a rotation}


The left-boundary production is highly structured rather than arbitrary.

Write

\[
\theta_j=\log_2 n_j-\ell_j\in[0,1),
\qquad
\lambda=\log_2 3,
\]

and

\[
\varepsilon_j
=
\log_2\left(1+\frac1{3n_j}\right)>0.
\]

Then

\[
\log_2 n_{j+1}
=
\log_2 n_j+\lambda-a_j+\varepsilon_j.
\]

Subtracting the odometer relation gives

\[
\boxed{
\theta_{j+1}
=
\theta_j+\lambda-\beta_j+\varepsilon_j.
}
\]

Therefore

\[
\boxed{
\sum_{j=0}^{k-1}\beta_j
=
k\lambda+\theta_0-\theta_k
+\sum_{j=0}^{k-1}\varepsilon_j.
}
\]

In chip language, the cumulative number of sites emitted at the left boundary is an irrational-rotation count, up to a bounded phase discrepancy and the small affine corrections \(\varepsilon_j\).

Consequently,

\[
\sum_{j<k}g_j
=
k(\log_2 3-1)
+\theta_0-\theta_k
+\sum_{j<k}\varepsilon_j.
\]

The unresolved problem is therefore entirely on the right boundary: proving that the cumulative excess valuation

\[
\sum_{j<k}(a_j-1)
\]

eventually beats this deterministic arrival count by the initial queue length \(\ell_0\).


\bigskip



\subsubsection*{5. What least action does\textemdash{}and does not\textemdash{}give}


For a genuine abelian network, the least-action principle says that a legal execution performs no more work than any finite complete execution, even when the comparison execution is illegal. Hence the existence of one finite complete execution would prove that every legal firing schedule terminates. [\hypertarget{cit:R38:3}{}\hyperlink{ref:R38:3}{3}, \hypertarget{cit:R38:4}{}\hyperlink{ref:R38:4}{4}]

In the present carrier system, however, constructing such a complete execution means specifying enough right-boundary absorptions to remove

\[
\ell_0+\sum g_j
\]

sites. That is exactly condition (S). Least action gives order independence \textbf{after} a complete firing vector has been found; it does not produce that firing vector.

This distinction is the central gap in the proposed proof.

A finite irreducible abelian network has a powerful all-input criterion: it halts on all inputs precisely when the eigenvalues of its production matrix lie inside the unit disk. But the Collatz carrier has an unbounded ordered tape and can create new sites, so it is not a finite network to which that theorem applies. Countably infinite abelian networks can even simulate Turing machines, so abelianity alone does not imply decidable or automatic stabilization. [\hypertarget{cit:R38:5}{}\hyperlink{ref:R38:5}{5}]

There is also no way to replace the tape by a fixed finite unary-input processor while retaining an ordinary output encoding of \(T(n)\): finite abelian processors compute weakly increasing functions of their input counts, whereas the Collatz shortcut map is not increasing. An ordered digit tape or an equivalent unbounded state is unavoidable. [\hypertarget{cit:R38:6}{}\hyperlink{ref:R38:6}{6}]


\bigskip



\subsubsection*{6. No nonnegative additive sandpile energy can work}


One might try to assign a positive weight to every chip or site so that every firing strictly decreases total energy. This entire class of proofs is impossible.

Allow the weight of a symbol to depend arbitrarily on its distance \(r\) from the right boundary. Write

\[
w_{\mathtt f}(r),\quad
w_{\mathtt t}(r),\quad
w_{\mathtt0}(r),\quad
w_{\mathtt1}(r),\quad
w_{\mathtt2}(r)\ge0.
\]

Suppose

\[
V(W)=\sum_{\text{sites }x}w_x(r_x)
\]

strictly decreases under every firing.

From the right-boundary rule

\[
\mathtt{t.}\longrightarrow\mathtt{2.}
\]

we obtain

\[
w_{\mathtt t}(0)>w_{\mathtt2}(0).
\]

Set

\[
\delta(r)=w_{\mathtt t}(r)-w_{\mathtt2}(r).
\]

The internal rule

\[
\mathtt{t2}\longrightarrow\mathtt{2t}
\]

at positions \(r+1,r\) requires

\[
w_{\mathtt t}(r+1)+w_{\mathtt2}(r)
>
w_{\mathtt2}(r+1)+w_{\mathtt t}(r),
\]

so

\[
\delta(r+1)>\delta(r).
\]

Hence \(\delta(r)>0\) for every \(r\).

But when a ternary \(2\) reaches the left boundary at distance \(m\), the rule

\[
\mathtt{/2}\longrightarrow\mathtt{/ft}
\]

requires

\[
w_{\mathtt2}(m)
>
w_{\mathtt f}(m+1)+w_{\mathtt t}(m)
\ge w_{\mathtt t}(m),
\]

which says

\[
\delta(m)<0.
\]

This is a contradiction.

Therefore:

\[
\boxed{\text{No nonnegative site-additive potential can decrease under every firing.}}
\]

The obstruction persists even with completely arbitrary position-dependent weights. Any successful stabilization proof must be nonlinear, block-based, or genuinely multiscale.


\bigskip



\subsubsection*{7. The critical all-ones glider}


There is an exact family showing why no fixed local carry window can provide uniform contraction.

Let

\[
n_0=2^m-1.
\]

For \(0\le j\le m-1\), as long as the displayed iteration is defined,

\[
\boxed{
n_j=3^j2^{m-j}-1.
}
\]

Indeed, for \(j\le m-2\),

\[
3n_j+1
=
2\left(3^{j+1}2^{m-j-1}-1\right),
\]

and the factor in parentheses is odd. Hence

\[
a_j=1
\]

and the formula follows by induction.

Two important consequences follow.


\paragraph*{No bounded-horizon descent}


For every fixed \(K\), choose \(m>K+1\). Then for \(1\le j\le K\),

\[
n_j=3^j2^{m-j}-1>2^m-1=n_0.
\]

Thus there is no universal bounded number of odd avalanches within which ordinary size must decrease.

In the carrier picture, the word consisting of a long block of binary \(1\)'s shadows the infinite all-ones configuration, which represents the \(2\)-adic fixed point \(-1\). The finite left boundary can remain invisible for arbitrarily long.


\paragraph*{Fixed binary and ternary windows both become blind}


For fixed \(r,s\), whenever

\[
s\le j\le m-r,
\]

we have

\[
n_j+1=3^j2^{m-j}
\]

divisible by \(3^s2^r\). Hence

\[
\boxed{
n_j\equiv-1\pmod{2^r3^s}.
}
\]

A long central part of the trajectory therefore appears completely stationary to every fixed collection of binary and ternary residue windows.

This is a particularly sharp obstruction to a finite-state sandpile certificate: any fixed amount of low-end binary carry data and terminal ternary carry data admits an arbitrarily long critical glider. An effective proof has to record an \textbf{unbounded distance from the \(2\)-adic state \(-1\)}, or use some other genuinely global quantity.


\bigskip



\subsubsection*{8. What the sandpile critical group sees about cycles}


The sandpile formulation gives a clean matrix for a putative odd cycle.

Suppose a cycle has odd values

\[
n_0,n_1,\ldots,n_{k-1}
\]

and valuations \(a_0,\ldots,a_{k-1}\), with indices taken modulo \(k\). Then

\[
2^{a_i}n_{i+1}-3n_i=1.
\]

Define the cyclic toppling matrix \(L_{\boldsymbol a}\) by

\[
(L_{\boldsymbol a})_{ii}=2^{a_{i-1}},
\qquad
(L_{\boldsymbol a})_{i,i-1}=-3
\]

with cyclic indices and all other entries zero. The cycle equations are

\[
L_{\boldsymbol a}
\begin{pmatrix}
n_0\\ \vdots\\ n_{k-1}
\end{pmatrix}
=
\begin{pmatrix}
1\\ \vdots\\1
\end{pmatrix}.
\]

Its determinant is

\[
\boxed{
\det L_{\boldsymbol a}=2^A-3^k,
\qquad
A=\sum_{i=0}^{k-1}a_i.
}
\]

Thus the familiar Collatz cycle denominator is exactly the determinant of a cyclic toppling matrix.

The matrix is a nonsingular \(M\)-matrix precisely when

\[
2^A>3^k.
\]

This condition is already necessary for a positive cycle. Sandpile theory therefore recovers the correct archimedean inequality, but a cycle exists only when

\[
\mathbf1\in L_{\boldsymbol a}\mathbb Z^k
\]

and the resulting integer solution is positive and has exactly the prescribed valuations. That lattice-integrality condition is the unresolved arithmetic part.

For the trivial cycle \(n=1\), \(k=1,a_0=2\), and

\[
L=[4-3]=[1].
\]

For a nontrivial valuation word, the critical group records the denominator \(2^A-3^k\), but it does not by itself prove that the source vector \(\mathbf1\) represents a nonzero class.

Likewise, on a divergent trajectory every finite prefix produces a perfectly stabilizable finite toppling system. The possible failure occurs only in the limit, through escape to newly created sites at infinity.


\bigskip



\subsubsection*{9. The exact missing theorem}


The proposed sandpile proof would be complete if one could prove the following, without assuming any form of Collatz convergence.


\begin{quote}
\textbf{Pointwise odometer theorem.}   For every odd \(n_0>1\), let \[
> n_{j+1}=\frac{3n_j+1}{2^{a_j}},
> \qquad
> a_j=v_2(3n_j+1),
> \] and \[
> \beta_j
> =
> \lfloor\log_2(3n_j+1)\rfloor
> -
> \lfloor\log_2 n_j\rfloor.
> \] Then for some \(k\), \[
> \sum_{j<k}(a_j-\beta_j)
> =
> \lfloor\log_2 n_0\rfloor.
> \]

\end{quote}


The proof would then be immediate:

\[
\ell_k
=
\ell_0-\sum_{j<k}(a_j-\beta_j)=0,
\]

so \(n_k=1\).

Conversely, if \(n_k=1\), the same equation shows that the claimed odometer identity holds. Hence the pointwise odometer theorem and the Collatz conjecture are equivalent.


\bigskip



\subsection*{Verdict}


The sandpile approach succeeds in producing:

\begin{enumerate}
\item an exact mixed binary\textendash{}ternary carrier network;
\item a strong local abelian/diamond property;
\item termination of every internal carry avalanche;
\item an exact growth-versus-deficit chip balance;
\item a cyclic toppling matrix with determinant \(2^A-3^k\);
\item rigorous obstructions to additive potentials, finite production matrices, bounded stopping horizons, and fixed residue-window certificates.
\end{enumerate}

It does \textbf{not} presently prove that every admissible finite configuration reaches the absorbing state \texttt{/.}. That assertion is exactly condition (S), hence exactly Collatz.

The central unresolved object is a critical one-dimensional queue:

\[
\ell_{j+1}
=
\ell_j
+\underbrace{(\beta_j-1)}_{\text{one possible birth}}
-\underbrace{(a_j-1)}_{\text{valuation absorption}}.
\]

A genuine completion would need a nonlinear multiscale burning argument that matches every initial and newly created binary site to a later valuation absorption while remaining sensitive to arbitrarily long approximations of the \(2\)-adic all-ones state. No such matching or burning certificate has been established here, so presenting this reduction as a full proof would be incorrect.
\subsection*{References}
\begin{enumerate}[leftmargin=2.2em]
\item \hypertarget{ref:R38:1}{}\hyperlink{cit:R38:1}{\texttt{turn287599view2}}
\item \hypertarget{ref:R38:2}{}\hyperlink{cit:R38:2}{\texttt{turn287599view3}}
\item \hypertarget{ref:R38:3}{}\hyperlink{cit:R38:3}{\texttt{turn960348view0}}
\item \hypertarget{ref:R38:4}{}\hyperlink{cit:R38:4}{\texttt{turn960348view2}}
\item \hypertarget{ref:R38:5}{}\hyperlink{cit:R38:5}{\texttt{turn287599view1}}
\item \hypertarget{ref:R38:6}{}\hyperlink{cit:R38:6}{\texttt{turn287599view4}}
\end{enumerate}

\clearpage
\section{Round 39: approach 34 --- Directed-polymer free energy}
\subsection*{Prompt}
\begin{quote}
Complete a full proof of the Collatz conjecture using the following approach:

\begin{enumerate}
\item Directed-polymer free energy
\end{enumerate}

\textbf{Speculative \textperiodcentered{} P}

Treat critical valuation words as polymers in a deterministic arithmetic environment. Establish a strictly negative quenched free energy after incorporating exact residue constraints \textemdash{} not merely the negative annealed drift seen by random words.

Take as long as you need.
\end{quote}
\subsection*{Response}

\subsection*{Result}


The directed-polymer program does \textbf{not} currently yield a full proof. More strongly, the proposed decisive statement\textemdash{}strictly negative quenched free energy after imposing exact residue constraints\textemdash{}is \textbf{false for the natural path-conditioned partition function}.

What does emerge is a rigorous and fairly sharp reformulation:

\begin{enumerate}
\item The ordinary annealed critical-word free energy is strictly negative, with an explicit exponent.
\item Exact \(3\)-adic residue conditioning makes a uniformly chosen residue almost always empty, but the residue generated by a candidate orbit is necessarily occupied. On that path-selected environment, the normalized quenched free energy is nonnegative.
\item All non-pinned cylinder mass can be proved to decay exponentially.
\item Collatz is exactly equivalent to extinction of the remaining \textbf{pinned ground-state partition}.
\item A pointwise occupation estimate, suggested by the free-energy calculation, would complete the proof\textemdash{}but that estimate is the unresolved all-seed synchronization theorem, not a consequence of annealed pressure.
\end{enumerate}

Thus the polymer idea survives only after replacing ordinary finite-temperature free energy by a pinned, zero-temperature problem.


\bigskip



\subsubsection*{1. Valuation words as polymers}


Use the accelerated map on positive odd integers

\[
U(n)=\frac{3n+1}{2^{\nu_2(3n+1)}}.
\]

For a word

\[
w=(a_0,\ldots ,a_{j-1}),\qquad a_i\ge 1,
\]

put

\[
A_t=\sum_{i<t}a_i,\qquad C_0=0,\qquad
C_{t+1}=3C_t+2^{A_t}.
\]

Then every realization of \(w\) satisfies

\[
2^{A_t}U^t(n)=3^t n+C_t. \tag{1}
\]

The word determines one odd starting residue class modulo \(2^{A_j+1}\):

\[
n\equiv
3^{-j}\bigl(2^{A_j}-C_j\bigr)
\pmod {2^{A_j+1}}. \tag{2}
\]

Let \(R_w\) be its least positive representative. Thus every finite word is realizable, and its positive realizations are exactly

\[
R_w+\ell 2^{A_j+1},\qquad \ell\ge 0. \tag{3}
\]

The endpoint residue is

\[
\rho_j(w):=2^{-A_j}C_j\pmod {3^j}. \tag{4}
\]

In particular, exact endpoint constraints are genuinely \(3\)-adic. Your terminal-locality theorem further shows that divisibility of the endpoint \(X+1\) by \(3^r\) depends only on the last \(r\) valuations. [\hypertarget{cit:R39:1}{}\hyperlink{ref:R39:1}{1}]

Set

\[
\alpha=\log_2 3.
\]

The critical or strict-barrier language is

\[
\mathcal B_j=
\left\{
w:\ A_t<\alpha t\quad(1\le t\le j)
\right\}. \tag{5}
\]

For such a word the multiplicative coefficient \(3^t/2^{A_t}\) exceeds \(1\) at every prefix.


\bigskip



\subsubsection*{2. The annealed free energy is explicitly negative}


Under Haar measure on odd \(2\)-adic integers, the valuations are independent with

\[
\mathbb P(a_i=m)=2^{-m},\qquad m\ge1.
\]

This is the exact random-word model underlying the Syracuse probabilistic picture; Tao's almost-all theorem goes considerably further by obtaining approximate transport for associated \(3\)-adic random variables, but still in an averaged rather than pointwise setting. [\hypertarget{cit:R39:2}{}\hyperlink{ref:R39:2}{2}, \hypertarget{cit:R39:3}{}\hyperlink{ref:R39:3}{3}]

Define the annealed critical partition function

\[
Z_j^{\mathrm{ann}}
=
\sum_{w\in\mathcal B_j}2^{-A_j}
=
\mathbb P\left(A_t<\alpha t\text{ for all }t\le j\right). \tag{6}
\]

For one increment,

\[
M(\theta)
=
\mathbb E e^{\theta a}
=
\frac{e^\theta}{2-e^\theta}.
\]

The Cram\'er optimizer for mean \(\alpha<2\) is

\[
\theta_*=
\log\frac{2(\alpha-1)}{\alpha},
\]

and the rate is

\[
\begin{aligned}
I(\alpha)
&=\alpha\theta_*-\log M(\theta_*)\\
&=\log 3-
\left[
\alpha\log\alpha-(\alpha-1)\log(\alpha-1)
\right] \\
&=0.0549794728108\ldots . 
\end{aligned}
\tag{7}
\]

A cyclic-shift argument shows that requiring the inequality at every prefix changes only a subexponential factor: for any composition with \(A_j<\alpha j\), rotating just after the unique maximum of \(A_t-\alpha t\) gives a barrier word. Consequently,

\[
\boxed{
\lim_{j\to\infty}
\frac1j\log Z_j^{\mathrm{ann}}
=-I(\alpha)<0.
}
\tag{8}
\]

The unweighted critical-word entropy is

\[
|\mathcal B_j|
=
\exp\left(h(\alpha)j+o(j)\right), \tag{9}
\]

where

\[
h(\alpha)
=
\alpha\log\alpha-(\alpha-1)\log(\alpha-1)
=
1.043632815857\ldots .
\]

Thus

\[
e^{h(\alpha)}=2.83951373049\ldots <3
\]

and

\[
\frac{|\mathcal B_j|}{3^j}
=
e^{-I(\alpha)j+o(j)}. \tag{10}
\]

In base \(2\), the rate is

\[
I_2=\frac{I(\alpha)}{\log2}
=0.07931861277\ldots . \tag{11}
\]

So the annealed calculation is completely successful.


\bigskip



\subsubsection*{3. Why exact-residue quenching has the wrong sign}


For \(\xi\pmod {3^j}\), define the residue-fiber count

\[
N_j(\xi)
=
\#\{w\in\mathcal B_j:\rho_j(w)=\xi\}. \tag{12}
\]

Averaging over all residues gives

\[
\frac1{3^j}\sum_{\xi\bmod3^j}N_j(\xi)
=
\frac{|\mathcal B_j|}{3^j}
=
e^{-I(\alpha)j+o(j)}. \tag{13}
\]

Therefore a uniformly selected residue is almost always empty. This is the tempting ``negative quenched free energy.''

But a hypothetical critical trajectory does not select a uniform residue. It selects

\[
\xi_j=\rho_j(w_j)
\]

for its own prefix \(w_j\), and hence

\[
N_j(\xi_j)\ge1. \tag{14}
\]

Therefore

\[
\frac1j\log N_j(\xi_j)\ge0. \tag{15}
\]

The negative value in (13) says that \textbf{most environments contain no critical polymer}. It says nothing about an environment conditioned to contain a distinguished polymer.

The same obstruction persists for the weighted partition. Define

\[
W_j(\xi)
=
\sum_{\substack{w\in\mathcal B_j\\ \rho_j(w)=\xi}}
2^{-A_j},
\]

and remove the unavoidable probability \(3^{-j}\) of naming one particular residue:

\[
Q_j(\xi)=3^jW_j(\xi). \tag{16}
\]

Its uniform-residue average is exactly

\[
\frac1{3^j}\sum_{\xi}Q_j(\xi)
=
Z_j^{\mathrm{ann}}
=
e^{-I(\alpha)j+o(j)}. \tag{17}
\]

Nevertheless, for every \(w\in\mathcal B_j\),

\[
\begin{aligned}
Q_j(\rho_j(w))
&\ge3^j2^{-A_j}\\
&=\exp\left(j\log3-A_j\log2\right)\\
&>1. 
\end{aligned}
\tag{18}
\]

Hence a uniform estimate of the form

\[
Q_j(\xi)\le e^{-\varepsilon j}
\]

on every occupied critical environment is false at every depth.

Using the raw partition \(W_j\) instead does give a negative exponent, because a single term \(2^{-A_j}\) is exponentially small. But then one surviving infinite path is entirely compatible with negative free energy. Thus there is an unavoidable dichotomy:

\begin{itemize}
\item Give each surviving path subexponential weight, and a path-selected environment has nonnegative pressure.
\item Give each surviving path exponentially decaying weight, and negative pressure no longer implies extinction.
\end{itemize}

This is not a technical gap. It is a general no-go for ordinary positive partition functions in an endogenous, path-generated environment. Standard directed polymers instead interact with an external quenched disorder. [\hypertarget{cit:R39:4}{}\hyperlink{ref:R39:4}{4}, \hypertarget{cit:R39:5}{}\hyperlink{ref:R39:5}{5}]

Adding all prefix or suffix residues does not fix the issue: the distinguished path continues to satisfy its own exact constraints. Quenching on the starting integer is even more degenerate\textemdash{}the fiber then contains exactly its unique valuation path.


\bigskip



\subsubsection*{4. An exact pinned-polymer reduction}


There is, however, a corrected partition for which extinction really is equivalent to Collatz.

For an odd \(n\ge3\), define its first accelerated descent time

\[
\tau(n)=
\inf\{t\ge1:U^t(n)<n\},
\]

with \(\tau(n)=\infty\) if no such \(t\) exists. For \(s>1\), let

\[
\mathcal Z_j(s)
=
\sum_{\substack{n\ge3\ \mathrm{odd}\\ \tau(n)>j}}
n^{-s}. \tag{19}
\]


\paragraph*{Proposition 1}


For every \(s>1\),

\[
\boxed{
\text{Collatz}
\quad\Longleftrightarrow\quad
\mathcal Z_j(s)\longrightarrow0.
}
\tag{20}
\]


\paragraph*{Proof}


If Collatz holds, every odd \(n\ge3\) eventually reaches \(1<n\), so \(\tau(n)<\infty\). Since \(\sum n^{-s}<\infty\), dominated convergence gives \(\mathcal Z_j(s)\to0\).

Conversely, if \(\mathcal Z_j(s)\to0\), no \(n\ge3\) can have \(\tau(n)=\infty\). Thus every odd \(n>1\) has a later odd iterate smaller than itself. Strong induction on odd \(n\) then proves convergence to \(1\). \ensuremath{\blacksquare}

Now decompose \(\mathcal Z_j\) by valuation words.


\paragraph*{Lemma 2: least-realization pinning}


Suppose some realization

\[
n=R_w+\ell2^{A_j+1}
\]

of a length-\(j\) word satisfies

\[
U^t(n)\ge n\qquad(1\le t\le j).
\]

Then its least realization \(R_w\) also satisfies

\[
U^t(R_w)\ge R_w\qquad(1\le t\le j). \tag{21}
\]


\paragraph*{Proof}


For each prefix,

\[
U^t(x)-x
=
\frac{(3^t-2^{A_t})x+C_t}{2^{A_t}}. \tag{22}
\]

If \(3^t-2^{A_t}>0\), the right side is positive at every positive \(x\). If \(3^t-2^{A_t}<0\), it is a decreasing affine function of \(x\), so replacing \(n\) by the smaller \(R_w\) can only increase it. \ensuremath{\blacksquare}

Define the pinned words

\[
\mathscr P_j=
\left\{
w:\ R_w\ge3,\quad
U^t(R_w)\ge R_w\ (1\le t\le j)
\right\}
\]

and their ground-state partition

\[
\Pi_j(s)=\sum_{w\in\mathscr P_j}R_w^{-s}. \tag{23}
\]


\paragraph*{Theorem 3: exponential extinction of all unpinned cylinder mass}


For every \(s>1\),

\[
\boxed{
0\le
\mathcal Z_j(s)-\Pi_j(s)
\le
\frac{2^{-s}\zeta(s)}{(2^s-1)^j}.
}
\tag{24}
\]


\paragraph*{Proof}


By Lemma 2, a cylinder contributes only when its least representative is pinned. Its nonleast realizations satisfy

\[
R_w+\ell2^{A_j+1}\ge \ell2^{A_j+1},
\qquad \ell\ge1.
\]

Hence their total contribution is at most

\[
\begin{aligned}
\sum_{w}\sum_{\ell\ge1}
\left(\ell2^{A_j(w)+1}\right)^{-s}
&=
2^{-s}\zeta(s)
\sum_w2^{-sA_j(w)}\\
&=
2^{-s}\zeta(s)
\left(\sum_{a\ge1}2^{-sa}\right)^j\\
&=
\frac{2^{-s}\zeta(s)}{(2^s-1)^j}.
\end{aligned}
\]

Since \(s>1\), the final factor decays exponentially. \ensuremath{\blacksquare}

Therefore

\[
\boxed{
\text{Collatz}
\quad\Longleftrightarrow\quad
\Pi_j(s)\longrightarrow0.
}
\tag{25}
\]

This is the exact pinned-polymer version of the conjecture.


\bigskip



\subsubsection*{5. The pinned atoms are literally the counterexamples}


For an odd \(n\), let \(A_j(n)\) be the accumulated valuations along its actual orbit, and define its birth time

\[
b(n)=
\min\left\{
j\ge1:2^{A_j(n)+1}>n
\right\}. \tag{26}
\]

Once the cylinder modulus exceeds \(n\), the least positive representative of the orbit's word is exactly \(n\). Consequently,

\[
\boxed{
\Pi_j(s)
=
\sum_{\substack{n\ge3\\n\ \mathrm{odd}}}
n^{-s}
\mathbf 1_{\{\,b(n)\le j<\tau(n)\,\}}.
}
\tag{27}
\]

Dominated convergence gives the exact zero-temperature limit

\[
\boxed{
\lim_{j\to\infty}\Pi_j(s)
=
\sum_{\substack{n\ge3\ \mathrm{odd}\\\tau(n)=\infty}}
n^{-s}.
}
\tag{28}
\]

Thus:

\begin{itemize}
\item Finite-lifetime polymers eventually disappear.
\item New nonstationary representatives have exponentially small Dirichlet mass.
\item A number with no smaller iterate becomes a permanent pinned atom of weight \(n^{-s}\).
\end{itemize}

If Collatz is false and \(m\) is its least odd counterexample, then

\[
U^t(m)\ge m\qquad(t\ge0)
\]

by minimality. Therefore, once \(j\ge b(m)\),

\[
\Pi_j(s)\ge m^{-s}. \tag{29}
\]

Its free energy is consequently at least zero.

A theorem such as

\[
\Pi_j(s)\le Ce^{-cj} \tag{30}
\]

would indeed prove Collatz. But (28) shows exactly what such a theorem must do: it must exclude every permanent atom. Negative bulk pressure alone cannot do this.

Moreover, strict negative pressure is stronger than Collatz. Even if every \(\tau(n)\) is finite, the maximum descent times could grow so rapidly that \(\Pi_j(s)\to0\) only subexponentially.


\bigskip



\subsubsection*{6. The useful quenched target is an occupation theorem}


There is a second precise way to see what the free-energy calculation should target.

Let

\[
D_j(n)=A_j(n)-\alpha j,
\qquad x_j=U^j(n).
\]

Multiplying the step identities gives

\[
\frac{x_j}{n}
=
2^{-D_j(n)}
\prod_{i<j}
\left(1+\frac1{3x_i}\right).
\]

Therefore

\[
D_j(n)
=
\log_2\frac{n}{x_j}
+
\sum_{i<j}
\log_2\left(1+\frac1{3x_i}\right). \tag{31}
\]

Suppose \(n\) is a least counterexample. Then \(x_i\ge n\) for every \(i\), and hence

\[
D_j(n)
\le
j\log_2\left(1+\frac1{3n}\right)
\le
\frac{j}{3n\log 2}. \tag{32}
\]

For fixed \(c>0\), define the critical-strip occupation count

\[
O_c(n)
=
\#\{j\ge0:D_j(n)\le c\}. \tag{33}
\]

Equation (32) shows that a least counterexample satisfies

\[
O_c(n)\ge
\left\lfloor3cn\log2\right\rfloor+1. \tag{34}
\]

Thus the following would be sufficient.


\paragraph*{Pointwise quenched occupation criterion}


For some \(c>0\), prove uniformly that

\[
O_c(n)=o(n)
\qquad(n\to\infty). \tag{35}
\]

Then no sufficiently large least counterexample can exist; the remaining finite range can be checked, and strong induction completes Collatz.

The annealed polymer model strongly predicts much more. Since

\[
\mathbb P(D_j\le c)
=
\exp(-I(\alpha)j+o(j)),
\]

a union-bound heuristic over \(n\le N\) predicts that the final critical visit occurs near

\[
j\approx
\frac{\log_2N}{I_2}
=
12.6074\ldots\,\log_2N. \tag{36}
\]

But converting this prediction into a bound for \textbf{every} seed is exactly the annealed-to-quenched leap. Tao's result controls almost all starting values through approximate \(3\)-adic transport; it does not give such a uniform occupation bound. [\hyperlink{ref:R39:2}{2}]


\bigskip



\subsubsection*{7. Why terminal residue polymers do not close the gap}


The exact competitor decomposition in your previous work makes the remaining obstruction explicit.

Write \(w=uv\), where \(u\) has length \(d\) and \(v\) length \(k\). Put

\[
P=A(u),\qquad
\lambda=\frac{3^d}{2^P},\qquad
\beta=\frac{C_u}{2^P}.
\]

Let \(b\) be a competing inverse word of length \(k\), with exponent sum \(B\), and put

\[
c=A(v)-B>0.
\]

Exact endpoint compatibility is

\[
3^k\mid C_v-2^cC_b, \tag{37}
\]

and

\[
h(v,b,c)=
\frac{C_v-2^cC_b}{3^k}\in\mathbb Z.
\]

If \(Y\) is the competing ancestor and \(R\) the original start, then

\[
2^cY=\lambda R+\beta+h. \tag{38}
\]

This is the exact local-global decomposition proved in the existing manuscript. [\hypertarget{cit:R39:6}{}\hyperlink{ref:R39:6}{6}]

Define the zero-temperature score

\[
\mathcal E(u,v;b)
=
(\lambda-2^c)R+\beta+h. \tag{39}
\]

Then

\[
Y<R
\quad\Longleftrightarrow\quad
\mathcal E(u,v;b)<0. \tag{40}
\]

The terminal residue fiber determines \(h\). But the prefix determines

\[
(\lambda-2^c)R+\beta.
\]

The same terminal suffix and the same local slope data can have opposite archimedean signs when the prefix is changed; the exact example in the manuscript changes the relevant score by billions. [\hypertarget{cit:R39:7}{}\hyperlink{ref:R39:7}{7}]

Therefore a polymer environment that records only terminal exact residues has omitted the variable deciding descent.

At inverse temperature \(\gamma\), one could define

\[
Z_\gamma(u,v)
=
\sum_{b\ \mathrm{compatible}}
e^{-\gamma\mathcal E(u,v;b)}.
\]

As \(\gamma\to\infty\),

\[
-\frac1\gamma\log Z_\gamma(u,v)
\longrightarrow
\min_b\mathcal E(u,v;b). \tag{41}
\]

The required theorem is therefore a \textbf{zero-temperature min-plus statement}:

\[
\text{Every positive stationary retained ray eventually has }
\min_b\mathcal E(u,v;b)<0. \tag{42}
\]

That is precisely the previously identified synchronization theorem: terminal \(3\)-adic compatibility must align with the prefix-dependent archimedean surplus along one fixed positive integer. [\hypertarget{cit:R39:8}{}\hyperlink{ref:R39:8}{8}]

It remains unproved.


\bigskip



\subsubsection*{8. The strongest viable modification}


The directed-polymer route can be repaired into a concrete three-part program.


\paragraph*{Nonresonant sector}


Away from good rational approximations

\[
\frac{A}{j}\approx\log_2 3,
\]

seek a finite-state or min-plus transfer certificate forcing

\[
\min_b\mathcal E(u,v;b)\le-\delta
\]

for a uniform \(\delta>0\). This is the part most naturally accessible to a joint-spectral-radius or Bellman computation.


\paragraph*{Resonant sector}


Failure of a uniform margin should force long Christoffel or Sturmian blocks associated with continued-fraction convergents of \(\log_2 3\). An infinite valuation ray determines the \(2\)-adic value

\[
n=
-\sum_{i=0}^{\infty}
\frac{2^{A_i}}{3^{i+1}}
\qquad\text{in }\mathbb Z_2. \tag{43}
\]

One would need to prove that the resonant sequences produced by the first part make this value nonrational\textemdash{}and therefore not a positive integer. The Collatz conjugacy over mechanical/Sturmian words has been studied, but the central rationality question was not settled by that work. [\hypertarget{cit:R39:9}{}\hyperlink{ref:R39:9}{9}, \hypertarget{cit:R39:10}{}\hyperlink{ref:R39:10}{10}]


\paragraph*{Rigidity interface}


Finally, one must prove a dichotomy:

\[
\begin{aligned}
&\text{either infinitely many nonresonant blocks occur, giving descent;}\\
&\text{or the ray is eventually resonant enough for the \(2\)-adic}\\
&\text{irrationality theorem to apply.}
\end{aligned} \tag{44}
\]

This interface must retain the prefix surplus in (39). Merely classifying terminal suffixes does not suffice.


\bigskip



\subsection*{Conclusion}


The proposed proof cannot be completed by asserting negative exact-residue quenched free energy. For the natural normalized partition, that assertion is false on every critical environment selected by an actual path:

\[
Q_j(\rho_j(w))>1.
\]

The rigorous achievement of the approach is instead the exact decomposition

\[
\text{annealed/bulk mass}
\quad+\quad
\text{pinned least-realization atoms}.
\]

The bulk is exponentially subcritical. Collatz is equivalent to extinction of the pinned partition

\[
\Pi_j(s)
=
\sum_{w\in\mathscr P_j}R_w^{-s},
\]

or, equivalently, to a pointwise critical-strip occupation theorem. Proving that pinned extinction\textemdash{}or the zero-temperature synchronization inequality (42)\textemdash{}would be a complete proof. The current residue and free-energy estimates do not establish it.

As of August 2026, Collatz remains open; computation has verified convergence through \(2^{71}\), but no all-depth pinned-ground-state exclusion is known. [\hypertarget{cit:R39:11}{}\hyperlink{ref:R39:11}{11}, \hypertarget{cit:R39:12}{}\hyperlink{ref:R39:12}{12}]
\subsection*{References}
\begin{enumerate}[leftmargin=2.2em]
\item \hypertarget{ref:R39:1}{}\hyperlink{cit:R39:1}{\texttt{turn2file1 L415-L437}}
\item \hypertarget{ref:R39:2}{}\hyperlink{cit:R39:2}{\texttt{turn727320search0}}
\item \hypertarget{ref:R39:3}{}\hyperlink{cit:R39:3}{\texttt{turn727320search10}}
\item \hypertarget{ref:R39:4}{}\hyperlink{cit:R39:4}{\texttt{turn727320search3}}
\item \hypertarget{ref:R39:5}{}\hyperlink{cit:R39:5}{\texttt{turn727320search6}}
\item \hypertarget{ref:R39:6}{}\hyperlink{cit:R39:6}{\texttt{turn2file1 L523-L552}}
\item \hypertarget{ref:R39:7}{}\hyperlink{cit:R39:7}{\texttt{turn2file1 L555-L575}}
\item \hypertarget{ref:R39:8}{}\hyperlink{cit:R39:8}{\texttt{turn3file2 L147-L161}}
\item \hypertarget{ref:R39:9}{}\hyperlink{cit:R39:9}{\texttt{turn429563view1}}
\item \hypertarget{ref:R39:10}{}\hyperlink{cit:R39:10}{\texttt{turn804423search6}}
\item \hypertarget{ref:R39:11}{}\hyperlink{cit:R39:11}{\texttt{turn727320search5}}
\item \hypertarget{ref:R39:12}{}\hyperlink{cit:R39:12}{\texttt{turn727320search8}}
\end{enumerate}

\clearpage
\section{Round 40: approach 35 --- Dependent percolation domination}
\subsection*{Prompt}
\begin{quote}
Complete a full proof of the Collatz conjecture using the following approach:

\begin{enumerate}
\item Dependent percolation domination
\end{enumerate}

\textbf{Hybrid \textperiodcentered{} P}

At each depth, call a branch "open" if it remains growth-compatible and integrally realizable. Prove that the resulting highly dependent tree is dominated by a subcritical percolation process, perhaps after grouping levels into blocks.

Take as long as you need.
\end{quote}
\subsection*{Response}

\subsection*{Outcome}


I do not obtain a full proof of the Collatz conjecture from this approach. More decisively, the proposed subcritical domination statement is \textbf{false as written} under the natural forward-tree interpretation.

There is, however, a genuine subcritical calculation hiding in the idea. It gives an exponentially decaying \textbf{average density} of surviving branches. The obstruction is converting that annealed decay into a statement about every individual positive integer. I will make that distinction precise and formulate the exact strengthened estimate that would complete the proof.


\subsubsection*{1. The forward candidate tree}


Use the shortcut Collatz map

\[
T(n)=
\begin{cases}
n/2,&n\equiv0\pmod 2,\\[1mm]
(3n+1)/2,&n\equiv1\pmod 2.
\end{cases}
\]

For a binary word

\[
w=(\varepsilon _0,\ldots,\varepsilon _{k-1})\in\{0,1\}^k,
\qquad
S_j(w)=\sum_{i<j}\varepsilon_i,
\]

the corresponding affine iterate is

\[
T^j(n)=\frac{3^{S_j(w)}n+C_j(w)}{2^j},
\qquad C_j(w)\geq0.
\]

Its multiplicative coefficient is

\[
q_j(w)=\frac{3^{S_j(w)}}{2^j}.
\]

Call \(w\) growth-compatible when

\[
q_j(w)\geq1
\quad\text{for every }1\leq j\leq k.
\tag{1}
\]

Writing

\[
\beta=\log_3 2=0.6309297535\ldots,
\]

condition (1) is equivalent to

\[
S_j(w)\geq \beta j
\quad(1\leq j\leq k).
\tag{2}
\]


\paragraph*{Finite integrality does not prune this tree}


Every binary word of length \(k\) is the parity vector of exactly one residue class modulo \(2^k\), and hence of infinitely many positive integers. Thus every finite word satisfying (2) is already positively and integrally realizable. [\hypertarget{cit:R40:1}{}\hyperlink{ref:R40:1}{1}]

So, with this definition, ``integrally realizable'' supplies no additional thinning at finite depth.


\subsubsection*{2. A rigorous obstruction to subcritical domination}



\paragraph*{The all-ones ray}


Consider

\[
w_k=1^k.
\]

Every prefix satisfies

\[
q_j(w_k)=\left(\frac32\right)^j>1.
\]

Moreover, \(w_k\) is realized by the positive integer

\[
r_k=2^k-1.
\]

Indeed,

\[
T^j(2^k-1)=3^j2^{k-j}-1
\qquad(0\leq j\leq k),
\]

so its first \(k\) iterates are odd.

Consequently, the deterministic open tree contains the infinite ray

\[
1,\quad 11,\quad 111,\quad\ldots.
\tag{3}
\]

The nested representatives \(2^k-1\) do not stabilize at a positive integer. They converge 2-adically to \(-1\); equivalently, the unique 2-adic integer with parity sequence \(111\ldots\) is \(-1\). [\hypertarget{cit:R40:2}{}\hyperlink{ref:R40:2}{2}]


\paragraph*{No pathwise subcritical coupling}


Suppose the open arithmetic tree were contained, under a coupling, in a percolation cluster which is subcritical and therefore has no infinite ray almost surely. The deterministic ray (3) would then have to belong to that cluster almost surely, a contradiction.

This remains true after grouping any fixed number \(L\) of levels.

There is an even stronger local obstruction. Define the barrier surplus

\[
D_j(w)=S_j(w)-\beta j.
\]

If an open prefix \(w\) has

\[
D_k(w)\geq \beta L,
\]

then \textbf{every one} of its \(2^L\) binary continuations of length \(L\) remains open: even the all-zero continuation has

\[
D_{k+t}\geq D_k-\beta t\geq0
\qquad(0\leq t\leq L).
\]

For the prefix \(1^m\),

\[
D_m=(1-\beta)m,
\]

so for every fixed \(L\), sufficiently large all-one prefixes have a complete open binary cone of depth \(L\).

Therefore:

\[
\boxed{\text{No fixed block size admits a uniform conditional open probability }p<1.}
\]

This also shows why a Liggett\textendash{}Schonmann\textendash{}Stacey-type theorem cannot be applied directly. Such domination results concern genuine random fields satisfying uniform conditional-probability or finite-dependence hypotheses. Here the candidate tree is deterministic, its prefix events have unbounded dependence, and some states have deterministic full open cones. [\hypertarget{cit:R40:3}{}\hyperlink{ref:R40:3}{3}]


\subsubsection*{3. What the probabilistic calculation really proves}


Let \(B_k\) be the number of binary words satisfying (2). Then

\[
B_k
\leq
\sum_{m\geq\beta k}\binom{k}{m}
\leq
(k+1)2^{H_2(\beta)k},
\]

where

\[
H_2(x)=-x\log_2x-(1-x)\log_2(1-x).
\]

Conversely, fix \(p>\beta\). Under Bernoulli-\(p\) bits, the walk

\[
S_j-\beta j
\]

has positive drift and has a positive probability of never becoming negative. Intersecting with its typical set \(S_k=pk+o(k)\) gives

\[
B_k\geq 2^{(H_2(p)-o(1))k}.
\]

Letting \(p\downarrow\beta\) yields

\[
\lim_{k\to\infty}B_k^{1/k}
=
2^{H_2(\beta)}
=
1.9318131064\ldots.
\tag{4}
\]

Thus the number of open nodes grows exponentially. The tree is not combinatorially subcritical.

On the other hand, under a Haar-random 2-adic seed, all \(2^k\) parity words are equally likely. Therefore

\[
\mathbb P(\text{open through depth }k)
=
\frac{B_k}{2^k}
=
\bigl(0.9659065532\ldots\bigr)^{k+o(k)}.
\tag{5}
\]

This is the attractive ``subcritical'' number in the proposal. But it is a decay of \textbf{measure}, not of the number of candidate branches:

\[
1.931813\ldots>1,
\qquad
\frac{1.931813\ldots}{2}=0.965906\ldots<1.
\]

Equation (5) proves that a Haar-random 2-adic integer has finite coefficient stopping time almost surely. It does not exclude a single exceptional positive integer. The parity map is a bijection and 2-adic isometry, so the probabilistic model genuinely describes Haar-random 2-adic seeds; positive ordinary integers form a countable Haar-null subset. [\hyperlink{ref:R40:2}{2}]


\subsubsection*{4. The exact ordinary-integer condition}


For a parity word \(w\) of length \(k\), let

\[
r(w)\in\{0,\ldots,2^k-1\}
\]

be its unique canonical realizing residue.

Consider a nested ray

\[
w_1\prec w_2\prec\cdots
\]

and write \(r_k=r(w_k)\).


\paragraph*{Stationarity lemma}


The ray is the parity sequence of a positive ordinary integer \(n\) if and only if

\[
r_k=n
\quad\text{for all sufficiently large }k.
\tag{6}
\]

Indeed, if the ray belongs to \(n\), then \(r_k\equiv n\pmod{2^k}\). Once \(2^k>n\), the canonical representative is exactly \(n\). The converse is immediate.

Thus Collatz requires excluding not all infinite open rays\textemdash{}there are uncountably many 2-adic ones\textemdash{}but only the \textbf{eventually stationary positive rays}.

That is a tail property. No finite prefix, fixed suffix, or fixed block knows whether the canonical representatives will eventually stop changing. In particular, the all-ones ray is open forever but nonstationary:

\[
r_k=2^k-1\longrightarrow-1
\quad\text{2-adically}.
\]

This is the precise point at which the percolation analogy loses pointwise dynamical meaning.


\subsubsection*{5. A genuine subcritical theorem in the inverse formulation}


There is another interpretation of ``integrally realizable'' under which an exact subcritical calculation does occur.

Use the accelerated odd map

\[
U(x)=\frac{3x+1}{2^{\nu_2(3x+1)}}.
\]

A valuation word

\[
a=(a_1,\ldots,a_j),\qquad a_i\geq1,
\]

has cumulative sums

\[
A_t=a_1+\cdots+a_t.
\]

The growth barrier is

\[
2^{A_t}\leq3^t,
\qquad\text{hence}\qquad
A_t\leq \alpha t,
\quad
\alpha=\log_2 3.
\tag{7}
\]

Ignoring the prefix restrictions can only increase the number of words. The number of positive compositions of total \(A\) into \(j\) entries is \(\binom{A-1}{j-1}\), so

\[
\#\{\text{growth-compatible words of length }j\}
\leq
\sum_{A=j}^{\lfloor\alpha j\rfloor}
\binom{A-1}{j-1}
=
\binom{\lfloor\alpha j\rfloor}{j}.
\tag{8}
\]

Stirling's formula gives the exponential rate

\[
\Lambda
=
\frac{\alpha^\alpha}
     {(\alpha-1)^{\alpha-1}}
=
2.8395137305\ldots.
\tag{9}
\]

For a fixed word \(a\), its inverse path to an endpoint \(X\) is integral precisely when

\[
2^{A_j}X\equiv C_a\pmod{3^j}.
\tag{10}
\]

Because \(2\) is invertible modulo \(3^j\), equation (10) selects exactly one endpoint residue modulo \(3^j\). Hence, for a Haar-random \(X\in\mathbb Z_3\),

\[
\mathbb E\bigl[
\#\{\text{integral growth-compatible inverse words of length }j\}
\bigr]
\leq
\frac{\binom{\lfloor\alpha j\rfloor}{j}}{3^j}
=
\left(0.9465045768\ldots\right)^{j+o(j)}.
\tag{11}
\]

Since the right side is summable, Borel\textendash{}Cantelli gives:


\begin{quote}
For Haar-almost every 3-adic endpoint \(X\), there are growth-compatible integral inverse paths at only finitely many depths.

\end{quote}


This is a rigorous subcritical first-moment theorem. It requires no independence.

But once again it is annealed over \(X\). A hypothetical exceptional Collatz orbit supplies a correlated sequence of moving endpoints \(X_j\), chosen precisely so that one exceptional path exists at depth \(j\). Conditioning a subcritical branching model on such survival creates an ``immortal spine''; the annealed estimate (11) does not control that spine.


\subsubsection*{6. The fixed-endpoint version actually terminates}


The fixed-endpoint problem can be settled without probability.


\paragraph*{Proposition}


For every fixed positive integer \(X\), there are only finitely many pairs \((Y,j)\) such that

\[
T^j(Y)=X
\]

and the multiplicative coefficient of this segment is at least \(1\).


\paragraph*{Proof}


Write

\[
X=q_jY+E_j,
\qquad q_j\geq1,\quad E_j\geq0.
\]

If \(E_j=0\), the segment contains only even steps, in which case \(q_j=2^{-j}<1\). Hence \(E_j>0\), and therefore

\[
Y<X.
\]

There are only finitely many possible starting values \(Y\).

For a fixed \(Y<X\), the deterministic orbit of \(Y\) can hit \(X\) at most once unless \(X\) lies on a cycle. If \(X\) lies on a positive cycle of period \(p\), then its period identity has the form

\[
X=q_pX+E_p,
\qquad E_p>0,
\]

and consequently \(q_p<1\). Each additional turn around the cycle multiplies the coefficient by \(q_p\), so after finitely many turns the total coefficient is below \(1\). Therefore only finitely many lengths \(j\) qualify. \ensuremath{\blacksquare}

So the real problem is not extinction above a fixed root. It is obtaining an extinction bound \textbf{uniform enough for the endpoint to move with the orbit}.


\subsubsection*{7. The second gap: coefficient decline is not actual descent}


Even eliminating every positive stationary ray satisfying \(q_j\geq1\) forever would not yet prove Collatz.

The additive remainder can keep an iterate above its starting value after the coefficient has fallen below \(1\). For example, there are finite ``paradoxical'' segments with

\[
\frac{3^{S_j}}{2^j}<1
\qquad\text{but}\qquad
T^j(n)\geq n.
\]

Nontrivial cycles would necessarily have this property over one period. [\hypertarget{cit:R40:4}{}\hyperlink{ref:R40:4}{4}]

Thus the coefficient stopping time and the ordinary stopping time are distinct notions; current finite coefficient-tree descriptions do not prove either that every coefficient stopping time is finite or that the coefficient and classical stopping times agree. [\hypertarget{cit:R40:5}{}\hyperlink{ref:R40:5}{5}]

A complete proof using the proposed tree would need both:

\[
\begin{aligned}
&\text{(I) no positive stationary coefficient-barrier ray;}\\
&\text{(II) every least counterexample produces such a ray,}\\
&\hspace{1.4cm}\text{or else its first coefficient crossing gives actual descent.}
\end{aligned}
\]

The second statement is precisely the all-depth first-crossing problem that the previous exact computation verified only finitely far.


\subsubsection*{8. Why terminal locality does not create finite dependence}


Your existing terminal-locality framework gives an exact identity for a compatible competing ancestor:

\[
Y-R=\frac{\Sigma(u,\tau,c)+h(\tau,b)}{2^c}.
\tag{12}
\]

Here \(h\) is determined by the terminal \(3\)-adic compatibility fibre, but the surplus \(\Sigma\) depends on the entire prefix. Thus even perfect suffix mixing does not determine whether the competitor lies below the original seed. [\hypertarget{cit:R40:6}{}\hyperlink{ref:R40:6}{6}]

The required qualitative statement is exactly that every positive stationary retained ray eventually has a compatible competitor with negative total score. That synchronization theorem remains the stated all-depth gap. [\hypertarget{cit:R40:7}{}\hyperlink{ref:R40:7}{7}]

The explicit length-\(42\) strict-barrier word with optimal competitor cutoff \(814\) also shows that a small fixed cutoff or bounded suffix-only compression cannot supply the missing block domination. [\hypertarget{cit:R40:8}{}\hyperlink{ref:R40:8}{8}] The conclusion of that computation is explicitly that the terminal congruence data must be coupled to unbounded prefix height information. [\hypertarget{cit:R40:9}{}\hyperlink{ref:R40:9}{9}]


\subsubsection*{9. A repaired percolation criterion that really would prove Collatz}


The correct pointwise analogue of subcritical domination is the following.

For \(N,k\geq1\), define

\[
E_k(N)=
\left\{
1<n\leq N:
T^i(n)\geq n\text{ for every }1\leq i\leq k
\right\}.
\]


\paragraph*{Uniform quenched extinction criterion}


Suppose there exist constants

\[
C<\infty,\qquad a<\infty,\qquad 0<\rho<1
\]

such that

\[
|E_k(N)|\leq C N^a\rho^k
\qquad
\text{for every }N,k.
\tag{13}
\]

Then the Collatz conjecture follows.


\paragraph*{Proof}


Fix \(N\). The right side of (13) tends to zero as \(k\to\infty\). Since \(|E_k(N)|\) is an integer, it must eventually equal zero. Therefore every \(n\leq N\) eventually reaches a value below itself.

As \(N\) is arbitrary, every \(n>1\) has finite stopping time. Strong induction now proves that every positive integer reaches \(1\): once \(n\) reaches \(m<n\), the induction hypothesis applies to \(m\). \ensuremath{\blacksquare}

Estimate (13) is a genuine pointwise, subcritical-percolation theorem. It is much stronger than

\[
\frac{|E_k(N)|}{N}\longrightarrow0
\]

in a density sense. Density estimates permit a single nested exceptional ray; (13) does not.

An equivalent tree formulation would be to construct positive weights \(W(w)\), a block length \(L\), and \(\rho<1\) such that

\[
\sum_{\substack{u\succeq w\\|u|-|w|=L\\u\text{ open}}}
W(u)
\leq
\rho W(w)
\tag{14}
\]

for every open node, while \(W(w_k)\) stays bounded below along every eventually stationary positive ray. Iterating (14) would make such a ray impossible.

The weight cannot depend only on a bounded suffix or on the multiplicative surplus. It must detect at least:

\[
\text{canonical-seed stabilization}
\;+\;
\text{terminal }3\text{-adic fibre}
\;+\;
\text{prefix archimedean surplus}.
\]

Constructing that weight is exactly the unresolved synchronization problem, rather than a consequence of ordinary dependent-percolation domination.


\subsection*{Final assessment}


The approach yields two rigorous results:

\[
\boxed{
\begin{aligned}
&\text{Haar mass of forward growth-compatible cylinders decays like }
0.9659065^k,\\
&\text{the annealed number of accelerated integral inverse words decays like }
0.9465046^j.
\end{aligned}}
\]

But neither is pointwise. Meanwhile, the deterministic forward tree contains the infinite all-ones ray and arbitrarily deep complete open cones, so it cannot be dominated pathwise by any subcritical percolation process, at any fixed block scale.

Therefore Approach 35, in its stated form, does not complete a Collatz proof. Its viable repaired form is the uniform quenched estimate (13), or the stationary-ray weighted certificate (14). Proving either would indeed give a full proof, but the needed estimate is precisely the global arithmetic synchronization step not supplied by percolation or finite integrality alone.
\subsection*{References}
\begin{enumerate}[leftmargin=2.2em]
\item \hypertarget{ref:R40:1}{}\hyperlink{cit:R40:1}{\texttt{turn631843view0}}
\item \hypertarget{ref:R40:2}{}\hyperlink{cit:R40:2}{\texttt{turn631843view1}}
\item \hypertarget{ref:R40:3}{}\hyperlink{cit:R40:3}{\texttt{turn167226search6}}
\item \hypertarget{ref:R40:4}{}\hyperlink{cit:R40:4}{\texttt{turn753536view0}}
\item \hypertarget{ref:R40:5}{}\hyperlink{cit:R40:5}{\texttt{turn100079view2}}
\item \hypertarget{ref:R40:6}{}\hyperlink{cit:R40:6}{\texttt{turn2file1 L227-L260}}
\item \hypertarget{ref:R40:7}{}\hyperlink{cit:R40:7}{\texttt{turn2file1 L324-L338}}
\item \hypertarget{ref:R40:8}{}\hyperlink{cit:R40:8}{\texttt{turn2file0 L43-L94}}
\item \hypertarget{ref:R40:9}{}\hyperlink{cit:R40:9}{\texttt{turn2file0 L172-L181}}
\end{enumerate}

\clearpage
\section{Round 41: approach 36 --- An algorithmic Lovasz local lemma}
\subsection*{Prompt}
\begin{quote}
Complete a full proof of the Collatz conjecture using the following approach:

\begin{enumerate}
\item An algorithmic Lovasz local lemma
\end{enumerate}

\textbf{Speculative \textperiodcentered{} P}

Construct local "death events" that cover every infinite critical word. A dependency-graph argument or resampling certificate could show that no assignment avoids all death events, even though each individual condition is local.

Take as long as you need.
\end{quote}
\subsection*{Response}

\subsubsection*{Verdict}


I did not obtain a valid unconditional proof of the Collatz conjecture from this approach. The literal Lov\'asz-local-lemma step has the \textbf{opposite logical direction} from the one required. After repairing that defect, the approach becomes a precise and potentially computer-verifiable \textbf{ranked covering certificate}, but constructing that certificate requires exactly the local\textendash{}global synchronization theorem that remains unproved.

The analysis does, however, yield several rigorous conclusions:

\begin{enumerate}
\item A direct Moser\textendash{}Tardos or Shearer argument cannot prove that the death events cover every critical word.
\item Purely suffix-local death events are insufficient: the sign deciding descent depends on the entire prefix.
\item The symbolic critical-word space contains exact infinite survivors that correspond to negative or genuinely \(2\)-adic seeds rather than positive integers.
\item There is a corrected ``dual witness-tree'' certificate which, if found, would prove the conjecture after separately handling coefficient crossings.
\end{enumerate}

What follows is the complete derivation.


\bigskip



\subsubsection*{1. Exact formulation of a death event}


Use the accelerated map on positive odd integers

\[
U(x)=\frac{3x+1}{2^{\nu_2(3x+1)}}.
\]

For a valuation word

\[
w=(a_0,\ldots ,a_{j-1}),\qquad a_i\ge 1,
\]

put

\[
A_t=\sum_{i<t}a_i,
\qquad
C(w)=\sum_{i=0}^{j-1}3^{j-1-i}2^{A_i}.
\]

If \(w\) is followed from \(n\), ending at \(X\), then

\[
2^{A_j}X=3^j n+C(w).
\tag{1}
\]

Moreover, \(w\) is realized by exactly one odd residue class modulo \(2^{A_j+1}\). Denote its least positive representative by \(R(w)\). These affine and realization statements are exact. [\hypertarget{cit:R41:1}{}\hyperlink{ref:R41:1}{1}]

Call \(w\) \textbf{critical} when

\[
2^{A_t}\le 3^t
\qquad (1\le t\le j).
\tag{2}
\]

The strict version uses \(<\). These are precisely the finite words whose multiplicative coefficient has not yet become contractive.

Now split

\[
w=u\tau,
\]

where \(\tau\) has length \(k\). Let \(Z\) be the state after following the prefix \(u\) from \(R=R(w)\). Consider a competing inverse word \(b\) of length \(k\), with

\[
c=A(\tau)-A(b)>0.
\]

Compatibility of \(b\) with the same endpoint is equivalent to

\[
C(b)\equiv 2^{-c}C(\tau)\pmod {3^k}.
\tag{3}
\]

When this holds, define

\[
h(\tau,b)=
\frac{C(\tau)-2^cC(b)}{3^k}\in\mathbb Z.
\tag{4}
\]

If \(Y\) is the ancestor obtained through \(b\), then

\[
2^cY=Z+h(\tau,b).
\tag{5}
\]

Defining the prefix surplus

\[
\Sigma(u,\tau,c)=Z-2^cR,
\tag{6}
\]

we obtain the exact comparison

\[
Y-R=\frac{\Sigma(u,\tau,c)+h(\tau,b)}{2^c}.
\tag{7}
\]

Thus a genuine death certificate consists of a compatible \(b\) for which

\[
0<Y<R.
\tag{8}
\]

For a least counterexample, such a \(Y\) is impossible: \(Y\) merges into the nonterminating orbit while being smaller than the least counterexample. The compatibility term \(h\) is terminal and \(3\)-adic, while \(\Sigma\) depends on the entire prefix. [\hypertarget{cit:R41:2}{}\hyperlink{ref:R41:2}{2}]


\bigskip



\subsubsection*{2. The polarity obstruction for the Lov\'asz local lemma}


Let \(\Omega\) be a space of infinite word assignments, and let \(D_i\) be the proposed death events. The required covering statement is

\[
\Omega=\bigcup_iD_i,
\]

or equivalently

\[
\bigcap_i D_i^{\,c}=\varnothing.
\tag{9}
\]

Under any full-support probability measure, one would need

\[
\mathbb P\!\left(\bigcap_iD_i^{\,c}\right)=0.
\tag{10}
\]

But the Lov\'asz local lemma is designed to prove the reverse conclusion. Under its probability and dependency hypotheses, it gives

\[
\mathbb P\!\left(\bigcap_iD_i^{\,c}\right)>0,
\tag{11}
\]

and algorithmic versions construct an assignment avoiding every \(D_i\). [\hypertarget{cit:R41:3}{}\hyperlink{ref:R41:3}{3}, \hypertarget{cit:R41:4}{}\hyperlink{ref:R41:4}{4}]

Consequently:


\begin{quote}
\textbf{Polarity proposition.}   Applying the algorithmic Lov\'asz local lemma directly to the Collatz death events, with the death events treated as bad events, would construct a symbolic word that survives all of them. It cannot prove that no survivor exists.

\end{quote}


Applying the LLL to the complementary events \(D_i^c\) does not repair matters. It would produce one assignment lying in every \(D_i\), whereas the required statement is that \textbf{every} assignment lies in at least one \(D_i\).

Nor can one take the contrapositive of an LLL numerical condition. The failure of an inequality such as \(epd<1\) does not imply that the corresponding constraint system is unsatisfiable. Even Shearer-type criteria concern positivity of avoidance probabilities based on prescribed dependency information; failure of those criteria is not a refutation of a particular arithmetic event system.

For computable infinite event systems, effective versions of the LLL can even produce computable infinite avoiders under suitable slack. Thus satisfying an infinite algorithmic-LLL hypothesis here would certify a symbolic survivor rather than eliminate one. [\hypertarget{cit:R41:5}{}\hyperlink{ref:R41:5}{5}]

This does not make probabilistic ideas useless, but the random object must be changed. One could randomize the \textbf{choice of certificates} and use the LLL to select a family covering all finite states. The eventual proof would then be a deterministic covering certificate, not an LLL avoidance conclusion about Collatz words.


\bigskip



\subsubsection*{3. Why the death events are not genuinely local}


Equation (3) is terminal-local. Equation (7) is not.

The exact terminal-locality theorem says that, for a suffix \(v\) of length \(r\),

\[
3^r\mid X+1
\quad\Longleftrightarrow\quad
3^r\mid C(v)+2^{A(v)}.
\tag{12}
\]

The prefix and starting value disappear completely from this congruence. There is also an exact right-to-left carry recursion deciding it. [\hypertarget{cit:R41:6}{}\hyperlink{ref:R41:6}{6}]

However, descent requires the sign of

\[
\Sigma+h,
\]

not merely the existence or value of \(h\). Two strict-barrier words of length \(23\) were found with the same final suffix

\[
(3,1,1,2,1,2),
\]

and the same local length, exponent, and slope data, but their prefix surpluses are

\[
Z_1-16R_1=13,
\qquad
Z_2-16R_2=-6\,228\,506\,163.
\tag{13}
\]

Thus the part of the score omitted by terminal locality can move by billions while the entire terminal compatibility problem remains unchanged. [\hypertarget{cit:R41:7}{}\hyperlink{ref:R41:7}{7}]

This produces a dependency dilemma.

If an event is allowed to inspect only the final \(r\) valuations, it has a bounded dependency neighborhood, but it cannot decide the sign of the descent score.

If an event includes the prefix information required to calculate \(\Sigma\), then an event at depth \(j\) depends on

\[
a_0,\ldots ,a_{j-1}.
\]

Events at two different depths therefore share almost all their variables. In the natural variable-dependency graph, the degree is unbounded and the events are nested rather than weakly dependent.

Grouping levels into blocks does not remove this problem: the prefix surplus is propagated through every preceding block. A successful block construction would need a finite state summarizing the prefix strongly enough to control \(\Sigma\). Producing such a finite summary is already the missing synchronization theorem.

The obstruction is not merely hypothetical. An explicit strict-barrier word of length \(42\) has optimal competing-ancestor cutoff exactly \(814\): no compatible competitor kills every realization beginning at \(813\), while one does kill every realization beginning at \(814\). This rules out the simplest bounded-small-cutoff version of the approach. [\hypertarget{cit:R41:8}{}\hyperlink{ref:R41:8}{8}]


\bigskip



\subsubsection*{4. Positive-integral realizability is a global tail condition}


There is a second locality obstruction that is independent of the descent score.

For an infinite valuation word \(a=(a_0,a_1,\ldots)\), let \(w_j\) be its prefix of length \(j\), and let

\[
R_j=R(w_j)
\]

be the least positive representative of its realization class. Because the realization classes are nested,

\[
R_{j+1}\equiv R_j\pmod {2^{A_j+1}}.
\tag{14}
\]

Since \(A_j\to\infty\), these congruences determine a unique \(2\)-adic integer

\[
\xi(a)\in\mathbb Z_2.
\]


\paragraph*{Proposition}


The infinite word \(a\) is the valuation word of one fixed positive odd integer \(n\) if and only if

\[
R_j=n
\]

for all sufficiently large \(j\).


\paragraph*{Proof}


If \(n\) realizes every prefix, then

\[
n\equiv R_j\pmod {2^{A_j+1}}.
\]

For sufficiently large \(j\), the modulus \(2^{A_j+1}\) exceeds \(n\). Since both \(n\) and \(R_j\) are positive representatives smaller than the modulus, they must be equal.

Conversely, if \(R_j=n\) eventually, then \(n\) lies in every sufficiently long realization class and therefore realizes every finite prefix. \ensuremath{\blacksquare}

Thus ordinary positive realizability is equivalent to \textbf{eventual stabilization} of a sequence of prefix representatives. It cannot be detected from any bounded window of valuation letters.

Every finite word has positive realizations, so no finite prefix can distinguish between:

\begin{itemize}
\item a word ultimately generated by one fixed positive integer, and
\item a word whose positive representatives continually lift to larger residue-class representatives and converge only \(2\)-adically.
\end{itemize}


\paragraph*{An exact spurious critical ray}


Consider

\[
a=(1,2,1,2,\ldots).
\]

It is strictly critical. Indeed,

\[
A_{2m}=3m,
\qquad
A_{2m+1}=3m+1,
\]

and hence

\[
2^{A_{2m}}=8^m<9^m=3^{2m},
\]

while

\[
2^{A_{2m+1}}=2\cdot8^m
   <3\cdot9^m=3^{2m+1}.
\]

For one block \((1,2)\),

\[
U^2(x)=\frac{9x+5}{8}.
\tag{15}
\]

The algebraic fixed point is \(x=-5\). Indeed,

\[
-5\longmapsto -7\longmapsto -5,
\]

with valuations \(1,2\). Therefore the unique \(2\)-adic seed associated with the infinite word is the ordinary negative integer \(-5\).

For the prefix \((1,2)^m\), the realization class is

\[
n\equiv -5\pmod {2^{3m+1}},
\]

so its least positive representative is

\[
R_{2m}=2^{3m+1}-5.
\tag{16}
\]

These representatives tend to infinity rather than stabilizing. Thus this is an exact infinite critical symbolic ray, but not a positive Collatz counterexample.

It also avoids every nontrivial terminal-to-\(1\) suffix event arising from (12). A suffix ending in \(2\) already fails the one-level condition because

\[
1+2^2\not\equiv0\pmod3.
\]

A suffix of length at least \(2\) ending in \(1\) has final pair \((2,1)\). The exact level-two classification says that for an even penultimate exponent, the final exponent must be \(5\bmod 6\), not \(1\bmod6\). [\hypertarget{cit:R41:9}{}\hyperlink{ref:R41:9}{9}]

Therefore terminal death events alone cannot cover the symbolic critical tree. They must be supplemented by a global condition detecting whether the nested \(2\)-adic representative stabilizes at a positive integer.


\bigskip



\subsubsection*{5. A finite-scale entropy obstruction}


There is also a quantitative reason not to expect compatible local competitors to be plentiful.

Let

\[
\alpha=\log_2 3.
\]

Suppose a length-\(k\) competitor is allowed to have total exponent at most

\[
M=\lfloor\alpha k\rfloor.
\]

The number of positive exponent words of length \(k\) and total exponent \(B\) is

\[
\binom{B-1}{k-1}.
\]

Hence the total number of candidate words with \(B\le M\) is

\[
\sum_{B=k}^{M}\binom{B-1}{k-1}
=
\binom{M}{k}.
\tag{17}
\]

Each candidate contributes at most one value of \(C(b)\bmod 3^k\). Therefore its image occupies at most

\[
\binom{M}{k}
\]

of the \(3^k\) possible residue classes.

By Stirling's formula,

\[
\binom{\lfloor\alpha k\rfloor}{k}^{1/k}
\longrightarrow
\gamma:=
\frac{\alpha^\alpha}{(\alpha-1)^{\alpha-1}}
\approx 2.8395137<3.
\tag{18}
\]

Consequently, the proportion of residue classes hit by these candidates is at most

\[
\left(\frac{\gamma}{3}+o(1)\right)^k
=
(0.9465046+o(1))^k.
\tag{19}
\]

Thus, at a fixed near-critical scale, the possible affine constants are exponentially sparse inside the full \(3^k\)-residue space.

This is not a proof that an actual Collatz endpoint avoids them: actual target residues are highly correlated with the prefix. But it shows precisely why a naive ``every word has many independent local killers'' estimate is unavailable. Any proof must exploit a special correlation between the prefix surplus, the terminal residue, and eventual ordinary stabilization.

That correlation is the synchronization problem.


\bigskip



\subsubsection*{6. The corrected replacement: a ranked covering certificate}


There is a logically sound version of the proposed strategy. It is closer to a cyclic termination proof or a finite refutation system than to the ordinary LLL.

Let \(\mathcal T\) be the finitely branching tree of critical prefixes. For a fixed finite verification threshold \(M\), call a node \(w\) \textbf{\(M\)-dead} when an exact certificate proves that every positive realization \(n\ge M\) of \(w\) has a smaller positive integer whose orbit merges into its orbit.

A certificate may use:

\begin{itemize}
\item a direct iterate below \(n\);
\item a compatible competing ancestor \(0<Y<n\);
\item an exact residue obstruction;
\item a previously verified descendant class.
\end{itemize}

The following is the appropriate dual analogue of a witness-tree theorem.


\paragraph*{Ranked-cover theorem}


Suppose the nondead nodes of \(\mathcal T\) admit a finite-state labeling

\[
q:\mathcal T\setminus\mathcal D_M\longrightarrow Q
\]

and a rank

\[
\rho:Q\longrightarrow W
\]

into a well-founded ordered set \(W\). Suppose that for every state \(q\in Q\) there is a finite prefix code \(\mathcal P_q\) of extension blocks such that every feasible infinite continuation from a node \(w\) with label \(q\) begins with some \(v\in\mathcal P_q\), and every such feasible leaf \(wv\) satisfies at least one of:

\[
\begin{array}{ll}
\textup{(D)} & wv\text{ is }M\text{-dead};\\[1mm]
\textup{(L)} & R(wv)>R(w);\\[1mm]
\textup{(R)} & \rho(q(wv))<\rho(q(w)).
\end{array}
\tag{20}
\]

Then no nondead infinite critical ray can be realized by a fixed positive integer \(n\ge M\).


\paragraph*{Proof}


Assume such a ray exists, generated by a positive integer \(n\). By the stabilization proposition,

\[
R(w_j)=n
\]

for every sufficiently long prefix \(w_j\).

Begin the ranked-cover decomposition after this stabilization point. Alternative (D) cannot occur because the ray is assumed nondead. Alternative (L) cannot occur because the representatives have stabilized. Therefore every successive covered block must satisfy (R).

This produces an infinite strictly descending sequence

\[
\rho(q_0)>\rho(q_1)>\rho(q_2)>\cdots
\]

in the well-founded order \(W\), a contradiction. \ensuremath{\blacksquare}

This theorem handles the spurious \(2\)-adic rays correctly. Such a ray may repeatedly take lift transitions (L), whereas a ray coming from one fixed positive integer cannot do so forever.

It is also algorithmic in the relevant sense. A finite certificate would contain:

\begin{enumerate}
\item the state abstraction \(Q\);
\item the finite block covers \(\mathcal P_q\);
\item exact arithmetic verifications of every D/L/R leaf;
\item the well-founded rank inequalities.
\end{enumerate}

A small checker could verify the whole object without trusting the search procedure.


\bigskip



\subsubsection*{7. Where an algorithmic LLL could legitimately enter}


One could use Moser\textendash{}Tardos as a \textbf{certificate-discovery algorithm}.

For each abstract state \(q\), introduce variables selecting candidate extension blocks and candidate arithmetic certificates. Define a flaw to mean that some feasible continuation from \(q\) is not covered by any D/L/R leaf. If these flaws have sufficiently small probabilities and sufficiently sparse dependencies, an algorithmic LLL could select a complete covering table.

But two estimates would first have to be proved:

\[
\Pr(\text{a feasible continuation is uncovered})\ll1,
\]

and

\[
\deg(\text{the certificate-selection dependency graph})<\infty
\]

with adequate quantitative margin.

The local\textendash{}global score shows why neither estimate is presently available. The validity of a D-leaf depends on \(\Sigma+h\), and \(\Sigma\) carries the entire prefix. The unbounded \(3\)-adic carry runs similarly prevent a fixed-modulus automaton from deciding all relevant transitions; the existing carry analysis identifies the missing object as a finite-basis or synchronization invariant. [\hypertarget{cit:R41:10}{}\hyperlink{ref:R41:10}{10}]

Thus the LLL could search for a finite ranked cover \textbf{after} the correct state space is discovered. It does not supply that state space by itself.


\bigskip



\subsubsection*{8. Exact remaining theorem}


The arithmetic core can now be stated without probabilistic language.


\begin{quote}
\textbf{Stationary-ray synchronization theorem.}   Let \(a=(a_i)_{i\ge0}\) be an infinite critical valuation ray whose nested realization representatives eventually stabilize at an ordinary positive integer \(n>1\). Then some finite prefix \(w=u\tau\) admits a competing word \(b\) such that \[
> c=A(\tau)-A(b)>0,
> \] \[
> 3^{|\tau|}
> \mid C(\tau)-2^cC(b),
> \] and, with \[
> h=\frac{C(\tau)-2^cC(b)}{3^{|\tau|}},
> \] \[
> 0<Z+h<2^c n.
> \tag{21}
> \] Equivalently, the associated compatible ancestor \[
> Y=\frac{Z+h}{2^c}
> \] satisfies \[
> 0<Y<n.
> \]

\end{quote}


This is exactly the sought death event along a fixed positive stationary ray. The existing structural work reduces the missing comparison to the sign of

\[
\Sigma(u,\tau,c)+h_{\max}(\tau,B),
\]

but does not control that sign at all depths. [\hypertarget{cit:R41:11}{}\hyperlink{ref:R41:11}{11}]

Even this theorem would need to be paired with an all-depth coefficient-crossing result. Current computation supports the claim that every nontrivial first coefficient crossing forces descent, but it has only been checked finitely; it is not presently an induction theorem. [\hypertarget{cit:R41:12}{}\hyperlink{ref:R41:12}{12}]

Hence a complete proof through this architecture requires two new results:

\[
\boxed{\text{first-crossing elimination}}
\]

and

\[
\boxed{\text{a finite ranked cover proving stationary-ray synchronization}.}
\]

The algorithmic Lov\'asz local lemma may be useful for discovering the second certificate, but the ordinary LLL avoidance conclusion cannot replace it. Until those two boxed statements are proved, the argument is a rigorous reduction and certificate framework rather than a proof of the Collatz conjecture.
\subsection*{References}
\begin{enumerate}[leftmargin=2.2em]
\item \hypertarget{ref:R41:1}{}\hyperlink{cit:R41:1}{\texttt{turn2file0 L83-L135}}
\item \hypertarget{ref:R41:2}{}\hyperlink{cit:R41:2}{\texttt{turn2file0 L289-L377}}
\item \hypertarget{ref:R41:3}{}\hyperlink{cit:R41:3}{\texttt{turn993955search5}}
\item \hypertarget{ref:R41:4}{}\hyperlink{cit:R41:4}{\texttt{turn993955search2}}
\item \hypertarget{ref:R41:5}{}\hyperlink{cit:R41:5}{\texttt{turn993955search1}}
\item \hypertarget{ref:R41:6}{}\hyperlink{cit:R41:6}{\texttt{turn2file0 L202-L287}}
\item \hypertarget{ref:R41:7}{}\hyperlink{cit:R41:7}{\texttt{turn3file0 L260-L280}}
\item \hypertarget{ref:R41:8}{}\hyperlink{cit:R41:8}{\texttt{turn3file0 L310-L424}}
\item \hypertarget{ref:R41:9}{}\hyperlink{cit:R41:9}{\texttt{turn2file0 L258-L285}}
\item \hypertarget{ref:R41:10}{}\hyperlink{cit:R41:10}{\texttt{turn2file1 L942-L972}}
\item \hypertarget{ref:R41:11}{}\hyperlink{cit:R41:11}{\texttt{turn2file0 L441-L455}}
\item \hypertarget{ref:R41:12}{}\hyperlink{cit:R41:12}{\texttt{turn2file0 L451-L453}}
\end{enumerate}

\clearpage
\section{Round 42: approach 37 --- Coding-theoretic minimum distance}
\subsection*{Prompt}
\begin{quote}
Complete a full proof of the Collatz conjecture using the following approach:

\begin{enumerate}
\item Coding-theoretic minimum distance
\end{enumerate}

\textbf{Speculative \textperiodcentered{} P}

View the valuation/carry words generated by different seed classes as a code. Prove that sustained-growth words lie in a region whose distance from the realizable code exceeds the covering radius.

Take as long as you need.
\end{quote}
\subsection*{Response}

\subsection*{Coding-theoretic audit of the Collatz map}


I did not obtain a valid full proof. More importantly, the proposed argument has two structural defects that prevent it from working literally:

\begin{enumerate}
\item The set of finite realizable parity or valuation words is the entire ambient word space, so it has no redundancy and hence no useful minimum distance.
\item A set cannot have distance from a code exceeding that code's covering radius: by definition, \[
   R(C)=\max_{x}d(x,C).
   \]
\end{enumerate}

There is, however, a mathematically meaningful repair. The correct object is a \textbf{terminated nonlinear code}, and the relevant invariant is not ordinary minimum distance but a zero-syndrome or list-decoding estimate. I will derive that formulation precisely and identify the exact theorem that would finish Collatz.


\bigskip



\subsubsection*{1. The finite realizability code is trivial}


Use the shortcut map \[
T(n)=
\begin{cases}
n/2,&n\equiv0\pmod2,\\[2mm]
(3n+1)/2,&n\equiv1\pmod2.
\end{cases}
\]

For \(K\geq1\), define \[
Q_K(n)=\bigl(n\bmod2,T(n)\bmod2,\ldots,T^{K-1}(n)\bmod2\bigr)
       \in\{0,1\}^K.
\]


\paragraph*{Theorem 1: finite parity saturation}


For all integers \(m,n\), \[
Q_K(m)=Q_K(n)
\quad\Longleftrightarrow\quad
m\equiv n\pmod{2^K}.
\]

Consequently, \(Q_K\) induces a bijection \[
\mathbb Z/2^K\mathbb Z\longrightarrow\{0,1\}^K.
\]


\paragraph*{Proof}


The assertion is immediate for \(K=1\).

Suppose it is known for \(K-1\). If \(m,n\) have the same initial parity \(b\), then \[
T(m)-T(n)=
\begin{cases}
(m-n)/2,&b=0,\\[1mm]
3(m-n)/2,&b=1.
\end{cases}
\] Since \(3\) is odd, \[
T(m)\equiv T(n)\pmod{2^{K-1}}
\quad\Longleftrightarrow\quad
m\equiv n\pmod{2^K}.
\] Applying the induction hypothesis to the remaining \(K-1\) parity symbols proves the result. \ensuremath{\blacksquare}

Thus every finite binary word is realized by one residue class modulo \(2^K\), and hence by infinitely many positive integers. This well-known finite bijection extends to a bijective 2-adic parity transform \(Q:\mathbb Z_2\to\{0,1\}^{\mathbb N}\), which is in fact a 2-adic isometry. [\hypertarget{cit:R42:1}{}\hyperlink{ref:R42:1}{1}]

Therefore, if \[
\mathcal C_K
=
\{Q_K(n):n\in\mathbb N\},
\] then \[
\boxed{\mathcal C_K=\{0,1\}^K.}
\]

Its Hamming minimum distance is \(1\), and its covering radius is \(0\).


\bigskip



\subsubsection*{2. Accelerated valuation words have the same defect}


On odd integers let \[
U(n)=\frac{3n+1}{2^{a(n)}},
\qquad
a(n)=v_2(3n+1).
\]

Given a finite valuation word \[
a=(a_0,\ldots,a_{r-1}),\qquad a_i\geq1,
\] write \[
A_i=\sum_{t=0}^{i-1}a_t,\qquad A=A_r.
\]

This word corresponds to the shortcut parity word \[
1\,0^{a_0-1}\,1\,0^{a_1-1}\cdots
1\,0^{a_{r-1}-1}\,1,
\] which has length \(A+1\). By Theorem 1, this binary word is realized by a unique residue class modulo \(2^{A+1}\). Hence:


\paragraph*{Corollary 2}


Every finite word in \(\mathbb Z_{\geq1}^r\) occurs as the first \(r\) accelerated valuations of infinitely many positive odd integers.

More explicitly, define \[
C(a)=\sum_{i=0}^{r-1}3^{r-1-i}2^{A_i}.
\] Then \[
2^A U^r(n)=3^r n+C(a).
\] The condition that \(U^r(n)\) be odd gives the unique congruence \[
\boxed{
n\equiv
3^{-r}\bigl(2^A-C(a)\bigr)
\pmod{2^{A+1}}.
}
\]

So adding the valuation symbols does not create a code with redundancy. Adding deterministic carry symbols does not help either: it merely decorates each already-realizable word by the carries generated by its realizing seed.

This gives a general no-go statement.


\paragraph*{Finite-local coding no-go}


Any code whose information consists only of:

\begin{itemize}
\item a finite parity or valuation word, and
\item finitely many carries or residues determined from that word and a realizing residue class,
\end{itemize}

contains an actual codeword above every finite parity or valuation word.

Therefore, every nonempty finite sustained-growth language intersects the realizable code. Its distance from that code is zero.


\bigskip



\subsubsection*{3. The covering-radius formulation is formally impossible}


For any code \(C\) in a metric space \(X\), its covering radius is \[
R(C)=\sup_{x\in X}d(x,C).
\] For every nonempty \(G\subseteq X\), \[
d(G,C)=\inf_{g\in G}d(g,C)\leq R(C).
\]

Thus the requested inequality \[
d(G,C)>R(C)
\] cannot hold.

The appropriate replacement is a \textbf{syndrome-exclusion statement} \[
G\cap C=\varnothing,
\] or quantitatively a list-decoding estimate bounding \(|G\cap C|\).


\bigskip



\subsubsection*{4. The correct repair: terminate the code}


Let \[
\Phi_K=Q_K^{-1}:\{0,1\}^K\longrightarrow
\{0,\ldots,2^K-1\}.
\]

If \(p=(p_0,\ldots,p_{K-1})\) and \[
s_i=p_0+\cdots+p_i,
\] then the inverse parity formula gives \[
\Phi_K(p)\equiv
-\sum_{i=0}^{K-1}p_i2^i3^{-s_i}
\pmod{2^K},
\] where the least nonnegative representative is taken.

For \(L\leq K\), define the terminated code \[
\boxed{
\mathcal C_{L,K}
=
\{p\in\{0,1\}^K:\Phi_K(p)<2^L\}.
}
\]

It has exactly \(2^L\) words. Equivalently, the binary digits \(L,\ldots,K-1\) of \(\Phi_K(p)\) are zero.

Define its nonlinear syndrome \[
\sigma_{L,K}(p)
=
\left\lfloor\frac{\Phi_K(p)}{2^L}\right\rfloor
\in\{0,\ldots,2^{K-L}-1\}.
\] Then \[
p\in\mathcal C_{L,K}
\quad\Longleftrightarrow\quad
\sigma_{L,K}(p)=0.
\]

This is the genuine source of redundancy: a positive seed has a \textbf{finite binary expansion}, so all sufficiently high decoded seed bits vanish.


\paragraph*{Ordinary minimum distance still fails}


For \(K\geq2\), \[
Q_K(1)=1010\cdots,
\qquad
Q_K(4)=0010\cdots.
\] The two orbits merge at \(2\), so their parity words differ in exactly one position. Therefore \[
d_{\min}(\mathcal C_{L,K})=1
\] whenever \(L\geq3\).

Likewise, for the accelerated map, \[
U(3)=5,\qquad v_2(3\cdot3+1)=1,
\] while \[
U(13)=5,\qquad v_2(3\cdot13+1)=3.
\] Their valuation sequences differ only in the first symbol and then coincide.

Thus neither block Hamming distance nor asymptotic relative distance is the right invariant. The proper coding-theoretic notion is analogous to \textbf{active column distance} or \textbf{noncatastrophicity}: can a growth-constrained infinite input produce an eventually zero decoded output?


\bigskip



\subsubsection*{5. The coefficient-growth cone and its exact entropy}


For \(p\in\{0,1\}^K\), put \[
S_k(p)=\sum_{i=0}^{k-1}p_i.
\] The multiplicative coefficient after \(k\) shortcut steps is \[
q_k(p)=\frac{3^{S_k(p)}}{2^k}.
\]

Define the sustained coefficient-growth cone \[
\mathcal G_K
=
\left\{
p:
3^{S_k(p)}\geq2^k
\text{ for every }1\leq k\leq K
\right\}.
\]

Writing \[
\vartheta=\log_3 2=0.6309297535\ldots,
\] this condition is \[
S_k(p)\geq\vartheta k
\qquad(1\leq k\leq K).
\]

Standard ballot estimates give \[
|\mathcal G_K|
=
2^{H_2(\vartheta)K+o(K)},
\] where \[
H_2(x)=-x\log_2x-(1-x)\log_2(1-x).
\] Numerically, \[
H_2(\vartheta)
=
0.9499555272\ldots.
\] Hence the entropy deficit is \[
\boxed{
\delta
=
1-H_2(\vartheta)
=
0.0500444728\ldots.
}
\]


\paragraph*{Exact counting before termination becomes visible}


For \(K\leq L\), each parity word of length \(K\) is produced by exactly \(2^{L-K}\) seeds in \([0,2^L)\). Therefore \[
\#\left\{
0\leq n<2^L:
Q_K(n)\in\mathcal G_K
\right\}
=
2^{L-K}|\mathcal G_K|.
\] Consequently, \[
\#\{\text{coefficient-growth survivors}\}
=
2^{L-\delta K+o(K)}
\qquad(K\leq L).
\]

This is an exact ensemble estimate, not a randomness heuristic.

The hard regime begins only when \(K>L\), because the zero-tail parity checks then start constraining the word. If the nonlinear syndrome behaved like a random balanced map, one would still expect \[
|\mathcal C_{L,K}\cap\mathcal G_K|
\approx
2^{L-\delta K}.
\] That predicts extinction near \[
K\approx\frac L\delta
=
19.98222668\ldots\,L.
\]

But no deterministic estimate of this strength follows from bijectivity or entropy alone.

The first genuinely sufficient coding theorem would be:


\paragraph*{Coefficient list-decoding conjecture}


There exist \(C,\eta>0\) such that \[
\boxed{
|\mathcal C_{L,K}\cap\mathcal G_K|
\leq
C\,2^{L-\eta K}
}
\tag{CLD}
\] uniformly for \(K\geq L\).

For fixed \(L\), the right side eventually becomes smaller than \(1\), excluding every fixed seed from maintaining \(q_k\geq1\) forever.

I did not prove (CLD). It is precisely a deterministic pointwise-mixing theorem for the high decoded seed bits.


\bigskip



\subsubsection*{6. Adding the \(3\)-adic endpoint gives a genuine sparse code}


For a valuation word \(a=(a_0,\ldots,a_{r-1})\), define \[
\Psi_r(a)
=
2^{-A}C(a)\pmod{3^r}.
\] The affine identity shows that \[
U^r(n)\equiv\Psi_r(a)\pmod{3^r}.
\]

Thus \(\Psi_r(a)\) is the \(3\)-adic endpoint syndrome determined by the valuation/carry word. It lies in \[
(\mathbb Z/3^r\mathbb Z)^\times,
\] a set of size \(2\cdot3^{r-1}\).

Let \[
\alpha=\log_2 3=1.5849625007\ldots.
\] A coefficient-growth-compatible valuation word has total cost approximately \[
A\leq\alpha r.
\]

The number of length-\(r\) positive compositions with total \(A\leq\alpha r\) is \[
N_r
=
\sum_{A=r}^{\lfloor\alpha r\rfloor}
\binom{A-1}{r-1}.
\] Stirling's formula gives \[
N_r
=
\exp\bigl(f(\alpha)r+o(r)\bigr),
\] where \[
f(\beta)
=
\beta\log\beta-(\beta-1)\log(\beta-1).
\] Numerically, \[
f(\alpha)=1.0436328159\ldots,
\qquad
e^{f(\alpha)}=2.8395137305\ldots<3.
\]

Moreover, \[
\frac{f(\alpha)}{\log3}
=
H_2(\log_3 2)
=
0.9499555272\ldots.
\]

So the binary growth cone and the \(3\)-adic low-cost endpoint code have exactly the same asymptotic rate deficit: \[
1-\frac{f(\alpha)}{\log3}
=
\delta.
\]

This is the most promising genuine coding structure in the proposal.


\paragraph*{But sparsity is only an average statement}


Only an exponentially small fraction, roughly \[
\left(\frac{2.8395}{3}\right)^r,
\] of all \(3\)-adic endpoint syndromes can be represented by low-cost valuation words.

Unfortunately, an actual growth-compatible orbit comes equipped with such a word. It therefore automatically selects one of these exceptional endpoint syndromes. Knowing that most syndromes are uncovered does not exclude the one dynamically selected syndrome.

This is exactly the distinction between:

\begin{itemize}
\item average covering of residue space, and
\item pointwise avoidance along one exceptional orbit.
\end{itemize}


\bigskip



\subsubsection*{7. The weighted covering radius lies on the wrong side of the threshold}


Define the weighted coset-leader cost \[
\ell_r(y)
=
\min\{A(a):\Psi_r(a)=y\}
\] and the weighted covering radius \[
R_r=\max_{y\in(\mathbb Z/3^r\mathbb Z)^\times}\ell_r(y).
\]

Let \(\beta_*\) solve \[
f(\beta_*)=\log3.
\] Numerically, \[
\boxed{
\beta_*=1.6417941208\ldots
}
\] whereas \[
\alpha=1.5849625007\ldots.
\]

Counting forces \[
\liminf_{r\to\infty}\frac{R_r}{r}\geq\beta_*>\alpha.
\]

Thus low-cost words below the growth threshold do not have enough cardinality to cover all endpoint syndromes. Ordinary covering-radius theory therefore cannot guarantee the existence of an alternative low-cost representative in every endpoint coset.

More importantly, for an actual sustained-growth word, its endpoint coset already has a representative of cost at most \(\alpha r\)\textemdash{}namely the word itself. Covering-radius information about the other cosets does not create a contradiction.


\bigskip



\subsubsection*{8. A genuine pointwise lower bound at coefficient-record times}


There is nevertheless a useful unconditional ``active distance'' estimate.

Let \[
x_{i+1}=U(x_i),\qquad
a_i=v_2(3x_i+1),
\qquad
A_j=\sum_{i=0}^{j-1}a_i,
\qquad
q_j=\frac{3^j}{2^{A_j}}.
\]

Suppose \(j\) is a record time: \[
q_j\geq q_i\qquad(0\leq i<j).
\]

For \(1\leq r\leq j\), put \[
S_r=A_j-A_{j-r},
\qquad
u_r=\frac{3^r}{2^{S_r}}=\frac{q_j}{q_{j-r}}.
\] Then \(u_r\geq1\).

Reindexing the exact affine carry formula gives \[
\boxed{
x_j=q_jx_0+\frac13\sum_{r=1}^j u_r.
}
\tag{1}
\]

Indeed, \[
x_j
=
\frac{3^j}{2^{A_j}}x_0
+
\sum_{i=0}^{j-1}
\frac{3^{j-1-i}}{2^{A_j-A_i}},
\] and the substitution \(r=j-i\) yields (1).

In particular, \[
x_j>\frac j3.
\]

A stronger bound comes from packing the previous orbit states below \(x_j\).

Since \(u_r\geq1\), the affine formula for the suffix from \(x_{j-r}\) to \(x_j\) gives \[
x_{j-r}<x_j.
\] Assume the segment contains no repetition, as it must in a divergent orbit. Put \(y=x_j\). Then \(x_0,\ldots,x_{j-1}\) are distinct odd numbers below \(y\), and each satisfies \(U(x_i)<y\).

Among integers \(x<y\):

\begin{itemize}
\item the states with \(a(x)=1\) and \(U(x)<y\) number at most \[
  \frac y6+O(1);
  \]
\item the states with \(a(x)=2\) number at most \[
  \frac y8+O(1);
  \]
\item the states with \(a(x)=3\) number at most \[
  \frac y{16}+O(1);
  \]
\item generally, for \(a\geq2\), the capacity is \(y/2^{a+1}+O(1)\).
\end{itemize}

Also \(q_j\geq1\), so \[
A_j\leq\alpha j.
\]

To maximize the number of distinct states under this valuation budget, one takes all available \(a=1\) and \(a=2\) states and then part of the \(a=3\) class. The resulting fractional-knapsack calculation gives \[
j
\leq
\frac{11}{24(3-\alpha)}\,y
+
O(\log y).
\] Equivalently, \[
\boxed{
x_j
\geq
\frac{24(3-\log_2 3)}{11}\,j
-
O(\log x_j).
}
\tag{2}
\] The constant is \[
\frac{24(3-\log_2 3)}{11}
=
3.0873545439\ldots.
\]

This is a real unconditional pointwise packing bound. But it is only linear in \(j\), corresponding to a ternary output length of order \(\log j\). A full coding proof would need a much stronger active-distance theorem, roughly linear growth of the output digit length in the input length.


\bigskip



\subsubsection*{9. Why excluding coefficient-growth words would still not prove Collatz}


There are two separate problems.


\paragraph*{9.1 Additive carry can rescue a multiplicatively contracting prefix}


The exact iterate is \[
T^k(n)=q_k n+r_k,
\qquad r_k>0.
\] It is possible that \[
q_k<1
\quad\text{but}\quad
T^k(n)\geq n
\] because \(r_k\) compensates for the coefficient loss.

Therefore, excluding words with \(q_k\geq1\) for every prefix does not exclude a minimal counterexample.


\paragraph*{9.2 Nontrivial cycles lie outside the sustained-growth cone}


For an accelerated cycle of length \(r\) with valuation sum \(A\), \[
x_r=x_0
\] implies \[
(2^A-3^r)x_0=C(a)>0.
\] Consequently, \[
2^A>3^r,
\qquad
A>\alpha r.
\]

Thus a positive cycle has multiplicative coefficient \[
\frac{3^r}{2^A}<1
\] over a full period. It lies outside the proposed low-valuation sustained-growth region.

So even a proof of (CLD) would not be a full proof of Collatz.

For cycles, coding or additive-combinatorial formulations become zero-avoidance problems modulo \(2^A-3^r\). Tao's Littlewood\textendash{}Offord analysis explains why ordinary high-dimensional covering or equidistribution estimates do not exclude the single zero syndrome once the modulus is large; a genuinely arithmetic separation theorem is needed. [\hypertarget{cit:R42:2}{}\hyperlink{ref:R42:2}{2}]


\bigskip



\subsubsection*{10. The coding theorem that really would prove Collatz}


For \(L,K\geq1\), define the actual-barrier list \[
\mathcal B_{L,K}
=
\left\{
1<n<2^L:
T^t(n)\geq n
\text{ for every }1\leq t\leq K
\right\}.
\]

Under the parity encoding, this is a subset of the zero-syndrome terminated code: \[
Q_K(\mathcal B_{L,K})\subseteq\mathcal C_{L,K}.
\]

The following statement would be sufficient.


\paragraph*{Actual-barrier list-decoding theorem}


There exist absolute constants \(C,\eta>0\) such that \[
\boxed{
|\mathcal B_{L,K}|
\leq
C\,2^L e^{-\eta K}
}
\tag{ABLD}
\] for all \(L,K\).


\paragraph*{Proof that (ABLD) implies Collatz}


Suppose Collatz is false, and choose the least positive integer \(n>1\) whose orbit does not reach \(1\).

No iterate can be smaller than \(n\): if \(T^t(n)=m<n\), then by minimality the orbit of \(m\) reaches \(1\), and hence so does the orbit of \(n\).

Choose \(L\) with \(n<2^L\). Then \[
n\in\mathcal B_{L,K}
\qquad\text{for every }K.
\] But (ABLD) gives \[
|\mathcal B_{L,K}|
\leq C2^Le^{-\eta K}<1
\] for sufficiently large \(K\), a contradiction. \ensuremath{\blacksquare}

This theorem automatically excludes both divergent orbits and nontrivial cycles. It is the correct coding-theoretic completion target.

However, I have not proved (ABLD). It is substantially stronger than an entropy estimate: it requires the zero-tail syndrome of one fixed exceptional integer interval to behave pseudorandomly against the actual carry-rescued barrier language. A March 2026 structural Collatz reduction describes an analogous remaining obstruction as the passage from exact map-level residue balance to deterministic balance along a single orbit. [\hypertarget{cit:R42:3}{}\hyperlink{ref:R42:3}{3}]


\bigskip



\subsection*{Final assessment}


The approach in its original form cannot yield a proof:

\[
\boxed{
\text{finite realizable parity/valuation code}
=
\text{entire finite word space}.
}
\]

Moreover, ``distance greater than the covering radius'' is impossible by definition, and ordinary minimum distance degenerates because Collatz trajectories can merge.

The repaired formulation is:

\[
\boxed{
\text{Collatz as zero-syndrome list decoding for a terminated,
nonlinear mixed }2\text{-adic}/3\text{-adic code}.
}
\]

The analysis gives three exact advances:

\[
\delta
=
1-H_2(\log_3 2)
=
0.0500444728\ldots
\] is the common binary and ternary rate deficit;

\[
\beta_*
=
1.6417941208\ldots
>
\log_2 3
\] shows that ordinary weighted covering lies beyond the growth threshold; and the record-time packing estimate \[
x_j
\geq
3.0873545439\ldots\,j-O(\log x_j)
\] is an unconditional pointwise active-distance bound.

But these estimates do not establish the needed actual-barrier list-decoding theorem (ABLD). Presenting them as a completed Collatz proof would conceal precisely the pointwise arithmetic-mixing statement that remains unproved.
\subsection*{References}
\begin{enumerate}[leftmargin=2.2em]
\item \hypertarget{ref:R42:1}{}\hyperlink{cit:R42:1}{\texttt{turn207311view0}}
\item \hypertarget{ref:R42:2}{}\hyperlink{cit:R42:2}{\texttt{turn400590view0}}
\item \hypertarget{ref:R42:3}{}\hyperlink{cit:R42:3}{\texttt{turn207311view1}}
\end{enumerate}

\clearpage
\section{Round 43: approach 38 --- List decoding of valuation profiles}
\subsection*{Prompt}
\begin{quote}
Complete a full proof of the Collatz conjecture using the following approach:

\begin{enumerate}
\item List decoding of valuation profiles
\end{enumerate}

\textbf{Speculative \textperiodcentered{} B}

Given a coarse growth profile, bound how many residue classes modulo 2\textasciicircum{}k can realize it while also satisfying specified endpoint conditions modulo 3\textasciicircum{}r. Show that the list size eventually becomes zero \textemdash{} or find a uniquely surviving nested sequence that could encode a counterexample.

Take as long as you need.
\end{quote}
\subsection*{Response}

\subsection*{Verdict}


I do \textbf{not} have a valid full proof of the Collatz conjecture from this route. The list-decoding formulation can be made completely rigorous, and it yields a useful new pruning certificate, but the hoped-for extinction theorem does not follow. In fact:

\begin{enumerate}
\item The endpoint map is not injective, even after fixing the length, total valuation, and strict growth barrier.
\item Vanishing density or bounded list size does not imply an empty list.
\item After optimal decoding, the remaining assertion is precisely the unresolved synchronization theorem from the earlier affine-compression framework: an infinite retained ray whose \(2\)-adic representatives stabilize at a positive integer must eventually admit a smaller positive compatible ancestor. [\hypertarget{cit:R43:1}{}\hyperlink{ref:R43:1}{1}]
\end{enumerate}

What follows is the strongest rigorous version of approach 38 that I could obtain.


\bigskip



\subsubsection*{1. The exact valuation code}


Let

\[
U(x)=\frac{3x+1}{2^{\nu_2(3x+1)}}
\]

on positive odd integers. For a word

\[
w=(a_0,\ldots,a_{j-1}),\qquad a_i\geq 1,
\]

write

\[
A_0=0,\qquad A_t=\sum_{i<t}a_i,\qquad A=A_j,
\]

and

\[
C_w=\sum_{i=0}^{j-1}3^{j-1-i}2^{A_i}.
\]

Then every realization satisfies

\[
2^A U^j(n)=3^j n+C_w.
\tag{1}
\]

The exact seed cylinder and endpoint syndrome are

\[
n\equiv \rho_2(w)
:=3^{-j}(2^A-C_w)
\pmod {2^{A+1}},
\tag{2}
\]

and

\[
U^j(n)\equiv \rho_3(w)
:=C_w2^{-A}
\pmod {3^j}.
\tag{3}
\]

Every word has infinitely many positive realizations, and every positive odd endpoint in the class (3) has a positive inverse ancestor along \(w\). [\hypertarget{cit:R43:2}{}\hyperlink{ref:R43:2}{2}]

Thus a valuation word is naturally a codeword with:

\[
\text{message }w,\qquad
\text{\(2\)-adic coordinate }\rho_2(w),\qquad
\text{\(3\)-adic syndrome }\rho_3(w).
\]


\bigskip



\subsubsection*{2. Exact list-size formula}


Let \(\mathcal P\) be any collection of length-\(j\) words, representing a coarse growth profile. For example, \(\mathcal P\) might impose intervals

\[
L_t\leq A_t\leq U_t.
\]

Fix an endpoint condition

\[
\rho_3(w)\equiv c\pmod{3^r},
\qquad r\leq j,
\]

and take

\[
K\geq 1+\max_{w\in\mathcal P}A(w).
\]

Let \(\mathscr L_K(\mathcal P;c,r)\) be the set of residue classes modulo \(2^K\) whose first \(j\) valuations form a word in \(\mathcal P\) and whose endpoint satisfies the specified \(3^r\)-condition.


\paragraph*{Exact list formula}


\[
\boxed{
\left|\mathscr L_K(\mathcal P;c,r)\right|
=
\sum_{\substack{w\in\mathcal P\\
\rho_3(w)\equiv c\; (3^r)}}
2^{K-A(w)-1}.
}
\tag{4}
\]


\paragraph*{Proof}


A word of total valuation \(A\) determines one residue class modulo \(2^{A+1}\). It therefore has exactly

\[
2^{K-A-1}
\]

lifts modulo \(2^K\). Distinct valuation words of the same length give disjoint seed cylinders, since the first \(j\) valuations of a fixed integer are deterministic. Summing over the words with the required endpoint syndrome proves (4). \ensuremath{\blacksquare}

Dividing by the number \(2^{K-1}\) of odd classes modulo \(2^K\) gives

\[
\boxed{
\frac{|\mathscr L_K(\mathcal P;c,r)|}{2^{K-1}}
=
\sum_{\substack{w\in\mathcal P\\
\rho_3(w)\equiv c\; (3^r)}}
2^{-A(w)}.
}
\tag{5}
\]

Thus normalized list size is exactly the Haar measure of the union of the corresponding \(2\)-adic cylinders.


\bigskip



\subsubsection*{3. A rigorous lift-width list-decoding theorem}


For fixed length \(j\) and fixed total valuation \(A\), define the real lift

\[
\Lambda(w):=\frac{C_w}{3^j}
=\sum_{i=0}^{j-1}\frac{2^{A_i}}{3^{i+1}}.
\tag{6}
\]

Suppose that all words in a profile \(\mathcal P\) have the same total \(A\).

For a full endpoint syndrome \(c\bmod 3^j\), compatibility is

\[
C_w\equiv 2^A c\pmod{3^j}.
\tag{7}
\]

Hence the possible constants are

\[
C_w=2^Ac+q3^j,
\qquad q\in\mathbb Z.
\tag{8}
\]

Distinct words of fixed \((j,A)\) have distinct constants \(C_w\): otherwise they would have the same seed cylinder in (2), and a positive integer in that cylinder would realize two different valuation words.


\paragraph*{Lift-width theorem}


Suppose

\[
\Lambda(w)\in[I_-,I_+]
\qquad\text{for every }w\in\mathcal P.
\]

Then every full endpoint syndrome has list size at most

\[
\boxed{
N_{\mathcal P}(c)
\leq
\lfloor I_+-I_-\rfloor+1.
}
\tag{9}
\]

More generally, for a condition only modulo \(3^r\),

\[
\boxed{
N_{\mathcal P}(c;r)
\leq
\left\lfloor
3^{j-r}(I_+-I_-)
\right\rfloor+1.
}
\tag{10}
\]


\paragraph*{Proof}


For the full syndrome, all compatible \(C_w\) lie in one arithmetic progression of step \(3^j\). Since

\[
3^jI_-\leq C_w\leq3^jI_+,
\]

there are at most

\[
\left\lfloor
\frac{3^j(I_+-I_-)}{3^j}
\right\rfloor+1
\]

possible values. The modulus-\(3^r\) version has progression step \(3^r\). \ensuremath{\blacksquare}

If the profile gives componentwise bounds

\[
L_i\leq A_i\leq U_i,
\]

then a completely explicit width is

\[
W(\mathcal P)
=
\sum_{i=0}^{j-1}
\frac{2^{U_i}-2^{L_i}}{3^{i+1}},
\tag{11}
\]

and therefore

\[
N_{\mathcal P}(c)
\leq \lfloor W(\mathcal P)\rfloor+1.
\tag{12}
\]

In particular,

\[
W(\mathcal P)<1
\quad\Longrightarrow\quad
\text{unique decoding for every occupied endpoint syndrome.}
\tag{13}
\]

This is the cleanest rigorous list-decoding theorem produced by the approach.


\bigskip



\subsubsection*{4. The lift index is exactly the seed}


There is an especially useful interpretation of (8).

Let \(E\) be an actual positive endpoint in the syndrome \(c\). Then

\[
2^A E=3^jR+C_w,
\]

where \(R\) is the corresponding seed representative. Thus

\[
C_w=2^AE-3^jR,
\]

or

\[
\boxed{
\Lambda(w)=\frac{2^AE}{3^j}-R.
}
\tag{14}
\]

Consequently, the integer lift \(q\) in (8) is precisely

\[
q=-R.
\]

So list decoding of \(C_w\bmod3^j\) is literally list decoding of possible seed representatives.

If the profile constrains \(\Lambda(w)\in[I_-,I_+]\), then the possible seeds lie in

\[
\boxed{
R\in
\left[
\frac{2^AE}{3^j}-I_+,\,
\frac{2^AE}{3^j}-I_-
\right]\cap\mathbb Z.
}
\tag{15}
\]

When \(I_+-I_-<1\), there is at most one possible seed.

This gives an exact decoder, but unique decoding is not yet extinction: the unique seed might be the hypothetical counterexample itself.


\bigskip



\subsubsection*{5. The entropy calculation}


Ignore prefix restrictions temporarily and fix \(j\) and \(A\). The number of valuation words is

\[
\binom{A-1}{j-1}.
\tag{16}
\]

There are

\[
\varphi(3^j)=2\cdot3^{j-1}
\]

possible endpoint-unit syndromes. Therefore the average full-syndrome list size is

\[
\frac{\binom{A-1}{j-1}}{2\cdot3^{j-1}}.
\tag{17}
\]

Put

\[
\alpha=\log_2 3=1.5849625007\ldots
\]

and suppose \(A=\alpha j+O(1)\). Stirling's formula gives

\[
\binom{A-1}{j-1}
=
\exp\!\left(jf(\alpha)+O(\log j)\right),
\]

where

\[
f(\beta)
=
\beta\log\beta-(\beta-1)\log(\beta-1).
\]

Numerically,

\[
e^{f(\alpha)}
=
2.8395137305\ldots<3.
\]

Thus the average list size behaves like

\[
\frac{1}{\sqrt j}
\left(0.9465045768\ldots\right)^j.
\tag{18}
\]

This is a genuine coding-rate gap of

\[
\log 3-f(\alpha)
=
0.0549794728\ldots
\]

nats, or about \(0.07932\) bits per accelerated step.

But the interpretation is crucial:


\begin{quote}
The rate gap says that most \(3\)-adic syndromes are empty. It does not say that the syndrome selected by an actual orbit is empty.

\end{quote}


Every valuation word automatically occupies its own syndrome. An exceptional orbit can therefore travel through a nested sequence of atypical occupied syndromes even though the average fibre tends to zero.

This is the same qualitative gap between probabilistic and pointwise Collatz results. Tao's proof controls a skew random walk on \(3\)-adic cyclic groups sufficiently well to obtain an almost-all statement, but it does not provide uniform extinction for one adversarial nested cylinder. [\hypertarget{cit:R43:3}{}\hyperlink{ref:R43:3}{3}]

A July 2026 preprint studying essentially the same start/end code likewise obtains necessary residue-rate conditions and numerical diagnostics, while explicitly distinguishing these from a proof. [\hypertarget{cit:R43:4}{}\hyperlink{ref:R43:4}{4}]


\bigskip



\subsubsection*{6. Fixed-total injectivity is false}


A natural hope was:


\begin{quote}
Among strict-barrier words of fixed length and fixed total valuation, the endpoint syndrome determines the word.

\end{quote}


This is false.

Consider

\[
\begin{aligned}
w_1={}&(1,1,1,1,1,1,1,1,1,2,1,1,3,3,1,1,2,1,1,5,1,3),\\
w_2={}&(1,2,1,2,1,2,2,1,2,1,1,3,1,2,1,2,1,2,2,1,2,1).
\end{aligned}
\tag{19}
\]

Both have

\[
|w_1|=|w_2|=22,
\qquad
A(w_1)=A(w_2)=34.
\]

Their cumulative valuations are respectively

\[
\begin{aligned}
&(1,2,3,4,5,6,7,8,9,11,12,13,16,19,20,21,23,24,25,30,31,34),\\
&(1,3,4,6,7,9,11,12,14,15,16,19,20,22,23,25,26,28,30,31,33,34).
\end{aligned}
\]

Both are componentwise at most

\[
\bigl(\lfloor t\log_2 3\rfloor\bigr)_{t=1}^{22}
=
(1,3,4,6,7,9,11,12,14,15,17,19,20,22,23,25,26,28,30,31,33,34),
\]

and hence satisfy

\[
2^{A_t}<3^t
\qquad(1\leq t\leq22).
\]

Their constants are

\[
C_{w_1}=41\,256\,549\,401,
\]

and

\[
C_{w_2}=166\,780\,787\,837.
\]

The decisive identity is

\[
\boxed{
C_{w_2}-C_{w_1}
=
125\,524\,238\,436
=
4\cdot3^{22}.
}
\tag{20}
\]

Since their total valuations agree,

\[
\rho_3(w_1)=\rho_3(w_2)
=
9\,770\,112\,830
\pmod{3^{22}}.
\tag{21}
\]

Their seed representatives modulo \(2^{35}\) are

\[
R(w_1)=22\,528\,613\,375,
\]

and

\[
R(w_2)=22\,528\,613\,371.
\tag{22}
\]

The common canonical endpoint is

\[
E=41\,151\,172\,439.
\]

For every common compatible endpoint \(X\),

\[
n_i=\frac{2^{34}X-C_{w_i}}{3^{22}},
\]

so (20) gives

\[
\boxed{n_2=n_1-4.}
\tag{23}
\]

Thus \(w_1\) is unconditionally dominated by the equal-slope, larger-intercept word \(w_2\).

This is not a counterexample to Collatz; it is an exact descending diamond and therefore a useful pruning certificate. It does, however, disprove fixed-total unique decoding.


\paragraph*{Persistent collision family}


For every common suffix \(s\) of length \(t\),

\[
C_{w_2s}-C_{w_1s}
=
3^t(C_{w_2}-C_{w_1})
=
4\cdot3^{22+t}.
\tag{24}
\]

Hence \(w_1s\) and \(w_2s\) continue to have identical endpoint syndromes, identical total valuations, and seed representatives differing by \(4\).

In particular, appending \(1^t\) preserves the strict barrier because

\[
34+t<(22+t)\log_2 3
\qquad(t\geq0).
\]

Therefore there are fixed-total strict-barrier collisions at every depth at least \(22\).

The exact arithmetic is reproducible here:

[Exact collision verifier](sandbox:/mnt/data/list\_decoding\_collision\_certificate.py)


\bigskip



\subsubsection*{7. Bellman decoding repairs the collision\textemdash{}but does not prove extinction}


For a fixed endpoint, an inverse word \(w\) gives

\[
n_w(X)=\frac{2^{A(w)}X-C_w}{3^{|w|}}.
\]

The optimal inverse decoder therefore:

\begin{enumerate}
\item minimizes \(A(w)-|w|\log_2 3\), equivalently maximizes the inverse slope;
\item among equal slopes, maximizes \(C_w\).
\end{enumerate}

This is exactly the Bellman order previously derived from least-counterexample minimality. [\hypertarget{cit:R43:5}{}\hyperlink{ref:R43:5}{5}]

The depth-\(22\) pair is handled correctly: \(w_2\) wins the equal-slope intercept tie and \(w_1\) is discarded.

But the optimally decoded tree still has infinite rays. A particularly clean one is

\[
(1,1,2)^\infty.
\]

One block has

\[
U^3(x)=\frac{27x+19}{16},
\]

so its \(2\)-adic fixed point is

\[
x=-\frac{19}{11}.
\]

Every finite prefix of this ray is Bellman-retained, even though the limiting \(2\)-adic point is not a positive ordinary integer. [\hypertarget{cit:R43:6}{}\hyperlink{ref:R43:6}{6}]

Thus:

\[
\text{nonempty nested list}
\;\not\Longrightarrow\;
\text{positive counterexample}.
\]

The correct ordinary-integer test is that the least seed representatives \(R_j\) must eventually stabilize. [\hypertarget{cit:R43:7}{}\hyperlink{ref:R43:7}{7}]


\bigskip



\subsubsection*{8. Density tending to zero is not enough}


The recurrent subsystem already isolated in the compression framework gives a particularly sharp warning. Its dynamics are

\[
m_{t+1}=\frac{81m_t+\varepsilon_t}{64},
\qquad \varepsilon_t\in\{0,1\},
\]

with admissible residues

\[
m_t\equiv0\pmod{64}\quad(\varepsilon_t=0),
\qquad
m_t\equiv15\pmod{64}\quad(\varepsilon_t=1).
\]

At depth \(r\), the \(2^r\) binary words give \(2^r\) distinct residue classes modulo \(64^r\). Consequently, their density is exactly

\[
32^{-r}\longrightarrow0,
\]

but their cardinality is never zero. Termination is equivalent to the least positive representative \(M_r\) tending to infinity\textemdash{}not to the density tending to zero. [\hypertarget{cit:R43:8}{}\hyperlink{ref:R43:8}{8}]

Eventually periodic survivor itineraries are nonpositive \(2\)-adic rationals, but an aperiodic survivor is not excluded; it must merely have very high factor complexity. [\hypertarget{cit:R43:9}{}\hyperlink{ref:R43:9}{9}]

This is almost exactly the list-decoding obstruction:

\[
\text{exponentially sparse lists}
\quad\text{can still contain an infinite nested ray.}
\]


\bigskip



\subsubsection*{9. The two-place small-rectangle formulation}


For a word \(w\), let \(R(w)\) be its least positive seed representative and define

\[
E(w)=\frac{3^jR(w)+C_w}{2^A}.
\tag{25}
\]

Then

\[
0<E(w)<2\cdot3^j.
\]

Indeed, changing \(R\) by \(2^{A+1}\) changes the endpoint by \(2\cdot3^j\), and the positivity of all backward iterates makes \(E(w)\) the canonical positive odd endpoint representative.

Now suppose a fixed positive integer \(n\) has valuation prefixes \(w_j\). For all sufficiently large \(j\),

\[
R(w_j)=n.
\tag{26}
\]

Also,

\[
U(x)+1
\leq \frac32(x+1),
\]

so

\[
U^j(n)+1
\leq
\left(\frac32\right)^j(n+1).
\tag{27}
\]

Since the right side is eventually below \(2\cdot3^j\), the canonical endpoint representative eventually equals the actual endpoint:

\[
E(w_j)=U^j(n).
\tag{28}
\]

Therefore a positive counterexample would give an infinite sequence of code points satisfying

\[
\boxed{
R(w_j)=n,\qquad
E(w_j)\leq(n+1)\left(\frac32\right)^j.
}
\tag{29}
\]

On a critical block \(A\sim j\log_2 3\), both ambient coordinate ranges have size comparable to \(3^j\), whereas (29) lies in an exponentially tiny anchored rectangle.

A random-code calculation predicts approximately

\[
\binom{A-1}{j-1}
\frac{n}{2^{A+1}}
\frac{(n+1)(3/2)^j}{2\cdot3^j}
\asymp
n(n+1)
\left(\frac{2.8395137\ldots}{6}\right)^j,
\]

where

\[
\frac{2.8395137\ldots}{6}
=
0.4732522\ldots<1.
\tag{30}
\]

So the random model strongly predicts eventual emptiness.

But (30) is only a first-moment heuristic. The actual orbit selects one exceptional nested syndrome sequence; no known discrepancy or spectral estimate is strong enough to exclude a single code point in a rectangle whose expected occupancy is already below \(1\).

Finite computation through \(2^{71}\) can discharge any explicit bounded range, but it does not supply this all-depth worst-case estimate. [\hypertarget{cit:R43:10}{}\hyperlink{ref:R43:10}{10}]


\bigskip



\subsubsection*{10. The exact theorem still needed}


The list-decoding route would prove Collatz if one could establish the following.


\paragraph*{Positive stationary list-extinction theorem}


Let \(a=(a_0,a_1,\ldots)\) be an infinite valuation ray, and let \(R_j\) and \(E_j\) be its canonical seed and endpoint representatives. Suppose that

\[
R_j=m>1
\]

for all sufficiently large \(j\). Then at some finite depth \(j\), there must exist a competing word \(b\), of some length \(r\) and total valuation \(B\), such that

\[
E_j\equiv C_b2^{-B}\pmod{3^r},
\]

and

\[
\boxed{
0<
\frac{2^B E_j-C_b}{3^r}
<m.
}
\tag{31}
\]


\paragraph*{Why this proves Collatz}


Assume a counterexample exists and let \(m\) be the least positive odd counterexample. Its valuation representatives eventually stabilize at \(m\). Equation (31) would produce a positive odd integer \(z<m\) whose orbit reaches the later point \(E_j\) of the orbit of \(m\). But then \(z\) and \(m\) coalesce, contradicting minimality of \(m\). The coalescence principle is exact. [\hypertarget{cit:R43:11}{}\hyperlink{ref:R43:11}{11}]

In the earlier local-global notation, (31) is the assertion that some terminal fibre satisfies

\[
\Sigma(u,\tau,c)+h_{\max}(\tau,B)<0.
\]

That is exactly the synchronization theorem already identified as missing. [\hypertarget{cit:R43:12}{}\hyperlink{ref:R43:12}{12}]


\bigskip



\subsection*{Final assessment}


Approach 38 produces a mathematically useful framework:

\[
\boxed{
\text{profile width}
\;\Longrightarrow\;
\text{explicit list-size bound}
}
\]

and

\[
\boxed{
\text{endpoint lift index}
=
-\text{seed representative}.
}
\]

It also supplies an exact depth-\(22\) fixed-total collision and an infinite family of associated pruning certificates.

But it does not currently yield a full proof. The natural injectivity lemma is false, and optimal/Bellman decoding leaves infinite \(2\)-adic rays. The unresolved point is not counting average lists; it is excluding one \textbf{positive stationary nested list}. That requires a worst-case synchronization between:

\[
\text{the \(2\)-adic seed representative},
\qquad
\text{the \(3\)-adic endpoint syndrome},
\qquad
\text{and the archimedean ancestor inequality}.
\]

Neither the entropy gap, the lift-width theorem, Tao-type average transport, nor the presently known finite certificates proves that synchronization.
\subsection*{References}
\begin{enumerate}[leftmargin=2.2em]
\item \hypertarget{ref:R43:1}{}\hyperlink{cit:R43:1}{\texttt{turn4file1 L301-L315}}
\item \hypertarget{ref:R43:2}{}\hyperlink{cit:R43:2}{\texttt{turn4file0 L12-L94}}
\item \hypertarget{ref:R43:3}{}\hyperlink{cit:R43:3}{\texttt{turn215668view1}}
\item \hypertarget{ref:R43:4}{}\hyperlink{cit:R43:4}{\texttt{turn215668view0}}
\item \hypertarget{ref:R43:5}{}\hyperlink{cit:R43:5}{\texttt{turn4file0 L117-L136}}
\item \hypertarget{ref:R43:6}{}\hyperlink{cit:R43:6}{\texttt{turn4file0 L226-L252}}
\item \hypertarget{ref:R43:7}{}\hyperlink{cit:R43:7}{\texttt{turn4file0 L208-L235}}
\item \hypertarget{ref:R43:8}{}\hyperlink{cit:R43:8}{\texttt{turn5file0 L18-L78}}
\item \hypertarget{ref:R43:9}{}\hyperlink{cit:R43:9}{\texttt{turn5file0 L80-L99}}
\item \hypertarget{ref:R43:10}{}\hyperlink{cit:R43:10}{\texttt{turn946066search5}}
\item \hypertarget{ref:R43:11}{}\hyperlink{cit:R43:11}{\texttt{turn4file0 L96-L136}}
\item \hypertarget{ref:R43:12}{}\hyperlink{cit:R43:12}{\texttt{turn4file1 L301-L311}}
\end{enumerate}

\clearpage
\section{Round 44: approach 39 --- An information-capacity bound}
\subsection*{Prompt}
\begin{quote}
Complete a full proof of the Collatz conjecture using the following approach:

\begin{enumerate}
\item An information-capacity bound
\end{enumerate}

\textbf{Speculative \textperiodcentered{} P}

A seed supplies k low binary bits to determine a length-k parity prefix. Simultaneously satisfying increasingly precise ternary endpoint conditions consumes additional information. Formalize this as a communication or data-processing inequality and prove the combined demands exceed the information available from one seed.

Take as long as you need.
\end{quote}
\subsection*{Response}

\subsection*{Result}


A valid full proof does \textbf{not} emerge from the proposed information-capacity argument. The obstruction is exact: the binary itinerary and the ternary endpoint data are not independent messages. The latter is a deterministic function of the former, so the proposed additive information budget

\[
\text{binary itinerary cost}+\text{ternary endpoint cost}
\]

double-counts the same information.

What the approach does yield is:

\begin{enumerate}
\item an exact no-go theorem for a Shannon/data-processing proof;
\item an exponential sparsity theorem for growth-compatible valuation words;
\item a precise zero-error inequality which, if established, would imply Collatz;
\item a proof that counting, entropy, or measure alone cannot establish that inequality.
\end{enumerate}

A July 2026 preprint by Oliver Kramer independently formalizes almost exactly the same ``small \(2\)-adic start representative plus small \(3\)-adic endpoint representative'' diagnostic. It proves that both normalized residue rates must vanish for any code generated by a fixed positive integer, while explicitly noting that this is a diagnostic rather than a verification method. [\hypertarget{cit:R44:1}{}\hyperlink{ref:R44:1}{1}, \hypertarget{cit:R44:2}{}\hyperlink{ref:R44:2}{2}]


\subsubsection*{1. The exact finite-word identities}


Use the shortcut Collatz map

\[
T(n)=
\begin{cases}
n/2,&n\equiv0\pmod 2,\\[2mm]
(3n+1)/2,&n\equiv1\pmod 2.
\end{cases}
\]

For a length-\(k\) parity word

\[
e=(e_0,\ldots ,e_{k-1})\in\{0,1\}^k,
\qquad
s_i=\sum_{h<i}e_h,\qquad q=s_k,
\]

iteration gives

\[
2^kT^k(n)=3^q n+C(e),
\tag{1}
\]

where

\[
C(e)=\sum_{i=0}^{k-1}e_i\,2^i3^{q-s_{i+1}}.
\tag{2}
\]

Every parity word occurs for exactly one residue class modulo \(2^k\); equivalently, the map

\[
n\bmod 2^k\longmapsto
\bigl(n,T(n),\ldots,T^{k-1}(n)\bigr)\bmod2
\]

is a bijection. This is the Terras\textendash{}Everett parity-vector theorem. [\hypertarget{cit:R44:3}{}\hyperlink{ref:R44:3}{3}]

Now let \(r\le q\). Reducing (1) modulo \(3^r\) gives

\[
T^k(n)\equiv 2^{-k}C(e)\pmod{3^r}.
\tag{3}
\]

Thus the endpoint modulo \(3^r\) is determined completely by the parity word.

The same statement is particularly clean for the accelerated odd map

\[
U(x)=\frac{3x+1}{2^{v_2(3x+1)}}.
\]

Let

\[
a_i=v_2(3x_i+1),\qquad
A_i=\sum_{h<i}a_h,\qquad A=A_j.
\]

Then

\[
2^A x_j=3^j n+C_w,
\tag{4}
\]

where

\[
C_w=\sum_{i=0}^{j-1}3^{j-1-i}2^{A_i}.
\tag{5}
\]

Consequently, for every \(r\le j\),

\[
x_j\equiv 2^{-A}C_w\pmod{3^r}.
\tag{6}
\]

Equation (6) is the exact information-theoretic obstruction.


\subsubsection*{2. The data-processing inequality goes in the wrong direction}


Take a seed \(N\) uniformly distributed modulo \(2^k\). Let

\[
E=E_k(N)
\]

be its parity word. By the parity-vector theorem,

\[
H(E)=H(N)=k
\]

when entropy is measured in bits.

Let

\[
Y=T^k(N)\bmod 3^r,
\]

on a family of words having at least \(r\) odd steps. Equation (3) says that \(Y=f(E)\) for a deterministic function \(f\). Hence

\[
H(Y\mid E)=0
\]

and therefore

\[
H(E,Y)=H(E)+H(Y\mid E)=k.
\tag{7}
\]

It is not

\[
k+r\log_2 3.
\]

Equivalently, the Markov chain is

\[
N\longrightarrow E\longrightarrow Y,
\]

and

\[
I(N;E,Y)
 =I(N;E)+I(N;Y\mid E)
 =I(N;E).
\tag{8}
\]

The conditional term is exactly zero.

Suppose instead that only a coarse growth profile \(G=g(E)\) is recorded. Then the endpoint may distinguish words inside a profile class, but still

\[
H(Y\mid G)\le H(E\mid G)
\tag{9}
\]

and

\[
H(G,Y)\le H(E)\le k.
\tag{10}
\]

So ternary endpoint data can consume only the uncertainty left inside the parity-profile list. It cannot create a second independent information demand.

This disproves the central additive-budget premise of Approach 39.


\subsubsection*{3. What the capacity calculation really proves}


Set

\[
\alpha=\log_2 3=1.5849625007\ldots
\]

and consider accelerated valuation words

\[
w=(a_0,\ldots ,a_{j-1}),\qquad a_i\ge1,
\]

satisfying the noncontracting prefix inequalities

\[
A_t=\sum_{i<t}a_i\le \alpha t
\qquad(1\le t\le j).
\tag{11}
\]

Let \(\mathcal B_j\) denote this barrier language.

Ignoring the prefix conditions and retaining only \(A_j\le\lfloor\alpha j\rfloor\) gives

\[
|\mathcal B_j|
 \le
 \sum_{A=j}^{\lfloor\alpha j\rfloor}
 \binom{A-1}{j-1}
 =
 \binom{\lfloor\alpha j\rfloor}{j}.
\tag{12}
\]

Stirling's formula therefore gives

\[
\log_2|\mathcal B_j|
 \le
 \beta j+O(\log j),
\]

where

\[
\beta
 =\alpha H_2(1/\alpha)
 =1.5056438879\ldots
\tag{13}
\]

and \(H_2\) is binary entropy. Since there are \(3^j=2^{\alpha j}\) possible endpoint residues modulo \(3^j\), the apparent capacity gap is

\[
\gamma=\alpha-\beta
       =0.0793186128\ldots
\tag{14}
\]

bits per accelerated step.

For \(b\in\mathbb Z/3^j\mathbb Z\), define the endpoint list

\[
L_j(b)=
\left\{
w\in\mathcal B_j:
2^{-A_j}C_w\equiv b\pmod{3^j}
\right\}.
\]

Then

\[
\sum_{b\bmod 3^j}|L_j(b)|=|\mathcal B_j|,
\]

so

\[
\frac1{3^j}\sum_{b\bmod3^j}|L_j(b)|
\le
2^{-\gamma j+O(\log j)}.
\tag{15}
\]

This is a rigorous exponential capacity theorem.

But its interpretation is:


\begin{quote}
Almost every ternary residue has no growth-compatible preimage.

\end{quote}


It does \textbf{not} say:


\begin{quote}
Every residue selected by a hypothetical orbit has no preimage.

\end{quote}


Indeed, every word \(w\in\mathcal B_j\) selects its own residue

\[
b=2^{-A_j}C_w\bmod3^j,
\]

and for that adaptively selected residue,

\[
|L_j(b)|\ge1.
\]

The average list size can be exponentially below one while a nested sequence of singleton lists survives forever. Average list decoding is therefore insufficient; a \textbf{zero-error, worst-case} statement is required.


\subsubsection*{4. Why measure-zero and entropy-zero results cannot finish Collatz}


The capacity estimate shows that exceptional words form a set of zero natural symbolic measure. That is entirely compatible with containing one or countably many integer-generated paths.

The same obstruction is familiar in the \(2\)-adic formulation. The \(3x+1\) and \(5x+1\) maps have the same full-shift symbolic dynamics and are topologically and measure-theoretically conjugate on \(\mathbb Z_2\). But the ordinary integers form a countable, Haar-measure-zero subset, so those ergodic facts give no conclusion about dynamics on the integers. [\hypertarget{cit:R44:4}{}\hyperlink{ref:R44:4}{4}]

Likewise, Tao's theorem that almost all Collatz orbits attain almost bounded values is an almost-everywhere result in logarithmic density; it does not eliminate a pointwise exceptional seed. [\hypertarget{cit:R44:5}{}\hyperlink{ref:R44:5}{5}]

There is an even simpler information-theoretic explanation. For a hypothetical fixed counterexample \(n\), the random variable ``seed'' is constant. Consequently,

\[
H(n)=H(E_k(n))=H(T^k(n)\bmod3^r)=0.
\]

One cannot use the entropy of a uniform ensemble of \(2^k\) seeds to rule out one deterministic member of that ensemble.


\subsubsection*{5. Drift alone misses the additive carry}


For the accelerated map define

\[
q_j=\frac{3^j}{2^{A_j}},
\qquad
h_j=\frac{x_j}{q_j}.
\]

A direct calculation gives

\[
h_{j+1}
 =h_j+\frac{1}{3q_j}.
\tag{16}
\]

Thus \(h_j\) is increasing. Even when \(q_j<1\), the accumulated carry \(h_j\) can keep

\[
x_j=q_jh_j
\]

above its starting value.

So a divergent or infinite-stopping orbit need not give an infinite word satisfying the barrier inequalities (11). It may cross repeatedly into the nominally contracting region while additive carries prevent descent.

This is not merely a logical possibility. Rozier and Terracol call a finite segment paradoxical when its coefficient is below one but its endpoint is at least its starting value. They prove that an integer of infinite stopping time would generate infinitely many such segments, and consequently that finiteness of all paradoxical segments would imply the Collatz conjecture. [\hypertarget{cit:R44:6}{}\hyperlink{ref:R44:6}{6}, \hypertarget{cit:R44:7}{}\hyperlink{ref:R44:7}{7}]

Therefore a proof that eliminates only sustained noncontracting coefficient words would not prove Collatz. It must also control the carry term \(C(e)\).


\subsubsection*{6. A positive countermodel to rate-only reasoning}


Consider the shortcut \(5x+1\) map. It has the positive cycle

\[
13\longmapsto33\longmapsto83\longmapsto208
\longmapsto104\longmapsto52\longmapsto26
\longmapsto13.
\tag{17}
\]

Its parity word is

\[
1110000,
\]

with three odd steps in seven iterations. The multiplicative coefficient is

\[
\frac{5^3}{2^7}=\frac{125}{128}<1,
\]

yet the endpoint equals the starting value.

Equivalently, for the accelerated \(5x+1\) map the valuation block is

\[
(1,1,5),
\]

and repeating this block produces a fixed positive seed \(13\), a fixed endpoint \(13\), and arbitrarily precise \(5\)-adic endpoint congruences. The \(5x+1\) variant is known to have cycles beginning at \(13\) and \(17\), while its binary symbolic dynamics is the same full shift as for \(3x+1\). [\hypertarget{cit:R44:8}{}\hyperlink{ref:R44:8}{8}, \hyperlink{ref:R44:4}{4}]

Thus:


\begin{quote}
A negative asymptotic coding rate, exact integrality, a contracting multiplicative coefficient, and arbitrarily precise endpoint congruences can coexist with a uniquely surviving positive nested sequence.

\end{quote}


Any proof based only on information rates, alphabet sizes, data processing, or generic symbolic dynamics would incorrectly rule out (17). A successful theorem must use arithmetic special to the multiplier \(3\).


\subsubsection*{7. The exact zero-error reformulation}


The shortcut representation gives a precise replacement for the false Shannon inequality.

For a parity word \(e\) of length \(k\), let \(q=|e|_1\), and define \(C(e)\) by (2). Let

\[
\rho(e)\in\{1,\ldots ,2^k\}
\]

be the least positive representative satisfying

\[
3^q\rho(e)+C(e)\equiv0\pmod{2^k}.
\tag{18}
\]

Define

\[
y(e)=\frac{3^q\rho(e)+C(e)}{2^k}.
\tag{19}
\]

Then every positive integer producing the word \(e\) has the form

\[
n=\rho(e)+t2^k,
\qquad t\ge0,
\tag{20}
\]

and its endpoint is

\[
T^k(n)=y(e)+t3^q.
\tag{21}
\]

One can also prove inductively that

\[
1\le y(e)\le3^q.
\tag{22}
\]

Now assume that the multiplicative coefficient is contracting:

\[
3^q<2^k.
\]

Put

\[
D(e)=2^k-3^q>0.
\]

Subtracting (20) from (21) gives

\[
T^k(n)-n
 =
\bigl(y(e)-\rho(e)\bigr)-tD(e).
\tag{23}
\]

The right side decreases with \(t\). Consequently:

\[
\boxed{
\text{Some seed in the residue class is paradoxical}
\iff y(e)\ge\rho(e).
}
\tag{24}
\]

Using (19), this is equivalent to

\[
\boxed{
C(e)\ge
\bigl(2^k-3^q\bigr)\rho(e).
}
\tag{25}
\]

This is the exact zero-error bottleneck.

A theorem of the following form would complete the proposed route:


\begin{quote}
\textbf{Zero-error arithmetic inequality.}   Apart from finite segments belonging to the trivial \(1\leftrightarrow2\) cycle, there exists \(K\) such that every parity word \(e\) of length \(k\ge K\) satisfying \(3^{|e|_1}<2^k\) also satisfies \[
> C(e)<
> \bigl(2^k-3^{|e|_1}\bigr)\rho(e).
> \tag{ZE}
> \]

\end{quote}


If (ZE) held, there would be no sufficiently long paradoxical segment. For each of the finitely many shorter lengths there are finitely many words, and equation (23) permits only finitely many \(t\). Thus only finitely many paradoxical segments would exist, and Rozier\textendash{}Terracol's theorem would imply Collatz. [\hyperlink{ref:R44:7}{7}, \hyperlink{ref:R44:6}{6}]

But (ZE) is not an entropy inequality. It is a worst-case correlation inequality between:

\[
\text{the carry }C(e)
\quad\text{and}\quad
\text{the unusually small start representative }\rho(e).
\]

Counting tells us that such correlations are rare. It does not prove that they never occur.


\subsubsection*{8. A class that can be excluded completely}


There is one useful zero-error consequence.

Suppose an accelerated valuation block \(w\) of length \(j\), total valuation \(A\), repeats periodically. Its unique periodic realization satisfies

\[
2^A n=3^j n+C_w,
\]

so

\[
n=\frac{C_w}{2^A-3^j}.
\tag{26}
\]

If the periodic word is noncontracting,

\[
A\le j\log_2 3,
\]

then \(2^A<3^j\), since no positive powers of \(2\) and \(3\) are equal. Because \(C_w>0\), equation (26) gives

\[
n<0.
\]

Therefore:


\begin{quote}
Every periodic infinite noncontracting valuation code represents a negative rational \(2\)-adic state, not a positive integer.

\end{quote}


Hence any positive integer survivor of the noncontracting tree would necessarily be genuinely aperiodic. This closes the periodic/finite-state case, but information capacity does not force an arbitrary surviving path to become periodic.


\subsubsection*{Final assessment}


The proposed argument fails at one identifiable equation:

\[
H(\text{word},\text{endpoint})
\neq
H(\text{word})+r\log_2 3.
\]

The correct identity is

\[
H(\text{word},\text{endpoint})
=
H(\text{word}),
\]

because the endpoint is a deterministic hash of the word.

The approach nevertheless establishes exponential sparsity:

\[
\frac{|\text{growth-compatible endpoint residues}|}{3^j}
\le
2^{-0.0793186\,j+O(\log j)}.
\]

That is an average-case or measure-zero theorem. A full proof would require the zero-error inequality (ZE), or a comparable arithmetic uncertainty principle showing that no nontrivial infinite nested code can have simultaneously:

\[
\text{a bounded positive start representative},
\qquad
\text{an archimedean-compatible endpoint},
\qquad
\text{and the required carry/drift profile}.
\]

Proving that statement would be a genuine Collatz breakthrough. It does not follow from Shannon entropy, data processing, list-size averages, or dimension bounds; those tools leave open exactly one adaptively chosen surviving sequence.
\subsection*{References}
\begin{enumerate}[leftmargin=2.2em]
\item \hypertarget{ref:R44:1}{}\hyperlink{cit:R44:1}{\texttt{turn761298view1}}
\item \hypertarget{ref:R44:2}{}\hyperlink{cit:R44:2}{\texttt{turn761298view4}}
\item \hypertarget{ref:R44:3}{}\hyperlink{cit:R44:3}{\texttt{turn468024view0}}
\item \hypertarget{ref:R44:4}{}\hyperlink{cit:R44:4}{\texttt{turn306460view4}}
\item \hypertarget{ref:R44:5}{}\hyperlink{cit:R44:5}{\texttt{turn468024view4}}
\item \hypertarget{ref:R44:6}{}\hyperlink{cit:R44:6}{\texttt{turn468024view2}}
\item \hypertarget{ref:R44:7}{}\hyperlink{cit:R44:7}{\texttt{turn468024view1}}
\item \hypertarget{ref:R44:8}{}\hyperlink{cit:R44:8}{\texttt{turn468024view3}}
\end{enumerate}

\clearpage
\section{Round 45: approach 40 --- Mahler-function transcendence of parity series}
\subsection*{Prompt}
\begin{quote}
Complete a full proof of the Collatz conjecture using the following approach:

\begin{enumerate}
\item Mahler-function transcendence of parity series
\end{enumerate}

\textbf{Speculative \textperiodcentered{} P}

Associate to an orbit a generating series F(z) = sum over j >= 0 of epsilon\_j z\textasciicircum{}j, or a carry-weighted variant. Derive functional equations from the Collatz recursion. Show that a counterexample would force F to be both highly transcendental and algebraic or meromorphic of a prohibited type.

Take as long as you need.
\end{quote}
\subsection*{Response}

\subsection*{Result of the Mahler-function approach}


Carrying the proposal through produces a sharp analytic classification of Collatz orbits, but it does \textbf{not} furnish a full proof. The missing implication is:

\[
\text{“the seed is a positive integer”}
\quad\Longrightarrow\quad
\text{“its parity/carry series satisfies a finite Mahler equation.”}
\]

That implication does not follow from the Collatz recursion. In fact, in the precise formulation below, it is equivalent to excluding divergent orbits. A second, independent obstruction remains because a hypothetical nontrivial cycle would give a perfectly rational\textemdash{}and hence Mahler\textemdash{}series.

The current literature still treats Collatz as unproved; for example, a February 2026 verification paper explicitly describes a proof as out of reach with current methods. [\hypertarget{cit:R45:1}{}\hyperlink{ref:R45:1}{1}]

What follows is the strongest rigorous theorem yielded by the proposed approach.


\bigskip



\subsubsection*{1. Two orbit series}


Use the shortcut Collatz map

\[
T(x)=
\begin{cases}
x/2,&x\equiv0\pmod2,\\[2mm]
(3x+1)/2,&x\equiv1\pmod2.
\end{cases}
\]

For a fixed \(n\geq1\), write

\[
x_j=T^j(n),\qquad
\varepsilon_j=x_j\bmod2\in\{0,1\},\qquad
s_j=\sum_{i=0}^{j-1}\varepsilon_i ,
\]

where \(s_0=0\). Then

\[
2x_{j+1}=3^{\varepsilon_j}x_j+\varepsilon_j.
\tag{1}
\]

The ordinary parity series is

\[
F_n(z)=\sum_{j\geq0}\varepsilon_jz^j.
\tag{2}
\]

The more useful carry-weighted series is

\[
G_n(z)=\sum_{j\geq0}
\frac{\varepsilon_j}{3^{s_{j+1}}}z^j.
\tag{3}
\]

If \(d_m\) is the position of the \((m+1)\)-st odd term, then

\[
G_n(z)=\sum_{m\geq0}\frac{z^{d_m}}{3^{m+1}}.
\tag{4}
\]

Thus its support is the parity word, while its nonzero coefficients successively equal

\[
\frac13,\frac1{3^2},\frac1{3^3},\ldots.
\]


\paragraph*{Exact inverse identity}


Define

\[
h_k=\frac{2^kx_k}{3^{s_k}}.
\]

Equation (1) gives

\[
h_{j+1}
=h_j+\varepsilon_j\frac{2^j}{3^{s_{j+1}}},
\]

and hence

\[
\boxed{
\frac{2^kx_k}{3^{s_k}}
=
n+\sum_{j=0}^{k-1}
\varepsilon_j\frac{2^j}{3^{s_{j+1}}}.
}
\tag{5}
\]

In \(\mathbb Q_2\), the left side tends to \(0\), because its 2-adic valuation is at least \(k\). Therefore

\[
\boxed{
G_n(2)=-n\qquad\text{in }\mathbb Q_2.
}
\tag{6}
\]

This is the standard inverse parity transform: every infinite binary word is the parity sequence of exactly one 2-adic seed, and finite words determine one residue class modulo the corresponding power of \(2\). [\hypertarget{cit:R45:2}{}\hyperlink{ref:R45:2}{2}, \hypertarget{cit:R45:3}{}\hyperlink{ref:R45:3}{3}]

The analogous unweighted identity is

\[
F_n(2)=Q(n):=\sum_{j\geq0}\varepsilon_j2^j
\qquad\text{in }\mathbb Z_2.
\tag{7}
\]


\bigskip



\subsubsection*{2. The functional equations are shift equations}


The parity series satisfies

\[
\boxed{
F_n(z)=\varepsilon_0+zF_{T(n)}(z).
}
\tag{8}
\]

For the weighted series,

\[
\boxed{
G_n(z)
=
\frac{\varepsilon_0}{3}
+
z\,3^{-\varepsilon_0}G_{T(n)}(z).
}
\tag{9}
\]

Thus

\[
G_n(z)=
\begin{cases}
zG_{T(n)}(z),&n\text{ even},\\[2mm]
\dfrac13+\dfrac z3G_{T(n)}(z),&n\text{ odd}.
\end{cases}
\]

After \(k\) steps,

\[
\boxed{
G_n(z)
=
\sum_{j=0}^{k-1}
\frac{\varepsilon_j}{3^{s_{j+1}}}z^j
+
z^k3^{-s_k}G_{x_k}(z).
}
\tag{10}
\]

This is not a Mahler equation. A \(q\)-Mahler equation relates one fixed function at dilated arguments,

\[
A_0(z)f(z)+A_1(z)f(z^q)+\cdots+A_d(z)f(z^{q^d})=0.
\tag{11}
\]

Equations (8)\textendash{}(10) instead relate \(F_n\) or \(G_n\) to a new tail function belonging to the next state of the orbit.


\bigskip



\subsection*{3. The analytic rigidity theorem}


Let \(R_n\) denote the complex radius of convergence of \(G_n\).


\subsubsection*{Theorem}


For every positive integer \(n\), the following conditions are equivalent:

\[
\begin{aligned}
&\text{the orbit of }n\text{ is eventually periodic};\\
&F_n(z)\text{ is rational};\\
&G_n(z)\text{ is rational};\\
&G_n(z)\text{ is algebraic over }\mathbb C(z);\\
&G_n(z)\text{ is D-finite};\\
&G_n(z)\text{ is }q\text{-Mahler for at least one }q\geq2;\\
&R_n<2.
\end{aligned}
\tag{12}
\]

In the Mahler condition, ``at least one \(q\)'' can equivalently be replaced by ``every \(q\geq2\).''

Moreover:

\[
\boxed{
R_n=
\begin{cases}
\sqrt3,
&\text{if the orbit reaches }1,\\[1mm]
\text{a number in }(\sqrt3,2),
&\text{if it reaches a nontrivial positive cycle},\\[1mm]
\text{a number in }[2,3],
&\text{if the orbit is nonperiodic}.
\end{cases}}
\tag{13}
\]

Consequently,

\[
\boxed{
n\text{ satisfies Collatz}
\quad\Longleftrightarrow\quad
R_n=\sqrt3.
}
\tag{14}
\]


\paragraph*{Proof}



\subparagraph*{Eventual periodicity and \(F_n\)}


If the orbit is eventually periodic, then its parity sequence is eventually periodic, so \(F_n\) is a polynomial plus a periodic geometric tail and is rational.

Conversely, suppose \(F_n\) is rational. Its coefficient sequence satisfies a constant-coefficient linear recurrence from some point onward. Since every coefficient belongs to the finite set \(\{0,1\}\), the finite recurrence-state vector must eventually repeat. Hence \((\varepsilon_j)\) is eventually periodic.

If the parity tail has period \(p\), then \(x_N\) and \(x_{N+p}\) have the same infinite parity sequence. The 2-adic parity transform is injective, so \(x_N=x_{N+p}\). Thus the orbit itself is eventually periodic.


\subparagraph*{Eventual periodicity implies rationality of \(G_n\)}


Suppose the eventual parity period has length \(p\) and contains \(r\) ones. Then, sufficiently far along the sequence,

\[
\varepsilon_{j+p}=\varepsilon_j,
\qquad
s_{j+p+1}=s_{j+1}+r,
\]

so the coefficients \(a_j=[z^j]G_n\) satisfy

\[
a_{j+p}=3^{-r}a_j.
\tag{15}
\]

Therefore the tail of \(G_n\) has denominator

\[
1-\frac{z^p}{3^r},
\]

and \(G_n\) is rational.

Conversely, if \(G_n\) is rational, its coefficients satisfy a linear recurrence. By the Skolem\textendash{}Mahler\textendash{}Lech theorem, their zero set is, up to finitely many exceptions, a finite union of arithmetic progressions. [\hypertarget{cit:R45:4}{}\hyperlink{ref:R45:4}{4}] Since

\[
[z^j]G_n=0\quad\Longleftrightarrow\quad\varepsilon_j=0,
\]

the parity sequence is eventually periodic.


\subparagraph*{Algebraic and D-finite rigidity}


Every algebraic power series is D-finite. Conversely, the coefficients of \(G_n\) lie in

\[
\{0\}\cup\{3^{-m}:m\geq1\},
\]

which is contained in the finitely generated multiplicative group \(\langle 1/3\rangle\), together with zero. A theorem extending work of B\'ezivin states that a D-finite power series in characteristic zero whose coefficients lie in a finitely generated multiplicative group, together with zero, must be rational. [\hypertarget{cit:R45:5}{}\hyperlink{ref:R45:5}{5}]

Hence

\[
G_n\text{ algebraic}
\Longleftrightarrow
G_n\text{ D-finite}
\Longleftrightarrow
G_n\text{ rational}.
\tag{16}
\]


\subparagraph*{The radius in the nonperiodic case}


Assume the orbit is not eventually periodic. It cannot visit a bounded set infinitely often, since some state would repeat. Hence

\[
x_j\longrightarrow\infty.
\tag{17}
\]

Taking logarithms in (1) gives

\[
\log x_k-\log n
=
s_k\log3-k\log2+
\sum_{\substack{0\leq j<k\\\varepsilon_j=1}}
\log\left(1+\frac1{3x_j}\right).
\tag{18}
\]

Because \(x_j\to\infty\), the summands in the final sum tend to zero, so that sum is \(o(k)\). Since \(\log x_k\geq0\),

\[
\liminf_{k\to\infty}\frac{s_k}{k}
\geq
\frac{\log2}{\log3}.
\tag{19}
\]

For every nonzero coefficient of \(G_n\),

\[
\left|
\frac{\varepsilon_j}{3^{s_{j+1}}}
\right|^{1/j}
=
3^{-s_{j+1}/j}.
\]

Equation (19) therefore gives

\[
\limsup_{j\to\infty}|[z^j]G_n|^{1/j}
\leq\frac12,
\]

and hence

\[
R_n\geq2.
\tag{20}
\]

There are infinitely many odd terms in a positive orbit: an eventually all-even positive integer would have to be divisible by arbitrarily high powers of \(2\). Also \(s_{j+1}\leq j+1\), so along the nonzero coefficients

\[
|[z^j]G_n|^{1/j}
\geq3^{-(j+1)/j}.
\]

Thus

\[
R_n\leq3.
\tag{21}
\]

This proves \(R_n\in[2,3]\) for every nonperiodic positive orbit.


\subparagraph*{The radius in the periodic case}


If the cycle has period \(p\) and \(r\) odd entries, equation (15) gives

\[
R_n=3^{r/p}.
\tag{22}
\]

For a point \(x\) on the cycle, equation (5) over one period gives

\[
\frac{2^px}{3^r}
=
x+
\sum_{j=0}^{p-1}
\varepsilon_j\frac{2^j}{3^{s_{j+1}}}
>x.
\]

Hence

\[
3^r<2^p,
\qquad
\frac rp<\frac{\log2}{\log3},
\qquad
R_n<2.
\tag{23}
\]

For the trivial cycle \(1\leftrightarrow2\), \(p=2,r=1\), so

\[
R_n=3^{1/2}=\sqrt3.
\]

Now consider a nontrivial positive cycle. Every member is greater than \(1\). At an odd step,

\[
\frac{x_{j+1}}{x_j}
=
\frac{3x_j+1}{2x_j}
<2,
\]

whereas an even step has ratio \(1/2\). Multiplying these ratios around the cycle gives

\[
1
<
2^r\left(\frac12\right)^{p-r}
=
2^{2r-p}.
\]

Therefore

\[
r>\frac p2.
\tag{24}
\]

Combining (23) and (24),

\[
\frac12<\frac rp<\frac{\log2}{\log3},
\]

and consequently

\[
\sqrt3<R_n<2.
\tag{25}
\]

This proves the three-way radius classification (13).


\subparagraph*{Mahler rigidity}


Every rational function is \(q\)-Mahler for every \(q\geq2\).

Conversely, suppose \(G_n\) is \(q\)-Mahler. We have just shown that every positive-orbit \(G_n\) has radius strictly greater than \(1\). A theorem of Bell, Coons and Rowland states that a nonrational \(q\)-Mahler function has the unit circle as a natural boundary. [\hypertarget{cit:R45:6}{}\hyperlink{ref:R45:6}{6}] But \(G_n\) is holomorphic in a disk larger than the unit disk. Therefore it cannot be a nonrational Mahler function, and so it is rational.

This completes the proof of the theorem. \ensuremath{\blacksquare}


\bigskip



\subsubsection*{4. What happens to the direct parity series}


For a nonperiodic positive orbit, \(F_n\) has infinitely many coefficients equal to \(1\), so its radius is exactly \(1\). It is not rational, because its coefficient sequence is not eventually periodic.

The P\'olya\textendash{}Carlson theorem now gives

\[
\boxed{
\text{if the orbit is nonperiodic, then }|z|=1
\text{ is a natural boundary for }F_n.
}
\tag{26}
\]

The theorem states that an integer-coefficient power series of radius \(1\) is rational or has the unit circle as a natural boundary. [\hypertarget{cit:R45:7}{}\hyperlink{ref:R45:7}{7}] Thus a divergent counterexample would indeed force \(F_n\) to be transcendental and non-D-finite.

This establishes the ``highly transcendental'' half of the proposed strategy.

It does \textbf{not} produce a contradiction. Nonrational Mahler series with natural boundaries exist. The elementary example

\[
L(z)=\sum_{m\geq0}z^{2^m}
\]

satisfies the genuine Mahler equation

\[
L(z)=z+L(z^2),
\]

yet is nonrational and has the unit circle as a natural boundary.

More importantly, a nontrivial-cycle counterexample would have an eventually periodic parity sequence, so both \(F_n\) and \(G_n\) would be rational. A transcendence argument cannot exclude that kind of counterexample.


\bigskip



\subsection*{5. An exact generating-function formulation of Collatz}


The two members of the trivial cycle have weighted series

\[
G_1(z)
=
\frac13+\frac{z^2}{3^2}+\frac{z^4}{3^3}+\cdots
=
\frac1{3-z^2},
\tag{27}
\]

and

\[
G_2(z)=\frac z{3-z^2}.
\tag{28}
\]

Therefore:


\subsubsection*{Exact criterion}


For every \(n\geq1\),

\[
\boxed{
n\text{ reaches }1
\quad\Longleftrightarrow\quad
(3-z^2)G_n(z)\in\mathbb Q[z].
}
\tag{29}
\]


\paragraph*{Proof}


If \(x_N\in\{1,2\}\), equation (10) expresses \(G_n\) as a polynomial plus a rational multiple of either \(z^NG_1\) or \(z^NG_2\). Multiplication by \(3-z^2\) therefore gives a polynomial.

Conversely, suppose \((3-z^2)G_n\) is a polynomial. Writing \(a_j=[z^j]G_n\), we have, for all sufficiently large \(j\),

\[
3a_j=a_{j-2}.
\tag{30}
\]

Hence

\[
a_j=0\quad\Longleftrightarrow\quad a_{j-2}=0,
\]

so the parity support is eventually periodic with period \(2\).

It cannot be eventually \(00\,00\,00\cdots\), since a positive integer cannot remain even forever. It also cannot be eventually \(11\,11\,11\cdots\): in that case

\[
\frac{a_j}{a_{j-2}}
=
3^{-(s_{j+1}-s_{j-1})}
=
\frac1{9},
\]

contradicting (30), which requires the ratio \(1/3\).

Thus the tail must alternate:

\[
10\,10\,10\cdots
\quad\text{or}\quad
01\,01\,01\cdots.
\]

The unique positive 2-adic states with those parity tails are \(1\) and \(2\). Equivalently, if the alternating tail begins at an odd \(x\), then

\[
T^2(x)=\frac{3x+1}{4}=x,
\]

so \(x=1\); if it begins at an even \(x\), then

\[
T^2(x)=\frac{3x+2}{4}=x,
\]

so \(x=2\). Therefore the orbit reaches the trivial cycle. \ensuremath{\blacksquare}

There is a corresponding unweighted criterion:

\[
\boxed{
n\text{ reaches }1
\quad\Longleftrightarrow\quad
(1-z^2)F_n(z)\in\mathbb Z[z].
}
\tag{31}
\]

These are exact reformulations, not proofs that the displayed polynomial conditions always hold.


\bigskip



\subsection*{6. Why the recursion does not supply a Mahler equation}



\subsubsection*{6.1 Finite closure is equivalent to periodicity}


Equations (8) and (9) produce the tail family

\[
F_n,F_{x_1},F_{x_2},\ldots
\]

or

\[
G_n,G_{x_1},G_{x_2},\ldots.
\]

If two tail functions coincide, their parity supports coincide. The inverse parity transform is injective, so the corresponding states coincide. Therefore

\[
\{G_{x_j}:j\geq0\}\text{ is finite}
\quad\Longleftrightarrow\quad
(x_j)\text{ is eventually periodic}.
\tag{32}
\]

Thus eliminating the tail functions to obtain a finite functional system works precisely when the orbit already repeats.


\subsubsection*{6.2 No universal Mahler equation can come from local recursion}


Fix \(q\geq2\). Suppose polynomials independent of the seed satisfied

\[
\sum_{i=0}^d A_i(z)F_n(z^{q^i})=B(z)
\tag{33}
\]

for every positive integer \(n\).

Every finite binary word is the initial parity word of some positive integer\textemdash{}indeed, of one residue class modulo the appropriate power of \(2\). [\hyperlink{ref:R45:2}{2}] The coefficient of \(z^M\) in (33) depends on only finitely many parity bits. Consequently, (33) would have to hold for every infinite binary sequence.

Taking the all-zero sequence gives \(B=0\). Taking the sequence with a single \(1\) at position \(r\) gives

\[
\sum_{i=0}^d A_i(z)z^{rq^i}=0
\tag{34}
\]

for every \(r\). For sufficiently large \(r\), the polynomial supports of the terms in (34) are disjoint. Hence every \(A_i\) is zero.

The same argument applies to \(G_n\), since a word containing a single \(1\) at position \(r\) gives \(G(z)=z^r/3\).

Therefore no fixed, nontrivial, finite Mahler equation can be derived solely from the local Collatz rules and hold uniformly for all seeds.


\bigskip



\subsection*{7. There is a genuine Mahler function\textemdash{}but in the wrong variable}


Let

\[
Q(x)=\sum_{j\geq0}\varepsilon_j(x)2^j
\]

be the parity-coordinate map, and let

\[
\Phi=Q^{-1}.
\]

The inverse map satisfies

\[
\Phi(2m)=2\Phi(m),
\qquad
\Phi(2m+1)=\frac{2\Phi(m)-1}{3}.
\tag{35}
\]

These are established functional equations for the 2-adic conjugacy. [\hyperlink{ref:R45:3}{3}]

For ordinary nonnegative integers \(m\), put

\[
b_m=\Phi(m),
\qquad
B(z)=\sum_{m\geq0}b_mz^m.
\]

Splitting the sum into even and odd indices gives

\[
\begin{aligned}
B(z)
&=
\sum_{m\geq0}2b_mz^{2m}
+
\sum_{m\geq0}\frac{2b_m-1}{3}z^{2m+1}\\
&=
\boxed{
\frac{2(3+z)}3B(z^2)
-\frac{z}{3(1-z^2)}.
}
\end{aligned}
\tag{36}
\]

This is a genuine inhomogeneous 2-Mahler equation.

Unfortunately, a Collatz seed is not obtained by evaluating \(B\) at a fixed algebraic argument.

Let

\[
q_k=\sum_{j=0}^{k-1}\varepsilon_j2^j,
\]

the integer represented by the first \(k\) parity bits. Then

\[
b_{q_k}
=
-\sum_{j=0}^{k-1}
\varepsilon_j\frac{2^j}{3^{s_{j+1}}},
\]

so (5) becomes

\[
\boxed{
n-b_{q_k}
=
\frac{2^kx_k}{3^{s_k}}.
}
\tag{37}
\]

Thus

\[
b_{q_k}\longrightarrow n
\qquad\text{2-adically}.
\tag{38}
\]

The seed is therefore a limit of \textbf{coefficients of \(B\) sampled along a variable, parity-dependent sequence of indices}. Mahler's method controls functional values such as \(B(\alpha)\) and finite-dimensional Mahler systems. It does not convert an arbitrary nested coefficient path

\[
q_0,q_1,q_2,\ldots
\]

into an algebraic-value problem. That coefficient path contains exactly the original Collatz parity information.


\bigskip



\subsection*{8. Using both dilations \(2\) and \(3\) does not close the argument}


For a nonperiodic orbit, all \(x_j\) are distinct. Consider its orbit-set series

\[
A_n(z)=\sum_{j\geq0}z^{x_j}.
\]

Splitting the orbit states into even and odd ones yields

\[
\boxed{
2A_n(z^2)-A_n(z)-A_n(-z)
-zA_n(z^3)+zA_n(-z^3)
=
2z^{2n}.
}
\tag{39}
\]

This genuinely contains the two multiplicatively independent dilations \(z\mapsto z^2\) and \(z\mapsto z^3\).

However, it is one \textbf{mixed} equation, not separate 2-Mahler and 3-Mahler equations. The theorem of Adamczewski and Bell says that a power series which is separately \(2\)-Mahler and \(3\)-Mahler must be rational. [\hypertarget{cit:R45:8}{}\hyperlink{ref:R45:8}{8}] It does not apply to (39).

A simple counterexample demonstrates the distinction. Let

\[
H(z)=\sum_{a,b\geq0}z^{2^a3^b}.
\]

Then

\[
\boxed{
H(z)=z+H(z^2)+H(z^3)-H(z^6).
}
\tag{40}
\]

Its coefficients belong to \(\{0,1\}\), and it is nonrational because the infinite set of \(3\)-smooth exponents is not eventually periodic. Thus a mixed \(2/3\)-dilation identity can coexist with a nonrational series.


\bigskip



\subsection*{9. The exact stopping point}


The Mahler-parity program proves the following complete conditional statement:

\[
\boxed{
\bigl[\text{for every }n,\ G_n
\text{ is }q\text{-Mahler}\bigr]
\iff
\bigl[\text{every positive Collatz orbit is eventually periodic}\bigr].
}
\tag{41}
\]

It also proves

\[
\boxed{
\text{Collatz}
\iff
\bigl[(3-z^2)G_n(z)\in\mathbb Q[z]
\text{ for every }n\geq1\bigr]
\iff
\bigl[R_n=\sqrt3
\text{ for every }n\geq1\bigr].
}
\tag{42}
\]

For a hypothetical divergent orbit:

\[
F_n\text{ has the unit circle as a natural boundary},
\]

while

\[
G_n\text{ has }R_n\geq2
\text{ and is transcendental, non-D-finite, and non-Mahler}.
\]

For a hypothetical nontrivial positive cycle:

\[
F_n,G_n\text{ are rational and Mahler},
\qquad
\sqrt3<R_n<2.
\]

So the proposed contradiction cannot cover both counterexample mechanisms:

\begin{enumerate}
\item \textbf{Divergence} would have to be excluded by proving a new Mahler or D-finite regularity theorem for \(G_n\). By (12), that theorem would already imply eventual periodicity.

\item \textbf{Nontrivial cycles} lie entirely inside the rational/Mahler class and require a separate pole- or denominator-exclusion argument.
\end{enumerate}

Finally, because \(F_n(2)=Q(n)\) in \(\mathbb Z_2\), deriving rationality of \(F_n\) merely from \(n\in\mathbb Z_{>0}\) is precisely the assertion that \(Q(n)\) is a rational 2-adic number. This is the positive-integer part of Lagarias' Periodicity Conjecture; full Collatz requires the stronger assertion that the eventual period is specifically the alternating \(1,0\) tail. These equivalences are explicitly recorded in the literature on the 2-adic parity conjugacy. [\hypertarget{cit:R45:9}{}\hyperlink{ref:R45:9}{9}, \hyperlink{ref:R45:3}{3}]

Therefore, the Mahler-function approach yields a rigorous and unusually clean analytic reformulation\textemdash{}especially the radius trichotomy (13)\textemdash{}but the hoped-for algebraicity implication is not supplied by the recursion. As stated, inserting it would be circular rather than a completed proof.
\subsection*{References}
\begin{enumerate}[leftmargin=2.2em]
\item \hypertarget{ref:R45:1}{}\hyperlink{cit:R45:1}{\texttt{turn924229view4}}
\item \hypertarget{ref:R45:2}{}\hyperlink{cit:R45:2}{\texttt{turn788726view1}}
\item \hypertarget{ref:R45:3}{}\hyperlink{cit:R45:3}{\texttt{turn788726view3}}
\item \hypertarget{ref:R45:4}{}\hyperlink{cit:R45:4}{\texttt{turn268785search1}}
\item \hypertarget{ref:R45:5}{}\hyperlink{cit:R45:5}{\texttt{turn924229view1}}
\item \hypertarget{ref:R45:6}{}\hyperlink{cit:R45:6}{\texttt{turn924229view0}}
\item \hypertarget{ref:R45:7}{}\hyperlink{cit:R45:7}{\texttt{turn268785search0}}
\item \hypertarget{ref:R45:8}{}\hyperlink{cit:R45:8}{\texttt{turn924229view2}}
\item \hypertarget{ref:R45:9}{}\hyperlink{cit:R45:9}{\texttt{turn788726view2}}
\end{enumerate}

\clearpage
\section{Round 46: approach 41 --- Christol-Cobham algebraicity trap}
\subsection*{Prompt}
\begin{quote}
Complete a full proof of the Collatz conjecture using the following approach:

\begin{enumerate}
\item Christol-Cobham algebraicity trap
\end{enumerate}

\textbf{Speculative \textperiodcentered{} P}

Reduce an orbit-generating series modulo suitable primes. Algebraicity over a finite field would imply automaticity; automaticity in two independent bases would force periodic structure. The real work is deriving algebraicity from exact affine carries rather than assuming it.

Take as long as you need.
\end{quote}
\subsection*{Response}

\subsection*{Christol\textendash{}Cobham analysis of a Collatz orbit}


Carrying this program through does \textbf{not} produce a full proof of the Collatz conjecture. It does produce an exact characterization showing where the approach becomes circular:


\begin{quote}
For a fixed positive integer \(n\), algebraicity of its parity series in both characteristics \(2\) and \(3\) is equivalent to eventual periodicity of its Collatz orbit.

\end{quote}


Thus the proposed algebraicity statement is not presently a consequence of the affine carry formulas; it is essentially the ``no divergent orbit'' part of Collatz written in finite-field language. Moreover, a hypothetical nontrivial cycle already has a rational parity series, so Christol\textendash{}Cobham cannot by itself exclude cycles.

There is, however, a better value-indexed generating series with an exact mixed \(2\)-Cartier/\(3\)-substitution equation. For it one can rigorously exclude infinite ultimately periodic ranges. I give that result below, followed by a proof that the reduced functional equation still leaves an arbitrary Cartier section and therefore does not force algebraicity.


\bigskip



\subsubsection*{1. Exact affine identities}


Use the shortcut Collatz map

\[
T(x)=
\begin{cases}
x/2,&x\equiv0\pmod 2,\\[2mm]
(3x+1)/2,&x\equiv1\pmod 2.
\end{cases}
\]

Fix \(n\geq1\), and write

\[
x_j=T^j(n),\qquad
\varepsilon_j=x_j\bmod 2\in\{0,1\},
\qquad
s_j=\sum_{i=0}^{j-1}\varepsilon_i .
\]

Then

\[
2x_{j+1}=3^{\varepsilon_j}x_j+\varepsilon_j .
\tag{1}
\]

Iterating (1) gives

\[
2^k x_k=3^{s_k}n+C_k,
\tag{2}
\]

where

\[
C_k
=
\sum_{j=0}^{k-1}
\varepsilon_j\,2^j\,3^{\,s_k-s_{j+1}},
\tag{3}
\]

and equivalently

\[
C_{k+1}=3^{\varepsilon_k}C_k+\varepsilon_k2^k,
\qquad C_0=0.
\tag{4}
\]

These are the exact affine carries one would like to convert into algebraicity.

Define

\[
E_n(z)=\sum_{j\geq0}\varepsilon_jz^j,
\qquad
X_n(z)=\sum_{j\geq0}x_jz^j.
\]

Since \(3^{\varepsilon_j}=1+2\varepsilon_j\), equation (1) gives

\[
\frac{2(X_n(z)-n)}{z}
=
X_n(z)+2\bigl(E_n\odot X_n\bigr)(z)+E_n(z),
\tag{5}
\]

where

\[
(E_n\odot X_n)(z)=\sum_{j\geq0}\varepsilon_jx_jz^j
\]

is the Hadamard product.

Equation (5) is exact, but it is not an algebraic equation for \(E_n\): it contains the independent series \(X_n\) and their coefficientwise product. In characteristic \(2\), it collapses merely to

\[
\overline X_n=\overline E_n,
\]

which says only that \(x_j\bmod2=\varepsilon_j\).


\bigskip



\subsubsection*{2. The exact Christol\textendash{}Cobham criterion}


Christol's theorem says that a power series over a finite field of characteristic \(p\) is algebraic precisely when its coefficient sequence is \(p\)-automatic. A finite-field series is rational precisely when its coefficients are ultimately periodic. Cobham's theorem says that a sequence automatic in two multiplicatively independent bases is ultimately periodic. Since \(2\) and \(3\) are multiplicatively independent, these theorems apply directly to \(E_n\). [\hypertarget{cit:R46:1}{}\hyperlink{ref:R46:1}{1}, \hypertarget{cit:R46:2}{}\hyperlink{ref:R46:2}{2}]


\paragraph*{Theorem 1}


For every \(n\geq1\), the following conditions are equivalent.

\begin{enumerate}
\item The reduction of \(E_n(z)\) modulo \(2\) is algebraic over \(\mathbb F_2(z)\), and its reduction modulo \(3\) is algebraic over \(\mathbb F_3(z)\).
\item The parity sequence \((\varepsilon_j)_{j\geq0}\) is ultimately periodic.
\item The orbit \((T^j(n))_{j\geq0}\) is eventually periodic.
\end{enumerate}


\paragraph*{Proof}


Assume condition 1. By Christol's theorem, \((\varepsilon_j)\) is simultaneously \(2\)-automatic and \(3\)-automatic. Cobham's theorem therefore makes it ultimately periodic. Thus \(1\Rightarrow2\).

If \((\varepsilon_j)\) is ultimately periodic, then

\[
E_n(z)=P(z)+\frac{z^NQ(z)}{1-z^L}
\]

for suitable polynomials \(P,Q\) and integers \(N,L\). Hence \(E_n\) is rational in every characteristic. Thus \(2\Rightarrow1\).

It remains to relate parity periodicity to orbit periodicity. The \(2\)-adic parity map

\[
Q(x)=\sum_{j\geq0}\bigl(T^j(x)\bmod2\bigr)2^j
\]

is a bijection \(\mathbb Z_2\to\mathbb Z_2\) and conjugates \(T\) with the binary shift:

\[
Q\circ T=\mathcal S\circ Q.
\tag{6}
\]

Every binary sequence is the parity sequence of exactly one \(2\)-adic integer. In particular, an ultimately periodic parity tail is fixed by some power of the shift. By the injectivity of \(Q\), the corresponding Collatz iterate is fixed by the same power of \(T\). Hence the ordinary integer orbit is eventually periodic. The converse is immediate. [\hypertarget{cit:R46:3}{}\hyperlink{ref:R46:3}{3}, \hypertarget{cit:R46:4}{}\hyperlink{ref:R46:4}{4}]

This proves the theorem. \ensuremath{\blacksquare}


\paragraph*{Immediate consequences}


If the orbit of \(n\) is not eventually periodic, then

\[
E_n(z)\bmod2
\quad\text{or}\quad
E_n(z)\bmod3
\]

must be transcendental over its corresponding rational function field.

On the other hand, if \(n\) lies in a hypothetical nontrivial cycle, then \(E_n(z)\) is periodic and therefore rational over every finite field. Consequently:


\begin{quote}
Dual finite-field algebraicity excludes divergent orbits, but it cannot distinguish the trivial cycle from a nontrivial cycle.

\end{quote}


This is an unavoidable obstruction, not merely a weakness of the proof technique.


\bigskip



\subsubsection*{3. Why periodicity does not finish Collatz}


Suppose a periodic parity word has length \(L\), contains \(q\) ones, and its ones occur at

\[
0\leq d_0<d_1<\cdots<d_{q-1}<L.
\]

Iterating one period gives

\[
T^L(x)=\frac{3^q x+C_w}{2^L},
\tag{7}
\]

where

\[
C_w=\sum_{i=0}^{q-1}3^{q-1-i}2^{d_i}.
\tag{8}
\]

A positive cycle therefore requires

\[
(2^L-3^q)x=C_w.
\tag{9}
\]

Christol\textendash{}Cobham gives only that some word \(w\) repeats. It provides no reason that (9) cannot have a positive integral solution. Indeed every periodic word is automatic in every base, regardless of whether its cycle equation is integrally realizable.

Thus the proposed route would still need an independent theorem saying that the only positive integral realization of (9) compatible with its own parity word is the word \(10\), corresponding to

\[
1\longmapsto2\longmapsto1.
\]

That is precisely the nontrivial-cycle problem.


\bigskip



\subsubsection*{4. The index mismatch in the algebraicity step}


For a sequence \(a=(a_j)\), Christol's theorem concerns finiteness of its \(p\)-kernel

\[
\mathcal K_p(a)=
\left\{
(a_{p^e m+r})_{m\geq0}:
e\geq0,\ 0\leq r<p^e
\right\}.
\tag{10}
\]

Thus algebraicity in characteristic \(p\) requires control of the sequence at the decimated times

\[
p^e m+r.
\]

The Collatz carry identities (1)\textendash{}(4), however, connect adjacent times \(j\) and \(j+1\), while the factors \(2\) and \(3\) act on the \textbf{values} \(x_j\), not on the time index \(j\).

There is no identity in (1)\textendash{}(4) relating, for example,

\[
(\varepsilon_{3m})_{m\geq0},
\qquad
(\varepsilon_{3m+1})_{m\geq0},
\qquad
(\varepsilon_{3m+2})_{m\geq0}
\]

to finitely many transforms of the original sequence. The occurrence of the integer \(3\) in \(3x+1\) does not create base-\(3\) automaticity in the orbit-time variable.

This distinction is easy to overlook:

\[
\boxed{\text{arithmetic base acting on values}}
\quad\neq\quad
\boxed{\text{automatic base acting on coefficient indices}}.
\]


\bigskip



\subsubsection*{5. Local affine carries permit arbitrary parity data}


The obstruction is stronger than the absence of an evident functional equation.

The \(2\)-adic parity map is onto: every infinite binary sequence \(S=(s_j)\) occurs as the parity sequence of a unique \(2\)-adic seed. Its inverse is

\[
x
=
-\sum_{j\geq0}
s_j\,2^j\,3^{-\sigma_j(S)}
\qquad\text{in }\mathbb Z_2,
\tag{11}
\]

where

\[
\sigma_j(S)=\sum_{i=0}^j s_i.
\]

Furthermore, every finite binary word of length \(k\) is realized by exactly one residue class modulo \(2^k\), and hence by infinitely many positive integers. [\hypertarget{cit:R46:5}{}\hyperlink{ref:R46:5}{5}, \hyperlink{ref:R46:3}{3}]

There are uncountably many binary sequences but only countably many algebraic series over \(\mathbb F_p(z)\). Therefore identities valid for all \(2\)-adic Collatz seeds\textemdash{}including the affine carry formulas\textemdash{}cannot by themselves imply algebraicity.

A successful argument would have to use the genuinely global restriction that one fixed \(2\)-adic seed is an ordinary positive integer. No bounded carry window or finite parity-prefix calculation can see that distinction: every such finite prefix already occurs for positive integers.


\bigskip



\subsubsection*{6. An explicit carry exhibiting the wrong automatic base}


There is a concrete finite-word manifestation of the same problem.

Let

\[
m=\sum_{j=0}^{L-1}\eta_j2^j,
\qquad \eta_{L-1}=1,
\]

encode a binary word, and let \(q=\sum_j\eta_j\). Define its affine carry

\[
\mathcal C(m)
=
\sum_{j=0}^{L-1}
\eta_j2^j3^{\,q-\sum_{i=0}^{j}\eta_i}.
\tag{12}
\]

Modulo \(3\), every summand vanishes except the one corresponding to the last \(1\). Hence

\[
\mathcal C(m)\equiv 2^{L-1}
=
2^{\lfloor\log_2m\rfloor}
\pmod3.
\tag{13}
\]

The sequence on the right is constant on each dyadic interval

\[
[2^r,2^{r+1})
\]

and alternates between \(1\) and \(2\) modulo \(3\). It is \(2\)-automatic: an automaton only needs to record the parity of the binary length. It is not ultimately periodic, because its alternating constant blocks have unbounded lengths. Therefore Cobham implies that it is not \(3\)-automatic, and Christol implies that its generating series over \(\mathbb F_3\) is transcendental.

So even after reduction modulo \(3\), the exact affine carry naturally remembers \textbf{base \(2\)}, not base \(3\). This is the opposite of what the proposed algebraicity step needs.


\bigskip



\subsection*{7. Repair: index the series by orbit values}


The time-indexed parity series misaligns with multiplication by \(2\) and \(3\). A better object is the indicator series of the orbit's range.

Assume temporarily that the orbit of \(n\) is not eventually periodic. It is then injective. Put

\[
a_m=
\begin{cases}
1,&m=T^j(n)\text{ for some }j\geq0,\\
0,&\text{otherwise},
\end{cases}
\]

and

\[
A_n(z)=\sum_{m\geq0}a_mz^m.
\tag{14}
\]

Every visited \(m\neq n\) has exactly one visited predecessor. Its possible predecessors are

\[
2m
\]

and, when \(m\equiv2\pmod3\),

\[
\frac{2m-1}{3}.
\]

The two alternatives cannot both occur in an injective orbit. Consequently,

\[
a_m-\mathbf 1_{\{m=n\}}
=
a_{2m}
+
\mathbf 1_{\{m\equiv2\pmod3\}}
a_{(2m-1)/3}.
\tag{15}
\]

Introduce the binary Cartier sections

\[
\Lambda_0 A(z)=\sum_{m\geq0}a_{2m}z^m,
\qquad
\Lambda_1 A(z)=\sum_{m\geq0}a_{2m+1}z^m.
\]

Summing (15) gives the exact functional equation

\[
\boxed{
A_n(z)-z^n
=
\Lambda_0A_n(z)
+
z^2\bigl(\Lambda_1A_n\bigr)(z^3).
}
\tag{16}
\]

This is much closer to the intended Christol\textendash{}Cobham structure: division by \(2\) becomes a Cartier section, while multiplication by \(3\) becomes substitution by \(z^3\).

Replacing \(z\) by \(z^2\) and eliminating the sections yields, over characteristic zero,

\[
\boxed{
2A_n(z^2)-2z^{2n}
=
A_n(z)+A_n(-z)
+
zA_n(z^3)-zA_n(-z^3).
}
\tag{17}
\]

This equation has a nontrivial consequence.


\bigskip



\subsubsection*{8. An injective orbit cannot have an ultimately periodic range}



\paragraph*{Lemma 2}


Suppose \(A(z)=\sum a_mz^m\), with \(a_m\in\{0,1\}\), satisfies (17). If \((a_m)\) is ultimately periodic, then it has finite support.

In particular, the range of an infinite injective Collatz orbit is not ultimately periodic as a subset of the positive integers.


\paragraph*{Proof}


Suppose instead that \((a_m)\) is ultimately periodic and contains infinitely many ones. Let \(M\) be its least eventual period. Then \(A(z)\) is rational, with poles at roots of unity whose orders divide \(M\).

Fourier decomposition of the period-\(M\) tail shows that the least common multiple of the orders of these poles is \(M\).


\paragraph*{Case 1: \(M\) is odd}


Let

\[
\delta>0
\]

be the mean number of ones per period. As \(z\to1^{-}\),

\[
A(z)\sim\frac{\delta}{1-z},
\qquad
A(z^2)\sim\frac{\delta}{2(1-z)},
\qquad
A(z^3)\sim\frac{\delta}{3(1-z)}.
\]

Because \(M\) is odd, \(-1\) is not a pole of \(A\), so \(A(-z)\) and \(A(-z^3)\) remain bounded as \(z\to1\).

Taking the principal parts at \(z=1\) in (17), its left side has coefficient

\[
2\cdot\frac{\delta}{2}=\delta,
\]

while the right side has coefficient

\[
\delta+\frac{\delta}{3}=\frac{4\delta}{3}.
\]

Thus

\[
\delta=\frac{4\delta}{3},
\]

contradicting \(\delta>0\).


\paragraph*{Case 2: \(M\) is even}


Put

\[
t=v_2(M)\geq1.
\]

Since the least common multiple of the pole orders is \(M\), there is a pole \(\alpha\) whose order \(d\) satisfies

\[
v_2(d)=t.
\]

Choose a square root \(\beta\) of \(\alpha\) having order \(2d\). Then \(A(z^2)\) has a pole at \(z=\beta\).

On the other hand, each of

\[
\beta,\quad -\beta,\quad \beta^3,\quad -\beta^3
\]

has order whose \(2\)-adic valuation is \(t+1\). None can therefore be a pole of \(A\), since every pole order divides \(M\) and hence has \(2\)-adic valuation at most \(t\).

Thus at \(z=\beta\), the left side of (17) has a pole, while every term on the right side is analytic. This is impossible.

Both cases give contradictions. Hence an ultimately periodic solution must have finite support. \ensuremath{\blacksquare}


\bigskip



\subsubsection*{9. A second exact algebraicity criterion}


Combining Lemma 2 with Christol and Cobham gives another characterization.


\paragraph*{Theorem 3}


For each \(n\geq1\), define \(A_n(z)\) by the range indicator (14), whether or not the orbit is finite. Then the following are equivalent:

\begin{enumerate}
\item \(A_n(z)\bmod2\) is algebraic over \(\mathbb F_2(z)\) and \(A_n(z)\bmod3\) is algebraic over \(\mathbb F_3(z)\).
\item The orbit of \(n\) is eventually periodic.
\item \(A_n(z)\) is a polynomial.
\end{enumerate}


\paragraph*{Proof}


If the orbit is eventually periodic, its range is finite, so \(A_n\) is a polynomial.

Conversely, suppose both reductions are algebraic. Christol makes \((a_m)\) both \(2\)-automatic and \(3\)-automatic; Cobham makes it ultimately periodic. If the orbit were infinite, it would be injective, and Lemma 2 would force \(A_n\) to have finite support, a contradiction. Therefore the orbit is eventually periodic. \ensuremath{\blacksquare}

Thus changing from the parity series to the range series produces a stronger exact functional equation, but dual algebraicity is again equivalent to eventual cycling rather than a separately established fact.


\bigskip



\subsection*{10. The reduced equation still does not force algebraicity}


Equation (16) might appear sufficiently rigid to imply algebraicity in characteristic \(2\). It does not.

Work over \(\mathbb F_2[[z]]\), and write

\[
E=\Lambda_0A,
\qquad
O=\Lambda_1A.
\]

Reconstruction from the even and odd sections gives

\[
A(z)=E(z^2)+zO(z^2).
\]

In characteristic \(2\),

\[
E(z^2)=E(z)^2,
\qquad
O(z^2)=O(z)^2,
\]

so

\[
A=E^2+zO^2.
\tag{18}
\]

Reducing (16) modulo \(2\) and substituting (18) gives

\[
\boxed{
E^2+E
=
z^n+zO^2+z^2O(z^3).
}
\tag{19}
\]

Now choose \textbf{any}

\[
O(z)\in\mathbb F_2[[z]].
\]

The right side \(R(z)\) of (19) has zero constant coefficient. Writing

\[
E(z)=\sum_{m\geq0}e_mz^m,
\qquad e_0=0,
\]

the Artin\textendash{}Schreier equation \(E^2+E=R\) is solved recursively by

\[
e_m=R_m \quad(m\text{ odd}),
\tag{20}
\]

and

\[
e_{2k}=R_{2k}+e_k
\quad(k\geq1).
\tag{21}
\]

Thus every choice of the odd Cartier section \(O\) gives a solution \(E\), and hence a solution \(A\), of the finite-field reduction of (16).

There are uncountably many choices of \(O\), while only countably many series are algebraic over \(\mathbb F_2(z)\). Therefore:


\begin{quote}
The exact reduced orbit-set equation has uncountably many transcendental solutions. It does not force algebraicity.

\end{quote}


These solutions need not lift to genuine collision-free positive Collatz orbits. That is precisely the missing global information. The mod-\(2\) equation has forgotten the nonnegative integer equality and the condition that the set is one single forward orbit.

In characteristic \(3\), substitution by \(z^3\) becomes Frobenius, but the binary Cartier sections \(\Lambda_0,\Lambda_1\) remain uncontrolled. Thus each characteristic simplifies one side of the mixed \(2/3\) equation while leaving the other side open:

\[
\begin{array}{c|c|c}
\text{characteristic}&\text{simplified operation}&\text{uncontrolled operation}\\ \hline
2&z\mapsto z^2&z\mapsto z^3\\
3&z\mapsto z^3&\text{binary Cartier sections}.
\end{array}
\]

This is the finite-field form of the carry obstruction.


\bigskip



\subsubsection*{11. Why a two-Mahler repair also stops here}


There is a characteristic-zero analogue: a power series satisfying both a finite-order \(2\)-Mahler equation and a finite-order \(3\)-Mahler equation must be rational. [\hypertarget{cit:R46:6}{}\hyperlink{ref:R46:6}{6}]

Equation (17), however, is not a pair of independent Mahler equations. It mixes

\[
A(z^2),\quad A(z^3),\quad A(-z),\quad A(-z^3)
\]

in one relation. Repeatedly substituting \(z^2\) and \(z^3\) creates an expanding family of mixed arguments rather than a finite-dimensional \(2\)- or \(3\)-Mahler module. Equation (19) identifies the underlying reason: one Cartier section remains free.


\bigskip



\subsection*{12. Exact endpoint of approach 41}


The proposed proof chain was

\[
\text{exact affine carries}
\Longrightarrow
\text{algebraicity modulo two primes}
\Longrightarrow
\text{automaticity in two bases}
\Longrightarrow
\text{periodicity}
\Longrightarrow
\text{Collatz}.
\]

The rigorous outcome is instead

\[
\text{dual algebraicity of }E_n
\iff
\text{eventual periodicity of the orbit},
\tag{22}
\]

and also

\[
\text{dual algebraicity of }A_n
\iff
\text{eventual periodicity of the orbit}.
\tag{23}
\]

The middle Christol\textendash{}Cobham step is completely valid. The failure occurs on both sides:

\begin{enumerate}
\item \textbf{Before Christol:} the affine carry equations do not force algebraicity. They act on orbit values, while Christol's kernels decimate the time or value index. The strongest natural value-series equation retains a free Cartier section after reduction.
\item \textbf{After Cobham:} eventual periodicity allows every hypothetical nontrivial positive cycle. Periodic series are rational in every characteristic, so algebraicity cannot exclude them.
\end{enumerate}

Consequently, a full proof along this route would still require two genuinely new theorems:

\[
\boxed{\text{Global integer-lifting rigidity}}
\]

showing that the free finite-field Cartier sections cannot arise from one positive integral orbit, and

\[
\boxed{\text{Cycle classification}}
\]

showing that the only positive periodic parity word satisfying its affine cycle equation is \(10\).

The first statement is essentially the no-divergent-orbit component of Collatz; the second is the no-nontrivial-cycle component. The Christol\textendash{}Cobham formulation gives a clean and rigorous restatement of those obstacles, together with the exact range equation (16) and the periodic-range exclusion lemma, but it does not supply the missing global rigidity needed for a complete proof.
\subsection*{References}
\begin{enumerate}[leftmargin=2.2em]
\item \hypertarget{ref:R46:1}{}\hyperlink{cit:R46:1}{\texttt{turn320968view0}}
\item \hypertarget{ref:R46:2}{}\hyperlink{cit:R46:2}{\texttt{turn304222view3}}
\item \hypertarget{ref:R46:3}{}\hyperlink{cit:R46:3}{\texttt{turn304222view1}}
\item \hypertarget{ref:R46:4}{}\hyperlink{cit:R46:4}{\texttt{turn304222view2}}
\item \hypertarget{ref:R46:5}{}\hyperlink{cit:R46:5}{\texttt{turn304222view0}}
\item \hypertarget{ref:R46:6}{}\hyperlink{cit:R46:6}{\texttt{turn253926view5}}
\end{enumerate}

\clearpage
\section{Round 47: approach 42 --- Aperiodic tilings of Collatz spacetime}
\subsection*{Prompt}
\begin{quote}
Complete a full proof of the Collatz conjecture using the following approach:

\begin{enumerate}
\item Aperiodic tilings of Collatz spacetime
\end{enumerate}

\textbf{Speculative \textperiodcentered{} B}

Build local Wang tiles whose horizontal direction performs binary iteration and whose vertical direction performs ternary conversion. Prove that every tiling with a finite-support boundary becomes periodic, or find an aperiodic finite-boundary tiling corresponding to an infinite orbit.

Take as long as you need.
\end{quote}
\subsection*{Response}

\subsection*{Aperiodic tilings of Collatz spacetime: exact construction and obstruction}


After making the tiling construction completely local, the proposed final step splits in two:

\begin{enumerate}
\item For the unrestricted arithmetic tile set, ``every finite-boundary tiling becomes periodic'' is \textbf{false}. There is an explicit finite-boundary aperiodic tiling whose columns are the ternary expansions of \(1,2,4,8,\ldots\).
\item After imposing the boundary geometry that makes the tiling encode an actual positive-integer Collatz orbit, eventual periodicity is \textbf{equivalent to eventual cycling of that orbit}. One must then separately exclude nontrivial positive cycles.
\end{enumerate}

Thus I do not obtain a valid full proof of Collatz. What follows is a complete rigorous reduction, together with several no-go theorems showing why the natural local-periodicity and height-function arguments do not close the remaining gap.


\subsubsection*{1. An honest Wang tile system}


Use the shortcut Collatz map

\[
T(n)=
\begin{cases}
n/2,&n\equiv0\pmod2,\\[2mm]
(3n+1)/2,&n\equiv1\pmod2.
\end{cases}
\]

This is equivalent to the usual \(3n+1\) formulation.

For \(x\in\{0,1,2,3,4,5\}\), define a Wang tile with north and south colors in \(\{0,1\}\), and east and west colors in \(\{0,1,2\}\), by

\[
N=\left\lfloor\frac{x}{3}\right\rfloor,\qquad
E=x\bmod 3,\qquad
W=\left\lfloor\frac{x}{2}\right\rfloor,\qquad
S=x\bmod 2.
\]

Explicitly:

\[
\begin{array}{c|cccc}
x&N&E&W&S\\ \hline
0&0&0&0&0\\
1&0&1&0&1\\
2&0&2&1&0\\
3&1&0&1&1\\
4&1&1&2&0\\
5&1&2&2&1
\end{array}
\]

Every tile satisfies the local arithmetic law

\[
\boxed{\,3N+E=2W+S\,}. \tag{1}
\]

Conversely, for every pair \((N,E)\in\{0,1\}\times\{0,1,2\}\), equation (1) determines a unique pair \((W,S)\). Thus the system is northeast-deterministic.

These six tiles are the Collatz tiles of Cook\textendash{}St\'erin\textendash{}Woods. They simultaneously implement multiplication by \(3\) in binary and division by \(2\) in ternary; the earlier St\'erin\textendash{}Woods quasi-cellular automaton used the same arithmetic rule but required a nonlocal bootstrap operation. [\hypertarget{cit:R47:1}{}\hyperlink{ref:R47:1}{1}, \hypertarget{cit:R47:2}{}\hyperlink{ref:R47:2}{2}, \hypertarget{cit:R47:3}{}\hyperlink{ref:R47:3}{3}]


\paragraph*{Localizing the nonlocal operation}


Introduce one control color \(\sigma\) and two more tiles

\[
D_0=(N,E,W,S)=(0,\sigma,\sigma,0),
\]

\[
D_1=(N,E,W,S)=(1,\sigma,2,0).
\]

The control signal \(\sigma\) moves across trailing binary zeroes using \(D_0\). When it encounters the first binary \(1\), tile \(D_1\) consumes the control signal and emits the ternary carry \(2\). This localizes both deletion of trailing binary zeroes and the bootstrap needed for the \(3n+1\) branch. These are exactly the two control tiles displayed in the published eight-tile construction. [\hypertarget{cit:R47:4}{}\hyperlink{ref:R47:4}{4}, \hyperlink{ref:R47:1}{1}]

Given a finite binary word along the north boundary and a control ray along the east boundary, every interior tile is now determined locally. No nonlocal update rule remains.


\bigskip



\subsubsection*{2. A weighted discrete Stokes theorem}


Consider a \(w\times h\) rectangular tiled region. Number columns \(i=0,\ldots,w-1\) from east to west and rows \(j=0,\ldots,h-1\) from south to north.

Define the four boundary numbers by

\[
[N]_{2}=\sum_{i=0}^{w-1}N_i2^i,\qquad
[S]_{2}=\sum_{i=0}^{w-1}S_i2^i,
\]

\[
[E]_{3}=\sum_{j=0}^{h-1}E_j3^j,\qquad
[W]_{3}=\sum_{j=0}^{h-1}W_j3^j.
\]

For an arithmetic tile at \((i,j)\), multiply (1) by \(2^i3^j\), then sum over the rectangle. Every interior edge cancels: a north edge receives the same weight as the matching south edge above it, and a west edge receives the same weight as the matching east edge to its west. Therefore

\[
\boxed{\;
3^h[N]_{2}+[E]_{3}
=
2^w[W]_{3}+[S]_{2}.
\;} \tag{2}
\]

This is the rectangular identity appearing in the tile construction, now viewed as a discrete divergence theorem. [\hyperlink{ref:R47:1}{1}]


\paragraph*{Control tiles as point charges}


Assign the numerical value \(0\) to \(\sigma\) solely for this calculation, and define the local defect

\[
d=3N+E-2W-S.
\]

Then

\[
d(D_0)=0,\qquad d(D_1)=3-4=-1.
\]

Repeating the Stokes calculation gives

\[
3^h[N]_{2}+[E]_{3}
-
2^w[W]_{3}-[S]_{2}
=
-\sum_{D_1\text{ at }(i,j)}2^i3^j,
\]

or equivalently

\[
\boxed{\;
2^w[W]_{3}+[S]_{2}
=
3^h[N]_{2}+[E]_{3}
+
\sum_{D_1\text{ at }(i,j)}2^i3^j.
\;} \tag{3}
\]

Thus each odd Collatz event is literally a positive point charge of weight \(2^i3^j\).


\subsubsection*{3. Recovering the exact Collatz affine formula}


Let

\[
x_t=T^t(n),\qquad
\epsilon_t=x_t\bmod2,\qquad
s_k=\sum_{t=0}^{k-1}\epsilon_t.
\]

The one-step identity is

\[
2x_{t+1}=3^{\epsilon_t}x_t+\epsilon_t.
\]

Iterating gives

\[
\boxed{\;
2^k x_k
=
3^{s_k}n+C_k,
\;} \tag{4}
\]

where

\[
\boxed{\;
C_k=
\sum_{t=0}^{k-1}
\epsilon_t\,2^t\,3^{\,s_k-s_{t+1}}.
\;} \tag{5}
\]

Equation (5) is exactly the sum of the \(D_1\)-charges in (3), after indexing the control staircase by Collatz time.

For example, for \(n=75\), the first seven parities are

\[
(1,1,0,1,0,0,0),
\]

so \(s_7=3\), \(x_7=16\), and

\[
C_7=3^2+2\cdot3+2^3=9+6+8=23.
\]

Consequently

\[
2^7\cdot16=3^3\cdot75+23=2048.
\]

The tiling therefore encodes the complete affine arithmetic of the orbit, not merely its parity pattern.


\bigskip



\subsubsection*{4. The unrestricted periodicity claim is false}


There is a particularly simple finite-support boundary for the six arithmetic tiles alone:

\begin{itemize}
\item all vertical boundary digits are \(0\);
\item the horizontal boundary consists of one \(1\), followed by zeroes.
\end{itemize}

The resulting columns encode, in base \(3\),

\[
1,2,4,8,16,\ldots,2^j,\ldots
\]

The finite square assemblies with this property are explicitly described in Cook\textendash{}St\'erin\textendash{}Woods. Because the boundary seeds are nested and the tile rule is deterministic, the squares have a compatible infinite-quadrant limit. Its boundary has only one nonzero digit. [\hypertarget{cit:R47:5}{}\hyperlink{ref:R47:5}{5}, \hyperlink{ref:R47:4}{4}]

This limiting tiling is not eventually translation-periodic in the iteration direction.

Indeed, suppose a translation by \((p,q)\), with \(p>0\), eventually preserved the tiling. Comparing sufficiently distant columns would imply that the ternary expansion of \(2^{j+p}\) agrees, after a vertical shift by \(q\), with that of \(2^j\), apart from a bounded number of low ternary digits. Hence

\[
2^{j+p}-3^q2^j
\]

would remain bounded as \(j\to\infty\). But

\[
2^{j+p}-3^q2^j
=
2^j(2^p-3^q).
\]

It can remain bounded only if

\[
2^p=3^q,
\]

which is impossible for nonzero integers \(p,q\). A purely vertical period is also impossible because every nonzero column has a finite, nonperiodic ternary support.

Therefore:

\[
\boxed{\text{Finite-support boundary does not imply periodicity for these Wang tiles.}}
\]

This tiling represents base conversion and repeated multiplication by \(2\), not a Collatz orbit. It refutes the unrestricted rigidity lemma but is not a Collatz counterexample.


\bigskip



\subsubsection*{5. Restricting to Collatz-admissible tilings}


For \(n\in\mathbb N\), let \(\mathcal A(n)\) be the direct limit of the northeast-seeded square assemblies whose:

\begin{itemize}
\item north boundary is the binary expansion of \(n\), least significant bit first, padded by zeroes;
\item east boundary is the fixed control color \(\sigma\).
\end{itemize}

Because the eight-tile system is northeast-deterministic, these finite assemblies are compatible. The \(m\times m\) assembly computes \(T^m(n)\) on its western boundary in ternary. [\hyperlink{ref:R47:1}{1}, \hyperlink{ref:R47:4}{4}]

A parity block of length \(p\), containing \(r\) odd steps, moves the control staircase by the geometric vector

\[
p\,\mathrm{West}+r\,\mathrm{South}.
\]

This is the translation vector used in the cyclic version of the Collatz spacetime construction. [\hypertarget{cit:R47:6}{}\hyperlink{ref:R47:6}{6}]


\paragraph*{Proposition: orbit periodicity equals time-periodicity}


For a positive integer \(n\), the following are equivalent:

\[
T^{K+p}(n)=T^K(n)
\]

for some \(K\geq0,p\geq1\), and

\[
\mathcal A(n)
\]

is eventually periodic under the corresponding time-translation vector.

\textbf{Proof.}

If \(T^{K+p}(n)=T^K(n)\), determinism implies that all later iterates, parity decisions, carries and tiles repeat with period \(p\). Hence the spacetime beyond time \(K\) repeats.

Conversely, if two complete time cuts in the tiling agree, then the binary or ternary integers decoded on those cuts agree. Therefore the corresponding Collatz iterates are equal. \ensuremath{\blacksquare}

So the desired theorem


\begin{quote}
every Collatz-admissible finite-boundary tiling becomes time-periodic

\end{quote}


is exactly the assertion that every positive Collatz orbit eventually enters some cycle.

That is still weaker than the full Collatz conjecture: it would not exclude a nontrivial positive cycle.


\paragraph*{Periodic words and rational \(2\)-adic cycles}


Let \(u=(u_0,\ldots,u_{\ell-1})\) be a parity word with \(r\) ones, and let

\[
C_u=\sum_{t=0}^{\ell-1}
u_t2^t3^{\,r-s_{t+1}}.
\]

A cycle supported by the repeated word \(u^\infty\) must satisfy

\[
2^\ell z=3^r z+C_u,
\]

so

\[
\boxed{\;
z_u=\frac{C_u}{2^\ell-3^r}.
\;} \tag{6}
\]

The denominator is odd, so this always defines a rational \(2\)-adic integer. Thus \textbf{every periodic parity word already gives a periodic rational \(2\)-adic tiling}. The difficult question is whether \(z_u\) is a positive ordinary integer. The cyclic Collatz spacetime construction makes precisely this distinction through the appearance of an all-zero finite-support cut. [\hypertarget{cit:R47:7}{}\hyperlink{ref:R47:7}{7}, \hyperlink{ref:R47:6}{6}]

For \(u=(1,0)\),

\[
C_u=1,\qquad 2^2-3=1,
\]

and \(z_u=1\), giving the trivial \(1\leftrightarrow2\) cycle.

Therefore the full conjecture requires both:

\[
\begin{aligned}
&\text{Aperiodic exclusion: every finite-support orbit tiling eventually cycles;}\\
&\text{Cycle rigidity: the only positive finite-support periodic tiling is }1\leftrightarrow2.
\end{aligned}
\]

Their conjunction is precisely Collatz.


\bigskip



\subsubsection*{6. Every finite staircase is realizable}


Let

\[
\epsilon=(\epsilon_0,\ldots,\epsilon_{k-1})\in\{0,1\}^k
\]

be arbitrary. Define \(s_k\) and \(C_k\) from this word as in (5), without yet assuming that it comes from an orbit.

For an integer to realize this word, equation (4) requires

\[
3^{s_k}n+C_k\equiv0\pmod{2^k}.
\]

Since \(3^{s_k}\) is invertible modulo \(2^k\), there is exactly one solution class

\[
\boxed{\;
n\equiv-3^{-s_k}C_k\pmod{2^k}.
\;} \tag{7}
\]

Choosing a positive representative gives a genuine positive integer whose first \(k\) Collatz parities are exactly \(\epsilon\). Compatibility for every prefix follows by reducing (7) modulo the corresponding lower power of \(2\).

Hence:

\[
\boxed{\text{Every finite binary control staircase occurs in a positive-integer Collatz tiling.}}
\]

This standard parity-vector bijection is also explicit in the Collatz CQCA analysis: all finite parity vectors are feasible, with exactly one realizing residue class modulo \(2^k\). [\hyperlink{ref:R47:6}{6}]

This has an important consequence. No finite forbidden patch can certify that a parity staircase is the beginning of an impossible infinite orbit. Every finite initial patch is arithmetically legitimate.


\bigskip



\subsubsection*{7. Compactness loses finite support}


Let \(\mathbb Z_2\) denote the \(2\)-adic integers. For an infinite parity sequence

\[
\epsilon=(\epsilon_0,\epsilon_1,\ldots),
\]

the compatible congruences (7) converge \(2\)-adically to

\[
\boxed{\;
n_\epsilon
=
-\sum_{t=0}^{\infty}
\epsilon_t2^t3^{-s_{t+1}}
\in\mathbb Z_2.
\;} \tag{8}
\]

The map

\[
Q:\mathbb Z_2\longrightarrow\{0,1\}^{\mathbb N},
\qquad
Q(z)_t=T^t(z)\bmod2,
\]

is bijective, and satisfies

\[
Q(Tz)=\operatorname{shift}(Q(z)).
\]

Thus the \(2\)-adic Collatz system is conjugate to the full one-sided binary shift. Arbitrary aperiodic parity sequences consequently generate valid aperiodic \(2\)-adic Collatz tilings. [\hyperlink{ref:R47:6}{6}]

Ordinary positive integers are dense in \(\mathbb Z_2\). More concretely, for every \(z\in\mathbb Z_2\) and every \(m\), choose a positive integer

\[
n_m\equiv z\pmod{2^m}.
\]

Then \(n_m\) and \(z\) have the same first \(m\) parity decisions, and their spacetime tilings agree on every fixed finite window once \(m\) is large enough. Therefore

\[
\boxed{\;
\overline{\{\mathcal A(n):n\in\mathbb N_{>0}\}}
=
\{\mathcal A(z):z\in\mathbb Z_2\}.
\;} \tag{9}
\]

This is a decisive obstruction to a standard tiling compactness proof:

\begin{itemize}
\item finite-support boundaries are not a closed condition;
\item limits of finite-support tilings may have infinite \(2\)-adic boundaries;
\item such limits may be completely aperiodic, even if every approximating positive-integer orbit were eventually periodic.
\end{itemize}

Adding finitely many extra tile colors does not repair this by itself. If every positive-integer tiling admitted a lift to a finite local certificate layer, compactness would supply lifts to its \(2\)-adic limit points as well. A successful proof must use the eventual-zero binary tail as a genuinely global hypothesis.


\bigskip



\subsubsection*{8. The arithmetic layer has a confluent normal form}


The six arithmetic tiles can be interpreted as local commutation relations between binary and ternary digit-appending maps.

Define

\[
B_i(x)=2x+i,\qquad i\in\{0,1\},
\]

and

\[
A_j(x)=3x+j,\qquad j\in\{0,1,2\}.
\]

For a tile \((N,E,W,S)\), equation (1) gives

\[
A_E\!\circ B_N(x)
=
6x+3N+E
=
6x+2W+S
=
B_S\!\circ A_W(x).
\]

Thus every tile is the rewriting relation

\[
\boxed{\;
A_EB_N\longleftrightarrow B_SA_W.
\;} \tag{10}
\]

Orient the relation so that every \(A\)-letter is moved past every \(B\)-letter. The number of inverted \(A,B\) pairs strictly decreases, so rewriting terminates.

A word containing \(r\) binary append maps and \(s\) ternary append maps represents an affine function

\[
x\longmapsto2^r3^s x+c,
\qquad
0\leq c<2^r3^s.
\]

For any fixed ordering of the bases, mixed-radix expansion gives a unique digit word for \(c\). Therefore the sorted normal form is unique, and the rewriting system is confluent.

Consequently, every finite pure arithmetic tiling is exactly a certificate that two mixed-radix expressions represent the same affine map. There is no hidden finite-patch inconsistency analogous to the frustration used by strongly aperiodic Wang tile sets. All difficulty lies in the infinite control path and the finite-support boundary condition.


\bigskip



\subsubsection*{9. No additive local Lyapunov height exists}


A natural attempt is to assign potentials \(v_0,v_1\) to binary edge colors and \(h_0,h_1,h_2\) to ternary edge colors, and require

\[
v_N+h_E\geq h_W+v_S \tag{11}
\]

on every arithmetic tile.

Suppose these potentials take values in any ordered abelian group. Set

\[
d=v_1-v_0.
\]

From tile \(1=(0,1,0,1)\),

\[
h_1-h_0\geq d.
\]

From tile \(3=(1,0,1,1)\),

\[
h_0\geq h_1.
\]

Hence \(d\leq0\).

From tile \(2=(0,2,1,0)\),

\[
h_2\geq h_1.
\]

From tile \(4=(1,1,2,0)\),

\[
d\geq h_2-h_1.
\]

Hence \(d\geq0\). Therefore

\[
d=0,\qquad h_0=h_1=h_2.
\]

Thus also \(v_0=v_1\).

We have proved:

\[
\boxed{\text{Every one-sided additive edge potential on the six tiles is trivial.}} \tag{12}
\]

Reversing all inequalities gives the same conclusion. In particular, there is no translation-invariant local height whose curl has one sign and which could force finite-boundary stabilization.

Indeed, even the exact linear conservation law is unique: if

\[
\alpha N+\beta E+\gamma W+\delta S=0
\]

on all six tiles, direct substitution gives

\[
(\alpha,\beta,\gamma,\delta)
=
\lambda(3,1,-2,-1).
\]

So equation (1) spans the entire space of linear tilewise conservation laws.


\bigskip



\subsubsection*{10. Perturbing the \(2\)-\(3\) flux also fails}


Another natural idea is to deform the exact equality \(3N+E=2W+S\). For \(a,b>1\), define

\[
\Delta_{a,b}=bN+E-aW-S.
\]

On tiles \(0,\ldots,5\), the six values are

\[
0,\quad
0,\quad
2-a,\quad
b-a-1,\quad
b+1-2a,\quad
b+1-2a.
\]

Therefore

\[
\Delta_{a,b}\geq0\quad\text{on all tiles}
\]

holds exactly when

\[
a\leq2,\qquad b\geq a+1,
\]

and

\[
\Delta_{a,b}\leq0\quad\text{on all tiles}
\]

holds exactly when

\[
a\geq2,\qquad b\leq a+1.
\]

To preserve the critical binary\textendash{}ternary scaling, one must keep

\[
\frac{\log b}{\log a}=\frac{\log3}{\log2}.
\]

Write

\[
p=\log_2 3,\qquad b=a^p.
\]

The function

\[
f(a)=a^p-a-1
\]

is strictly increasing for \(a\geq1\), because

\[
f'(a)=pa^{p-1}-1>0,
\]

and \(f(2)=0\). Hence

\[
a<2\implies a^p<a+1,
\]

while

\[
a>2\implies a^p>a+1.
\]

Both one-sided sign conditions consequently fail away from \(a=2\). At \(a=2,b=3\), every residual is exactly zero.

Thus:

\[
\boxed{\text{There is no strict base-compatible exponential Stokes inequality.}} \tag{13}
\]

The exact \(2\)-\(3\) conservation law cannot be perturbed into a Lyapunov inequality at the critical Collatz slope.


\bigskip



\subsubsection*{11. The precise missing rigidity theorem}


Equation (8) turns the entire tiling problem into one explicit statement.

For every infinite binary word \(\epsilon\), define

\[
s_t=\epsilon_0+\cdots+\epsilon_{t-1}
\]

and

\[
\Psi(\epsilon)
=
-\sum_{t=0}^{\infty}
\epsilon_t2^t3^{-s_{t+1}}
\in\mathbb Z_2.
\]

The north boundary of the associated tiling has finite binary support exactly when

\[
\Psi(\epsilon)\in\mathbb N.
\]

The Collatz conjecture is therefore equivalent to

\[
\boxed{
\Psi(\epsilon)\in\mathbb N_{>0}
\quad\Longrightarrow\quad
\epsilon\text{ is eventually }(1,0,1,0,\ldots).
} \tag{FBR}
\]

Indeed:

\begin{itemize}
\item given \(n>0\), its parity sequence \(\epsilon=Q(n)\) satisfies \(\Psi(\epsilon)=n\);
\item eventual parity \((1,0)^\infty\) means that the orbit has reached the unique \(2\)-cycle \(1\leftrightarrow2\);
\item conversely, reaching \(1\) produces that eventual parity sequence.
\end{itemize}

So \((\mathrm{FBR})\) is the exact finite-boundary rigidity theorem required by the proposed approach.

None of the preceding tiling arguments proves \((\mathrm{FBR})\). In particular:

\begin{itemize}
\item unrestricted finite-boundary periodicity is disproved by the powers-of-two tiling;
\item every finite parity staircase is realizable;
\item ordinary-integer tilings are dense among all \(2\)-adic tilings;
\item every periodic parity word already gives a periodic rational \(2\)-adic tiling;
\item additive edge heights are necessarily trivial;
\item the exact \(2\)-\(3\) flux admits no strict critical deformation.
\end{itemize}

Therefore no valid full proof, and no aperiodic \textbf{Collatz-admissible} finite-boundary counterexample, emerges from this approach. The tiling construction itself is complete, but its remaining global statement is precisely Collatz in the form \((\mathrm{FBR})\). A successful continuation would need a genuinely nonlocal invariant that detects the eventual-zero binary boundary and is coercive against unbounded active width; it cannot be a finite forbidden-patch rule, a compactness argument, or an edge-separable height of the forms above.
\subsection*{References}
\begin{enumerate}[leftmargin=2.2em]
\item \hypertarget{ref:R47:1}{}\hyperlink{cit:R47:1}{\texttt{turn197406view1}}
\item \hypertarget{ref:R47:2}{}\hyperlink{cit:R47:2}{\texttt{turn914720view0}}
\item \hypertarget{ref:R47:3}{}\hyperlink{cit:R47:3}{\texttt{turn914720view1}}
\item \hypertarget{ref:R47:4}{}\hyperlink{cit:R47:4}{\texttt{turn392583view0}}
\item \hypertarget{ref:R47:5}{}\hyperlink{cit:R47:5}{\texttt{turn197406view2}}
\item \hypertarget{ref:R47:6}{}\hyperlink{cit:R47:6}{\texttt{turn914720view3}}
\item \hypertarget{ref:R47:7}{}\hyperlink{cit:R47:7}{\texttt{turn914720view2}}
\end{enumerate}

\clearpage
\section{Round 48: approach 43 --- A topological-entropy gap}
\subsection*{Prompt}
\begin{quote}
Complete a full proof of the Collatz conjecture using the following approach:

\begin{enumerate}
\item A topological-entropy gap
\end{enumerate}

\textbf{Hybrid \textperiodcentered{} P}

Let X\_int be the subshift of valuation words realized by positive integers and X\_crit the subshift with nonnegative long-term drift. Prove that their finite-word intersection has exponentially shrinking branching and no infinite path.

Entropy alone cannot prove emptiness, so this needs a hereditary or compactness property turning exponential scarcity into extinction.

Take as long as you need.
\end{quote}
\subsection*{Response}

\subsection*{Result}


Carrying the proposed construction through reveals an exact obstruction: \textbf{approach 43 cannot be completed as stated}, because the two objects called \(X_{\mathrm{int}}\) and \(X_{\mathrm{crit}}\) cannot simultaneously be subshifts with the claimed entropy properties.

This is not merely the familiar warning that ``entropy does not imply emptiness.'' More specifically:

\[
\boxed{\text{integer realizability has the full finite language,}}
\]

and

\[
\boxed{\text{asymptotic nonnegative drift also has the full finite language.}}
\]

Closing either set to obtain compactness therefore produces the full shift and destroys the entropy gap. Keeping the entropy gap requires a nonclosed, anchored condition, for which compactness cannot promote finite realizations to a positive-integer realization.

Current literature continues to describe the pointwise Collatz conjecture and the exclusion of nontrivial cycles as unresolved; Tao's current revision proves an almost-all result rather than a pointwise one. [\hypertarget{cit:R48:1}{}\hyperlink{ref:R48:1}{1}, \hypertarget{cit:R48:2}{}\hyperlink{ref:R48:2}{2}]

What follows is a rigorous no-go theorem for the formulation, followed by the strongest viable repair.


\bigskip



\subsubsection*{1. Shortcut and valuation codings}


Use the shortcut map

\[
T(n)=
\begin{cases}
n/2,&n\equiv0\pmod2,\\[2mm]
(3n+1)/2,&n\equiv1\pmod2.
\end{cases}
\]

Write

\[
\varepsilon_k(n)=T^k(n)\pmod2,\qquad
s_m(n)=\sum_{k=0}^{m-1}\varepsilon_k(n).
\]

The multiplicative coefficient after \(m\) shortcut steps is

\[
q_m(n)=\frac{3^{s_m(n)}}{2^m},
\]

and its logarithmic drift is

\[
\Phi_m(n)=s_m(n)\log 3-m\log 2.
\]

For the accelerated odd map

\[
U(x)=\frac{3x+1}{2^{a(x)}},
\qquad
a(x)=v_2(3x+1),
\]

a valuation word \(w=(a_0,\ldots ,a_{j-1})\) corresponds to the shortcut parity block

\[
1\,0^{a_0-1}\,1\,0^{a_1-1}\cdots
1\,0^{a_{j-1}-1}\,1.
\]

The final \(1\) ensures that every listed valuation is exact.


\bigskip



\subsection*{2. Integer realizability imposes no finite-word restriction}



\subsubsection*{Proposition 1}


Every finite binary word occurs as the initial parity word of infinitely many positive integers.


\paragraph*{Proof}


For every \(m\), the parity-vector map

\[
n\bmod 2^m
   \longmapsto
   \bigl(\varepsilon_0(n),\ldots,\varepsilon_{m-1}(n)\bigr)
\]

is a bijection from \(\mathbb Z/2^m\mathbb Z\) to \(\{0,1\}^m\). Consequently, a prescribed parity word determines exactly one residue class modulo \(2^m\), and that class contains infinitely many positive integers. This is the classical Terras\textendash{}Everett finite parity correspondence. [\hypertarget{cit:R48:3}{}\hyperlink{ref:R48:3}{3}, \hypertarget{cit:R48:4}{}\hyperlink{ref:R48:4}{4}]

The same conclusion holds for valuation words. If

\[
A_j=a_0+\cdots+a_{j-1},
\]

then the parity block displayed above has length \(A_j+1\). It therefore determines one residue class modulo \(2^{A_j+1}\), every positive representative of which realizes precisely the valuation word \(w\). \ensuremath{\blacksquare}

Hence, if

\[
X_{\mathrm{int}}
 =\{(a_0(n),a_1(n),\ldots):n\text{ a positive odd integer}\},
\]

then

\[
\mathcal L(X_{\mathrm{int}})
   =\bigcup_{j\geq0}\mathbb N_{\geq1}^{\,j}.
\]

In words:

\[
\boxed{\mathcal L(X_{\mathrm{int}})\text{ is the full valuation language.}}
\]

The set \(X_{\mathrm{int}}\) is countable and shift-invariant, but it is not closed. Its closure is the full shift on the valuation alphabet. Equivalently, in the parity representation, the positive-integer parity sequences form a dense subset of the full binary shift. The extension to \(\mathbb Z_2\) makes this exact: every infinite binary sequence is the parity sequence of exactly one \(2\)-adic integer, and the parity map is a \(2\)-adic isometric conjugacy with the binary shift. [\hypertarget{cit:R48:5}{}\hyperlink{ref:R48:5}{5}, \hypertarget{cit:R48:6}{}\hyperlink{ref:R48:6}{6}]

Thus the natural \(X_{\mathrm{int}}\) is \textbf{not a subshift} in the standard sense, while its subshift closure is the full shift.


\bigskip



\subsection*{3. Long-term criticality is also not a subshift}


Put

\[
\rho=\frac{\log 2}{\log 3}
      =0.630929753571\ldots .
\]

A natural shortcut definition of nonnegative long-term drift is

\[
X_{\mathrm{crit}}
 =
 \left\{\varepsilon\in\{0,1\}^{\mathbb N}:
 \liminf_{m\to\infty}
 \frac{s_m\log3-m\log2}{m}\geq0
 \right\}.
\]

Equivalently,

\[
\liminf_{m\to\infty}\frac{s_m}{m}\geq\rho.
\]


\subsubsection*{Proposition 2}


Every finite binary word occurs in an element of \(X_{\mathrm{crit}}\).


\paragraph*{Proof}


Given any finite word \(u\), append \(1^\infty\). The contribution of \(u\) disappears in the limiting average, and the resulting sequence has limiting drift

\[
\log3-\log2=\log(3/2)>0.
\]

Therefore every cylinder intersects \(X_{\mathrm{crit}}\). \ensuremath{\blacksquare}

Thus

\[
\mathcal L(X_{\mathrm{crit}})=\{0,1\}^{<\infty}.
\]

Moreover \(X_{\mathrm{crit}}\) is not closed. For example,

\[
0^r1^\infty\in X_{\mathrm{crit}}
\]

for every finite \(r\), but

\[
0^r1^\infty\longrightarrow 0^\infty,
\]

and \(0^\infty\) has drift \(-\log2\).

The valuation version has the same problem. Set

\[
\alpha=\log_2 3=1.584962500721\ldots .
\]

A long-term critical valuation sequence might be defined by

\[
\limsup_{j\to\infty}\frac{A_j}{j}\leq\alpha.
\]

But every finite valuation word can be extended by \(1^\infty\), whose limiting valuation average is \(1<\alpha\). Hence this set is again dense, nonclosed, and has the full finite language.

Therefore

\[
\boxed{
\overline{X_{\mathrm{int}}}
 =
\overline{X_{\mathrm{crit}}}
 =
\text{the full shift}.
}
\]

In particular, the finite-word intersection does not shrink at all:

\[
\mathcal L(X_{\mathrm{int}})
 \cap
\mathcal L(X_{\mathrm{crit}})
 =
\text{the full language}.
\]


\bigskip



\subsection*{4. Where the numerical entropy gap really occurs}


There is nevertheless a genuine numerical gap. It applies to words satisfying a drift constraint \textbf{at their endpoint}, not to the language of an asymptotic-drift subshift.

Define

\[
\mathcal C_m=
\left\{
u\in\{0,1\}^m:
|u|_1\log3-m\log2\geq0
\right\}.
\]

Then

\[
|\mathcal C_m|
 =
\sum_{s=\lceil \rho m\rceil}^{m}\binom ms.
\]

Stirling's formula gives

\[
|\mathcal C_m|
 =
\exp\bigl(mH(\rho)+O(\log m)\bigr),
\]

where

\[
H(p)=-p\log p-(1-p)\log(1-p).
\]

Numerically,

\[
H(\rho)=0.658458995327\ldots
\]

or, in bits,

\[
\frac{H(\rho)}{\log2}
 =0.949955527188\ldots .
\]

Thus the entropy deficit is

\[
g=\log2-H(\rho)
  =0.034688185232\ldots
\]

nats per shortcut step, or

\[
\frac g{\log2}
 =0.050044472812\ldots
\]

bits per shortcut step. Equivalently,

\[
\frac{|\mathcal C_m|}{2^m}
 =
\exp(-gm+O(\log m))
 \approx 0.9659065532^m
\]

up to polynomial factors.

This is a real exponential scarcity. But the \textbf{absolute} number of critical words is

\[
|\mathcal C_m|\approx 1.9318131065^m,
\]

which still grows exponentially. It is not an extinction estimate.

For accelerated valuation words, the endpoint-critical count is similarly explicit:

\[
\begin{aligned}
M_j
&=
\#\left\{(a_0,\ldots,a_{j-1})\in\mathbb N_{\geq1}^j:
\sum_i a_i\leq \lfloor\alpha j\rfloor\right\}\\
&=
\sum_{A=j}^{\lfloor\alpha j\rfloor}
   \binom{A-1}{j-1}\\
&=
\binom{\lfloor\alpha j\rfloor}{j}.
\end{aligned}
\]

Consequently,

\[
M_j
 =
\exp\left(
j\left[
\alpha\log\alpha-(\alpha-1)\log(\alpha-1)
\right]+O(\log j)
\right),
\]

with exponential base

\[
2.8395137305\ldots .
\]

Since

\[
\frac{2.8395137305\ldots}{3}
 =0.9465045768\ldots ,
\]

there is a gap relative to the \(3^j\) endpoint residue space. But Proposition 1 says that \textbf{every one of these critical finite words is realized by positive integers}. Thus intersecting with finite integer realizability removes nothing.

There are therefore only two interpretations:

\begin{enumerate}
\item ``Critical word'' means extendable to a sequence with nonnegative asymptotic drift. Then every finite word is critical and the entropy is full.

\item ``Critical word'' means its own endpoint drift is nonnegative. Then there is an entropy gap, but these words are not the language of a subshift.
\end{enumerate}

Indeed, a word \(0^r1^M\) can have positive total drift for large \(M\), while containing the arbitrarily noncritical subword \(0^r\). A subshift language must be closed under taking subwords.


\bigskip



\subsection*{5. The explicit compactness failure: the \(2\)-adic ghost \(-1\)}


The obstruction is visible in one elementary family.

For \(m\geq1\), let

\[
n_m=2^m-1.
\]

Then for \(0\leq k\leq m\),

\[
T^k(n_m)=3^k2^{m-k}-1.
\]

This follows immediately by induction:

\[
T\!\left(3^k2^{m-k}-1\right)
 =
\frac{3(3^k2^{m-k}-1)+1}{2}
 =
3^{k+1}2^{m-k-1}-1.
\]

Thus \(n_m\) has parity prefix \(1^m\), and

\[
T^m(n_m)=3^m-1.
\]

So every finite prefix of the maximally growing path

\[
1^\infty
\]

is realized by a positive integer. Nevertheless,

\[
2^m-1\longrightarrow -1
\qquad\text{in }\mathbb Z_2,
\]

and the infinite all-ones parity sequence belongs to the \(2\)-adic integer \(-1\), not to a positive integer. This correspondence is also obtained directly from the inverse parity transform. [\hyperlink{ref:R48:5}{5}, \hypertarget{cit:R48:7}{}\hyperlink{ref:R48:7}{7}]

This is exactly what compactness does to the proposed argument:

\[
\text{positive finite realizations}
\quad\stackrel{\text{compactness}}{\longrightarrow}\quad
\text{\(2\)-adic realization},
\]

not

\[
\text{positive finite realizations}
\quad\longrightarrow\quad
\text{one positive integer realizing every prefix}.
\]

For an infinite parity path \(u\), let \(r_m(u)\in[0,2^m)\) be its unique realizing residue modulo \(2^m\). The corresponding \(2\)-adic point is an ordinary nonnegative integer precisely when \(r_m(u)\) is eventually constant. For \(1^\infty\),

\[
r_m=2^m-1,
\]

so the residues never stabilize.

Hence the missing property is not compactness but

\[
\boxed{\text{eventual stabilization of the canonical binary residue.}}
\]

Entropy controls the number of possible residues. It does not control whether one nested residue sequence eventually stabilizes.


\bigskip



\subsection*{6. Natural heredity conflicts with criticality}


Suppose \(n\) has a nonconvergent orbit. Then \(2n\) also has a nonconvergent orbit, because

\[
T(2n)=n.
\]

At the symbolic level,

\[
\varepsilon(2n)=0\,\varepsilon(n).
\]

Therefore any naturally backward-hereditary bad set contains

\[
0^k\varepsilon(n)
\qquad(k\geq0).
\]

If this set is closed, it must also contain the limit

\[
0^\infty.
\]

But \(0^\infty\) has drift \(-\log2\). Thus:


\subsubsection*{Proposition 3}


There is no nonempty closed set of parity sequences that is simultaneously

\begin{itemize}
\item backward-hereditary under the guaranteed even inverse branch, and
\item contained in the uniformly nonnegative-drift set.
\end{itemize}

A still more general topological observation makes the conflict clear.


\subsubsection*{Prefix-replacement lemma}


Let \(K\subseteq\{0,1\}^{\mathbb N}\) be nonempty and closed. Suppose that whenever \(x\in K\), replacing any finite prefix of \(x\) by an arbitrary finite binary word leaves the resulting sequence in \(K\). Then

\[
K=\{0,1\}^{\mathbb N}.
\]


\paragraph*{Proof}


Fix \(x\in K\). For an arbitrary target \(z\), form

\[
z_0z_1\cdots z_{m-1}x.
\]

Every such sequence lies in \(K\), and as \(m\to\infty\) it converges to \(z\). Closedness gives \(z\in K\). \ensuremath{\blacksquare}

Long-term drift is insensitive to finite prefixes. Therefore its natural heredity, combined with closedness, forces the full shift rather than extinction.


\bigskip



\subsection*{7. Divergent orbits really would be critical\textemdash{}but cycles would not}


The drift condition is relevant to divergent nonperiodic trajectories.

Suppose a positive orbit is nonperiodic and nonconvergent. It must tend to infinity: otherwise it visits some finite interval infinitely often, and hence repeats an integer, producing a cycle.

For the shortcut map,

\[
\log T(x)-\log x=
\begin{cases}
-\log2,&x\text{ even},\\[1mm]
\log3-\log2+\log\left(1+\dfrac1{3x}\right),
 &x\text{ odd}.
\end{cases}
\]

Summing gives

\[
\log x_m-\log x_0
 =
s_m\log3-m\log2
 +
\sum_{\substack{k<m\\x_k\text{ odd}}}
 \log\left(1+\frac1{3x_k}\right).
\]

Since \(x_k\to\infty\), the final sum divided by \(m\) tends to zero. Because \(\log x_m/m\geq0\),

\[
\liminf_{m\to\infty}
\frac{s_m\log3-m\log2}{m}\geq0.
\]

Thus every divergent nonperiodic orbit would indeed produce an asymptotically critical parity sequence.

But a cycle behaves differently. For a period-\(m\) cycle with \(s\) odd steps, summing the same identity around the cycle gives

\[
0
 =
s\log3-m\log2
 +
\sum_{\substack{k<m\\x_k\text{ odd}}}
 \log\left(1+\frac1{3x_k}\right).
\]

The correction sum is strictly positive, so

\[
s\log3-m\log2<0.
\]

Therefore a multiplicative nonnegative-drift set excludes every positive cycle, including any hypothetical nontrivial one. Proving that no positive integer lies in \(X_{\mathrm{int}}\cap X_{\mathrm{crit}}\) would exclude divergence, but it would \textbf{not} exclude nontrivial cycles.

One can instead define drift using the actual height change, in which a cycle has zero drift. But actual height drift depends on the affine correction

\[
\log\left(1+\frac1{3x}\right),
\]

not on the valuation word alone. Such a system must carry archimedean height or affine-carry information and is no longer a valuation subshift.


\bigskip



\subsection*{8. The genuine averaged inverse-pressure gap}


There is one strong entropy calculation that survives the audit and is likely the useful core of a repaired argument.

For a binary word \(u\) of length \(m\), with \(s=|u|_1\), the corresponding affine iterate has the form

\[
2^m y=3^s x+C_u.
\]

Reading \(u\) as an inverse word ending at \(y\),

\[
x=\frac{2^m y-C_u}{3^s}.
\]

For fixed \(u\), integrality of \(x\) requires \(y\) to belong to exactly one residue class modulo \(3^s\). Hence, when \(y\) is averaged uniformly modulo a sufficiently high power of \(3\), the probability that \(u\) is inverse-integral is exactly \(3^{-s}\).

Let \(N_m(y)\) count inverse-integral words with nonnegative drift. Ignoring positivity, which only enlarges the count,

\[
\mathbb E_y N_m(y)
 =
\sum_{s=\lceil\rho m\rceil}^{m}
 \binom ms3^{-s}.
\]

For \(s=pm\), the exponent is

\[
H(p)-p\log3.
\]

On \(p\geq\rho\), this is maximized at \(p=\rho\), and

\[
\rho\log3=\log2.
\]

Therefore

\[
\boxed{
\mathbb E_yN_m(y)
 =
\exp\bigl(-gm+O(\log m)\bigr)
}
\]

with

\[
g=\log2-H(\rho)
 =0.034688185232\ldots>0.
\]

So the averaged number of critical inverse words really does decay like

\[
0.9659065532^m
\]

up to polynomial factors.

Equivalently, the tilted averaged transfer factor is

\[
F(t)=2^{-t}\bigl(1+3^{t-1}\bigr).
\]

Its minimum occurs at

\[
t_*
 =
1+\frac{\log(\rho/(1-\rho))}{\log3}
 =
1.4880771321\ldots ,
\]

and

\[
F(t_*)=e^{-g}=0.9659065532\ldots<1.
\]

This is the authentic topological-pressure gap.

It is, however, only an average over endpoints. It is not uniform. The family

\[
x_m=2^m-1,\qquad y_m=3^m-1
\]

provides an exceptional endpoint for the word \(1^m\) at every depth, and the \(3\)-adic limit is again \(-1\). A single orbit could in principle keep moving through exceptional endpoint classes even though those classes have exponentially shrinking density.

Thus Haar measure, entropy, Hausdorff dimension, or Borel\textendash{}Cantelli arguments cannot exclude the exceptional positive-integer trajectory by themselves.


\bigskip



\subsection*{9. The exact theorem a repaired approach must establish}


The appropriate replacement is not an ordinary entropy bound but a \textbf{uniform arithmetic negative-pressure certificate}.


\subsubsection*{Weighted extinction lemma}


Let \(\mathcal T_m\) be a finite set of arithmetic survivor states of depth \(m\), and let \(\Phi(v)\) be the accumulated logarithmic drift. Suppose that for some \(t,\delta>0\),

\[
Z_m=
\sum_{v\in\mathcal T_m} e^{t\Phi(v)}
 \leq C e^{-\delta m}.
\]

Suppose also that every hypothetical divergent orbit supplies states \(v_m\in\mathcal T_m\) with

\[
\liminf_{m\to\infty}\frac{\Phi(v_m)}m\geq0.
\]

Then no divergent orbit exists.


\paragraph*{Proof}


Choose

\[
0<\eta<\frac{\delta}{t}.
\]

For all sufficiently large \(m\),

\[
\Phi(v_m)\geq-\eta m.
\]

The single term corresponding to \(v_m\) gives

\[
Z_m\geq e^{-t\eta m}.
\]

But

\[
e^{-t\eta m}>Ce^{-\delta m}
\]

for sufficiently large \(m\), a contradiction. \ensuremath{\blacksquare}

A local block form would be even better. One would seek a positive state weight \(W\), a block length \(L\), and \(\theta<1\) such that

\[
\sum_{v\longrightarrow u\atop\text{\(L\)-step survivor}}
 e^{t(\Phi(u)-\Phi(v))}W(u)
 \leq
 \theta W(v)
\]

for every nonterminal state \(v\), while along any genuine positive-integer orbit

\[
\log W(v_m)=o(m).
\]

Iterating would give the required negative pressure.

The state cannot consist only of a valuation suffix. It must record enough information to control:

\[
\begin{array}{ll}
\text{\(2\)-adic side:}&
\text{whether canonical seed residues stabilize;}\\[1mm]
\text{\(3\)-adic side:}&
\text{endpoint congruences and exceptional proximity to }-1;\\[1mm]
\text{real side:}&
\text{height and affine carry, so that cycles and escape are distinguished.}
\end{array}
\]

The terminal-suffix locality theorem from the preceding work is valuable for making the \(3\)-adic component local. But the \(2\)-adic stabilization and archimedean nonescape conditions remain indispensable.


\bigskip



\subsection*{Final verdict}


The requested proof does not close, because the central proposed object has incompatible requirements:

\[
\boxed{
\begin{array}{c}
\text{keeping actual positive integers gives a dense, nonclosed set;}\\[1mm]
\text{closing it gives the full shift;}\\[1mm]
\text{keeping long-term drift gives a dense, nonclosed set;}\\[1mm]
\text{closing it also gives the full shift.}
\end{array}}
\]

The numerical entropy deficit

\[
0.0500444728\ldots\text{ bits per shortcut step}
\]

is genuine, but it is a large-deviation or averaged inverse-pressure gap\textemdash{}not an entropy gap for the intersection of two compact subshifts.

Compactness does not produce a positive-integer counterexample; it produces \(2\)-adic ghosts such as

\[
-1\longleftrightarrow1^\infty.
\]

Accordingly, a correct version of approach 43 must prove a \textbf{uniform weighted arithmetic-pressure inequality} that simultaneously controls binary residue stabilization, ternary endpoint carries, and real height. Establishing that inequality would indeed exclude divergent trajectories, but a separate component would still be required to exclude nontrivial cycles.
\subsection*{References}
\begin{enumerate}[leftmargin=2.2em]
\item \hypertarget{ref:R48:1}{}\hyperlink{cit:R48:1}{\texttt{turn161280view0}}
\item \hypertarget{ref:R48:2}{}\hyperlink{cit:R48:2}{\texttt{turn161280view6}}
\item \hypertarget{ref:R48:3}{}\hyperlink{cit:R48:3}{\texttt{turn161280view1}}
\item \hypertarget{ref:R48:4}{}\hyperlink{cit:R48:4}{\texttt{turn161280view5}}
\item \hypertarget{ref:R48:5}{}\hyperlink{cit:R48:5}{\texttt{turn161280view2}}
\item \hypertarget{ref:R48:6}{}\hyperlink{cit:R48:6}{\texttt{turn161280view4}}
\item \hypertarget{ref:R48:7}{}\hyperlink{cit:R48:7}{\texttt{turn161280view3}}
\end{enumerate}

\clearpage
\section{Round 49: approach 44 --- Schmidt-game avoidance in Z\_2}
\subsection*{Prompt}
\begin{quote}
Complete a full proof of the Collatz conjecture using the following approach:

\begin{enumerate}
\item Schmidt-game avoidance in Z\_2
\end{enumerate}

\textbf{Speculative \textperiodcentered{} P}

Define the 2-adic set of infinite parity sequences with nonnegative archimedean drift. Prove that it is avoidable by a winning strategy whose forbidden set contains every embedded positive integer \textemdash{} or conversely that it is winning and therefore large enough to motivate disproof searches.

Take as long as you need.
\end{quote}
\subsection*{Response}

\subsection*{Verdict}


After formalizing the proposal, it does \textbf{not} produce a full proof of Collatz. Instead, it yields a fairly sharp no-go theorem for this route:

\begin{enumerate}
\item The natural nonnegative-drift set is winning for some fixed Schmidt parameters and losing for others; neither it nor its complement is Schmidt-winning in the usual uniform sense.
\item A winning strategy produces or avoids a \textbf{game outcome}. It does not prove a universal statement about the countable dense set of parity codes arising from positive integers.
\item Every positive cycle\textemdash{}including a hypothetical nontrivial one\textemdash{}has strictly \textbf{negative} asymptotic multiplicative drift. Thus the drift set detects divergent orbits but is completely blind to the cycle half of Collatz.
\end{enumerate}

The useful outcome is an exact classification of the game and a precise identification of the missing arithmetic ingredient.


\bigskip



\subsubsection*{1. The exact 2-adic critical set}


Use the shortcut Collatz map

\[
T(x)=
\begin{cases}
x/2,&x\equiv0\pmod 2,\\[2mm]
(3x+1)/2,&x\equiv1\pmod2.
\end{cases}
\]

For \(x\in\mathbb Z_2\), put

\[
\varepsilon_j(x)=T^j(x)\pmod2,\qquad
s_k(x)=\sum_{j=0}^{k-1}\varepsilon_j(x).
\]

The parity-code map is

\[
Q(x)=\sum_{j=0}^{\infty}\varepsilon_j(x)2^j.
\]

It is a bijective 2-adic isometry and conjugates \(T\) to the one-sided binary shift. Moreover, a length-\(k\) parity word determines exactly one residue class modulo \(2^k\). [\hypertarget{cit:R49:1}{}\hyperlink{ref:R49:1}{1}]

Set

\[
\theta=\frac{\log 2}{\log 3}
      =0.6309297535\ldots
\]

and define the drift-critical set in parity space by

\[
\mathcal D
=
\left\{
y=\sum_{j\ge0}\varepsilon_j2^j:
\liminf_{k\to\infty}\frac{s_k}{k}\ge\theta
\right\}.
\]

Equivalently,

\[
\mathcal D
=
\left\{
y:
\liminf_{k\to\infty}
\frac{s_k\log3-k\log2}{k}\ge0
\right\}.
\]

This normalized lower drift is the correct pure-parity condition forced by divergence.


\bigskip



\subsubsection*{2. Divergent orbits are exactly the positive-integer points in \(\mathcal D\)}


Let \(n\) be a positive integer and \(x_j=T^j(n)\). At each step,

\[
\log x_{j+1}-\log x_j
=
\varepsilon_j\log3-\log2
+
\varepsilon_j\log\left(1+\frac1{3x_j}\right).
\]

Summing gives

\[
\log x_k-\log n
=
s_k\log3-k\log2+E_k,
\tag{1}
\]

where

\[
E_k=
\sum_{\substack{0\le j<k\\ \varepsilon_j=1}}
\log\left(1+\frac1{3x_j}\right)\ge0.
\]


\paragraph*{Proposition 1}


If the orbit of \(n\) is nonperiodic, then \(x_j\to\infty\), and consequently

\[
Q(n)\in\mathcal D.
\]


\paragraph*{Proof}


A nonperiodic orbit contains no repeated value. Since only finitely many positive integers lie below any fixed bound, it eventually leaves every bounded set. Hence \(x_j\to\infty\).

The summands defining \(E_k\) therefore tend to zero, so by Ces\`aro averaging,

\[
\frac{E_k}{k}\longrightarrow0.
\]

From (1),

\[
\frac{s_k\log3-k\log2}{k}
=
\frac{\log x_k-\log n-E_k}{k}.
\]

Because \(x_k\ge1\),

\[
\liminf_{k\to\infty}
\frac{s_k\log3-k\log2}{k}\ge0.
\]

Thus \(Q(n)\in\mathcal D\). \ensuremath{\blacksquare}

Now suppose \(n\) belongs to a positive cycle of period \(p\), containing \(s\) odd steps. Applying (1) over one period gives

\[
0=s\log3-p\log2+E_p.
\]

Every positive cycle has an odd step, so \(E_p>0\). Therefore

\[
s\log3-p\log2<0,
\qquad
\frac{s}{p}<\theta.
\tag{2}
\]

Thus \textbf{no positive periodic orbit belongs to \(\mathcal D\)}.

Let

\[
P=Q(\mathbb N_{>0})\subset\mathbb Z_2.
\]

We have proved the exact equivalence

\[
\boxed{
P\cap\mathcal D=\varnothing
\quad\Longleftrightarrow\quad
\text{there is no divergent positive Collatz orbit}.
}
\tag{3}
\]

So the proposed drift set captures the divergent-orbits conjecture exactly. It does \textbf{not} capture the nontrivial-cycles conjecture.


\bigskip



\subsubsection*{3. Size of the critical set}


Write

\[
H_2(t)=-t\log_2t-(1-t)\log_2(1-t).
\]


\paragraph*{Theorem 2}


The set \(\mathcal D\) has the following properties:

\[
\mu(\mathcal D)=0,
\]

\[
\dim_{\mathrm H}\mathcal D
=
H_2(\theta)
=
0.9499555271\ldots,
\]

and

\[
\overline{\mathcal D}
=
\overline{\mathcal D^{\,c}}
=
\mathbb Z_2.
\]

It is also meagre.


\paragraph*{Proof of density}


Every finite binary word can be followed by an all-ones tail, producing a point of \(\mathcal D\). It can also be followed by an all-zeros tail, producing a point of \(\mathcal D^c\). Therefore both sets meet every cylinder.

In particular, although \(\mathcal D\) is shift-invariant, it is not a closed subshift: its closure is the full binary shift.


\paragraph*{Proof of Haar nullity}


Under Haar measure on \(\mathbb Z_2\), the binary digits are independent fair bits, so almost surely

\[
\frac{s_k}{k}\longrightarrow\frac12<\theta.
\]

Hence \(\mu(\mathcal D)=0\).


\paragraph*{Proof of the dimension formula}


Fix \(\delta>0\). A point of \(\mathcal D\) eventually has

\[
s_k\ge(\theta-\delta)k.
\]

The number of length-\(k\) words with at least \((\theta-\delta)k\) ones is

\[
2^{\,kH_2(\theta-\delta)+o(k)}.
\]

Since a length-\(k\) cylinder has diameter \(2^{-k}\),

\[
\dim_{\mathrm H}\mathcal D
\le H_2(\theta-\delta).
\]

Letting \(\delta\downarrow0\) gives the upper bound \(H_2(\theta)\).

For the reverse inequality, put Bernoulli measure of weight \(\theta\) on the digits. Almost every sequence then has limiting one-frequency \(\theta\), and its local dimension is

\[
-\theta\log_2\theta-(1-\theta)\log_2(1-\theta)
=H_2(\theta).
\]

The mass-distribution principle gives the lower bound. \ensuremath{\blacksquare}

Thus the critical set is probabilistically negligible but has only about

\[
1-H_2(\theta)=0.0500444728\ldots
\]

of Hausdorff codimension.

At depth \(k\), there are roughly

\[
2^{0.9499555\,k}
\]

critical-looking words. This explains why long finite growth-compatible prefixes are abundant without implying that one fixed positive integer realizes an infinite one.


\bigskip



\subsubsection*{4. Exact Schmidt-game classification}


On \(\mathbb Z_2\), balls are binary cylinders. It is therefore natural first to use dyadic Schmidt parameters.

In the game \(G(a,b)\):

\begin{itemize}
\item Alice appends \(a\) binary digits on each turn.
\item Bob then appends \(b\) binary digits.
\item Their nested cylinders determine one infinite sequence.
\end{itemize}

Recent work on ordinary real binary expansions also finds parameter-dependent winning and losing regions for digit-frequency targets. In the ultrametric setting, the cylinder structure makes the classification exact. [\hypertarget{cit:R49:2}{}\hyperlink{ref:R49:2}{2}, \hypertarget{cit:R49:3}{}\hyperlink{ref:R49:3}{3}]


\paragraph*{Theorem 3}


For positive integers \(a,b\),

\[
\boxed{
\mathcal D\text{ is Alice-winning in }G(a,b)
\iff
\frac{a}{a+b}\ge\theta.
}
\tag{4}
\]

Moreover,

\[
\boxed{
\mathcal D^c\text{ is Alice-winning in }G(a,b)
\iff
\frac{b}{a+b}<\theta.
}
\tag{5}
\]


\paragraph*{Proof}


Suppose first that

\[
\frac{a}{a+b}\ge\theta.
\]

Alice writes \(1\) in all of her positions. Whatever Bob does, after \(N\) complete rounds at least \(aN\) among approximately \((a+b)N\) digits are ones. The finitely many fluctuations within a block disappear asymptotically, so

\[
\liminf\frac{s_k}{k}\ge\frac{a}{a+b}\ge\theta.
\]

Hence Alice forces \(\mathcal D\).

Conversely, if

\[
\frac{a}{a+b}<\theta,
\]

Bob writes \(0\) in all his positions. Even if every Alice digit is \(1\),

\[
\limsup\frac{s_k}{k}
\le\frac{a}{a+b}<\theta.
\]

Thus Bob prevents \(\mathcal D\). This proves (4).

For (5), Alice writes only zeros. Even if Bob writes only ones,

\[
\limsup\frac{s_k}{k}\le\frac{b}{a+b}.
\]

She therefore forces \(\mathcal D^c\) when this number is strictly below \(\theta\).

Conversely, if

\[
\frac{b}{a+b}\ge\theta,
\]

Bob writes only ones. Alice's positions can contribute zeros, but the lower one-frequency is still at least \(b/(a+b)\), so Bob forces the outcome into \(\mathcal D\). \ensuremath{\blacksquare}


\paragraph*{Corollary 4}


Neither \(\mathcal D\) nor \(\mathcal D^c\) is Schmidt-winning in the usual sense of being \(\alpha\)-winning for one fixed \(\alpha\) and every \(\beta\).


\paragraph*{Proof}


Fix Alice's contraction \(\alpha\), and set

\[
A=\left\lceil-\log_2\alpha\right\rceil.
\]

On each turn Alice fixes at most \(A\) new binary digits. Choose

\[
\beta=2^{-b}
\]

with \(b\) arbitrarily large. Bob then fixes exactly \(b\) digits per turn.

To defeat the target \(\mathcal D\), Bob writes zeros and chooses \(b\) so that

\[
\frac{A}{A+b}<\theta.
\]

To defeat the target \(\mathcal D^c\), Bob writes ones and chooses \(b\) so that

\[
\frac{b}{A+b}\ge\theta.
\]

Thus no fixed \(\alpha\) works against every \(\beta\). \ensuremath{\blacksquare}

So the proposed dichotomy\textemdash{}


\begin{quote}
either the critical set is avoidable, or it is winning\textemdash{}

\end{quote}


is false for ordinary Schmidt winning. The answer depends on the game parameters.


\bigskip



\subsubsection*{5. The positive integers are invisible to the game}


The set

\[
P=Q(\mathbb N_{>0})
\]

is countable and dense in \(\mathbb Z_2\). Density follows because every finite parity word is realized by exactly one residue class modulo \(2^k\), and hence by infinitely many positive integers. [\hyperlink{ref:R49:1}{1}]

This creates a decisive quantifier mismatch.

A Schmidt-winning statement has the form

\[
\exists\text{ Alice strategy }\sigma\;
\forall\text{ Bob plays }\tau,\qquad
z(\sigma,\tau)\in\mathcal D^c.
\tag{6}
\]

The desired Collatz-divergence statement is

\[
\forall p\in P,\qquad p\in\mathcal D^c.
\tag{7}
\]

Statement (6) does not imply (7). To make the implication, one would need

\[
\forall p\in P\quad
\exists\tau_p
\quad
z(\sigma,\tau_p)=p.
\tag{8}
\]

But Alice may exclude \(p\) on her first move.

In fact, avoiding all embedded positive integers is game-theoretically trivial.


\paragraph*{Proposition 5}


The set

\[
\mathbb Z_2\setminus P
\]

is \(1/2\)-winning.


\paragraph*{Proof}


Enumerate

\[
P=\{p_1,p_2,p_3,\ldots\}.
\]

Every 2-adic ball has exactly two child balls of half its radius. On Alice's \(j\)-th move:

\begin{itemize}
\item if the current ball does not contain \(p_j\), play arbitrarily;
\item if it contains \(p_j\), choose the child that does not contain \(p_j\).
\end{itemize}

Once \(p_j\) has been excluded, all later balls remain disjoint from it. The final outcome avoids every \(p_j\). \ensuremath{\blacksquare}

Therefore, under the natural interpretation that the ``forbidden set contains every embedded positive integer,'' such a winning strategy exists \textbf{without using Collatz at all}.

It proves only that Alice can steer the game to a noninteger 2-adic point.


\bigskip



\subsubsection*{6. The finite-prefix shadowing obstruction}


There is an even more concrete failure.

Let

\[
y\in\mathbb Z_2
\]

be any parity sequence\textemdash{}critical or otherwise\textemdash{}and put

\[
x=Q^{-1}(y)\in\mathbb Z_2.
\]

For every \(k\), choose a positive integer \(n_k\) with

\[
n_k\equiv x\pmod{2^k}.
\]

Since \(Q\) is an isometry,

\[
Q(n_k)\equiv Q(x)=y\pmod{2^k}.
\]

Thus \(n_k\) has the same first \(k\) Collatz parities as \(y\).

Consequently:


\begin{quote}
Every infinite critical game path is shadowed to arbitrary finite depth by positive integers.

\end{quote}


But the integer \(n_k\) may change with \(k\), usually escaping to infinity. Nothing in the compact space \(\mathbb Z_2\) forces the witnesses to stabilize to one fixed ordinary integer.

This is why playing the game directly on \(P\) does not repair the argument. The subspace \(P\) is not complete; its completion is all of \(\mathbb Z_2\). Nested integer cylinders can converge 2-adically to a noninteger ``ghost'' point.

The game supplies

\[
\forall k\;\exists n_k,
\]

whereas a Collatz counterexample requires

\[
\exists n\;\forall k.
\]

No Schmidt-game compactness principle interchanges these quantifiers.


\bigskip



\subsubsection*{7. Cycles are on the wrong side of the game}


For completeness, there is an exact affine version of (2).

For a parity prefix \(\varepsilon_0,\ldots,\varepsilon_{k-1}\), define

\[
C_0=0
\]

and

\[
C_{j+1}=
\begin{cases}
C_j,&\varepsilon_j=0,\\
3C_j+2^j,&\varepsilon_j=1.
\end{cases}
\]

Then

\[
2^kT^k(n)=3^{s_k}n+C_k.
\tag{9}
\]

For a positive cycle of period \(p\),

\[
(2^p-3^s)n=C_p.
\tag{10}
\]

Since \(C_p>0\),

\[
2^p>3^s,
\]

which again gives

\[
\frac{s}{p}<\theta.
\]

Thus even a perfect proof of

\[
P\cap\mathcal D=\varnothing
\]

would leave every possible nontrivial cycle untouched.

The known \(1\leftrightarrow2\) shortcut cycle has parity density \(1/2\), and any hypothetical nontrivial positive cycle would lie on the same, negative-drift side of the game.


\bigskip



\subsubsection*{8. What a repaired game would have to retain}


A nonconvergent positive orbit has a least element. Start the orbit from that least element \(n\). Then

\[
T^k(n)\ge n
\qquad\text{for every }k.
\]

Using (9), this is equivalent to

\[
\bigl(3^{s_k}-2^k\bigr)n+C_k\ge0
\qquad\text{for every }k.
\tag{11}
\]

A valid arithmetic game would therefore have to retain simultaneously:

\begin{enumerate}
\item the parity cylinder;
\item the exact carry \(C_k\);
\item a \textbf{single common positive integer} \(n\);
\item all inequalities (11).
\end{enumerate}

Ordinary Schmidt play retains only the residue of \(n\) modulo \(2^k\). It permits a different positive representative at every depth.

To express the missing condition explicitly, for a parity sequence \(y\), let

\[
r_k(y)\in\{0,1,\ldots,2^k-1\}
\]

be the least representative of

\[
Q^{-1}(y)\pmod{2^k}.
\]

Then

\[
y\in P
\quad\Longleftrightarrow\quad
r_k(y)\text{ is eventually constant and positive}.
\tag{12}
\]

Consequently the exact missing theorem is

\[
\boxed{
y\in\mathcal D
\quad\Longrightarrow\quad
r_k(y)\text{ is not eventually constant}.
}
\tag{13}
\]

By (3), statement (13) is precisely the divergent-orbits conjecture.

For the full Collatz conjecture one must additionally prove that no nontrivial periodic parity word satisfies the divisibility condition arising from (10):

\[
2^p-3^s\mid C_p.
\tag{14}
\]

Neither (13) nor (14) follows from Schmidt winning, Hausdorff dimension, entropy, density, category, or ordinary 2-adic compactness.


\bigskip



\subsection*{Final conclusion}


Approach 44 can be resolved rigorously, but negatively:

\[
\boxed{
\begin{aligned}
&\dim_{\mathrm H}\mathcal D
   =H_2(\log_3 2)
   \approx0.9499555,\\
&\mathcal D\text{ is Haar-null, meagre, and dense},\\
&\mathcal D\text{ is winning or losing according to the players' digit shares},\\
&\text{neither }\mathcal D\text{ nor }\mathcal D^c
  \text{ is Schmidt-winning uniformly},\\
&P\cap\mathcal D=\varnothing
  \text{ is exactly the divergent-orbits conjecture},\\
&\text{all positive cycles lie in }\mathcal D^c.
\end{aligned}
}
\]

The pure \(\mathbb Z_2\) Schmidt game therefore cannot complete Collatz. Its compactness creates abundant infinite critical \textbf{ghost paths}, while allowing the positive-integer witnesses at finite depths to escape archimedeanly.

A viable descendant would need to be a height-augmented or adelic arithmetic game preserving one fixed integer and the exact affine carry. Once that common-witness condition is imposed, the problem is no longer a Schmidt-game avoidance problem; it becomes the substantive Collatz arithmetic encoded in (11)\textendash{}(14).
\subsection*{References}
\begin{enumerate}[leftmargin=2.2em]
\item \hypertarget{ref:R49:1}{}\hyperlink{cit:R49:1}{\texttt{turn557195view2}}
\item \hypertarget{ref:R49:2}{}\hyperlink{cit:R49:2}{\texttt{turn444467view0}}
\item \hypertarget{ref:R49:3}{}\hyperlink{cit:R49:3}{\texttt{turn444467view1}}
\end{enumerate}

\clearpage
\section{Round 50: approach 45 --- Arithmetic fractal intersection}
\subsection*{Prompt}
\begin{quote}
Complete a full proof of the Collatz conjecture using the following approach:

\begin{enumerate}
\item Arithmetic fractal intersection
\end{enumerate}

\textbf{Speculative \textperiodcentered{} B}

Estimate the Hausdorff dimensions of the critical-drift set, the 3-adic carry-compatible set, and the image of positive integers under parity conjugacy. A sufficiently strong transversality theorem could show their triple intersection is empty.

Take as long as you need.
\end{quote}
\subsection*{Response}

\subsection*{Outcome}


Carrying the arithmetic-fractal program through does \textbf{not} produce a valid proof of the Collatz conjecture. Instead, it reveals a structural obstruction:


\begin{quote}
Ordinary Hausdorff dimension is incapable of distinguishing an empty triple intersection from one containing a single positive-integer counterexample.

\end{quote}


The relevant dimensions can be computed sharply, and the carry coding can be made exact. But the required conclusion reduces to a new \textbf{finite-height adelic transversality theorem}, not to a standard fractal-intersection theorem. That theorem remains unproved, and positive cycles require a separate divisibility argument.

This is consistent with the current status of Collatz as an open problem as of August 15, 2026. [\hypertarget{cit:R50:1}{}\hyperlink{ref:R50:1}{1}, \hypertarget{cit:R50:2}{}\hyperlink{ref:R50:2}{2}]

What follows is the complete rigorous analysis, including the strongest corrected reduction I can obtain.


\bigskip



\subsection*{1. The critical-drift parity fractal}


Use the shortcut Collatz map

\[
T(n)=
\begin{cases}
n/2,&n\equiv0\pmod2,\\[2mm]
(3n+1)/2,&n\equiv1\pmod2.
\end{cases}
\]

For \(n\in\mathbb N\), write

\[
\epsilon_j(n)=T^j(n)\pmod 2,\qquad
S_k(n)=\sum_{j=0}^{k-1}\epsilon_j(n).
\]

Let

\[
\alpha=\log_3 2=\frac{\log 2}{\log 3}
       =0.630929753571\ldots .
\]

The formal multiplicative part of \(k\) shortcut steps is

\[
\frac{3^{S_k}}{2^k}.
\]

Thus the critical odd-step frequency is \(\alpha\).

Equip \(\{0,1\}^{\mathbb N}\cong\mathbb Z_2\) with the binary ultrametric

\[
d_2(\epsilon,\epsilon')=2^{-m},
\]

where \(m\) is the first position at which the sequences differ. Define

\[
\mathcal C_\alpha
=
\left\{
\epsilon\in\{0,1\}^{\mathbb N}:
\liminf_{k\to\infty}\frac{S_k(\epsilon)}k\ge\alpha
\right\}.
\]


\subsubsection*{Theorem 1}


\[
\boxed{
\dim_H\mathcal C_\alpha
=
H_2(\alpha)
}
\]

where

\[
H_2(p)=-p\log_2p-(1-p)\log_2(1-p).
\]

Numerically,

\[
\boxed{
\dim_H\mathcal C_\alpha
=
0.949955527188\ldots
}
\]


\paragraph*{Proof: upper bound}


Fix \(\delta>0\). For all sufficiently large \(k\), any sequence in \(\mathcal C_\alpha\) belongs to a length-\(k\) cylinder whose word has at least \((\alpha-\delta)k\) ones.

The number of such words is

\[
N_k(\delta)
=
\sum_{j\ge(\alpha-\delta)k}\binom{k}{j}.
\]

Since \(\alpha-\delta>1/2\) for sufficiently small \(\delta\),

\[
N_k(\delta)
=
2^{\,k(H_2(\alpha-\delta)+o(1))}.
\]

Each cylinder has diameter \(2^{-k}\). Hence for every

\[
s>H_2(\alpha-\delta),
\]

the total \(s\)-dimensional cost is

\[
N_k(\delta)2^{-ks}
=
2^{-k(s-H_2(\alpha-\delta)+o(1))}
\longrightarrow0.
\]

Therefore

\[
\dim_H\mathcal C_\alpha
\le H_2(\alpha-\delta).
\]

Letting \(\delta\downarrow0\) gives

\[
\dim_H\mathcal C_\alpha\le H_2(\alpha).
\]


\paragraph*{Proof: lower bound}


Put the Bernoulli measure \(\mu_\alpha\) on binary sequences, with

\[
\mu_\alpha(\epsilon_j=1)=\alpha.
\]

By the strong law,

\[
\frac{S_k}{k}\longrightarrow\alpha
\]

for \(\mu_\alpha\)-almost every sequence, so \(\mu_\alpha(\mathcal C_\alpha)=1\).

For a typical sequence,

\[
-\frac1k\log_2\mu_\alpha([\epsilon_0\ldots\epsilon_{k-1}])
\longrightarrow H_2(\alpha).
\]

Thus \(\mu_\alpha\) has exact local dimension \(H_2(\alpha)\), and the mass-distribution principle gives

\[
\dim_H\mathcal C_\alpha\ge H_2(\alpha).
\]

This proves the theorem. \ensuremath{\blacksquare}


\bigskip



\subsubsection*{Why a divergent orbit lies in this set}


For an actual positive orbit \(x_j=T^j(n)\),

\[
\frac{x_{j+1}}{x_j}
=
\begin{cases}
1/2,&\epsilon_j=0,\\[1mm]
\frac32\left(1+\frac{1}{3x_j}\right),&\epsilon_j=1.
\end{cases}
\]

Consequently,

\[
\log x_k-\log n
=
S_k\log3-k\log2
+
\sum_{\substack{0\le j<k\\\epsilon_j=1}}
\log\left(1+\frac1{3x_j}\right).
\tag{1}
\]

If a positive orbit is nonperiodic but does not tend to infinity, it visits some finite set infinitely often; one integer then repeats, forcing periodicity. Therefore every nonperiodic counterexample necessarily satisfies

\[
x_k\longrightarrow\infty.
\]

In that case,

\[
\frac1k
\sum_{\substack{j<k\\\epsilon_j=1}}
\log\left(1+\frac1{3x_j}\right)
\longrightarrow0
\]

by Ces\`aro convergence. Dividing (1) by \(k\) and using \(\liminf \log x_k/k\ge0\), we obtain

\[
\liminf_{k\to\infty}\frac{S_k}{k}\ge\frac{\log2}{\log3}=\alpha.
\]

Thus:

\[
\boxed{
\text{Every divergent positive Collatz orbit has parity sequence in }
\mathcal C_\alpha.
}
\tag{2}
\]

This argument does \textbf{not} include positive cycles. Indeed, over a positive cycle of length \(k\) with \(s\) odd steps, (1) gives

\[
0=s\log3-k\log2+\text{a strictly positive term},
\]

so

\[
\frac{s}{k}<\alpha.
\tag{3}
\]

The critical-drift fractal detects divergent orbits, not nontrivial cycles.


\bigskip



\subsection*{2. Accelerated valuation words}


It is more natural to express the 3-adic carries using the odd-to-odd Syracuse map

\[
U(x)=\frac{3x+1}{2^{a(x)}},
\qquad
a(x)=v_2(3x+1),
\]

defined on positive odd integers.

For an orbit

\[
x_{j+1}=U(x_j),
\]

write

\[
a_j=v_2(3x_j+1),\qquad
A_k=a_0+\cdots+a_{k-1}.
\]

Then

\[
2^{A_k}x_k=3^k x_0+C_k,
\tag{4}
\]

where

\[
C_k
=
\sum_{j=0}^{k-1}3^{k-1-j}2^{A_j},
\qquad A_0=0.
\tag{5}
\]

Taking logarithms step by step gives

\[
\log x_k-\log x_0
=
k\log3-A_k\log2
+
\sum_{j=0}^{k-1}
\log\left(1+\frac1{3x_j}\right).
\tag{6}
\]

Set

\[
\beta=\log_2 3=\frac1\alpha
=1.584962500721\ldots .
\]

For a divergent orbit, the last average in (6) tends to zero and therefore

\[
\boxed{
\limsup_{k\to\infty}\frac{A_k}{k}\le\beta.
}
\tag{7}
\]

Define the critical valuation set

\[
\mathcal E_\beta
=
\left\{
(a_0,a_1,\ldots)\in\mathbb N_{\ge1}^{\mathbb N}:
\limsup_{k\to\infty}\frac{A_k}{k}\le\beta
\right\}.
\]

Equip valuation sequences with the metric

\[
d_3(a,b)=3^{-m},
\]

where \(m\) is their first differing coordinate. This is the correct scale for the 3-adic carry code below.


\bigskip



\subsubsection*{Theorem 2}


\[
\boxed{
\dim_H\mathcal E_\beta
=
\frac{g(\beta)}{\log3}
=
H_2(\alpha)
=
0.949955527188\ldots
}
\]

where

\[
g(\lambda)
=
\lambda\log\lambda-(\lambda-1)\log(\lambda-1).
\]


\paragraph*{Proof: counting compositions}


The number of positive integer words

\[
(a_0,\ldots,a_{k-1})
\]

with total \(A_k=A\) is

\[
\binom{A-1}{k-1}.
\]

Therefore the number with \(A_k\le m\) is

\[
\sum_{A=k}^{m}\binom{A-1}{k-1}
=
\binom{m}{k}.
\tag{8}
\]

For \(m=\lfloor\lambda k\rfloor\), Stirling's formula gives

\[
\frac1k\log\binom{\lambda k}{k}
\longrightarrow
g(\lambda).
\tag{9}
\]

Covering \(\mathcal E_\beta\) by cylinders with total at most \((\beta+\delta)k\) gives

\[
\dim_H\mathcal E_\beta
\le
\frac{g(\beta)}{\log3}.
\]

For the lower bound, use the geometric distribution

\[
\mathbb P(a_j=m)
=
\alpha(1-\alpha)^{m-1}.
\tag{10}
\]

Its mean is \(1/\alpha=\beta\), and its entropy is

\[
-\sum_{m\ge1}\alpha(1-\alpha)^{m-1}
\log\!\left(\alpha(1-\alpha)^{m-1}\right)
=
g(\beta).
\]

The strong law places almost every such sequence in \(\mathcal E_\beta\), while the cylinder local dimension is \(g(\beta)/\log3\). This establishes the lower bound.

Finally,

\[
\frac{g(\beta)}{\log3}
=
\frac{H(\alpha)}{\log2}
=
H_2(\alpha),
\]

where \(H\) denotes natural-log binary entropy. \ensuremath{\blacksquare}


\bigskip



\subsection*{3. The exact 3-adic carry fractal}


For \(a=(a_0,a_1,\ldots)\), define

\[
\Psi(a)
=
\sum_{r=0}^{\infty}
3^r2^{-(a_0+\cdots+a_r)}
\in\mathbb Z_3.
\tag{11}
\]

It satisfies the self-similar recursion

\[
\boxed{
\Psi(a)
=
2^{-a_0}\bigl(1+3\Psi(\sigma a)\bigr).
}
\tag{12}
\]

This is the same 3-adic random-series structure that occurs in Tao's Syracuse random variables and their fine-scale mixing theory. Tao's estimates concern the unconditioned geometric valuation law; the critical lower-tail constraint \(A_k/k\le\beta\) is substantially thinner. [\hypertarget{cit:R50:3}{}\hyperlink{ref:R50:3}{3}, \hypertarget{cit:R50:4}{}\hyperlink{ref:R50:4}{4}]

Define

\[
\mathcal K_\beta=\Psi(\mathcal E_\beta)\subset\mathbb Z_3.
\]

This is a natural candidate for the ``3-adic carry-compatible'' fractal.


\bigskip



\subsubsection*{Theorem 3}


\[
\boxed{
\frac{\log2}{\log3}
\le
\dim_H\mathcal K_\beta
\le
H_2(\alpha).
}
\]

Equivalently,

\[
\boxed{
0.630929753571\ldots
\le
\dim_H\mathcal K_\beta
\le
0.949955527188\ldots .
}
\]


\paragraph*{Upper bound}


If two valuation sequences agree in their first \(m\) entries, the corresponding sums (11) differ by a multiple of \(3^m\). Thus

\[
|\Psi(a)-\Psi(b)|_3\le3^{-m}=d_3(a,b).
\]

Hence \(\Psi\) is \(1\)-Lipschitz, so

\[
\dim_H\mathcal K_\beta
\le
\dim_H\mathcal E_\beta
=
H_2(\alpha).
\]


\paragraph*{Lower bound}


Restrict to valuation digits \(a_j\in\{1,2\}\). If two such sequences first differ at position \(m\), then after removing the common prefix the difference between their two first-level branches is

\[
\frac12(1+3u)-\frac14(1+3v)
=
\frac14+3z,
\]

which is a 3-adic unit. Therefore their images under \(\Psi\) differ by exactly \(3^{-m}\): the coding is isometric on \(\{1,2\}^{\mathbb N}\).

Under the fair Bernoulli measure on \(\{1,2\}\),

\[
\frac{A_k}{k}\longrightarrow\frac32<\beta.
\]

Hence almost every such sequence belongs to \(\mathcal E_\beta\). Its image has dimension

\[
\frac{\log2}{\log3}.
\]

This proves the lower bound. \ensuremath{\blacksquare}

The precise value of \(\dim_H\mathcal K_\beta\) is not needed for the main obstruction. It is already large enough that \(\mathcal K_\beta\) cannot be regarded as a negligible carry set.


\bigskip



\subsection*{4. The positive integers under parity conjugacy}


Define the parity conjugacy

\[
Q(x)
=
\sum_{j=0}^{\infty}
\bigl(T^j(x)\bmod2\bigr)2^j,
\qquad x\in\mathbb Z_2.
\tag{13}
\]

The map \(Q:\mathbb Z_2\to\mathbb Z_2\) is a bijective isometry conjugating the shortcut Collatz map to the binary shift. In particular, each finite parity word corresponds to one residue class modulo \(2^k\). [\hypertarget{cit:R50:5}{}\hyperlink{ref:R50:5}{5}]

Let

\[
\mathcal P=Q(\mathbb N)\subset\mathbb Z_2.
\]

Then:

\[
\boxed{\dim_H\mathcal P=0}
\]

because \(\mathcal P\) is countable.

But \(\mathcal P\) is also dense. Indeed, every binary cylinder contains \(Q(n)\) for infinitely many positive integers \(n\). Consequently,

\[
\boxed{
\underline{\dim}_{B}\mathcal P
=
\overline{\dim}_{B}\mathcal P
=
1.
}
\]

At level \(k\), all \(2^k\) binary cylinders meet \(\mathcal P\).

Thus \(\mathcal P\) simultaneously has

\[
\text{Hausdorff dimension }0
\quad\text{and}\quad
\text{full box dimension }1.
\tag{14}
\]

This discrepancy is central.


\bigskip



\subsubsection*{The natural adelic graph is also dense}


Since the critical future lies in \(\mathbb Z_2\) and the carry information lies in \(\mathbb Z_3\), the literal ``triple intersection'' is not initially defined in one space.

A natural common ambient space is

\[
\mathbb Z_2\times\mathbb Z_3,
\]

using the positive-integer graph

\[
\mathcal D_+
=
\{(Q(n),n):n\in\mathbb N\}.
\tag{15}
\]

This set is countable, so

\[
\dim_H\mathcal D_+=0.
\]

Nevertheless it is dense in \(\mathbb Z_2\times\mathbb Z_3\). To see this, prescribe

\[
Q(n)\equiv q\pmod{2^k},
\qquad
n\equiv y\pmod{3^r}.
\]

The first condition is equivalent, through the isometry \(Q\), to one residue condition

\[
n\equiv c_q\pmod{2^k}.
\]

The Chinese remainder theorem combines this with the prescribed residue modulo \(3^r\), producing a positive integer \(n\).

Therefore every product cylinder meets \(\mathcal D_+\).

With the standard maximum product metric, it has full box dimension:

\[
\boxed{
\dim_H\mathcal D_+=0,
\qquad
\dim_B\mathcal D_+=2.
}
\tag{16}
\]


\bigskip



\subsection*{5. Why Hausdorff transversality cannot prove emptiness}


The proposed strategy now encounters four separate obstructions.


\subsubsection*{Obstruction 1: dimension zero does not mean empty}


Any intersection with \(\mathcal P\) or \(\mathcal D_+\) is countable. Hence

\[
\dim_H\bigl((\mathcal C_\alpha\times\mathcal K_\beta)
\cap\mathcal D_+\bigr)=0
\]

whether the intersection is empty, contains one divergent counterexample, or contains infinitely many counterexamples.

Adding or deleting one point leaves every Hausdorff dimension unchanged.

Thus a conclusion of the form

\[
\dim_H(\text{triple intersection})=0
\]

contains no information relevant to Collatz.


\bigskip



\subsubsection*{Obstruction 2: generic translations are irrelevant}


A Marstrand-type theorem may imply that for almost every translation \(t\),

\[
(\mathcal C_\alpha\times\mathcal K_\beta)
\cap(\mathcal D_++t)=\varnothing.
\]

But Collatz corresponds to one fixed arithmetic placement, namely \(t=0\).

A generic theorem is allowed to have a measure-zero exceptional set of translations. The arithmetic placement can lie in precisely that exceptional set.

Indeed, \(\mathcal D_+\) meets every finite product cylinder, so no argument based only on finite-scale topological separation can rule out \(t=0\).


\bigskip



\subsubsection*{Obstruction 3: the two main fractals are not dimensionally transverse}


We have

\[
\dim_H\mathcal C_\alpha
=
0.949955527188\ldots
\]

and

\[
\dim_H\mathcal K_\beta
\ge
0.630929753571\ldots .
\]

Thus

\[
\dim_H\mathcal C_\alpha
+
\dim_H\mathcal K_\beta
>
1.5808.
\]

If one artificially identified the two local coordinates in a one-dimensional ambient space, the usual heuristic would predict a large intersection, not emptiness.

Even using the upper estimate

\[
\dim_H\mathcal K_\beta\le0.949955527188\ldots,
\]

the critical and carry sets are individually very large.


\bigskip



\subsubsection*{Obstruction 4: dimensions do not see the sign of the affine term}


Consider the accelerated \(3x-1\) map

\[
U_-(x)=\frac{3x-1}{2^{v_2(3x-1)}}.
\]

It has the positive cycle

\[
5\longmapsto7\longmapsto5,
\]

because

\[
\frac{3\cdot5-1}{2}=7,
\qquad
\frac{3\cdot7-1}{4}=5.
\]

Its valuation word is \((1,2)\), whose mean is

\[
\frac{1+2}{2}=\frac32<\beta.
\]

Replacing \(+1\) by \(-1\) merely changes the signs of the affine carry translations. Multiplication by \(-1\) is a 3-adic isometry, so all the Hausdorff-dimension calculations above remain unchanged.

Therefore any argument depending only on:

\begin{itemize}
\item contraction ratios,
\item valuation frequencies,
\item local entropy,
\item 2-adic or 3-adic Hausdorff dimension,
\end{itemize}

would give the same conclusion for \(3x+1\) and \(3x-1\), even though the latter has this nontrivial positive cycle.

The missing information is the joint \textbf{Archimedean sign and height} of the same integer representative.


\bigskip



\subsection*{6. Exact terminal-suffix carry locality}


The 3-adic carry relation nevertheless gives a useful exact theorem.

Let

\[
v=(a_{k-r},\ldots,a_{k-1})
\]

be the final \(r\) valuations of a length-\(k\) orbit segment. Put

\[
B(v)=a_{k-r}+\cdots+a_{k-1}.
\]

Let \(C_v\) be the affine constant for this block, so that

\[
2^{B(v)}x_k
=
3^r x_{k-r}+C_v.
\tag{17}
\]

Reducing modulo \(3^r\),

\[
x_k
\equiv
2^{-B(v)}C_v
\pmod{3^r}.
\tag{18}
\]

Expanding \(C_v\) gives

\[
\boxed{
x_k
\equiv
\sum_{t=0}^{r-1}
3^t
2^{-(a_{k-1}+\cdots+a_{k-1-t})}
\pmod{3^r}.
}
\tag{19}
\]

Hence \(x_k\bmod3^r\) depends only on the last \(r\) valuation exponents.

In particular,

\[
\boxed{
3^r\mid x_k+1
\iff
3^r\mid C_v+2^{B(v)}.
}
\tag{20}
\]

This is precisely the finite terminal-suffix locality statement underlying the proposed carry fractal.


\bigskip



\subsection*{7. The orientation obstruction}


Equation (19) has an important directional feature.

At time \(k\), the endpoint residue is coded by the reversed word

\[
(a_{k-1},a_{k-2},\ldots,a_0).
\]

At time \(k+1\), the relevant word is

\[
(a_k,a_{k-1},\ldots,a_0).
\]

The new digit is \textbf{prepended}, not appended. Consequently these words do not form nested prefixes in the coding (11). There is generally no single point

\[
\Psi(a_0,a_1,\ldots)
\]

whose successive 3-adic truncations are the actual endpoint residues.

This invalidates the naive static interpretation in which one divergent orbit automatically produces one point of \(\mathcal K_\beta\).

The forward estimate

\[
A_k\le(\beta+o(1))k
\]

controls the total of the first \(k\) exponents. It does not imply that all reversed subprefixes satisfy the same barrier.


\bigskip



\subsubsection*{A partial repair using record minima}


Define

\[
D_k=A_k-\beta k.
\]

Suppose \(k\) is a record-minimum time:

\[
D_k\le D_j
\qquad(0\le j\le k).
\]

Then for every \(1\le r\le k\),

\[
D_k\le D_{k-r},
\]

and hence

\[
A_k-A_{k-r}\le\beta r.
\tag{21}
\]

Thus the entire reversed valuation word at time \(k\) satisfies the prefix barrier

\[
a_{k-1}+\cdots+a_{k-r}\le\beta r
\qquad(1\le r\le k).
\tag{22}
\]

Moreover, from (17),

\[
x_k
=
\frac{3^r}{2^{A_k-A_{k-r}}}x_{k-r}
+
\frac{C_v}{2^{A_k-A_{k-r}}}
>
x_{k-r},
\]

because (21) implies

\[
3^r\ge2^{A_k-A_{k-r}}.
\]

Therefore every record-minimum time for \(D_k\) is a record-maximum time for \(x_k\).

This is a genuine synchronization between Archimedean growth and reversed 3-adic carry words.

However, a divergent orbit need not have infinitely many record minima of \(D_k\): it is logically possible that \(D_k\) is eventually bounded below while the positive correction sum in (6) diverges slowly. So this does not close the argument.


\bigskip



\subsection*{8. Canonical finite-height representatives}


The correct replacement for ordinary fractal transversality is obtained by retaining the actual Archimedean representatives of each symbolic cylinder.

Let

\[
w=(a_0,\ldots,a_{k-1})
\]

be any positive valuation word, with total \(A=A_k\) and affine constant \(C_w\).

Define \(r_3(w)\) to be the unique odd integer satisfying

\[
1\le r_3(w)<2\cdot3^k
\]

and

\[
2^A r_3(w)\equiv C_w\pmod{3^k}.
\tag{23}
\]

Now define

\[
r_2(w)
=
\frac{2^A r_3(w)-C_w}{3^k}.
\tag{24}
\]


\subsubsection*{Proposition 4}


For every valuation word \(w\):

\begin{enumerate}
\item \(r_2(w)\) is a positive odd integer;
\item \(r_2(w)<2^{A+1}\);
\item \(r_2(w)\) realizes exactly the valuation word \(w\);
\item its endpoint after those \(k\) odd steps is \(r_3(w)\);
\item all positive realizations of \(w\) are exactly
\end{enumerate}

\[
x_0=r_2(w)+t2^{A+1},
\qquad
x_k=r_3(w)+t(2\cdot3^k),
\qquad t\ge0.
\tag{25}
\]


\paragraph*{Proof}


The congruence (23) makes (24) integral.

Backward recursion from \(r_3(w)\),

\[
y_i=\frac{2^{a_i}y_{i+1}-1}{3},
\]

is integral at every stage because the full congruence modulo \(3^k\) implies the corresponding suffix congruence modulo \(3^{k-i}\). Each \(y_i\) is positive and odd. Thus \(y_0=r_2(w)\) realizes \(w\).

To prove the upper bound, replace \(r_3(w)\) by

\[
r_3(w)-2\cdot3^k<0.
\]

The corresponding initial representative is \(r_2(w)-2^{A+1}\). From the affine identity,

\[
3^k\bigl(r_2(w)-2^{A+1}\bigr)
=
2^A\bigl(r_3(w)-2\cdot3^k\bigr)-C_w<0.
\]

Hence

\[
r_2(w)<2^{A+1}.
\]

Adding \(2\cdot3^k\) to an admissible endpoint adds \(2^{A+1}\) to the initial value, giving (25). \ensuremath{\blacksquare}


\bigskip



\subsubsection*{Consequence for one fixed counterexample}


Suppose \(n\) is an actual odd seed and \(w_k\) is its first \(k\) valuation exponents.

Since \(A_k\ge k\), for all sufficiently large \(k\),

\[
n<2^{A_k+1}.
\]

The seed \(n\) is then already the least positive representative in its word cylinder, so

\[
\boxed{
r_2(w_k)=n,\qquad r_3(w_k)=x_k.
}
\tag{26}
\]

Thus the canonical coordinates are eventually not artificial: they are exactly the seed and the endpoint of the orbit.


\bigskip



\subsection*{9. The small adelic corner forced by a counterexample}


For every odd Syracuse step,

\[
x_{j+1}
=
\frac{3x_j+1}{2^{a_j}}
\le
\frac{3x_j+1}{2}.
\]

Hence

\[
x_{j+1}+1\le\frac32(x_j+1),
\]

and therefore

\[
x_k+1\le
\left(\frac32\right)^k(n+1).
\tag{27}
\]

Let \(\eta>0\). A divergent counterexample would satisfy, for all sufficiently large \(k\),

\[
A_k\le(\beta+\eta)k,
\tag{28}
\]

\[
r_2(w_k)=n\le e^{\eta k},
\tag{29}
\]

and, by (27),

\[
r_3(w_k)=x_k
\le
\left(\frac32\right)^ke^{\eta k}.
\tag{30}
\]

Define the level-\(k\) small-corner set

\[
\mathscr S_k(\eta)
=
\left\{
\begin{array}{l}
w\in\mathbb N_{\ge1}^k:\\[1mm]
A(w)\le(\beta+\eta)k,\\
r_2(w)\le e^{\eta k},\\
r_3(w)\le(3/2)^ke^{\eta k}
\end{array}
\right\}.
\tag{31}
\]

Then:

\[
\boxed{
\text{Every divergent counterexample supplies a nested word }
w_k\in\mathscr S_k(\eta)
\text{ for every sufficiently large }k.
}
\tag{32}
\]

This is the rigorous finite-height formulation of arithmetic transversality.


\bigskip



\subsection*{10. The exact missing theorem}


The required result is not a Hausdorff-dimension theorem. It is the following deterministic statement.


\subsubsection*{Adelic small-corner extinction theorem}


There exists \(\eta>0\) such that no infinite valuation sequence has all sufficiently long prefixes in \(\mathscr S_k(\eta)\).

A stronger version would assert

\[
\boxed{
\mathscr S_k(\eta)=\varnothing
\quad\text{for all sufficiently large }k.
}
\tag{AC}
\]

Either version would exclude divergent positive orbits.

The crucial content is the simultaneous smallness of:

\[
r_2(w)\quad\text{inside a range of size }2^{A(w)}
\]

and

\[
r_3(w)\quad\text{inside a range of size }3^k.
\]

Ordinary dimensions discard both canonical heights.


\bigskip



\subsection*{11. Why the conjectured transversality is heuristically strong}


For total valuation

\[
A\sim\lambda k,
\]

the number of words is approximately

\[
\exp(g(\lambda)k).
\]

If the canonical representatives behaved as independent uniform points in their natural ranges, the probability that

\[
r_2(w)\le e^{\eta k}
\]

would be approximately

\[
\frac{e^{\eta k}}{2^A},
\]

while the probability that

\[
r_3(w)\le(3/2)^ke^{\eta k}
\]

would be approximately

\[
\frac{(3/2)^ke^{\eta k}}{3^k}
=
2^{-k}e^{\eta k}.
\]

The expected number of survivors with \(A\sim\lambda k\) would therefore have exponential rate

\[
g(\lambda)-\lambda\log2-\log2+2\eta.
\tag{33}
\]

At the critical endpoint \(\lambda=\beta\),

\[
\beta\log2=\log3,
\]

so the rate at \(\eta=0\) is

\[
g(\beta)-\log6.
\]

Numerically,

\[
g(\beta)=1.043632815857\ldots ,
\]

and

\[
\boxed{
\log6-g(\beta)
=
0.748126653371\ldots>0.
}
\tag{34}
\]

Thus the independent model predicts exponentially fast extinction.

By comparison, using the initial coordinate alone gives only

\[
\log3-g(\beta)
=
0.054979472811\ldots .
\]

The endpoint condition supplies a much larger entropy gap.


\bigskip



\subsubsection*{A rigorous generic version}


One can make this heuristic into a theorem after independently Haar-translating the two canonical coordinates at every scale.

For each \(k,A\), randomly translate the \(r_2\)-coordinates in the cyclic group of odd residues modulo \(2^{A+1}\), and independently translate the \(r_3\)-coordinates among odd residues modulo \(2\cdot3^k\).

The expected number of translated words satisfying the three inequalities in (31) is bounded by

\[
\sum_{A\le(\beta+\eta)k}
\binom{A-1}{k-1}
O\!\left(\frac{e^{\eta k}}{2^A}\right)
O\!\left(2^{-k}e^{\eta k}\right).
\tag{35}
\]

For sufficiently small fixed \(\eta>0\), (34) implies that this is \(O(e^{-ck})\) for some \(c>0\). The sum over \(k\) converges, so Borel\textendash{}Cantelli gives eventual extinction almost surely.

Therefore:

\[
\boxed{
\text{The desired small-corner extinction holds for generic independent
translations.}
}
\]

But the Collatz system uses the fixed, correlated, zero translation. The affine identity

\[
2^A r_3(w)=3^k r_2(w)+C_w
\]

couples the two coordinates exactly. Proving that this deterministic correlation still behaves transversely is the whole missing problem.


\bigskip



\subsection*{12. Cycles remain a separate obstruction}


Even a proof of \((AC)\) would only exclude divergent orbits.

A positive odd cycle with valuation word \(w\) satisfies

\[
x_k=x_0=n.
\]

Using (4),

\[
2^A n=3^k n+C_w,
\]

so

\[
\boxed{
(2^A-3^k)n=C_w.
}
\tag{36}
\]

Since \(C_w>0\),

\[
2^A>3^k,
\qquad
\frac Ak>\beta.
\]

Thus cycle words lie on the opposite side of the drift threshold and are invisible to \(\mathcal E_\beta\).

To finish Collatz after proving adelic small-corner extinction, one would additionally need:


\subsubsection*{Cycle divisibility theorem}


For every positive valuation word \(w\), except repetitions of the trivial word \((2)\),

\[
\boxed{
2^{A(w)}-3^{|w|}
\nmid C_w.
}
\tag{CYC}
\]

For \(w=(2)\),

\[
A=2,\qquad C_w=1,
\qquad
2^2-3=1,
\]

which yields the fixed odd orbit \(1\mapsto1\).

No Hausdorff-dimension argument proves \((CYC)\), because it is an exact divisibility assertion about a single finite word.


\bigskip



\subsection*{13. Precise conditional reduction}


Combining the analysis gives a clean theorem.


\subsubsection*{Conditional theorem}


Assume:

\begin{enumerate}
\item the adelic small-corner extinction theorem \((AC)\);
\item the cycle divisibility theorem \((CYC)\).
\end{enumerate}

Then every positive integer reaches \(1\) under Collatz iteration.


\paragraph*{Proof}


Suppose a positive orbit never reaches \(1\).

If it is eventually periodic, it gives a nontrivial positive valuation cycle, contradicting \((CYC)\).

Otherwise it cannot visit any bounded set infinitely often, since repetition of an integer would force periodicity. Hence it tends to infinity. Equations (28)\textendash{}(30) then produce an infinite nested sequence

\[
w_k\in\mathscr S_k(\eta)
\]

for all sufficiently large \(k\), contradicting \((AC)\).

Therefore every positive orbit reaches \(1\). \ensuremath{\blacksquare}


\bigskip



\subsection*{Final assessment}


The arithmetic-fractal approach yields three sharp facts:

\[
\boxed{
\dim_H\mathcal C_\alpha
=
H_2(\log_3 2)
=
0.949955527188\ldots
}
\]

\[
\boxed{
0.630929753571\ldots
\le
\dim_H\mathcal K_\beta
\le
0.949955527188\ldots
}
\]

and

\[
\boxed{
\dim_HQ(\mathbb N)=0,
\qquad
\dim_BQ(\mathbb N)=1.
}
\]

These estimates do \textbf{not} imply empty intersection. The image of the positive integers is countable but dense, and Hausdorff dimension cannot detect whether one exceptional arithmetic point survives.

The valid replacement is not ordinary fractal transversality but the deterministic statement that the canonical pair

\[
\bigl(r_2(w),r_3(w)\bigr)
\]

cannot repeatedly enter the exponentially small Archimedean corner

\[
r_2\le e^{\eta k},
\qquad
r_3\le(3/2)^ke^{\eta k},
\qquad
A\le(\beta+\eta)k.
\]

There is a strong entropy margin of

\[
\log6-g(\beta)=0.748126653371\ldots,
\]

and generic randomized translations satisfy the required extinction. What is missing is a deterministic anti-concentration or discrepancy theorem for the exact affine carries \(C_w\). Positive cycles additionally require the independent divisibility theorem \((CYC)\).

Accordingly, this route has been reduced to two explicit arithmetic statements, but it has not produced a verified full proof of the Collatz conjecture.
\subsection*{References}
\begin{enumerate}[leftmargin=2.2em]
\item \hypertarget{ref:R50:1}{}\hyperlink{cit:R50:1}{\texttt{turn737387search4}}
\item \hypertarget{ref:R50:2}{}\hyperlink{cit:R50:2}{\texttt{turn285695search4}}
\item \hypertarget{ref:R50:3}{}\hyperlink{cit:R50:3}{\texttt{turn285695search1}}
\item \hypertarget{ref:R50:4}{}\hyperlink{cit:R50:4}{\texttt{turn285695search7}}
\item \hypertarget{ref:R50:5}{}\hyperlink{cit:R50:5}{\texttt{turn737387search3}}
\end{enumerate}

\clearpage
\section{Round 51: approach 46 --- An adelic product-formula Lyapunov height}
\subsection*{Prompt}
\begin{quote}
Complete a full proof of the Collatz conjecture using the following approach:

\begin{enumerate}
\item An adelic product-formula Lyapunov height
\end{enumerate}

\textbf{Hybrid \textperiodcentered{} P}

Define local escape energies H\_infinity, H\_2, H\_3 whose changes under one accelerated step can be calculated exactly. Seek positive weights such that the product formula forces their total to decrease outside the trivial cycle.

Take as long as you need.
\end{quote}
\subsection*{Response}

\subsection*{Result}


Carrying the adelic calculation through does \textbf{not} produce a proof of Collatz. It produces a rigorous obstruction:


\begin{quote}
\textbf{A fixed one-step Lyapunov height assembled from standard local height functions at \(\infty,2,3\), with fixed weights, cannot prove the Collatz conjecture.}

\end{quote}


This is stronger than the failure of one particular weighting. There are explicit positive-integer orbit segments that gain an arbitrarily large archimedean factor while becoming arbitrarily indistinguishable at both the \(2\)-adic and \(3\)-adic places. Any successful repair must therefore use an unbounded itinerary/carry state, an adaptive block time, or information from primes \(p\ge 5\).

Below is a complete proof of that obstruction.


\bigskip



\subsubsection*{1. The accelerated map and its exact local changes}


Work with the Syracuse map on positive odd integers,

\[
T(n)=\frac{3n+1}{2^{a(n)}},
\qquad
a(n):=v_2(3n+1).
\]

The Collatz conjecture is equivalent to the assertion that every positive odd integer eventually reaches the fixed point \(1\).

The archimedean change is exactly

\[
\Delta_\infty(n)
=
\log T(n)-\log n
=
\log\left(3+\frac1n\right)-a(n)\log 2.
\tag{1}
\]

Two natural shifted local quantities are

\[
b(n):=v_2(n+1),
\qquad
r(n):=v_3(n+1).
\]

Since

\[
T(n)+1
=
\frac{3n+1+2^{a(n)}}{2^{a(n)}},
\]

their changes are

\[
b(Tn)
=
v_2\!\left(3n+1+2^{a(n)}\right)-a(n),
\tag{2}
\]

and

\[
r(Tn)
=
v_3\!\left(3n+1+2^{a(n)}\right).
\tag{3}
\]

In particular, whenever \(a(n)=1\),

\[
T(n)+1=\frac{3(n+1)}2.
\]

Consequently,

\[
b(Tn)=b(n)-1,
\qquad
r(Tn)=r(n)+1.
\tag{4}
\]

Writing

\[
n+1=2^{b(n)}3^{r(n)}u(n),
\qquad \gcd(u(n),6)=1,
\]

we also get

\[
u(Tn)=u(n)
\qquad\text{whenever }a(n)=1.
\tag{5}
\]

Thus a long \(a=1\) excursion simply transports one unit of valuation at each step from the \(2\)-adic side to the \(3\)-adic side while multiplying \(n+1\) by \(3/2\).


\bigskip



\subsubsection*{2. What the product formula actually says}


Normalize the finite absolute values by \(|p|_p=p^{-1}\). For the formal branch multiplier

\[
q_a=\frac{3}{2^a},
\]

one has

\[
\log |q_a|_\infty=\log 3-a\log 2,
\]

\[
\log |q_a|_2=a\log 2,
\]

\[
\log |q_a|_3=-\log 3,
\]

and hence

\[
\log |q_a|_\infty+\log |q_a|_2+\log |q_a|_3=0.
\tag{6}
\]

So equal product-formula weights give \textbf{conservation}, not descent. With arbitrary fixed weights,

\[
\lambda_\infty\log |q_a|_\infty
+\lambda_2\log |q_a|_2
+\lambda_3\log |q_a|_3
\]

equals

\[
(\lambda_\infty-\lambda_3)\log 3
+
(\lambda_2-\lambda_\infty)a\log 2.
\tag{7}
\]

Any negativity in (7) comes from manually biasing the places; it is not forced by the product formula.

More seriously, \(q_a\) is only the linear part of the branch. The actual multiplicative change is

\[
\frac{T(n)}n=\frac{3n+1}{2^{a(n)}n}.
\]

After one accelerated step, every orbit point is coprime to \(3\). Hence, on the invariant set \(\gcd(n,6)=1\), both \(n\) and \(T(n)\) are units in \(\mathbb Z_2\) and \(\mathbb Z_3\). Therefore

\[
\log\left|\frac{T(n)}n\right|_2
=
\log\left|\frac{T(n)}n\right|_3
=
0.
\]

The full product formula becomes

\[
\boxed{
\log\frac{T(n)}n
=
-\sum_{p\ge5}
\log\left|\frac{T(n)}n\right|_p
}.
\tag{8}
\]

Thus the entire archimedean change of the actual integer state is balanced by primes \(p\ge5\), not by \(2\) and \(3\).

Likewise,

\[
\log(n+1)-b(n)\log2-r(n)\log3
=
\sum_{p\ge5}v_p(n+1)\log p.
\tag{9}
\]

The right side is an independent ``outside-prime'' state. Removing it does not leave a closed three-place system.


\bigskip



\subsubsection*{3. A first no-go theorem: ordinary heights cannot decrease}


Suppose that \(V\) is a logarithmic height of the usual kind:

\[
V(n)=\gamma\log n+O(1),
\qquad \gamma>0.
\tag{10}
\]

This includes, on positive integers, conventional positive-degree Weil heights up to bounded error.


\paragraph*{Proposition 1}


No function satisfying (10) can obey

\[
V(T(n))<V(n)
\qquad\text{for every odd }n>1.
\tag{11}
\]


\paragraph*{Proof}


Let \(k\ge3\) be odd and set

\[
N_k=2^k-1.
\]

Define

\[
n_i=2^{k-i}3^i-1,
\qquad 0\le i\le k-1.
\]

For \(0\le i\le k-2\),

\[
3n_i+1
=
2\left(2^{k-i-1}3^{i+1}-1\right)
=
2n_{i+1},
\]

and \(n_{i+1}\) is odd. Hence \(a(n_i)=1\) and

\[
T(n_i)=n_{i+1}.
\]

At the final point,

\[
n_{k-1}=2\cdot3^{k-1}-1,
\]

and

\[
3n_{k-1}+1=2(3^k-1).
\]

Since \(k\) is odd, \(v_2(3^k-1)=1\). Therefore the last exponent is \(2\), and

\[
T^k(N_k)=M_k:=\frac{3^k-1}{2}.
\]

But

\[
\frac{M_k}{N_k}
=
\frac{3^k-1}{2(2^k-1)}
\sim
\frac12\left(\frac32\right)^k
\longrightarrow\infty.
\tag{12}
\]

If (11) held at every step, then \(V(M_k)<V(N_k)\). On the other hand,

\[
V(M_k)-V(N_k)
=
\gamma\log\frac{M_k}{N_k}+O(1)
\longrightarrow+\infty,
\]

a contradiction. \ensuremath{\blacksquare}

This already rules out every one-step Lyapunov function that differs from a positive multiple of \(\log n\) by a bounded finite-place correction.

The next result also permits unbounded local corrections with finitely many \(p\)-adic poles.


\bigskip



\subsection*{4. Arbitrarily fine adelic shadow returns}


The key construction has two freely selectable branch exponents:

\begin{itemize}
\item an entry exponent \(A\), controlling the \(2\)-adic limiting point;
\item an exit exponent \(C\), controlling the \(3\)-adic limiting point.
\end{itemize}

For even \(A\ge2\) and \(C\ge2\), define

\[
\alpha_A:=-\frac{2^A+1}{3}\in\mathbb Z_2,
\qquad
\beta_C:=-\frac{1}{2^{C+1}}\in\mathbb Z_3.
\tag{13}
\]


\paragraph*{Lemma 2: adelic shadow block}


Fix even \(A\ge2\), \(C\ge2\), and integers \(B,R\ge1\). There are arbitrarily large odd integers \(k\) for which

\[
P_{A,k}:=\frac{2^A(2^k-1)-1}{3}
\tag{14}
\]

and

\[
N_{C,k}:=\frac{3^{k+1}-1}{2^{C+1}}
\tag{15}
\]

are positive odd integers satisfying:

\begin{enumerate}
\item \(N_{C,k}\) lies in the accelerated Collatz orbit of \(P_{A,k}\);
\item the valuation word of this orbit segment is
\end{enumerate}

\[
\left(
A,
\underbrace{1,\ldots,1}_{k-1\text{ times}},
2,
C
\right);
\tag{16}
\]

\begin{enumerate}
\item the endpoints are simultaneously close at both finite places:
\end{enumerate}

\[
P_{A,k}\equiv N_{C,k}\pmod{2^B},
\qquad
P_{A,k}\equiv N_{C,k}\pmod{3^R}.
\tag{17}
\]


\paragraph*{Proof}


We first impose the \(2\)-adic congruence

\[
3^{k+2}
\equiv
3-2^{C+1}(2^A+1)
\pmod{2^{B+C+1}}.
\tag{18}
\]

The right side is congruent to \(3\pmod8\). For every \(m\ge3\), the subgroup generated by \(3\) modulo \(2^m\) consists precisely of the odd residue classes congruent to \(1\) or \(3\pmod8\). Indeed,

\[
\operatorname{ord}_{2^m}(3)=2^{m-2}.
\]

Thus (18) has a solution for \(k+2\), and every such solution has \(k\) odd.

Next impose

\[
2^{k+A+C+1}
\equiv
2^{A+C+1}+2^{C+1}-3
\pmod{3^{R+1}}.
\tag{19}
\]

The right side is a unit modulo \(3\). Since \(2\) is a primitive root modulo every \(3^m\), equation (19) has a solution modulo \(2\cdot3^R\).

The parity forced by (19) is again that \(k\) is odd. The solution class from (18) has modulus a power of \(2\), while the solution class from (19) has modulus \(2\cdot3^R\). Their residues agree modulo \(2\), so the Chinese remainder theorem gives arbitrarily large simultaneous solutions \(k\).

Now verify the dynamics.

Because \(A\) is even and \(k\) is odd,

\[
2^A(2^k-1)\equiv1\pmod3,
\]

so \(P_{A,k}\) is an integer. Moreover,

\[
3P_{A,k}+1
=
2^A(2^k-1),
\]

and \(2^k-1\) is odd. Hence

\[
P_{A,k}
\stackrel{A}{\longmapsto}
2^k-1.
\tag{20}
\]

For

\[
q_i=2^{k-i}3^i-1,
\qquad 0\le i\le k-1,
\]

we have, for \(i\le k-2\),

\[
3q_i+1
=
2\left(2^{k-i-1}3^{i+1}-1\right)
=
2q_{i+1},
\]

so the next \(k-1\) exponents are all \(1\).

At \(q_{k-1}=2\cdot3^{k-1}-1\),

\[
3q_{k-1}+1=2(3^k-1).
\]

Since \(k\) is odd, this has exact \(2\)-adic valuation \(2\), giving

\[
q_{k-1}
\stackrel{2}{\longmapsto}
M_k:=\frac{3^k-1}{2}.
\tag{21}
\]

From (18),

\[
3^{k+1}-1
\equiv
-\frac{2^{C+1}(2^A+1)}3
\pmod{2^{B+C+1}}.
\tag{22}
\]

The odd factor \((2^A+1)/3\) shows that

\[
v_2(3^{k+1}-1)=C+1.
\]

Therefore

\[
M_k
\stackrel{C}{\longmapsto}
\frac{3^{k+1}-1}{2^{C+1}}
=
N_{C,k}.
\tag{23}
\]

Dividing (22) by \(2^{C+1}\) gives

\[
N_{C,k}
\equiv
-\frac{2^A+1}{3}
\pmod{2^B}.
\]

For sufficiently large \(k\), equation (14) gives the same congruence for \(P_{A,k}\). Thus the first congruence in (17) holds.

Finally,

\[
3\left(2^{C+1}P_{A,k}+1\right)
=
2^{k+A+C+1}
-
2^{A+C+1}
-
2^{C+1}
+
3.
\]

Equation (19) therefore gives

\[
2^{C+1}P_{A,k}+1\equiv0\pmod{3^R}.
\]

Also,

\[
2^{C+1}N_{C,k}+1=3^{k+1}\equiv0\pmod{3^R}
\]

for large \(k\). Since \(2^{C+1}\) is invertible modulo \(3^R\),

\[
P_{A,k}\equiv N_{C,k}\pmod{3^R}.
\]

This proves all three assertions. \ensuremath{\blacksquare}


\bigskip



\subsubsection*{5. The archimedean gain of the shadow block}


For fixed \(A,C\),

\[
P_{A,k}\sim\frac{2^{A+k}}3,
\qquad
N_{C,k}\sim\frac{3^{k+1}}{2^{C+1}}.
\]

Hence

\[
\frac{N_{C,k}}{P_{A,k}}
\sim
\frac{3^{k+2}}{2^{A+C+k+1}}
=
\frac{9}{2^{A+C+1}}
\left(\frac32\right)^k.
\tag{24}
\]

In particular,

\[
\frac{N_{C,k}}{P_{A,k}}\longrightarrow\infty.
\tag{25}
\]

Now choose \(B=R=j\) and an increasing sequence of the corresponding integers \(k_j\). The endpoint congruences imply

\[
P_{A,k_j},N_{C,k_j}\longrightarrow\alpha_A
\qquad\text{in }\mathbb Q_2,
\tag{26}
\]

because

\[
P_{A,k_j}
=
\frac{2^{A+k_j}-2^A-1}{3}
\longrightarrow
-\frac{2^A+1}{3}.
\]

Similarly,

\[
P_{A,k_j},N_{C,k_j}\longrightarrow\beta_C
\qquad\text{in }\mathbb Q_3,
\tag{27}
\]

because

\[
N_{C,k_j}
=
\frac{3^{k_j+1}-1}{2^{C+1}}
\longrightarrow
-\frac1{2^{C+1}}.
\]

Thus the same genuine positive-integer orbit block has the following behavior:

\[
\begin{array}{c|c}
\text{place} & \text{endpoint behavior}\\
\hline
\infty & N_{C,k}/P_{A,k}\to\infty\\[2mm]
2 & P_{A,k},N_{C,k}\to\alpha_A\\[2mm]
3 & P_{A,k},N_{C,k}\to\beta_C.
\end{array}
\tag{28}
\]

This is the central obstruction.


\bigskip



\subsection*{6. No-go theorem for standard adelic local heights}



\paragraph*{Theorem 3}


Let

\[
V(n)=h_\infty(n)+h_2(n)+h_3(n)
\tag{29}
\]

on positive odd integers. Suppose:

\begin{enumerate}
\item for some \(\gamma>0\),
\end{enumerate}

\[
h_\infty(n)=\gamma\log n+O(1)
\qquad(n\to\infty);
\tag{30}
\]

\begin{enumerate}
\item \(h_2\) is locally bounded on \(\mathbb Q_2\) away from a finite set \(S_2\);
\item \(h_3\) is locally bounded on \(\mathbb Q_3\) away from a finite set \(S_3\).
\end{enumerate}

Then \(V\) cannot satisfy

\[
V(T(n))<V(n)
\qquad\text{for every positive odd }n>1.
\tag{31}
\]

This remains true with arbitrary fixed real or positive weights absorbed into the \(h_v\).


\paragraph*{Proof}


The points

\[
\alpha_A=-\frac{2^A+1}{3},
\qquad A\ge2\text{ even},
\]

are pairwise distinct. Since \(S_2\) is finite, choose an even \(A\) such that

\[
\alpha_A\notin S_2.
\]

Similarly, the points

\[
\beta_C=-2^{-(C+1)},
\qquad C\ge2,
\]

are pairwise distinct. Choose \(C\) such that

\[
\beta_C\notin S_3.
\]

Apply Lemma 2 with \(B=R=j\), obtaining orbit segments

\[
P_j\longmapsto\cdots\longmapsto N_j
\]

such that

\[
P_j,N_j\to\alpha_A\quad\text{in }\mathbb Q_2,
\]

and

\[
P_j,N_j\to\beta_C\quad\text{in }\mathbb Q_3.
\]

Because \(h_2\) and \(h_3\) are locally bounded at these points,

\[
h_2(N_j)-h_2(P_j)=O(1),
\qquad
h_3(N_j)-h_3(P_j)=O(1).
\tag{32}
\]

On the other hand, by (24),

\[
h_\infty(N_j)-h_\infty(P_j)
=
\gamma\log\frac{N_j}{P_j}+O(1)
\longrightarrow+\infty.
\tag{33}
\]

Therefore

\[
V(N_j)-V(P_j)\longrightarrow+\infty.
\tag{34}
\]

But strict descent at each intermediate accelerated step would imply

\[
V(N_j)<V(P_j),
\]

a contradiction. \ensuremath{\blacksquare}


\bigskip



\subsubsection*{7. What this theorem rules out}


The theorem covers all of the following.


\paragraph*{Conventional Weil-height constructions}


A standard local Weil function has only finitely many singular points and is locally bounded away from them. Thus no fixed positive combination of standard local heights at \(2\), \(3\), and \(\infty\) can decrease at every accelerated step.


\paragraph*{Finite lists of valuation expressions}


No height of the form

\[
\gamma\log n
+
\sum_{i=1}^{I}\lambda_i v_2(f_i(n))
+
\sum_{j=1}^{J}\mu_j v_3(g_j(n)),
\tag{35}
\]

where the \(f_i,g_j\) are fixed nonzero rational functions and the lists are finite, can work. Such functions have only finitely many \(2\)-adic and \(3\)-adic zeros and poles.

This includes arbitrary combinations involving quantities such as

\[
v_2(n+1),\quad
v_2(3n+1),\quad
v_2(3n+5),\quad
v_3(n+1),\quad
v_3(2n+1),
\]

and any finite extension of this list.


\paragraph*{Bounded finite-state corrections}


No function

\[
V(n)=\gamma\log n+\Phi(s(n)),
\tag{36}
\]

with \(s(n)\) taking only finitely many values, can work. The correction \(\Phi(s(n))\) is bounded, so Proposition 1 already applies.

Likewise, recording residues modulo fixed powers \(2^B\) and \(3^R\) is insufficient: Lemma 2 makes the beginning and end of a growing orbit block identical at any prescribed finite resolution.


\bigskip



\subsubsection*{8. Why adding more fixed poles does not solve it}


The obstruction points form the sets

\[
\mathcal A_2
=
\left\{
-\frac{2^A+1}{3}:A\ge2,\ A\text{ even}
\right\},
\tag{37}
\]

and

\[
\mathcal A_3
=
\left\{
-2^{-(C+1)}:C\ge2
\right\}.
\tag{38}
\]

The \(2\)-adic points in \(\mathcal A_2\) accumulate at

\[
-\frac13.
\]

More strongly, because \(2\) generates the unit groups modulo \(3^R\) for every \(R\), the set \(\mathcal A_3\) is dense in \(\mathbb Z_3^\times\).

Therefore, a separable local correction \(h_2+h_3\) capable of evading all the shadow blocks would have to satisfy at least one of the following pathological requirements:

\begin{itemize}
\item \(h_2\) fails to be locally bounded at every point of \(\mathcal A_2\); or
\item \(h_3\) fails to be locally bounded at every point of the dense set \(\mathcal A_3\).
\end{itemize}

Neither behavior occurs for a usual adelic local height, a finite sum of valuation terms, or a \(p\)-adic analytic potential.

The shadow points also have a simple dynamical interpretation. For the branch

\[
f_a(x)=\frac{3x+1}{2^a},
\]

one has

\[
f_A(\alpha_A)=-1,
\qquad
f_1(-1)=-1,
\qquad
f_2(-1)=-\frac12,
\qquad
f_C\!\left(-\frac12\right)=\beta_C.
\tag{39}
\]

Thus the word

\[
A,1,\ldots,1,2,C
\]

shadows the \(2\)-adic fixed point \(-1\) for a long time and then transfers that shadowing depth into a \(3\)-adic neighborhood. The Chinese-remainder construction closes the local shadow while allowing the real value to increase by \(\asymp(3/2)^k\).


\bigskip



\subsection*{9. The only viable repair}


A real-valued function that merely decreases is not enough; it must also be well-founded on the integer states. A useful block criterion is the following.


\paragraph*{Block-Lyapunov criterion}


Suppose there are:

\begin{itemize}
\item an adaptive integer \(\tau(n)\ge1\);
\item an unbounded state \(s(n)\);
\item a function
\end{itemize}

\[
V(n)=\log n+\Phi(s(n))
\]

whose sublevel sets on positive odd integers are finite;

such that for every \(n>1\),

\[
V\!\left(T^{\tau(n)}(n)\right)<V(n).
\tag{40}
\]

Then every orbit reaches \(1\).

Indeed, repeated block iteration stays inside the finite set

\[
\{m:V(m)\le V(n_0)\},
\]

while taking strictly decreasing \(V\)-values, which cannot continue indefinitely.

The shadow-block theorem imposes severe requirements on such a repair:

\begin{enumerate}
\item \textbf{The state must be unbounded.}   It cannot factor through any fixed modulus, finite valuation list, or fixed suffix length.

\item \textbf{The block endpoint must be adaptive.}   It must continue beyond arbitrarily long words of the form \[
   A,1^{k-1},2,C
   \] until the locally hidden archimedean gain has genuinely been repaid.

\item \textbf{The state must include outside-prime information or an equivalent carry invariant.}   Equation (8) shows that primes \(p\ge5\) carry exactly the multiplicative compensation omitted by the \(2\)- and \(3\)-adic terms.

\item \textbf{The required missing result is a debt-non-recycling theorem.}   A long positive-drift block must create an arithmetic debt at least as large as \[
   \log\frac{N}{P},
   \] and that debt must not be transferable indefinitely between deeper \(2\)-adic and \(3\)-adic shadows.
\end{enumerate}

Relative to the terminal-suffix locality framework developed in the earlier approaches, this pinpoints the necessary scaling: the \(3\)-adic suffix depth cannot be fixed. It must grow with the excursion length and with the exit exponent \(C\). The explicit blocks above defeat every fixed depth.


\bigskip



\subsubsection*{10. Relation to the probabilistic picture}


For a Haar-random odd \(2\)-adic starting point, the successive Syracuse exponents are independent \(\mathrm{Geom}(2)\) variables. Their mean is \(2\), producing average logarithmic drift

\[
\log3-2\log2=\log\frac34<0.
\]

This explains why an averaged Lyapunov principle is convincing. It does not give pointwise control of the exceptional positive integers embedded in \(\mathbb Z_2\). [\hypertarget{cit:R51:1}{}\hyperlink{ref:R51:1}{1}]

Tao's theorem proves that almost all Collatz orbits, in logarithmic density, attain values below every function tending to infinity, while explicitly stopping short of a uniform pointwise conclusion. [\hypertarget{cit:R51:2}{}\hyperlink{ref:R51:2}{2}] Recent \(2\)-adic/\(3\)-adic exponent-code work similarly describes its score as a finite-prefix diagnostic rather than a proof criterion. [\hypertarget{cit:R51:3}{}\hyperlink{ref:R51:3}{3}]


\bigskip



\subsection*{Verdict}


The literal approach

\[
\boxed{
\text{fixed }H_\infty,H_2,H_3
+\text{ fixed positive weights}
+\text{ one-step descent}
}
\]

is ruled out.

The obstruction is not merely that the first weighting tried fails. Explicit positive-integer orbit blocks have:

\[
\text{arbitrarily large real gain},
\]

while their endpoints are

\[
\text{arbitrarily close at both }2\text{ and }3.
\]

Consequently, no conventional finite-pole adelic height, no finite list of valuation energies, and no fixed finite-state correction can establish Collatz by strict one-step descent.

A surviving adelic strategy must become an \textbf{infinite-state, blockwise transversality argument involving the \(p\ge5\) leakage or an equivalent exact carry debt}. Establishing that debt cannot be recycled along one nested integer orbit is the unresolved theorem. Thus no full proof of the Collatz conjecture results from approach 46 as stated; what is complete is the no-go theorem sharply delimiting every viable repair.
\subsection*{References}
\begin{enumerate}[leftmargin=2.2em]
\item \hypertarget{ref:R51:1}{}\hyperlink{cit:R51:1}{\texttt{turn291036view1}}
\item \hypertarget{ref:R51:2}{}\hyperlink{cit:R51:2}{\texttt{turn985209view2}}
\item \hypertarget{ref:R51:3}{}\hyperlink{cit:R51:3}{\texttt{turn291036view3}}
\end{enumerate}

\clearpage
\section{Round 52: approach 47 --- Vojta or abc for uniform near-cycle exclusion}
\subsection*{Prompt}
\begin{quote}
Complete a full proof of the Collatz conjecture using the following approach:

\begin{enumerate}
\item Vojta or abc for uniform near-cycle exclusion
\end{enumerate}

\textbf{Hybrid \textperiodcentered{} P}

Applications of linear forms in logarithms already constrain exact cycles, and possible relationships with abc-type inequalities have been studied. The stronger target is to rule out arbitrarily long \textbf{near-cycles}, with uniform control of the radical of the affine numerator C\_j.

Take as long as you need.
\end{quote}
\subsection*{Response}

\subsection*{Result of the abc/Vojta route}


I carried the proposed argument through to the strongest form that I can justify. It does \textbf{not} currently give a full proof of Collatz, even assuming the classical \(abc\) conjecture or the standard higher-dimensional forms of Vojta's conjecture.

The obstruction is precise:

\begin{enumerate}
\item A first crossing of the critical line really does produce an additive near-cycle.
\item Baker's theorem makes the associated starting value only polynomially large.
\item Classical \(abc\) would then finish the argument if the affine numerator \(C_j\) had a uniform exponential saving in its radical.
\item That radical estimate is false for general first-crossing words: \(abc\) itself predicts almost maximal radical.
\item Expanding \(C_j\) into its constituent \(2\)-\(3\) units removes \(\operatorname{rad}(C_j)\), but produces a growing-dimensional \(abcde\ldots\) equation. A dimension-uniform inequality is false without exceptional families\textemdash{}the trivial Collatz cycle is already an explicit exception.
\item The corrected target is a \textbf{uniform moving-polynomial gcd theorem}, with periodic words as its exceptional set. Such a theorem would rule out all nonperiodic infinite orbits, but exact nontrivial cycles would still require a separate exclusion theorem.
\end{enumerate}

The exact reductions follow.


\bigskip



\subsubsection*{1. Affine formulation}


Use the odd accelerated map

\[
U(x)=\frac{3x+1}{2^{\nu _2(3x+1)}}.
\]

Let

\[
x_0=n,\qquad x_{i+1}=U(x_i),\qquad
a_i=\nu _2(3x_i+1),
\]

and put

\[
A_0=0,\qquad A_j=\sum_{i=0}^{j-1}a_i.
\]

Define \(C_0=0\) and

\[
C_{j+1}=3C_j+2^{A_j}.
\]

Induction gives the exact affine identity

\[
\boxed{2^{A_j}x_j=3^j n+C_j}
\tag{1}
\]

where

\[
\boxed{
C_j=\sum_{i=0}^{j-1}3^{j-1-i}2^{A_i}.
}
\tag{2}
\]

For \(j\geq 1\), \(C_j\) is odd and \(3\nmid C_j\).

Define

\[
D_j=2^{A_j}-3^j,\qquad d_j=x_j-n.
\]

Subtracting \(2^{A_j}n\) from (1) gives

\[
\boxed{C_j=D_jn+2^{A_j}d_j.}
\tag{3}
\]

An equivalent and particularly useful form is

\[
\boxed{C_j-3^jd_j=D_jx_j.}
\tag{4}
\]

Thus a near-cycle \(x_j=n+d_j\) forces the large number \(D_j\) to divide a word-dependent sparse exponential sum:

\[
D_j\mid C_j-3^jd_j.
\tag{5}
\]

This divisibility, rather than the radical of \(C_j\) alone, is the correct Diophantine target.


\bigskip



\subsubsection*{2. First crossings really are uniform near-cycles}


Write

\[
\rho=\log_2 3.
\]

Suppose \(j\) is a first crossing of the critical line:

\[
A_i\leq \rho i\quad(0\leq i<j),
\qquad
A_j>\rho j.
\tag{6}
\]

Then \(2^{A_i}\leq 3^i\) for every \(i<j\), so every term in (2) is at most \(3^{j-1}\). Consequently,

\[
C_j\leq j3^{j-1}.
\tag{7}
\]

Assume also that the orbit has not descended at time \(j\), so \(x_j\geq n\), or equivalently \(d_j\geq0\). Since \(D_j>0\), equation (3) gives

\[
0\leq 2^{A_j}d_j=C_j-D_jn\leq C_j.
\]

Using \(2^{A_j}>3^j\) and (7),

\[
\boxed{0\leq d_j<\frac j3.}
\tag{8}
\]

So every surviving first crossing is an additive near-cycle of displacement \(O(j)\). This part of the proposed approach works exactly.

There is also a useful overshoot bound. Since \(n\geq1\) and \(d_j\geq0\),

\[
D_j\leq C_j\leq j3^{j-1},
\]

and hence

\[
2^{A_j}
\leq 3^j\left(1+\frac j3\right).
\]

Therefore

\[
0<A_j-\rho j
\leq \log_2\left(1+\frac j3\right)=O(\log j).
\tag{9}
\]

A surviving first crossing is not merely on the contracting side; it lies logarithmically close to the critical line.


\bigskip



\subsubsection*{3. Baker reduces the starting value to polynomial size}


Let

\[
\Lambda_j=A_j\log 2-j\log 3>0.
\]

When \(A_j=O(j)\), a lower bound for a nonzero linear form in \(\log 2,\log 3\) gives constants \(c>0\) and \(\kappa>0\) such that

\[
\Lambda_j\geq c j^{-\kappa}.
\]

Since

\[
D_j=3^j(e^{\Lambda_j}-1),
\]

we obtain

\[
\boxed{
D_j\geq c\,3^j j^{-\kappa}.
}
\tag{10}
\]

If \(A_j\) is much larger than \(j\), this estimate is trivial because \(D_j\) is even larger; condition (9) already restricts a surviving first crossing to the relevant range. Linear-form estimates of this kind are the standard input in the work on possible Collatz cycles. [\hypertarget{cit:R52:1}{}\hyperlink{ref:R52:1}{1}, \hypertarget{cit:R52:2}{}\hyperlink{ref:R52:2}{2}]

Combining (7), (10), and \(D_jn\leq C_j\),

\[
n\leq \frac{C_j}{D_j}
\ll j^{\kappa+1}.
\tag{11}
\]

Together with (8),

\[
x_j=n+d_j\ll j^{\kappa+1}.
\tag{12}
\]

Thus a surviving first crossing would have all the following properties simultaneously:

\[
A_j=\rho j+O(\log j),\qquad
n,x_j=j^{O(1)},\qquad
0\leq x_j-n<\frac j3.
\tag{13}
\]

This is an extremely rigid near-cycle reduction.


\bigskip



\subsubsection*{4. What classical \(abc\) would need}


Let

\[
g=\gcd(n,C_j).
\]

Because \(3\nmid C_j\), dividing (1) by \(g\) produces a primitive \(abc\)-triple:

\[
3^j\frac ng+\frac{C_j}{g}
=
2^{A_j}\frac{x_j}{g}.
\tag{14}
\]

Classical \(abc\) would give, for every \(\varepsilon>0\),

\[
2^{A_j}\frac{x_j}{g}
\ll_\varepsilon
\operatorname{rad}(6n x_j C_j)^{1+\varepsilon}.
\tag{15}
\]

At a surviving first crossing, \(n\) and \(x_j\) are polynomial in \(j\). Therefore their radicals contribute only \(e^{O(\log j)}\). Consequently, (15) would contradict the exponential size \(3^j\) provided one could prove, for some fixed \(\delta>0\),

\[
\boxed{
\log\operatorname{rad}(C_j)
\leq (\log 3-\delta)j+O(\log j).
}
\tag{16}
\]

So the proposed radical strategy has a mathematically correct conditional core:


\begin{quote}
\textbf{A uniform exponential radical saving for every surviving first-crossing numerator would, together with \(abc\) and Baker's theorem, exclude all sufficiently long surviving first crossings.}

\end{quote}


The problem is that such a radical saving is not true for general first-crossing words.


\bigskip



\subsubsection*{5. A first-crossing family for which \(abc\) predicts maximal radical}


For \(j\geq2\), consider the valuation word

\[
w_j=(\underbrace{1,\ldots,1}_{j-1\text{ entries}},b_j),
\qquad
b_j=\lceil \rho j\rceil-j+1.
\tag{17}
\]

Then

\[
A_i=i\quad(0\leq i<j),
\qquad
A_j=\lceil\rho j\rceil.
\]

It is therefore a genuine first-crossing word. Its affine numerator is

\[
\begin{aligned}
C_j
&=\sum_{i=0}^{j-1}3^{j-1-i}2^i\\
&=\frac{3^j-2^j}{3-2}\\
&=\boxed{3^j-2^j}.
\end{aligned}
\tag{18}
\]

Apply \(abc\) to

\[
2^j+(3^j-2^j)=3^j.
\tag{19}
\]

The three terms are pairwise coprime, so \(abc\) predicts

\[
3^j
\ll_\varepsilon
\bigl(6\operatorname{rad}(3^j-2^j)\bigr)^{1+\varepsilon}.
\]

Hence

\[
\boxed{
\log\operatorname{rad}(C_j)
\geq
\frac{j\log3}{1+\varepsilon}-O_\varepsilon(1).
}
\tag{20}
\]

Because \(C_j<3^j\), allowing \(\varepsilon\to0\) says heuristically\textemdash{}and conditionally under \(abc\)\textemdash{}that

\[
\log\operatorname{rad}(C_j)
=
j\log3-o(j).
\tag{21}
\]

This is almost the largest possible radical. It is exactly the opposite of (16).

The associated finite orbit segment descends, so this does not logically disprove a radical saving restricted only to \textbf{surviving} first crossings. But it proves that such a saving cannot come from the shape of \(C_j\), from first crossing alone, or from \(abc\). It would have to use the full integrality and no-descent condition in an essentially new way.

Rozier's \(abc\)-based result succeeds for the much more special shortcut parity patterns having exactly one even term; there the orbit equation reduces to a controlled three-term equation such as \(1+3^kA=2^{j-k}B\). The paper explicitly identifies the unbounded number of terms in a general Collatz iterate as the obstacle to extending the argument. [\hypertarget{cit:R52:3}{}\hyperlink{ref:R52:3}{3}, \hypertarget{cit:R52:4}{}\hyperlink{ref:R52:4}{4}]


\bigskip



\subsubsection*{6. Expanding \(C_j\): the natural Vojta formulation}


Instead of treating \(C_j\) as one term, expand it using (2):

\[
3^jn+
3^{j-1}
+3^{j-2}2^{A_1}
+\cdots+
2^{A_{j-1}}
=
2^{A_j}x_j.
\tag{22}
\]

Consider the primitive projective tuple

\[
P_j=
\left[
3^jn:
3^{j-1}:
3^{j-2}2^{A_1}:
\cdots:
2^{A_{j-1}}:
-2^{A_j}x_j
\right].
\tag{23}
\]

Its logarithmic height satisfies

\[
h(P_j)\geq j\log3+\log n.
\tag{24}
\]

All the interior coordinates have prime factors only \(2\) and \(3\). Therefore the logarithmic projective radical is exactly

\[
R(P_j)
=
\log\operatorname{rad}(6n x_j).
\tag{25}
\]

This is the major advantage: \(\operatorname{rad}(C_j)\) has disappeared.

Moreover,

\[
U(x)=\frac{3x+1}{2^{a}}
\leq \frac{3x+1}{2}\leq2x
\]

for every positive odd \(x\). Hence

\[
x_j\leq2^j n
\tag{26}
\]

and

\[
R(P_j)
\leq j\log2+2\log n+\log6.
\tag{27}
\]

Suppose one had a dimension-uniform \(abcde\ldots\) inequality

\[
h(P_j)
\leq
(1+\varepsilon)R(P_j)+o(j)
\tag{28}
\]

for infinitely many prefixes of every nonperiodic orbit. Combining (24) and (27),

\[
j\log3+\log n
\leq
(1+\varepsilon)
\bigl(j\log2+2\log n+\log6\bigr)+o(j).
\]

Thus

\[
j\bigl(\log3-(1+\varepsilon)\log2\bigr)
\leq
(1+2\varepsilon)\log n+O(1)+o(j).
\tag{29}
\]

Choose

\[
0<\varepsilon<\log_2 3-1.
\]

Then the coefficient on the left is positive. For a fixed starting value \(n\), (29) is impossible as \(j\to\infty\).

This proves the following conditional result.


\paragraph*{Conditional Vojta reduction}


Assume that every infinite accelerated Collatz orbit whose valuation sequence is not eventually periodic has infinitely many prefixes satisfying the uniform inequality (28), with an error \(o(j)\). Then every positive accelerated Collatz orbit has eventually periodic valuation sequence.

This is a genuine full exclusion of nonperiodic divergent orbits under the proposed uniform Vojta statement.

For a finite prefix, the same calculation gives an exponential seed lower bound:

\[
\log n
\geq
\frac{\log3-(1+\varepsilon)\log2}{1+2\varepsilon}\,j-o(j).
\tag{30}
\]

This is precisely the kind of lower-bound hypothesis that Collatz parity-vector approaches seek.


\bigskip



\subsubsection*{7. Why standard Vojta does not supply (28)}


There are two fundamental problems.


\paragraph*{7.1 The dimension grows}


The tuple \(P_j\) lies in a hyperplane of dimension \(j\). Standard higher-dimensional \(abcd\) statements give

\[
h(P)\leq(1+\varepsilon)R(P)+C_{\varepsilon,j}
\]

outside a proper Zariski-closed subset, but both the constant and the exceptional subset depend on the dimension. Nothing in the standard formulation gives

\[
C_{\varepsilon,j}=o(j)
\]

or coherent exceptional sets as \(j\to\infty\).

Vojta's original higher-dimensional formulation is known to require modification; the surviving formulations explicitly allow proper Zariski-closed exceptional subsets, and the constants depend on the dimension. [\hypertarget{cit:R52:5}{}\hyperlink{ref:R52:5}{5}]


\paragraph*{7.2 The trivial Collatz orbit is an explicit exception}


For \(n=1\),

\[
a_i=2,\qquad A_i=2i,\qquad x_i=1.
\]

Equation (22) becomes

\[
\boxed{
3^j+
\sum_{i=0}^{j-1}3^{j-1-i}4^i
=
4^j.
}
\tag{31}
\]

Every coordinate has prime factors only \(2\) and \(3\), so

\[
R(P_j)=\log6,
\qquad
h(P_j)=j\log4.
\]

Thus no inequality of the form

\[
h(P_j)\leq(1+\varepsilon)R(P_j)+o(j)
\]

can hold for all Collatz staircase tuples. This remains true even though (31) has no proper zero subsum: all terms on the left are positive.

Therefore exceptional families are not a removable technicality. The trivial cycle forces one. Every hypothetical nontrivial Collatz cycle would generate an analogous exceptional family by repeating its valuation block.


\bigskip



\subsubsection*{8. Periodic exceptions can be analyzed exactly}


Let \(v\) be a valuation block of length \(r\), total valuation

\[
B=A(v),
\]

and affine numerator \(C_v\). The concatenation identity is

\[
C_{uv}=3^{|v|}C_u+2^{A(u)}C_v.
\tag{32}
\]

For the repeated word \(v^m\),

\[
\boxed{
C_{v^m}=C_vQ_m
}
\tag{33}
\]

where

\[
\boxed{
Q_m=
\sum_{k=0}^{m-1}2^{Bk}3^{r(m-1-k)}
=
\frac{2^{Bm}-3^{rm}}{2^B-3^r}.
}
\tag{34}
\]

Also,

\[
2^{Bm}-3^{rm}=(2^B-3^r)Q_m.
\tag{35}
\]

Suppose \(v^m\) carries \(n\) to \(n+d\). Equation (3) becomes

\[
C_vQ_m
=
(2^B-3^r)Q_m n+2^{Bm}d.
\]

Hence

\[
Q_m\mid 2^{Bm}d.
\]

But \(Q_m\) is odd, so

\[
\boxed{Q_m\mid d.}
\tag{36}
\]

Consequently,

\[
\boxed{|d|<Q_m\quad\Longrightarrow\quad d=0.}
\tag{37}
\]

Thus a long repeated-block near-cycle cannot have a small nonzero displacement. It must collapse to an exact cycle.

Moreover, when \(d=0\), dividing by \(Q_m\) gives

\[
C_v=(2^B-3^r)n.
\tag{38}
\]

So even if the near-cycle was initially observed only after \(m\) repetitions, one copy of \(v\) already fixes \(n\).

This is an unconditional and useful near-cycle exclusion theorem.


\bigskip



\subsubsection*{9. Eventually periodic valuations imply an actual cycle}


Suppose the valuation sequence is eventually periodic with block \(v\). Starting at the beginning of the periodic tail, let \(y_k\) be the orbit value after \(k\) copies of \(v\). Then

\[
2^B y_{k+1}=3^r y_k+C_v.
\tag{39}
\]

Set

\[
y_*=\frac{C_v}{2^B-3^r}.
\]

Since \(2^B-3^r\) is odd, \(y_*\in\mathbb Z_2\). Equation (39) gives

\[
y_{k+1}-y_*
=
\frac{3^r}{2^B}(y_k-y_*).
\tag{40}
\]

If \(y_0-y_*\neq0\), then

\[
\nu_2(y_k-y_*)
=
\nu_2(y_0-y_*)-Bk,
\]

which eventually becomes negative. But \(y_k\) and \(y_*\) are both \(2\)-adic integers, so their difference must remain \(2\)-adically integral. This is impossible.

Therefore

\[
y_0=y_*,
\]

and (39) gives \(y_{k+1}=y_k\). Hence:

\[
\boxed{\text{An eventually periodic valuation sequence on a positive integer orbit is a genuine cycle.}}
\tag{41}
\]

Combining this with the conditional Vojta reduction:


\begin{quote}
A dimension-uniform Vojta inequality outside coherent periodic exceptions would prove that every positive orbit eventually enters a cycle.

\end{quote}


It would not yet prove that this cycle is the trivial one.


\bigskip



\subsubsection*{10. The correct gcd target}


For a valuation word \(w\) of length \(j\), define the two moving polynomials

\[
G_w(X,Y)=X^{A_j}-Y^j
\tag{42}
\]

and

\[
F_{w,d}(X,Y)
=
\sum_{i=0}^{j-1}X^{A_i}Y^{j-1-i}
-
Y^jd.
\tag{43}
\]

At \((X,Y)=(2,3)\),

\[
G_w(2,3)=D_j
\]

and, using (4),

\[
F_{w,d_j}(2,3)
=
C_j-3^jd_j
=
D_jx_j.
\]

Therefore

\[
\boxed{
\gcd\bigl(G_w(2,3),F_{w,d_j}(2,3)\bigr)=|D_j|.
}
\tag{44}
\]

For a first-crossing near-cycle, Baker gives

\[
\log |D_j|=j\log3-O(\log j).
\tag{45}
\]

Thus the ideal Vojta-type theorem would be:


\paragraph*{Uniform Collatz gcd theorem}


For every \(\eta>0\), outside explicitly classified periodic or algebraic word families,

\[
\log\gcd\bigl(
G_w(2,3),F_{w,d}(2,3)
\bigr)
\leq
\eta j+O_\eta(\log(1+|d|)).
\tag{46}
\]

For \(0\leq d<j/3\), equations (44)\textendash{}(46) would immediately contradict (45).

This is a substantially better target than a radical upper bound for \(C_j\):

\begin{itemize}
\item It directly captures the divisibility forced by a near-cycle.
\item It naturally admits periodic exceptional families.
\item It resembles established gcd theorems for values of coprime polynomials at \(S\)-unit points.
\end{itemize}

Corvaja\textendash{}Zannier prove such small-gcd conclusions for \textbf{fixed} coprime polynomials evaluated on points in a finitely generated subgroup of a torus, outside a proper Zariski-closed set. [\hypertarget{cit:R52:6}{}\hyperlink{ref:R52:6}{6}] Moving-target variants require conditions such as fixed degree and coefficient height negligible compared with the point height. Here the evaluation point is fixed at \((2,3)\), while the degrees, exponent sets, and number of monomials in \(F_{w,d}\) grow with \(j\); those hypotheses do not apply. [\hypertarget{cit:R52:7}{}\hyperlink{ref:R52:7}{7}]

Proving (46), including a uniform classification of its exceptional words, is therefore a genuinely new theorem rather than an application of an existing Vojta result.


\bigskip



\subsubsection*{11. Exact cycles remain a separate obstruction}


Suppose \(v\) is an exact cycle block. Then

\[
2^B-3^r\mid C_v
\]

and

\[
n=\frac{C_v}{2^B-3^r}.
\tag{47}
\]

Repeating the block gives the geometric identities

\[
3^{rm}n+C_vQ_m=2^{Bm}n.
\tag{48}
\]

These are necessarily exceptional to the generic Vojta inequality: their expanded radical is bounded by \(\operatorname{rad}(6n)\), while their height grows exponentially in \(m\).

Compressing (48) into three terms does not help:

\[
3^{rm}+(2^B-3^r)Q_m=2^{Bm},
\tag{49}
\]

because the radical now contains \(Q_m\), which normally introduces exponentially many arithmetic contributions. Classical \(abc\) is entirely compatible with (49).

Thus an exceptional-set theorem saying ``all exceptions are periodic'' reduces the problem to exact cycles but cannot distinguish the trivial block \((2)\) from a hypothetical nontrivial cycle block.

Linear forms in logarithms give powerful bounds on possible cycles, and established work has excluded Collatz \(m\)-cycles with \(m\leq91\), but these results do not give an all-\(m\) exclusion. [\hypertarget{cit:R52:8}{}\hyperlink{ref:R52:8}{8}, \hypertarget{cit:R52:9}{}\hyperlink{ref:R52:9}{9}]


\bigskip



\subsection*{Final assessment}


The proposed idea splits into three statements:


\paragraph*{1. First-crossing near-cycle reduction \textemdash{} proved}


A surviving first crossing satisfies

\[
0\leq x_j-n<\frac j3,
\qquad
A_j-j\log_2 3=O(\log j),
\qquad
n,x_j=j^{O(1)}.
\]


\paragraph*{2. Uniform small radical of \(C_j\) \textemdash{} false in the natural unrestricted form}


The first-crossing family

\[
C_j=3^j-2^j
\]

has, under \(abc\),

\[
\operatorname{rad}(C_j)=3^{j-o(j)}.
\]

So \(abc\) predicts almost maximal radical rather than the exponential saving required by the compressed three-term argument.


\paragraph*{3. Correct replacement \textemdash{} a new uniform gcd/Vojta theorem}


The right statement is the moving-polynomial estimate

\[
\log\gcd\left(
2^{A_j}-3^j,\,
C_j-3^jd
\right)=o(j)
\]

outside coherent periodic exceptional families. Together with Baker and the first-crossing lemma, this would uniformly exclude long nonperiodic near-cycles. A dimension-uniform staircase \(abcde\ldots\) theorem with the same exceptional-set classification would rule out all nonperiodic infinite orbits.

But periodic exceptional families are unavoidable, and exact nontrivial cycles remain inside them. Their complete exclusion is not supplied by either classical \(abc\), standard Vojta, or existing linear-form estimates.

Therefore the approach does \textbf{not} presently yield a valid full proof of Collatz. Its strongest repaired form is:

\[
\boxed{
\begin{array}{c}
\text{uniform moving-gcd estimate}\\
+\ \text{coherent periodic-exception classification}\\
+\ \text{complete positive-cycle exclusion}\\[2mm]
\Longrightarrow\ \text{Collatz}.
\end{array}}
\]

The central proposed radical lemma cannot be used as stated; the genuinely viable research target is the uniform gcd estimate (46), with the periodic factorization lemma (33)\textendash{}(37) describing the unavoidable exceptional geometry.
\subsection*{References}
\begin{enumerate}[leftmargin=2.2em]
\item \hypertarget{ref:R52:1}{}\hyperlink{cit:R52:1}{\texttt{turn324824search6}}
\item \hypertarget{ref:R52:2}{}\hyperlink{cit:R52:2}{\texttt{turn324824search9}}
\item \hypertarget{ref:R52:3}{}\hyperlink{cit:R52:3}{\texttt{turn617796view3}}
\item \hypertarget{ref:R52:4}{}\hyperlink{cit:R52:4}{\texttt{turn324824search14}}
\item \hypertarget{ref:R52:5}{}\hyperlink{cit:R52:5}{\texttt{turn904164view0}}
\item \hypertarget{ref:R52:6}{}\hyperlink{cit:R52:6}{\texttt{turn870361view3}}
\item \hypertarget{ref:R52:7}{}\hyperlink{cit:R52:7}{\texttt{turn375000search21}}
\item \hypertarget{ref:R52:8}{}\hyperlink{cit:R52:8}{\texttt{turn617796view2}}
\item \hypertarget{ref:R52:9}{}\hyperlink{cit:R52:9}{\texttt{turn324824search1}}
\end{enumerate}

\clearpage
\section{Round 53: approach 48 --- The Subspace Theorem for trajectory identities}
\subsection*{Prompt}
\begin{quote}
Complete a full proof of the Collatz conjecture using the following approach:

\begin{enumerate}
\item The Subspace Theorem for trajectory identities
\end{enumerate}

\textbf{Speculative \textperiodcentered{} P}

Long trajectories create many relations involving powers of 2, powers of 3, and partial carry sums. Recast these as families of S-unit inequalities. If too many unusually close cancellations occur, the Subspace Theorem should force a fixed linear dependence and hence repetitive dynamics.

Take as long as you need.
\end{quote}
\subsection*{Response}

\subsection*{Result of the attempt}


The proposed route does \textbf{not} presently give a valid full proof of the Collatz conjecture. Pushing it carefully reveals two exact obstructions:

\begin{enumerate}
\item the trajectory relation has a number of \(S\)-unit summands that grows with the trajectory length;  
\item every natural fixed-dimensional compression lands exactly on\textemdash{}or above\textemdash{}the critical exponent of the Subspace Theorem, rather than gaining the required factor \(H^{-\varepsilon}\).
\end{enumerate}

Below is a rigorous derivation of those obstructions, together with an exact conditional theorem showing what new ``trajectory Subspace Theorem'' would suffice to finish the divergent-orbit part of Collatz.


\bigskip



\subsubsection*{1. The exact trajectory identity}


Work with the accelerated map on positive odd integers,

\[
U(x)=\frac{3x+1}{2^{\nu _2(3x+1)}}.
\]

Let

\[
x_{j+1}=U(x_j),\qquad   
a_j=\nu _2(3x_j+1),\qquad
A_0=0,\qquad
A_j=\sum_{i=0}^{j-1}a_i.
\]

Thus \(a_j\geq 1\) and

\[
2^{a_j}x_{j+1}=3x_j+1.
\]

Iterating gives

\[
\boxed{
2^{A_j}x_j=3^j x_0+C_j
}
\tag{1}
\]

where

\[
\boxed{
C_j=\sum_{i=0}^{j-1}3^{j-1-i}2^{A_i}.
}
\tag{2}
\]

Dividing (1) by \(3^j\), define

\[
h_j:=\frac{2^{A_j}x_j}{3^j}.
\]

Then

\[
\boxed{
h_j=x_0+\sum_{i=0}^{j-1}\frac{2^{A_i}}{3^{i+1}}.
}
\tag{3}
\]

Since \(A_j\geq j\), the left side tends to \(0\) in \(\mathbb Q_2\). Hence every infinite valuation trajectory satisfies the exact \(2\)-adic identity

\[
\boxed{
x_0=-\sum_{i=0}^{\infty}\frac{2^{A_i}}{3^{i+1}}
\quad\text{in }\mathbb Q_2.
}
\tag{4}
\]

This is the cleanest possible \(S\)-unit formulation: every summand is in the rank-two multiplicative group generated by \(2\) and \(3\).


\bigskip



\subsubsection*{2. Why the ordinary \(S\)-unit theorem does not apply uniformly}


The Evertse\textendash{}Schlickewei\textendash{}Schmidt theorem controls nondegenerate solutions of

\[
c_1y_1+\cdots+c_dy_d=1,
\qquad y_i\in\Gamma,
\]

when both the dimension \(d\) and the rank of \(\Gamma\) are fixed. Its bound depends explicitly on the dimension and rank. [\hypertarget{cit:R53:1}{}\hyperlink{ref:R53:1}{1}]

But (1) is the relation

\[
2^{A_j}x_j-3^jx_0-
\sum_{i=0}^{j-1}3^{j-1-i}2^{A_i}=0,
\tag{5}
\]

which has \(j+2\) summands. Its dimension grows with \(j\).

This is not a superficial technical issue. Fixed rank does not bound the length of a nondegenerate relation. The trivial orbit \(x_j=1\), for which \(a_j=2\) and \(A_j=2j\), produces

\[
\boxed{
4^j=3^j+\sum_{i=0}^{j-1}3^{j-1-i}4^i.
}
\tag{6}
\]

This is a nondegenerate relation of arbitrary length inside the rank-two group \(\langle 2,3\rangle\). It is merely the geometric identity

\[
4^j-3^j=(4-3)\sum_{i=0}^{j-1}4^i3^{j-1-i}.
\]

Thus a theorem saying ``long \(S\)-unit relations in fixed rank must be degenerate'' would already be false. What one would need is a theorem saying that every sufficiently long \textbf{trajectory-shaped} relation is generated by short identities such as \(4=3+1\). That is much stronger than the classical Subspace Theorem.

This same unbounded-summand obstruction appears in the standard cycle equation. Tao notes that bounded-exception valuation patterns can be compressed into boundedly many powers of \(2\) and \(3\), but general valuation words cannot. [\hypertarget{cit:R53:2}{}\hyperlink{ref:R53:2}{2}]


\bigskip



\subsubsection*{3. A fixed-dimensional compression theorem\textemdash{}and its exact failure}


One might try to group the \(j\) carry summands into a fixed number \(m\) of blocks. The following calculation shows why this loses precisely the divisibility gain that the Subspace Theorem would need.


\paragraph*{Proposition 1: the compression barrier}


Fix \(N\), and partition

\[
\{0,1,\ldots,N-1\}=I_1\sqcup\cdots\sqcup I_m
\]

into \(m\) nonempty sets. For each \(r\), put

\[
s_r=\min I_r,\qquad t_r=\max I_r,
\]

and define the grouped carry

\[
B_r=\sum_{i\in I_r}\frac{2^{A_i}}{3^{i+1}}.
\]

Then

\[
B_r=\frac{2^{A_{s_r}}K_r}{3^{t_r+1}},
\tag{7}
\]

where

\[
K_r=\sum_{i\in I_r}3^{t_r-i}2^{A_i-A_{s_r}}.
\tag{8}
\]

Moreover,

\[
\gcd(K_r,6)=1,
\qquad
K_r\geq 3^{t_r-s_r},
\tag{9}
\]

and consequently

\[
\boxed{
\prod_{r=1}^m K_r\geq 3^{N-m}.
}
\tag{10}
\]


\subparagraph*{Proof}


The term indexed by \(s_r\) in (8) is \(3^{t_r-s_r}\), which is odd, while every other term is even. Hence \(K_r\) is odd.

The term indexed by \(t_r\) is \(2^{A_{t_r}-A_{s_r}}\), which is nonzero modulo \(3\), while every earlier term is divisible by \(3\). Hence \(3\nmid K_r\).

The first term also gives \(K_r\geq3^{t_r-s_r}\). Finally,

\[
t_r-s_r\geq |I_r|-1,
\]

so

\[
\sum_{r=1}^m(t_r-s_r)
 \geq \sum_{r=1}^m(|I_r|-1)
 =N-m.
\]

This proves (10). \(\square\)

Now put

\[
Y=(Y_0,\ldots,Y_m)=(x_0,B_1,\ldots,B_m).
\]

At the \(2\)-adic place use the linearly independent forms

\[
L_{2,0}=Y_0+\cdots+Y_m,\qquad
L_{2,r}=Y_r\quad(1\leq r\leq m).
\]

At the real and \(3\)-adic places use the coordinate forms \(Y_0,\ldots,Y_m\).

Because

\[
Y_0+\cdots+Y_m=h_N=\frac{2^{A_N}x_N}{3^N},
\]

we have

\[
|L_{2,0}(Y)|_2=2^{-A_N}.
\]

Using (7) and \(\gcd(K_r,6)=1\),

\[
|B_r|_2=2^{-A_{s_r}},
\qquad
|B_r|_3=3^{t_r+1}.
\]

All powers of \(2\) and \(3\) contributed by the separate block coordinates cancel between the real, \(2\)-adic and \(3\)-adic places. One obtains the exact identity

\[
\boxed{
\prod_{v\in\{\infty,2,3\}}
\prod_{i=0}^{m}|L_{v,i}(Y)|_v
=
x_0^{(3)}\,2^{-A_N}\prod_{r=1}^mK_r,
}
\tag{11}
\]

where

\[
x_0^{(3)}=\frac{x_0}{3^{\nu_3(x_0)}}
\]

is the \(3\)-free part of \(x_0\).

Combining (10) and (11),

\[
\boxed{
\prod_{v,i}|L_{v,i}(Y)|_v
\geq
x_0^{(3)}
\exp\!\left(
N\log3-m\log3-A_N\log2
\right).
}
\tag{12}
\]

The affine Subspace Theorem requires a strict inequality of the form

\[
\prod_{v,i}|L_{v,i}(Y)|_v<H(Y)^{-\varepsilon}
\tag{13}
\]

for some fixed \(\varepsilon>0\). [\hypertarget{cit:R53:3}{}\hyperlink{ref:R53:3}{3}]

The important point is that (12) has no such power saving on a critical trajectory.


\bigskip



\subsubsection*{4. Why a hypothetical divergent orbit is necessarily critical}


Suppose Collatz is false and choose the least positive integer \(x_0\) whose orbit does not reach \(1\). It must be odd, and every odd iterate satisfies

\[
x_j\geq x_0.
\tag{14}
\]

If the accelerated orbit repeats, it is a nontrivial cycle. Otherwise, since it is an injective sequence of positive integers bounded below by \(x_0\), necessarily

\[
x_j\longrightarrow\infty.
\tag{15}
\]

From

\[
x_{j+1}
=\frac{3x_j}{2^{a_j}}
\left(1+\frac1{3x_j}\right)
\]

we obtain

\[
\log x_N-\log x_0
=
N\log3-A_N\log2+
\sum_{j=0}^{N-1}
\log\left(1+\frac1{3x_j}\right).
\]

Equivalently,

\[
A_N\log2
=
N\log3+
\sum_{j=0}^{N-1}
\log\left(1+\frac1{3x_j}\right)
-\log\frac{x_N}{x_0}.
\tag{16}
\]

Since \(x_N\geq x_0\), the last term is nonpositive. Since \(x_j\to\infty\),

\[
\frac1N
\sum_{j=0}^{N-1}
\log\left(1+\frac1{3x_j}\right)
\longrightarrow0.
\]

Therefore

\[
\boxed{
A_N\log2\leq N\log3+o(N).
}
\tag{17}
\]

Substitution into (12) gives, for every fixed \(m\),

\[
\prod_{v,i}|L_{v,i}(Y)|_v
\geq \exp(-o(N)).
\tag{18}
\]

On the other hand, one of the groups contains \(N-1\), so one coordinate \(B_r\) has reduced denominator \(3^N\). Hence

\[
H(Y)\geq3^N.
\]

Since \(A_N=O(N)\), one also has \(H(Y)\leq\exp(O(N))\). Thus

\[
\boxed{
\prod_{v,i}|L_{v,i}(Y)|_v
\geq H(Y)^{-o(1)}.
}
\tag{19}
\]

But the Subspace Theorem needs \(H(Y)^{-\varepsilon}\) for a \textbf{fixed} positive \(\varepsilon\).


\paragraph*{Consequence}


For a hypothetical minimal divergent trajectory, no grouping of all \(N\) carry terms into a fixed number \(m\) of positive \(S\)-unit coordinates can trigger the Subspace Theorem.

Letting \(m\) grow would avoid this particular obstruction, but then the projective dimension grows with \(N\), returning to the original problem.

Moving the integers \(K_r\) into the coefficients of moving hyperplanes also fails. For fixed \(m\), some group has span at least \(N/m\), and therefore

\[
h(K_r)\geq \left(\frac Nm-O(1)\right)\log3.
\]

This is comparable to the height of the point rather than \(o(h(Y))\). Moving-target versions of the Subspace Theorem explicitly require the target heights to be \(o\) of the point height. [\hypertarget{cit:R53:4}{}\hyperlink{ref:R53:4}{4}]


\bigskip



\subsubsection*{5. The strong-escape regime lies exactly at the critical exponent}


There is another possible plan: suppose the orbit has uniform positive multiplicative drift, so that for some \(\delta>0\),

\[
A_j\log2\leq j\log3-\delta j+O(1).
\tag{20}
\]

Then the real series in (3) converges. Put

\[
\eta
=
x_0+\sum_{i=0}^{\infty}\frac{2^{A_i}}{3^{i+1}}
>0.
\]

Writing

\[
P_N=2^{A_N}x_N,\qquad Q_N=3^N,
\]

we have

\[
\frac{P_N}{Q_N}=h_N\longrightarrow\eta.
\]

The omitted tail is positive, so

\[
\left|\eta-\frac{P_N}{Q_N}\right|
=
\sum_{i=N}^{\infty}\frac{2^{A_i}}{3^{i+1}}
\geq
\frac{2^{A_N}}{3Q_N}.
\tag{21}
\]

Meanwhile,

\[
|P_N|_2=2^{-A_N},
\qquad
|Q_N|_3=Q_N^{-1}.
\]

Multiplying (21) by the two local factors gives

\[
\boxed{
\left|\eta-\frac{P_N}{Q_N}\right|_\infty
|P_N|_2
|Q_N|_3
\geq
\frac1{3Q_N^2}.
}
\tag{22}
\]

Because \(P_N/Q_N\to\eta>0\), the projective height of \((P_N,Q_N)\) is comparable to \(Q_N\). Thus (22) is exactly the critical scale

\[
H(P_N,Q_N)^{-2},
\]

not

\[
H(P_N,Q_N)^{-2-\varepsilon}.
\]

The projective Subspace Theorem in dimension one requires precisely the latter strict power improvement. [\hyperlink{ref:R53:3}{3}]

There is an additional problem: nothing in the trajectory identity proves that the real limit \(\eta\) is algebraic. An algebraic coefficient would be needed even to form the relevant Subspace-Theorem linear form \(X-\eta Y\).

So the uniformly escaping case evades the method twice:

\[
\text{unknown algebraicity of }\eta,
\qquad
\text{and exact critical exponent even if }\eta\text{ were algebraic}.
\]


\bigskip



\subsubsection*{6. Why ``the dynamics becomes repetitive'' would still not finish Collatz}


Suppose a valuation word \(w=(a_0,\ldots,a_{p-1})\) repeats periodically. Let

\[
B=\sum_{i=0}^{p-1}a_i,
\qquad
C_w=\sum_{i=0}^{p-1}3^{p-1-i}2^{A_i}.
\]

The unique \(2\)-adic point with this periodic valuation word satisfies

\[
2^B x=3^p x+C_w,
\]

and hence

\[
\boxed{
(2^B-3^p)x=C_w.
}
\tag{23}
\]

If \(x>0\), then necessarily

\[
2^B>3^p.
\]

A nontrivial positive Collatz cycle is therefore already a fixed linear dependence of exactly the kind the proposed argument hopes the Subspace Theorem will discover.

Repeating the period produces infinitely many longer trajectory identities, but they are generated by (23). The conclusion ``all these points lie in a fixed subspace'' is true and supplies no contradiction.

Equation (23) is precisely the divisibility formulation of the weak Collatz conjecture: one must prove that the only positive integral solution arising from a legitimate valuation word is the trivial word \(a_0=2\), \(x=1\). Tao's exposition gives this equivalent cycle formulation and explains why the general carry sum has unboundedly many powers. [\hypertarget{cit:R53:5}{}\hyperlink{ref:R53:5}{5}, \hyperlink{ref:R53:2}{2}]

Thus even an ideal Subspace-Theorem argument forcing eventual periodicity would still leave the nontrivial-cycle problem.


\bigskip



\subsubsection*{7. The exact missing rigidity theorem}


The divergent-orbit half can be isolated cleanly.

Let

\[
\alpha=\log_2 3
\]

and, for a positive-integer word \(a=(a_0,a_1,\ldots)\), define

\[
\Phi(a)
=
-\sum_{i=0}^{\infty}\frac{2^{A_i}}{3^{i+1}}
\in\mathbb Z_2.
\]

Consider the following statement.


\begin{quote}
\textbf{Critical rationality-rigidity principle.}   If \[
> \Phi(a)\in\mathbb Q
> \qquad\text{and}\qquad
> \limsup_{N\to\infty}\frac{A_N}{N}\leq\log_2 3,
> \] then \(a\) is eventually periodic.

\end{quote}


This is a restricted, critical-drift version of the rationality/periodicity problem for the \(3x+1\) conjugacy. The broader assertion that rational \(2\)-adic inputs have eventually periodic parity sequences is Lagarias' Periodicity Conjecture. [\hypertarget{cit:R53:6}{}\hyperlink{ref:R53:6}{6}]


\paragraph*{Theorem 2: exact conditional completion}


Assume:

\begin{enumerate}
\item the critical rationality-rigidity principle above;
\item the weak Collatz conjecture: \(1\) is the only positive integral periodic point of \(U\).
\end{enumerate}

Then the full Collatz conjecture follows.


\subparagraph*{Proof}


Suppose otherwise and choose a least counterexample \(x_0\). Its accelerated odd orbit either repeats or does not.

If it repeats, it is a nontrivial positive integral cycle, contradicting assumption 2.

Otherwise, \(x_j\to\infty\), and equation (17) gives

\[
\limsup_{N\to\infty}\frac{A_N}{N}
\leq\log_2 3.
\]

Equation (4) gives

\[
\Phi(a)=x_0\in\mathbb Q.
\]

Assumption 1 therefore implies that \(a\) is eventually periodic. Let the eventual period have length \(p\), total exponent \(B\), and carry \(C>0\). The periodic tail state \(x>0\) satisfies

\[
(2^B-3^p)x=C.
\]

Consequently \(2^B>3^p\), so

\[
\frac Bp>\log_2 3.
\]

But for an eventually periodic sequence,

\[
\lim_{N\to\infty}\frac{A_N}{N}=\frac Bp,
\]

contradicting the preceding upper bound. Therefore no counterexample exists. \(\square\)

This conditional proof is complete, but both assumptions contain unresolved arithmetic: assumption 1 excludes divergent rational trajectories, while assumption 2 excludes nontrivial cycles.


\bigskip



\subsubsection*{8. A genuine partial rigidity result}


The Subspace-Theorem philosophy does work when the valuation sequence has sufficiently long periodic approximations.

Suppose \(a\) is the valuation sequence of an integer \(x_0\), and there are eventually periodic words \(b^{(r)}\) such that:

\[
a_0a_1\cdots a_{L_r-1}
=
b^{(r)}_0b^{(r)}_1\cdots b^{(r)}_{L_r-1},
\]

the periodic description of \(b^{(r)}\) has length \(d_r\), its total exponent is \(O(d_r)\), and

\[
\frac{L_r}{d_r}\longrightarrow\infty.
\]

Let \(z_r=\Phi(b^{(r)})\). The formula for an eventually periodic series shows

\[
H(z_r)\leq \exp(O(d_r)).
\]

Agreement through \(L_r\) valuations gives

\[
|x_0-z_r|_2\leq2^{-A_{L_r}}\leq2^{-L_r}.
\tag{24}
\]

Every \(z_r\) has odd denominator. For distinct rationals \(x_0\) and \(z_r=p_r/q_r\),

\[
|x_0-z_r|_2
=
2^{-\nu_2(x_0q_r-p_r)}
\geq
\frac{1}{(|x_0|+1)H(z_r)}.
\tag{25}
\]

Combining (24) and (25),

\[
2^{L_r}\leq (|x_0|+1)\exp(O(d_r)),
\]

which contradicts \(L_r/d_r\to\infty\). Hence eventually \(z_r=x_0\), and uniqueness of the valuation coding implies that \(a\) itself is eventually periodic.

So the approach rigorously excludes counterexamples with arbitrarily strong ``stammering'' or low-description periodic approximations. What it does not control is a high-complexity critical word whose new information grows linearly with its length.


\bigskip



\subsection*{Final assessment}


The proposed route achieves a sharp reduction, but not a full proof:

\[
\boxed{
\text{standard Subspace Theorem}
\quad\Longrightarrow\quad
\text{rigidity only with fixed dimension and an }\varepsilon\text{-margin}.
}
\]

A minimal divergent Collatz trajectory instead gives

\[
\boxed{
\text{dimension}\to\infty,
\qquad
\text{or fixed-dimensional compression with only }H^{-o(1)},
\qquad
\text{or an approximation exactly at }H^{-2}.
}
\]

The needed new result would be a \textbf{critical, growing-dimensional staircase Subspace Theorem}: it would have to exploit the nested exponent path

\[
(0,0),(A_1,1),(A_2,2),\ldots
\]

and prove that rational \(2\)-adic cancellation at the critical boundary forces eventual periodicity. Classical and moving-target Subspace Theorems do not supply that statement. Even after obtaining it, equation (23)\textemdash{}the nontrivial positive-cycle problem\textemdash{}would still have to be excluded.

Accordingly, presenting approach 48 as a completed Collatz proof would require asserting two unproved rigidity results. The calculations above identify them precisely rather than hiding them inside an unjustified ``the Subspace Theorem forces repetition'' step.
\subsection*{References}
\begin{enumerate}[leftmargin=2.2em]
\item \hypertarget{ref:R53:1}{}\hyperlink{cit:R53:1}{\texttt{turn993328view0}}
\item \hypertarget{ref:R53:2}{}\hyperlink{cit:R53:2}{\texttt{turn993328view3}}
\item \hypertarget{ref:R53:3}{}\hyperlink{cit:R53:3}{\texttt{turn353191view0}}
\item \hypertarget{ref:R53:4}{}\hyperlink{cit:R53:4}{\texttt{turn993328view2}}
\item \hypertarget{ref:R53:5}{}\hyperlink{cit:R53:5}{\texttt{turn640859view1}}
\item \hypertarget{ref:R53:6}{}\hyperlink{cit:R53:6}{\texttt{turn993328view1}}
\end{enumerate}

\clearpage
\section{Round 54: approach 49 --- Baker theory for repeated approximate returns}
\subsection*{Prompt}
\begin{quote}
Complete a full proof of the Collatz conjecture using the following approach:

\begin{enumerate}
\item Baker theory for repeated approximate returns
\end{enumerate}

\textbf{Hybrid \textperiodcentered{} P}

Cycle work controls one exceptionally small quantity |A log 2 - j log 3|. A divergent orbit with recurrent near-descents may create infinitely many related small forms. Prove that their shared coefficients or nesting make them collectively impossible, even when each individual approximation is allowed.

Take as long as you need.
\end{quote}
\subsection*{Response}

\subsection*{Baker rigidity for repeated Collatz returns}


After carrying this program through, the honest conclusion is:


\begin{quote}
\textbf{Baker theory plus nesting yields a strong collective exclusion theorem, but it does not complete the Collatz conjecture.}

\end{quote}


The obstruction is precise. Baker theory controls a contractive coefficient block once that block is also a real near-return. Two such blocks satisfy a much stronger determinant identity. However, a hypothetical divergent orbit is not presently known to produce sufficiently short contractive near-returns. In fact, there are arbitrarily long genuine finite Collatz trajectories satisfying every coefficient-barrier condition while producing no such returns. The missing ingredient is therefore an ordinary-integer/2-adic synchronization theorem, not a sharper lower bound for one logarithmic form.

Below is the complete valid argument and the exact theorem that would finish the proof.


\bigskip



\subsubsection*{1. Exact logarithmic form of a return}


Use the accelerated map on positive odd integers

\[
U(x)=\frac{3x+1}{2^{v_2(3x+1)}}.
\]

Let

\[
x_{i+1}=U(x_i),\qquad
a_i=v_2(3x_i+1).
\]

For a block beginning at \(x_s=X\) and ending at \(x_{s+r}=Y\), put

\[
A=\sum_{i=s}^{s+r-1}a_i,
\qquad
M=\min_{s\leq i<s+r}x_i.
\]

The usual affine identity is

\[
2^A Y=3^rX+C,
\qquad
C>0.
\]

More explicitly, valuation words determine exact affine maps and exact realization classes modulo \(2^{A+1}\). [\hypertarget{cit:R54:1}{}\hyperlink{ref:R54:1}{1}]

At each step,

\[
\frac{x_{i+1}}{x_i}
=
\frac{3}{2^{a_i}}
\left(1+\frac1{3x_i}\right).
\]

Therefore

\[
\frac YX
=
\frac{3^r}{2^A}
\prod_{i=s}^{s+r-1}
\left(1+\frac1{3x_i}\right).
\]

Define

\[
\Lambda=A\log 2-r\log 3
\]

and

\[
E=\sum_{i=s}^{s+r-1}
\log\left(1+\frac1{3x_i}\right).
\]

Taking logarithms gives the exact return identity

\[
\boxed{\Lambda=E-\log(Y/X).}
\tag{1}
\]

The coefficient \(3^r/2^A\) is contractive precisely when \(\Lambda>0\).

Let

\[
\rho=\log^+\frac XY
=
\max\left(0,\log\frac XY\right).
\]

Since \(\log(1+u)\leq u\),

\[
E\leq \frac r{3M}.
\]

Thus every contractive block satisfies

\[
\boxed{
0<\Lambda
\leq
\frac r{3M}+\rho.
}
\tag{2}
\]

For an \textbf{upper return}, meaning \(Y\geq X\), this simplifies to

\[
\boxed{
0< A\log2-r\log3
\leq \frac r{3M}.
}
\tag{3}
\]

More generally, if \(Y\geq X-\Delta\), then

\[
\rho
\leq
\log\left(1+\frac{\Delta}{Y}\right)
\leq \frac{\Delta}{M},
\]

and hence

\[
0<\Lambda\leq\frac{r/3+\Delta}{M}.
\tag{4}
\]

This is the exact point where a real near-return creates a small linear form.


\bigskip



\subsubsection*{2. What one application of Baker theory gives}


Rozier and Terracol use the following effective form of a theorem of Rhin:

\[
\left|u_0+u_1\log2+u_2\log3\right|
\geq H^{-13.3},
\qquad
H=\max(|u_1|,|u_2|)\geq2.
\tag{5}
\]

[\hypertarget{cit:R54:2}{}\hyperlink{ref:R54:2}{2}]

Apply this with

\[
u_0=0,\qquad u_1=A,\qquad u_2=-r.
\]

For an upper return, (3) gives \(\Lambda\leq r/3\). Consequently,

\[
A\log2=r\log3+\Lambda
\leq r\left(\log3+\frac13\right),
\]

so \(A<3r\). Therefore \(H\leq3r\), and Rhin's bound gives

\[
(3r)^{-13.3}\leq \Lambda\leq\frac r{3M}.
\]

Hence:


\paragraph*{Theorem 1 \textemdash{} isolated-return Baker bound}


Every contractive upper return of length \(r\) and block floor \(M\) satisfies

\[
\boxed{
M\leq 3^{12.3}r^{14.3}.
}
\tag{6}
\]

Equivalently,

\[
r\gg M^{1/14.3}.
\tag{7}
\]

This is unconditional.

It says that a large orbit value cannot participate in a very short paradoxical return. But a polynomial lower bound on \(r\) does not itself exclude a divergent orbit: the orbit may wait an arbitrarily long time before returning.

Even replacing \(13.3\) by an essentially optimal irrationality exponent would only move the natural threshold toward \(r\asymp\sqrt M\). It would not by itself force nonexistence.


\bigskip



\subsubsection*{3. Two small forms are much more rigid than one}


The main collective gain does not actually require a second invocation of Baker theory.

Consider two contractive near-return blocks with data

\[
(A_i,r_i,M_i,\rho_i),\qquad i=1,2,
\]

and put

\[
\Lambda_i=A_i\log2-r_i\log3,
\qquad
\varepsilon_i=\frac{r_i}{3M_i}+\rho_i.
\]

Then

\[
0<\Lambda_i\leq\varepsilon_i.
\]

Form the integer determinant

\[
D=A_1r_2-A_2r_1.
\]

The \(\log3\)-terms cancel exactly:

\[
r_2\Lambda_1-r_1\Lambda_2
=
D\log2.
\tag{8}
\]

If \(D\neq0\), then \(|D|\geq1\), and therefore

\[
\log2
\leq |D|\log2
=
|r_2\Lambda_1-r_1\Lambda_2|
\leq r_2\varepsilon_1+r_1\varepsilon_2.
\]

Thus:


\paragraph*{Theorem 2 \textemdash{} collective determinant rigidity}


Two contractive near-return blocks with nonproportional coefficient vectors satisfy

\[
\boxed{
\log2
\leq
r_2\left(\frac{r_1}{3M_1}+\rho_1\right)
+
r_1\left(\frac{r_2}{3M_2}+\rho_2\right).
}
\tag{9}
\]

For two upper returns,

\[
\boxed{
\log2
\leq
\frac{r_1r_2}{3}
\left(\frac1{M_1}+\frac1{M_2}\right),
}
\tag{10}
\]

unless

\[
A_1r_2=A_2r_1.
\]

At comparable height \(M\), two genuinely different coefficient directions therefore require

\[
r_1r_2\gg M.
\tag{11}
\]

In particular, two upper returns of lengths \(o(\sqrt M)\) cannot have different ratios \(A_i/r_i\).

This is stronger than applying a one-form Baker bound twice.


\bigskip



\subsubsection*{4. An infinite short-return chain is impossible}


Suppose there were an infinite sequence of contractive near-return blocks \(I_i\) such that

\[
\frac{\varepsilon_i}{r_i}\longrightarrow0
\tag{12}
\]

and

\[
r_{i+1}\varepsilon_i+r_i\varepsilon_{i+1}
\longrightarrow0.
\tag{13}
\]

By Theorem 2, for all sufficiently large \(i\),

\[
A_ir_{i+1}=A_{i+1}r_i.
\]

Because consecutive vectors are proportional, there are fixed coprime positive integers \(p,q\) and integers \(m_i\geq1\) such that

\[
(A_i,r_i)=m_i(p,q)
\]

for every sufficiently large \(i\). Hence

\[
\Lambda_i
=
m_i(p\log2-q\log3).
\]

Since \(\Lambda_i>0\),

\[
\delta:=p\log2-q\log3>0.
\]

But

\[
\frac{\Lambda_i}{r_i}
=
\frac{\delta}{q}
\leq
\frac{\varepsilon_i}{r_i}
\longrightarrow0,
\]

which is impossible.

Therefore:


\paragraph*{Corollary 3 \textemdash{} exclusion of collectively short returns}


There is no infinite chain of contractive near-returns satisfying (12)\textendash{}(13).

For upper returns, it is enough that

\[
M_i\longrightarrow\infty
\]

and

\[
\boxed{
r_ir_{i+1}
=
o\!\left(\min(M_i,M_{i+1})\right).
}
\tag{14}
\]

So an infinite chain of upper returns at height \(M_i\) cannot have consecutive lengths \(o(\sqrt{M_i})\).

This proves exactly the ``collectively impossible'' part proposed in the approach.


\bigskip



\subsubsection*{5. Continued fractions explain the \(\sqrt M\) threshold}


Let

\[
\alpha=\frac{\log3}{\log2}.
\]

For an upper return,

\[
0<
\frac Ar-\alpha
=
\frac{\Lambda}{r\log2}
\leq
\frac1{3M\log2}.
\tag{15}
\]

If

\[
2r^2<3M\log2,
\tag{16}
\]

then

\[
\left|\frac Ar-\alpha\right|<\frac1{2r^2}.
\]

By Legendre's theorem, the reduced fraction \(A/r\) must be a continued-fraction convergent of \(\alpha\), from above.

Thus returns shorter than roughly \(\sqrt M\) are forced into the sparse set of best rational approximants to \(\log_2 3\). This is consistent with recent computations: the coefficient ratios in the currently known paradoxical examples are convergents, semiconvergents, or one nearby Stern\textendash{}Brocot mediant. That observation is numerical rather than an all-length theorem, and the authors explicitly do not claim progress on Collatz or the coefficient-stopping-time conjecture. [\hypertarget{cit:R54:3}{}\hyperlink{ref:R54:3}{3}]

Continued fractions also show why the determinant theorem is close to optimal. Distinct best approximants naturally have determinant \(1\), and their errors occur at the reciprocal-product scale \(1/(r_1r_2)\). Merely having infinitely many related small forms is therefore arithmetically possible.


\bigskip



\subsubsection*{6. Consequences for a hypothetical divergent orbit}


For the shortcut map, write

\[
q_k=\frac{3^{s_k}}{2^k},
\qquad
h_k=\frac{x_k}{q_k}.
\]

Along a divergent orbit, the reciprocal correction sum converges, \(h_k\) approaches a finite positive limit, and after rebasing at a global minimum of \(q_k\), one obtains

\[
q_k\geq1\quad\text{for all }k,
\qquad
q_k\longrightarrow\infty.
\]

This barrier condition is equivalent to the divergent-orbit alternative, but it does not address nontrivial cycles. [\hypertarget{cit:R54:4}{}\hyperlink{ref:R54:4}{4}]

At accelerated odd times the global coefficient form is

\[
A_k\log2-k\log3=-\log q_k.
\]

Hence, on the normalized divergent orbit,

\[
A_k\log2-k\log3\longrightarrow-\infty.
\tag{17}
\]

This is crucial:


\begin{quote}
The global prefix forms of a divergent barrier orbit are not small positive forms. They eventually become large negative forms.

\end{quote}


Therefore Baker theory can be used only after one proves the existence of suitable \textbf{local} blocks on which the coefficient contracts while the integer value approximately returns.

There is also a useful tail formulation. Along a divergent orbit let

\[
H=\lim h_k
\]

and define

\[
\eta_s=\log\frac{H}{h_s}.
\]

For every upper return beginning at time \(s\),

\[
0<\Lambda\leq\eta_s,
\tag{18}
\]

because the correction accumulated over the block is at most the entire remaining correction tail.

Consequently every such return of length \(r\) obeys

\[
r\geq \frac13\,\eta_s^{-1/13.3},
\tag{19}
\]

and two nonproportional returns beginning at \(s_1,s_2\) satisfy

\[
\log2\leq r_2\eta_{s_1}+r_1\eta_{s_2}.
\tag{20}
\]

Thus any divergent counterexample must exhibit a severe ``return desert'': either contractive upper returns cease to occur, or their lengths grow at least at the scales dictated by (19)\textendash{}(20).

That is a meaningful structural restriction, but not yet a contradiction.


\bigskip



\subsubsection*{7. The recurrence implication is false in the symbolic relaxation}


The central missing implication would be:


\begin{quote}
A long barrier trajectory arising from one positive integer must contain short contractive approximate returns.

\end{quote}


This is not a consequence of coefficient inequalities, nesting, or finite realizability.

Consider the two valuation blocks

\[
u=(1,1,2),
\qquad
v=(1,2,1).
\]

Both have length \(3\), total exponent \(4\), and coefficient

\[
\frac{3^3}{2^4}=\frac{27}{16}>1.
\]

Their prefix coefficients are

\[
u:\quad \frac32,\ \frac94,\ \frac{27}{16},
\]

and

\[
v:\quad \frac32,\ \frac98,\ \frac{27}{16}.
\]

Every prefix coefficient is at least \(1\). Therefore \textbf{every concatenation of \(u\) and \(v\) is a barrier word}, and its coefficient grows by \(27/16\) per macroblock.

The exact affine maps are

\[
16U^3(x)=27x+19
\qquad\text{for }u,
\tag{21}
\]

and

\[
16U^3(x)=27x+23
\qquad\text{for }v.
\tag{22}
\]

Every finite valuation word is realized by infinitely many positive odd integers. [\hypertarget{cit:R54:5}{}\hyperlink{ref:R54:5}{5}]

Because every prefix coefficient is at least \(1\) and every affine intercept is positive, each finite realization remains strictly above its starting value throughout the prescribed word.

Now inspect an arbitrary factor of an infinite \(u/v\)-concatenation. It consists of

\begin{itemize}
\item a suffix of one macroblock,
\item some number \(N\) of complete macroblocks,
\item a prefix of one macroblock.
\end{itemize}

Its logarithmic coefficient form is

\[
\Lambda
=
N\log\frac{16}{27}+\beta,
\tag{23}
\]

where \(\beta\) belongs to a fixed finite set determined by the two boundary fragments. Since

\[
\log\frac{16}{27}<0,
\]

the condition \(\Lambda>0\) forces \(N\) to be bounded. Consequently:

\[
\{\Lambda>0:\Lambda\text{ occurs in a factor}\}
\]

is a finite set and therefore has a positive minimum.

It follows that these arbitrarily long barrier trajectories have \textbf{no logarithmic forms tending to zero from above}. By choosing their positive realization sufficiently large, none of their contractive factors can be an upper return, because (3) would require a fixed positive \(\Lambda\) to be at most \(r/(3M)\).

So there are:


\begin{quote}
Arbitrarily long genuine positive Collatz trajectory segments, with arbitrarily large starting values, satisfying every prefix barrier and containing no high-level contractive upper return.

\end{quote}


The periodic ray \(u^\infty\) corresponds to the negative rational fixed point

\[
16x=27x+19,
\qquad
x=-\frac{19}{11}.
\]

An aperiodic \(u/v\)-ray produces a nested \(2\)-adic realization. The issue is whether such a nested realization can stabilize to one positive ordinary integer. A \(2\)-adic ray represents a nonnegative ordinary integer precisely when its least residue representatives eventually stabilize. [\hypertarget{cit:R54:6}{}\hyperlink{ref:R54:6}{6}]

Thus the symbolic example does not disprove Collatz. It proves something more targeted:


\begin{quote}
\textbf{No argument using only coefficient forms, finite nesting, and arbitrarily long positive realizations can force the required repeated approximate returns.}

\end{quote}


One must exploit the fact that all depths arise from one fixed ordinary integer.


\bigskip



\subsubsection*{8. Why existing ``infinitely many paradoxical returns'' do not close the gap}


Rozier and Terracol proved that if an integer \(n\) has infinite stopping time, then continued-fraction approximations to \(\log 2/\log 3\) produce infinitely many paradoxical sequences starting from integers of the form

\[
2^k n.
\]

Their construction first performs \(k\) halving steps and then follows the orbit of \(n\). [\hypertarget{cit:R54:7}{}\hyperlink{ref:R54:7}{7}]

This gives infinitely many small forms, but it does not satisfy the high-floor hypothesis needed above:

\[
2^kn,\ 2^{k-1}n,\ldots,n,\ldots
\]

has block minimum at most \(n\), not \(2^kn\).

Therefore one cannot insert the growing first term \(2^kn\) into the bound

\[
\Lambda\leq\frac r{3M}.
\]

The correct \(M\) remains bounded by \(n\). The continued-fraction forms, typically of order \(1/r\), are entirely compatible with Rhin's lower bound \(r^{-13.3}\).

This is exactly how infinitely many individually allowed approximations evade Baker theory.


\bigskip



\subsubsection*{9. Repeated cores are the part Baker theory does solve}


When recurrence has bounded symbolic complexity, the method becomes effective.

For a fixed repeated core \(v^N\), the affine matrix can be diagonalized explicitly. Cyclotomic factors amplify the changing determinant, and a fixed or height-limited exit produces an exponentially small nonzero form in \(\log2,\log3\). Baker theory then gives an effective bound on \(N\). The same framework proves rational-center \(2\)-adic height growth and excludes pure powers and fixed-tail, bounded-defect exits. [\hypertarget{cit:R54:8}{}\hyperlink{ref:R54:8}{8}]

So the approach succeeds for:

\[
v^N,
\qquad
v^Nu
\quad\text{with fixed }u,
\]

and related bounded-grammar families.

The unresolved case is an aperiodic sequence with a moving exit and unbounded carry complexity. There, the affine numerators change dimension with the block length, so one no longer has a fixed algebraic number to which a standard Baker estimate applies.


\bigskip



\subsubsection*{10. Cycles remain logically separate}


For a positive accelerated cycle, \(Y=X\), and (1) gives

\[
\Lambda=E>0.
\]

Repeating the cycle \(N\) times produces

\[
(A_N,r_N)=N(A,r),
\qquad
\Lambda_N=N\Lambda.
\]

All coefficient vectors are proportional, so every determinant in Theorem 2 is zero. Moreover, the form becomes larger rather than smaller.

Thus collective near-return theory does not eliminate positive cycles.

Baker and continued-fraction methods have been very effective for cycles of bounded structural complexity; Simons and de Weger developed the general \(m\)-cycle framework, and later work excludes \(m\)-cycles through \(m=91\). It still does not constitute an all-cycle exclusion theorem. [\hypertarget{cit:R54:9}{}\hyperlink{ref:R54:9}{9}, \hypertarget{cit:R54:10}{}\hyperlink{ref:R54:10}{10}]


\bigskip



\subsubsection*{11. The exact conditional completion theorem}


The Baker-return route would become a full proof with the following additional result.


\paragraph*{Return-synchronization hypothesis}


Every divergent positive Collatz orbit contains contractive upper-return blocks \(I_i\) with lengths \(r_i\) and floors \(M_i\) such that

\[
M_i\longrightarrow\infty
\]

and

\[
r_ir_{i+1}
\left(
\frac1{M_i}+\frac1{M_{i+1}}
\right)
\longrightarrow0.
\tag{24}
\]


\paragraph*{Theorem 4 \textemdash{} conditional Collatz completion}


Assume:

\begin{enumerate}
\item the return-synchronization hypothesis;
\item there is no nontrivial positive Collatz cycle.
\end{enumerate}

Then every positive integer reaches \(1\).


\paragraph*{Proof}


A nonterminating positive orbit is either bounded or unbounded.

If bounded, determinism forces repetition, hence a positive cycle, contradicting assumption 2.

If unbounded and nonperiodic, every bounded set is eventually left permanently, so the orbit tends to infinity. Assumption 1 supplies upper returns satisfying (24). But Corollary 3 excludes exactly such a sequence. Therefore no unbounded nonterminating orbit exists.

Hence every orbit terminates in the trivial cycle. \(\square\)


\bigskip



\subsubsection*{12. Final verdict}


Approach 49 produces the following rigorous advances:

\[
\boxed{
M\leq3^{12.3}r^{14.3}
}
\]

for one contractive upper return, and the substantially stronger collective relation

\[
\boxed{
\log2
\leq
\frac{r_1r_2}{3}
\left(\frac1{M_1}+\frac1{M_2}\right)
}
\]

for two nonproportional upper returns.

It follows that no divergent orbit can possess an infinite chain of contractive upper returns with consecutive lengths \(o(\sqrt M)\). Repeated fixed cores and bounded exits are also effectively finite by Baker theory.

But this is not a full proof of Collatz. The proposed final inference\textemdash{}

\[
\text{divergence}
\Longrightarrow
\text{infinitely many collectively short contractive near-returns}
\]

\textemdash{}is neither presently proved nor a consequence of coefficients and nesting. The explicit \(u/v\) construction shows that the relaxed symbolic system supports arbitrarily long positive barrier trajectories with no such returns.

The precise missing theorem is therefore:

\[
\boxed{
\begin{minipage}{0.86\linewidth}
A barrier valuation ray whose nested \(2\)-adic realization is one fixed
positive ordinary integer must either become periodic or generate
contractive upper returns satisfying the collective shortness condition.
\end{minipage}
}
\]

Proving that carry-sensitive ordinary-integer synchronization theorem, together with all-cycle exclusion, would complete the Baker-return route. Recent work on paradoxical sequences likewise treats the coefficient-stopping-time and Collatz conclusions as unresolved, rather than supplying this synchronization step. [\hyperlink{ref:R54:3}{3}]
\subsection*{References}
\begin{enumerate}[leftmargin=2.2em]
\item \hypertarget{ref:R54:1}{}\hyperlink{cit:R54:1}{\texttt{turn6file0 L30-L76}}
\item \hypertarget{ref:R54:2}{}\hyperlink{cit:R54:2}{\texttt{turn417996search0}}
\item \hypertarget{ref:R54:3}{}\hyperlink{cit:R54:3}{\texttt{turn513292view0}}
\item \hypertarget{ref:R54:4}{}\hyperlink{cit:R54:4}{\texttt{turn6file2 L248-L323}}
\item \hypertarget{ref:R54:5}{}\hyperlink{cit:R54:5}{\texttt{turn6file1 L160-L183}}
\item \hypertarget{ref:R54:6}{}\hyperlink{cit:R54:6}{\texttt{turn3file4 L183-L192}}
\item \hypertarget{ref:R54:7}{}\hyperlink{cit:R54:7}{\texttt{turn494327view0}}
\item \hypertarget{ref:R54:8}{}\hyperlink{cit:R54:8}{\texttt{turn2file1 L339-L414}}
\item \hypertarget{ref:R54:9}{}\hyperlink{cit:R54:9}{\texttt{turn675823search4}}
\item \hypertarget{ref:R54:10}{}\hyperlink{cit:R54:10}{\texttt{turn675823search1}}
\end{enumerate}

\clearpage
\section{Round 55: approach 50 --- Geometry-of-numbers exclusion cones}
\subsection*{Prompt}
\begin{quote}
Complete a full proof of the Collatz conjecture using the following approach:

\begin{enumerate}
\item Geometry-of-numbers exclusion cones
\end{enumerate}

\textbf{Speculative \textperiodcentered{} B}

Embed exponent words as lattice points involving j, A\_j, and selected carry coordinates. Sustained growth confines them to a narrow cone, while integrality removes periodic sublattices. Show the remaining cone is lattice-free beyond a finite radius.

Take as long as you need.
\end{quote}
\subsection*{Response}

\subsubsection*{Outcome}


Carrying out the geometry-of-numbers construction reveals a rigorous obstruction: \textbf{the proposed ``lattice-free beyond a finite radius'' statement is false, even after imposing exact integrality, exact carry recurrences, and strict growth at every intermediate step}.

In fact, there are valid positive-integer trajectory prefixes of arbitrary length lying in a strip of width \(1\) around the critical line \(A=j\log _2 3\). Thus I cannot honestly present this approach as a completed proof. The Collatz conjecture remains unresolved in the current literature. [\hypertarget{cit:R55:1}{}\hyperlink{ref:R55:1}{1}]

What the approach does yield is:

\begin{enumerate}
\item an exact affine-lattice formulation;
\item an explicit counterexample to the proposed lattice-free lemma;
\item a proof that finitely many carry coordinates provide no independent determinant gain;
\item a corrected \textbf{nested-coset cusp criterion} exactly equivalent to Collatz.
\end{enumerate}


\bigskip



\subsection*{1. Exact trajectory lattices}


Work with the odd-only Syracuse map

\[
U(x)=\frac{3x+1}{2^{v_2(3x+1)}},
\qquad x\in 2\mathbb N+1.
\]

Let

\[
a_i=v_2(3U^i(n)+1),\qquad
A_k=\sum_{i=0}^{k-1}a_i,\qquad A_0=0.
\]

For a finite exponent word

\[
w=(a_0,\ldots ,a_{j-1})\in\mathbb Z_{>0}^{j},
\]

define its carries by

\[
C_0=0,\qquad C_{k+1}=3C_k+2^{A_k}.
\]

Equivalently,

\[
C_k=\sum_{i=0}^{k-1}3^{k-1-i}2^{A_i}.
\]

Repeated substitution gives the exact identity

\[
\boxed{\;
2^{A_k}U^k(n)=3^k n+C_k.
\;}
\tag{1}
\]

Thus, for fixed \(w\), its possible trajectory vectors are the integral points of the affine rank-one lattice

\[
\Lambda_w=
\left\{
(x_0,\ldots ,x_j)\in(2\mathbb Z+1)^{j+1}:
2^{a_i}x_{i+1}-3x_i=1
\right\}.
\tag{2}
\]

This is the natural geometry-of-numbers object.


\bigskip



\subsubsection*{Proposition 1: every finite exponent word is integrally realizable}


For every positive exponent word \(w\), there is a unique odd residue class modulo \(2^{A_j+1}\) whose first \(j\) Syracuse exponents are exactly \(w\).

More precisely, let \(r(w)\) be the unique representative with

\[
1\le r(w)<2^{A_j+1}
\]

satisfying

\[
\boxed{\;
r(w)\equiv
3^{-j}\bigl(2^{A_j}-C_j\bigr)
\pmod {2^{A_j+1}}.
\;}
\tag{3}
\]

Then the positive odd integers realizing \(w\) are exactly

\[
\boxed{\;
n=r(w)+t\,2^{A_j+1},
\qquad t\in\mathbb Z_{\ge 0}.
\;}
\tag{4}
\]


\paragraph*{Proof}


The final quotient in (1) is odd exactly when

\[
3^j n+C_j\equiv 2^{A_j}\pmod {2^{A_j+1}},
\]

which is equivalent to (3), since \(3\) is invertible modulo every power of \(2\).

It remains to show that this terminal congruence forces every earlier quotient to be odd. For \(k<j\),

\[
C_j
=
3^{j-k}C_k+
\sum_{i=k}^{j-1}3^{j-1-i}2^{A_i}.
\]

Since \(A_i\ge A_k+1\) for \(i>k\), reduction modulo \(2^{A_k+1}\) leaves only the \(i=k\) term:

\[
0\equiv
3^{j-k}(3^k n+C_k)
+
3^{j-1-k}2^{A_k}
\pmod {2^{A_k+1}}.
\]

After dividing by the odd factor \(3^{j-1-k}\), this yields

\[
3(3^k n+C_k)\equiv-2^{A_k}
\equiv2^{A_k}
\pmod {2^{A_k+1}}.
\]

Multiplying by the inverse of \(3\), which is odd, gives

\[
3^k n+C_k\equiv2^{A_k}\pmod {2^{A_k+1}}.
\]

Hence every intermediate quotient in (1) is an odd integer, proving that the exact valuations are \(a_0,\ldots ,a_{j-1}\). \ensuremath{\blacksquare}

This is the odd-only analogue of the standard parity-vector result: every finite parity word is realized by an entire congruence class of positive integers. [\hypertarget{cit:R55:2}{}\hyperlink{ref:R55:2}{2}]


\bigskip



\subsection*{2. The sustained-growth cone is saturated by lattice rays}


Let

\[
\alpha=\log_2 3.
\]

For a word \(w\), consider the growth cone

\[
\mathcal K_j=
\{(x_0,\ldots ,x_j):x_k\ge x_0\text{ for }1\le k\le j\}.
\]

From (1),

\[
U^k(n)-n
=
\frac{(3^k-2^{A_k})n+C_k}{2^{A_k}}.
\tag{5}
\]

Consequently:


\subsubsection*{Theorem 2: affine-ray obstruction}


Suppose

\[
A_k<\alpha k
\qquad(1\le k\le j).
\tag{6}
\]

Then \textbf{every positive integer realizing \(w\) grows strictly above its starting value through all \(j\) steps}:

\[
U^k(n)>n
\qquad(1\le k\le j).
\tag{7}
\]

Moreover, \(\Lambda_w\cap\operatorname{int}\mathcal K_j\) contains an infinite arithmetic ray.


\paragraph*{Proof}


Condition (6) means

\[
2^{A_k}<3^k.
\]

Since \(n>0\) and \(C_k>0\), equation (5) is strictly positive.

For the explicit ray, put \(r=r(w)\) and

\[
y_k=U^k(r).
\]

For

\[
n=r+t\,2^{A_j+1},
\]

equation (1) gives

\[
U^k(n)
=
y_k+t\,3^k2^{A_j+1-A_k}.
\tag{8}
\]

Thus the integral trajectory vectors form the ray

\[
\mathbf y+t\mathbf d,
\qquad
d_k=3^k2^{A_j+1-A_k},
\qquad t\ge0.
\]

Furthermore,

\[
d_k-d_0
=
2^{A_j+1-A_k}\bigl(3^k-2^{A_k}\bigr)>0,
\]

so its direction lies in the interior of the growth cone. \ensuremath{\blacksquare}

This is the opposite of lattice-freeness: every such exponent word produces an \textbf{infinite affine ray of exact integral trajectories inside the cone}.


\bigskip



\subsubsection*{A completely explicit family}


Take

\[
w=1^j=(1,\ldots ,1).
\]

Then \(A_k=k\) and

\[
C_k=3^k-2^k.
\]

Formula (3) gives

\[
r(1^j)=2^{j+1}-1.
\]

For \(1\le k\le j\),

\[
\boxed{\;
U^k(2^{j+1}-1)
=
3^k2^{j+1-k}-1.
\;}
\tag{9}
\]

Therefore

\[
U^k(2^{j+1}-1)-(2^{j+1}-1)
=
2^{j+1}\left(\left(\frac32\right)^k-1\right)>0.
\]

So arbitrary-length strictly growing exact prefixes exist even in closed form.

These prefixes do not combine into one positive orbit: their compatible residues converge \(2\)-adically to \(-1\). Indeed, the all-one infinite parity sequence corresponds to the \(2\)-adic integer \(-1\). [\hypertarget{cit:R55:3}{}\hyperlink{ref:R55:3}{3}]

That distinction\textemdash{}varying positive representatives versus one fixed positive seed\textemdash{}is exactly what a finite-dimensional cone fails to detect.


\bigskip



\subsection*{3. The obstruction remains in a width-one critical strip}


The preceding example lies well inside the expanding cone. The stronger obstruction is obtained by putting the exponent path immediately next to the critical boundary.

Define

\[
A_k=\lfloor \alpha k\rfloor,
\qquad
a_k=A_{k+1}-A_k.
\tag{10}
\]

Because \(1<\alpha<2\),

\[
a_k\in\{1,2\}.
\]

For every \(k\ge1\),

\[
0<\alpha k-A_k<1.
\tag{11}
\]

Thus the lattice points lie in a strip of width \(1\) around the critical line.

For each \(j\), let \(r_j\) be the canonical positive seed realizing the first \(j\) exponents from (10). Since

\[
2^{A_k}<3^k,
\]

Theorem 2 gives

\[
U^k(r_j)>r_j
\qquad(1\le k\le j).
\tag{12}
\]

The carry coordinate is also narrowly controlled. From

\[
\frac{3^i}{2}<2^{A_i}\le3^i
\]

we obtain

\[
\boxed{\;
\frac j2\,3^{j-1}
<
C_j
\le
j\,3^{j-1}.
\;}
\tag{13}
\]

Hence

\[
\log C_j
=
(j-1)\log3+\log j+O(1).
\]

We therefore have exact positive-integer trajectory prefixes of arbitrary depth satisfying simultaneously

\[
\left|A_j-\alpha j\right|<1,
\qquad
C_j\asymp j3^{j-1},
\qquad
U^k(r_j)>r_j\quad(k\le j).
\]

This directly disproves the proposed central statement:


\begin{quote}
``The remaining critical cone is lattice-free beyond a finite radius.''

\end{quote}


It is not lattice-free at any scale.


\bigskip



\subsection*{4. Why adding finitely many carry coordinates does not repair it}


The prefix congruences are \textbf{nested}, not independent.

For \(k<\ell\),

\[
C_\ell
=
3^{\ell-k}C_k+
2^{A_k}D_{k,\ell}
\tag{14}
\]

for an integer \(D_{k,\ell}\) determined by the suffix \(a_k,\ldots ,a_{\ell-1}\).

Consequently, the combined exact-integrality conditions at times

\[
1,2,\ldots ,j
\]

do not define a lattice of index roughly

\[
2^{A_1+A_2+\cdots+A_j}.
\]

They define just one congruence class modulo

\[
2^{A_j+1}.
\]

Equivalently, all \(j\) valuation conditions have total index

\[
2^{A_j}
\]

inside the odd integers. The earlier conditions are already encoded by the terminal one.

This destroys the determinant multiplication that a geometry-of-numbers exclusion would require.

There is also a general projection obstruction. Let

\[
L:\mathbb Z^{j+1}\longrightarrow\mathbb Z^m
\]

be any integer-linear collection of trajectory or carry coordinates. Applied to the ray from Theorem 2,

\[
L(\mathbf y+t\mathbf d)
=
L(\mathbf y)+tL(\mathbf d).
\]

Either:

\begin{itemize}
\item \(L(\mathbf d)\ne0\), in which case the projected cone still contains an unbounded arithmetic progression of lattice points; or
\item \(L(\mathbf d)=0\), in which case the projection collapses infinitely many distinct exact seeds and cannot determine whether one fixed positive seed persists.
\end{itemize}

Thus finitely many affine carry coordinates cannot convert the valid finite-prefix set into an eventually lattice-free set.

A large lattice determinant also says nothing about the location of an affine translate. The one-dimensional coset

\[
1+M\mathbb Z
\]

has covolume \(M\to\infty\) but always contains the bounded point \(1\). A hypothetical Collatz counterexample would manifest in exactly this way: a nested sequence of expanding-modulus affine lattices whose translates all contain one fixed small positive integer.


\bigskip



\subsection*{5. The correct object is a nested \(2\)-adic cusp}


For an infinite exponent word, let \(r_j\) be the canonical residue associated with its first \(j\) exponents.

Because every extension realizing \(j+1\) exponents also realizes the first \(j\),

\[
r_{j+1}\equiv r_j\pmod {2^{A_j+1}}.
\]

Consequently,

\[
\boxed{\;
r_{j+1}
=
r_j+t_j2^{A_j+1},
\qquad
0\le t_j<2^{a_j}.
\;}
\tag{15}
\]

In particular, \(r_j\) is nondecreasing.

Every infinite exponent word determines exactly one \(2\)-adic integer through these nested residue classes; equivalently, every infinite parity sequence corresponds to exactly one \(2\)-adic seed. [\hyperlink{ref:R55:3}{3}]

The crucial elementary dichotomy is:

\[
\boxed{\;
(r_j)\text{ is bounded}
\iff
(r_j)\text{ is eventually constant}
\iff
\text{the infinite word is realized by a positive ordinary integer}.
\;}
\tag{16}
\]

Indeed, bounded nondecreasing integers eventually stabilize. Conversely, if a positive integer \(n\) realizes the infinite word, then once

\[
2^{A_j+1}>n,
\]

its least positive representative modulo \(2^{A_j+1}\) is exactly \(n\).

So the positive-integer locus is not characterized by finite integrality. It is characterized by the \textbf{eventual vanishing of all lift digits \(t_j\)} in (15).

That is an infinite-support condition. No fixed collection of congruences or carry coordinates can see it.


\bigskip



\subsection*{6. An exact corrected criterion equivalent to Collatz}


Define the finite survivor minimum

\[
\mu_j
=
\min\left\{
n>1:
n\text{ odd and }U^k(n)\ge n
\text{ for }1\le k\le j
\right\}.
\tag{17}
\]

The set is nonempty: the seed \(2^{j+1}-1\) from (9) is one candidate.

The sequence \(\mu_j\) is nondecreasing.


\subsubsection*{Theorem 3: cusp criterion}


The Collatz conjecture is equivalent to

\[
\boxed{\;
\lim_{j\to\infty}\mu_j=\infty.
\;}
\tag{18}
\]


\paragraph*{Proof}


Assume first that (18) holds. Let \(n>1\) be odd. Choose \(j\) such that

\[
\mu_j>n.
\]

Then \(n\) cannot remain above itself through its first \(j\) odd iterates. Hence

\[
U^k(n)<n
\]

for some \(k\le j\). Strong induction on \(n\) now shows that every odd \(n\) reaches \(1\), and therefore every positive integer reaches \(1\).

Conversely, suppose Collatz holds. For every fixed odd \(n>1\), some iterate is \(1<n\), so it has a finite first descent below \(n\). For any fixed \(N\), there are only finitely many odd \(n\le N\); take the maximum of their first-descent times. Beyond that depth, no \(n\le N\) contributes to (17). Therefore \(\mu_j>N\) for all sufficiently large \(j\). Since \(N\) was arbitrary, \(\mu_j\to\infty\). \ensuremath{\blacksquare}

This identifies the precise theorem a repaired geometry-of-numbers approach would need:

\[
\boxed{
\begin{aligned}
&2^{a_i}x_{i+1}=3x_i+1,\qquad
x_i\in2\mathbb N+1,\\
&x_i\ge x_0>1\quad(1\le i\le j)
\end{aligned}
\quad\Longrightarrow\quad
x_0\ge g(j)
}
\tag{19}
\]

for some function \(g(j)\to\infty\).

But (19) is equivalent to Collatz; it is not a consequence of ordinary cone volume or lattice determinant estimates.


\bigskip



\subsubsection*{The survivor points do lie in a genuine cone}


There is still a useful cone bound. If

\[
x_i\ge x_0=n\ge3
\qquad(0\le i\le j),
\]

then

\[
\frac{x_j}{n}
=
\frac1{2^{A_j}}
\prod_{i=0}^{j-1}\left(3+\frac1{x_i}\right).
\]

Since \(x_j/n\ge1\),

\[
2^{A_j}
\le
\prod_{i=0}^{j-1}\left(3+\frac1{x_i}\right)
\le
\left(3+\frac1n\right)^j.
\]

Therefore

\[
\boxed{\;
A_j\le j\log_2\left(3+\frac1n\right)
\le
j\log_2\frac{10}{3}
\approx1.73697j.
\;}
\tag{20}
\]

Thus each finite survivor section has only finitely many exponent words.

For a divergent fixed orbit, \(x_i\to\infty\), so the average contribution of \(\log(1+1/(3x_i))\) vanishes and the upper slope approaches \(\log_2 3\). There is also a known parity-density rigidity result: L\'opez and Stoll prove that a divergent rational \(2\)-adic orbit must have lower parity density exactly \(\log 2/\log 3\). In odd-only coordinates, this forces \(A_j/j\) to approach \(\log_2 3\) along subsequences. [\hypertarget{cit:R55:4}{}\hyperlink{ref:R55:4}{4}]

This validates the idea of studying a critical cone for a divergent counterexample\textemdash{}but the width-one mechanical construction proves that \textbf{critical angular confinement plus finite integrality is still insufficient}.


\bigskip



\subsection*{7. What remains of the geometry-of-numbers heuristic}


There is a genuine entropy gap.

Fix

\[
A=\lfloor \alpha j\rfloor.
\]

The number of positive exponent words with total exponent \(A\) is

\[
\binom{A-1}{j-1}.
\]

Stirling's formula gives

\[
\binom{A-1}{j-1}
=
2^{(\sigma+o(1))j},
\]

where

\[
\sigma
=
\alpha H_2(1/\alpha)
=
\alpha\log_2\alpha
-
(\alpha-1)\log_2(\alpha-1)
\approx1.5056438879.
\]

The number of odd residue classes modulo \(2^{A+1}\) is \(2^A\), with exponential rate

\[
\alpha=\log_2 3\approx1.5849625007.
\]

Thus the exponent gap is

\[
\alpha-\sigma
\approx0.0793186128.
\tag{21}
\]

So critical words occupy an exponentially sparse set of residue classes. Heuristically, their least positive representatives should grow like roughly

\[
2^{0.0793186j}.
\]

That would be more than enough to prove escape from every fixed seed.

But sparsity does not imply exclusion. An exponentially sparse subset of residue classes can still always contain the class of \(27\), or of some other fixed integer. Ordinary lattice counting produces a main volume term plus an \(O(1)\) boundary or discrepancy error. Once the main term is below \(1\), the \(O(1)\) error can hide exactly the one exceptional seed that matters.

A successful repair would therefore require an unusually strong \textbf{inhomogeneous discrepancy theorem with no additive \(O(1)\) loss}, or an arithmetic theorem proving that the lift digits \(t_j\) in (15) are nonzero infinitely often on every nontrivial survivor path. That is substantially stronger than showing the cone has small volume.


\bigskip



\subsection*{Final verdict}


The proposed approach does not yield a full proof because its central geometric assertion is explicitly false:

\[
\boxed{\text{The sustained-growth critical cone contains exact lattice points at every depth.}}
\]

This remains true:

\begin{itemize}
\item inside a width-one strip around \(A=j\log_2 3\);
\item with exact affine carry identities;
\item with exact positive-integer realizability;
\item with strict growth at every prefix;
\item and after adding any fixed finite collection of affine carry coordinates.
\end{itemize}

The correct replacement is not an eventually lattice-free cone. It is the nested-coset cusp statement

\[
\mu_j\longrightarrow\infty,
\]

or equivalently the assertion that every nontrivial infinite growth-compatible branch has canonical representatives \(r_j\to\infty\). Proving that escape requires control of the infinite nesting\textemdash{}or eventual binary support\textemdash{}of one fixed seed. Finite-dimensional geometry of numbers, determinant growth, and deletion of periodic sublattices do not supply that missing control.
\subsection*{References}
\begin{enumerate}[leftmargin=2.2em]
\item \hypertarget{ref:R55:1}{}\hyperlink{cit:R55:1}{\texttt{turn720760view0}}
\item \hypertarget{ref:R55:2}{}\hyperlink{cit:R55:2}{\texttt{turn720760view1}}
\item \hypertarget{ref:R55:3}{}\hyperlink{cit:R55:3}{\texttt{turn720760view2}}
\item \hypertarget{ref:R55:4}{}\hyperlink{cit:R55:4}{\texttt{turn720760view3}}
\end{enumerate}

\clearpage
\section{Round 56: approach 51 --- Automorphic dynamics of the affine branch semigroup}
\subsection*{Prompt}
\begin{quote}
Complete a full proof of the Collatz conjecture using the following approach:

\begin{enumerate}
\item Automorphic dynamics of the affine branch semigroup
\end{enumerate}

\textbf{Speculative \textperiodcentered{} P}

Represent accelerated branches by matrices such as [[3, 1], [0, 2\textasciicircum{}a]] inside an adelic affine group. Use automorphic mixing, height recurrence, or nonescape of mass to control all integral products.

Take as long as you need.
\end{quote}
\subsection*{Response}

\subsection*{Outcome}


Carrying the automorphic program through does \textbf{not} currently produce a valid proof of the Collatz conjecture. Instead, it yields a fairly sharp no-go theorem for the approach in its unmarked form:


\begin{quote}
On the natural adelic automorphic quotient, the accelerated branch matrices are rational arithmetic elements, so their orbit is constant.   On a real modular quotient they become Hecke correspondences, but the seeded Collatz orbit collapses to its diagonal multiplier and loses the carry.   Retaining the carry requires an inverse congruence tower and a pointwise, growing-level contraction theorem much stronger than ordinary automorphic mixing.

\end{quote}


Current 2026 work continues to describe the pointwise Collatz conjecture as open; recent advances concern reductions, average behavior, and finite verification rather than a proof for every initial value. [\hypertarget{cit:R56:1}{}\hyperlink{ref:R56:1}{1}, \hypertarget{cit:R56:2}{}\hyperlink{ref:R56:2}{2}]

What follows is the complete derivation.


\bigskip



\subsubsection*{1. Exact affine branch products}


Work with the accelerated odd-to-odd map

\[
U(n)=\frac{3n+1}{2^{a(n)}},
\qquad
a(n)=v_2(3n+1),
\]

on positive odd integers.

For \(a\geq1\), put

\[
B_a=
\begin{pmatrix}
3&1\\
0&2^a
\end{pmatrix}.
\]

Its projective action is

\[
B_a\cdot x=\frac{3x+1}{2^a}.
\]

For a word

\[
w=(a_0,\ldots ,a_{j-1}),
\qquad
A_k=\sum_{i<k}a_i,\qquad A=A_j,
\]

write

\[
B_w=B_{a_{j-1}}\cdots B_{a_0}.
\]

A direct multiplication gives

\[
B_w=
\begin{pmatrix}
3^j&C_w\\
0&2^A
\end{pmatrix},
\]

where

\[
C_0=0,\qquad
C_{k+1}=3C_k+2^{A_k},
\]

and hence

\[
C_w=\sum_{k=0}^{j-1}3^{j-1-k}2^{A_k}.
\tag{1}
\]

If \(w\) is the initial valuation word of \(n\), then

\[
2^A U^j(n)=3^j n+C_w.
\tag{2}
\]

This is the exact affine identity that an automorphic formulation must preserve.


\paragraph*{Every finite word is genuinely realizable}


There is a stronger obstruction than merely saying that every \(B_w\) is an integral matrix.

\textbf{Proposition 1.} For every finite word \(w=(a_0,\ldots ,a_{j-1})\) of positive integers, there is exactly one odd residue class modulo \(2^{A+1}\) whose elements realize \(w\). It is determined by

\[
3^j n+C_w\equiv 2^A\pmod {2^{A+1}}.
\tag{3}
\]

Consequently, every finite word is the exact accelerated Collatz prefix of infinitely many positive integers.

\textbf{Proof.} Necessity follows from (2), since the endpoint \(U^j(n)\) must be odd.

Since \(3^j\) is invertible modulo \(2^{A+1}\), (3) has a unique solution. Because \(C_w\) is odd, that solution is odd.

It remains to prove that the final congruence forces all the intermediate valuations. Split

\[
w=(a_0,v),\qquad A=a_0+A(v).
\]

The carries satisfy

\[
C_w=3^{j-1}+2^{a_0}C_v.
\]

Thus (3) becomes

\[
3^{j-1}(3n+1)+2^{a_0}C_v
  \equiv 2^{a_0+A(v)}
  \pmod {2^{a_0+A(v)+1}}.
\]

When \(v\neq\varnothing\), \(C_v\) is odd and \(A(v)\geq1\). Reducing modulo \(2^{a_0+1}\) gives

\[
3^{j-1}(3n+1)\equiv 2^{a_0}\pmod {2^{a_0+1}}.
\]

Therefore

\[
v_2(3n+1)=a_0.
\]

After dividing by \(2^{a_0}\), the same congruence is obtained for the tail word \(v\). Induction completes the proof. \ensuremath{\blacksquare}

This means that no theorem which merely excludes ``most,'' ``generic,'' or ``sufficiently equidistributed'' finite integral products can prove Collatz. There are no locally forbidden branch products.


\bigskip



\subsubsection*{2. The branch group is a solvable Borel subgroup}


Normalize the branch as the affine transformation

\[
g_a(x)=\frac{3x+1}{2^a}.
\]

Let \(d_r(x)=rx\) and \(t_b(x)=x+b\). Then

\[
g_ag_{a+1}^{-1}=d_2,
\tag{4}
\]

and

\[
d_{2^a}g_a(x)=3x+1.
\tag{5}
\]

Writing \(h(x)=3x+1\), one obtains

\[
d_2hd_2^{-1}h^{-1}=t_1
\]

and then

\[
t_{-1}h=d_3.
\]

It follows that the group generated by the branches is exactly

\[
\mathbb Z[1/6]\rtimes \langle 2,3\rangle,
\tag{6}
\]

where the first component consists of translations and the second of multipliers \(2^r3^s\).

Thus its Zariski closure in \(\mathrm{PGL}_2\) is the upper Borel subgroup. It is solvable and metabelian, not Zariski-dense in \(\mathrm{PGL}_2\).

This already prevents a direct application of modern random-walk equidistribution theorems whose hypotheses require the support to generate a Zariski-dense subgroup. Those results also concern probability distributions of random products, rather than the one deterministic branch selected by an integer orbit. [\hypertarget{cit:R56:3}{}\hyperlink{ref:R56:3}{3}]

The failure is nevertheless stronger than the failure of a hypothesis: the natural automorphic orbit is exactly constant.


\bigskip



\subsubsection*{3. Adelic automorphic collapse}


Let

\[
G=\mathrm{PGL}_2,
\qquad
X=G(\mathbb Q)\backslash G(\mathbb A).
\]

Every \(B_w\) belongs to \(G(\mathbb Q)\). Therefore, at the arithmetic basepoint \(x_0=G(\mathbb Q)\),

\[
x_0B_w=x_0.
\tag{7}
\]

The same is true of a seed translation

\[
u_n=
\begin{pmatrix}
1&n\\
0&1
\end{pmatrix},
\qquad n\in\mathbb Z,
\]

since \(u_n\in G(\mathbb Q)\). Thus

\[
x_0u_n=x_0.
\]

So the direct diagonal adelic embedding identifies all rational branch products and all integer seeds with the same automorphic point.

The exact reason is visible in the following identity.


\paragraph*{Automorphic collapse identity}


Let \(n_j=U^j(n)\). Equation (2) is equivalent to

\[
\boxed{
B_wu_n=u_{n_j}
\begin{pmatrix}
3^j&0\\
0&2^A
\end{pmatrix}.
}
\tag{8}
\]

Indeed, both sides are

\[
\begin{pmatrix}
3^j&3^jn+C_w\\
0&2^A
\end{pmatrix}.
\]

The two unipotent factors \(u_n\) and \(u_{n_j}\) contain the seed and endpoint. Automorphy under integer translations removes both of them, leaving only

\[
d_{j,A}=
\begin{pmatrix}
3^j&0\\
0&2^A
\end{pmatrix}.
\]

Consequently, the standard automorphic quotient sees only the two exponents \(j,A\). It does not see \(n\), \(n_j\), \(C_w\), positivity, or the exact divisibility condition.


\paragraph*{The real modular quotient also loses the carry}


Acting on the upper half-plane,

\[
B_w\cdot(n+i)
  =\frac{3^j(n+i)+C_w}{2^A}
  =n_j+i\frac{3^j}{2^A}.
\]

Since \(n_j\in\mathbb Z\), modular translation gives

\[
\boxed{
[B_w\cdot(n+i)]
=
\left[i\frac{3^j}{2^A}\right].
}
\tag{9}
\]

Thus even before passing to all adeles, every modular or Maass observable evaluated on the actual seeded endpoint depends only on

\[
q_w=\frac{3^j}{2^A}.
\tag{10}
\]

Its Fourier modes cannot recover \(n_j\), because their value at the integer real part \(n_j\) is the same as at zero.


\paragraph*{The \(S\)-arithmetic quotient collapses completely}


Let

\[
S=\{\infty,2,3\},\qquad
\Gamma_S=\mathrm{PGL}_2(\mathbb Z[1/6]).
\]

The determinant of \(B_a\) is \(3\cdot2^a\), an \(S\)-unit. Hence

\[
B_a,\ B_w,\ u_n,\ d_{j,A}\in\Gamma_S.
\]

Therefore on

\[
\Gamma_S\backslash
\bigl(
\mathrm{PGL}_2(\mathbb R)
\times\mathrm{PGL}_2(\mathbb Q_2)
\times\mathrm{PGL}_2(\mathbb Q_3)
\bigr)
\]

the whole branch orbit is exactly the basepoint.

Quantitative \(S\)-arithmetic nondivergence is a genuine and powerful theory for flows on products of real and \(p\)-adic homogeneous spaces, but it cannot undo an exact arithmetic identification of the orbit with one point. [\hypertarget{cit:R56:4}{}\hyperlink{ref:R56:4}{4}]


\paragraph*{Rational conjugacy also erases the carry}


Every product is rationally conjugate to its diagonal part. Indeed,

\[
u_{\xi_w}^{-1}B_wu_{\xi_w}
=
\begin{pmatrix}
3^j&0\\
0&2^A
\end{pmatrix},
\qquad
\xi_w=\frac{C_w}{2^A-3^j}.
\tag{11}
\]

Thus trace, determinant, eigenvalues, and rational conjugacy class depend only on \(j,A\).

The carry survives only in the \textbf{integrality of the conjugator}. In particular, the word \(w\) gives a positive cycle precisely when

\[
\frac{C_w}{2^A-3^j}\in\mathbb Z_{>0}.
\tag{12}
\]

Automorphic conjugacy invariants discard exactly the datum needed to distinguish an integral cycle from a merely rational diagonalization.


\bigskip



\subsubsection*{4. The Hecke-correspondence repair is below the square-root scale}


One can avoid literal adelic constancy by keeping only the real quotient and treating \(B_w\) as a Hecke correspondence. This recovers individual matrices, but not enough pointwise control.

For \(j\geq1\), equation (1) shows

\[
C_w\equiv1\pmod2
\]

because its first term is odd and the remaining terms are even. Also,

\[
C_w\equiv2^{A_{j-1}}\not\equiv0\pmod3.
\]

Therefore the gcd of all entries of \(B_w\) is \(1\). Its Smith normal form is

\[
\operatorname{diag}(1,N),
\qquad
N=3^j2^A.
\tag{13}
\]

Hence every word with the same pair \((j,A)\) lies in the same primitive Hecke double coset

\[
\mathrm{GL}_2(\mathbb Z)
\begin{pmatrix}
1&0\\
0&N
\end{pmatrix}
\mathrm{GL}_2(\mathbb Z).
\]

The full cyclic Hecke correspondence has degree

\[
\psi(N)
=
N\prod_{p\mid N}\left(1+\frac1p\right)
=
2N.
\tag{14}
\]

By contrast, the number of words of length \(j\) and total exponent \(A\) is exactly

\[
M(j,A)=\binom{A-1}{j-1}.
\tag{15}
\]

These matrices are distinct: equality of their carries would give the same residue class in Proposition 1, while a seed cannot realize two different valuation words.

Now suppose the word has nonnegative multiplicative drift,

\[
2^A\leq3^j.
\tag{16}
\]

Then

\[
M(j,A)
  \leq 2^{A-1}
  <2^A
  \leq\sqrt{2^A3^j}
  =\sqrt N.
\tag{17}
\]

We have proved:


\begin{quote}
\textbf{Sub-square-root Hecke thinness.}   Every expanding \((j,A)\)-shell of Collatz branch products occupies fewer than \(\sqrt N\) points inside a Hecke correspondence of degree \(2N\).

\end{quote}


At the critical slope \(A\sim(\log_2 3)j\),

\[
M(j,A)=N^{\,0.474977\ldots+o(1)},
\]

whereas \(\sqrt N=N^{1/2}\). Equivalently,

\[
M(j,A)\sim
\frac{(2.839513\ldots)^j}{\sqrt j},
\qquad
\sqrt N\sim3^j.
\]

This is a sharp quantitative obstruction. Adelic mixing and Hecke-tree equidistribution naturally describe whole Hecke spheres or averaged families. The Collatz branch family is not only exponentially sparse inside the sphere; in the growth-critical region it lies below the natural square-root scale. Guilloux's adelic-mixing framework illustrates precisely the whole-Hecke-tree type of averaging that is available. [\hypertarget{cit:R56:5}{}\hyperlink{ref:R56:5}{5}]

Even perfect equidistribution of the full Hecke sphere would therefore say nothing pointwise about whether one specially selected branch survives forever.


\bigskip



\subsubsection*{5. Product-formula heights cannot distinguish the carry}


The diagonal multiplier is

\[
q_w=\frac{3^j}{2^A}.
\]

Its local logarithms are

\[
\log |q_w|_\infty=j\log3-A\log2,
\]

\[
\log |q_w|_2=A\log2,
\]

\[
\log |q_w|_3=-j\log3.
\]

Hence

\[
\log |q_w|_\infty+
\log |q_w|_2+
\log |q_w|_3=0.
\tag{18}
\]

This is just the product formula.

Every algebraic character of the upper Borel factors through its diagonal torus: the unipotent radical has no nontrivial algebraic characters. Consequently, all standard character-based cusp heights and Busemann functions depend on \(j,A\), not on \(C_w\).

Thus the two natural possibilities both fail:

\begin{itemize}
\item On the full \(S\)-arithmetic quotient, the product-formula displacement is an arithmetic lattice displacement and is trivial.
\item On a single local quotient, one sees escape in one direction but has discarded the compensating places and the carry that enforces integrality.
\end{itemize}

There is no globally automorphic linear combination of the three local diagonal drifts that is strictly negative: their canonical global sum is identically zero.


\bigskip



\subsubsection*{6. Uniform nonescape is explicitly false}


The word

\[
w=1^j
\]

has

\[
C_w=3^j-2^j
\]

and therefore

\[
B_1^j=
\begin{pmatrix}
3^j&3^j-2^j\\
0&2^j
\end{pmatrix}.
\]

Its real multiplier is

\[
q_j=\left(\frac32\right)^j\longrightarrow\infty,
\]

and its affine fixed point is \(-1\):

\[
U(-1)=-1,\qquad a(-1)=1.
\]

So there is an exact rational, integral, indefinitely escaping branch ray. It is not a positive counterexample, but positive integers shadow it for arbitrarily long times.

For \(L\geq1\), define

\[
n_L=2^{L+1}-1.
\]

Then for \(0\leq i\leq L\),

\[
\boxed{
U^i(n_L)=3^i2^{L+1-i}-1.
}
\tag{19}
\]

For \(i<L\),

\[
3U^i(n_L)+1
=
2\bigl(3^{i+1}2^{L-i}-1\bigr),
\]

and the factor in parentheses is odd. Thus the first \(L\) valuations are all \(1\).

Moreover,

\[
U^i(n_L)>n_L
\qquad(1\leq i\leq L).
\]

Hence positive integers exhibit arbitrarily long real cusp excursions with no descent.

They also satisfy

\[
n_L\equiv-1\pmod {2^{L+1}},
\]

so every fixed \(2\)-adic level identifies a sufficiently long positive shadow with the exceptional point \(-1\).

More generally, for any fixed modulus \(M\),

\[
n=M2^{L+1}-1
\]

satisfies

\[
U^i(n)\equiv-1\pmod M
\]

through the same long all-\(1\) segment. Thus no fixed-conductor automorphic form can distinguish the positive shadow from the rational exceptional ray for a uniformly bounded time.

This agrees with the known \(2\)-adic picture: the extended Collatz map is strongly ergodic and conjugate to a shift for Haar-almost every \(2\)-adic point, while the ordinary integers form a dense set of Haar measure zero. The pointwise conjecture concerns precisely this exceptional null set. [\hypertarget{cit:R56:6}{}\hyperlink{ref:R56:6}{6}]


\bigskip



\subsubsection*{7. Random affine dynamics gives the right heuristic, not the pointwise theorem}


For Haar-random odd \(n\in\mathbb Z_2\),

\[
\mathbb P(a(n)=k)=2^{-k},
\qquad k\geq1,
\]

and hence

\[
\mathbb E[a]=2.
\]

The average real logarithmic multiplier is therefore

\[
\mathbb E\log\frac{3}{2^a}
=
\log3-2\log2
=
\log\frac34<0.
\tag{20}
\]

So random affine products contract at the real place. This is exactly the setting in which Poisson-boundary theory for rational affine random walks is effective: the boundary is assembled from those real or \(p\)-adic fields in which the random products contract in mean. [\hypertarget{cit:R56:7}{}\hyperlink{ref:R56:7}{7}]

But the Collatz valuation word is not an independent random sample. It is the unique word selected by one prescribed \(2\)-adic boundary point \(n\). The all-\(1\) word is an explicit exceptional deterministic ray, and the positive numbers \(n_L\) shadow it for arbitrary durations.

Random mixing proves statements about harmonic measure or almost every boundary point. Collatz asks for every point of a countable measure-zero arithmetic set. There is no measure-theoretic implication from one statement to the other.


\bigskip



\subsubsection*{8. The strongest genuine automorphic height identity}


Although standard automorphic recurrence fails, there is a useful exact cusp-loss inequality.

For \(0<c\leq1\), attach to a positive odd integer \(n\) the upper-half-plane point

\[
z_c(n)=n+i(n+c).
\]

For one branch \(n'=U(n)\), with \(a=a(n)\),

\[
B_a z_c(n)
=
n'
+i\frac{3}{2^a}(n+c).
\]

The target section has height \(n'+c\). Their height ratio is

\[
\frac{n'+c}
     {(3/2^a)(n+c)}
=
\frac{3n+1+2^ac}{3(n+c)}.
\tag{21}
\]

For \(a=1\), this is at least \(1\) exactly when \(c\leq1\). For \(a\geq2\), it is automatically greater than \(1\). Therefore

\[
\frac{n'+c}
     {(3/2^a)(n+c)}
\geq1.
\]

Iterating gives the exact defect formula

\[
\boxed{
\log\frac{n_j+c}{q_j(n+c)}
=
\sum_{k=0}^{j-1}
\log
\frac{3n_k+1+2^{a_k}c}
     {3(n_k+c)}
\geq0.
}
\tag{22}
\]

For \(c=1\), this becomes the purely word-theoretic inequality

\[
\boxed{
C_w+2^A\geq3^j.
}
\tag{23}
\]

Indeed, defining

\[
D_j=C_j+2^{A_j}-3^j,
\]

one obtains

\[
D_{j+1}
=
3D_j+2^{A_j}(2^{a_j}-2),
\qquad D_0=0.
\tag{24}
\]

Hence \(D_j\geq0\), with equality exactly when

\[
a_0=\cdots=a_{j-1}=1.
\]

This is a genuine affine-cusp Lyapunov inequality. It says that deviations from the exceptional \(-1\) ray create a nonnegative cusp defect.

Unfortunately, the defect increment is approximately

\[
\frac{2^{a_k}-2}{3n_k}
\]

when \(n_k\) is large. Along an exponentially growing trajectory these increments may be summable. More decisively, the positive shadows \(n_L\) have \(D_j=0\) throughout their arbitrarily long all-\(1\) prefixes.

Thus (22) supplies monotonicity but not the divergent cumulative loss needed to force descent.


\bigskip



\subsubsection*{9. What an actual automorphic proof would have to establish}


To retain the seed and carry, one must replace the ordinary quotient by an inverse congruence tower such as

\[
Y=\varprojlim_r
\Gamma(2^r)\backslash \mathrm{PGL}_2(\mathbb R),
\]

together with the corresponding \(3\)-adic carry coordinates and an archimedean height.

At depth \(j\), exact validity of a word of total exponent \(A_j\) requires level at least

\[
2^{A_j+1}.
\]

Therefore no fixed automorphic representation or fixed conductor can follow an orbit indefinitely. The conductor must grow exponentially with the depth.

Moreover, ordinary mixing controls an averaged operator of the form

\[
\sum_w p_w R_{B_w}
\]

in an \(L^2\)-type norm. Collatz requires control of the worst surviving path, corresponding to a maximal or joint-spectral-radius operator

\[
\mathcal K f(x)
=
\sup_{\substack{a\geq1\\
\text{transition survives}}}
f(xB_a).
\tag{25}
\]

A sufficient new theorem would be the following.


\paragraph*{Marked-cusp contraction theorem}


There exist a proper weight \(W\geq1\), an integer \(m\geq1\), and \(0<\rho<1\) on the full marked congruence-height state space such that, whenever an exact \(m\)-step Collatz block remains above its marked initial value,

\[
W(xB_{a_{m-1}}\cdots B_{a_0})
\leq
\rho W(x).
\tag{26}
\]

This would prove Collatz immediately.

\textbf{Conditional proof.} Suppose \(n_0>1\) is the least odd counterexample. No iterate can be less than \(n_0\), since a smaller odd integer reaches \(1\) by minimality and would therefore pull \(n_0\) to \(1\). Thus the orbit gives an infinite survivor path.

Applying (26) successively gives

\[
W(x_{km})\leq\rho^kW(x_0)\longrightarrow0,
\]

contradicting \(W\geq1\). Hence no least counterexample exists, every odd positive integer reaches \(1\), and every even integer does after first dividing out its powers of \(2\). \ensuremath{\blacksquare}

The missing statement is therefore not ordinary mixing. It is a \textbf{uniform \(L^\infty\) contraction on a growing-level, marked congruence tower}, for a branch set that is below the square-root Hecke scale.

No such theorem follows from present automorphic mixing, nonescape of mass, or random affine boundary theory.


\bigskip



\subsection*{Final assessment}


The proposed approach has three possible implementations:

\begin{enumerate}
\item \textbf{Use the genuine adelic quotient.}   Then every branch product is rational and the orbit is constant.

\item \textbf{Use real or partial adelic quotients and Hecke correspondences.}   Then the carry becomes one exponentially sparse Hecke branch; the growth-compatible family has fewer than \(\sqrt N\) points at conductor \(N\), and endpoint automorphic observables retain only \(j,A\).

\item \textbf{Use a marked inverse congruence tower.}   This retains the arithmetic, but ordinary automorphic mixing no longer supplies the required pointwise theorem. One needs the unproved marked-cusp contraction (26).
\end{enumerate}

Therefore I do not have a valid full proof of Collatz from approach 51. The exact obstruction is equation (8): automorphic reduction removes the two integer translations containing the seed and endpoint and leaves only the diagonal drift. Any claimed proof that invokes mixing or nonescape after this reduction, without reconstructing the marked unipotent data at growing level, has discarded the essential Collatz condition.
\subsection*{References}
\begin{enumerate}[leftmargin=2.2em]
\item \hypertarget{ref:R56:1}{}\hyperlink{cit:R56:1}{\texttt{turn588450view1}}
\item \hypertarget{ref:R56:2}{}\hyperlink{cit:R56:2}{\texttt{turn322861academia12}}
\item \hypertarget{ref:R56:3}{}\hyperlink{cit:R56:3}{\texttt{turn275231view0}}
\item \hypertarget{ref:R56:4}{}\hyperlink{cit:R56:4}{\texttt{turn493036view0}}
\item \hypertarget{ref:R56:5}{}\hyperlink{cit:R56:5}{\texttt{turn275231view2}}
\item \hypertarget{ref:R56:6}{}\hyperlink{cit:R56:6}{\texttt{turn275231view3}}
\item \hypertarget{ref:R56:7}{}\hyperlink{cit:R56:7}{\texttt{turn275231view1}}
\end{enumerate}

\clearpage
\section{Round 57: approach 52 --- Property-tau for Collatz residue graphs}
\subsection*{Prompt}
\begin{quote}
Complete a full proof of the Collatz conjecture using the following approach:

\begin{enumerate}
\item Property-tau for Collatz residue graphs
\end{enumerate}

\textbf{Speculative \textperiodcentered{} P}

Study transition graphs modulo a family of moduli involving both powers of 2 and powers of 3. Prove uniform expansion. Combine expansion with the order structure of positive integers to show that an exceptional orbit cannot stay trapped in a shrinking family of bad residues.

Take as long as you need.
\end{quote}
\subsection*{Response}

\subsection*{Outcome of the property-\(\tau\) program}


After carrying this approach through, I cannot truthfully claim a full proof of Collatz. The obstruction is not merely that one difficult estimate remains: \textbf{the natural mixed Collatz residue graphs do not form a property-\(\tau\) family}, and the strongest genuine spectral gap that survives after quotienting out the \(2\)-adic memory averages over the wrong objects.

Below is a complete proof of those statements, together with an exact order-sensitive criterion that would repair the approach. As of August 2026, Collatz remains open; the current published computational frontier cited in recent work is convergence for every \(n<2^{71}\). [\hypertarget{cit:R57:1}{}\hyperlink{ref:R57:1}{1}, \hypertarget{cit:R57:2}{}\hyperlink{ref:R57:2}{2}]


\bigskip



\subsubsection*{1. The natural mixed residue graph}


Use the shortcut map

\[
T(n)=
\begin{cases}
n/2,&n\equiv0\pmod2,\\[1mm]
(3n+1)/2,&n\equiv1\pmod2.
\end{cases}
\]

For

\[
M=2^a3^b,\qquad a,b\geq1,
\]

define \(G_{a,b}\) on \(\mathbb Z/M\mathbb Z\) by putting an edge \(x\to y\) when there is an integer \(n\equiv x\pmod M\) with

\[
T(n)\equiv y\pmod M.
\]

There are two outgoing edges from each vertex, corresponding to the two lifts \(n\) and \(n+M\) modulo \(2M\).

For \(x\bmod 2^a\), let

\[
Q_a(x)=\bigl(\varepsilon_0(x),\ldots,\varepsilon_{a-1}(x)\bigr)
\]

be the first \(a\) parity bits of its orbit. The map

\[
Q_a:\mathbb Z/2^a\mathbb Z\longrightarrow\{0,1\}^a
\]

is a bijection. Indeed, for a prescribed word \(\varepsilon\), the usual affine identity

\[
2^aT^a(n)=3^{s_a}n+C_\varepsilon,
\qquad
s_a=\sum_{i<a}\varepsilon_i,
\]

determines one and only one \(n\bmod2^a\), since \(3^{s_a}\) is odd.

The Chinese remainder theorem therefore gives coordinates

\[
\Phi_{a,b}(x)
 =
(\varepsilon_0,\ldots,\varepsilon_{a-1};r),
\qquad r=x\bmod3^b.
\]


\paragraph*{Proposition 1. Exact skew-product form}


In these coordinates the edges are

\[
(\varepsilon_0,\ldots,\varepsilon_{a-1};r)
\longrightarrow
(\varepsilon_1,\ldots,\varepsilon_{a-1},\delta;
 F_{\varepsilon_0}(r)),
\qquad \delta\in\{0,1\},
\]

where

\[
F_0(r)=2^{-1}r,
\qquad
F_1(r)=2^{-1}(3r+1)
\pmod {3^b}.
\]


\subparagraph*{Proof}


Applying \(T\) deletes the first parity bit. The two lifts modulo \(2M\) have the same first \(a-1\) subsequent parity bits and opposite next parity bits, producing the free bit \(\delta\).

Modulo \(3^b\), division by \(2\) is invertible, and the residue update is exactly \(F_0\) or \(F_1\), according to \(\varepsilon_0\). \ensuremath{\blacksquare}

Thus \(G_{a,b}\) is a de Bruijn shift in the \(2\)-adic coordinate, skewed by two affine maps in the \(3\)-adic coordinate. The pure \(2^a\) modular Collatz graphs are known to be precisely binary de Bruijn graphs. [\hypertarget{cit:R57:3}{}\hyperlink{ref:R57:3}{3}]


\bigskip



\subsubsection*{2. Uniform expansion already fails}


There are two relevant meanings of ``expansion'': directed out-neighborhood growth and the Cheeger/spectral expansion of an undirected regularization. Both fail naturally.


\paragraph*{Theorem 2. No directed vertex expansion}


Let

\[
S_{a,b}
 =
\{(\varepsilon;r):r\equiv1\pmod3\}.
\]

Then

\[
N^+(S_{a,b})
 =
\{(\varepsilon;r):r\equiv2\pmod3\}
\]

and consequently

\[
|N^+(S_{a,b})|=|S_{a,b}|.
\]


\subparagraph*{Proof}


If \(r\equiv1\pmod3\), then

\[
F_0(r)\equiv 2r\equiv2\pmod3,
\]

while

\[
F_1(r)\equiv\frac12\equiv2\pmod3.
\]

Thus every outgoing neighbor has residue \(2\pmod3\).

Conversely, given any target \((\varepsilon';r')\) with \(r'\equiv2\pmod3\), choose the predecessor branch \(\varepsilon_0=0\) and put

\[
r=2r'\pmod{3^b}.
\]

Then \(r\equiv1\pmod3\) and \(F_0(r)=r'\). The predecessor parity coordinates can be chosen to produce the prescribed shifted word \(\varepsilon'\). \ensuremath{\blacksquare}

So there is no constant \(\eta>0\) such that

\[
|N^+(S)|\geq(1+\eta)|S|
\]

for all sets of size at most a fixed positive proportion of the graph.


\paragraph*{Theorem 3. The undirected Cheeger constants tend to zero as \(a\to\infty\)}


Let \(\widetilde G_{a,b}\) be the natural undirected symmetrization. Then, uniformly in \(b\),

\[
h(\widetilde G_{a,b})\ll a^{-1/2}.
\]


\subparagraph*{Proof}


Let

\[
w(\varepsilon)=\varepsilon_0+\cdots+\varepsilon_{a-1}
\]

and put

\[
m=\left\lfloor\frac{a-1}{2}\right\rfloor,
\qquad
S=\{(\varepsilon;r):w(\varepsilon)\leq m\}.
\]

Along an edge, the parity word is shifted and one bit is appended, so

\[
w(\varepsilon')-w(\varepsilon)=\delta-\varepsilon_0\in\{-1,0,1\}.
\]

Therefore a boundary vertex must lie in one of the two central layers

\[
w(\varepsilon)=m
\quad\text{or}\quad
w(\varepsilon)=m+1.
\]

Hence

\[
|\partial S|
 \leq
3^b\left(
\binom am+\binom a{m+1}
\right)
\ll
\frac{2^a3^b}{\sqrt a}.
\]

On the other hand,

\[
|S|
 =
3^b\sum_{j=0}^{m}\binom aj
\gg 2^a3^b.
\]

The symmetrized graphs have uniformly bounded degree, so the edge-boundary and vertex-boundary versions of the Cheeger constant differ only by an absolute factor. Thus

\[
h(\widetilde G_{a,b})\ll a^{-1/2}.
\]

\ensuremath{\blacksquare}

In fact, the de Bruijn component has explicit slowly varying modes. If \(X_j(\varepsilon)=(-1)^{\varepsilon_j}\), then the symmetrized shift acts on the span of \(X_0,\ldots,X_{a-1}\) like a path of length \(a\), whose spectral gap is of order \(a^{-2}\).

Thus any family that retains an unbounded amount of \(2\)-adic parity memory fails property-\(\tau\).


\bigskip



\subsubsection*{3. Passing to coprime moduli does not repair property-\(\tau\)}


One might try to use moduli \(m\) coprime to \(6\), so that both branches

\[
T_0(x)=x/2,
\qquad
T_1(x)=(3x+1)/2
\]

are permutations modulo \(m\). But the generated group is affine and solvable.

Indeed,

\[
\Gamma=\langle T_0,T_1\rangle
   \subset \operatorname{Aff}(\mathbb Q)
   =\mathbb Q\rtimes\mathbb Q^\times.
\]

The translation subgroup is abelian and the multiplier quotient is abelian, so \(\Gamma\) is metabelian and hence amenable. Monks' exact finite classification is

\[
G_m\cong C_m\rtimes M
\]

for \(m\) coprime to \(6\), with \(M\) generated by multiplication by \(2\) and \(3\). [\hypertarget{cit:R57:4}{}\hyperlink{ref:R57:4}{4}]


\paragraph*{Proposition 4. No separating affine congruence family has property-\(\tau\)}


Let \(\pi_i:\Gamma\to\Gamma_i\) be finite congruence quotients whose injectivity radii tend to infinity. Then the Cayley graphs

\[
\operatorname{Cay}(\Gamma_i,\pi_i(S))
\]

cannot have a uniform positive Cheeger constant for any fixed finite generating set \(S\).


\subparagraph*{Proof}


Since \(\Gamma\) is amenable, for every \(\epsilon>0\) there is a finite nonempty F\o{}lner set \(F\subset\Gamma\) satisfying

\[
|SF\setminus F|<\epsilon |F|.
\]

For sufficiently large \(i\), the quotient map is injective on \(F\cup SF\), and \(|\pi_i(F)|\leq|\Gamma_i|/2\). Therefore

\[
\frac{
|\partial\pi_i(F)|
}{
|\pi_i(F)|
}
=
\frac{|SF\setminus F|}{|F|}
<\epsilon.
\]

Since \(\epsilon\) is arbitrary, no uniform expansion constant exists. \ensuremath{\blacksquare}

This also rules out replacing one-step branches by any fixed finite collection of accelerated affine blocks: they remain inside the same solvable affine group. Letting the generator set grow with the modulus is no longer property-\(\tau\), and any resulting expansion constants do not automatically transfer to a fixed integer orbit.


\bigskip



\subsubsection*{4. A surprisingly strong \(3\)-adic spectral theorem}


There is nevertheless a genuine positive result after forgetting the \(2\)-adic memory.

Let

\[
U_b=(\mathbb Z/3^b\mathbb Z)^\times,
\qquad
N_b=|U_b|=2\cdot3^{b-1},
\]

and define on functions \(f:U_b\to\mathbb C\)

\[
(\mathcal L_{b,t}f)(x)
 =
f(x/2)+t\,f((3x+1)/2),
\qquad t\geq0.
\]

The parameter \(t\) weights odd branches.


\paragraph*{Theorem 5. Exact weighted characteristic polynomial}


For every \(b\geq1\),

\[
\boxed{
\det(zI-\mathcal L_{b,t})
 =
\bigl(z-(1+t)\bigr)
\frac{z^{N_b}-1}{z-1}.
}
\]

Consequently, for

\[
P_{b,t}=\frac{\mathcal L_{b,t}}{1+t},
\]

the constant eigenvalue is \(1\), while every other algebraic eigenvalue has modulus

\[
\frac1{1+t}.
\]

In particular, for equal branch weights \(t=1\), the entire nonconstant algebraic spectrum lies on the circle of radius \(1/2\).


\subparagraph*{Proof}


For \(0\leq j\leq b\), let \(V_j\) be the space of functions depending only on \(x\bmod3^j\), with \(V_0\) the constants. Then

\[
V_0\subset V_1\subset\cdots\subset V_b.
\]

Let

\[
Ef(x)=f(x/2),
\qquad
Of(x)=f((3x+1)/2).
\]

Clearly \(E(V_j)\subseteq V_j\). Moreover,

\[
(3x+1)/2\bmod3^j
\]

depends only on \(x\bmod3^{j-1}\); hence

\[
O(V_j)\subseteq V_{j-1}.
\]

Therefore, on the quotient \(V_j/V_{j-1}\), the operator \(\mathcal L_{b,t}=E+tO\) induces simply \(E\).

Multiplication by \(2^{-1}\) is one cycle on \(U_j\), because

\[
\operatorname{ord}_{3^j}(2)=2\cdot3^{j-1}=N_j.
\]

The order assertion follows from the lifting-the-exponent identity

\[
v_3\!\left(4^{3^{j-1}}-1\right)
 =
v_3(4-1)+v_3(3^{j-1})=j.
\]

Thus \(E\) on \(V_j\) has characteristic polynomial \(z^{N_j}-1\), and on \(V_j/V_{j-1}\) it has characteristic polynomial

\[
\frac{z^{N_j}-1}{z^{N_{j-1}}-1}.
\]

Multiplying these quotient factors telescopes to

\[
\frac{z^{N_b}-1}{z-1}.
\]

Finally, on constants,

\[
\mathcal L_{b,t}1=(1+t)1.
\]

This proves the formula. \ensuremath{\blacksquare}

This looks considerably stronger than the expansion originally requested. But it does \textbf{not} solve the pointwise problem.

First, \(\mathcal L_{b,t}\) is non-normal, so its algebraic spectrum is not by itself a uniform \(L^2\) property-\(\tau\) statement.

More importantly,

\[
P_{b,1}^k f(x)
 =
2^{-k}
\sum_{\varepsilon\in\{0,1\}^k}
f\!\left(
F_{\varepsilon_{k-1}}\circ\cdots\circ F_{\varepsilon_0}(x)
\right).
\]

It averages over all \(2^k\) parity continuations. A fixed integer follows exactly \textbf{one} continuation.


\bigskip



\subsubsection*{5. An explicit exceptional spine survives the spectral gap}


Consider the negative shortcut orbit

\[
-5\longmapsto-7\longmapsto-10\longmapsto-5.
\]

Its parity word is

\[
(110)^\infty.
\]

The leading multiplicative coefficient along the first three prefixes is

\[
\frac32,\qquad \frac94,\qquad \frac98,
\]

and each further block multiplies it by \(9/8>1\). Thus

\[
\frac{3^{s_k}}{2^k}\geq1
\qquad\text{for every }k.
\]

This gives an infinite critical-growth path in every finite \(3\)-adic residue graph, despite the exact algebraic spectral radius \(1/2\) off constants.

This is not a positive Collatz counterexample. Its role is to show that \textbf{average residue mixing and the existence of one forever-exceptional path are completely compatible}. The distinction between this path and a hypothetical positive divergent path is archimedean order, not modular expansion.

More generally, every infinite parity sequence is realized by exactly one \(2\)-adic integer. [\hypertarget{cit:R57:5}{}\hyperlink{ref:R57:5}{5}] Uniform branch averaging therefore averages over all \(2\)-adic points, while the positive integers form a highly exceptional countable subset.

There is an even more concrete interpretation. The two outgoing edges from \(x\bmod M\) correspond to the actual integers

\[
x\quad\text{and}\quad x+M,
\]

not to two choices available to the orbit of \(x\). Iterating all graph edges varies the higher binary digits of the initial integer. The positive representative \(x\in[0,M)\) is one distinguished boundary ray in that tree of lifts. Property-\(\tau\) gives no control over that one ray.


\bigskip



\subsubsection*{6. The exact order statement that would prove Collatz}


For \(k\geq1\), define

\[
\mathcal B_k
 =
\left\{
1\leq n<2^k:
T^j(n)\geq n\ \text{for every }1\leq j\leq k
\right\}.
\]

Let

\[
m_k=\min(\mathcal B_k\setminus\{1\}),
\]

with \(m_k=\infty\) if the set is empty.


\paragraph*{Theorem 6. Exact minimum-escape criterion}


The Collatz conjecture is equivalent to

\[
\boxed{m_k\longrightarrow\infty.}
\]


\subparagraph*{Proof}


Suppose Collatz is false and let \(n>1\) be its least counterexample. If some iterate satisfied

\[
T^j(n)<n,
\]

then \(T^j(n)\) would converge by minimality, forcing \(n\) to converge. Hence

\[
T^j(n)\geq n
\qquad\text{for every }j.
\]

For every sufficiently large \(k\), \(n<2^k\), so \(n\in\mathcal B_k\). Thus \(m_k\leq n\) infinitely often.

Conversely, suppose \(m_k\to\infty\). Fix \(n>1\). Choose \(k\) with \(2^k>n\) and \(m_k>n\). Then \(n\notin\mathcal B_k\), so some iterate \(T^j(n)<n\). Strong induction on \(n\) now proves that every positive integer reaches \(1\). \ensuremath{\blacksquare}

This exposes the precise logical gap:

\[
\frac{|\mathcal B_k|}{2^k}\longrightarrow0
\]

is not enough. A fixed integer \(n\) could remain in every \(\mathcal B_k\) while contributing vanishing density. Even an exponentially small bad set can contain one coherent nested spine.

A sufficient order-sensitive replacement for property-\(\tau\) would be, for example,

\[
|\mathcal B_k\cap[2,N]|
 \leq
C N^\theta\rho^k,
\qquad
\rho<1,
\]

uniformly for \(N<2^k\). For fixed \(N\), the right-hand side eventually becomes \(<1\).

An equivalent type of target is a height-weighted survivor estimate. For \(s>0\), put

\[
Z_k(s)
 =
\sum_{\substack{n\in\mathcal B_k\\n>1}}n^{-s}.
\]

Then

\[
\boxed{
Z_k(s)\longrightarrow0
\quad\Longrightarrow\quad
\text{Collatz}.
}
\]

Indeed, a least counterexample \(n\) would contribute the fixed positive amount \(n^{-s}\) to \(Z_k(s)\) for all sufficiently large \(k\).

This is the natural ``property-\(\tau\) with order'' statement. Ordinary property-\(\tau\) uses uniform counting measure and cannot see it.


\bigskip



\subsubsection*{7. An entropy\textendash{}height obstruction to the inverse-expansion bridge}


There is a second exact obstruction. It shows that even perfect residue equidistribution would not automatically supply smaller competing ancestors.

Use the odd accelerated map

\[
U(x)=\frac{3x+1}{2^{v_2(3x+1)}}.
\]

For a valuation word

\[
a=(a_1,\ldots,a_r),\qquad a_i\geq1,
\]

write

\[
A(a)=a_1+\cdots+a_r.
\]

Reversing this word from an endpoint \(X\) gives an affine expression

\[
z_a(X)
 =
\frac{2^{A(a)}X-C_a}{3^r}.
\]

For each fixed word \(a\), exact integrality at every reverse step selects at most one residue class

\[
X\bmod3^r.
\]

Let \(\mathcal R_r(L)\) be the endpoint residue classes that admit at least one reverse word of length \(r\) and total exponent at most \(L\).


\paragraph*{Theorem 7. Contracting inverse words occupy exponentially few residues}


For \(L\geq r\),

\[
\boxed{
|\mathcal R_r(L)|\leq\binom Lr.
}
\]

In particular, with

\[
\lambda=\log_2 3,
\qquad
L_r=\lfloor\lambda r\rfloor,
\]

one has

\[
\frac{|\mathcal R_r(L_r)|}{3^r}
\leq
\exp\bigl(-\kappa r+O(\log r)\bigr),
\]

where

\[
\kappa
 =
\log3-\lambda H(1/\lambda)
 =
0.0549794728\ldots>0
\]

and

\[
H(p)=-p\log p-(1-p)\log(1-p).
\]

Numerically, the exponential factor is approximately

\[
e^{-\kappa}=0.9465046.
\]


\subparagraph*{Proof}


The number of positive compositions of \(A\) into \(r\) parts is

\[
\binom{A-1}{r-1}.
\]

Therefore the number of words with \(A\leq L\) is

\[
\sum_{A=r}^{L}\binom{A-1}{r-1}
 =
\binom Lr.
\]

Each word contributes at most one endpoint class modulo \(3^r\), proving the first assertion.

Stirling's formula gives

\[
\log\binom{\lambda r+O(1)}r
 =
\lambda r\,H(1/\lambda)+O(\log r).
\]

After division by \(3^r\), the exponent is

\[
-\bigl(\log3-\lambda H(1/\lambda)\bigr)r+O(\log r).
\]

Finally,

\[
H(1/\lambda)<\log2
\]

because \(1/\lambda\neq1/2\), so

\[
\lambda H(1/\lambda)
<
\lambda\log2
=
\log3.
\]

Hence \(\kappa>0\). \ensuremath{\blacksquare}

The significance is that

\[
A(a)\leq\lambda r
\quad\Longleftrightarrow\quad
\frac{2^{A(a)}}{3^r}\leq1.
\]

So the endpoint residues admitting even a nonexpanding inverse affine slope form an exponentially small subset of \(\mathbb Z/3^r\mathbb Z\). To descend from a large orbit value \(X\) all the way below the original least counterexample \(n\), one needs

\[
\frac{2^{A(a)}}{3^r}\lesssim\frac nX,
\]

which is much stronger and gives an even smaller residue set.

Thus ideal equidistribution has the wrong implication: it says that robustly descending competitors are \textbf{rare}, not that every endpoint must possess one.

This aligns exactly with your earlier local-global decomposition. Terminal \(3\)-adic compatibility is genuinely suffix-local, as your terminal-locality theorem proves. [\hypertarget{cit:R57:6}{}\hyperlink{ref:R57:6}{6}] But whether a compatible competitor is positive and below the original input depends on an archimedean affine cutoff. [\hypertarget{cit:R57:7}{}\hyperlink{ref:R57:7}{7}] Property-\(\tau\) controls the first datum, while the required contradiction lives in the second.


\bigskip



\subsubsection*{8. Final verdict}


The proposed proof cannot be completed as stated for three independent reasons:

\begin{enumerate}
\item The natural mixed graphs have exact directed nonexpanding sets and de Bruijn bottlenecks, so the required uniform expansion theorem is false.

\item The only literal group-quotient version lies in a solvable affine group, and every separating congruence family fails property-\(\tau\).

\item Even the surviving exact \(3\)-adic algebraic spectral gap averages over all parity continuations. It does not control the single continuation corresponding to a fixed positive integer, and compatible shrinking inverse paths are themselves exponentially sparse.
\end{enumerate}

A recent independent reduction reaches the same conceptual boundary from another direction: it proves almost perfect map-level residue balance but leaves a deterministic, pointwise orbit-visitation statement as the unresolved step. [\hyperlink{ref:R57:2}{2}]

The repaired version of approach 52 is therefore not ordinary property-\(\tau\), but the following order-weighted contraction estimate:

\[
\boxed{
\exists\,s>0,\ C<\infty,\ \rho<1
\quad\text{such that}\quad
\sum_{\substack{2\leq n<2^k\\
T^j(n)\geq n\ \forall\,1\leq j\leq k}}
n^{-s}
\leq C\rho^k.
}
\]

This estimate would give a complete proof immediately by Theorem 6. Establishing it would require a genuinely pointwise height-residue mechanism\textemdash{}stronger than expansion, and capable of distinguishing the distinguished positive lift from the exponentially many other \(2\)-adic continuations.
\subsection*{References}
\begin{enumerate}[leftmargin=2.2em]
\item \hypertarget{ref:R57:1}{}\hyperlink{cit:R57:1}{\texttt{turn596684view0}}
\item \hypertarget{ref:R57:2}{}\hyperlink{cit:R57:2}{\texttt{turn521129view0}}
\item \hypertarget{ref:R57:3}{}\hyperlink{cit:R57:3}{\texttt{turn494958academia12}}
\item \hypertarget{ref:R57:4}{}\hyperlink{cit:R57:4}{\texttt{turn596684view1}}
\item \hypertarget{ref:R57:5}{}\hyperlink{cit:R57:5}{\texttt{turn596684view3}}
\item \hypertarget{ref:R57:6}{}\hyperlink{cit:R57:6}{\texttt{turn1file0 L29-L56}}
\item \hypertarget{ref:R57:7}{}\hyperlink{cit:R57:7}{\texttt{turn1file0 L60-L94}}
\end{enumerate}

\clearpage
\section{Round 58: approach 53 --- A thin-group affine sieve}
\subsection*{Prompt}
\begin{quote}
Complete a full proof of the Collatz conjecture using the following approach:

\begin{enumerate}
\item A thin-group affine sieve
\end{enumerate}

\textbf{Speculative \textperiodcentered{} B}

The semigroup of branch matrices may form a thin orbit in an algebraic group. Sieve its orbit for the divisibility properties required by a cycle or critical trajectory. An unexpectedly surviving local pattern could identify a counterexample family.

Take as long as you need.
\end{quote}
\subsection*{Response}

\subsection*{Outcome}


I do \textbf{not} obtain a valid full proof of the Collatz conjecture from the thin-group affine-sieve approach. The obstruction is structural, not merely a missing numerical estimate:

\begin{enumerate}
\item The branch matrices generate a \textbf{solvable arithmetic affine group}, not a semisimple thin group with congruence expansion.
\item There are arbitrarily long, primitive, positive rational Collatz cycles whose denominators avoid any prescribed finite collection of primes\textemdash{}even while every point of the cycle is arbitrarily large and the average drift is arbitrarily close to critical.
\item Every finite critical branch word is realized by positive integers, so finite congruence sieving cannot eliminate divergent paths either.
\end{enumerate}

The Collatz conjecture remains open as of August 2026. [\hypertarget{cit:R58:1}{}\hyperlink{ref:R58:1}{1}]

What follows is a rigorous analysis of the approach, including an explicit local-global obstruction theorem.


\bigskip



\subsubsection*{1. The branch-matrix orbit}


Work with the accelerated map on odd integers,

\[
U(n)=\frac{3n+1}{2^{\nu _2(3n+1)}}.
\]

For a prescribed valuation \(a\geq1\), write

\[
U_a(x)=\frac{3x+1}{2^a},
\qquad
M_a=
\begin{pmatrix}
3&1\\
0&2^a
\end{pmatrix}.
\]

For a word

\[
w=(a_0,\ldots ,a_{j-1}),
\qquad
A_i=\sum_{k<i}a_k,
\qquad
A=A_j,
\]

the product is

\[
M_w=M_{a_{j-1}}\cdots M_{a_0}
 =
\begin{pmatrix}
3^j&C_w\\
0&2^A
\end{pmatrix},
\]

where

\[
C_w=\sum_{i=0}^{j-1}3^{j-1-i}2^{A_i}.
\]

Thus

\[
U_w(x)=\frac{3^jx+C_w}{2^A}. \tag{1}
\]

The positive branch semigroup is free: if two words have the same product, the diagonal entries first determine \(j\) and \(A\), and then the first differing partial sum \(A_i\) gives a uniquely smallest 2-adic valuation in the difference of the two \(C\)-values.

So the orbit is combinatorially sparse. But that does \textbf{not} make its generated group a thin group of the kind required by the standard affine sieve.


\bigskip



\subsubsection*{2. The generated group is solvable and essentially arithmetic}


Let

\[
S(x)=\frac{x}{2},
\qquad
T(x)=3x+1.
\]

Then \(U_a=S^aT\). Conversely,

\[
U_{a+1}U_a^{-1}=S,
\qquad
S^{-a}U_a=T.
\]

Hence the group generated by all branches is exactly

\[
\Gamma=\langle S,T\rangle.
\]

In affine-pair notation, an element is written \((\lambda,t)\) for the map \(x\mapsto\lambda x+t\). The commutator gives

\[
STS^{-1}T^{-1}(x)=x-\frac12.
\]

Conjugating this translation by powers of \(2\) and \(3\) produces every translation in \(\mathbb Z[1/6]\). Therefore

\[
\boxed{
\Gamma=
\mathbb Z[1/6]\rtimes \langle 2,3\rangle .
}
\tag{2}
\]

Its Zariski closure is the full one-dimensional affine group

\[
G=\mathbb G_a\rtimes\mathbb G_m.
\]

This group is solvable, and

\[
[G,G]=\mathbb G_a\neq G.
\]

Consequently \(G^\circ\) is not perfect. In fact, relative to its own \(S\)-arithmetic closure, \(\Gamma\) is essentially arithmetic\textemdash{}of index at most two after adjoining the scale \(-1\)\textemdash{}rather than thin.

This matters decisively. Salehi Golsefidy and Varj\'u prove that the congruence Cayley graphs of such a finitely generated linear group form an expander family precisely when the connected component of its Zariski closure is perfect. That condition fails here. [\hypertarget{cit:R58:2}{}\hyperlink{ref:R58:2}{2}]

The standard affine sieve is built from exactly this sort of congruence expansion or level-of-distribution input; its usual output concerns primes or almost-primes in matrix orbits. [\hypertarget{cit:R58:3}{}\hyperlink{ref:R58:3}{3}]

Thus:

\[
\boxed{\text{The standard thin-group affine-sieve machinery does not apply to }\Gamma.}
\]

One might still attempt a specialized sieve on fixed \((j,A)\)-fibres, but that faces the stronger obstruction below.


\bigskip



\subsubsection*{3. The exact cycle condition}


From (1), the rational fixed point associated to \(w\) is

\[
x_w=\frac{C_w}{D_w},
\qquad
D_w=2^A-3^j. \tag{3}
\]

Since \(C_w>0\), a positive fixed point requires \(D_w>0\).

Set

\[
g_w=\gcd(C_w,D_w),
\qquad
p_w=\frac{C_w}{g_w},
\qquad
q_w=\frac{D_w}{g_w}.
\]

Then \(x_w=p_w/q_w\) is reduced. Cyclically rotating \(w\) produces the other points of an exact rational periodic orbit. Indeed, if \(C_i\) is the numerator attached to the \(i\)-th rotation, then

\[
2^{a_i}C_{i+1}=3C_i+D_w. \tag{4}
\]

Because \(D_w\) is coprime to \(6\), equation (4) shows that

\[
\gcd(C_i,D_w)=g_w
\]

for every rotation. Dividing (4) by \(g_w\) gives

\[
2^{a_i}p_{i+1}=3p_i+q_w,
\]

with every \(p_i\) odd. Hence the prescribed valuations are exact.

It follows that

\[
\boxed{
w\text{ gives a positive integer cycle}
\iff
D_w>0\text{ and }q_w=1.
}
\tag{5}
\]

Equivalently,

\[
D_w\mid C_w.
\]

In matrix language, this says \(M_w\) becomes scalar in

\[
\operatorname{PGL}_2(\mathbb Z/D_w\mathbb Z).
\]

For a prime \(\ell\),

\[
x_w\in\mathbb Z_\ell
\iff
\ell\nmid q_w.
\]

Thus the prime factors of \(q_w\) are exactly the local obstructions to integrality.

A simple example is

\[
w=(1,1,3),\qquad
A=5,\quad C_w=19,\quad D_w=5.
\]

It gives the exact rational cycle

\[
\frac{19}{5}
\longmapsto
\frac{31}{5}
\longmapsto
\frac{49}{5}
\longmapsto
\frac{19}{5}.
\tag{6}
\]

Repeating a word does not change its reduced fixed point: if \(w=v^m\), both \(C_w\) and \(D_w\) acquire the same geometric factor. Therefore rational periodic templates survive at every repeated length.

The existence of rational periodic orbits for every periodic symbolic pattern, and the legality of every finite symbolic sequence, are well-known features of the rational and 2-adic Collatz extensions. [\hypertarget{cit:R58:4}{}\hyperlink{ref:R58:4}{4}]


\bigskip



\subsubsection*{4. Arbitrarily rough positive rational pseudo-cycles}


The following result gives a precise obstruction to every finite-prime affine sieve.


\paragraph*{Theorem}


Let \(H\geq1\), and let \(\mathcal S\) be any finite set of primes. Then there exist arbitrarily long primitive valuation words \(w\) such that:

\begin{enumerate}
\item \(w\) generates an exact positive rational Collatz cycle;
\item every point of that cycle is greater than \(H\);
\item its reduced denominator satisfies \(q_w>1\);
\item no prime in \(\mathcal S\) divides \(q_w\);
\item its average exponent satisfies \[
   \frac{A}{j}=\log_2 3+O\!\left(\frac1j\right).
   \]
\end{enumerate}

Thus the cycles may be arbitrarily long, arbitrarily high, arbitrarily close to critical drift, and locally integral at every prescribed finite set of primes.


\paragraph*{Proof}


Put

\[
\theta=\log_2 3.
\]

Choose constants

\[
0<\eta _1<\eta _2<\log_2\!\left(1+\frac1H\right). \tag{7}
\]

Let

\[
M=\operatorname{lcm}\{p-1:p\in\mathcal S,\ p\geq5\},
\]

with \(M=1\) if the set is empty.

Since \(\theta\) is irrational, the fractional parts \(\{\theta n\}\) are dense in \([0,1]\). Choose arbitrarily large \(n\) satisfying

\[
\frac{1-\eta _2}{M}
<
\{\theta n\}
<
\frac{1-\eta _1}{M}.
\]

Define

\[
j=Mn,
\qquad
A=1+M\lfloor\theta n\rfloor.
\]

Then

\[
\delta:=A-\theta j
   =1-M\{\theta n\}
   \in(\eta _1,\eta _2). \tag{8}
\]

In particular \(A>\theta j\), so

\[
D=2^A-3^j
  =3^j(2^\delta-1)>0. \tag{9}
\]

Let \(L=A-j\), and choose the word

\[
w=(\underbrace{1,\ldots ,1}_{j-1\text{ times}},L+1). \tag{10}
\]

It is primitive because it contains exactly one exponent different from \(1\). Its total exponent is \(A\), and

\[
C_w=3^j-2^j. \tag{11}
\]

Therefore its smallest cyclic value is

\[
x_w
 =\frac{3^j-2^j}{2^A-3^j}
 =\frac{1-(2/3)^j}{2^\delta-1}. \tag{12}
\]

By (7) and sufficiently large \(j\),

\[
1-\left(\frac23\right)^j
>
H(2^{\eta _2}-1),
\]

and hence \(x_w>H\). The first \(j-1\) branches have exponent \(1\), so

\[
x_k+1=\left(\frac32\right)^k(x_w+1),
\qquad 0\leq k<j.
\]

Thus every other point of the rational cycle is larger than \(x_w\), and all points exceed \(H\).

Now consider local divisibility. If \(p\in\mathcal S\) and \(p\geq5\), then \(p-1\mid M\), so

\[
j\equiv0\pmod{p-1},
\qquad
A\equiv1\pmod{p-1}.
\]

Consequently,

\[
D=2^A-3^j\equiv2-1\equiv1\pmod p. \tag{13}
\]

Also \(D\) is odd, and

\[
D\equiv2^A\not\equiv0\pmod3.
\]

Hence no prime in \(\mathcal S\) divides \(D\), and therefore none divides \(q_w\).

It remains to prove that \(q_w>1\). Suppose instead that \(D\mid C_w\). Since

\[
C_w+D=2^A-2^j=2^j(2^L-1)
\]

and \(D\) is odd, this would imply

\[
D\mid 2^L-1. \tag{14}
\]

But from (8)\textendash{}(9),

\[
\frac{D}{2^L}
 =
\frac{2^A-3^j}{2^{A-j}}
 =
2^j\bigl(1-2^{-\delta}\bigr)
 \geq
2^j\bigl(1-2^{-\eta _1}\bigr).
\]

For sufficiently large \(j\), this is greater than \(1\). Hence

\[
D>2^L>2^L-1,
\]

contradicting (14). Thus \(D\nmid C_w\), and \(q_w>1\).

Finally, by (8),

\[
\frac{A}{j}
=\theta+\frac{\delta}{j}
=\log_2 3+O\!\left(\frac1j\right).
\]

This proves the theorem. \ensuremath{\blacksquare}


\bigskip



\subsubsection*{5. A concrete survivor}


Take

\[
j=13,\qquad A=21,\qquad
w=(1,1,\ldots ,1,9),
\]

with twelve \(1\)'s. Then

\[
D=2^{21}-3^{13}=502829,
\]

and

\[
C=3^{13}-2^{13}=1586131.
\]

Here

\[
\gcd(C,D)=1,
\]

and \(502829\) is prime. Thus

\[
x=\frac{1586131}{502829}
  \approx3.1544143238.
\]

The first twelve branches have exponent \(1\), followed by one branch of exponent \(9\), which returns to \(x\). Every point is greater than \(3.154\).

Therefore this primitive 13-step rational cycle is:

\begin{itemize}
\item positive;
\item entirely above \(3\);
\item in the nontrivial cycle slope window, since \[
  \frac{21}{13}\approx1.61538;
  \]
\item integral in \(\mathbb Z_p\) for every prime \(p<502829\);
\item globally nonintegral only because of the single prime \(502829\).
\end{itemize}

So even an affine sieve using \textbf{every prime below half a million} would leave this exact cycle-like object untouched.

The general theorem makes the least obstructing prime arbitrarily large.


\bigskip



\subsubsection*{6. Why this defeats a finite local sieve}


Define

\[
\mathcal B(w)
=
\left\{
p:
v_p(C_w)<v_p(D_w)
\right\}.
\]

These are precisely the prime divisors of \(q_w\). An integer cycle occurs exactly when

\[
\mathcal B(w)=\varnothing.
\]

The theorem proves that for every \(B\) there are arbitrarily long, primitive, positive near-critical rational cycles satisfying

\[
\min\mathcal B(w)>B.
\]

Thus there is no fixed finite set of local congruence obstructions that detects every nonintegral cycle candidate.

This does not logically exclude a completely new sieve whose level grows essentially as far as the moving modulus

\[
D_w=2^A-3^j.
\]

But such a theorem would have to establish, pointwise, that for every nontrivial word there is some prime\textemdash{}possibly comparable to \(D_w\)\textemdash{}for which

\[
p^{v_p(D_w)}\nmid C_w. \tag{15}
\]

That is almost exactly the integer-cycle problem itself. Standard affine-sieve estimates are density estimates and do not provide this kind of universal, moving-modulus exclusion.


\bigskip



\subsubsection*{7. The obstruction for divergent trajectories}


Cycles are only half of Collatz. A full proof must also exclude a positive orbit that escapes to infinity.

For a finite word \(w\), an odd integer \(n\) realizes exactly that valuation prefix precisely when

\[
3^j n+C_w\equiv2^A\pmod{2^{A+1}}.
\]

Therefore the exact prefix cylinder is the unique odd residue class

\[
\boxed{
n\equiv
3^{-j}(2^A-C_w)
\pmod{2^{A+1}}.
}
\tag{16}
\]

Every finite valuation word is consequently realized by infinitely many positive integers. This matches the known symbolic-dynamics fact that every finite pattern is legal; moreover, exponentially many finite residue classes can remain without a verified decrease even at large depth. [\hypertarget{cit:R58:5}{}\hyperlink{ref:R58:5}{5}]

The simplest example is the all-ones word

\[
w_j=(1,\ldots ,1).
\]

Here \(A=j\) and \(C=3^j-2^j\). Formula (16) gives

\[
n\equiv-1\pmod{2^{j+1}}.
\]

The least positive representative is

\[
n_j=2^{j+1}-1.
\]

Thus, for every depth \(j\), there is a positive integer whose first \(j\) accelerated valuations are all \(1\), producing multiplicative growth approximately \((3/2)^j\).

But the nested residue classes converge 2-adically to

\[
-1\in\mathbb Z_2,
\]

not to a positive integer. This is exactly the phenomenon that purely 2-adic or finite-symbolic dynamics cannot distinguish from an exceptional positive orbit; the real and 2-adic places must be controlled simultaneously. [\hyperlink{ref:R58:4}{4}]

So an estimate saying ``almost all finite words die'' is insufficient. One must prove that no single bounded positive representative survives along an infinite nested branch.


\bigskip



\subsubsection*{8. A useful two-place rigidity identity}


There is one exact bridge between rational periodic templates and positive integer cylinders.

Let

\[
x_w=\frac pq
\]

be the reduced rational fixed point of \(w\), and let \(r_m\) be an integer realizing the repeated prefix \(w^m\). The total exponent is \(mA\). Comparing the fixed-point equation with the prefix congruence gives

\[
\boxed{
qr_m\equiv p\pmod{2^{mA}}.
}
\tag{17}
\]

Therefore, if \(q>1\),

\[
|qr_m-p|\geq2^{mA},
\]

once the difference is nonzero. In particular,

\[
r_m\longrightarrow\infty.
\]

Thus a nonintegral rational pseudo-cycle cannot be shadowed for arbitrarily many repetitions by one fixed positive integer. Its positive cylinder representatives eventually escape exponentially.

This points to the correct repair of the approach: one must track both

\begin{itemize}
\item the odd-prime denominator \(q_w\), and
\item the archimedean size of the least positive 2-adic cylinder representative.
\end{itemize}

A congruence sieve sees only the first. The second is what distinguishes the rational and 2-adic decoys from actual positive integers.


\bigskip



\subsubsection*{9. The two missing global theorems}


A completed hybrid argument would need genuinely new results of the following form.


\paragraph*{Integral fixed-point rigidity}


For every primitive valuation word \(w\) with \(D_w>0\),

\[
q_w=1
\quad\Longrightarrow\quad
w=(2),
\]

up to repetition and rotation.

This is precisely the exclusion of nontrivial positive cycles.


\paragraph*{Cylinder-height escape}


For every infinite critical valuation word, if \(r_j\) denotes the least positive representative of its length-\(j\) exact cylinder, then

\[
r_j\longrightarrow\infty.
\]

Indeed, if the \(r_j\) remained bounded, the nested congruence classes would eventually have a common fixed positive representative, giving a positive integer with that infinite critical itinerary.

Neither statement follows from affine-sieve expansion, and the first cannot be reduced to any fixed finite collection of prime tests because of the rough pseudo-cycle theorem.


\bigskip



\subsection*{Final assessment}


The proposed premise\textemdash{}


\begin{quote}
the Collatz branch semigroup forms a thin orbit in a group to which standard affine-sieve expansion applies\textemdash{}

\end{quote}


is false. Its generated group has solvable, non-perfect Zariski closure, so the central expander input is unavailable. More importantly, the family

\[
w=(1^{j-1},A-j+1)
\]

produces arbitrarily long primitive positive rational cycles that are:

\[
\text{arbitrarily high}
\;+\;
\text{arbitrarily near critical drift}
\;+\;
\text{integral at every prescribed finite set of primes}.
\]

Hence a finite local affine sieve cannot distinguish them from an integer counterexample. The exact missing ingredient is a global archimedean\textendash{}profinite rigidity theorem, not a stronger ordinary sieve estimate.

Accordingly, claiming a full Collatz proof from approach 53 would conceal the unresolved step inside an unjustified local-to-global assertion. The rigorous result of the investigation is instead a no-go theorem for the standard thin-group affine-sieve strategy and an explicit family of its surviving local patterns.
\subsection*{References}
\begin{enumerate}[leftmargin=2.2em]
\item \hypertarget{ref:R58:1}{}\hyperlink{cit:R58:1}{\texttt{turn318263search12}}
\item \hypertarget{ref:R58:2}{}\hyperlink{cit:R58:2}{\texttt{turn356521view1}}
\item \hypertarget{ref:R58:3}{}\hyperlink{cit:R58:3}{\texttt{turn356521view3}}
\item \hypertarget{ref:R58:4}{}\hyperlink{cit:R58:4}{\texttt{turn356521view0}}
\item \hypertarget{ref:R58:5}{}\hyperlink{cit:R58:5}{\texttt{turn356521view2}}
\end{enumerate}

\clearpage
\section{Round 59: approach 54 --- Bourgain-Gamburd flattening for branch products}
\subsection*{Prompt}
\begin{quote}
Complete a full proof of the Collatz conjecture using the following approach:

\begin{enumerate}
\item Bourgain-Gamburd flattening for branch products
\end{enumerate}

\textbf{Speculative \textperiodcentered{} P}

Show that products of admissible inverse branches become equidistributed modulo suitable primes unless they lie in a structured subgroup. Then classify the structured case and rule it out by positivity or cycle equations.

Take as long as you need.
\end{quote}
\subsection*{Response}

\subsection*{Outcome}


I do \textbf{not} obtain a valid proof of the Collatz conjecture from this approach. The obstruction is exact rather than merely technical:


\begin{quote}
Every accelerated inverse-branch product lies in the structured subgroup from the outset.

\end{quote}


More precisely, all branch matrices lie in a solvable affine/Borel subgroup, every genuine trajectory is a moving conjugate of a diagonal cocycle, and the branch semigroup remains free inside that solvable group. Consequently, the Bourgain\textendash{}Gamburd alternative ``flatten unless trapped in a proper subgroup'' degenerates: \textbf{all the mass is permanently trapped in the proper subgroup}.

Current references still describe Collatz as unresolved; Tao's major result is an almost-all-orbits theorem rather than pointwise convergence for every starting value. [\hypertarget{cit:R59:1}{}\hyperlink{ref:R59:1}{1}, \hypertarget{cit:R59:2}{}\hyperlink{ref:R59:2}{2}]

What follows is a complete proof of the obstruction, a classification of the structured subgroup, an exact spectral no-go theorem, and the strongest surviving fibrewise-flattening statement.


\bigskip



\subsubsection*{1. Exact branch-product algebra}


Work with the accelerated odd Collatz map

\[
U(n)=\frac{3n+1}{2^{v_2(3n+1)}}.
\]

If \(a=v_2(3n+1)\), the corresponding inverse branch is

\[
T_a(x)=\frac{2^a x-1}{3},
\qquad a\geq 1.
\]

Represent it by

\[
B_a=
\begin{pmatrix}
2^a&-1\\
0&3
\end{pmatrix}.
\]

For a word

\[
w=(a_0,\ldots ,a_{j-1}),
\]

write

\[
A_i=\sum_{k<i}a_k,\qquad A=A_j,
\]

and define

\[
C_w=\sum_{i=0}^{j-1}3^{j-1-i}2^{A_i}.
\]


\paragraph*{Proposition 1: product formula}


\[
B_w:=B_{a_0}\cdots B_{a_{j-1}}
=
\begin{pmatrix}
2^A&-C_w\\
0&3^j
\end{pmatrix}.
\]

Hence the inverse branch determined by \(w\) is

\[
T_w(x)=\frac{2^A x-C_w}{3^j}.
\]


\paragraph*{Proof}


Suppose a prefix of length \(\ell\) has exponent sum \(S\) and carry \(C\). Appending an exponent \(a\) gives

\[
\begin{pmatrix}
2^S&-C\\
0&3^\ell
\end{pmatrix}
\begin{pmatrix}
2^a&-1\\
0&3
\end{pmatrix}
=
\begin{pmatrix}
2^{S+a}&-\bigl(3C+2^S\bigr)\\
0&3^{\ell+1}
\end{pmatrix}.
\]

Thus

\[
C_{\mathrm{new}}=3C+2^S,
\]

which iterates to the stated formula. \ensuremath{\blacksquare}


\bigskip



\subsubsection*{2. Reduction modulo primes always lands in an affine group}


Let \(p\neq 2,3\). In \(\operatorname{PGL}_2(\mathbb F_p)\), normalize \(B_a\) to the affine transformation

\[
g_a(x)=\frac{2^a}{3}x-\frac13.
\]

Write an affine transformation as a pair \((r,t)\), meaning \(x\mapsto rx+t\), with multiplication

\[
(r,t)(s,u)=(rs,ru+t).
\]

Then

\[
g_a=\left(\frac{2^a}{3},-\frac13\right).
\]

Therefore every branch product lies in

\[
\operatorname{Aff}(\mathbb F_p)
=
\mathbb F_p\rtimes\mathbb F_p^\times,
\]

the upper-triangular Borel subgroup of \(\operatorname{PGL}_2(\mathbb F_p)\).

This group is solvable. The classical Bourgain\textendash{}Gamburd expansion mechanism concerns non-elementary or Zariski-dense subgroups of semisimple groups such as \(\mathrm{SL}_2\); superstrong-approximation formulations similarly require a semisimple Zariski closure. [\hypertarget{cit:R59:3}{}\hyperlink{ref:R59:3}{3}, \hypertarget{cit:R59:4}{}\hyperlink{ref:R59:4}{4}] The affine group is instead the standard solvable \(ax+b\) group, treated separately from the simple-group growth setting. [\hypertarget{cit:R59:5}{}\hyperlink{ref:R59:5}{5}, \hypertarget{cit:R59:6}{}\hyperlink{ref:R59:6}{6}]

Thus the first necessary hypothesis for ordinary Bourgain\textendash{}Gamburd flattening fails for \textbf{every} Collatz branch word and \textbf{every} auxiliary prime.


\bigskip



\subsubsection*{3. The generated subgroup can be classified exactly}


Let \(E\subseteq\mathbb N_{\geq1}\) be the permitted set of exponents. Define

\[
G_{p,E}=\langle g_a:a\in E\rangle
\]

and

\[
H_{p,E}
=
\left\langle
2^{b-a},\,\frac{2^a}{3}
:
a,b\in E
\right\rangle
\leq \mathbb F_p^\times.
\]


\paragraph*{Proposition 2: structured-subgroup classification}


Assume there exist \(a,b\in E\) with

\[
2^{b-a}\not\equiv1\pmod p.
\]

Then

\[
G_{p,E}
=
\mathbb F_p\rtimes H_{p,E}.
\]


\paragraph*{Proof}


First,

\[
g_a^{-1}g_b=(2^{b-a},0).
\]

Hence the group contains a nontrivial dilation \(D_u=(u,0)\), where \(u=2^{b-a}\neq1\).

Let

\[
g_a=(r,t),\qquad t=-\frac13\neq0.
\]

Then

\[
D_u g_aD_u^{-1}g_a^{-1}
=
(1,(u-1)t).
\]

This is a nonzero translation. Since the additive group of the prime field \(\mathbb F_p\) is generated by any nonzero element, \(G_{p,E}\) contains every translation \((1,x)\).

Multiplying the generators by suitable translations isolates all their multiplier parts, so the multiplier subgroup is precisely \(H_{p,E}\). Therefore

\[
G_{p,E}=\mathbb F_p\rtimes H_{p,E}.
\qquad\Box
\]

For the complete exponent set \(E=\mathbb N_{\geq1}\),

\[
H_{p,E}=\langle 2,3\rangle\leq\mathbb F_p^\times.
\]

This is a complete classification, but it does not produce a useful exceptional family: the ``structured case'' is the entire branch system.


\bigskip



\subsubsection*{4. Fixed exponent sum gives an even stronger obstruction}


For fixed length \(j\) and fixed total exponent \(A\), Proposition 1 gives

\[
B_w=
\begin{pmatrix}
2^A&-C_w\\
0&3^j
\end{pmatrix}.
\]

Modulo \(p\), every such word has the same multiplier

\[
r_{j,A}=2^A3^{-j}.
\]

Thus all products in the \((j,A)\)-fibre lie in a single coset of the translation subgroup:

\[
\{(r_{j,A},t):t\in\mathbb F_p\}.
\]

They cannot become equidistributed in the full affine group.

More quantitatively, suppose a probability measure on

\[
\mathbb F_p\rtimes H
\]

is supported on at most \(m\) multiplier fibres. Its support has size at most \(mp\), and therefore

\[
\|\mu\|_2\geq \frac1{\sqrt{mp}}.
\]

The uniform measure on the entire group has norm

\[
\|u_G\|_2=\frac1{\sqrt{p|H|}}.
\]

Consequently,

\[
\frac{\|\mu\|_2}{\|u_G\|_2}
\geq
\sqrt{\frac{|H|}{m}}.
\]

For a fixed \((j,A)\)-fibre, \(m=1\). Even perfect uniformity of the carry coordinate leaves an \(L^2\)-distance factor of at least \(\sqrt{|H|}\) from full-group uniformity.


\bigskip



\subsubsection*{5. Genuine orbit products are diagonal coboundaries}


The obstruction becomes still more rigid along an actual Collatz trajectory.

Let

\[
n_{i+1}=U(n_i),
\qquad
2^{a_i}n_{i+1}=3n_i+1.
\]

Define

\[
L_x=
\begin{pmatrix}
1&x\\
0&1
\end{pmatrix},
\qquad
D_a=
\begin{pmatrix}
2^a&0\\
0&3
\end{pmatrix}.
\]

Then

\[
L_{n_i}D_{a_i}L_{-n_{i+1}}
=
\begin{pmatrix}
2^{a_i}&3n_i-2^{a_i}n_{i+1}\\
0&3
\end{pmatrix}
=
B_{a_i}.
\]

Therefore the product telescopes:

\[
B_{a_0}\cdots B_{a_{j-1}}
=
L_{n_0}
\begin{pmatrix}
2^{A_j}&0\\
0&3^j
\end{pmatrix}
L_{-n_j}.
\]

Equivalently,

\[
2^{A_j}n_j=3^j n_0+C_{w_j}.
\]

So every realized trajectory is a diagonal cocycle conjugated by the moving endpoints \(n_i\). It does not merely occasionally enter a torus-like structured configuration: it is always of that form.

This also shows why adding finitely many automaton states does not automatically repair the problem. If a state-dependent lift has labels

\[
\widetilde B_{s\to t,a}=P_s^{-1}B_aP_t,
\]

then along a composable path,

\[
\widetilde B_{s_0\to s_1,a_0}\cdots
\widetilde B_{s_{j-1}\to s_j,a_{j-1}}
=
P_{s_0}^{-1}B_wP_{s_j}.
\]

The internal conjugations telescope. Noncomposable lifted generators might generate a large semisimple group, but their expansion would concern words that do not correspond to Collatz paths.


\bigskip



\subsubsection*{6. ``Structured'' does not mean periodic: the semigroup is free}


Although the matrices are triangular, distinct valuation words remain arithmetically distinguishable.

For \(w=(a_0,\ldots ,a_{j-1})\),

\[
C_w=3^{j-1}+2^{a_0}C_{(a_1,\ldots ,a_{j-1})}.
\]

The tail carry is odd. Hence

\[
a_0=v_2(C_w-3^{j-1}).
\]

After recovering \(a_0\), set

\[
C^{(1)}
=
\frac{C_w-3^{j-1}}{2^{a_0}}
\]

and repeat. This recovers \(a_0,\ldots ,a_{j-2}\); the total \(A\) recovers \(a_{j-1}\).

Thus

\[
w\longmapsto (j,A,C_w)
\]

is injective. The upper-triangular semigroup contains a free branch semigroup of enormous combinatorial complexity. Solvability of the ambient algebraic group therefore does not reduce words to a finite collection of periodic or virtually abelian patterns.


\bigskip



\subsubsection*{7. There is no uniform spectral gap even on the multiplier quotient}


Consider the geometric branch measure

\[
\mu_p=\sum_{a\geq1}2^{-a}\delta_{g_a}.
\]

Project it to the multiplier subgroup:

\[
\nu_p
=
\sum_{a\geq1}2^{-a}
\delta_{2^a/3}.
\]

Let

\[
r=\operatorname{ord}_p(2).
\]

Choose a multiplicative character \(\chi\) of the multiplier subgroup whose restriction to \(\langle2\rangle\) satisfies

\[
\chi(2)=z=e^{2\pi i/r}.
\]

Then

\[
\widehat{\nu_p}(\chi)
=
\chi(3)^{-1}
\sum_{a\geq1}2^{-a}z^a
=
\chi(3)^{-1}\frac{z}{2-z}.
\]

Therefore

\[
\left|\widehat{\nu_p}(\chi)\right|
=
\frac1{|2-z|}
=
\frac1{\sqrt{1+8\sin^2(\pi/r)}}.
\]

As \(r\to\infty\),

\[
\left|\widehat{\nu_p}(\chi)\right|
=
1-\frac{4\pi^2}{r^2}+O(r^{-4}).
\]

The orders \(\operatorname{ord}_p(2)\) are unbounded: otherwise every prime would divide the fixed finite product

\[
\prod_{r\leq R}(2^r-1)
\]

for some \(R\), which is impossible.

Hence there is no prime-uniform spectral gap for this branch measure, even after quotienting out the translations.

There are genuine mixing theorems for random \(ax+b\) processes modulo many primes, but they are action-level, averaged random-walk statements rather than Bourgain\textendash{}Gamburd expansion of the full branch group. [\hypertarget{cit:R59:7}{}\hyperlink{ref:R59:7}{7}] The explicit character above identifies the slow mode that any Collatz application would have to overcome.


\bigskip



\subsubsection*{8. What does flatten: the carry coordinate inside one fibre}


The surviving promising phenomenon is \textbf{fibrewise translation mixing}.

Fix \(j,A,p\), and choose uniformly among the

\[
N_{j,A}=\binom{A-1}{j-1}
\]

positive compositions of \(A\) into \(j\) parts. Let \(\mu_{j,A,p}\) be the distribution of \(C_w\bmod p\).

This distribution can be computed exactly. Let \(F_{\ell,s}(c)\) count prefixes of length \(\ell\), total exponent \(s\), and carry \(c\bmod p\). Then

\[
F_{0,0}(0)=1
\]

and, because appending an exponent changes \(C\) to \(3C+2^s\),

\[
F_{\ell+1,t}(y)
=
\sum_{s<t}
F_{\ell,s}\!\left(3^{-1}(y-2^s)\right).
\]

For

\[
(j,A)=(20,32),
\qquad
N_{20,32}=\binom{31}{19}=141\,120\,525,
\]

exact integer computation gives:


\begin{center}
\begin{tabular}{@{}p{0.235\linewidth}p{0.235\linewidth}p{0.235\linewidth}p{0.235\linewidth}@{}}
\hline
\(p\) & \(\min_c \mu(c)/(1/p)\) & \(\max_c \mu(c)/(1/p)\) & \(\sqrt p\,\|\mu\|_2\) \\
\hline
101 & 0.99853645 & 1.00170485 & 1.00000018 \\
251 & 0.99859281 & 1.00153643 & 1.00000017 \\
509 & 0.99701519 & 1.00330192 & 1.00000069 \\
1009 & 0.99326643 & 1.00647948 & 1.00000180 \\
\hline
\end{tabular}
\end{center}


The value of \(\sqrt p\,\|\mu\|_2\) is exactly \(1\) for a uniform distribution. Thus the carry is strikingly close to uniform modulo these auxiliary primes.

This is strong evidence for an affine fibre-flattening theorem. But it has the wrong logical direction for extinction: in every row of the table, \textbf{every residue occurs}, and occurs many times. Stronger uniformity says that prescribed congruences have approximately \(N_{j,A}/p\) solutions; it does not say they have none.


\bigskip



\subsubsection*{9. The cycle case reduces exactly to the original cycle equation}


Set

\[
D_{j,A}=2^A-3^j.
\]

A word \(w\) has a positive fixed point precisely when

\[
n=T_w(n)=\frac{2^A n-C_w}{3^j},
\]

or equivalently,

\[
C_w=D_{j,A}n.
\]

Thus positive cycles require

\[
D_{j,A}>0,
\qquad
D_{j,A}\mid C_w.
\]

This condition is not merely necessary; it already constructs the cycle.


\paragraph*{Proposition 3: divisibility is equivalent to a positive cycle}


Suppose \(D=2^A-3^j>0\). Then

\[
D\mid C_w
\]

if and only if the cyclic rotations of \(w\) generate a positive accelerated Collatz cycle.


\paragraph*{Proof}


Write \(w=(a_0,v)\), and let \(w'=(v,a_0)\) be the first cyclic rotation. Then

\[
C_w=3^{j-1}+2^{a_0}C_v,
\]

whereas

\[
C_{w'}=3C_v+2^{A-a_0}.
\]

Consequently,

\[
2^{a_0}C_{w'}
=
3C_w+(2^A-3^j)
=
3C_w+D.
\]

Because \(D\) is odd, \(D\mid C_w\) implies \(D\mid C_{w'}\). Iterating proves divisibility for every cyclic rotation.

Define

\[
n_i=\frac{C_{\operatorname{rot}^i(w)}}D.
\]

Every \(C_{\operatorname{rot}^i(w)}\) and \(D\) is positive and odd, so each \(n_i\) is a positive odd integer. Dividing the rotation identity by \(D\) gives

\[
2^{a_i}n_{i+1}=3n_i+1.
\]

Since \(n_{i+1}\) is odd, \(a_i\) is exactly the required \(2\)-adic valuation. Thus these integers form an accelerated Collatz cycle.

The converse follows immediately by multiplying the cycle equations around the word. \ensuremath{\blacksquare}

For the word

\[
w=(2,\ldots ,2),
\]

one has

\[
C_w
=
\sum_{i=0}^{j-1}3^{j-1-i}4^i
=
4^j-3^j
=
D,
\]

giving the trivial cycle \(n=1\).

Therefore the proposed instruction ``classify the structured case and rule it out by the cycle equations'' becomes exactly:

\[
2^A-3^j\mid C_w
\quad\Longrightarrow\quad
w=(2,\ldots ,2)
\ \text{up to repetition and rotation}.
\]

That statement is itself the complete nontrivial-cycle problem. The subgroup classification does not make it easier.


\bigskip



\subsubsection*{10. Eventual periodicity also gives only a cycle}


Suppose that, from some point onward, a block of \(r\) exponents with total \(A\) repeats. The corresponding block map is

\[
F(x)=\frac{3^r x+C}{2^A}.
\]

Let \(x_m=F^m(x_0)\), and put

\[
D=2^A-3^r.
\]

Then

\[
D x_{m+1}-C
=
\frac{3^r}{2^A}(D x_m-C),
\]

so

\[
2^{Am}(D x_m-C)
=
3^{rm}(D x_0-C).
\]

All \(x_m\) are integers. Since \(2\) and \(3\) are coprime,

\[
2^{Am}\mid D x_0-C
\]

for every \(m\). Hence

\[
D x_0-C=0.
\]

Thus \(x_0\) is already a fixed point of the repeated block. An eventually periodic valuation sequence realized by a positive integer cannot encode a genuinely divergent orbit; it encodes a cycle.

So even a hypothetical inverse theorem saying ``nonflattening forces periodicity'' would still leave Proposition 3's divisibility problem to solve.


\bigskip



\subsubsection*{11. Why auxiliary-prime flattening cannot handle a divergent orbit}


For an actual prefix \(w_j\),

\[
2^{A_j}n_j=3^j n_0+C_{w_j}.
\]

Modulo a prime \(p\neq2,3\), this merely relates the two moving endpoint residues \(n_0,n_j\) to \(C_{w_j}\). There is no forbidden residue for the carry: the unknown endpoint absorbs it.

The genuinely restrictive modulus is \(3^j\). Given an odd endpoint \(X\), a word \(v\) is an integral inverse path precisely when

\[
C_v\equiv2^{A(v)}X\pmod{3^j}.
\]

At \(p=3\), however, the matrices

\[
B_a\bmod3
=
\begin{pmatrix}
2^a&-1\\
0&0
\end{pmatrix}
\]

are singular. There is no finite-group Bourgain\textendash{}Gamburd theory to invoke. Products successively lose information rather than expand.

This is also exactly where the entropy is on the wrong side. For

\[
A\sim\alpha j,
\]

the number of words satisfies

\[
\binom{A-1}{j-1}
=
\exp\!\left(
j\bigl[
\alpha\log\alpha-(\alpha-1)\log(\alpha-1)
\bigr]
+o(j)
\right).
\]

At the critical Collatz drift

\[
\lambda=\log_2 3,
\]

the exponent is

\[
g(\lambda)
=
\lambda\log\lambda-(\lambda-1)\log(\lambda-1)
=
1.043632815\ldots ,
\]

whereas

\[
\log 3=1.098612288\ldots .
\]

Thus

\[
\binom{\lambda j+O(1)}{j}
\asymp
(2.8395137\ldots)^j,
\]

while the relevant \(3\)-adic modulus has size

\[
3^j.
\]

The ratio decays as

\[
\left(\frac{2.8395137\ldots}{3}\right)^j
=
(0.9465045\ldots)^j.
\]

Therefore a critical fibre is too small to cover the \(3^j\) residue classes. Average equidistribution modulo smaller auxiliary primes cannot rule out one exceptional word hitting the required \(3\)-adic class at every depth.

That is precisely the distributional-to-pointwise gap.


\bigskip



\subsubsection*{12. The exact missing statement}


Assume a nonconvergent orbit exists, and choose the least positive odd starting value \(n_0\) among all such orbits. Its odd iterates satisfy

\[
n_j\geq n_0
\]

for every \(j\).

Let \(w_j\) be its length-\(j\) valuation prefix, with total exponent \(A_j\). Suppose there were another length-\(j\) word \(v\), with the same total exponent, satisfying

\[
C_v=C_{w_j}+q3^j,
\qquad
1\leq q<n_0.
\]

Then

\[
\frac{2^{A_j}n_j-C_v}{3^j}
=
\frac{2^{A_j}n_j-C_{w_j}}{3^j}-q
=
n_0-q.
\]

The congruence ensures that all inverse steps are integral; positivity of the final predecessor ensures positivity of the intermediate inverse steps. Hence \(n_0-q\) is a smaller positive integer whose orbit merges into the nonconvergent tail at \(n_j\), contradicting minimality.

Accordingly, the missing theorem would have to be something like:


\begin{quote}
\textbf{Ordered \(3\)-adic collision-or-periodicity theorem.}   Every infinite positive critical valuation word either has a prefix \(w_j\) admitting another word \(v\) with \[
> A(v)=A(w_j),\qquad
> C_v-C_{w_j}\in\{3^j,2\cdot3^j,\ldots,(n_0-1)3^j\},
> \] or the valuation sequence is eventually periodic.

\end{quote}


The collision alternative gives a smaller counterexample; the periodic alternative gives a cycle.

Bourgain\textendash{}Gamburd flattening modulo auxiliary primes proves neither the ordered sign condition \(C_v>C_w\) nor equality modulo the singular modulus \(3^j\). The numerical fibre mixing above supplies no bridge to this statement.


\bigskip



\subsection*{Final assessment}


Approach 54 yields three rigorous conclusions:

\begin{enumerate}
\item \textbf{The usual Bourgain\textendash{}Gamburd mechanism is algebraically unavailable.} Every branch product lies in a solvable affine subgroup, and every actual trajectory is a moving conjugate of a diagonal cocycle.

\item \textbf{Translation-coordinate flattening is real and very strong.} Exact computation shows near-perfect equidistribution of \(C_w\bmod p\) in fixed \((j,A)\)-fibres for auxiliary primes well below the fibre size.

\item \textbf{That flattening is not pointwise enough.} The relevant obstruction lives modulo \(3^j\), where the matrices are singular and the critical word count is exponentially smaller than the modulus. The structured periodic case reduces exactly to \[
   2^A-3^j\mid C_w,
   \] which is the nontrivial-cycle problem itself.
\end{enumerate}

Therefore this investigation does \textbf{not} complete a proof of Collatz. It identifies a precise no-go for the literal Bourgain\textendash{}Gamburd plan and isolates the genuinely new ingredient required: an ordered, pointwise \(3\)-adic collision-or-rigidity theorem. Standard affine flattening, even in an optimally strong form, does not supply that theorem.
\subsection*{References}
\begin{enumerate}[leftmargin=2.2em]
\item \hypertarget{ref:R59:1}{}\hyperlink{cit:R59:1}{\texttt{turn878669search5}}
\item \hypertarget{ref:R59:2}{}\hyperlink{cit:R59:2}{\texttt{turn878669search27}}
\item \hypertarget{ref:R59:3}{}\hyperlink{cit:R59:3}{\texttt{turn228330view0}}
\item \hypertarget{ref:R59:4}{}\hyperlink{cit:R59:4}{\texttt{turn228330view4}}
\item \hypertarget{ref:R59:5}{}\hyperlink{cit:R59:5}{\texttt{turn228330view2}}
\item \hypertarget{ref:R59:6}{}\hyperlink{cit:R59:6}{\texttt{turn228330view3}}
\item \hypertarget{ref:R59:7}{}\hyperlink{cit:R59:7}{\texttt{turn875948view0}}
\end{enumerate}

\clearpage
\section{Round 60: approach 55 --- Sum-product growth modulo 3\textasciicircum{}r}
\subsection*{Prompt}
\begin{quote}
Complete a full proof of the Collatz conjecture using the following approach:

\begin{enumerate}
\item Sum-product growth modulo 3\textasciicircum{}r
\end{enumerate}

\textbf{Hybrid \textperiodcentered{} P}

The possible carry residues generated by many prefixes should expand rapidly under additions and multiplication by powers of 2 and 3. Prove that a critical survivor set cannot remain inside the small collection of residues needed for repeated 3-adic rescue.

Take as long as you need.
\end{quote}
\subsection*{Response}

\subsection*{Result of the sum\textendash{}product attack}


I do \textbf{not} obtain a valid full proof of the Collatz conjecture from Approach 55. The obstruction is now precise: the proposed principle


\begin{quote}
``a critical survivor set must expand modulo \(3^r\)''

\end{quote}


is false without a substantial extra hypothesis. The carry dynamics are \(3\)-adically \textbf{contracting}, and a hypothetical obstruction initially supplies a single nested, graph-like ray\textemdash{}not the large, multiscale nonconcentrated family required by sum\textendash{}product theorems.

What follows is the rigorous reduction, several exact new lemmas, an explicit failure of naive residue injectivity, and the exact amplification theorem that would be needed to repair the approach.


\bigskip



\subsubsection*{1. Exact carry dynamics}


For a valuation word

\[
w=(a_0,\ldots ,a_{r-1}),\qquad a_i\ge 1,
\]

write

\[
A_k=\sum_{i<k}a_i,\qquad Q_k=2^{A_k},
\]

and define

\[
C_0=0,\qquad C_{k+1}=3C_k+Q_k.
\]

The standard affine identity is

\[
2^{A_r}U^r(n)=3^rn+C_r.
\]

For concatenated words \(w=uv\),

\[
C_{uv}=3^{|v|}C_u+2^{A(u)}C_v.
\]

These identities, including the exact terminal-suffix locality theorem, have already been proved in the existing framework. [\hypertarget{cit:R60:1}{}\hyperlink{ref:R60:1}{1}] [\hypertarget{cit:R60:2}{}\hyperlink{ref:R60:2}{2}]

Set

\[
B_k=C_k+Q_k.
\]

Then

\[
\boxed{
B_{k+1}=3B_k+(2^{a_k}-2)Q_k,
\qquad
Q_{k+1}=2^{a_k}Q_k.
}
\tag{1}
\]

The target-one terminal condition is simply

\[
\boxed{3^r\mid B_r.}
\tag{2}
\]

The strict prefix barrier is

\[
\boxed{Q_k<3^k\qquad(1\le k\le r).}
\tag{3}
\]

For a shortest nontrivial word satisfying (2)\textendash{}(3), the existing minimality argument gives a strict terminal record,

\[
A_i-i\log_2 3<A_r-r\log_2 3<0\qquad(i<r),
\]

as well as

\[
A_r=\lfloor r\log_2 3\rfloor,
\qquad a_0=1,\qquad a_{r-1}>1.
\tag{4}
\]

The corresponding \(3x-1\) trajectory has endpoint \(1\), and if

\[
m_0=\frac{B_r}{3^r}=1+\Delta,
\qquad
h=\#\{i:a_i>1\},
\]

then the activation identity gives

\[
\boxed{0<\Delta<h.}
\tag{5}
\]

These minimal-word and activation statements have been independently audited. [\hypertarget{cit:R60:3}{}\hyperlink{ref:R60:3}{3}]

Thus the terminal problem is equivalent to proving that the arithmetic divisibility in (2) cannot occur for a nontrivial strict-record word.


\bigskip



\subsubsection*{2. The favorable entropy gap}


Put

\[
\alpha=\log_2 3,\qquad \beta=\alpha-1=\log_2(3/2).
\]

Writing \(b_i=a_i-1\ge0\), the barrier implies

\[
\sum_{i<r}b_i\le \lfloor\beta r\rfloor.
\]

Therefore the number \(N_r\) of barrier words is at most

\[
N_r
\le
\binom{r+\lfloor\beta r\rfloor}{r}.
\]

By Stirling's formula,

\[
N_r\le r^{O(1)}\lambda^r,
\]

where

\[
\lambda=
\frac{\alpha^\alpha}{\beta^\beta}
=
2.839513730497\ldots <3.
\tag{6}
\]

Consequently,

\[
\frac{N_r}{3^r}
\le
r^{O(1)}
\left(\frac{2.839514}{3}\right)^r.
\tag{7}
\]

So an equidistributed residue map would have exponentially fewer than one word per residue on average. This is the genuine reason the approach initially looks promising.

But average scarcity is not zero-fibre rigidity: an exceptional structured residue can contain one word forever. The all-ones word is exactly such an exception.


\bigskip



\subsubsection*{3. Terminality is nested, not independent}



\paragraph*{Lemma 1 \textemdash{} Every suffix is terminal}


Let \(w=uv\), with \(|u|=n\) and \(|v|=s\). Then

\[
B_{uv}=3^sC_u+2^{A(u)}B_v.
\tag{8}
\]

If \(3^{n+s}\mid B_{uv}\), reduction modulo \(3^s\) gives

\[
3^s\mid B_v.
\]

Thus:

\[
\boxed{\text{Every suffix of a terminal word is terminal at its own depth.}}
\tag{9}
\]

This is an important correction to the probabilistic picture. The \(r\) rescue conditions are not \(r\) roughly independent events. They form one triangular tower of exact divisibilities.

If

\[
B_v=3^s m_v,
\]

then (8) becomes

\[
B_{uv}
=
3^s\bigl(C_u+2^{A(u)}m_v\bigr).
\]

Hence

\[
\boxed{
3^{n+s}\mid B_{uv}
\iff
C_u+2^{A(u)}m_v\equiv0\pmod{3^n}.
}
\tag{10}
\]

Equation (10) is the exact point at which a sum\textendash{}product projection enters:

\[
a+cb\equiv0\pmod{3^n},
\]

with

\[
a=C_u,\qquad c=2^{A(u)},\qquad b=m_v.
\]

This is the strongest clean formulation of Approach 55.


\bigskip



\subsubsection*{4. The rescue relation is a graph}


Fix a prefix \(u\) of length \(r-1\), and append an exponent \(a\). From (1),

\[
B_{ua}=3C_u+Q_u(1+2^a).
\]

Terminality is equivalent to

\[
2^a
\equiv
-1-3C_uQ_u^{-1}
\pmod{3^r}.
\tag{11}
\]

Now \(2\) generates \((\mathbb Z/3^r\mathbb Z)^\times\), with

\[
\operatorname{ord}_{3^r}(2)=2\cdot3^{r-1}.
\]

Thus (11) determines \textbf{at most one} class

\[
a\pmod{2\cdot3^{r-1}}.
\]

The barrier permits only \(a=O(r)\), vastly shorter than the period \(2\cdot3^{r-1}\). Therefore:

\[
\boxed{
\text{For every admissible prefix, there is at most one admissible terminal extension.}
}
\tag{12}
\]

So the rescue set is not a thick subset of a residue ring. It is the graph of an extremely sparse function from prefixes to exponents.

That is almost the opposite of the hypothesis required for an ordinary sum\textendash{}product argument.


\bigskip



\subsubsection*{5. The branch maps contract \(3\)-adically}


Normalize by the unit \(Q_k\):

\[
z_k=B_kQ_k^{-1}\pmod{3^r}.
\]

Equation (1) becomes

\[
\boxed{
z_{k+1}
=
2^{-a_k}\bigl(3z_k+2^{a_k}-2\bigr).
}
\tag{13}
\]

For two states using the same branch \(a\),

\[
z_{k+1}-z'_{k+1}
=
3\cdot2^{-a}(z_k-z'_k).
\]

Hence

\[
v_3(z_{k+1}-z'_{k+1})
=
v_3(z_k-z'_k)+1.
\tag{14}
\]

Each individual branch is therefore a strict \(3\)-adic contraction. Branching over many possible values of \(a\) may enlarge a union of images, but a single survivor ray becomes progressively \textbf{more concentrated}, not less.

This destroys a direct ``iteration implies expansion'' argument.


\bigskip



\subsubsection*{6. Two exact structured survivor rays}



\paragraph*{The trivial ray}


For \(w=1^r\),

\[
C_w=3^r-2^r,
\qquad
B_w=3^r.
\]

Thus \(1^r\) is terminal and satisfies the barrier at every depth.


\paragraph*{A nontrivial cycle ray}


For the accelerated \(3x-1\) map

\[
U_-(m)=\frac{3m-1}{2^{v_2(3m-1)}},
\]

one has

\[
5\xrightarrow{\,a=1\,}7
\xrightarrow{\,a=2\,}5.
\]

Consequently, the infinite exponent ray

\[
(1,2,1,2,\ldots)
\]

is periodic. It also satisfies every prefix barrier:

\[
A_{2n}=3n<2n\log_2 3,
\]

and

\[
A_{2n+1}=3n+1<(2n+1)\log_2 3.
\]

Therefore a nontrivial barrier survivor can remain trapped in a two-state structured component forever. It does not have to expand.

The target-one condition excludes this particular cycle, but the example proves that any successful theorem must be \textbf{component-sensitive}. ``Barrier plus persistence implies expansion'' is false.


\bigskip



\subsubsection*{7. Naive residue injectivity is also false}


Define

\[
\Phi_r(w)=B_w\bmod3^r.
\]

Computationally, \(\Phi_r\) appeared injective on all barrier words through depth \(18\). That pattern breaks at depth \(19\).

Consider

\[
w_1=(1^{15},2,3,1,3)
\]

and

\[
w_2=(1,2,1,2,1,2,1,2,1,2,2,2,1,1,2,2,1,1,3).
\]

Both satisfy

\[
A_k\le\lfloor k\log_2 3\rfloor
\qquad(1\le k\le19).
\]

Exact evaluation gives

\[
3^{19}=1\,162\,261\,467,
\]

\[
B(w_1)=1\,183\,691\,739,
\]

and

\[
B(w_2)=5\,832\,737\,607.
\]

Moreover,

\[
B(w_2)-B(w_1)
=
4\cdot3^{19}.
\]

Thus

\[
\boxed{
\Phi_{19}(w_1)=\Phi_{19}(w_2)=21\,430\,272.
}
\tag{15}
\]

Neither residue is zero, so this does not produce a terminal counterexample. It does prove that a simple ``unique \(3\)-adic coding of barrier words'' theorem is false.

The stronger zero-fibre statement

\[
\Phi_r^{-1}(0)=\{1^r\}
\tag{16}
\]

remains consistent with all computation. Exact reverse enumeration has verified (16) through \(r=78\); at depth \(78\), more than \(642\) million reverse states were checked. [\hypertarget{cit:R60:4}{}\hyperlink{ref:R60:4}{4}]


\bigskip



\subsubsection*{8. Why available sum\textendash{}product theorems do not apply directly}


The current \(p\)-adic sum\textendash{}product results are genuinely powerful, but their hypotheses are not supplied by one Collatz survivor.

Shen's recent theorem proves expansion for sufficiently nonconcentrated subsets of \(\mathbb Z_p^\times\). Its ABC form controls projections

\[
a+cb
\]

when the \(b\)-set and coefficient set \(c\) have positive multiscale dimension and when one has a dense graph of pairs \((a,b)\). [\hypertarget{cit:R60:5}{}\hyperlink{ref:R60:5}{5}]

The Collatz situation fails these hypotheses in three separate ways.


\paragraph*{8.1 One bad word gives one incidence}


At a cut \(w=uv\), equation (10) gives one pair

\[
(C_u,m_v)
\]

and one coefficient \(2^{A(u)}\). It does not produce a dense graph in \(A\times B\).

A theorem about large sets cannot exclude one exceptional edge unless that edge can first be \textbf{amplified} into many related edges.


\paragraph*{8.2 The coefficient set has dimension zero}


At all cuts of a length-\(r\) word, the possible coefficients are

\[
2^{A(u)},\qquad 0\le A(u)\le \alpha r.
\]

There are only \(O(r)\) of them.

At a modulus \(3^n\) with \(n\asymp r\), their normalized dimension is

\[
\frac{\log O(r)}{\log 3^n}
=
O\!\left(\frac{\log r}{r}\right)
\longrightarrow0.
\]

Shen's ABC theorem requires a coefficient set with a fixed positive multiscale dimension. The Collatz coefficients are far too sparse at the full terminal scale.

One can lower the modulus to \(3^{O(\log r)}\), where \(O(r)\) coefficients become substantial. But that controls only \(O(\log r)\) ternary digits, not the required \(r\) digits.

This is not a technical distinction. There are strict-record words of length approximately \(3^q\) that imitate terminality through exactly \(q\) digits. Thus every bounded-depth or logarithmic-depth argument has explicit adversarial examples. [\hypertarget{cit:R60:6}{}\hyperlink{ref:R60:6}{6}]


\paragraph*{8.3 The full barrier set is still too small for large-set estimates}


Finite-valuation-ring estimates such as Pham's require sets whose cardinalities are very close to the size \(3^r\) of the ambient ring\textemdash{}for example, hypotheses of the form

\[
|A|\gg 3^{r-1/3}.
\]

The barrier language has only about

\[
2.8395^r=3^{0.94996\ldots r}
\]

elements before any terminal restrictions are imposed, so these large-set estimates do not apply. [\hypertarget{cit:R60:7}{}\hyperlink{ref:R60:7}{7}]


\bigskip



\subsubsection*{9. The exact missing amplification theorem}


The sum\textendash{}product strategy could be repaired by proving the following genuinely new statement.


\paragraph*{Component-sensitive amplification lemma}


There exist constants

\[
\kappa,\sigma,\eta>0
\]

such that any nontrivial strict-record terminal word \(w\) of length \(r\), whose reverse trajectory lies in the \(3x-1\) component rooted at \(1\), produces some scale

\[
n\ge\kappa r
\]

and sets

\[
A,B\subset\mathbb Z/3^n\mathbb Z,
\qquad
\Theta\subset(\mathbb Z/3^n\mathbb Z)^\times
\]

with the following properties:

\begin{enumerate}
\item \(B\) and \(\Theta\) are nonconcentrated at every \(3\)-adic scale and have size at least \(3^{\sigma n}\).

\item There is a dense graph \(G\subset A\times B\).

\item For every \(\theta\in\Theta\), the rescue projections

\end{enumerate}
\[
   \{a+\theta b:(a,b)\in G\}
   \]

remain inside a set of size at most

\[
   3^{\eta n}|A|,
   \]

where \(\eta\) is smaller than the expansion exponent supplied by the \(p\)-adic ABC theorem.

\begin{enumerate}
\item Periodic rays and rays lying in non-\(1\) components are excluded from this amplification assertion.
\end{enumerate}


\paragraph*{Conditional proof}


Assume the lemma. Shen's ABC theorem applied to \(A,B,\Theta,G\) produces some \(\theta\in\Theta\) for which

\[
|\{a+\theta b:(a,b)\in G\}|
\ge
3^{\chi n}|A|
\]

for a fixed \(\chi>\eta\).

But the rescue property gives the opposite estimate

\[
|\{a+\theta b:(a,b)\in G\}|
\le
3^{\eta n}|A|.
\]

This is a contradiction. Therefore no nontrivial strict-record terminal word exists.

That would establish the desired all-depth zero-fibre theorem

\[
\boxed{
3^r\mid C_w+2^{A(w)},\quad
2^{A_k}<3^k\ \forall k
\quad\Longrightarrow\quad
w=1^r.
}
\tag{17}
\]

The missing part is not the application of sum\textendash{}product. It is the amplification lemma that turns \textbf{one nested Collatz ray} into a large nonconcentrated incidence configuration.

The examples above show that no unconditional version of that lemma can hold:

\begin{itemize}
\item \(1^\infty\) is a singleton terminal ray;
\item \((1,2)^\infty\) is a nontrivial periodic barrier ray;
\item the terminal extension is a graph with outdegree at most one;
\item fixed-depth terminal mimicry persists to arbitrarily large word lengths;
\item residue coding is not globally injective.
\end{itemize}

Component sensitivity is therefore indispensable.


\bigskip



\subsubsection*{10. Why even (17) would not yet be a full Collatz proof}


Within the current framework, terminal exclusion is only one of three logically separate requirements. A complete argument must also:

\begin{enumerate}
\item exclude nontrivial positive \(3x+1\) cycles;

\item prove that every hypothetical unbounded orbit yields arbitrarily long \textbf{nontrivial, consecutive, aligned} terminal/barrier blocks;

\item apply the all-depth exclusion (17).
\end{enumerate}

The existing audit explicitly identifies these three obligations and notes that the cycle word has the opposite full-word slope from the terminal barrier, so terminal exclusion by itself does not settle cycles. [\hypertarget{cit:R60:8}{}\hyperlink{ref:R60:8}{8}]

The synchronization problem is also stated explicitly in the existing structural write-up: terminal locality is proved, but no theorem yet forces one positive integer on an infinite retained ray to supply the required aligned inverse path. [\hypertarget{cit:R60:9}{}\hyperlink{ref:R60:9}{9}]


\bigskip



\subsection*{Final assessment}


Approach 55 produces a clean and valuable reduction:

\[
\boxed{
\text{terminal Collatz obstruction}
\quad\Longrightarrow\quad
C_u+2^{A(u)}m_v\equiv0\pmod{3^{|u|}},
}
\]

which is exactly an affine \(a+cb\) incidence over \(\mathbb Z/3^n\mathbb Z\).

But the proposed expansion step is not presently valid. A counterexample candidate is one nested, contracting, graph-like ray, while available sum\textendash{}product theorems require large nonconcentrated sets and positive-dimensional coefficient families. Periodic \(3x-1\) rays demonstrate that this structured obstruction is real.

Therefore I cannot honestly state a full Collatz proof from this route. The exact new theorem needed is the \textbf{component-sensitive amplification lemma} above, followed by the still-separate cycle and synchronization theorems. Without those statements, invoking sum\textendash{}product would replace the central Collatz difficulty with an unproved assumption rather than resolve it.
\subsection*{References}
\begin{enumerate}[leftmargin=2.2em]
\item \hypertarget{ref:R60:1}{}\hyperlink{cit:R60:1}{\texttt{turn8file0 L100-L137}}
\item \hypertarget{ref:R60:2}{}\hyperlink{cit:R60:2}{\texttt{turn8file0 L204-L235}}
\item \hypertarget{ref:R60:3}{}\hyperlink{cit:R60:3}{\texttt{turn7file0 L14-L61}}
\item \hypertarget{ref:R60:4}{}\hyperlink{cit:R60:4}{\texttt{turn9file0 L29-L53}}
\item \hypertarget{ref:R60:5}{}\hyperlink{cit:R60:5}{\texttt{turn369625view0}}
\item \hypertarget{ref:R60:6}{}\hyperlink{cit:R60:6}{\texttt{turn7file1 L118-L187}}
\item \hypertarget{ref:R60:7}{}\hyperlink{cit:R60:7}{\texttt{turn232272view1}}
\item \hypertarget{ref:R60:8}{}\hyperlink{cit:R60:8}{\texttt{turn5file0 L74-L96}}
\item \hypertarget{ref:R60:9}{}\hyperlink{cit:R60:9}{\texttt{turn8file0 L48-L51}}
\end{enumerate}

\clearpage
\section{Round 61: approach 56 --- A cross-prime uncertainty principle}
\subsection*{Prompt}
\begin{quote}
Complete a full proof of the Collatz conjecture using the following approach:

\begin{enumerate}
\item A cross-prime uncertainty principle
\end{enumerate}

\textbf{Speculative \textperiodcentered{} P}

A counterexample prefix is sharply localized in a 2-adic seed cylinder and a 3-adic endpoint cylinder. Develop an uncertainty inequality saying a nonzero arithmetic measure cannot remain simultaneously localized at both primes over arbitrarily many scales.

Take as long as you need.
\end{quote}
\subsection*{Response}

\subsection*{Outcome}


Carrying the proposed uncertainty principle through exactly does \textbf{not} complete a proof of Collatz. The obstruction is stronger than a missing estimate:


\begin{quote}
\textbf{Dyadic seed localization and ternary endpoint localization are not competing localizations. For every fixed Collatz branch, they are exact affine images of one another.}

\end{quote}


Consequently, the natural cross-prime localization operator has norm \(1\), not \(<1\). The only valid uncertainty inequality also involving the real place gives \textbf{uniqueness or sparsity}, never automatic emptiness.

Below is the exact calculation, a sharp no-go theorem for the proposed principle, and the strongest corrected form that survives.


\bigskip



\subsubsection*{1. Exact branch notation}


Use the accelerated map on positive odd integers

\[
U(n)=\frac{3n+1}{2^{\nu_2(3n+1)}}.
\]

Let

\[
w=(a_0,\ldots ,a_{r-1}),\qquad a_i\ge 1
\]

be a proposed valuation word, and put

\[
A_0=0,\qquad A_k=\sum_{i<k}a_i,
\]

\[
C_0=0,\qquad C_{k+1}=3C_k+2^{A_k}.
\]

For every exact realization \(n,x_1,\ldots,x_r\) of \(w\),

\[
2^{A_k}x_k=3^k n+C_k. \tag{1}
\]

This is the established affine identity in the existing framework. [\hypertarget{cit:R61:1}{}\hyperlink{ref:R61:1}{1}]

Write \(A=A_r\) and \(C=C_r\).


\bigskip



\subsection*{2. The exact cylinder-isomorphism theorem}



\subsubsection*{Theorem 1}


For every word \(w\), there is exactly one odd residue class

\[
\alpha_w\pmod {2^{A+1}}
\]

whose elements realize \(w\), namely

\[
\boxed{\quad
\alpha_w\equiv 3^{-r}(2^A-C)\pmod {2^{A+1}}.
\quad} \tag{2}
\]

There is also exactly one odd endpoint class

\[
\beta_w\pmod {2\cdot 3^r},
\]

characterized by

\[
\boxed{\quad
\beta_w\equiv C\,2^{-A}\pmod {3^r},
\qquad
\beta_w\equiv1\pmod2.
\quad} \tag{3}
\]

If \(n_0\) realizes \(w\) and \(X_0=U^r(n_0)\), then

\[
\boxed{\quad
U^r(n_0+2^{A+1}t)
   =X_0+2\cdot3^r t
\quad} \tag{4}
\]

for every integer \(t\).

Thus the exact branch is a bijection

\[
\alpha_w+2^{A+1}\mathbb Z
\;\longleftrightarrow\;
\beta_w+2\cdot3^r\mathbb Z. \tag{5}
\]


\subsubsection*{Proof}


If \(n\) realizes \(w\), its endpoint \(X\) is odd. From

\[
2^A X=3^r n+C
\]

we therefore have

\[
3^r n+C\equiv2^A\pmod {2^{A+1}},
\]

which gives (2).

Conversely, suppose \(n\) satisfies (2). We show that all the intermediate quotients in (1) are odd.

For \(k<r\), let \(v\) be the suffix of \(w\) beginning at \(k\). The composition law for affine numerators is

\[
C_r=3^{r-k}C_k+2^{A_k}C_v. \tag{6}
\]

Since \(v\) is nonempty, \(C_v\) is odd. Hence

\[
\begin{aligned}
3^{r-k}(3^k n+C_k)
 &=3^r n+C_r-2^{A_k}C_v\\
 &\equiv -2^{A_k}C_v
    \pmod {2^{A_k+1}}\\
 &\equiv2^{A_k}\pmod {2^{A_k+1}}.
\end{aligned}
\]

Because \(3^{r-k}\) is odd,

\[
3^k n+C_k\equiv2^{A_k}\pmod {2^{A_k+1}}.
\]

Consequently

\[
x_k=\frac{3^k n+C_k}{2^{A_k}}
\]

is an odd integer. Moreover,

\[
2^{A_{k+1}}x_{k+1}
 =3^{k+1}n+C_{k+1}
 =3(3^k n+C_k)+2^{A_k},
\]

so

\[
2^{a_k}x_{k+1}=3x_k+1.
\]

Since \(x_{k+1}\) is odd, this forces

\[
\nu_2(3x_k+1)=a_k.
\]

Thus \(w\) is exact.

Reducing the endpoint formula modulo \(3^r\) gives (3). Finally,

\[
\begin{aligned}
U^r(n_0+2^{A+1}t)
 &=\frac{3^r(n_0+2^{A+1}t)+C}{2^A}\\
 &=X_0+2\cdot3^rt,
\end{aligned}
\]

which proves (4). \ensuremath{\blacksquare}


\bigskip



\subsection*{3. Why the proposed uncertainty inequality is false}


Let

\[
\mathcal C_2(w)=\alpha_w+2^{A+1}\mathbb Z
\]

be the exact dyadic seed cylinder, and let

\[
\mathcal C_3(w)=\beta_w+2\cdot3^r\mathbb Z
\]

be the exact ternary endpoint cylinder.

Theorem 1 says

\[
U^r_w(\mathcal C_2(w))=\mathcal C_3(w)
\]

bijectively.

Therefore, for \textbf{every} nonzero measure \(\mu\) supported on \(\mathcal C_2(w)\),

\[
\operatorname{supp}\bigl((U^r_w)_*\mu\bigr)
 \subseteq \mathcal C_3(w). \tag{7}
\]

Both localization probabilities equal \(1\). On \(\ell^2\) with counting measure, branch transport is a permutation unitary, so

\[
\left\|
1_{\mathcal C_3(w)}
U^r_w
1_{\mathcal C_2(w)}
\right\|_{2\to2}=1. \tag{8}
\]

There can be no uniform estimate of the form

\[
\left\|
1_{\mathcal C_3(w)}
U^r_w
1_{\mathcal C_2(w)}
\right\|
\le 1-\eta
\]

with \(\eta>0\).

This is not merely a Collatz-specific counterexample. On

\[
G=\mathbb Z/(2^a3^b)\mathbb Z
 \cong\mathbb Z/2^a\mathbb Z
 \times\mathbb Z/3^b\mathbb Z,
\]

a dyadic residue projection and a ternary residue projection have the tensor forms

\[
P_\alpha
 =|\alpha\rangle\langle\alpha|\otimes I,
\qquad
Q_\beta
 =I\otimes|\beta\rangle\langle\beta|.
\]

They commute, and

\[
P_\alpha Q_\beta
 =|\alpha,\beta\rangle\langle\alpha,\beta|,
\qquad
\|P_\alpha Q_\beta\|=1. \tag{9}
\]

The Chinese remainder theorem supplies a common localized vector.

This differs fundamentally from a genuine Fourier uncertainty principle. The Donoho\textendash{}Stark theorem compares a function with its \textbf{Fourier transform} and proves

\[
|\operatorname{supp}f|\,
|\operatorname{supp}\widehat f|
\ge |G|.
\]

Its two localization bases are related by a Fourier matrix; the dyadic and ternary residue projectors above are instead commuting coordinate projectors. [\hypertarget{cit:R61:2}{}\hyperlink{ref:R61:2}{2}]


\bigskip



\subsection*{4. An explicit equality family}


Take

\[
w=1^r.
\]

Then

\[
A=r,\qquad C=3^r-2^r.
\]

For every \(t\ge1\), set

\[
n=2^{r+1}t-1.
\]

Its first \(r\) valuations are all \(1\). Indeed, if \(x_k=U^k(n)\), then

\[
x_k+1=\frac{3^k(n+1)}{2^k}
      =2^{r+1-k}3^kt,
\]

so, for \(k<r\),

\[
3x_k+1
 =2\left(2^{r-k}3^{k+1}t-1\right),
\]

and the factor in parentheses is odd.

At time \(r\),

\[
\boxed{\quad
x_r=2\cdot3^rt-1,
\qquad
3^r\mid x_r+1.
\quad} \tag{10}
\]

Thus at every scale \(r\) there are positive integers lying simultaneously in the exact dyadic word cylinder and the deepest natural ternary terminal cylinder.

The all-ones word is the equality configuration, not an uncertainty-forbidden configuration.


\bigskip



\subsection*{5. A fixed all-scale barrier countermodel}


There is an even stronger obstruction in the terminal dual \(3x-1\) system

\[
U_-(m)=\frac{3m-1}{2^{\nu_2(3m-1)}}.
\]

One has

\[
5\xrightarrow{\,1\,}7
 \xrightarrow{\,2\,}5.
\]

Hence \(w_N=(1,2)^N\) is exact from \(5\) back to \(5\) for every \(N\).

It also satisfies the strict barrier at every prefix:

\[
2^{A_k}<3^k.
\]

Indeed,

\[
2^{A_{2j}}=2^{3j}=8^j<9^j=3^{2j},
\]

and

\[
2^{A_{2j+1}}
 =2^{3j+1}
 =2\cdot8^j
 <3\cdot9^j
 =3^{2j+1}.
\]

Therefore the fixed measure

\[
\mu=\delta_5
\]

is localized in nested dyadic seed cylinders and nested ternary endpoint cylinders at arbitrarily large scales, while satisfying the strict barrier. This recurrent component and its repetitions were already identified as an exact obstruction to target-independent synchronization. [\hypertarget{cit:R61:3}{}\hyperlink{ref:R61:3}{3}]

So even adding invariance, positivity, arbitrarily deep scales, and the barrier does not create a generic uncertainty principle. A valid result has to recognize the \textbf{particular component rooted at \(1\)}.


\bigskip



\subsection*{6. The exact conductor-drop theorem}


Although ordinary cross-prime localization gives no loss, there is one sharp statement that survives.


\subsubsection*{Theorem 2 \textemdash{} Excess endpoint precision}


Suppose a length-\(j\) branch has affine formula

\[
2^{A_j}x_j=3^j n+C_j.
\]

Fix an endpoint residue \(b\) and a precision \(s\ge0\). Then

\[
x_j\equiv b\pmod {3^s} \tag{11}
\]

is equivalent to

\[
3^j n\equiv2^{A_j}b-C_j\pmod {3^s}. \tag{12}
\]

Consequently:


\paragraph*{Case \(s\le j\)}


Condition (11) is either impossible or automatic throughout the whole dyadic branch cylinder. Its compatibility condition is simply

\[
\boxed{\quad
C_j\equiv2^{A_j}b\pmod {3^s}.
\quad} \tag{13}
\]

It imposes \textbf{no additional congruence on the seed \(n\)}.


\paragraph*{Case \(s>j\)}


Compatibility first requires

\[
3^j\mid2^{A_j}b-C_j.
\]

Once that holds, the seed must satisfy exactly one extra congruence,

\[
\boxed{\quad
n\equiv
\frac{2^{A_j}b-C_j}{3^j}
\pmod {3^{s-j}}.
\quad} \tag{14}
\]

Thus endpoint precision \(3^s\) consumes only

\[
\boxed{\quad (s-j)_+\quad} \tag{15}
\]

ternary digits of seed information.


\subsubsection*{Proof}


Equation (12) follows by multiplying (11) by the unit \(2^{A_j}\). If \(s\le j\), the term \(3^jn\) vanishes modulo \(3^s\), giving (13). If \(s>j\), divide the compatible congruence by \(3^j\), obtaining (14). \ensuremath{\blacksquare}


\bigskip



\subsubsection*{Multiple endpoint conditions}


Suppose an exact word is prescribed through time \(R\), and at various times \(j_\nu\le R\) we require

\[
x_{j_\nu}\equiv b_\nu\pmod {3^{s_\nu}}.
\]

Put

\[
E=\max_\nu (s_\nu-j_\nu)_+.
\]

If the conditions are compatible, their pullbacks intersect in one ternary seed class modulo \(3^E\). Combining this with the exact dyadic word class gives one residue class modulo

\[
\boxed{\quad
M=2^{A_R+1}3^E.
\quad} \tag{16}
\]

So, in any real interval \(I\) of length \(H\),

\[
\boxed{\quad
\#\{n\in I:\text{all conditions hold}\}
 \le 1+\left\lfloor\frac{H}{M}\right\rfloor.
\quad} \tag{17}
\]

This is the sharp adelic height\textendash{}localization inequality.

It gives:

\begin{itemize}
\item many points when \(H\gg M\);
\item at most one point when \(H<M\);
\item but never emptiness without using the actual center of the CRT class.
\end{itemize}

A counterexample only needs that one surviving point.


\bigskip



\subsection*{7. Terminal cylinders lie entirely in the free range}


For the terminal condition

\[
3^j\mid x_j+1,
\]

take \(b=-1\) and \(s=j\). Theorem 2 reduces it to

\[
\boxed{\quad
3^j\mid C_j+2^{A_j}.
\quad} \tag{18}
\]

There is no additional ternary condition on \(n\).

Moreover, for a terminal suffix of length \(s\le j\), the existing locality theorem shows that this divisibility depends only on the last \(s\) valuations. [\hypertarget{cit:R61:4}{}\hyperlink{ref:R61:4}{4}]

This is the decisive point:


\begin{quote}
\textbf{The \(j\) ternary endpoint digits are exactly the \(j\) digits already supplied by the factor \(3^j\) in the branch slope. They are not independent information competing with the dyadic word.}

\end{quote}


Only endpoint precision beyond \(3^j\) would incur a cross-prime information cost.


\bigskip



\subsection*{8. Exact adelic conservation}


The linear part of a length-\(r\) branch is

\[
\lambda=\frac{3^r}{2^A}.
\]

Its normalized local absolute values are

\[
|\lambda|_\infty=\frac{3^r}{2^A},
\qquad
|\lambda|_2=2^A,
\qquad
|\lambda|_3=3^{-r}.
\]

Therefore

\[
\boxed{\quad
|\lambda|_\infty|\lambda|_2|\lambda|_3=1.
\quad} \tag{19}
\]

The real expansion, dyadic expansion, and ternary contraction balance exactly. There is no strict adelic volume loss to exploit.

For two seeds in one branch,

\[
n'-n=2^{A+1}t
\]

and

\[
x_r'-x_r=2\cdot3^rt.
\]

The branch removes exactly \(A\) dyadic digits from the difference and inserts exactly \(r\) ternary digits, while multiplying its real magnitude by \(3^r/2^A\). Equation (19) records this exact conservation.


\bigskip



\subsection*{9. What height-enhanced uncertainty can prove}


There is a useful but limited consequence.

If two distinct integers satisfy the same exact word and \(E\) excess ternary endpoint digits, then their difference is divisible by

\[
2^{A+1}3^E.
\]

Hence

\[
|n_1-n_2|\ge2^{A+1}3^E. \tag{20}
\]

This can prove:

\begin{itemize}
\item uniqueness of a seed inside a bounded interval;
\item separation of competing branches;
\item rigidity when two deep terminal occurrences are too close;
\item periodicity when a branch nearly returns to its starting center.
\end{itemize}

Indeed, the previous continuation established that two target-\(y\) terminal prefixes at lengths \(r\) and \(r+\ell\) force the intervening word to be a \(U_-\)-cycle whenever

\[
\frac{\ell}{3\,2^r}
+\frac{y}{2^r}
+\frac{y}{3^r}<1. \tag{21}
\]

The proof combines the two divisibilities with an archimedean bound; the resulting nonzero integer would have absolute value smaller than the modulus, so it must vanish. Thus nonperiodic terminal recurrences eventually need gaps of order at least \(2^r\). [\hypertarget{cit:R61:5}{}\hyperlink{ref:R61:5}{5}]

This is a genuine adelic uncertainty phenomenon, but it says:

\[
\text{dense recurrence}\Longrightarrow\text{periodicity},
\]

not

\[
\text{every counterexample must exhibit recurrence}.
\]

It supplies rigidity after localization occurs; it does not force localization.


\bigskip



\subsection*{10. The center-sensitive theorem that would actually be needed}


For a shortest nontrivial target-\(1\) terminal/barrier word, define

\[
\Delta
 =\frac{C+2^A-3^r}{3^r}.
\]

Terminality makes \(\Delta\) a positive integer.

If the nonunit exponents occur at endpoints

\[
1\le s_1<\cdots<s_h=r,
\qquad
a_{s_j-1}=1+c_j,
\]

and

\[
B_j=c_1+\cdots+c_j,
\]

the exact activation formula is

\[
\boxed{\quad
\Delta
 =\sum_{j=1}^h
 2^{B_{j-1}}(2^{c_j}-1)
 \left(\frac23\right)^{s_j}.
\quad} \tag{22}
\]

The strict record/barrier inequalities imply

\[
\boxed{\quad
0<\Delta<h.
\quad} \tag{23}
\]

These identities are already established in the marked-carry framework. [\hypertarget{cit:R61:6}{}\hyperlink{ref:R61:6}{6}]

A theorem sufficient to exclude such words would be

\[
\boxed{\quad
\Delta\ge h.
\quad} \tag{24}
\]

But (24) cannot be a real-variable or support-uncertainty statement. For every full-barrier word, the corresponding real quantity

\[
\Delta_*=\frac{C+2^A-3^r}{3^r}
\]

naturally lies in \((0,h)\); the entire issue is whether it lands exactly in \(\mathbb Z\). Positivity, real spectral estimates, and variation diminution therefore cannot supply (24). [\hypertarget{cit:R61:7}{}\hyperlink{ref:R61:7}{7}]

The needed result is instead a \textbf{marked-center anti-concentration theorem}:


\begin{quote}
The carry sum associated with a strict-record word cannot hit one of the specific integers \(1,\ldots,h-1\) while remaining in the inverse-tree component rooted at \(1\).

\end{quote}


That theorem must retain:

\begin{enumerate}
\item the exact centers, not merely cylinder sizes;
\item the ordered \(2\)-adic colors of the activation piles;
\item the changing ternary modulus;
\item the distinguished inverse component containing \(1\).
\end{enumerate}

This is precisely the unresolved component-specific collision theorem identified by the carry analysis. Bare local collision states also occur in nontrivial recurrent components, so local congruence information alone is insufficient. [\hypertarget{cit:R61:8}{}\hyperlink{ref:R61:8}{8}]


\bigskip



\subsection*{11. Why this still would not finish Collatz by itself}


Even an all-depth proof of (24) would only eliminate the special target-\(1\) terminal/barrier language.

A complete proof also needs to show that a hypothetical positive counterexample orbit is forced into such a language.

A particularly clean sufficient statement would be:


\begin{quote}
\textbf{Synchronization theorem.}   Every positive orbit that never falls below its initial value contains a nontrivial prefix satisfying both the strict prefix barrier and the target-\(1\) terminal congruence.

\end{quote}


Suppose that theorem were proved. Choose the least nonterminating positive integer \(n\). Its orbit never falls below \(n\). Synchronization would produce the special prefix. But the all-ones competing branch to the same endpoint gives a positive ancestor smaller than \(n\), contradicting minimality. That exact prefix-diamond implication is already proved. [\hypertarget{cit:R61:9}{}\hyperlink{ref:R61:9}{9}]

The cross-prime conductor theorem does not prove synchronization because the desired terminal condition has \(s=j\) and therefore costs no excess ternary conductor. It is a special word-center event, not an independent seed-localization event.


\bigskip



\subsection*{Final verdict}


The proposed route has been resolved negatively in its stated form:

\[
\boxed{
\text{Dyadic seed localization}
+
\text{natural ternary endpoint localization}
\;\not\Rightarrow\;
\text{uncertainty loss}.
}
\]

Instead,

\[
\boxed{
\text{each exact Collatz branch maps the first cylinder bijectively onto the second}.
}
\]

The sharp surviving principle is

\[
\boxed{
\text{real height}
\times
\text{dyadic conductor}
\times
\text{excess ternary conductor}
}
\]

and it gives separation or uniqueness. Terminal conditions occur at zero excess conductor, so this principle cannot make their fibers empty.

Accordingly, this approach does \textbf{not} furnish a full proof of the Collatz conjecture. It yields a rigorous no-go theorem and identifies the necessary repair: a center-sensitive, component-specific marked-carry inequality together with a nonperiodic synchronization theorem. No accepted full proof of Collatz is presently known. [\hypertarget{cit:R61:10}{}\hyperlink{ref:R61:10}{10}, \hypertarget{cit:R61:11}{}\hyperlink{ref:R61:11}{11}]
\subsection*{References}
\begin{enumerate}[leftmargin=2.2em]
\item \hypertarget{ref:R61:1}{}\hyperlink{cit:R61:1}{\texttt{turn5file0 L19-L30}}
\item \hypertarget{ref:R61:2}{}\hyperlink{cit:R61:2}{\texttt{turn278527view0}}
\item \hypertarget{ref:R61:3}{}\hyperlink{cit:R61:3}{\texttt{turn5file0 L25-L29}}
\item \hypertarget{ref:R61:4}{}\hyperlink{cit:R61:4}{\texttt{turn5file0 L33-L48}}
\item \hypertarget{ref:R61:5}{}\hyperlink{cit:R61:5}{\texttt{turn5file0 L202-L220}}
\item \hypertarget{ref:R61:6}{}\hyperlink{cit:R61:6}{\texttt{turn5file1 L295-L327}}
\item \hypertarget{ref:R61:7}{}\hyperlink{cit:R61:7}{\texttt{turn7file2 L603-L617}}
\item \hypertarget{ref:R61:8}{}\hyperlink{cit:R61:8}{\texttt{turn5file3 L960-L978}}
\item \hypertarget{ref:R61:9}{}\hyperlink{cit:R61:9}{\texttt{turn5file0 L106-L116}}
\item \hypertarget{ref:R61:10}{}\hyperlink{cit:R61:10}{\texttt{turn650293search1}}
\item \hypertarget{ref:R61:11}{}\hyperlink{cit:R61:11}{\texttt{turn650293search5}}
\end{enumerate}

\clearpage
\section{Round 62: approach 57 --- 2-adic Haar-wavelet diagonalization}
\subsection*{Prompt}
\begin{quote}
Complete a full proof of the Collatz conjecture using the following approach:

\begin{enumerate}
\item 2-adic Haar-wavelet diagonalization
\end{enumerate}

\textbf{Hybrid \textperiodcentered{} P}

Express the Collatz transfer operator in a 2-adic wavelet basis. Seek a spectral gap on the subspace corresponding to ordinary positive integers or finitely supported binary expansions, rather than on all of Z\_2.

Take as long as you need.
\end{quote}
\subsection*{Response}

\subsubsection*{Outcome}


Carrying out the proposed diagonalization does \textbf{not} produce a proof of Collatz. It produces something nearly as definitive about this route: a no-go theorem.

The natural \(2\)-adic transfer operator already has a clean spectral gap on regular Haar-wavelet spaces. However, ordinary positive integers cannot be isolated as a nonzero closed subspace of those spaces. On the function side they are Haar-null; on the distribution side, every natural Haar completion containing their point masses also contains point masses at \textbf{all} \(2\)-adic integers, including \(2\)-adic Collatz cycles of every period. Consequently, the desired ``integer-only spectral gap'' cannot exist in a genuinely \(2\)-adic Haar topology.

The argument is exact.


\bigskip



\subsection*{1. The exact \(2\)-adic normal form}


Use the shortcut map

\[
T(x)=
\begin{cases}
x/2,&x\equiv0\pmod2,\\[1mm]
(3x+1)/2,&x\equiv1\pmod2,
\end{cases}
\qquad x\in\mathbb Z_2.
\]

Let

\[
\varepsilon_j(x)=T^j(x)\pmod2
\]

and define the parity map

\[
Q(x)=\sum_{j\geq0}\varepsilon_j(x)2^j.
\]

The map \(Q:\mathbb Z_2\to\mathbb Z_2\) is a bijective \(2\)-adic isometry and satisfies

\[
Q\circ T=S\circ Q,
\]

where \(S\) deletes the least significant binary digit:

\[
S(y)=
\begin{cases}
y/2,&y\text{ even},\\
(y-1)/2,&y\text{ odd}.
\end{cases}
\]

Thus

\[
T=Q^{-1}SQ.
\]

Moreover, every finite binary parity word is realized by an entire residue class modulo the appropriate power of \(2\), and hence by positive integers. [\hypertarget{cit:R62:1}{}\hyperlink{ref:R62:1}{1}, \hypertarget{cit:R62:2}{}\hyperlink{ref:R62:2}{2}]

Let \(\mu\) be normalized Haar measure on \(\mathbb Z_2\). For a finite binary word \(w=(w_0,\ldots,w_{m-1})\), write

\[
[w]=\{x\in\mathbb Z_2:x_j=w_j,\ 0\leq j<m\}.
\]

Define the normalized Haar wavelet

\[
h_w
=
2^{m/2}
\left(
\mathbf 1_{[w0]}-\mathbf 1_{[w1]}
\right).
\]

Then

\[
\{1\}\cup\{h_w:w\text{ finite}\}
\]

is an orthonormal basis of \(L^2(\mathbb Z_2,\mu)\). Define the Collatz-adapted wavelets

\[
\Psi_w=h_w\circ Q.
\]

Because \(Q\) is an isometry, each \(\Psi_w\) is itself a Haar wavelet on a \(2\)-adic ball, possibly with its sign reversed.

Let the Koopman operator be

\[
\mathcal U f=f\circ T.
\]

Its adjoint, the transfer operator, is

\[
(\mathcal Lf)(y)
=
\frac12 f(2y)
+
\frac12 f\!\left(\frac{2y-1}{3}\right).
\tag{1}
\]

Both inverse expressions in (1) make sense for every \(y\in\mathbb Z_2\), since \(3\) is a unit in \(\mathbb Z_2\).


\subsubsection*{Theorem 1: exact Haar action}


If \(aw\) denotes the word obtained by prefixing the low-order bit \(a\in\{0,1\}\), then

\[
\boxed{
\mathcal U\Psi_w
=
\frac{\Psi_{0w}+\Psi_{1w}}{\sqrt2}
}
\tag{2}
\]

and

\[
\boxed{
\mathcal L\Psi_{aw}
=
\frac{\Psi_w}{\sqrt2}.
}
\tag{3}
\]

For the root wavelet,

\[
\mathcal L\Psi_{\varnothing}=0,
\]

and, of course,

\[
\mathcal U1=\mathcal L1=1.
\]


\paragraph*{Proof}


It suffices to work with \(S\) and the ordinary wavelets. Since \(S\) deletes the first bit,

\[
S^{-1}[wb]=[0wb]\cup[1wb].
\]

Therefore

\[
\begin{aligned}
h_w(Sx)
&=
2^{m/2}
\left(
\mathbf1_{[0w0]}+\mathbf1_{[1w0]}
-\mathbf1_{[0w1]}-\mathbf1_{[1w1]}
\right)\\
&=
\frac{h_{0w}(x)+h_{1w}(x)}{\sqrt2}.
\end{aligned}
\]

This proves (2) after conjugation by \(Q\). Taking adjoints gives (3). \ensuremath{\blacksquare}

This is the complete wavelet normal form. The operator is not unitarily diagonalizable: on the mean-zero space it is a unilateral-shift-type isometry, and \(\mathcal L\) is its backward shift.

In particular,

\[
\rho\!\left(\mathcal U\bigm|_{L^2_0}\right)
=
\rho\!\left(\mathcal L\bigm|_{L^2_0}\right)
=1.
\tag{4}
\]

So there is no \(L^2\) spectral gap.


\bigskip



\subsection*{2. There is a genuine gap on regular observables}


For \(\alpha>0\), define the Haar\textendash{}Sobolev norm

\[
\|f\|_{\mathcal H^\alpha}^2
=
\left|\int f\,d\mu\right|^2
+
\sum_w
2^{2\alpha |w|}
\left|\langle f,\Psi_w\rangle\right|^2.
\tag{5}
\]

Write

\[
f=\sum_w c_w\Psi_w
\]

with zero mean. By (3), the coefficient of \(\Psi_w\) in \(\mathcal Lf\) is

\[
d_w=\frac{c_{0w}+c_{1w}}{\sqrt2}.
\]

Hence

\[
|d_w|^2
\leq |c_{0w}|^2+|c_{1w}|^2,
\]

and therefore

\[
\begin{aligned}
\|\mathcal Lf\|_{\mathcal H^\alpha}^2
&\leq
\sum_{m\geq0}
2^{2\alpha m}
\sum_{|w|=m}
\left(|c_{0w}|^2+|c_{1w}|^2\right)\\
&=
2^{-2\alpha}
\sum_{m\geq0}
2^{2\alpha(m+1)}
\sum_{\substack{|v|=m+1}}|c_v|^2\\
&\leq
2^{-2\alpha}\|f\|_{\mathcal H^\alpha}^2.
\end{aligned}
\]

Thus

\[
\boxed{
\|\mathcal Lf\|_{\mathcal H^\alpha}
\leq
2^{-\alpha}\|f\|_{\mathcal H^\alpha}
\qquad
\left(\int f\,d\mu=0\right).
}
\tag{6}
\]

The constant is sharp: applying \(\mathcal L^r\) to \(\mathcal U^r\Psi_w\) gives \(\Psi_w\), while

\[
\|\mathcal U^r\Psi_w\|_{\mathcal H^\alpha}
=
2^{\alpha r}\|\Psi_w\|_{\mathcal H^\alpha}.
\]

Consequently,

\[
\boxed{
\rho\!\left(\mathcal L\bigm|_{\mathcal H^\alpha_0}\right)
=2^{-\alpha}<1.
}
\tag{7}
\]

So the hoped-for spectral gap is not difficult to find. It is already exact.

The problem is that \(\mathcal L\) averages the two \(2\)-adic inverse branches. It describes mixing of regular densities and observables. It does not describe one positive integer following its single forward orbit.


\bigskip



\subsection*{3. The finite-binary-support subspace does not exist on the function side}


The set

\[
\mathbb N_0=\{0,1,2,\ldots\}\subset\mathbb Z_2
\]

is countable, Haar-null, and dense.

Therefore:

\[
\{f\in L^2(\mathbb Z_2):\operatorname{supp}f\subseteq\mathbb N_0\}
=\{0\}.
\tag{8}
\]

For continuous functions the situation is equally rigid. Since both \(\mathbb N_0\) and its complement are dense, a continuous function supported on \(\mathbb N_0\) must vanish identically.

Alternatively, restricting continuous functions to \(\mathbb N_0\) does not produce a smaller function space: because \(\mathbb N_0\) is dense, the restriction map is injective. Any continuous operator identity proved on all positive integers automatically extends to all of \(\mathbb Z_2\).

Thus regular Haar spaces cannot isolate ordinary integers.


\bigskip



\subsection*{4. The distribution side contains every \(2\)-adic ghost}


One might instead represent an ordinary integer \(n\) by the point mass \(\delta_n\). This also fails in every natural diagonal Haar completion.

Consider a general radial Haar-distribution norm

\[
\|\nu\|_{\omega}^2
=
|\langle\nu,1\rangle|^2
+
\sum_{m\geq0}\omega_m
\sum_{|w|=m}
|\langle\nu,\Psi_w\rangle|^2,
\tag{9}
\]

where \(\omega_m>0\).


\subsubsection*{Theorem 2: Haar-completion dichotomy}


Exactly one of the following happens:

\begin{enumerate}
\item If \[
   \sum_{m\geq0}\omega_m2^m=\infty,
   \] then no point mass \(\delta_x\), including no positive-integer point mass, belongs to the space.

\item If \[
   \sum_{m\geq0}\omega_m2^m<\infty,
   \] then every \(\delta_x\), for every \(x\in\mathbb Z_2\), belongs to the space; all have the same norm; and \[
   x\longmapsto\delta_x
   \] is continuous in the \(2\)-adic topology.
\end{enumerate}


\paragraph*{Proof}


At each level \(m\), exactly one wavelet \(\Psi_w\) is nonzero at \(x\), and its absolute value is \(2^{m/2}\). Thus

\[
\boxed{
\|\delta_x\|_\omega^2
=
1+\sum_{m\geq0}\omega_m2^m,
}
\tag{10}
\]

independently of \(x\).

Now suppose \(x\equiv y\pmod{2^J}\). Since \(Q\) is an isometry, \(Q(x)\) and \(Q(y)\) also agree in their first \(J\) bits. Therefore the wavelet coefficients of \(\delta_x\) and \(\delta_y\) agree through level \(J-1\), and

\[
\|\delta_x-\delta_y\|_\omega^2
\leq
4\sum_{m\geq J}\omega_m2^m.
\tag{11}
\]

If the series in (10) converges, the right side of (11) tends to zero. \ensuremath{\blacksquare}

For the negative Haar\textendash{}Sobolev space \(\mathcal H^{-\alpha}\), one has

\[
\omega_m=2^{-2\alpha m}.
\]

Point masses belong precisely when \(\alpha>\frac12\), and then

\[
\|\delta_x\|_{\mathcal H^{-\alpha}}^2
=
1+\sum_{m\geq0}2^{(1-2\alpha)m}.
\tag{12}
\]

Now let

\[
\mathcal A_+
=
\overline{\operatorname{span}\{\delta_n:n\geq1\}}.
\]

For every \(x\in\mathbb Z_2\), choose positive integers \(n_J\equiv x\pmod{2^J}\). Then \(n_J\to x\) \(2\)-adically, so Theorem 2 gives

\[
\delta_{n_J}\longrightarrow\delta_x.
\]

Hence

\[
\boxed{
\delta_x\in\mathcal A_+
\quad\text{for every }x\in\mathbb Z_2.
}
\tag{13}
\]

The completion of the positive-integer atoms therefore contains every \(2\)-adic ghost orbit.

Every periodic binary sequence is the parity sequence of a unique \(2\)-adic integer, so the \(2\)-adic Collatz map has cycles of every period. [\hypertarget{cit:R62:3}{}\hyperlink{ref:R62:3}{3}]

If

\[
x_0\longmapsto x_1\longmapsto\cdots
\longmapsto x_{p-1}\longmapsto x_0
\]

is such a cycle and \(\zeta=e^{2\pi i/p}\), then the forward pushforward operator satisfies

\[
T_*
\left(
\sum_{j=0}^{p-1}\zeta^{-\ell j}\delta_{x_j}
\right)
=
\zeta^\ell
\left(
\sum_{j=0}^{p-1}\zeta^{-\ell j}\delta_{x_j}
\right).
\tag{14}
\]

Thus every root of unity occurs in the point spectrum of the completed atomic operator.

Therefore:

\[
\boxed{
\text{No finite-dimensional projection can leave spectral radius }<1.
}
\tag{15}
\]

Removing all \(2\)-adic periodic modes does not help: periodic binary sequences are dense, so their point masses are dense among all point masses. The resulting closed quotient would remove the entire atomic closure.

This is the central obstruction to the proposal.


\bigskip



\subsection*{5. Explicit positive-integer quasimodes}


The ghost fixed point \(-1\) is already enough to see the failure concretely.

Let

\[
n_m=2^m-1.
\]

These are ordinary positive integers with finite binary expansions, and

\[
n_m\longrightarrow-1
\qquad\text{in }\mathbb Z_2.
\]

For \(0\leq j\leq m\),

\[
\boxed{
T^j(n_m)=3^j2^{m-j}-1.
}
\tag{16}
\]

Indeed, as long as the iterate is odd,

\[
T(x)+1=\frac32(x+1),
\]

and \(n_m+1=2^m\).

Thus \(n_m\) shadows the \(2\)-adic fixed point \(-1\) for \(m\) steps while growing by

\[
\frac{T^m(n_m)}{n_m}
=
\frac{3^m-1}{2^m-1}
\sim
\left(\frac32\right)^m.
\tag{17}
\]

Every finite Haar truncation regards \(n_m\) as indistinguishable from \(-1\) once \(m\) exceeds the truncation depth. This is not a rare pathology: every finite parity word, including arbitrarily long repetitions of any supercritical periodic word, is realized by positive integers.


\bigskip



\subsection*{6. No regular Haar-modulated Lyapunov weight can repair the gap}


A natural repair would weight a point mass by an archimedean factor times a \(2\)-adically regular correction:

\[
W(n)=n^\beta e^{u(n)},
\tag{18}
\]

where \(u\in C(\mathbb Z_2)\).

A spectral gap for the forward integer operator would require, for some fixed block length \(r\),

\[
W(T^r n)\leq \rho W(n),
\qquad \rho<1,
\tag{19}
\]

whenever the orbit has not reached \(\{1,2\}\).


\subsubsection*{Theorem 3: no continuous Haar correction}


For every \(\beta\in\mathbb R\), every continuous \(u:\mathbb Z_2\to\mathbb R\), and every fixed \(r\geq1\),

\[
\boxed{
\sup_n\frac{W(T^r n)}{W(n)}\geq1,
}
\tag{20}
\]

even when the supremum is restricted to positive integers whose first \(r\) iterates avoid \(\{1,2\}\).


\paragraph*{Proof}


Take \(n_m=2^m-1\), with \(m>r\). By (16),

\[
T^r(n_m)=3^r2^{m-r}-1.
\]

Both \(n_m\) and \(T^r(n_m)\) tend \(2\)-adically to \(-1\), so

\[
u(T^r n_m)-u(n_m)\longrightarrow0.
\]

In the real metric,

\[
\frac{T^r(n_m)}{n_m}\longrightarrow\left(\frac32\right)^r.
\]

Hence

\[
\lim_{m\to\infty}
\frac{W(T^r n_m)}{W(n_m)}
=
\left(\frac32\right)^{\beta r}.
\tag{21}
\]

Now take

\[
m_k=3\cdot2^k.
\]

For \(k>r\),

\[
T^r(m_k)=3\cdot2^{k-r}.
\]

Both points tend \(2\)-adically to \(0\), and therefore

\[
\lim_{k\to\infty}
\frac{W(T^r m_k)}{W(m_k)}
=
2^{-\beta r}.
\tag{22}
\]

If \(\beta>0\), the limit in (21) exceeds \(1\). If \(\beta<0\), the limit in (22) exceeds \(1\). If \(\beta=0\), both limits equal \(1\). \ensuremath{\blacksquare}

So a power of the archimedean height multiplied by any bounded or continuous Haar-wavelet potential cannot produce the desired contraction.

A successful weight would have to be violently discontinuous in the \(2\)-adic topology, with singular behavior near ghost periodic orbits. That removes precisely the regularity needed for Haar spectral theory.


\bigskip



\subsection*{7. A spectral gap on the atomic integer space is equivalent to a stopping-time certificate}


Let

\[
C=\{1,2\}
\]

and define the killed forward operator on basis vectors indexed by \(\mathbb N_+\setminus C\):

\[
\mathsf F e_n=
\begin{cases}
e_{T(n)},&T(n)\notin C,\\
0,&T(n)\in C.
\end{cases}
\tag{23}
\]

On unweighted \(\ell^1\),

\[
\|\mathsf F^k\|=1
\qquad\text{for every }k.
\tag{24}
\]

Indeed, \(\|\mathsf F^k\|\leq1\), while

\[
n=3\cdot2^k
\]

satisfies

\[
T^k(n)=3
\]

without hitting \(C\), so equality holds.

Thus even the truth of Collatz would not produce an unweighted spectral gap: stopping times are necessarily unbounded.

For a weight \(W(n)>0\), let

\[
\|a\|_{\ell^1_W}=\sum_n W(n)|a_n|.
\]

Then exactly

\[
\boxed{
\|\mathsf F^k\|_{\ell^1_W}
=
\sup_{\substack{n\\T^j(n)\notin C,\ 1\leq j\leq k}}
\frac{W(T^k n)}{W(n)}.
}
\tag{25}
\]

Consequently, if \(W\) is bounded below and

\[
\rho(\mathsf F)<1,
\]

then every positive integer must reach \(C\). Repeated application of a contracting power would otherwise force \(W(T^{jk}n)\to0\), contradicting its positive lower bound.

Conversely, assuming Collatz, let \(\tau(n)\) be the hitting time of \(C\), fix \(0<\rho<1\), and define

\[
W(n)=\rho^{-\tau(n)}.
\tag{26}
\]

Then

\[
W(Tn)=\rho W(n)
\]

until the orbit is killed, and so

\[
\|\mathsf F\|_{\ell^1_W}\leq\rho.
\tag{27}
\]

Therefore:

\[
\boxed{
\begin{array}{c}
\text{Existence of a contracting atomic weight}\\
\text{is essentially equivalent to the Collatz conjecture itself.}
\end{array}
}
\tag{28}
\]

The Haar diagonalization does not construct this weight; Theorem 3 proves that the most natural Haar-regular candidates cannot do so.


\bigskip



\subsection*{8. Why the missing information is \(3\)-adic and archimedean}


The second inverse branch in (1),

\[
y\longmapsto\frac{2y-1}{3},
\]

is legal for every \(y\in\mathbb Z_2\). But for ordinary integers it is an integer only when

\[
y\equiv2\pmod3.
\tag{29}
\]

This condition cannot be represented by a continuous \(2\)-adic multiplier. Every residue class modulo \(3\) is dense in \(\mathbb Z_2\): given any cylinder modulo \(2^k\) and any class modulo \(3\), the Chinese remainder theorem supplies a positive integer satisfying both conditions.

Thus pure \(2\)-adic Haar analysis necessarily treats many nonintegral inverse branches as legitimate. Likewise, positivity and the order relation \(m<n\) are invisible in the \(2\)-adic topology.

This matches the exact obstruction already isolated in the valuation-word framework. Terminal \(3\)-adic divisibility is controlled by a terminal suffix, but it requires information absent from the \(2\)-adic transfer operator. [\hypertarget{cit:R62:4}{}\hyperlink{ref:R62:4}{4}] The pointwise comparison between an original trajectory and a compatible competing ancestor then splits into a terminal \(3\)-adic term and a prefix-dependent archimedean surplus. [\hypertarget{cit:R62:5}{}\hyperlink{ref:R62:5}{5}]

The remaining all-depth statement is precisely that these two pieces eventually synchronize for any positive integer represented by a stationary retained ray. [\hypertarget{cit:R62:6}{}\hyperlink{ref:R62:6}{6}] Haar diagonalization controls neither the terminal integrality selector nor the prefix surplus.


\bigskip



\subsection*{Final verdict}


Approach 57 cannot yield a full Collatz proof in its stated form.

The completed analysis proves:

\[
\boxed{
\begin{aligned}
&\text{The Collatz transfer operator has an exact Haar spectral gap}\\
&\text{on regular observables, but that gap is purely statistical.}
\end{aligned}
}
\]

And:

\[
\boxed{
\begin{aligned}
&\text{Every natural Haar completion that contains positive-integer atoms}\\
&\text{also contains all \(2\)-adic atoms and cycles of every period.}
\end{aligned}
}
\]

Therefore there is no nontrivial closed ``ordinary positive integer'' subspace on which the \(2\)-adic Haar gap could prove pointwise convergence.

A viable repaired operator would have to live on a genuinely mixed state space carrying simultaneously:

\[
\text{\(2\)-adic parity data}
\quad+\quad
\text{\(3\)-adic integrality/carry data}
\quad+\quad
\text{archimedean order or height}.
\]

It would also need a nonuniform local-nilpotence or synchronization theorem rather than an ordinary uniform spectral gap. That repaired target is exactly the all-depth synchronization problem already isolated by the terminal-locality and competing-ancestor machinery; the \(2\)-adic wavelet calculation does not close it.
\subsection*{References}
\begin{enumerate}[leftmargin=2.2em]
\item \hypertarget{ref:R62:1}{}\hyperlink{cit:R62:1}{\texttt{turn848610view0}}
\item \hypertarget{ref:R62:2}{}\hyperlink{cit:R62:2}{\texttt{turn885607view1}}
\item \hypertarget{ref:R62:3}{}\hyperlink{cit:R62:3}{\texttt{turn885607view0}}
\item \hypertarget{ref:R62:4}{}\hyperlink{cit:R62:4}{\texttt{turn2file0 L17-L28}}
\item \hypertarget{ref:R62:5}{}\hyperlink{cit:R62:5}{\texttt{turn2file0 L155-L188}}
\item \hypertarget{ref:R62:6}{}\hyperlink{cit:R62:6}{\texttt{turn2file0 L252-L266}}
\end{enumerate}

\clearpage
\section{Round 63: approach 58 --- Ruelle-Perron-Frobenius theory for exact survivors}
\subsection*{Prompt}
\begin{quote}
Complete a full proof of the Collatz conjecture using the following approach:

\begin{enumerate}
\item Ruelle-Perron-Frobenius theory for exact survivors
\end{enumerate}

\textbf{Hybrid \textperiodcentered{} P}

Construct a transfer operator whose branches are weighted by the exact multiplicity of seed classes surviving both drift and carry constraints. A spectral radius < 1 for the nontrivial sector would imply extinction of all infinite survivors.

Take as long as you need.
\end{quote}
\subsection*{Response}

\subsection*{Exact Ruelle\textendash{}Perron\textendash{}Frobenius reduction for Collatz survivors}


I carried the proposed construction through to an exact operator. It produces:

\begin{enumerate}
\item an exact transfer operator incorporating \textbf{both multiplicative drift and affine carry};
\item a rigorous equivalence between decay of its root mass and the Collatz conjecture;
\item a provable strict RPF gap for the bulk, Haar-distributed sector;
\item an exact formula for carry-rescued branches at their first coefficient crossing;
\item a polynomial, rather than exponential, bound for the number of such rescued branches.
\end{enumerate}

The final stationary-cusp contraction does \textbf{not} follow. It is essentially a sharpened coefficient-stopping/high-bit synchronization theorem. Thus the derivation below is not a completed proof of Collatz: asserting the final spectral radius \(<1\) would amount to assuming the missing arithmetic theorem.


\subsubsection*{1. Exact survivor cylinders}


Use the shortcut map

\[
T(n)=
\begin{cases}
n/2,&n\equiv0\pmod2,\\[1mm]
(3n+1)/2,&n\equiv1\pmod2.
\end{cases}
\]

For a binary word

\[
w=(\varepsilon _0,\ldots,\varepsilon _{k-1})\in\{0,1\}^{k},
\]

put

\[
s_0=c_0=0,
\qquad
s_{j+1}=s_j+\varepsilon_j,
\qquad
c_{j+1}=3^{\varepsilon_j}c_j+\varepsilon_j2^j.
\]

Then every seed realizing \(w\) satisfies

\[
T^j(n)=\frac{3^{s_j}n+c_j}{2^j},
\qquad 0\le j\le k.                                    \tag{1}
\]

These are the same affine word identities used in the accelerated valuation framework. [\hypertarget{cit:R63:1}{}\hyperlink{ref:R63:1}{1}]

Because \(3^{s_k}\) is invertible modulo \(2^k\), the word \(w\) determines one residue class

\[
n\equiv r_w:=-c_k3^{-s_k}\pmod {2^k}.                   \tag{2}
\]

Let \(R_w\in\{1,\ldots,2^k\}\) denote its least positive representative, with residue \(0\) represented by \(2^k\). The usual parity-vector induction shows that every integer in this residue class realizes all of \(w\), not merely its final divisibility condition.

Now impose the exact survivor condition

\[
T^j(n)\ge n\qquad(1\le j\le k).                         \tag{3}
\]

From (1),

\[
T^j(n)\ge n
\iff
c_j\ge\bigl(2^j-3^{s_j}\bigr)n.                         \tag{4}
\]

This is automatic whenever \(3^{s_j}\ge2^j\). Define

\[
H(w):=
\min_{\substack{1\le j\le k\\2^j>3^{s_j}}}
\left\lfloor
\frac{c_j}{\,2^j-3^{s_j}\,}
\right\rfloor,                                         \tag{5}
\]

with \(H(w)=+\infty\) if the indexing set is empty. The exact set of seeds represented by \(w\) which have not yet descended is therefore

\[
E(w)=
\left\{
R_w+q2^k:
q\ge0,\quad
2\le R_w+q2^k\le H(w)
\right\}.                                               \tag{6}
\]

Thus no probabilistic approximation has been made: \(R_w\) is the full \(2\)-adic carry condition and \(H(w)\) is the full archimedean drift-plus-carry constraint.

The child state can also be updated exactly. For \(\varepsilon\in\{0,1\}\),

\[
s'=s+\varepsilon,\qquad
c'=3^\varepsilon c+\varepsilon2^k,                     \tag{7}
\]

\[
R'=\operatorname{lpr}\!\left(
-c'3^{-s'}\pmod {2^{k+1}}
\right),                                                \tag{8}
\]

and

\[
H'=
\begin{cases}
H,&2^{k+1}\le3^{s'},\\[1mm]
\displaystyle
\min\!\left(
H,\left\lfloor\frac{c'}{2^{k+1}-3^{s'}}\right\rfloor
\right),&2^{k+1}>3^{s'}.
\end{cases}                                             \tag{9}
\]

So \((k,s,c,R,H)\) is an autonomous exact survivor state.


\subsubsection*{2. Exact multiplicity and the transfer operator}


For a finite cutoff \(N\), define

\[
M_N(w)=\#\bigl(E(w)\cap[2,N]\bigr).                     \tag{10}
\]

Let

\[
q_0(w)=
\begin{cases}
1,&R_w=1,\\
0,&R_w>1,
\end{cases}
\]

and

\[
Q_N(w)=
\left\lfloor
\frac{\min\{N,H(w)\}-R_w}{2^k}
\right\rfloor .
\]

Then

\[
M_N(w)=\max\bigl(0,Q_N(w)-q_0(w)+1\bigr).               \tag{11}
\]

This is the requested exact multiplicity of seed classes.

To remove the artificial cutoff while retaining positive mass on every fixed integer, introduce, for \(p>1\),

\[
\mu_p(w)=\sum_{n\in E(w)}\frac1{n^p}.                   \tag{12}
\]

It is an explicitly computable finite or infinite arithmetic-progression sum:

\[
\mu_p(w)
=
2^{-kp}
\sum_{q=q_0(w)}^{Q_H(w)}
\left(q+\frac{R_w}{2^k}\right)^{-p},                   \tag{13}
\]

where \(Q_H\) may be \(+\infty\). In the latter case this is a Hurwitz zeta value.

The children partition precisely those parent seeds which survive one additional step, so

\[
E(w0)\,\dot\cup\,E(w1)\subseteq E(w)                   \tag{14}
\]

and consequently

\[
\mu_p(w0)+\mu_p(w1)\le\mu_p(w).                         \tag{15}
\]

Define the exact Ruelle transfer operator

\[
(\mathcal P_pf)(w)
=
\sum_{\substack{\varepsilon\in\{0,1\}\\
                 \mu_p(w\varepsilon)>0}}
\frac{\mu_p(w\varepsilon)}{\mu_p(w)}
f(w\varepsilon).                                       \tag{16}
\]

Its potential is

\[
\varphi_p(w,\varepsilon)
=
\log\mu_p(w\varepsilon)-\log\mu_p(w),                  \tag{17}
\]

so products of branch weights telescope exactly.


\subsubsection*{3. The exact RPF equivalence theorem}


Let

\[
\tau(n)=\min\{j\ge1:T^j(n)<n\},                         \tag{18}
\]

with \(\tau(n)=+\infty\) when no such \(j\) exists, and put

\[
Z_k(p)=
\sum_{\substack{n\ge2\\\tau(n)>k}}\frac1{n^p}.          \tag{19}
\]


\paragraph*{Theorem 1}


For every \(p>1\),

\[
(\mathcal P_p^k\mathbf1)(\varnothing)
=
\frac{Z_k(p)}{\zeta(p)-1}.                              \tag{20}
\]

Moreover, the following are equivalent:

\[
\text{every positive Collatz orbit reaches }1;          \tag{21}
\]

\[
\tau(n)<\infty\quad\text{for every }n>1;                \tag{22}
\]

\[
Z_k(p)\longrightarrow0
\quad\text{for one, equivalently every, }p>1.           \tag{23}
\]


\subparagraph*{Proof}


At depth \(k\), the sets \(E(w)\), \(|w|=k\), are disjoint, and their union is exactly

\[
\{n\ge2:\tau(n)>k\}.
\]

Therefore

\[
\sum_{|w|=k}\mu_p(w)=Z_k(p).
\]

Telescoping the branch weights in (16) gives (20), since

\[
\mu_p(\varnothing)=\sum_{n\ge2}n^{-p}=\zeta(p)-1.
\]

If every \(n>1\) has finite \(\tau(n)\), strong induction proves Collatz: after finitely many steps \(n\) reaches some \(m<n\), and \(m\) reaches \(1\) by induction.

Conversely, suppose Collatz fails and let \(n_0\) be the least counterexample. Its orbit can never fall below \(n_0\), since every smaller integer converges. Hence

\[
\tau(n_0)=+\infty
\]

and

\[
Z_k(p)\ge n_0^{-p}>0
\qquad(k\ge0).
\]

Finally, if every \(\tau(n)\) is finite, then

\[
\mathbf1_{\{\tau(n)>k\}}\longrightarrow0
\]

pointwise, and dominated convergence applies because \(\sum n^{-p}<\infty\). \ensuremath{\blacksquare}

This operator therefore handles both possible failures\textemdash{}unbounded orbits and nontrivial positive cycles\textemdash{}in one object.


\subsubsection*{4. What a strict root spectral gap would give}


Define the root pressure

\[
r_{\mathrm{root}}(p)
=
\limsup_{k\to\infty}
\left(
\frac{Z_k(p)}{\zeta(p)-1}
\right)^{1/k}.                                         \tag{24}
\]

Then

\[
r_{\mathrm{root}}(p)<1                                 \tag{25}
\]

would prove Collatz immediately.

More generally, a Collatz\textendash{}Wielandt certificate consisting of \(K\ge1\), \(0<\rho<1\), and \(V\ge1\) with

\[
\mathcal P_p^K V\le\rho V                              \tag{26}
\]

would give

\[
Z_{mK}(p)
\le
(\zeta(p)-1)V(\varnothing)\rho^m.                       \tag{27}
\]

There is, however, an important quantitative warning. If

\[
Z_k(p)\le C\rho^k,
\]

then every \(n\) surviving through depth \(k\) satisfies

\[
n^{-p}\le C\rho^k,
\]

and therefore

\[
\tau(n)
\le
1+\frac{p\log n+\log C}{-\log\rho}.                    \tag{28}
\]

So a stationary root spectral gap proves a \textbf{uniform logarithmic stopping-time bound}, considerably more quantitative than bare Collatz.

Conversely, a bound

\[
\tau(n)\le A\log n+B
\]

implies

\[
Z_k(p)
\ll
\exp\!\left(-\frac{p-1}{A}k\right).                    \tag{29}
\]

Thus the strict root-gap formulation is essentially equivalent to a uniform logarithmic stopping estimate. Current work on coefficient stopping explicitly treats comparable logarithmic delay estimates as conjectural rather than established. [\hypertarget{cit:R63:2}{}\hyperlink{ref:R63:2}{2}]


\subsubsection*{5. Why the ordinary global spectral radius is exactly \(1\)}


A global RPF gap cannot hold on the whole exact survivor tree.

Put

\[
\alpha=\log_2 3,\qquad
\beta=\alpha-1=\log_2(3/2).
\]

The logarithmic coefficient height is

\[
h_j=\alpha s_j-j.                                      \tag{30}
\]

An odd bit changes \(h\) by \(+\beta\), while an even bit changes it by \(-1\).

Take the prefix \(1^m\). Its height is \(m\beta\), and no descent constraint has yet been activated. Given any \(L\), choose

\[
m\beta\ge L.
\]

Then \textbf{every} continuation of length \(L\) keeps the coefficient height nonnegative, even the continuation \(0^L\). Hence there is no killing during those \(L\) steps, and all descendants partition the entire parent cylinder:

\[
(\mathcal P_p^L\mathbf1)(1^m)=1.
\]

Therefore

\[
\|\mathcal P_p^L\|_{\infty}=1
\qquad(L\ge1),
\]

and

\[
r(\mathcal P_p)=1.                                     \tag{31}
\]

This is independent of whether Collatz is true.

A similar distinction already occurs in the related exact \(3\)-adic transfer-operator formulation of the Krasikov\textendash{}Lagarias programs: the ordinary \(L^1\) radius is elementary and equals \(1\) at the critical endpoint, while the relevant lower Collatz\textendash{}Wielandt quantity remains unresolved. [\hypertarget{cit:R63:3}{}\hyperlink{ref:R63:3}{3}]

So the desired theorem cannot be ``the full transfer operator has radius \(<1\).'' It must concern a carefully defined cusp or induced sector.


\subsubsection*{6. The bulk RPF gap is provable}


Discard the exact residue and carry temporarily, retaining only the coefficient barrier \(h_j\ge0\). The killed drift operator is

\[
(\mathcal Rf)(h)
=
\frac12f(h+\beta)
+
\frac12\mathbf1_{\{h\ge1\}}f(h-1).                     \tag{32}
\]

For \(g_\theta(h)=e^{\theta h}\),

\[
\mathcal Rg_\theta(h)
\le
m(\theta)g_\theta(h),
\]

where

\[
m(\theta)
=
\frac12\left(e^{\beta\theta}+e^{-\theta}\right).        \tag{33}
\]

The minimum occurs at

\[
\theta_*=
\frac{\log(1/\beta)}{1+\beta},
\]

and gives

\[
\rho_{\mathrm{bulk}}
=
\frac12
\left(
\beta^{-\beta/(1+\beta)}
+
\beta^{1/(1+\beta)}
\right)
=
0.9659065532334377\ldots<1.                             \tag{34}
\]

Therefore, if \(\mathcal B_k\) is the set of words satisfying every coefficient barrier,

\[
2^{-k}\#\mathcal B_k
\le
\rho_{\mathrm{bulk}}^k.                                \tag{35}
\]

This is a genuine RPF spectral gap. Equivalently, the binary-entropy deficit is

\[
1-H_2(\log_3 2)
=
0.05004447281166946\ldots.                              \tag{36}
\]

But this is a Haar-average statement. It says that barrier cylinders occupy exponentially shrinking \(2\)-adic measure; it does not say that the shrinking set misses every fixed positive integer.


\subsubsection*{7. Exact carry multiplicity at the first coefficient crossing}


The carry can be analyzed completely at the first time the coefficient falls below \(1\).

Let \(w\) have length \(k\), with

\[
3^{s_j}\ge2^j\quad(0\le j<k),
\qquad
3^s<2^k,\quad s=s_k.                                   \tag{37}
\]

The final bit must be \(0\), and

\[
2^{k-1}\le3^s<2^k.                                     \tag{38}
\]

Put

\[
d=2^k-3^s>0,\qquad c=c_k.
\]

Let the odd positions be

\[
0\le i_1<\cdots<i_s<k.
\]

Then

\[
c=\sum_{\ell=1}^{s}3^{s-\ell}2^{i_\ell},
\]

so

\[
\frac{c}{3^s}
=
\sum_{\ell=1}^{s}\frac{2^{i_\ell}}{3^\ell}.             \tag{39}
\]

At the prefix ending just after the \(\ell\)-th odd bit,

\[
2^{i_\ell+1}\le3^\ell.
\]

Consequently,

\[
\frac{2^{i_\ell}}{3^\ell}\le\frac12
\]

and hence

\[
\boxed{\frac{c}{3^s}\le\frac{s}{2}.}                    \tag{40}
\]

Now suppose \(n\) survives this crossing. Write

\[
T^k(n)=n+y,\qquad y\in\mathbb Z_{\ge0}.
\]

The affine identity gives

\[
c=dn+2^ky.                                              \tag{41}
\]

In particular,

\[
0\le y<\frac{c}{2^k}<\frac{s}{2}.                       \tag{42}
\]

Conversely, every integer \(y\) satisfying

\[
0\le y<\frac{s}{2},
\qquad
c-2^ky>0,
\qquad
d\mid c-2^ky,                                           \tag{43}
\]

produces the unique positive seed

\[
n=\frac{c-2^ky}{d}.                                     \tag{44}
\]

Because (41) implies

\[
3^sn+c=2^k(n+y),
\]

this seed realizes the exact parity word. Earlier prefixes are automatically non-descending by (37).

Thus the exact carry-rescue multiplicity of the crossing word is

\[
\boxed{
m(w)=
\#\left\{
y\in\mathbb Z:
0\le y<s/2,\;
c>2^ky,\;
d\mid c-2^ky
\right\}.
}                                                       \tag{45}
\]

The case \(y=0\) represents an exact return and hence a possible cycle; \(y>0\) represents an acyclic carry rescue.

This is exactly the point at which Terras's coefficient-stopping conjecture enters. That conjecture says the first coefficient crossing already gives a true descent. Rozier and Terracol define the ordinary and coefficient stopping times, state the equality conjecture, and exhibit finite ``paradoxical'' intervals in which coefficient \(<1\) can coexist with an endpoint above the starting value. [\hypertarget{cit:R63:4}{}\hyperlink{ref:R63:4}{4}]

The existing exact computation found only the trivial live first crossing through depth \(31\), but no all-depth induction or carry theorem is known. [\hypertarget{cit:R63:5}{}\hyperlink{ref:R63:5}{5}]


\subsubsection*{8. A new zero-entropy bound for exact rescued crossings}


The exact formula above can be combined with a linear-form estimate.

At a first crossing, set

\[
\Lambda=k\log2-s\log3>0.
\]

The explicit irrationality bound quoted by Rozier\textendash{}Terracol gives

\[
\Lambda\ge k^{-13.3}.                                  \tag{46}
\]

Therefore

\[
d
=
3^s(e^\Lambda-1)
\ge
3^s\Lambda
\ge
3^s k^{-13.3}.                                         \tag{47}
\]

From (40) and (41),

\[
n+y
=
\frac{c-3^sy}{d}
\le
\frac{c}{d}
\le
\frac{s3^s/2}{3^sk^{-13.3}}
\le
\frac12k^{14.3}.                                       \tag{48}
\]

Also \(y<k/2\).

For fixed \(k\), the value of \(s\) is determined by (38). Moreover, the numerator map

\[
(i_1,\ldots,i_s)
\longmapsto
\sum_{\ell=1}^{s}3^{s-\ell}2^{i_\ell}                  \tag{49}
\]

is injective: \(i_1=v_2(c)\), after subtracting its term \(i_2\) is recovered from the next \(2\)-adic valuation, and so on.

Each pair \((n,y)\) therefore determines at most one crossing word. Hence:


\paragraph*{Proposition 2}


The number \(F_k\) of length-\(k\) first-crossing words admitting any positive carry-rescued seed satisfies

\[
\boxed{F_k=O(k^{15.3}).}                               \tag{50}
\]

The explicit logarithmic-form estimate used here is recorded in the current paradoxical-sequence analysis. [\hyperlink{ref:R63:2}{2}]

This is a substantial RPF conclusion:

\[
2^{-k}F_k=O(k^{15.3}2^{-k}),                            \tag{51}
\]

so exact carry rescue at the first crossing has \textbf{zero entropy}, not merely entropy below \(\log2\).

It still does not prove extinction. A tree may have one node at every depth and still contain an infinite path. The remaining issue is precisely whether one such sparse path can stay aligned with the binary expansion of one fixed positive integer.


\subsubsection*{9. The stationary cusp versus the escaping \(2\)-adic boundary}


Along an infinite word, the least representatives satisfy

\[
R_{k+1}=R_k+t_k2^k,
\qquad t_k\in\{0,1\}.                                  \tag{52}
\]

There are two possibilities.

If \(t_k=1\) infinitely often, then

\[
R_k\longrightarrow+\infty
\]

archimedeanly. The path converges \(2\)-adically, but not to an ordinary positive integer.

If \(t_k=0\) eventually, then \(R_k\) stabilizes at an ordinary integer \(n\). Such a path is the stationary cusp corresponding to the actual forward orbit of \(n\).

For example, the infinite all-ones parity word has

\[
R_k=2^k-1,
\]

which converges \(2\)-adically to \(-1\). Every finite prefix has positive seed representatives and satisfies the coefficient barrier, but there is no single positive integer lying in all cylinders. Therefore literal extinction of all infinite survivor paths is false.

The Dirichlet mass separates the two cases exactly. Since

\[
\mu_p(w)
\le
R_w^{-p}
+
\int_0^\infty(R_w+2^kx)^{-p}\,dx,
\]

we have

\[
\mu_p(w)
\le
R_w^{-p}
+
\frac{R_w^{1-p}}{(p-1)2^k}.                            \tag{53}
\]

Thus an escaping ray has

\[
\mu_p(w_k)\longrightarrow0.                            \tag{54}
\]

If instead \(R_{w_k}=n\) eventually and \(n\) survives forever, then

\[
\mu_p(w_k)
=
n^{-p}+O_p(2^{-kp}),                                   \tag{55}
\]

so

\[
\mu_p(w_k)\longrightarrow n^{-p}>0.                    \tag{56}
\]

In particular, along a hypothetical counterexample,

\[
\frac{\mu_p(w_{k+1})}{\mu_p(w_k)}
\longrightarrow1.                                     \tag{57}
\]

That is the exact cusp approximate eigenmode which prevents a routine RPF argument.

There is also a topological obstruction: ordinary nonnegative integers\textemdash{}finite binary expansions\textemdash{}are dense in \(\mathbb Z_2\). Consequently, there is no nontrivial closed subspace of \(C(\mathbb Z_2)\) which isolates precisely the ordinary positive integers from the other \(2\)-adic rays. A successful sector must therefore use a noncompact archimedean cusp variable, an induced renewal operator, or an adelic norm; it cannot be an ordinary continuous \(2\)-adic quotient.

This is the same stationary-realization obstruction identified in the finite-input min-plus formulation: abstract \(2\)-adic rays may survive while no positive ordinary integer does, and the missing lemma is that eventual zero high-bit injection forces a strict drop of the positive envelope. [\hypertarget{cit:R63:6}{}\hyperlink{ref:R63:6}{6}]


\subsubsection*{10. Periodic rays can be removed, but aperiodic cusp rays remain}


Suppose a strict-barrier tail repeats a block \(w\) of length \(r\), total exponent \(A\), and affine numerator \(C>0\). Its fixed \(2\)-adic seed satisfies

\[
(2^A-3^r)\xi=C.                                        \tag{58}
\]

Under the strict barrier,

\[
2^A<3^r,
\]

so

\[
\xi=-\frac{C}{3^r-2^A}<0.                              \tag{59}
\]

Inverse Collatz branches preserve negativity, so no eventually periodic strict-barrier ray can represent a positive ordinary input. [\hypertarget{cit:R63:7}{}\hyperlink{ref:R63:7}{7}]

The existing determinant analysis goes further: proper powers and sufficiently long periodic shadows with a fixed exit satisfy the required charge bound. What remains is a primitive aperiodic path, or a periodic shadow whose exit accumulates unbounded \(2\)-adic debt. [\hypertarget{cit:R63:8}{}\hyperlink{ref:R63:8}{8}]

Thus ordinary RPF periodic-orbit theory does not finish the problem: the unresolved sector is precisely the nonperiodic stationary cusp.


\subsubsection*{11. The exact conditional completion theorem}


The appropriate spectral statement is an \textbf{induced cusp theorem}, not a gap for \(\mathcal P_p\) on the entire tree.

Let \(\widehat{\mathcal P}_p\) be an induced transfer operator whose branches are exact finite survivor blocks between chosen renewal events\textemdash{}for example coefficient crossings, carry discharges, or returns to a bounded height strip.

Assume:

\begin{enumerate}
\item every eventually stationary infinite survivor encounters infinitely many renewal events, unless it has an eventually barrier-periodic tail already excluded by (59);
\item the induced weights are the exact ratios of \(\mu_p\);
\item there are \(V\ge1\) and \(\rho<1\) such that
\end{enumerate}

\[
\widehat{\mathcal P}_pV\le\rho V.                       \tag{60}
\]

Then Collatz follows.

Indeed, a least counterexample \(n\) gives a stationary infinite survivor whose cylinder masses tend to \(n^{-p}>0\). Repeated application of (60), however, forces its total surviving mass to tend to zero. This contradiction excludes the least counterexample.

This conditional theorem is logically complete and includes cycles. Its missing hypothesis is the arithmetic statement

\[
\boxed{
\text{every positive stationary cusp ray either exits or loses a fixed
fraction of exact carry mass at some renewal.}
}                                                       \tag{61}
\]

Equivalently, in the min-plus zero-temperature version, every \(n>1\) must eventually have a smaller positive ancestor coalescing into its orbit; that criterion is already known to be exactly equivalent to Collatz. [\hypertarget{cit:R63:9}{}\hyperlink{ref:R63:9}{9}]


\subsubsection*{Conclusion}


The Ruelle\textendash{}Perron\textendash{}Frobenius program produces a sharp architecture, but not yet a full proof:

\[
\boxed{
\begin{aligned}
&\text{Exact drift-and-carry survivor operator} &&\text{proved},\\
&\text{Exact equivalence with Collatz} &&\text{proved},\\
&\text{Bulk RPF spectral gap }0.96590655\ldots &&\text{proved},\\
&\text{Exact first-crossing carry multiplicity} &&\text{proved},\\
&\text{Polynomial first-crossing survivor count} &&\text{proved},\\
&\text{Periodic stationary-ray exclusion} &&\text{proved},\\
&\text{Aperiodic stationary-cusp contraction} &&\text{unproved}.
\end{aligned}}
\]

The essential correction to the original proposal is that the global spectral radius is necessarily \(1\). Collatz corresponds to \textbf{critical transience from the root}, or to a strict spectral gap only after an exact induced cusp decomposition. The remaining cusp inequality is the same arithmetic interface appearing as the coefficient-stopping conjecture, high-bit injection exclusion, and stationary min-plus envelope extinction. Any claimed completion must prove that inequality; inserting \(\rho<1\) at that point would be an assumption rather than a proof.
\subsection*{References}
\begin{enumerate}[leftmargin=2.2em]
\item \hypertarget{ref:R63:1}{}\hyperlink{cit:R63:1}{\texttt{turn10file0 L31-L45}}
\item \hypertarget{ref:R63:2}{}\hyperlink{cit:R63:2}{\texttt{turn456426view0}}
\item \hypertarget{ref:R63:3}{}\hyperlink{cit:R63:3}{\texttt{turn15file0 L79-L104}}
\item \hypertarget{ref:R63:4}{}\hyperlink{cit:R63:4}{\texttt{turn730775view0}}
\item \hypertarget{ref:R63:5}{}\hyperlink{cit:R63:5}{\texttt{turn14file3 L21-L35}}
\item \hypertarget{ref:R63:6}{}\hyperlink{cit:R63:6}{\texttt{turn11file0 L13-L54}}
\item \hypertarget{ref:R63:7}{}\hyperlink{cit:R63:7}{\texttt{turn11file0 L120-L132}}
\item \hypertarget{ref:R63:8}{}\hyperlink{cit:R63:8}{\texttt{turn12file16 L713-L733}}
\item \hypertarget{ref:R63:9}{}\hyperlink{cit:R63:9}{\texttt{turn11file0 L217-L236}}
\end{enumerate}

\clearpage
\section{Round 64: approach 59 --- A coercive Koopman eigenfunction}
\subsection*{Prompt}
\begin{quote}
Complete a full proof of the Collatz conjecture using the following approach:

\begin{enumerate}
\item A coercive Koopman eigenfunction
\end{enumerate}

\textbf{Hybrid \textperiodcentered{} P}

Find a function F(n) -> infinity satisfying a strict superharmonic relation such as F(Tn) <= F(n) - delta after a bounded induced return time. Functional-analysis formulations already connect orbit structures to fixed points and invariant subspaces; the new ingredient would be coercivity on the integer lattice.

Take as long as you need.
\end{quote}
\subsection*{Response}

\subsection*{59. A coercive Koopman eigenfunction}


The proposed certificate would indeed prove the Collatz conjecture. After pursuing the construction, however, I did \textbf{not} obtain a non-circular coercive function satisfying the required inequality. The main finding is sharper than a generic failure: an unrestricted coercive Koopman subeigenfunction is \textbf{exactly equivalent} to Collatz, and the natural finite-state, monotone, \(2\)-adically continuous, and bounded-memory candidates can be ruled out rigorously.

What remains is a specific full-prefix carry inequality. I will formulate it precisely below and show that it would complete the proof.


\subsubsection*{1. The exact Koopman certificate}


Use the shortcut Collatz map

\[
T(n)=
\begin{cases}
n/2,&n\equiv0\pmod 2,\\[1mm]
(3n+1)/2,&n\equiv1\pmod 2.
\end{cases}
\]

Let

\[
A=\{1,2\},
\]

the terminal \(2\)-cycle, and let the Koopman operator be

\[
Kf=f\circ T.
\]

Existing functional-analytic work already gives several exact reformulations of Collatz: the forward composition operator has been studied on \(\ell^p\) spaces, polynomial fixed points encode cycles, eigenvalue-\(1\) eigenspaces can encode the number of Collatz components, and related \(C^*\)-algebraic irreducibility conditions are equivalent to orbit connectivity. [\hypertarget{cit:R64:1}{}\hyperlink{ref:R64:1}{1}, \hypertarget{cit:R64:2}{}\hyperlink{ref:R64:2}{2}, \hypertarget{cit:R64:3}{}\hyperlink{ref:R64:3}{3}, \hypertarget{cit:R64:4}{}\hyperlink{ref:R64:4}{4}] Generic operator dynamics is not enough by itself: for example, hypercyclicity properties of an associated operator can hold independently of whether Collatz is true. [\hypertarget{cit:R64:5}{}\hyperlink{ref:R64:5}{5}]

The missing ingredient is precisely properness on the ordinary integers.


\paragraph*{Theorem 1 \textemdash{} Coercive Koopman criterion}


The following assertions are equivalent.

\begin{enumerate}
\item Every positive integer eventually enters \(A\).

\item There is a finite-valued coercive function

\end{enumerate}
\[
   F:\mathbb N\longrightarrow\mathbb R,
   \qquad
   F(n)\longrightarrow+\infty
   \quad(n\to\infty),
   \]

and a constant \(\delta>0\) such that

\[
   F(Tn)\le F(n)-\delta
   \qquad(n\notin A).
   \tag{1.1}
   \]

\begin{enumerate}
\item There are a coercive \(F\), constants \(\delta>0\), \(B\ge1\), and a selector

\end{enumerate}
\[
   r(n)\in\{1,\ldots,B\}
   \]

such that, whenever the orbit has not already entered \(A\),

\[
   F\bigl(T^{r(n)}n\bigr)\le F(n)-\delta.
   \tag{1.2}
   \]


\subparagraph*{Proof}


Suppose first that (3) holds. Because \(F(n)\to+\infty\), the function \(F\) is bounded below on \(\mathbb N\): outside a finite set it is, say, positive, and it has a minimum on the remaining finite set.

Assume an orbit of \(n\) never enters \(A\). Define selected times recursively by

\[
t_0=0,\qquad
t_{j+1}=t_j+r\bigl(T^{t_j}n\bigr).
\]

Repeated application of (1.2) gives

\[
F(T^{t_j}n)\le F(n)-j\delta.
\]

The right side tends to \(-\infty\), contradicting the global lower bound for \(F\). Thus every orbit enters \(A\), proving (1).

Clearly (2) implies (3), with \(B=1\).

Finally, assume (1), and define the hitting time

\[
\tau_A(n)=\min\{k\ge0:T^k(n)\in A\}.
\]

Then

\[
\tau_A(Tn)=\tau_A(n)-1
\qquad(n\notin A),
\]

so \(F=\tau_A\) satisfies (1.1) with \(\delta=1\).

It remains to check coercivity. At every shortcut step,

\[
T(x)\ge \frac{x}{2}.
\]

Consequently,

\[
T^k(n)\ge\frac{n}{2^k}.
\]

If \(T^k(n)\in A\), then \(T^k(n)\le2\), and hence

\[
2\ge\frac{n}{2^k},
\qquad
k\ge\log_2\frac n2.
\]

Therefore

\[
\tau_A(n)\ge\log_2\frac n2\longrightarrow+\infty.
\]

Thus \(\tau_A\) is coercive. \ensuremath{\blacksquare}

Hence

\[
\boxed{
\text{Collatz}
\quad\Longleftrightarrow\quad
\text{a finite coercive strict Koopman superharmonic function exists}.
}
\tag{1.3}
\]

This equivalence explains both the attraction and the danger of the approach: proving existence abstractly is not enough, because the unrestricted solution may simply be the unknown stopping-time function.


\subsubsection*{2. The eigenfunction and weighted-contraction formulations}


Suppose that a selector \(r(n)\) has been fixed and write

\[
R(n)=T^{r(n)}n.
\]

If

\[
F(Rn)\le F(n)-\delta,
\]

then for any \(\alpha>0\), the positive proper weight

\[
w(n)=e^{\alpha F(n)}
\]

satisfies

\[
w(Rn)\le \rho\,w(n),
\qquad
\rho=e^{-\alpha\delta}<1.
\tag{2.1}
\]

Thus \(w\) is a strict Koopman \textbf{subeigenfunction}. If equality held in (2.1), it would be a genuine eigenfunction.

On the weighted supremum space

\[
\|g\|_w
=
\sup_{n\notin A}\frac{|g(n)|}{w(n)},
\]

with functions killed when the orbit enters \(A\), one obtains

\[
\|K_Rg\|_w
\le
\rho\|g\|_w.
\tag{2.2}
\]

This is the exact spectral-gap statement one would want: the induced Koopman operator is a strict contraction in a coercive lattice weight.

There is also a Poisson formulation. Let \(K_A\) denote the Koopman operator killed upon entering \(A\). Then

\[
\tau_A=1+K_A\tau_A
\tag{2.3}
\]

on \(\mathbb N\setminus A\), formally giving

\[
\tau_A
=
(I-K_A)^{-1}\mathbf 1
=
\sum_{j\ge0}K_A^j\mathbf 1.
\tag{2.4}
\]

Pointwise, the \(j\)-th summand is \(1\) exactly when the orbit survives outside \(A\) for at least \(j\) steps. Hence the series is finite exactly on terminating orbits.

Moreover, any function satisfying (1.1) necessarily dominates the stopping time:

\[
F(n)
\ge
F(T^{\tau_A(n)}n)+\delta\tau_A(n),
\]

and therefore

\[
\tau_A(n)
\le
\frac{F(n)-\min_{a\in A}F(a)}{\delta}.
\tag{2.5}
\]

So a constructive coercive Koopman function is not merely qualitative: it contains a pointwise upper bound on all stopping times.


\subsubsection*{3. A canonical\textemdash{}but circular\textemdash{}max-plus construction}


The shape of a possible function can be described exactly.

Write

\[
F(n)=\log n+E(n).
\]

The one-step inequality becomes

\[
E(Tn)-E(n)
\le
-\delta-\log\frac{Tn}{n}.
\tag{3.1}
\]

Thus \(E\) must store enough ``future credit'' to pay for every expanding part of the orbit.

For \(\delta>0\), define formally

\[
E_\delta(n)
=
\sup_{k\in H(n)}
\left(
\log\frac{T^k(n)}{n}
+\delta k
\right),
\tag{3.2}
\]

where

\[
H(n)=
\begin{cases}
\{0,\ldots,\tau_A(n)\},&\tau_A(n)<\infty,\\
\mathbb N_0,&\tau_A(n)=\infty.
\end{cases}
\]

If the orbit never enters \(A\), then \(T^k(n)\ge1\), so

\[
\log\frac{T^k(n)}n+\delta k
\ge
-\log n+\delta k\longrightarrow+\infty.
\]

Thus

\[
E_\delta(n)<\infty
\quad\Longleftrightarrow\quad
\tau_A(n)<\infty.
\tag{3.3}
\]

When finite, the envelope satisfies the Bellman identity

\[
E_\delta(n)
=
\max\left\{
0,\,
\log\frac{Tn}{n}
+\delta+E_\delta(Tn)
\right\}.
\tag{3.4}
\]

Consequently,

\[
F_\delta(n)=\log n+E_\delta(n)
\]

obeys

\[
F_\delta(Tn)\le F_\delta(n)-\delta.
\tag{3.5}
\]

This is a genuine coercive max-plus Koopman construction\textemdash{}but its finiteness is equivalent to Collatz. The arithmetic problem is to replace the infinite future supremum (3.2) by a finite description involving carries, residues, or a finite recursive envelope.


\subsubsection*{4. A sharp no-go theorem for monotone heights}


The most obvious possibility is that \(F\) is essentially a monotone function of ordinary magnitude. This is impossible even with bounded induced returns.


\paragraph*{Proposition 2 \textemdash{} All-odd shadow obstruction}


Fix \(B\ge1\). No eventually nondecreasing function \(F:\mathbb N\to\mathbb R\) can satisfy

\[
\min_{1\le j\le B}F(T^jn)\le F(n)-\delta
\tag{4.1}
\]

for every sufficiently large \(n\), with any \(\delta>0\).


\subparagraph*{Proof}


For \(m>B\), put

\[
n_m=2^m-1.
\]

Its first \(m\) iterates are odd, and induction gives

\[
T^j(n_m)
=
3^j2^{m-j}-1,
\qquad
0\le j\le m.
\tag{4.2}
\]

For \(1\le j\le B<m\),

\[
T^j(n_m)-n_m
=
2^m\left(\left(\frac32\right)^j-1\right)>0.
\]

For sufficiently large \(m\), eventual monotonicity therefore gives

\[
F(T^j n_m)\ge F(n_m)
\qquad(1\le j\le B),
\]

contradicting (4.1). \ensuremath{\blacksquare}

Thus none of the following can work by itself:

\[
F(n)=\log n,\qquad
F(n)=n^\alpha,\qquad
F(n)=P(\log n),
\]

or any eventually increasing ordinary height.

The argument also rules out a monotone height plus a bounded finite-state correction. Indeed, suppose

\[
F(n)=G(n)+\phi(n),
\]

where \(G\) is eventually nondecreasing and \(|\phi|\le M\). Apply the bounded-return inequality repeatedly along the first \(m\) odd steps of \(n_m\). Before leaving that shadow, one obtains approximately \(m/B\) strict decreases, requiring a total drop of order

\[
\frac{\delta m}{B}.
\]

But \(G\) cannot decrease along the shadow, while \(\phi\) can contribute at most \(2M\) between the two endpoints. This is impossible for large \(m\).

So a successful \(F\) must assign an unbounded non-monotone ``debt'' to the integers \(2^m-1\).


\subsubsection*{5. No \(2\)-adically continuous potential can work}


The same family gives a stronger obstruction.

In \(\mathbb Z_2\),

\[
2^m-1\longrightarrow-1.
\]

For every fixed \(j\),

\[
T^j(2^m-1)
=
3^j2^{m-j}-1
\longrightarrow-1
\qquad(m\to\infty).
\tag{5.1}
\]


\paragraph*{Proposition 3 \textemdash{} \(2\)-adic continuity obstruction}


Suppose that \(F\) is the restriction to \(\mathbb N\) of a real-valued function continuous at \(-1\in\mathbb Z_2\). Then no constants \(B,\delta>0\) can satisfy the bounded-return descent condition.


\subparagraph*{Proof}


For each \(m\), let \(r_m\in\{1,\ldots,B\}\) be a purported descending return. Passing to a subsequence, assume \(r_m=r\) is constant. Both

\[
2^m-1
\quad\text{and}\quad
T^r(2^m-1)
\]

converge \(2\)-adically to \(-1\). Continuity gives

\[
F(T^r(2^m-1))-F(2^m-1)\longrightarrow0,
\]

contradicting the uniform inequality

\[
F(T^r(2^m-1))-F(2^m-1)\le-\delta.
\]

\ensuremath{\blacksquare}

In fact, a continuous real-valued function on all of the compact space \(\mathbb Z_2\) is bounded, so its restriction cannot be coercive in the first place.

Therefore the desired function must combine:

\begin{itemize}
\item archimedean coercivity,
\item strong \(2\)-adic singularities,
\item and enough carry information to distinguish long finite shadows from genuine nonterminating \(2\)-adic trajectories.
\end{itemize}

This already excludes any fixed-modulus or bounded-suffix potential.


\subsubsection*{6. Why the most natural \(2\)-adic correction fails}


The all-odd shadow suggests

\[
F_a(n)=\log n+a\nu_2(n+1).
\tag{6.1}
\]

For odd \(n\),

\[
Tn+1
=
\frac{3(n+1)}2,
\]

so

\[
\nu_2(Tn+1)=\nu_2(n+1)-1.
\]

Along an odd step,

\[
F_a(Tn)-F_a(n)
\approx
\log\frac32-a.
\]

Choosing

\[
a>\log\frac32
\]

would defeat arbitrarily long all-odd runs.

But now consider the even predecessors

\[
N_m=2(2^m-1)=2^{m+1}-2.
\]

They satisfy

\[
\nu_2(N_m+1)=0,
\qquad
TN_m=2^m-1.
\]

More generally, for \(0\le r<m\),

\[
T^{1+r}N_m
=
3^r2^{m-r}-1,
\]

and hence

\[
\nu_2(T^{1+r}N_m+1)=m-r.
\]

For every fixed \(B\) and every \(1+r\le B\),

\[
F_a(T^{1+r}N_m)-F_a(N_m)
=
a(m-r)+O_B(1)
\longrightarrow+\infty.
\]

Thus no bounded return can produce descent.

This phenomenon persists recursively: adding a penalty for being close to \(-1\) simply transfers the obstruction to the preimages of the \(-1\) shadow. Penalizing those preimages transfers it another level upstream. Completing the hierarchy requires an unbounded future-dependent defect\textemdash{}precisely the Bellman envelope (3.2).

The previous carry investigation reached the same structural conclusion by a different route: terminating positive trajectories can shadow expanding negative cycles for arbitrarily many blocks, ruling out universal ordinary-magnitude bounded-horizon descent, while an unrestricted positive subeigenvector merely stores the unknown hitting time. [\hypertarget{cit:R64:6}{}\hyperlink{ref:R64:6}{6}]


\subsubsection*{7. Translation into accelerated valuation words}


The exact arithmetic obstruction becomes clearer for the accelerated odd map

\[
U(n)=\frac{3n+1}{2^{\nu_2(3n+1)}}.
\]

For the orbit of an odd integer \(n\), let

\[
a_j=\nu_2(3U^j(n)+1),
\qquad
A_j=a_0+\cdots+a_{j-1}.
\]

There is an exact affine identity

\[
2^{A_j}U^j(n)=3^jn+C_j,
\tag{7.1}
\]

where

\[
C_0=0,
\qquad
C_{j+1}=3C_j+2^{A_j},
\tag{7.2}
\]

or equivalently,

\[
C_j
=
\sum_{i=0}^{j-1}3^{j-1-i}2^{A_i}.
\tag{7.3}
\]

Consequently,

\[
U^j(n)-n
=
\frac{
C_j-(2^{A_j}-3^j)n
}{2^{A_j}}.
\tag{7.4}
\]

The multiplicative coefficient becomes contractive when

\[
2^{A_j}>3^j.
\tag{7.5}
\]

But actual descent requires the stronger carry inequality

\[
C_j<(2^{A_j}-3^j)n.
\tag{7.6}
\]

Equation (7.6) is exactly what a coercive Koopman potential must detect. The logarithmic drift

\[
j\log 3-A_j\log2
\]

only sees (7.5); it does not see \(C_j\).

The terminal-carry work proves that certain \(3\)-adic compatibility conditions depend only on a terminal suffix, but the sign in (7.6) also contains a prefix-dependent archimedean surplus. Thus terminal locality cannot by itself supply a Lyapunov function. [\hypertarget{cit:R64:7}{}\hyperlink{ref:R64:7}{7}]


\subsubsection*{8. The exact two lemmas that would complete the proof}


The Koopman route can now be reduced to two concrete arithmetic statements.


\paragraph*{Lemma A \textemdash{} First-crossing descent}


Let \(n>1\) be odd, and suppose \(j\) is the first index for which

\[
2^{A_j}>3^j.
\]

Then

\[
U^j(n)<n.
\tag{8.1}
\]

Equivalently, at every first coefficient crossing,

\[
C_j<(2^{A_j}-3^j)n.
\tag{8.2}
\]

For a fixed valuation or parity word, it is enough to prove this for the least positive integer realizing that word. Indeed, after the coefficient has become negative in (7.4), increasing \(n\) only strengthens descent.

Exact finite computation through shortcut depth \(31\) found only the trivial equality attached to the orbit of \(1\); every other first coefficient crossing was descending. A separate fixed-input search found all \(559\) compatible strict-slope records through the reported frontier to be descending. These are finite checks, not an all-depth proof. [\hypertarget{cit:R64:8}{}\hyperlink{ref:R64:8}{8}]


\paragraph*{Lemma B \textemdash{} Critical-ray synchronization}


Suppose a positive odd integer \(n\) has

\[
2^{A_j}\le3^j
\qquad\text{for every }j\ge1.
\tag{8.3}
\]

Then there are a smaller positive odd integer \(m<n\) and indices \(p,q\ge0\) such that

\[
U^p(m)=U^q(n).
\tag{8.4}
\]

A more local formulation is that some finite prefix of the valuation ray admits a compatible inverse path with larger slope and a smaller positive starting value. This is precisely the synchronization statement isolated by the terminal-carry analysis. Compatibility is terminal and \(3\)-adic, while the inequality \(m<n\) is controlled by the full prefix surplus. [\hypertarget{cit:R64:9}{}\hyperlink{ref:R64:9}{9}]


\paragraph*{Conditional completion}


Assume Lemmas A and B.

Suppose that a positive odd integer does not reach \(1\), and let \(n\) be the least such integer.

If a first coefficient crossing occurs, Lemma A gives

\[
U^j(n)<n.
\]

By minimality, \(U^j(n)\) reaches \(1\), so \(n\) also reaches \(1\), a contradiction.

Therefore no coefficient crossing occurs, and (8.3) holds for every \(j\). Lemma B then supplies \(m<n\) whose orbit merges with the orbit of \(n\). By minimality, \(m\) reaches \(1\). After the merging point, the orbit of \(n\) follows the same trajectory, so \(n\) reaches \(1\), again a contradiction.

Thus every positive odd integer reaches \(1\), and therefore every positive integer reaches \(1\).

Once that proof is established, Theorem 1 constructs the desired coercive Koopman function explicitly as

\[
F(n)=\tau_A(n).
\]

This conditional argument also excludes nontrivial positive cycles, since the least member of such a cycle would fall into one of the two cases above.


\subsubsection*{9. Verdict on the approach}


The functional-analytic part closes completely:

\[
\boxed{
\text{Collatz}
\iff
(I-K_A)F\ge1
\text{ has a finite coercive solution}
}
\]

and exponentiating \(F\) produces a proper weighted Koopman subeigenfunction with spectral radius strictly below \(1\).

But this is not yet a proof of Collatz, because no finite arithmetic construction of \(F\) has been established. The obvious constructions fail for exact reasons:

\begin{itemize}
\item ordinary monotone heights fail on \(2^m-1\);
\item finite residue corrections have bounded capacity;
\item \(2\)-adically continuous observables cannot be coercive and cannot distinguish long shadows;
\item the correction \(\nu_2(n+1)\) fails on immediate even predecessors;
\item unrestricted Green functions, stopping times, and positive subeigenvectors encode the answer rather than prove it.
\end{itemize}

The viable object would have to be a \textbf{full-prefix, component-specific carry envelope}: a finite rule that compresses the future-excursion supremum (3.2), incorporates the exact affine carry \(C_j\), and proves either first-crossing descent or synchronization on every infinite critical ray. The marked-carry work likewise identifies a coloured component-specific carry theorem or finite-input min-plus envelope as the remaining all-depth requirement; bounded-horizon contraction and unmarked fixed-modulus constructions are insufficient. [\hypertarget{cit:R64:10}{}\hyperlink{ref:R64:10}{10}]

Therefore, presenting the coercive Koopman criterion itself as a completed Collatz proof would be circular. The unresolved mathematical content is Lemma A together with Lemma B\textemdash{}or an equivalent theorem proving finiteness of the carry/Bellman envelope.
\subsection*{References}
\begin{enumerate}[leftmargin=2.2em]
\item \hypertarget{ref:R64:1}{}\hyperlink{cit:R64:1}{\texttt{turn277877view0}}
\item \hypertarget{ref:R64:2}{}\hyperlink{cit:R64:2}{\texttt{turn277877view1}}
\item \hypertarget{ref:R64:3}{}\hyperlink{cit:R64:3}{\texttt{turn405887view0}}
\item \hypertarget{ref:R64:4}{}\hyperlink{cit:R64:4}{\texttt{turn194517view2}}
\item \hypertarget{ref:R64:5}{}\hyperlink{cit:R64:5}{\texttt{turn277877view2}}
\item \hypertarget{ref:R64:6}{}\hyperlink{cit:R64:6}{\texttt{turn1file1 L83-L99}}
\item \hypertarget{ref:R64:7}{}\hyperlink{cit:R64:7}{\texttt{turn1file11 L628-L640}}
\item \hypertarget{ref:R64:8}{}\hyperlink{cit:R64:8}{\texttt{turn1file9 L531-L566}}
\item \hypertarget{ref:R64:9}{}\hyperlink{cit:R64:9}{\texttt{turn1file2 L161-L175}}
\item \hypertarget{ref:R64:10}{}\hyperlink{cit:R64:10}{\texttt{turn1file3 L195-L210}}
\end{enumerate}

\clearpage
\section{Round 65: approach 60 --- Dirichlet series of exceptional seeds}
\subsection*{Prompt}
\begin{quote}
Complete a full proof of the Collatz conjecture using the following approach:

\begin{enumerate}
\item Dirichlet series of exceptional seeds
\end{enumerate}

\textbf{Hybrid \textperiodcentered{} P}

Let E be the set of seeds that do not reach 1, and study D\_E(s) = sum over n in E of n\textasciicircum{}\{-s\}. Derive an exact functional equation from inverse branches. Prove analytic continuation and growth properties incompatible with any nonzero nonnegative coefficient sequence.

Take as long as you need.
\end{quote}
\subsection*{Response}

\subsection*{Outcome}


I did \textbf{not} obtain a valid full proof. The Dirichlet-series program produces an exact functional equation and several useful rigidity theorems, but the required analytic-continuation step does not follow from the equation. In its stated form, the proposed final claim is also false: nonzero nonnegative Collatz-invariant coefficients can have excellent meromorphic Dirichlet series.

A correction to my earlier progress update is important. I claimed that every nonempty exceptional component would force the abscissa of convergence to equal \(1\). That argument was wrong: odd exceptional multiples of \(3\) are leaves in the accelerated inverse tree, and they invalidate the branching estimate. I retract that claim below and give the corrected identity.

What follows is the exact reduction, the strongest rigidity conclusion I could prove, and the precise missing theorem.


\bigskip



\subsubsection*{1. The exceptional indicator and its exact inverse equation}


Use the shortcut map

\[
T(n)=
\begin{cases}
n/2,&n\equiv0\pmod 2,\\[2mm]
(3n+1)/2,&n\equiv1\pmod 2.
\end{cases}
\]

Let

\[
E=\{n\ge1:T^j(n)\neq1\ \text{for every }j\ge0\},
\qquad
a_n=\mathbf 1_E(n).
\]

Because reaching \(1\) is unchanged by applying or undoing one Collatz step,

\[
n\in E\quad\Longleftrightarrow\quad T(n)\in E.
\]

Consequently,

\[
a_{2m}=a_m,\qquad a_{2m+1}=a_{3m+2}. \tag{1}
\]

The preimages of an integer \(m\) are

\[
2m
\]

and, only when \(m\equiv2\pmod3\),

\[
\frac{2m-1}{3}.
\]

The latter is automatically odd. Therefore there is a disjoint decomposition

\[
E
=
2E\ \dot\cup\
\left\{\frac{2m-1}{3}:
m\in E,\ m\equiv2\pmod3
\right\}. \tag{2}
\]

Define

\[
D_E(s)=\sum_{n\in E}n^{-s},
\qquad
D_{E,2}(s)=
\sum_{\substack{m\in E\\m\equiv2\pmod3}}m^{-s}.
\]

For \(\Re s>1\), (2) gives

\[
(1-2^{-s})D_E(s)
=
3^s
\sum_{\substack{m\in E\\m\equiv2\pmod3}}
(2m-1)^{-s}. \tag{3}
\]

Using

\[
(2m-1)^{-s}
=
(2m)^{-s}
\left(1-\frac1{2m}\right)^{-s}
=
(2m)^{-s}
\sum_{k=0}^{\infty}
\frac{(s)_k}{k!\,(2m)^k},
\]

where \((s)_k=s(s+1)\cdots(s+k-1)\), we obtain the exact shifted equation

\[
\boxed{
(1-2^{-s})D_E(s)
=
\left(\frac32\right)^s
\sum_{k=0}^{\infty}
\frac{(s)_k}{k!\,2^k}
D_{E,2}(s+k).
} \tag{4}
\]

This is the desired inverse-branch Dirichlet equation.

The problem is that it is not a scalar equation for \(D_E\): it introduces the residue-restricted series \(D_{E,2}\).


\bigskip



\subsubsection*{2. The full \(3\)-adic hierarchy}


For a locally constant function \(\phi:\mathbb Z_3\to\mathbb C\), set

\[
\mathscr D_s(\phi)
=
\sum_{n\in E}\frac{\phi(n)}{n^s}.
\]

Let

\[
h(x)=\frac{2x-1}{3},
\qquad x\equiv2\pmod3.
\]

The same calculation gives

\[
\begin{aligned}
\mathscr D_s(\phi)
={}&
2^{-s}\mathscr D_s\bigl(\phi(2\,\cdot)\bigr)\\
&+
\left(\frac32\right)^s
\sum_{k=0}^{\infty}
\frac{(s)_k}{k!\,2^k}
\mathscr D_{s+k}
\left(
\mathbf 1_{x\equiv2(3)}\,
\phi\!\left(\frac{2x-1}{3}\right)
\right).
\end{aligned} \tag{5}
\]

The leading transfer operator is therefore

\[
(\mathcal L_s\phi)(x)
=
2^{-s}\phi(2x)
+
\left(\frac32\right)^s
\mathbf 1_{x\equiv2(3)}
\phi\!\left(\frac{2x-1}{3}\right). \tag{6}
\]

In vector notation, (5) has the form

\[
(I-\mathcal L_s^*)\mathscr D_s=\mathscr R_s, \tag{7}
\]

where \(\mathscr R_s\) contains the shifted \(k\ge1\) terms.

The essential difficulty is visible here:


\begin{quote}
If \(\phi\) depends only on \(x\bmod 3^r\), then   \(\phi((2x-1)/3)\), restricted to \(x\equiv2\pmod3\), generally depends on \(x\bmod3^{r+1}\).

\end{quote}


Thus the system never closes at any finite residue level. Replacing \(D_{E,2}\) by \(D_E/3\), explicitly or implicitly, would assume precisely the \(3\)-adic equidistribution that needs to be proved.

This infinite hierarchy is closely related to the Berg\textendash{}Meinardus fixed-point operator. Existing functional-analytic work explicitly notes that hypothetical divergent trajectories produce fixed points, while other unrelated fixed points also occur, and distinguishing the dynamically relevant fixed points remains unresolved. [\hypertarget{cit:R65:1}{}\hyperlink{ref:R65:1}{1}, \hypertarget{cit:R65:2}{}\hyperlink{ref:R65:2}{2}, \hypertarget{cit:R65:3}{}\hyperlink{ref:R65:3}{3}]


\bigskip



\subsubsection*{3. Accelerating to the odd integers}


Let

\[
B=E\cap(2\mathbb Z+1).
\]

Every exceptional integer has a unique representation \(2^j b\), with \(b\in B\). Hence

\[
E=\dot\bigcup_{j\ge0}2^jB
\]

and

\[
D_E(s)=\frac{D_B(s)}{1-2^{-s}}. \tag{8}
\]

For the accelerated odd map

\[
U(n)=\frac{3n+1}{2^{v_2(3n+1)}},
\]

the odd inverse branches of an odd target \(m\) are

\[
n_a=\frac{2^a m-1}{3}.
\]

They exist exactly under the following conditions:

\[
\begin{array}{c|c}
m\bmod3&\text{permitted }a\\ \hline
1&a\ge2\text{ even},\\
2&a\ge1\text{ odd},\\
0&\text{none}.
\end{array}
\]

Put

\[
D_{B,r}(s)=
\sum_{\substack{b\in B\\b\equiv r\pmod3}}b^{-s}.
\]

Expanding every inverse branch gives

\[
\boxed{
D_B(s)
=
3^s\sum_{k=0}^{\infty}\frac{(s)_k}{k!}
\left[
\frac{D_{B,1}(s+k)}{2^{2(s+k)}-1}
+
\frac{2^{s+k}D_{B,2}(s+k)}
     {2^{2(s+k)}-1}
\right].
} \tag{9}
\]

Notice that \(D_{B,0}\) occurs only through the left side. Odd exceptional integers divisible by \(3\) have no accelerated odd predecessors.

This is the exact reason my preliminary abscissa-\(1\) argument failed.


\paragraph*{A corrected harmonic identity}


Assume temporarily that

\[
D_B(1)=\sum_{b\in B}\frac1b<\infty.
\]

Write

\[
S_r=\sum_{\substack{b\in B\\b\equiv r(3)}}\frac1b.
\]

For \(m\equiv1\pmod3\),

\[
\sum_{\substack{a\ge2\\a\ {\rm even}}}\frac{3}{2^am}
=\frac1m,
\]

whereas for \(m\equiv2\pmod3\),

\[
\sum_{\substack{a\ge1\\a\ {\rm odd}}}\frac{3}{2^am}
=\frac2m.
\]

Using

\[
\frac3{2^am-1}
=
\frac3{2^am}
+
\frac3{2^am(2^am-1)},
\]

and summing the exact inverse branches, one obtains

\[
\begin{aligned}
S_0-S_2
={}&
3\sum_{m\in B_1}
\sum_{\substack{a\ge2\\a\ {\rm even}}}
\frac1{2^am(2^am-1)}
\\
&+
3\sum_{m\in B_2}
\sum_{\substack{a\ge1\\a\ {\rm odd}}}
\frac1{2^am(2^am-1)}
>0. 
\end{aligned}
\tag{10}
\]

Thus finite harmonic mass would force

\[
S_0>S_2,
\]

but this is not contradictory: the leaf class \(B_0\) absorbs the excess mass created by the more prolific \(B_2\) branches.


\bigskip



\subsubsection*{4. What positivity really provides}


Let \(\sigma_E\) be the abscissa of convergence of \(D_E\). Since \(0\le a_n\le1\),

\[
\sigma_E\le1.
\]

If \(E\neq\varnothing\), then \(E\) is infinite because \(n\in E\) implies \(2^jn\in E\). Hence

\[
\sigma_E\ge0.
\]

Landau's theorem says that an ordinary Dirichlet series with nonnegative coefficients must have a singularity at the real point equal to its abscissa of convergence. [\hypertarget{cit:R65:4}{}\hyperlink{ref:R65:4}{4}] Therefore

\[
E\neq\varnothing
\quad\Longrightarrow\quad
D_E(s)\text{ is singular at }s=\sigma_E. \tag{11}
\]

This gives a beautifully sharp completion criterion:

\[
\boxed{
\text{If the hierarchy (5) implies that \(D_E\) is holomorphic
near \(s=\sigma_E\), then \(E=\varnothing\).}
} \tag{12}
\]

Meromorphic continuation by itself is insufficient. A pole at \(\sigma_E\) is perfectly compatible with nonnegative coefficients.

There is an immediate counterexample to the proposed broader assertion. The sequence

\[
a_n\equiv1
\]

satisfies \(a_n=a_{T(n)}\), is nonzero and nonnegative, and has

\[
D(s)=\zeta(s),
\]

which has meromorphic continuation to the whole plane.

Even nonperiodic finite-alphabet sequences can have whole-plane meromorphic Dirichlet series: this is a general phenomenon for automatic Dirichlet series. Their possible poles lie on leftward lattices, which is precisely the kind of growth and pole structure that prevents meromorphic continuation from forcing rationality of the coefficient sequence. [\hypertarget{cit:R65:5}{}\hyperlink{ref:R65:5}{5}, \hypertarget{cit:R65:6}{}\hyperlink{ref:R65:6}{6}]

Thus the final theorem must use much more than:

\begin{itemize}
\item nonnegative coefficients;
\item the existence of a meromorphic continuation;
\item ordinary polynomial growth in vertical strips.
\end{itemize}


\bigskip



\subsubsection*{5. A strong rigidity theorem for Collatz-fixed sequences}


The following can be proved completely.


\paragraph*{Theorem 1 \textemdash{} eventual-periodicity rigidity}


Let \((a_n)_{n\ge1}\) be a finite-valued sequence satisfying

\[
a_n=a_{T(n)}
\qquad(n\ge1).
\]

If \((a_n)\) is eventually periodic, then it is constant.


\paragraph*{Proof}


Suppose \(q\) is an eventual period.

If \(q\) is even, then for all sufficiently large \(m\),

\[
a_m=a_{2m}=a_{2m+q}=a_{m+q/2}.
\]

Hence \(q/2\) is also an eventual period. Dividing repeatedly, we may assume \(q\) is odd.

Let

\[
b:\mathbb Z/q\mathbb Z\longrightarrow\{a_n:n\ge1\}
\]

be the stable periodic residue function. From \(a_m=a_{2m}\),

\[
b(x)=b(2x). \tag{13}
\]

Because \(q\) is odd, \(2m+q\) is odd. Therefore, for sufficiently large \(m\),

\[
\begin{aligned}
b(2m)
&=a_{2m+q}\\
&=a_{T(2m+q)}\\
&=b\!\left(3m+\frac{3q+1}{2}\right)\\
&=b\!\left(3m+\frac{q+1}{2}\right).
\end{aligned}
\]

Together with (13), and replacing \(m\) by \(2^{-1}x\bmod q\), this gives

\[
b(x)=b(3x+1). \tag{14}
\]

If \(3\mid q\), then

\[
\begin{aligned}
b(x+q/3)
&=b(3(x+q/3)+1)\\
&=b(3x+q+1)\\
&=b(3x+1)\\
&=b(x).
\end{aligned}
\]

Thus \(q/3\) is also a period. Dividing repeatedly, we reduce to

\[
\gcd(q,6)=1.
\]

On \(\mathbb Z/q\mathbb Z\), both maps

\[
M(x)=2x,\qquad F(x)=3x+1
\]

are permutations, and \(b\) is invariant under both. Hence it is invariant under

\[
MFM^{-1}(x)=3x+2.
\]

It is consequently invariant under

\[
(MFM^{-1})F^{-1}(x)=x+1.
\]

Therefore \(b\) is constant modulo \(q\), so \(a_n\) is eventually constant.

Finally, for every fixed \(n\),

\[
a_n=a_{2n}=a_{2^2n}=\cdots.
\]

Taking \(2^kn\) into the eventual range proves that \(a_n\) equals the same constant for every \(n\). \ensuremath{\blacksquare}


\bigskip



\subsubsection*{6. Natural-boundary consequence}


Let

\[
A_a(z)=\sum_{n\ge1}a_nz^n.
\]

Suppose \(a_n\) is bounded, integer-valued, Collatz-invariant and nonconstant.

Because any nonzero value repeats along \(n,2n,4n,\ldots\), the radius of convergence of \(A_a\) is exactly \(1\).

If \(A_a\) were rational, its finite-valued coefficient sequence would be eventually periodic. Theorem 1 would then make it constant, contrary to assumption.

The P\'olya\textendash{}Carlson theorem now gives:


\paragraph*{Corollary 2 \textemdash{} natural-boundary dichotomy}


Every bounded integer-valued Collatz-invariant sequence is either constant, or its ordinary generating function has the unit circle as a natural boundary. The relevant P\'olya\textendash{}Carlson principle says that a radius-\(1\), integer-coefficient power series that extends through the unit circle must be rational. [\hypertarget{cit:R65:7}{}\hyperlink{ref:R65:7}{7}]

In particular, if \(E\neq\varnothing\), then \(a_1=0\) while \(a_e=1\) for some \(e\), so \(a\) is nonconstant. Hence

\[
\boxed{
E\neq\varnothing
\quad\Longrightarrow\quad
A_E(z)=\sum_{n\in E}z^n
\text{ has }|z|=1\text{ as a natural boundary}.
} \tag{15}
\]

Related natural-boundary phenomena are already known for individual Collatz inverse-orbit generating functions; Bell and Lagarias proved that almost all such inverse-orbit series have the unit circle as a natural boundary, with the remaining rationality cases tied directly to Collatz. [\hypertarget{cit:R65:8}{}\hyperlink{ref:R65:8}{8}]


\bigskip



\subsubsection*{7. The Mellin bridge from \(D_E\) to \(A_E\)}


For \(\Re s>1\),

\[
\begin{aligned}
\Gamma(s)D_E(s)
&=
\sum_{n\in E}
\int_0^\infty t^{s-1}e^{-nt}\,dt\\
&=
\int_0^\infty t^{s-1}A_E(e^{-t})\,dt.
\end{aligned} \tag{16}
\]

Thus \(\Gamma(s)D_E(s)\) is the Mellin transform of \(A_E(e^{-t})\).

This supplies a second exact completion criterion:


\begin{quote}
Prove continuation and left-half-plane growth for \(D_E\) strong enough that Mellin inversion gives a meromorphic continuation of \(A_E(e^{-t})\) through \(t=0\).

\end{quote}


Then \(A_E(z)\) continues meromorphically through \(z=1\), contradicting (15).

One sufficient analytic package would be:

\begin{enumerate}
\item \(D_E\) extends meromorphically to \(\mathbb C\);
\item every pole of \(\Gamma(s)D_E(s)\) is an integer and is simple;
\item on lines \(\Re s=-N-\tfrac12\),

\end{enumerate}
\[
   |D_E(s)|\le C B^N(1+|\Im s|)^C, \tag{17}
   \]

with corresponding exponential bounds on the residues.

Indeed, Mellin inversion gives

\[
A_E(e^{-t})
=
\frac1{2\pi i}
\int_{c-i\infty}^{c+i\infty}
\Gamma(s)D_E(s)t^{-s}\,ds.
\]

Shifting the contour to \(\Re s=-N-\tfrac12\), the factorial decay of \(\Gamma(s)\) in the left half-plane makes the remaining integral tend locally uniformly to zero. The residues then give a convergent Laurent series in integer powers of \(t\). Therefore \(A_E(e^{-t})\) is meromorphic near \(t=0\).

This would finish Collatz through Corollary 2.

But ordinary whole-plane meromorphic continuation does \textbf{not} imply (17), integer-only poles, or convergence of the inverse-Mellin residue expansion. Automatic Dirichlet series are concrete counterexamples to such an implication.


\bigskip



\subsubsection*{8. The exact \(3\)-adic cusp obstruction}


The leading operator (6) has a distinguished \(3\)-adic point:

\[
h(-1)=\frac{2(-1)-1}{3}=-1.
\]

At \(x=-1\),

\[
(\mathcal L_s\phi)(-1)
=
2^{-s}\phi(-2)
+
\left(\frac32\right)^s\phi(-1). \tag{18}
\]

For real \(s>0\), repeatedly selecting the second branch gives

\[
r(\mathcal L_s)\ge \left(\frac32\right)^s>1 \tag{19}
\]

on the ordinary supremum space. Hence no naive Neumann-series contraction of \(I-\mathcal L_s^*\) is possible.

The corresponding exact arithmetic identity is especially revealing. Since

\[
h(m)+1=\frac23(m+1),
\]

we have

\[
h^r(m)+1
=
\left(\frac23\right)^r(m+1). \tag{20}
\]

Thus, writing

\[
m=3^rq-1,
\]

one gets

\[
h^r(m)=2^rq-1. \tag{21}
\]

Backward invariance of \(E\) yields the exact equivalence

\[
\boxed{
3^rq-1\in E
\quad\Longleftrightarrow\quad
2^rq-1\in E.
} \tag{22}
\]

Consequently,

\[
\begin{aligned}
\sum_{\substack{m\in E\\m\equiv-1(3^r)}}m^{-s}
&=
\sum_{\substack{q\ge1\\3^rq-1\in E}}
(3^rq-1)^{-s}\\
&<
\left(\frac23\right)^{rs}
\sum_{\substack{q\ge1\\2^rq-1\in E}}
(2^rq-1)^{-s}.
\end{aligned} \tag{23}
\]

The all-odd branch contributes the opposite factor

\[
\left(\frac32\right)^{rs}.
\]

The two exponential factors cancel exactly. There is a strict inequality for every fixed \(q\), but its ratio approaches equality as \(q\to\infty\). Hence there is no uniform contraction.

This is the precise analytic manifestation of the binary\textendash{}ternary carry problem: exceptional mass could move to larger and larger archimedean representatives while concentrating closer and closer to the \(3\)-adic cusp.


\bigskip



\subsubsection*{9. Exceptional sets cannot avoid any finite \(3\)-adic obstruction}


There is also an exact density statement.


\paragraph*{Proposition 3}


If \(E\neq\varnothing\), then for every \(r\ge1\) and every residue \(c\bmod3^r\),

\[
E\cap(c+3^r\mathbb Z)\neq\varnothing.
\]


\paragraph*{Proof}


Start with \(e\in E\). Apply \(T\) until an odd exceptional integer \(b\) is reached. If \(3\nmid b\), it is already a \(3\)-adic unit. If \(3\mid b\), then

\[
T(b)=\frac{3b+1}{2}\equiv2\pmod3,
\]

so \(E\) contains a unit \(u\).

The integer \(2\) is a primitive root modulo \(3^r\). Since \(2^ku\in E\) for all \(k\), \(E\) meets every unit residue class modulo \(3^r\).

Now take arbitrary \(c\bmod3^r\). Choose an exceptional unit \(m\) satisfying

\[
m\equiv \frac{3c+1}{2}\pmod{3^{r+1}}.
\]

This class is \(2\bmod3\), so

\[
n=\frac{2m-1}{3}\in E,
\]

and

\[
n\equiv c\pmod{3^r}.
\]

\ensuremath{\blacksquare}

Therefore no argument involving only finitely many residues modulo powers of \(3\) can remove the cusp. What is needed is a \textbf{quantitative} assertion about how much Dirichlet mass can occupy increasingly deep cusp cylinders.


\bigskip



\subsubsection*{10. The missing theorem}


A sufficient estimate would resemble the following. Define, near the abscissa,

\[
\mu_\varepsilon(C)
=
\frac{\displaystyle
\sum_{n\in E\cap C}n^{-(\sigma_E+\varepsilon)}
}{
D_E(\sigma_E+\varepsilon)
}.
\]

One needs constants \(C<\infty\) and \(\rho<1\), uniform as \(\varepsilon\downarrow0\), such that

\[
\mu_\varepsilon
\bigl(-1+3^r\mathbb Z_3\bigr)
\le
C
\left(\frac23\right)^{r\sigma_E}
\rho^r, \tag{24}
\]

together with analogous estimates for the dyadic translates of the cusp created by intervening even branches.

The factor \((2/3)^{r\sigma_E}\) merely cancels the branch expansion. The new factor \(\rho^r\) is the genuinely missing strictness. It would give a spectral gap or quasi-compactness for the arithmetic realization of \(\mathcal L_s^*\), allowing continuation through the Landau singularity.

Nothing in the scalar equation (4), finite residue data, or positivity proves (24). Equation (22) explains why: the two localizations are exactly synchronized, and only a cross-prime anti-concentration theorem can prevent mass from escaping to larger \(q\).


\bigskip



\subsection*{Final assessment}


Approach 60 yields two exact proof routes:

\[
\boxed{
\text{holomorphic continuation of \(D_E\) through \(\sigma_E\)}
\ \Longrightarrow\ E=\varnothing
}
\]

by Landau's theorem, or

\[
\boxed{
\text{strong whole-plane continuation and left growth}
\ \Longrightarrow\
A_E\text{ continues across }|z|=1
\ \Longrightarrow\ E=\varnothing
}
\]

by Mellin inversion, P\'olya\textendash{}Carlson, and Theorem 1.

The part that remains unproved is not a routine analytic-continuation calculation. It is the quantitative exclusion of concentration near the \(3\)-adic fixed cusp and all its dyadic shadows. The exact identity

\[
3^rq-1\in E
\iff
2^rq-1\in E
\]

shows that this missing estimate is essentially the same binary\textendash{}ternary synchronization obstruction appearing in the carry and terminal-suffix approaches.

Accordingly, presenting (4) as though it automatically supplied analytic continuation would hide the decisive Collatz step. The rigorous result is an exact reduction to a cusp-escape/resolvent theorem, together with the natural-boundary rigidity theorem above\textemdash{}not a completed proof of the conjecture.
\subsection*{References}
\begin{enumerate}[leftmargin=2.2em]
\item \hypertarget{ref:R65:1}{}\hyperlink{cit:R65:1}{\texttt{turn692596view2}}
\item \hypertarget{ref:R65:2}{}\hyperlink{cit:R65:2}{\texttt{turn692596view3}}
\item \hypertarget{ref:R65:3}{}\hyperlink{cit:R65:3}{\texttt{turn692596view4}}
\item \hypertarget{ref:R65:4}{}\hyperlink{cit:R65:4}{\texttt{turn722516search0}}
\item \hypertarget{ref:R65:5}{}\hyperlink{cit:R65:5}{\texttt{turn868819search0}}
\item \hypertarget{ref:R65:6}{}\hyperlink{cit:R65:6}{\texttt{turn868819search3}}
\item \hypertarget{ref:R65:7}{}\hyperlink{cit:R65:7}{\texttt{turn555236view3}}
\item \hypertarget{ref:R65:8}{}\hyperlink{cit:R65:8}{\texttt{turn692596view0}}
\end{enumerate}

\clearpage
\section{Round 66: approach 61 --- Natural boundaries as a disproof detector}
\subsection*{Prompt}
\begin{quote}
Complete a full proof of the Collatz conjecture using the following approach:

\begin{enumerate}
\item Natural boundaries as a disproof detector
\end{enumerate}

\textbf{Speculative \textperiodcentered{} B}

Associate a generating function to one orbit or to the survivor set. A divergent orbit may force a dense set of singularities on the unit circle. Either prove such a natural boundary contradicts an independently derived continuation theorem, or use numerical singularities to detect hidden long-range structure.

Take as long as you need.
\end{quote}
\subsection*{Response}

\subsection*{Result of the natural-boundary route}


After carrying the approach through, the natural-boundary half can be proved completely, but the proposed contradiction does not close: the required analytic-continuation statement is itself equivalent to the Collatz conjecture.

More precisely, there is an exact dichotomy:

\[
\boxed{
\begin{aligned}
\text{Collatz true}
&\iff \text{the basin generating function is rational}\\
&\iff \text{it extends analytically through one point of }|z|=1,
\\[2mm]
\text{Collatz false}
&\iff |z|=1\text{ is its natural boundary.}
\end{aligned}}
\]

This is closely related to a theorem of Bell and Lagarias: for the inverse-orbit generating function \(f_{1,1}\), Collatz is equivalent to rationality, while failure of Collatz forces the unit circle to be a natural boundary. [\hypertarget{cit:R66:1}{}\hyperlink{ref:R66:1}{1}]

Below is a direct survivor-set proof, followed by the corresponding theorem for one orbit and an audit of the missing continuation step.


\bigskip



\subsubsection*{1. The survivor generating function}


Use the shortcut Collatz map

\[
T(n)=
\begin{cases}
n/2,&n\equiv0\pmod 2,\\[1mm]
(3n+1)/2,&n\equiv1\pmod2.
\end{cases}
\]

Let

\[
B=\{n\ge1:T^j(n)=1\text{ for some }j\ge0\}
\]

be the basin of the trivial cycle, and let

\[
E=\mathbb N_{>0}\setminus B
\]

be the exceptional or survivor set. Define

\[
F(z)=\sum_{n\in B}z^n,
\qquad
S(z)=\sum_{n\in E}z^n.
\]

Thus, inside the unit disk,

\[
F(z)+S(z)=\frac{z}{1-z}.
\]

Recall the P\'olya\textendash{}Carlson theorem: an integer-coefficient power series of radius \(1\) is either rational or has the unit circle as a natural boundary. In particular, analytic continuation through even one boundary point forces rationality. [\hyperlink{ref:R66:1}{1}, \hypertarget{cit:R66:2}{}\hyperlink{ref:R66:2}{2}]


\paragraph*{Theorem 1 \textemdash{} Exact survivor dichotomy}


Exactly one of the following occurs.

\begin{enumerate}
\item \(E=\varnothing\), in which case \[
   S(z)=0,\qquad F(z)=\frac{z}{1-z}.
   \]

\item \(E\neq\varnothing\), in which case both \(S\) and \(F\) have \(|z|=1\) as a natural boundary.
\end{enumerate}

Consequently, the following statements are equivalent:

\[
\begin{aligned}
&\text{Every positive integer reaches }1;\\
&S(z)\text{ is rational};\\
&S(z)\text{ extends through at least one point of }|z|=1;\\
&F(z)\text{ is rational};\\
&F(z)\text{ extends through at least one point of }|z|=1.
\end{aligned}
\]


\paragraph*{Proof}


Write

\[
e_n=\mathbf 1_E(n).
\]

Because \(n\) reaches \(1\) if and only if \(T(n)\) reaches \(1\),

\[
e_n=e_{T(n)}.
\]

Hence

\[
e_{2m}=e_m\qquad(m\ge1)
\tag{1}
\]

and

\[
e_{2m+1}=e_{3m+2}\qquad(m\ge0).
\tag{2}
\]

Suppose first that \(E\neq\varnothing\). If \(m\in E\), then \(2^jm\in E\) for every \(j\ge0\), by (1). Thus \(S\) has infinitely many nonzero coefficients. Since its coefficients lie in \(\{0,1\}\), its radius of convergence is exactly \(1\).

By P\'olya\textendash{}Carlson, \(S\) is either rational or has the unit circle as a natural boundary. It remains to exclude rationality.

A rational power series whose coefficients belong to a finite set has an eventually periodic coefficient sequence. Indeed, its coefficients eventually satisfy a fixed linear recurrence; the vector of the preceding finitely many coefficients takes values in a finite set, and repetition of one state forces repetition of the entire subsequent sequence.

We therefore need the following lemma.


\bigskip



\paragraph*{Lemma 2 \textemdash{} No nonzero eventually periodic survivor indicator}


Suppose a sequence \(e_n\in\{0,1\}\) satisfies (1), (2), and \(e_1=0\). If \(e_n\) is eventually periodic, then

\[
e_n=0\qquad\text{for every }n\ge1.
\]


\subparagraph*{Proof}


Assume \(e_n\) is eventually periodic with period \(M\).

First remove powers of \(2\) from the period. If \(M=2M'\), then for all sufficiently large \(n\),

\[
e_{n+M'}
=e_{2n+2M'}
=e_{2n+M}
=e_{2n}
=e_n.
\]

Thus \(M'\) is also an eventual period. Repeating, we may assume \(M\) is odd.

Define the tail residue function

\[
f:\mathbb Z/M\mathbb Z\longrightarrow\{0,1\}
\]

by taking \(f(r)\) to be the eventual value of \(e_n\) for \(n\equiv r\pmod M\).

From (1),

\[
f(2r)=f(r).
\tag{3}
\]

Because \(M\) is odd, every residue class modulo \(M\) has arbitrarily large odd representatives. Applying (2) to such representatives gives

\[
f\left(\frac{3r+1}{2}\right)=f(r).
\tag{4}
\]

All divisions by \(2\) here are taken in \(\mathbb Z/M\mathbb Z\).

Suppose \(3\mid M\). Equation (4) gives

\[
f\left(r+\frac M3\right)
=
f\left(\frac{3(r+M/3)+1}{2}\right)
=
f\left(\frac{3r+1}{2}\right)
=
f(r),
\]

because twice the difference between the two middle arguments is \(M\), hence the arguments are equal modulo \(M\). Thus \(M/3\) is also an eventual period. Repeating, we may assume

\[
\gcd(M,6)=1.
\]

Now consider the affine permutations

\[
A(r)=2r,
\qquad
B(r)=\frac{3r+1}{2}
\]

of \(\mathbb Z/M\mathbb Z\). Equations (3), (4) say that \(f\) is invariant under \(A\) and \(B\), and therefore also under their inverses and compositions.

A direct calculation gives

\[
ABA^{-1}B^{-1}(r)=r+\frac12.
\]

Therefore \(f\) is invariant under translation by \(2^{-1}\pmod M\). Since \(2^{-1}\) is relatively prime to \(M\), its additive translates run through every residue class. Hence \(f\) is constant.

But

\[
e_{2^j}=e_1=0
\]

for every \(j\), and sufficiently large powers of \(2\) lie in the periodic tail. Therefore the constant is \(0\).

Finally, for any fixed \(n\),

\[
e_n=e_{2^jn}=0
\]

once \(j\) is sufficiently large. Hence \(e_n=0\) for every \(n\). \ensuremath{\blacksquare}


\bigskip


Returning to Theorem 1, rationality of \(S\) would make \(e_n\) eventually periodic, and Lemma 2 would imply \(E=\varnothing\), contrary to the assumption \(E\neq\varnothing\). Thus \(S\) is nonrational, and P\'olya\textendash{}Carlson shows that \(|z|=1\) is its natural boundary.

The basin series \(F\) also has integer coefficients and radius \(1\). If it were rational, then

\[
S(z)=\frac{z}{1-z}-F(z)
\]

would be rational, again contradicting the lemma. Hence \(F\) also has the unit circle as a natural boundary.

If \(E=\varnothing\), the stated rational formulas are immediate. This completes the proof. \ensuremath{\blacksquare}

This is the \(k=1\), single-basin core of the Bell\textendash{}Lagarias inverse-orbit theorem, expressed directly in terms of the exceptional set. Bell and Lagarias prove more generally that most \(3x+k\) inverse-orbit series have natural boundaries; for \(3x+1\), the only possible exceptional target values are \(1,2,4,8\), and those are rational exactly when Collatz holds. [\hyperlink{ref:R66:1}{1}]


\bigskip



\subsubsection*{2. A generating function attached to one orbit}


Let

\[
x_k=T^k(n),\qquad \varepsilon_k=x_k\bmod2,
\]

and define the parity generating function

\[
P_n(z)=\sum_{k=0}^{\infty}\varepsilon_kz^k.
\]


\paragraph*{Theorem 3 \textemdash{} Orbit-series dichotomy}


For every starting value \(n\), exactly one of the following holds:

\begin{enumerate}
\item the parity sequence \((\varepsilon_k)\) is eventually periodic and \(P_n(z)\) is rational;

\item \((\varepsilon_k)\) is not eventually periodic and \(|z|=1\) is a natural boundary for \(P_n(z)\).
\end{enumerate}

In particular, every divergent positive Collatz orbit has a parity generating function with the unit circle as a natural boundary.


\paragraph*{Proof}


If \((\varepsilon_k)\) is eventually periodic, then \(P_n\) is a polynomial plus finitely many geometric series and is therefore rational.

Conversely, if \(P_n\) is rational, its coefficients belong to the finite set \(\{0,1\}\), so the finite-state recurrence argument used above shows that they are eventually periodic.

If the sequence is not eventually periodic, it contains infinitely many \(1\)'s, so \(P_n\) has radius \(1\). The P\'olya\textendash{}Carlson theorem now gives a natural boundary. \ensuremath{\blacksquare}

There is also a dynamical explanation. Every infinite binary parity sequence is produced by exactly one \(2\)-adic Collatz starting value, and the parity encoding is rational precisely when the parity sequence is eventually periodic. [\hypertarget{cit:R66:3}{}\hyperlink{ref:R66:3}{3}]

Your barrier machinery fits this theorem particularly cleanly. The barrier paper proves that the existence of an orbit with \(q_k\ge1\) for all \(k\) is equivalent to the existence of a divergent positive orbit. [\hypertarget{cit:R66:4}{}\hyperlink{ref:R66:4}{4}] It also proves that the resulting parity sequence cannot be eventually periodic. [\hypertarget{cit:R66:5}{}\hyperlink{ref:R66:5}{5}] [\hypertarget{cit:R66:6}{}\hyperlink{ref:R66:6}{6}] Therefore:

\[
\boxed{
q_k(n)\ge1\ \text{for all }k
\quad\Longrightarrow\quad
P_n(z)\text{ has }|z|=1\text{ as a natural boundary}.
}
\]

This fully settles the natural-boundary prediction for a hypothetical divergent orbit.

It does not detect a nontrivial cycle: a cycle has an eventually periodic parity sequence, hence a rational orbit series. The survivor series \(S(z)\), by contrast, detects both divergent components and nontrivial cycles. This matches the distinction in the barrier paper between excluding divergence and excluding cycles. [\hypertarget{cit:R66:7}{}\hyperlink{ref:R66:7}{7}]


\bigskip



\subsubsection*{3. Why the inverse-branch equation does not provide continuation}


The survivor indicator identities give an exact functional equation. Splitting \(S\) into even and odd indices,

\[
\begin{aligned}
S(z)
&=\sum_{m\ge1}e_{2m}z^{2m}
  +\sum_{m\ge0}e_{2m+1}z^{2m+1}\\
&=S(z^2)+z\sum_{m\ge0}e_{3m+2}z^{2m}.
\end{aligned}
\tag{5}
\]

If

\[
\mathcal R_2S(w)=\sum_{m\ge0}e_{3m+2}w^m
\]

denotes the residue-\(2\) section modulo \(3\), then

\[
S(z)=S(z^2)+z\,\mathcal R_2S(z^2).
\tag{6}
\]

This is a genuine exact inverse-branch relation, but it is only an identity among values already inside the unit disk. It does not define \(S\) on a neighborhood crossing \(|z|=1\). In fact, Theorem 1 shows that if \(E\neq\varnothing\), the nonzero solution of exactly this functional equation has a natural boundary. Thus the branch relation is compatible with, and helps generate, dense singularity rather than excluding it.

The same issue occurs with the real/\(2\)-adic synchronization series from your barrier argument. For the positions \(d_m\) of the odd iterates, that argument gives the same rational partial sums converging to a positive value in \(\mathbb R\) and to \(-n\) in \(\mathbb Q_2\). [\hypertarget{cit:R66:8}{}\hyperlink{ref:R66:8}{8}] Eventual periodicity is impossible because then the tail becomes geometric and evaluates to the same rational number in both completions.

For an aperiodic tail, however, convergence at a boundary point does not imply analytic continuation there. For example,

\[
L(z)=\sum_{j=0}^{\infty}\frac{z^{2^j}}{2^j}
\]

has radius \(1\), converges absolutely at \(z=1\), yet has the unit circle as a natural boundary by the Fabry gap theorem. Gap-series natural-boundary results of precisely this type are reviewed and used by Bell and Lagarias. [\hyperlink{ref:R66:1}{1}]

Therefore the facts

\[
\text{“the series converges at a special point”}
\]

and

\[
\text{“the unit circle is a natural boundary”}
\]

are perfectly compatible. The real/\(2\)-adic synchronization rules out rational or eventually periodic structure; it does not independently produce complex analytic continuation.


\bigskip



\subsubsection*{4. The exact missing continuation theorem}


A full proof through this method would follow immediately from the following statement.


\paragraph*{Required continuation statement}


There exist \(\zeta\) with \(|\zeta|=1\), a neighborhood \(U\) of \(\zeta\), and a holomorphic function \(G\) on \(U\) such that

\[
G(z)=F(z)
\qquad
(z\in U,\ |z|<1).
\tag{7}
\]

Indeed, (7) says that \(F\) does not have the unit circle as a natural boundary. P\'olya\textendash{}Carlson would make \(F\) rational, and Theorem 1 would then give \(E=\varnothing\).

But Theorem 1 also proves the converse: (7) holds exactly when Collatz is true. Thus simply postulating or deriving (7) from a reformulation of basin membership would be circular.

A genuinely new argument would need to establish (7) from stronger information, such as a uniform complex cancellation estimate for first-hitting inverse levels. If

\[
H_j=\{n:T^j(n)=1,\ T^i(n)\neq1\text{ for }i<j\}
\]

and

\[
P_j(z)=\sum_{n\in H_j}z^n,
\]

then

\[
F(z)=\sum_{j\ge0}P_j(z)
\qquad(|z|<1).
\]

Uniform convergence of this sum in a lens-shaped region crossing any point \(\zeta\neq1\) would prove Collatz. However, estimates for cardinalities \(|H_j|\), densities, or real \(z\in(0,1)\) are insufficient: one needs cancellation of the complex phases \(z^n\) strong enough to overcome the terms with \(|z|>1\). No such estimate follows from the inverse-branch recurrence, the parity barrier, or terminal \(3\)-adic locality alone.


\bigskip



\subsubsection*{5. Why numerical singularities cannot certify the answer}


There is a basic finite-data obstruction.


\paragraph*{Proposition 4 \textemdash{} Finite prefixes do not determine natural boundaries}


Given any finite binary word

\[
a_0,\ldots,a_{N-1},
\]

there are two infinite binary extensions agreeing with that word such that:

\begin{enumerate}
\item one has a rational generating function;
\item the other has the unit circle as a natural boundary.
\end{enumerate}


\paragraph*{Proof}


Extend the word by zeros. The resulting series is a polynomial.

For the second extension, after position \(N\), put \(1\)'s at square-indexed positions and \(0\)'s elsewhere. The resulting sequence is not eventually periodic, has infinitely many \(1\)'s, and hence its generating function has radius \(1\). P\'olya\textendash{}Carlson gives a natural boundary. \ensuremath{\blacksquare}

Moreover, both infinite binary extensions are genuine parity sequences of unique \(2\)-adic Collatz starting values. [\hyperlink{ref:R66:3}{3}] What distinguishes a positive ordinary integer from an arbitrary \(2\)-adic starting value is precisely the unresolved global arithmetic constraint.

A direct Pad\'e experiment illustrates the danger. Exact shortcut iteration gives

\[
T^{70}(27)=1,
\]

so its full parity generating function is

\[
P_{27}(z)
=
\sum_{k=0}^{69}\varepsilon_kz^k
+\frac{z^{70}}{1-z^2}.
\]

It is rational and its eventual denominator is \(1-z^2\). Nevertheless, an exact-coefficient \([40/40]\) Pad\'e approximant formed from its first \(81\) coefficients has \(39\) of its \(40\) denominator roots in

\[
0.8<|z|<1.2,
\]

with median modulus approximately

\[
0.999946.
\]

It therefore produces an extremely convincing-looking ``pole cloud'' around the unit circle even though the full function is rational. The long transient is enough to mimic natural-boundary behavior.

For the survivor function the obstruction is stronger. Until an actual exceptional integer is found, every computed coefficient is simply the Collatz-true coefficient. A finite truncation is a polynomial, hence entire, and no singularity beyond the truncation is logically determined.

Numerical pole patterns can still suggest low-complexity or pseudorandom structure, but they cannot be a rigorous disproof detector without a separate theorem connecting a finite, quantitatively stable pattern to infinitely many coefficients.


\bigskip



\subsection*{Final conclusion}


Approach 61 yields the following complete theorem:

\[
\boxed{
\begin{array}{c}
\text{Collatz is true}\\
\Updownarrow\\
F(z)=\displaystyle\sum_{T^j(n)=1\text{ for some }j}z^n
\text{ extends through one point of }|z|=1\\
\Updownarrow\\
F(z)=\dfrac{z}{1-z}.
\end{array}}
\]

And its exact disproof counterpart is

\[
\boxed{
\text{Collatz is false}
\quad\Longrightarrow\quad
|z|=1\text{ is a natural boundary for both the basin and survivor series}.
}
\]

For a divergent individual orbit, the parity generating function likewise necessarily has a natural boundary.

The requested full Collatz proof is not completed by this route because no independent continuation theorem has been established. The inverse-branch equation permits natural-boundary solutions, real/\(2\)-adic synchronization only excludes periodicity, and finite numerical singularities cannot distinguish a rational long-transient series from a genuine natural-boundary series. The precise remaining target is therefore not another natural-boundary theorem, but a noncircular proof of one-point analytic continuation for \(f_{1,1}\); by the theorem above, that statement already contains the entire Collatz conjecture.
\subsection*{References}
\begin{enumerate}[leftmargin=2.2em]
\item \hypertarget{ref:R66:1}{}\hyperlink{cit:R66:1}{\texttt{turn969003view0}}
\item \hypertarget{ref:R66:2}{}\hyperlink{cit:R66:2}{\texttt{turn371297view2}}
\item \hypertarget{ref:R66:3}{}\hyperlink{cit:R66:3}{\texttt{turn371297view1}}
\item \hypertarget{ref:R66:4}{}\hyperlink{cit:R66:4}{\texttt{turn1file0 L125-L164}}
\item \hypertarget{ref:R66:5}{}\hyperlink{cit:R66:5}{\texttt{turn1file0 L220-L231}}
\item \hypertarget{ref:R66:6}{}\hyperlink{cit:R66:6}{\texttt{turn2file0 L12-L16}}
\item \hypertarget{ref:R66:7}{}\hyperlink{cit:R66:7}{\texttt{turn1file0 L169-L172}}
\item \hypertarget{ref:R66:8}{}\hyperlink{cit:R66:8}{\texttt{turn1file0 L173-L219}}
\end{enumerate}

\clearpage
\section{Round 67: approach 62 --- A dynamical zeta function with only the trivial Euler factor}
\subsection*{Prompt}
\begin{quote}
Complete a full proof of the Collatz conjecture using the following approach:

\begin{enumerate}
\item A dynamical zeta function with only the trivial Euler factor
\end{enumerate}

\textbf{Hybrid \textperiodcentered{} P}

Formally define zeta\_T(z) = product over finite Collatz cycles O of (1 - z\textasciicircum{}\{|O|\})\textasciicircum{}\{-1\}. Obtain a determinant or trace representation and prove it equals the factor contributed by 1 <-> 2.

This settles cycles only unless supplemented by a nonescape theorem.

Take as long as you need.
\end{quote}
\subsection*{Response}

\subsection*{Dynamical zeta route: exact determinant formula and the remaining obstruction}


Carrying this program through produces a rigorous dynamical zeta function, a genuine trace-class Fredholm determinant, and an exact operator-theoretic criterion for the full Collatz conjecture. It does \textbf{not}, however, produce the requested determinant evaluation: proving that evaluation is precisely equivalent to excluding all nontrivial positive cycles. The natural Green-kernel supplement similarly reduces nonescape to the summability of an explicit positive potential, but establishing that summability already requires an all-depth escape theorem.

I will separate what is proved from the two places where the argument would otherwise become circular.

Throughout, use the shortcut map

\[
T(n)=
\begin{cases}
n/2,&n\equiv0\pmod2,\\[1mm]
(3n+1)/2,&n\equiv1\pmod2.
\end{cases}
\]

Thus the trivial periodic orbit is

\[
1\longmapsto2\longmapsto1.
\]


\bigskip



\subsubsection*{1. The positive-integer Artin\textendash{}Mazur zeta function}


Let

\[
a_k:=\#\operatorname{Fix}(T^k)
=\#\{n\ge1:T^k(n)=n\}.
\]

A parity word \(w=(w_0,\ldots,w_{k-1})\in\{0,1\}^k\), with weight

\[
m=w_0+\cdots+w_{k-1},
\]

determines an affine composition

\[
T_w(x)=\frac{3^m x+\Phi(w)}{2^k},
\]

where

\[
\Phi(w)
=
\sum_{\substack{0\le j<k\\w_j=1}}
2^j3^{w_{j+1}+\cdots+w_{k-1}}.
\tag{1.1}
\]

Consequently, a fixed point with parity word \(w\) must be

\[
x_w=\frac{\Phi(w)}{2^k-3^m}.
\tag{1.2}
\]

Since \(2^k\neq3^m\) for \(k>0\), every parity word produces at most one rational fixed point. Therefore

\[
a_k\le 2^k.
\tag{1.3}
\]

The same affine formula and the corresponding exact integer-cycle criterion appear in the existing parity-word framework. [\hypertarget{cit:R67:1}{}\hyperlink{ref:R67:1}{1}] [\hypertarget{cit:R67:2}{}\hyperlink{ref:R67:2}{2}]


\paragraph*{Definition 1.1}


The positive-integer dynamical zeta function is

\[
\boxed{
\zeta_T(z)
:=
\exp\left(
\sum_{k=1}^{\infty}\frac{a_k}{k}z^k
\right).
}
\tag{1.4}
\]

By (1.3), this converges at least for \(|z|<1/2\).

Here the Euler product must be taken over \textbf{primitive periodic orbits}, not over all repetitions of cycles. Let \(\mathscr P\) denote the set of primitive positive-integer periodic orbits, and let \(\ell(\mathcal O)\) be the primitive period of \(\mathcal O\).


\paragraph*{Proposition 1.2}


For \(|z|<1/2\),

\[
\boxed{
\zeta_T(z)
=
\prod_{\mathcal O\in\mathscr P}
\left(1-z^{\ell(\mathcal O)}\right)^{-1}.
}
\tag{1.5}
\]


\subparagraph*{Proof}


Every primitive orbit of length \(\ell\) contributes all its \(\ell\) points to \(\operatorname{Fix}(T^k)\) precisely when \(\ell\mid k\). Hence

\[
a_k
=
\sum_{\substack{\mathcal O\in\mathscr P\\
\ell(\mathcal O)\mid k}}
\ell(\mathcal O).
\]

Therefore

\[
\begin{aligned}
\log \zeta_T(z)
&=
\sum_{k\ge1}\frac{z^k}{k}
\sum_{\ell(\mathcal O)\mid k}\ell(\mathcal O)\\
&=
\sum_{\mathcal O\in\mathscr P}
\sum_{r\ge1}
\frac{z^{r\ell(\mathcal O)}}{r}\\
&=
-\sum_{\mathcal O\in\mathscr P}
\log\left(1-z^{\ell(\mathcal O)}\right).
\end{aligned}
\]

Exponentiating gives (1.5). \ensuremath{\blacksquare}

The trivial orbit has length \(2\), so its Euler factor is

\[
\boxed{(1-z^2)^{-1}.}
\tag{1.6}
\]

Thus the desired zeta identity is

\[
\boxed{\zeta_T(z)=\frac1{1-z^2}.}
\tag{1.7}
\]

As a formal consequence of Proposition 1.2, (1.7) is already equivalent to saying that \(1\leftrightarrow2\) is the only positive-integer periodic orbit. The determinant construction below makes this equivalence analytic and trace-theoretic, but does not make it easier by itself.


\bigskip



\subsubsection*{2. A genuine trace-class Collatz operator}


Let

\[
\mathcal H=\ell^2(\mathbb N).
\]

Define the Koopman composition operator

\[
(Cf)(n)=f(Tn).
\tag{2.1}
\]

Every \(m\) has at most two positive predecessors:

\[
2m
\]

and, when it is integral,

\[
\frac{2m-1}{3}.
\]

Therefore, writing \(d(m)=\#T^{-1}(m)\le2\),

\[
\|Cf\|_2^2
=
\sum_{n\ge1}|f(Tn)|^2
=
\sum_{m\ge1}d(m)|f(m)|^2
\le 2\|f\|_2^2.
\tag{2.2}
\]

Thus

\[
\|C\|\le\sqrt2.
\]

For \(0<t<1\), let \(M_t\) be the diagonal operator

\[
(M_tf)(n)=t^n f(n).
\tag{2.3}
\]

It is trace class because

\[
\|M_t\|_1=\sum_{n\ge1}t^n=\frac{t}{1-t}.
\tag{2.4}
\]

Define

\[
\boxed{L_t:=M_tC.}
\tag{2.5}
\]

Then \(L_t\) is trace class.

For \(k\ge1\), put

\[
S_k(n):=\sum_{j=0}^{k-1}T^j(n).
\tag{2.6}
\]

A direct induction gives

\[
(L_t^k f)(n)
=
t^{S_k(n)}f(T^k n).
\tag{2.7}
\]


\paragraph*{Proposition 2.1: exact trace formula}


For every \(k\ge1\),

\[
\boxed{
\operatorname{tr}(L_t^k)
=
\sum_{\substack{n\ge1\\T^k n=n}}
t^{S_k(n)}.
}
\tag{2.8}
\]


\subparagraph*{Proof}


In the standard basis \(e_n\),

\[
(L_t^k e_n)(n)
=
t^{S_k(n)}
\mathbf 1_{\{T^kn=n\}}.
\]

Since \(L_t^k\) is trace class, its trace is the sum of its diagonal entries. \ensuremath{\blacksquare}

Thus there is no symbolic or probabilistic contamination: the trace counts exactly the positive-integer periodic points, with a strictly positive height weight.


\bigskip



\subsubsection*{3. Fredholm determinant and weighted Euler product}


Define

\[
\boxed{
\mathcal D_t(z):=\det(I-zL_t).
}
\tag{3.1}
\]

Because \(L_t\) is trace class, \(\mathcal D_t(z)\) is an entire function of \(z\). Near \(z=0\),

\[
\log \mathcal D_t(z)
=
-\sum_{k=1}^{\infty}
\frac{z^k}{k}\operatorname{tr}(L_t^k).
\tag{3.2}
\]

For a primitive orbit \(\mathcal O\), define its additive height

\[
H(\mathcal O):=\sum_{n\in\mathcal O}n.
\tag{3.3}
\]


\paragraph*{Theorem 3.1: weighted Collatz determinant formula}


As an analytic identity near \(z=0\),

\[
\boxed{
\mathcal D_t(z)
=
\prod_{\mathcal O\in\mathscr P}
\left(
1-z^{\ell(\mathcal O)}t^{H(\mathcal O)}
\right).
}
\tag{3.4}
\]

The Fredholm determinant is the entire continuation of this local Euler product.


\subparagraph*{Proof}


Suppose \(\mathcal O\) has primitive length \(\ell\). For \(k=r\ell\), each of its \(\ell\) points contributes

\[
t^{rH(\mathcal O)}
\]

to \(\operatorname{tr}(L_t^{r\ell})\). Its total contribution to (3.2) is therefore

\[
-\sum_{r\ge1}
\frac{z^{r\ell}}{r\ell}
\ell t^{rH(\mathcal O)}
=
-\sum_{r\ge1}
\frac{
(z^\ell t^{H(\mathcal O)})^r
}{r}.
\]

This equals

\[
\log\left(
1-z^\ell t^{H(\mathcal O)}
\right).
\]

Summing over primitive orbits proves (3.4). \ensuremath{\blacksquare}

For the trivial orbit \(\{1,2\}\),

\[
\ell=2,\qquad H=1+2=3,
\]

so its determinant factor is

\[
\boxed{1-z^2t^3.}
\tag{3.5}
\]

Equivalently, its weighted zeta factor is

\[
(1-z^2t^3)^{-1}.
\]


\bigskip



\subsubsection*{4. Exact equivalence with the positive-cycle problem}


The Fredholm determinant representation is not merely compatible with the cycle problem; its desired evaluation is exactly equivalent to solving that problem.


\paragraph*{Theorem 4.1}


Fix any \(t\in(0,1)\). The following are equivalent:

\[
\mathcal D_t(z)=1-z^2t^3,
\tag{4.1}
\]

and

\[
\text{\(T\) has no positive periodic orbit other than \(1\leftrightarrow2\).}
\tag{4.2}
\]


\subparagraph*{Proof}


If the trivial cycle is the only primitive orbit, (3.4) immediately gives (4.1).

Conversely, suppose a nontrivial primitive orbit \(\mathcal O\) exists, with length \(\ell\). Its contribution to

\[
\operatorname{tr}(L_t^\ell)
\]

is

\[
\ell\,t^{H(\mathcal O)}>0.
\]

All trace contributions are nonnegative. Thus this term cannot be canceled by another orbit. Consequently the logarithmic Taylor series of \(\mathcal D_t\) differs from that of \(1-z^2t^3\). \ensuremath{\blacksquare}

In trace form, the required identity is

\[
\boxed{
\operatorname{tr}(L_t^k)
=
\begin{cases}
0,&k\ \text{odd},\\[1mm]
2t^{3k/2},&k\ \text{even}.
\end{cases}
}
\tag{4.3}
\]

Indeed, for \(k=2r\), the two trivial periodic points each have orbit sum \(3r\).

Taking \(t\uparrow1\) coefficientwise gives

\[
\mathcal D_t(z)
\longrightarrow
\exp\left(
-\sum_{k\ge1}\frac{a_k}{k}z^k
\right)
=
\zeta_T(z)^{-1}.
\tag{4.4}
\]

Therefore

\[
\mathcal D_t(z)=1-z^2t^3
\quad\text{for one, equivalently every, }t\in(0,1)
\]

would imply

\[
\zeta_T(z)=\frac1{1-z^2}.
\]

The obstacle is that trace-classness supplies convergence, not the numerical values (4.3).


\bigskip



\subsubsection*{5. What the determinant identity means arithmetically}


For a parity word \(w\) of length \(K\) and weight \(m\), put

\[
\Delta(w):=2^K-3^m.
\tag{5.1}
\]

The fixed point is

\[
N(w)=\frac{\Phi(w)}{\Delta(w)}.
\tag{5.2}
\]

A minimal word represents a positive integer cycle exactly when

\[
\Delta(w)>0,
\qquad
\Delta(w)\mid\Phi(w).
\tag{5.3}
\]

Divisibility propagates under cyclic rotation, and the resulting rotated rational points are all integers with the prescribed parities. [\hypertarget{cit:R67:3}{}\hyperlink{ref:R67:3}{3}]

Hence proving

\[
\mathcal D_t(z)=1-z^2t^3
\]

amounts to proving

\[
\boxed{
\Delta(w)\nmid\Phi(w)
}
\tag{5.4}
\]

for every nontrivial candidate word, apart from repetitions and rotations of \(10\).

The elementary product identity around a positive cycle gives

\[
\frac{2^K}{3^m}
=
\prod_{i=1}^{m}
\left(1+\frac1{3b_i}\right),
\tag{5.5}
\]

where \(b_i\) are the odd cycle values. It follows immediately that

\[
K\le2m,
\]

with equality only for the trivial orbit. Thus every nontrivial candidate lies in the narrow strip

\[
m\log_2 3<K<2m.
\tag{5.6}
\]

These facts, including uniqueness of the trivial boundary case \(K=2m\), are recorded in the parity-word analysis. [\hypertarget{cit:R67:4}{}\hyperlink{ref:R67:4}{4}]

Recent Christoffel-word work sharpens the extremal structure of these candidate words, but explicitly leaves hypothetical integer cycles in the critical strip as strictly suboptimal parity patterns rather than excluding them altogether. [\hypertarget{cit:R67:5}{}\hyperlink{ref:R67:5}{5}] Existing rigorous circuit results also exclude only finite circuit ranges, such as Hercher's exclusion through 91 circuits. [\hypertarget{cit:R67:6}{}\hyperlink{ref:R67:6}{6}]

The determinant does not create analytic cancellation capable of overcoming (5.4), because every periodic-point contribution to every trace is positive. It repackages the divisibility problem into the statement that certain positive trace coefficients vanish.

This agrees with the prior arithmetic audit: the exact decoder identifies a possible completion, but the unresolved step is an all-scale modular anti-alignment theorem, not a further uniqueness argument. [\hypertarget{cit:R67:7}{}\hyperlink{ref:R67:7}{7}]


\bigskip



\subsubsection*{6. Why the full \(2\)-adic determinant gives the wrong zeta function}


There is a tempting alternative: extend \(T\) to \(\mathbb Z_2\) and use the parity-vector conjugacy. Bernstein and Lagarias proved that the \(2\)-adic Collatz map is conjugate to the binary shift. [\hypertarget{cit:R67:8}{}\hyperlink{ref:R67:8}{8}]

The \(k\)-th iterate of the binary shift has exactly \(2^k\) fixed points, one for every periodic binary block of length \(k\). Consequently its zeta function is

\[
\exp\left(
\sum_{k\ge1}\frac{2^k}{k}z^k
\right)
=
\frac1{1-2z}.
\tag{6.1}
\]

Thus

\[
\boxed{
\zeta_{T:\mathbb Z_2\to\mathbb Z_2}(z)
=
\frac1{1-2z},
}
\tag{6.2}
\]

not \((1-z^2)^{-1}\).

This is not a defect of the \(2\)-adic conjugacy. Every periodic parity word genuinely defines a rational \(2\)-adic periodic point. What distinguishes a positive-integer cycle is exactly the divisibility condition

\[
2^K-3^m\mid\Phi(w).
\]

Therefore any determinant built on the whole \(2\)-adic space counts far too many cycles. Projecting onto ordinary positive integers is the entire arithmetic difficulty.


\bigskip



\subsubsection*{7. A zeta function cannot detect escape}


Even a complete proof of

\[
\zeta_T(z)=\frac1{1-z^2}
\]

would settle only the cycle half of Collatz.

For example, the map

\[
F(n)=n+1
\]

has no periodic points, so its dynamical zeta function is identically \(1\), yet every orbit escapes to infinity.

More generally, one may attach arbitrarily many acyclic infinite rays or rooted trees to a dynamical system without changing any periodic orbit and hence without changing its zeta function. The periodic-orbit data cannot distinguish a globally terminating system from one containing an escaping ray.

For Collatz, a counterexample avoiding \(1\leftrightarrow2\) is either a nontrivial positive cycle or an injective infinite orbit. The zeta function rules out only the first possibility.


\bigskip



\subsubsection*{8. Exact ``zeta plus nonescape'' theorem}


Let

\[
A:=\{1,2\},
\qquad
X:=\mathbb N\setminus A.
\]

The following gives a precise completion of the proposed architecture.


\paragraph*{Theorem 8.1}


Fix \(t\in(0,1)\). The Collatz conjecture is equivalent to the conjunction of:

\[
\mathcal D_t(z)=1-z^2t^3,
\tag{8.1}
\]

and the existence of a function

\[
\pi:X\longrightarrow(0,\infty)
\]

such that

\[
\sum_{n\in X}\pi(n)<\infty
\tag{8.2}
\]

and

\[
\pi(n)\le\pi(Tn)
\qquad
(n,Tn\in X).
\tag{8.3}
\]


\subparagraph*{Proof: Collatz implies the two conditions}


The determinant identity follows because the trivial cycle is the only cycle.

For \(n\in X\), let

\[
\tau(n):=\min\{k\ge1:T^k n\in A\}
\]

be the hitting time of \(A\). Choose \(q\in(0,1/2)\) and put

\[
\pi(n)=q^{\tau(n)}.
\tag{8.4}
\]

If \(n,Tn\in X\), then

\[
\tau(Tn)=\tau(n)-1,
\]

so

\[
\pi(Tn)=q^{-1}\pi(n)>\pi(n).
\]

It remains to prove summability. Every point has at most two predecessors, so the number of integers with hitting time \(r\) is at most

\[
2\cdot2^r=2^{r+1}.
\]

Hence

\[
\sum_{n\in X}\pi(n)
\le
\sum_{r\ge1}2^{r+1}q^r
=
2\sum_{r\ge1}(2q)^r
<\infty.
\]


\subparagraph*{Proof: the two conditions imply Collatz}


Suppose an orbit starting at \(n\) never enters \(A\).

If it repeats, it contains a nontrivial positive cycle, contradicting (8.1).

It must therefore consist of infinitely many distinct integers. By (8.3),

\[
\pi(T^j n)\ge\pi(n)>0
\qquad(j\ge0).
\]

Since these are infinitely many distinct states,

\[
\sum_{m\in X}\pi(m)
\ge
\sum_{j\ge0}\pi(T^j n)
=
\infty,
\]

contradicting (8.2). \ensuremath{\blacksquare}

Thus the proposed route can be completed to the exact theorem

\[
\boxed{
\text{Collatz}
\iff
\text{trivial zeta factor}
+
\text{positive proper nonescape weight}.
}
\tag{8.5}
\]

The problem is now concentrated in constructing \(\pi\) without using the unknown hitting times.


\bigskip



\subsubsection*{9. A stronger single transfer-operator certificate}


There is a more economical formulation in which one positive vector excludes both escape and nontrivial cycles.

Define the killed inverse-transfer operator \(P:\ell^1(X)\to\ell^1(X)\) by

\[
(Pu)(m)
=
\sum_{\substack{n\in X\\Tn=m}}u(n).
\tag{9.1}
\]

With the convention that \(u\) is zero outside \(X\),

\[
(Pu)(m)
=
u(2m)
+
\mathbf 1_{m\equiv2\pmod3}
u\left(\frac{2m-1}{3}\right).
\tag{9.2}
\]

Since every \(n\) has exactly one forward image,

\[
\|Pu\|_1\le\|u\|_1.
\]


\paragraph*{Theorem 9.1: exact killed-transfer criterion}


The Collatz conjecture is equivalent to the existence of

\[
\pi\in\ell^1(X),
\qquad
\pi(n)>0,
\]

such that

\[
\boxed{P\pi\le\pi.}
\tag{9.3}
\]


\subparagraph*{Proof: Collatz implies the certificate}


Again take

\[
\pi(n)=q^{\tau(n)}
\]

with \(0<q<1/2\). If \(Tn=m\in X\), then

\[
\tau(n)=\tau(m)+1,
\]

so every predecessor contributes \(q\pi(m)\). There are at most two predecessors, hence

\[
(P\pi)(m)\le2q\pi(m)<\pi(m).
\]

Summability was proved in Theorem 8.1.


\subparagraph*{Proof: the certificate implies Collatz}


For every edge \(n\mapsto m\) within \(X\), \(\pi(n)\) is one of the positive summands in \(P\pi(m)\). Thus

\[
\pi(n)\le(P\pi)(m)\le\pi(m).
\tag{9.4}
\]

An infinite injective orbit would therefore contain infinitely many distinct states of \(\pi\)-mass at least \(\pi(n)>0\), contradicting \(\pi\in\ell^1\).

It remains to exclude a nontrivial cycle. Along a cycle, (9.4) forces all \(\pi\)-values to be equal. Choose an odd cycle value \(n>1\), and put

\[
m=Tn=\frac{3n+1}{2}.
\]

The point \(m\) has both the odd predecessor \(n\) and the even predecessor

\[
2m=3n+1\in X.
\]

Consequently,

\[
(P\pi)(m)
\ge
\pi(n)+\pi(2m)
>
\pi(n)
=
\pi(m),
\]

contradicting \(P\pi\le\pi\). \ensuremath{\blacksquare}

This theorem is stronger than the requested zeta architecture:

\[
\boxed{
P\pi\le\pi,\quad
0<\pi\in\ell^1(X)
}
\]

would prove both the absence of nontrivial cycles and nonescape, so the separate zeta identity would become redundant.

It also shows why merely asserting that such a vector exists is not a proof. Under Collatz, the vector can be manufactured from the hitting-time function. An independently checkable proof requires a finite-state, residue-class, automaton, algebraic, or otherwise explicit description of \(\pi\). This matches the earlier infinite-matrix diagnosis: an unrestricted positive subeigenvector is equivalent in strength to the unknown totality statement, while a finitely presented one would be a genuine certificate. [\hypertarget{cit:R67:9}{}\hyperlink{ref:R67:9}{9}]

Writing

\[
H(n):=-\log\pi(n)
\]

turns (9.4) into

\[
H(Tn)\le H(n).
\tag{9.5}
\]

Moreover, \(\pi\in\ell^1\) implies every sublevel set of \(H\) is finite. Thus \(H\) is a coercive nonincreasing Lyapunov height. The branching term in (9.3), not merely the edgewise inequality, is what excludes cycles.


\bigskip



\subsubsection*{10. The Green-kernel attempt and its exact failure point}


The natural way to construct such a \(\pi\) is to start with a positive source

\[
b\in\ell^1(X),\qquad b(n)>0,
\]

and form the killed Green potential

\[
\boxed{
G_b:=\sum_{j=0}^{\infty}P^j b.
}
\tag{10.1}
\]

Whenever this converges in \(\ell^1\),

\[
PG_b
=
\sum_{j\ge1}P^j b
=
G_b-b
<
G_b,
\tag{10.2}
\]

so Theorem 9.1 proves Collatz.

The summability condition can be calculated exactly. Let

\[
\tau_A(n)
=
\min\{k\ge1:T^k n\in A\},
\]

with \(\tau_A(n)=\infty\) when the orbit never enters \(A\). Since \(P\) transports mass forward and kills it when it reaches \(A\), Tonelli's theorem gives

\[
\begin{aligned}
\|G_b\|_1
&=
\sum_{j\ge0}\|P^j b\|_1\\
&=
\sum_{n\in X}b(n)\tau_A(n).
\end{aligned}
\tag{10.3}
\]

Thus

\[
\boxed{
G_b\in\ell^1
\iff
\sum_{n\in X}b(n)\tau_A(n)<\infty.
}
\tag{10.4}
\]

This identifies the precise obstruction. If even one \(n\) has \(\tau_A(n)=\infty\), then every strictly positive source \(b\) makes the right-hand side infinite. Therefore proving Green-kernel summability already proves nonescape.

The algebra of the Green construction is complete; the missing theorem is an all-depth bound on its total mass.

This is also visible through the Fredholm resolvent. For sufficiently small \(z\),

\[
(I-zL_t)^{-1}
=
\sum_{k\ge0}z^kL_t^k,
\]

and

\[
\bigl((I-zL_t)^{-1}e_a\bigr)(n)
=
\sum_{\substack{k\ge0\\T^k n=a}}
z^k t^{S_k(n)}.
\tag{10.5}
\]

The determinant \(\mathcal D_t\) is the denominator of this resolvent. Reachability of \(a\) is encoded in individual matrix entries or Fredholm minors\textemdash{}the numerators\textemdash{}not in the determinant alone. An escaping orbit can alter all the relevant hitting entries while contributing no periodic-orbit factor to \(\mathcal D_t\).


\bigskip



\subsubsection*{11. Finite determinant approximations}


Let \(P_N\) be the projection onto

\[
\operatorname{span}\{e_1,\ldots,e_N\}
\]

and define

\[
L_{t,N}:=P_NL_tP_N.
\]

Then

\[
\mathcal D_{t,N}(z)
:=
\det(I-zL_{t,N})
\tag{11.1}
\]

is an ordinary finite determinant. Because the finite directed graph has out-degree at most one, its trees contribute no determinant factor, while its directed cycles contribute

\[
\mathcal D_{t,N}(z)
=
\prod_{\substack{\mathcal O\in\mathscr P\\
\mathcal O\subseteq[1,N]}}
\left(
1-z^{\ell(\mathcal O)}t^{H(\mathcal O)}
\right).
\tag{11.2}
\]

Since \(L_t\) is trace class,

\[
\|L_{t,N}-L_t\|_1\longrightarrow0,
\]

and hence

\[
\mathcal D_{t,N}(z)\longrightarrow\mathcal D_t(z)
\]

locally uniformly in \(z\).

This gives a rigorous computational approximation scheme. However, showing

\[
\mathcal D_{t,N}(z)=1-z^2t^3
\]

for each finite \(N\) merely says that no nontrivial cycle is contained below \(N\). Passing to all \(N\) requires a theorem excluding new cycles at unbounded heights.


\bigskip



\subsubsection*{12. Final status of Approach 62}


The following parts of the proposed route are complete:

\[
\zeta_T(z)
=
\exp\left(
\sum_{k\ge1}\frac{\#\operatorname{Fix}(T^k)}{k}z^k
\right),
\]

\[
\zeta_T(z)
=
\prod_{\mathcal O\in\mathscr P}
(1-z^{\ell(\mathcal O)})^{-1},
\]

\[
\mathcal D_t(z)
=
\det(I-zM_tC),
\]

and

\[
\mathcal D_t(z)
=
\prod_{\mathcal O\in\mathscr P}
\left(
1-z^{\ell(\mathcal O)}t^{H(\mathcal O)}
\right).
\]

Moreover,

\[
\boxed{
\mathcal D_t(z)=1-z^2t^3
\iff
\text{there is no nontrivial positive cycle}.
}
\]

The exact nonescape supplement is

\[
\boxed{
\exists\,0<\pi\in\ell^1(X):
\quad
P\pi\le\pi.
}
\]

That condition by itself is equivalent to the full Collatz conjecture.

What remains unproved is therefore not a functional-analytic technicality. It is one of the following equivalent arithmetic constructions:

\[
\operatorname{tr}(L_t^k)
=
\begin{cases}
0,&k\text{ odd},\\
2t^{3k/2},&k\text{ even},
\end{cases}
\quad\text{for every }k,
\tag{12.1}
\]

together with an independent nonescape theorem; or, more efficiently,

\[
\boxed{
\text{construct explicitly }
0<\pi\in\ell^1(X)
\text{ satisfying }
\sum_{Tn=m}\pi(n)\le\pi(m)
\text{ for every }m\in X.
}
\tag{12.2}
\]

The existing compression framework isolates the same infinitary obstruction: finite trace and carry statements can terminate particular recurrent subsystems, but a separate global synchronization theorem is needed to force every positive stationary ray into a descending configuration. [\hypertarget{cit:R67:10}{}\hyperlink{ref:R67:10}{10}]

Accordingly, the dynamical zeta approach gives an exact and clean reformulation\textemdash{}arguably the sharpest operator formulation of the ``cycle plus escape'' dichotomy\textemdash{}but the determinant's trivial-factor evaluation and the summable Green-kernel estimate have not been proved here. Replacing either one by the desired conclusion would amount to assuming the corresponding half of the Collatz conjecture.
\subsection*{References}
\begin{enumerate}[leftmargin=2.2em]
\item \hypertarget{ref:R67:1}{}\hyperlink{cit:R67:1}{\texttt{turn3file0 L23-L55}}
\item \hypertarget{ref:R67:2}{}\hyperlink{cit:R67:2}{\texttt{turn3file0 L73-L83}}
\item \hypertarget{ref:R67:3}{}\hyperlink{cit:R67:3}{\texttt{turn3file0 L56-L83}}
\item \hypertarget{ref:R67:4}{}\hyperlink{cit:R67:4}{\texttt{turn3file0 L84-L120}}
\item \hypertarget{ref:R67:5}{}\hyperlink{cit:R67:5}{\texttt{turn542179view0}}
\item \hypertarget{ref:R67:6}{}\hyperlink{cit:R67:6}{\texttt{turn542179view3}}
\item \hypertarget{ref:R67:7}{}\hyperlink{cit:R67:7}{\texttt{turn4file1 L1060-L1068}}
\item \hypertarget{ref:R67:8}{}\hyperlink{cit:R67:8}{\texttt{turn982505search1}}
\item \hypertarget{ref:R67:9}{}\hyperlink{cit:R67:9}{\texttt{turn6file0 L12-L22}}
\item \hypertarget{ref:R67:10}{}\hyperlink{cit:R67:10}{\texttt{turn4file0 L211-L231}}
\end{enumerate}

\clearpage
\section{Round 68: approach 63 --- A uniform determinant involution}
\subsection*{Prompt}
\begin{quote}
Complete a full proof of the Collatz conjecture using the following approach:

\begin{enumerate}
\item A uniform determinant involution
\end{enumerate}

\textbf{Established/Hybrid \textperiodcentered{} P for cycles}

The determinant formulation asks for det M\_k = 1 - x\textasciicircum{}2 for every k. Find a sign-reversing involution on all determinant permutations except those representing the trivial two-cycle. That would provide an elegant cycle-exclusion theorem.

Take as long as you need.
\end{quote}
\subsection*{Response}

\subsection*{Outcome}


The proposed sign-reversing involution does \textbf{not} yield a full proof. More strongly, it cannot exist as a nontrivial cancellation mechanism: the determinant expansion is already cancellation-free. Every nonzero determinant permutation is exactly a disjoint union of genuine Collatz cycles. Consequently, an involution pairing away hypothetical nontrivial-cycle terms exists only after one has already proved that those terms do not exist.

This gives a rigorous resolution of Approach 63, including an exact determinant factorization, a no-go theorem for the proposed involution, and a precise matrix expression for the still-missing nonescape statement.

The matrices below are the Collatz matrices introduced by Alves\textendash{}Gra\c{c}a\textendash{}Dias\textendash{}Sousa Ramos and restated by Kauffman\textendash{}Lopes; their determinant criterion concerns uniqueness of the finite periodic orbit, not exclusion of divergent trajectories. [\hypertarget{cit:R68:1}{}\hyperlink{ref:R68:1}{1}]


\subsubsection*{1. The exact determinant expansion}


Use the shortcut map

\[
T(n)=
\begin{cases}
n/2,&n\equiv0\pmod 2,\\[2mm]
(3n+1)/2,&n\equiv1\pmod 2.
\end{cases}
\]

For \(k\ge2\), let

\[
A_k(i,j)=
\begin{cases}
1,&T(i)=j\le k,\\
0,&\text{otherwise},
\end{cases}
\qquad
M_k(x)=I_k+xA_k.
\]

Thus row \(i\) of \(M_k\) has its diagonal \(1\), together with an \(x\) in column \(T(i)\) whenever \(T(i)\le k\).

Let \(\mathcal C_k\) denote the set of directed \(T\)-cycles contained in \(\{1,\ldots ,k\}\).


\paragraph*{Theorem 1 \textemdash{} Functional-digraph determinant formula}


For every \(k\ge2\),

\[
\boxed{
\det M_k(x)
=
\prod_{C\in\mathcal C_k}
\left(1-(-x)^{|C|}\right).
}
\tag{1}
\]


\paragraph*{Proof}


Write the Leibniz expansion

\[
\det M_k(x)
=
\sum_{\sigma\in S_k}
\operatorname{sgn}(\sigma)
\prod_{i=1}^k M_k(i,\sigma(i)).
\]

A term indexed by \(\sigma\) is nonzero only when

\[
\sigma(i)\in\{i,T(i)\}
\qquad (1\le i\le k).
\tag{2}
\]

Consider a nontrivial cycle

\[
(i_1\,i_2\,\cdots\,i_\ell)
\]

in the disjoint-cycle decomposition of \(\sigma\). Because its length is at least two,

\[
\sigma(i_j)\ne i_j,
\]

so condition (2) forces

\[
\sigma(i_j)=T(i_j)=i_{j+1},
\]

with indices interpreted cyclically. Therefore every nontrivial permutation cycle appearing in a nonzero determinant term is an actual \(T\)-cycle.

Conversely, given any collection \(S\subseteq\mathcal C_k\), define \(\sigma_S\) by applying \(T\) on the vertices belonging to cycles in \(S\), and fixing all other vertices. Distinct cycles of a function are disjoint, so \(\sigma_S\) is a permutation and its determinant term is nonzero.

This gives a bijection

\[
\{\text{nonzero determinant permutations}\}
\longleftrightarrow
\mathcal P(\mathcal C_k).
\]

If

\[
L(S)=\sum_{C\in S}|C|,
\]

then

\[
\operatorname{sgn}(\sigma_S)
=
\prod_{C\in S}(-1)^{|C|-1}
=
(-1)^{L(S)-|S|}
\]

and the product of matrix entries is \(x^{L(S)}\). Hence

\[
\det M_k(x)
=
\sum_{S\subseteq\mathcal C_k}
\prod_{C\in S}
\left((-1)^{|C|-1}x^{|C|}\right),
\]

which factors as

\[
\prod_{C\in\mathcal C_k}
\left(1+(-1)^{|C|-1}x^{|C|}\right)
=
\prod_{C\in\mathcal C_k}
\left(1-(-x)^{|C|}\right).
\]

This proves (1). \(\square\)


\subsubsection*{2. The multivariate strengthening}


The lack of cancellation becomes even clearer if every row-edge receives an independent variable. Put

\[
M_k(\mathbf y)_{ij}
=
\delta_{ij}+y_i\mathbf 1_{\{T(i)=j\le k\}}.
\]

The same proof gives

\[
\boxed{
\det M_k(\mathbf y)
=
\prod_{C\in\mathcal C_k}
\left(
1-\prod_{i\in C}(-y_i)
\right).
}
\tag{3}
\]

Thus each cycle contributes its own distinct squarefree monomial

\[
\prod_{i\in C}y_i.
\]

There is no hidden collision of determinant terms before specializing all \(y_i\) to the same variable \(x\).


\subsubsection*{3. The trivial-cycle factor and the actual pairing}


The cycle

\[
C_0=(1\,2)
\]

is always present. Its factor in (1) is

\[
1-(-x)^2=1-x^2.
\]

Consequently,

\[
\boxed{
\det M_k(x)
=
(1-x^2)\,
Q_k(x),
}
\tag{4}
\]

where

\[
Q_k(x)
=
\prod_{\substack{C\in\mathcal C_k\\ C\ne C_0}}
\left(1-(-x)^{|C|}\right).
\tag{5}
\]

There is a canonical two-term pairing that extracts the factor \(1-x^2\): pair a cycle collection \(S\) not containing \(C_0\) with \(S\cup\{C_0\}\). Their contributions are

\[
w(S)
\quad\text{and}\quad
-x^2w(S),
\]

so their sum is

\[
(1-x^2)w(S).
\]

This pairing proves divisibility by \(1-x^2\). It does \textbf{not} prove that the quotient is \(1\); the quotient is exactly the product (5) over nontrivial cycles.


\subsubsection*{4. No-go theorem for the requested involution}


A sign-reversing involution proving a polynomial identity must pair terms of exactly opposite polynomial weight:

\[
w(\iota(\sigma))=-w(\sigma).
\tag{6}
\]

Merely reversing the sign of the permutation while changing the power of \(x\) does not cancel the terms.


\paragraph*{Theorem 2 \textemdash{} Involution obstruction}


Suppose that \(\mathcal C_k\) contains at least one nontrivial cycle. Then the nonconstant terms of \(Q_k(x)\) cannot be partitioned into weight-preserving, sign-reversing pairs.


\paragraph*{Proof}


Let

\[
\mathcal C_k^*=\mathcal C_k\setminus\{C_0\}
\]

and assume \(\mathcal C_k^*\ne\varnothing\). In the subset expansion of \(Q_k\), consider the subset containing \textbf{all} nontrivial cycles:

\[
S_*=\mathcal C_k^*.
\]

Its degree is

\[
D_*=\sum_{C\in\mathcal C_k^*}|C|.
\]

Every other subset omits at least one cycle of positive length, and therefore has degree strictly smaller than \(D_*\). Thus the term indexed by \(S_*\) is the unique term of highest degree in \(Q_k\). Its coefficient is \(+1\) or \(-1\).

It has no term of the same degree with the opposite coefficient, so it cannot be paired under (6). \(\square\)

Therefore:

\[
\boxed{
\text{The requested involution exists nontrivially only if }
\mathcal C_k^*=\varnothing.
}
\]

But when \(\mathcal C_k^*=\varnothing\), there is nothing left to pair. In that case the quotient \(Q_k\) is already \(1\).

So the proposed involution is equivalent to cycle exclusion in the following vacuous sense:

\[
\begin{aligned}
&\text{there is a cancellation involution leaving only }
1-x^2\\
&\qquad\Longleftrightarrow
\text{there are no nontrivial determinant terms}\\
&\qquad\Longleftrightarrow
\text{there are no nontrivial Collatz cycles}.
\end{aligned}
\]

It does not supply an independent mechanism for proving the last statement.


\subsubsection*{5. Exact interpretation of the recurrence in \(k\)}


Formula (1) also completely explains the attempted induction

\[
\det M_k=\det M_{k-1}.
\]

Any cycle contained in \(\{1,\ldots ,k\}\) but not in \(\{1,\ldots ,k-1\}\) must contain \(k\). Since two cycles of a function cannot share a vertex, at most one new cycle can occur. Therefore

\[
\boxed{
\frac{\det M_k(x)}{\det M_{k-1}(x)}
=
\begin{cases}
1,&k\text{ is not the maximum of a Collatz cycle},\\[2mm]
1-(-x)^\ell,
 &k\text{ is the maximum of a cycle of length }\ell.
\end{cases}
}
\tag{7}
\]

Thus proving \(\det M_k=\det M_{k-1}\) is exactly the same as proving that \(k\) is not the maximum of a nontrivial cycle. The Laplace-expansion difficulties found in the earlier matrix work are not artifacts of the chosen calculation; they are the cycle problem itself. Kauffman\textendash{}Lopes similarly identify the remaining minor with a closed inverse Collatz orbit. [\hypertarget{cit:R68:2}{}\hyperlink{ref:R68:2}{2}, \hypertarget{cit:R68:3}{}\hyperlink{ref:R68:3}{3}]


\subsubsection*{6. Arithmetic information from the maximum}


Although it does not finish cycle exclusion, the determinant-change interpretation immediately recovers the sharp first congruence restriction.


\paragraph*{Proposition 3}


If \(K>2\) is the maximum element of a positive Collatz cycle, then

\[
\boxed{K\equiv8\pmod{18}.}
\tag{8}
\]


\paragraph*{Proof}


The maximum \(K\) cannot be odd, because for odd \(K>1\),

\[
T(K)=\frac{3K+1}{2}>K.
\]

Hence \(K\) is even and

\[
T(K)=K/2.
\]

Let \(a\) be the predecessor of \(K\) in the cycle. An even predecessor would be \(2K>K\), impossible. Therefore \(a\) is odd and

\[
a=\frac{2K-1}{3}.
\tag{9}
\]

Let \(b\) be the predecessor of \(a\). The even candidate is

\[
2a=\frac{4K-2}{3}>K
\]

because \(K>2\). It is therefore also impossible, so

\[
b=\frac{2a-1}{3}
\]

must be an integer. This requires

\[
a\equiv2\pmod3.
\]

Using (9),

\[
\frac{2K-1}{3}\equiv2\pmod3,
\]

so

\[
2K-1\equiv6\pmod9,
\qquad
2K\equiv7\pmod9,
\]

and hence

\[
K\equiv8\pmod9.
\]

Since \(K\) is even,

\[
K\equiv8\pmod{18}.
\]

\(\square\)

This explains why the exceptional determinant indices occur in the residue class \(8\bmod18\). Further inverse branching is precisely where the unrestricted cycle problem reappears.


\subsubsection*{7. Why the determinant identity is not the full Collatz conjecture}


The full conjecture has two logically separate obstructions:

\begin{enumerate}
\item a nontrivial finite cycle;
\item an unbounded nonperiodic orbit.
\end{enumerate}

The determinant identity sees only the first. Paparella's corrected matrix formulation explicitly distinguishes convergent trajectories, nontrivial cycles, and divergent trajectories, and calls the assertion excluding nontrivial cycles the \textbf{aperiodic Collatz conjecture}. [\hypertarget{cit:R68:4}{}\hyperlink{ref:R68:4}{4}]

Here is a decisive countermodel to any inference from the determinant identity to convergence. Define a different map \(F:\mathbb N\to\mathbb N\) by

\[
F(1)=2,\qquad F(2)=1,\qquad F(n)=n+1\quad(n\ge3).
\]

Its only finite cycle is \(1\leftrightarrow2\). Therefore its corresponding matrices satisfy

\[
\det(I+xA_k^F)=1-x^2
\qquad\text{for every }k.
\]

Nevertheless,

\[
F^m(3)=3+m\longrightarrow\infty.
\]

Thus even the complete family of determinant identities is compatible with divergence.


\subsubsection*{8. The blindness extends to Fredholm determinants}


The limitation is not peculiar to finite truncations.

Let \(L_w\) be a weighted transition operator

\[
L_we_n=w_ne_{T(n)}
\]

with weights chosen so that \(L_w\) is trace class. Then

\[
L_w^me_n
=
\left(\prod_{j=0}^{m-1}w_{T^j(n)}\right)e_{T^m(n)},
\]

and therefore

\[
\operatorname{tr}(L_w^m)
=
\sum_{T^m(n)=n}
\prod_{j=0}^{m-1}w_{T^j(n)}.
\]

Only periodic points occur in this trace. Inserting this into the Fredholm determinant formula gives

\[
\boxed{
\det(I-zL_w)
=
\prod_{C\text{ finite cycle}}
\left(
1-z^{|C|}
\prod_{n\in C}w_n
\right).
}
\tag{10}
\]

Every infinite acyclic orbit contributes exactly \(1\). Thus every ordinary trace/Fredholm determinant of this transition type is intrinsically blind to nonperiodic escape.


\subsubsection*{9. The correct matrix quantity for detecting escape}


There is, however, a determinant-compatible object that records escape.

Remove the trivial vertices and let

\[
V_k=\{3,\ldots ,k\},
\qquad
C_k(i,j)=\mathbf 1_{\{T(i)=j\}},
\quad i,j\in V_k.
\]

Define the upper-boundary vector

\[
u_k(i)=
\begin{cases}
1,&T(i)>k,\\
0,&\text{otherwise}.
\end{cases}
\]

Assuming cycle exclusion, \(C_k\) is nilpotent. For \(n\in V_k\), define

\[
E_{n,k}(z)
=
z\,e_n^{\mathsf T}(I-zC_k)^{-1}u_k.
\tag{11}
\]

Since

\[
(I-zC_k)^{-1}
=
\sum_{m\ge0}z^mC_k^m,
\]

we obtain

\[
E_{n,k}(z)
=
\sum_{m\ge0}
z^{m+1}e_n^{\mathsf T}C_k^mu_k.
\]

Because the dynamics is deterministic, this polynomial has an especially simple form. Let \(r\) be the first time \(T^r(n)\) leaves \(V_k\). Then

\[
\boxed{
E_{n,k}(z)=
\begin{cases}
z^r,&T^r(n)>k,\\
0,&T^r(n)\in\{1,2\}.
\end{cases}
}
\tag{12}
\]

It is itself a determinant coefficient. By the matrix determinant lemma and \(\det(I-zC_k)=1\),

\[
\boxed{
E_{n,k}(z)
=
z[\lambda]\,
\det\!\left(
I-zC_k+\lambda u_ke_n^{\mathsf T}
\right).
}
\tag{13}
\]

Consequently, the full Collatz conjecture is exactly equivalent to the conjunction

\[
\boxed{
\begin{aligned}
&\det M_k(x)=1-x^2
&&\text{for every }k,\\
&\forall n\ge3\ \exists k\ge n:
E_{n,k}(z)=0.
\end{aligned}
}
\tag{14}
\]

The first line excludes cycles. The second says that each orbit reaches \(\{1,2\}\) before escaping some sufficiently large finite interval.

But (12) also shows why another sign-reversing involution will not prove the second line: when \(E_{n,k}\ne0\), it consists of a \textbf{single positive monomial}. There are no competing path terms to cancel.

Moreover, analytic convergence is insufficient. For an unbounded orbit, the first-exit exponent \(r=r(k)\) tends to infinity, so

\[
E_{n,k}(z)=z^{r(k)}\longrightarrow0
\qquad (|z|<1),
\]

even though \(E_{n,k}\) is never eventually zero. What is needed is exact eventual vanishing, not a determinant limit or analytic continuation statement.


\subsection*{Final assessment}


Approach 63 can be completed, but its completed conclusion is a no-go theorem rather than a Collatz proof:

\[
\boxed{
\det M_k(x)
=
(1-x^2)
\prod_{\substack{C\ne(1\,2)\\ \max C\le k}}
\left(1-(-x)^{|C|}\right).
}
\]

The determinant permutations do not possess hidden cancellations. They are already in bijection with subsets of genuine cycles. A sign-reversing involution eliminating nontrivial-cycle terms would therefore be either impossible\textemdash{}if a nontrivial cycle exists\textemdash{}or empty\textemdash{}if none exists.

The remaining full-conjecture statement is the exact nonescape condition in (14). Neither the original determinants nor any trace/Fredholm determinant of the transition operator can establish it, because infinite acyclic trajectories make no closed-walk contribution.
\subsection*{References}
\begin{enumerate}[leftmargin=2.2em]
\item \hypertarget{ref:R68:1}{}\hyperlink{cit:R68:1}{\texttt{turn305709view2}}
\item \hypertarget{ref:R68:2}{}\hyperlink{cit:R68:2}{\texttt{turn779723view1}}
\item \hypertarget{ref:R68:3}{}\hyperlink{cit:R68:3}{\texttt{turn779723view2}}
\item \hypertarget{ref:R68:4}{}\hyperlink{cit:R68:4}{\texttt{turn472150view0}}
\end{enumerate}

\clearpage
\section{Round 69: approach 64 --- Uniform nilpotency after removing the trivial basin}
\subsection*{Prompt}
\begin{quote}
Complete a full proof of the Collatz conjecture using the following approach:

\begin{enumerate}
\item Uniform nilpotency after removing the trivial basin
\end{enumerate}

\textbf{Speculative \textperiodcentered{} P}

Truncate the Collatz adjacency operator at height N, remove vertices already connected to 1, and study the remainder. Prove a nilpotency bound independent of N, together with control of edges crossing the truncation boundary.

Take as long as you need.
\end{quote}
\subsection*{Response}

\subsubsection*{Verdict}


The proposed central lemma is false under the natural truncation. This is not merely an unresolved technical gap: for every \(L\), one can explicitly construct a cutoff whose residual adjacency operator contains a Jordan chain of length at least \(L\). Thus no nilpotency exponent independent of \(N\) can exist.

As of August 16, 2026, current references still classify Collatz as unsolved. Tao's strongest general theorem is an almost-everywhere result in logarithmic density, not a pointwise convergence theorem. [\hypertarget{cit:R69:1}{}\hyperlink{ref:R69:1}{1}, \hypertarget{cit:R69:2}{}\hyperlink{ref:R69:2}{2}]

What follows is a complete disproof of Approach 64 as stated, followed by the strongest repair I can derive: a quantitative logarithmic-nilpotency criterion that would prove Collatz, an unconditional lower bound showing how sharp such a criterion must be, and an adaptive-block version that avoids the obstruction.


\bigskip



\subsubsection*{1. The truncated residual operator}


Use the shortcut Collatz map, matching the trivial two-cycle \(1\leftrightarrow2\):

\[
T(n)=
\begin{cases}
n/2,&n\equiv0\pmod 2,\\[2mm]
(3n+1)/2,&n\equiv1\pmod 2.
\end{cases}
\]

For \(N\ge2\), let

\[
B_N=
\left\{
n\le N:
T^r(n)=1\text{ for some }r\ge0
\text{ and }T^j(n)\le N\text{ for }0\le j\le r
\right\}
\]

be the basin of \(1\) visible entirely within the truncated graph \([1,N]\). Set

\[
R_N=\{1,\ldots,N\}\setminus B_N.
\]

On the vector space with basis \(\{e_n:n\in R_N\}\), define the compressed forward adjacency operator

\[
A_Ne_n=
\begin{cases}
e_{T(n)},&T(n)\in R_N,\\
0,&T(n)\notin R_N.
\end{cases}
\]

Transposing this convention gives the usual adjacency orientation and does not change nilpotency.

For every \(k\),

\[
A_N^ke_n=e_{T^k(n)}
\]

provided the first \(k\) iterates remain in \(R_N\); otherwise \(A_N^ke_n=0\).

Write

\[
\nu(N)=\inf\{m\ge1:A_N^m=0\},
\]

with \(\nu(N)=\infty\) when \(A_N\) is not nilpotent.

Because \(R_N\) is finite and every vertex has at most one forward edge, we immediately have:


\paragraph*{Lemma 1}


\(A_N\) is nilpotent if and only if the residual truncated graph contains no directed cycle.

Indeed, a directed cycle supplies vectors surviving every power of \(A_N\). Conversely, a finite directed graph without a cycle has a finite maximal directed-path length.

Consequently, proving merely that every individual \(A_N\) is nilpotent would exclude nontrivial finite Collatz cycles, but it would not exclude an unbounded trajectory. An unbounded trajectory can appear through longer and longer finite nilpotent chains.

The proposed \textbf{uniform} nilpotency was intended to close this gap. Unfortunately, uniformity is impossible.


\bigskip



\subsubsection*{2. Explicit obstruction to uniform nilpotency}



\paragraph*{Theorem 2 \textemdash{} Arbitrarily long residual chains}


For every integer \(L\ge2\), there is a cutoff

\[
N_L=2\cdot3^{L-1}-1
\]

such that

\[
A_{N_L}^{\,L-1}\neq0.
\]

In particular,

\[
\nu(N_L)\ge L,
\qquad
\sup_N\nu(N)=\infty.
\]


\paragraph*{Proof}


For \(0\le j\le L\), define

\[
x_j=3^j2^{L-j}-1.
\]

For \(j<L\), \(x_j\) is odd. Therefore

\[
\begin{aligned}
T(x_j)
&=\frac{3(3^j2^{L-j}-1)+1}{2}\\
&=3^{j+1}2^{L-j-1}-1\\
&=x_{j+1}.
\end{aligned}
\]

Moreover,

\[
x_{j+1}-x_j
=
3^j2^{L-j-1}>0,
\]

so the sequence is strictly increasing:

\[
x_0<x_1<\cdots<x_L.
\]

Its penultimate and final elements are

\[
x_{L-1}=2\cdot3^{L-1}-1=N_L,
\qquad
x_L=3^L-1>N_L.
\]

Thus

\[
x_0\longmapsto x_1\longmapsto\cdots
\longmapsto x_{L-1}\longmapsto x_L
\]

is a Collatz trajectory segment whose first \(L\) vertices lie in \([1,N_L]\), while the next vertex exits the truncation.

None of \(x_0,\ldots,x_{L-1}\) belongs to \(B_{N_L}\): its unique forward path increases through the displayed vertices and exits \([1,N_L]\) before reaching \(1\). Hence all these vertices belong to \(R_{N_L}\), and

\[
A_{N_L}^{\,L-1}e_{x_0}=e_{x_{L-1}}\neq0.
\]

Therefore \(\nu(N_L)\ge L\). \ensuremath{\blacksquare}

Since

\[
L=1+\log_3\left(\frac{N_L+1}{2}\right),
\]

this already gives

\[
\nu(N_L)
\ge
\frac{\log_2N_L}{\log_2 3}+O(1)
=
0.630929\ldots\,\log_2N_L+O(1).
\]


\paragraph*{Example}


Taking \(L=5\) gives

\[
31\longmapsto47\longmapsto71
\longmapsto107\longmapsto161
\longmapsto242.
\]

For the cutoff \(N=161\), the first five vertices are residual, and the fifth step crosses the boundary. Hence

\[
A_{161}^{4}e_{31}=e_{161}\neq0,
\]

so \(\nu(161)\ge5\).

This family alone rules out every \(N\)-independent nilpotency exponent.


\bigskip



\subsubsection*{3. Boundary control does not make the obstruction perturbative}


Let the boundary-exit operator be

\[
E_Ne_n=
\begin{cases}
e_{T(n)},&T(n)>N,\\
0,&T(n)\le N.
\end{cases}
\]

For the preceding family,

\[
E_{N_L}A_{N_L}^{\,L-1}e_{x_0}=e_{x_L}.
\]

Thus, in every ordinary \(\ell^p\) basis norm,

\[
\left\|E_{N_L}A_{N_L}^{\,L-1}\right\|\ge1.
\]

The long residual chain does not become small as it approaches the truncation boundary. It terminates in a full-strength boundary crossing.

There is also large transient amplification. Define the finite partial Green operator

\[
G_{N,L}=\sum_{j=0}^{L-1}A_N^j.
\]

Then

\[
G_{N_L,L}e_{x_0}
=
\sum_{j=0}^{L-1}e_{x_j},
\]

and the \(x_j\) are distinct, so in \(\ell^2\),

\[
\|G_{N_L,L}\|
\ge
\left\|G_{N_L,L}e_{x_0}\right\|_2
=
\sqrt L.
\]

Thus even if every finite truncation were individually nilpotent, the family would possess arbitrarily long Jordan-type transients. Eigenvalues alone would see only zero, while the direct limit could still behave like a unilateral shift carrying mass toward infinity.

Any successful boundary argument must therefore follow \textbf{iterated boundary exits across increasing scales}. It cannot treat the boundary as an error whose norm tends uniformly to zero.


\bigskip



\subsubsection*{4. The other interpretation of ``remove the basin'' is equivalent to Collatz}


Suppose instead that ``already connected to \(1\)'' means the actual global basin

\[
B_\infty=\{n:T^r(n)=1\text{ for some }r\ge0\},
\]

and let

\[
R_\infty=\mathbb N\setminus B_\infty.
\]

Compress \(T\) to \(R_\infty\cap[1,N]\), obtaining operators \(\widetilde A_N\).

Then uniform nilpotency of these operators is exactly equivalent to the Collatz conjecture.


\paragraph*{Proposition 3}


The following are equivalent:

\begin{enumerate}
\item Every positive integer reaches \(1\).
\item There is a constant \(K\) such that \[
   \widetilde A_N^K=0
   \qquad\text{for every }N.
   \]
\end{enumerate}


\paragraph*{Proof}


If Collatz holds, then \(R_\infty=\varnothing\), so every \(\widetilde A_N\) is the zero operator.

Conversely, suppose \(n\in R_\infty\).

Because \(R_\infty\) is forward invariant, the complete trajectory of \(n\) lies in \(R_\infty\).

If the trajectory is eventually periodic, there is a nontrivial cycle. Choosing \(N\) large enough to contain it makes \(\widetilde A_N\) nonnilpotent.

If the trajectory is not eventually periodic, its elements are distinct. For every \(k\), let

\[
N_k=\max_{0\le j\le k}T^j(n).
\]

Then all the first \(k+1\) trajectory vertices belong to \(R_\infty\cap[1,N_k]\), so

\[
\widetilde A_{N_k}^ke_n=e_{T^k(n)}\neq0.
\]

No uniform exponent exists. \ensuremath{\blacksquare}

Therefore:

\begin{itemize}
\item Using the \textbf{local truncated basin}, uniform nilpotency is explicitly false.
\item Using the \textbf{true global basin}, uniform nilpotency simply repackages the entire conjecture into the definition of the removed set.
\end{itemize}

Neither interpretation supplies the proposed proof.


\bigskip



\subsubsection*{5. A viable quantitative replacement}


Although constant nilpotency is impossible, a sufficiently strong logarithmic upper bound would prove Collatz.


\paragraph*{Theorem 4 \textemdash{} Logarithmic nilpotency criterion}


Suppose there are constants \(c,C\) such that, for every sufficiently large \(N\),

\[
\nu(N)\le c\log_2N+C,
\]

where

\[
c<\frac1{\log_2(3/2)}
=
1.709511\ldots.
\]

Then every positive integer reaches \(1\).


\paragraph*{Proof}


First, the assumed finite bound excludes every nontrivial cycle, since a cycle would give \(\nu(N)=\infty\) for sufficiently large \(N\).

Suppose nevertheless that \(n\) has an unbounded trajectory. Write

\[
x_j=T^j(n),
\qquad
M_k=\max_{0\le j\le k}x_j.
\]

Since the trajectory never reaches \(1\), every \(x_j\) lies in the residual set \(R_{M_k}\). All first \(k\) edges stay below \(M_k\), and therefore

\[
A_{M_k}^ke_n=e_{x_k}\neq0.
\]

Hence

\[
\nu(M_k)\ge k+1.
\]

On the other hand,

\[
T(x)+1\le\frac32(x+1)
\]

for every positive integer \(x\). Equality holds when \(x\) is odd; the even case is strictly smaller. Consequently,

\[
x_j+1\le\left(\frac32\right)^j(n+1),
\]

and therefore

\[
M_k+1\le\left(\frac32\right)^k(n+1).
\]

Thus

\[
\log_2M_k
\le
k\log_2(3/2)+O_n(1).
\]

Combining this with the assumed nilpotency estimate gives

\[
\begin{aligned}
k+1
&\le \nu(M_k)\\
&\le c\log_2M_k+C\\
&\le c\,k\log_2(3/2)+O_n(1).
\end{aligned}
\]

But

\[
c\log_2(3/2)<1,
\]

so the last inequality is impossible for sufficiently large \(k\). Hence no unbounded trajectory exists. Nontrivial bounded trajectories were already excluded because they would be cycles. Therefore every trajectory reaches \(1\). \ensuremath{\blacksquare}

This is a genuine sufficient theorem. It converts Approach 64 into a quantitative target:

\[
\boxed{\nu(N)<(1.709511\ldots-\varepsilon)\log_2N+O(1).}
\]

Unfortunately, the required bound remains unproved.


\bigskip



\subsubsection*{6. The obstruction is actually at least \(\log_2N\)}


The simple all-odd construction only forced coefficient \(1/\log_2 3\). A more balanced parity construction gives a stronger unconditional lower bound.


\paragraph*{Theorem 5}


For the locally truncated residual operators,

\[
\boxed{
\limsup_{N\to\infty}
\frac{\nu(N)}{\log_2N}\ge1.
}
\]


\paragraph*{Proof}


Let

\[
\alpha=\log_3 2.
\]

For an integer \(m\ge2\), define

\[
s_j=\lceil \alpha j\rceil,
\qquad 0\le j\le m,
\]

and define parity bits

\[
\varepsilon_j=s_{j+1}-s_j\in\{0,1\},
\qquad 0\le j<m.
\]

The number of odd steps among the first \(j\) positions is \(s_j\). Hence the multiplicative factor after \(j\) steps is

\[
q_j=\frac{3^{s_j}}{2^j}.
\]

From

\[
\alpha j\le s_j<\alpha j+1
\]

we obtain

\[
1\le q_j<3.
\]

For any integer \(r\) realizing this parity prefix, the exact affine iteration formula gives

\[
T^j(r)
=
q_j\left(
r+
\sum_{\substack{0\le i<j\\ \varepsilon_i=1}}
\frac{2^i}{3^{s_{i+1}}}
\right).
\]

Because

\[
s_{i+1}\ge\alpha(i+1),
\]

we have

\[
3^{s_{i+1}}\ge 2^{i+1},
\qquad
\frac{2^i}{3^{s_{i+1}}}\le\frac12.
\]

Therefore, for \(0\le j\le m\),

\[
r\le T^j(r)<3\left(r+\frac j2\right).
\]

So the first \(m\) iterates remain within an additive \(O(m)\) and multiplicative factor \(3\) of the seed.

Now append

\[
d=1+
\left\lceil
\log_{3/2}\bigl(3(m+1)\bigr)
\right\rceil
\]

consecutive odd parity bits, and let

\[
K=m+d.
\]

Every binary parity word of length \(K\) is realized by exactly one residue class modulo \(2^K\). To see this inductively, a parity prefix of length \(K\) determines one class modulo \(2^K\); its two lifts modulo \(2^{K+1}\) have \(K\)-th iterates differing by an odd number, so exactly one realizes either chosen next parity.

Let \(r\in[0,2^K)\) realize the resulting length-\(K\) word. The word begins with two odd bits, whereas the trajectory of \(1\) begins odd-even, so \(r\ne1\); in particular \(r\ge3\).

Write \(x_j=T^j(r)\). During the appended all-odd tail,

\[
x_{m+t}+1
=
\left(\frac32\right)^t(x_m+1).
\]

At the penultimate tail position,

\[
\begin{aligned}
x_{K-1}+1
&=
\left(\frac32\right)^{d-1}(x_m+1)\\
&\ge
3(m+1)(r+1).
\end{aligned}
\]

This exceeds every earlier value, since those values are less than

\[
3\left(r+\frac m2\right).
\]

Set

\[
N=x_{K-1}.
\]

Then

\[
x_0,x_1,\ldots,x_{K-1}\le N,
\]

while the final prescribed odd step gives

\[
x_K>x_{K-1}=N.
\]

Thus these \(K\) vertices form a residual path inside the cutoff, and

\[
\nu(N)\ge K.
\]

For the upper size estimate, minimality in the definition of \(d\) gives

\[
\left(\frac32\right)^{d-1}
<
\frac92(m+1).
\]

Also,

\[
x_m+1<3\left(r+\frac m2\right)+1.
\]

Since \(r<2^K\), it follows that

\[
N\le C(m+1)2^K
\]

for an absolute constant \(C\). Therefore

\[
\log_2N\le K+O(\log m).
\]

Finally,

\[
K=m+O(\log m),
\]

so

\[
\frac{\nu(N)}{\log_2N}
\ge
\frac{K}{K+O(\log m)}
\longrightarrow1.
\]

This proves the theorem. \ensuremath{\blacksquare}

Consequently, any estimate of the form

\[
\nu(N)\le c\log_2N+O(1)
\]

must have \(c\ge1\). Combining Theorems 4 and 5, the potentially useful coefficient window is

\[
\boxed{
1\le c<\frac1{\log_2(3/2)}
=
1.709511\ldots.
}
\]

That is a substantially sharper formulation of the operator problem than uniform nilpotency.


\bigskip



\subsubsection*{7. The natural repair: adaptive dyadic blocking}


The explicit long chains arise because a number with \(L\) low odd parity bits can remain inside the raw adjacency graph for \(L\) consecutive steps. A natural way to remove this artifact is to group together a number of steps comparable to the current bit length.

Let

\[
\ell(n)=\lfloor\log_2n\rfloor+1
\]

and define the jump map

\[
J(n)=T^{\ell(n)}(n).
\]

For the obstruction seed \(2^L-1\), all its first \(L\) steps are odd, so

\[
J(2^L-1)=3^L-1.
\]

Thus the entire length-\(L\) raw Jordan chain is compressed into a single jump.

The relevant uniform statement becomes:


\paragraph*{Adaptive extinction conjecture}


There exists an integer \(B\) such that for every \(n>1\),

\[
\min_{1\le j\le B}J^j(n)<n.
\]

This statement would prove Collatz immediately by strong induction: once the jump orbit falls below \(n\), the lower integer reaches \(1\) by induction, and each jump is a finite collection of ordinary Collatz steps.

This exact bit-length jump viewpoint has been studied by Eliahou, Fromentin, and Simonetto. They present computational and heuristic evidence and conjecture that, for \(n\ge2^{500}\), at most four such jumps are needed to fall below \(n\); their result is evidence and a conjecture, not a proof. [\hypertarget{cit:R69:3}{}\hyperlink{ref:R69:3}{3}]

It also has a clean operator formulation. On basis vectors indexed by pairs \((a,x)\) with \(x\ge a>1\), define

\[
\mathcal Ke_{a,x}
=
\begin{cases}
e_{a,J(x)},&J(x)\ge a,\\
0,&J(x)<a.
\end{cases}
\]

Then the adaptive extinction conjecture is precisely the diagonal uniform-nilpotency statement

\[
\mathcal K^Be_{n,n}=0
\qquad\text{for all }n>1.
\]

Notice that this is not global nilpotency of \(\mathcal K\), which would be unnecessarily strong. It is uniform extinction on the diagonal starting states.


\bigskip



\subsubsection*{8. The exact arithmetic lemma still needed}


For a binary parity word

\[
w=(\varepsilon_0,\ldots,\varepsilon_{k-1}),
\]

put

\[
s_j=\sum_{i<j}\varepsilon_i,
\qquad s=s_k.
\]

The exact endpoint formula is

\[
T^k(n)
=
\frac{3^sn+C(w)}{2^k},
\]

where

\[
C(w)
=
\sum_{j=0}^{k-1}
\varepsilon_j\,2^j\,3^{\,s-s_{j+1}}.
\]

The word \(w\) is realized by one residue class modulo \(2^k\).

A sequence of \(B\) non-descending jumps would therefore give integers

\[
x_0=n,\quad x_1,\ldots,x_B
\]

and parity blocks \(w_0,\ldots,w_{B-1}\) satisfying

\[
|w_i|=\ell(x_i),
\]

\[
x_{i+1}
=
\frac{3^{s_i}x_i+C(w_i)}{2^{|w_i|}},
\]

and

\[
x_i\ge x_0
\qquad(1\le i\le B).
\]

Thus the precise missing theorem is:


\paragraph*{Adaptive block-extinction lemma}


There is an absolute \(B\) for which no positive integer \(x_0>1\) and no admissible sequence of parity blocks \(w_0,\ldots,w_{B-1}\) can simultaneously satisfy

\[
|w_i|=\ell(x_i),\qquad
x_{i+1}
=
\frac{3^{s_i}x_i+C(w_i)}{2^{|w_i|}},
\qquad
x_i\ge x_0
\]

for every \(1\le i\le B\).

Proving this lemma would complete the adaptive-nilpotency proof of Collatz. The hard point is not average multiplicative drift. It is showing that the unique residue classes realizing several adaptively chosen, nearly critical parity blocks cannot remain compatible with the changing top binary digits of their endpoints.


\bigskip



\subsubsection*{Conclusion}


Approach 64, literally interpreted, cannot yield a proof:

\[
\boxed{\text{Uniform nilpotency of the raw truncated operators is false.}}
\]

The explicit family

\[
3^j2^{L-j}-1
\]

produces arbitrarily long residual chains and full-strength boundary exits.

The strongest repaired conclusions are:

\[
\limsup_{N\to\infty}
\frac{\nu(N)}{\log_2N}\ge1,
\]

while

\[
\nu(N)
\le
c\log_2N+O(1),
\qquad
c<1.709511\ldots,
\]

would prove Collatz. Hence the meaningful raw-operator target lies in the narrow coefficient interval

\[
1\le c<1.709511\ldots.
\]

The more structurally promising replacement is uniform \textbf{adaptive} nilpotency after grouping \(\ell(n)\) steps at a time. That formulation survives the explicit no-go family and would prove Collatz, but its adaptive block-extinction lemma remains the unresolved arithmetic step. Presenting the original uniform-nilpotency claim as a completed Collatz proof would therefore be incorrect.
\subsection*{References}
\begin{enumerate}[leftmargin=2.2em]
\item \hypertarget{ref:R69:1}{}\hyperlink{cit:R69:1}{\texttt{turn673822view2}}
\item \hypertarget{ref:R69:2}{}\hyperlink{cit:R69:2}{\texttt{turn673822view3}}
\item \hypertarget{ref:R69:3}{}\hyperlink{cit:R69:3}{\texttt{turn673822view0}}
\end{enumerate}

\clearpage
\section{Round 70: approach 65 --- Compute the operator commutant directly}
\subsection*{Prompt}
\begin{quote}
Complete a full proof of the Collatz conjecture using the following approach:

\begin{enumerate}
\item Compute the operator commutant directly
\end{enumerate}

\textbf{Established/Hybrid \textperiodcentered{} P}

Mori's equivalence reduces the conjecture to C*(T\_1, T\_2)' = C I. Try to solve the matrix-coefficient equations for an operator commuting with both generators and show every solution is scalar, using only their sparse arithmetic structure.

Take as long as you need.
\end{quote}
\subsection*{Response}

\subsubsection*{Result of the direct computation}


The commutant can be computed completely. The computation eliminates \textbf{all off-diagonal operators}, but the remaining diagonal equations are precisely the original Collatz reachability problem.

Thus this approach gives an exact structural theorem,

\[
C^*(T_1,T_2)'
   \cong \ell^\infty(\{\text{Collatz orbit-components}\}),
\]

but it does \textbf{not} prove that the commutant is scalar without separately proving that there is only one orbit-component. Claiming scalarity after the calculation below would therefore be circular.

Mori works on \(H=\ell^2(\mathbb N)\), with standard basis \((e_n)_{n\ge1}\), and defines

\[
T_1e_n=
\begin{cases}
e_{3n+1},&n\text{ odd},\\
0,&n\text{ even},
\end{cases}
\qquad
T_2e_n=
\begin{cases}
0,&n\text{ odd},\\
e_{n/2},&n\text{ even}.
\end{cases}
\]

He proves that irreducibility of the generated \(C^*\)-algebra, equivalently \(C^*(T_1,T_2)'=\mathbb CI\), is equivalent to the Collatz conjecture. [\hypertarget{cit:R70:1}{}\hyperlink{ref:R70:1}{1}]

Let

\[
F(n)=
\begin{cases}
3n+1,& n\text{ odd},\\
n/2,&n\text{ even}.
\end{cases}
\]


\bigskip



\subsection*{1. The elementary operator relations}


The adjoints are

\[
T_2^*e_n=e_{2n},
\]

and

\[
T_1^*e_m=
\begin{cases}
e_{(m-1)/3},&m\equiv4\pmod 6,\\
0,&m\not\equiv4\pmod6.
\end{cases}
\]

Consequently,

\[
T_1^*T_1=P_{\mathrm{odd}},
\qquad
T_2^*T_2=P_{\mathrm{even}},
\]

while

\[
T_1T_1^*=P_{\{m:m\equiv4\pmod6\}},
\qquad
T_2T_2^*=I.
\]

Suppose

\[
A\in C^*(T_1,T_2)',
\qquad
a_{m,n}:=\langle Ae_n,e_m\rangle.
\]

Since \(A\) commutes with \(P_{\mathrm{odd}}\) and \(P_{\mathrm{even}}\),

\[
a_{m,n}=0
\quad\text{whenever }m\not\equiv n\pmod2.
\tag{1}
\]

Commutation with \(T_2^*\) gives

\[
A e_{2n}=T_2^*Ae_n,
\]

hence

\[
a_{2m,2n}=a_{m,n}.
\tag{2}
\]

In particular, iterating (1)\textendash{}(2),

\[
a_{m,n}\ne0
\quad\Longrightarrow\quad
v_2(m)=v_2(n),
\tag{3}
\]

and, after removing the common power of \(2\),

\[
a_{m,n}
=
a_{m/2^r,n/2^r},
\qquad r=v_2(m)=v_2(n).
\tag{4}
\]

For odd \(m,n\), commutation with \(T_1\) gives

\[
a_{3m+1,\,3n+1}=a_{m,n}.
\tag{5}
\]

There are also the corresponding range-block vanishing conditions; for example,

\[
a_{m,n}=0
\]

whenever exactly one of \(m,n\) is congruent to \(4\bmod6\).

Equations (1)\textendash{}(5) are the first sparse matrix-coefficient equations. To solve all of them simultaneously, it is useful to pass to finite parity cylinders.


\bigskip



\subsection*{2. Every finite parity cylinder is a commutant projection}


For a word

\[
\alpha=(i_0,\ldots,i_{k-1})\in\{1,2\}^k,
\]

put

\[
V_\alpha=T_{i_{k-1}}\cdots T_{i_1}T_{i_0}.
\]

On a basis vector \(e_n\), the operator \(V_\alpha\) is nonzero exactly when the first \(k\) Collatz branches of \(n\) are the prescribed branches:

\[
i_j=
\begin{cases}
1,&F^j(n)\text{ odd},\\
2,&F^j(n)\text{ even}.
\end{cases}
\]

In that case,

\[
V_\alpha e_n=e_{F^k(n)}.
\]

Thus \(V_\alpha\) is a partial permutation, and

\[
P_\alpha:=V_\alpha^*V_\alpha
\]

is the diagonal projection onto the set

\[
E_\alpha
 =
 \left\{
 n:
 \text{the first \(k\) branches of \(n\) equal }\alpha
 \right\}.
\]

Because \(A\) commutes with the whole generated \(C^*\)-algebra,

\[
AP_\alpha=P_\alpha A.
\tag{6}
\]

Suppose \(m\in E_\alpha\) and \(n\notin E_\alpha\). Then

\[
\begin{aligned}
a_{m,n}
 &=\langle Ae_n,e_m\rangle\\
 &=\langle P_\alpha Ae_n,e_m\rangle\\
 &=\langle AP_\alpha e_n,e_m\rangle\\
 &=0.
\end{aligned}
\]

Therefore

\[
a_{m,n}\ne0
\quad\Longrightarrow\quad
m,n\text{ have the same first \(k\) parity decisions}
\tag{7}
\]

for every \(k\).

So a nonzero off-diagonal coefficient could exist only between two distinct positive integers with exactly the same infinite Collatz parity itinerary.


\bigskip



\subsection*{3. Two positive integers cannot have the same infinite itinerary}


Assume \(m,n\) have the same first \(k\) parity decisions. Let

\[
o_k=\#\{0\le j<k:F^j(m)\text{ is odd}\},
\]

and

\[
e_k=\#\{0\le j<k:F^j(m)\text{ is even}\}.
\]

These counts are the same for \(m\) and \(n\). Since the additive \(+1\) terms cancel when the two trajectories follow the same branches,

\[
F^k(m)-F^k(n)
 =
 \frac{3^{o_k}}{2^{e_k}}(m-n).
\tag{8}
\]

The left side is an integer. Since \(3^{o_k}\) is odd, (8) implies

\[
2^{e_k}\mid m-n.
\tag{9}
\]

But an odd Collatz value is always followed by an even value. Hence there cannot be two consecutive odd steps, and consequently

\[
e_k\ge \left\lfloor\frac{k}{2}\right\rfloor.
\tag{10}
\]

Letting \(k\to\infty\), equation (9) says that \(m-n\) is divisible by arbitrarily large powers of \(2\). Therefore

\[
m-n=0.
\]

We have proved:

\[
\boxed{
\text{The infinite parity-itinerary map }
n\longmapsto
\bigl(F^j(n)\bmod2\bigr)_{j\ge0}
\text{ is injective on }\mathbb N.
}
\tag{11}
\]

Mori's later ``total uniqueness'' framework likewise identifies itinerary injectivity as the mechanism that turns reducing subspaces into invariant subsets; for the affine Collatz branches he proves this uniqueness property directly. [\hypertarget{cit:R70:2}{}\hyperlink{ref:R70:2}{2}, \hypertarget{cit:R70:3}{}\hyperlink{ref:R70:3}{3}]

Combining (7) with (11),

\[
m\ne n\quad\Longrightarrow\quad a_{m,n}=0.
\]

Thus:

\[
\boxed{\text{Every operator in }C^*(T_1,T_2)'\text{ is diagonal.}}
\tag{12}
\]

This completely solves the off-diagonal portion of the commutant problem.


\bigskip



\subsection*{4. The remaining diagonal equations}


Write

\[
Ae_n=\lambda_n e_n,
\qquad
(\lambda_n)\in\ell^\infty(\mathbb N).
\]

From \(AT_2=T_2A\), for every even \(n\),

\[
\lambda_{n/2}=\lambda_n.
\tag{13}
\]

From \(AT_1=T_1A\), for every odd \(n\),

\[
\lambda_{3n+1}=\lambda_n.
\tag{14}
\]

Together,

\[
\boxed{\lambda_{F(n)}=\lambda_n\qquad(n\ge1).}
\tag{15}
\]

Conversely, suppose \(\lambda\in\ell^\infty(\mathbb N)\) satisfies (15), and define

\[
M_\lambda e_n=\lambda_n e_n.
\]

Then for even \(n\),

\[
M_\lambda T_2e_n
  =\lambda_{n/2}e_{n/2}
  =\lambda_ne_{n/2}
  =T_2M_\lambda e_n,
\]

and for odd \(n\),

\[
M_\lambda T_1e_n
 =\lambda_{3n+1}e_{3n+1}
 =\lambda_ne_{3n+1}
 =T_1M_\lambda e_n.
\]

The same relation along the reverse edges shows directly that \(M_\lambda\) commutes with \(T_1^*\) and \(T_2^*\) as well. Therefore \(M_\lambda\in C^*(T_1,T_2)'\).

This proves the exact classification.


\subsubsection*{Exact commutant theorem}


Define \(m\sim n\) when their forward trajectories eventually meet:

\[
m\sim n
\quad\Longleftrightarrow\quad
\exists r,s\ge0\quad F^r(m)=F^s(n).
\tag{16}
\]

Then

\[
\boxed{
C^*(T_1,T_2)'
 =
 \left\{
 M_\lambda:
 \lambda\in\ell^\infty(\mathbb N),\
 \lambda_{F(n)}=\lambda_n
 \right\}
}
\tag{17}
\]

and hence

\[
\boxed{
C^*(T_1,T_2)'
 \cong
 \ell^\infty(\mathbb N/{\sim}).
}
\tag{18}
\]

In words: \textbf{the commutant consists exactly of bounded functions on the set of Collatz orbit-components.}


\bigskip



\subsection*{5. Scalarity is exactly the Collatz conjecture}


A diagonal invariant sequence \(\lambda\) is constant precisely on each \(\sim\)-component. Therefore

\[
C^*(T_1,T_2)'=\mathbb CI
\]

if and only if \(\mathbb N/{\sim}\) has exactly one element.

The forward orbit of \(1\) is

\[
1\longmapsto4\longmapsto2\longmapsto1.
\]

Thus an integer belongs to the component of \(1\) exactly when its forward orbit reaches \(1\), \(2\), or \(4\), equivalently when it eventually reaches \(1\). Hence

\[
\mathbb N/{\sim}\text{ has one component}
\quad\Longleftrightarrow\quad
\text{every positive integer reaches }1.
\tag{19}
\]

So the direct calculation reproduces Mori's equivalence in the strongest possible coefficient-level form:

\[
\boxed{
C^*(T_1,T_2)'=\mathbb CI
\quad\Longleftrightarrow\quad
\text{Collatz}.
}
\tag{20}
\]

But it does not make the right-hand side easier: it identifies it exactly with the surviving diagonal sector.


\bigskip



\subsection*{6. The explicit obstruction to finishing the argument}


Suppose there is an orbit-component \(C\) that does not contain \(1\). Define

\[
P_Ce_n=
\begin{cases}
e_n,&n\in C,\\
0,&n\notin C.
\end{cases}
\]

Because \(C\) is closed under both forward and reverse Collatz edges,

\[
P_CT_i=T_iP_C,
\qquad
P_CT_i^*=T_i^*P_C
\]

for \(i=1,2\). Thus

\[
P_C\in C^*(T_1,T_2)'.
\]

It is a nonzero, nonidentity projection. Every one of the sparse matrix equations is satisfied automatically.

Conversely, every nonconstant commutant element distinguishes at least two such components.

Therefore an attempted final argument of the form


\begin{quote}
``The sparse recurrence equations force all \(\lambda_n\) to equal \(\lambda_1\)''

\end{quote}


must establish that every \(n\) is connected to \(1\). That is already the Collatz conjecture.

More sharply, for every fixed \(n\),

\[
n\sim1
\quad\Longleftrightarrow\quad
\lambda_n=\lambda_1
\text{ for every }M_\lambda\in C^*(T_1,T_2)'.
\tag{21}
\]

Indeed, if \(n\not\sim1\), the component projection \(P_{[n]}\) has diagonal value \(1\) at \(n\) and \(0\) at \(1\).

Thus even the pointwise statement ``all commutant operators agree at \(n\) and \(1\)'' is equivalent to proving Collatz for that particular seed \(n\).


\bigskip



\subsection*{7. What the sparse structure proves around \(e_1\)}


There is a useful strengthening of the usual \(e_1\)-anchor.

Set

\[
W=T_2^2T_1.
\]

Then

\[
We_n=
\begin{cases}
e_{(3n+1)/4},&n\equiv1\pmod4,\\
0,&n\not\equiv1\pmod4.
\end{cases}
\tag{22}
\]

In particular,

\[
We_1=e_1.
\]

Moreover,

\[
\operatorname{Fix}(W)=\mathbb Ce_1.
\tag{23}
\]

To see this, write \(x=\sum c_ne_n\) and suppose \(Wx=x\). For \(m>1\), the coefficient of \(e_m\) in \(Wx\) can come from at most one index,

\[
n=\frac{4m-1}{3}.
\]

When this is an integer, \(n>m\), and the fixed-vector equation gives \(c_m=c_n\). Repeating either eventually reaches an index with no predecessor, forcing the coefficient to vanish, or produces infinitely many distinct indices carrying the same nonzero coefficient, contradicting \(x\in\ell^2\). Hence \(c_m=0\) for every \(m>1\), while \(c_1\) is unrestricted.

Consequently, for every \(A\in C^*(T_1,T_2)'\),

\[
W(Ae_1)=A(We_1)=Ae_1,
\]

so

\[
Ae_1\in\mathbb Ce_1.
\tag{24}
\]

The sparse relations therefore anchor every commutant operator at \(e_1\), and the equations propagate the same scalar throughout the entire component of \(1\). They cannot propagate it into a separate component without proving that the separate component does not exist.


\bigskip



\subsection*{8. A direct arithmetic sieve for a hypothetical exceptional component}


The operator words do impose nontrivial restrictions on a hypothetical exceptional component.

Assume \(C\) is a component not containing \(1\), and let

\[
n=\min C.
\]

Then:

First, \(n\) must be odd. Otherwise \(T_2e_n=e_{n/2}\) places the smaller integer \(n/2\) in \(C\).

Second,

\[
n\not\equiv1\pmod4,
\]

because in that case

\[
T_2^2T_1e_n=e_{(3n+1)/4},
\]

and

\[
\frac{3n+1}{4}<n
\qquad(n>1),
\]

again contradicting minimality. Since \(n\) is odd,

\[
n\equiv3\pmod4.
\tag{25}
\]

Third,

\[
n\not\equiv2\pmod3.
\]

Indeed, when \(n\) is odd and \(n\equiv2\pmod3\),

\[
T_1^*T_2^*e_n
 =
 e_{(2n-1)/3},
\]

and

\[
\frac{2n-1}{3}<n.
\]

Combining these restrictions,

\[
\boxed{n\equiv3\text{ or }7\pmod{12}.}
\tag{26}
\]

Longer operator words give increasingly refined residue exclusions. But extinction of all residue classes at all depths would mean that every component contains a descending path to \(1\); this is another exact encoding of Collatz, not a consequence yet obtained from the commutant formalism.


\bigskip



\subsection*{Final assessment}


The direct commutant computation succeeds completely as an operator classification:

\[
\boxed{
C^*(T_1,T_2)'
=
\operatorname{diag}
\left\{
\lambda\in\ell^\infty(\mathbb N):
\lambda\circ F=\lambda
\right\}
\cong
\ell^\infty(\mathbb N/{\sim}).
}
\]

It proves:

\begin{enumerate}
\item Every off-diagonal matrix coefficient vanishes.
\item Every commutant operator is diagonal.
\item Its diagonal is constant exactly on Collatz orbit-components.
\item Every hypothetical exceptional component produces an explicit nontrivial commutant projection.
\item Scalarity is equivalent, component by component and globally, to Collatz convergence.
\end{enumerate}

Therefore this approach does \textbf{not} presently yield a full proof. Its exact terminal obstruction is not an unresolved analytic estimate or operator lemma; it is the existence or nonexistence of an exceptional Collatz component itself. The commutant machinery removes all operator-theoretic ambiguity and leaves the original arithmetic problem intact.
\subsection*{References}
\begin{enumerate}[leftmargin=2.2em]
\item \hypertarget{ref:R70:1}{}\hyperlink{cit:R70:1}{\texttt{turn177415view0}}
\item \hypertarget{ref:R70:2}{}\hyperlink{cit:R70:2}{\texttt{turn659391view1}}
\item \hypertarget{ref:R70:3}{}\hyperlink{cit:R70:3}{\texttt{turn659391view2}}
\end{enumerate}

\clearpage
\section{Round 71: approach 66 --- Exclude peripheral spectrum of a weighted Collatz shift}
\subsection*{Prompt}
\begin{quote}
Complete a full proof of the Collatz conjecture using the following approach:

\begin{enumerate}
\item Exclude peripheral spectrum of a weighted Collatz shift
\end{enumerate}

\textbf{Hybrid \textperiodcentered{} B}

Choose weights so that a cycle produces a unit-modulus eigenvalue and a divergent orbit produces approximate eigenvectors on the unit circle. Prove the nontrivial spectral radius is strictly below one \textemdash{} or detect approximate eigenvectors computationally as counterexample signatures.

Take as long as you need.
\end{quote}
\subsection*{Response}

\subsection*{Verdict}


I do \textbf{not} obtain a full proof. The weighted-shift detector can be constructed rigorously, but the proposed final step\textemdash{}proving a strict spectral-radius gap after removing only the trivial cycle\textemdash{}is false for the natural classes of weights. In fact, unit-circle approximate eigenvectors arise from completely terminating Collatz trajectories.

This is not merely the familiar observation that an operator reformulation may be equivalent to Collatz. It is a concrete no-go theorem for the approximate-spectrum strategy. The strongest repair I found replaces peripheral approximate spectrum by the \textbf{multiplicity of an actual eigenvalue}; that gives an exact and unusually simple spectral characterization, but proving the eigenvalue is simple remains precisely the missing component-connectivity theorem.

Operator formulations already in the literature associate cycles and hypothetical divergent trajectories with special fixed points, while noting that unrelated fixed points may also occur. Mori's operator-algebra formulation makes the conjecture equivalent to the absence of nontrivial reducing subspaces. [\hypertarget{cit:R71:1}{}\hyperlink{ref:R71:1}{1}, \hypertarget{cit:R71:2}{}\hyperlink{ref:R71:2}{2}]


\bigskip



\subsubsection*{1. The weighted Collatz shift}


Use the shortcut map

\[
T(n)=
\begin{cases}
n/2,&n\ \text{even},\\[2mm]
(3n+1)/2,&n\ \text{odd}.
\end{cases}
\]

Its trivial cycle is \(1\leftrightarrow2\).

On \(H=\ell ^2(\mathbb N)\), let

\[
S_we_n=w(n)e_{T(n)} .
\]

Since every \(m\) has at most two preimages,

\[
\|S_w\|^2
=
\sup_{m\ge1}\sum_{T(n)=m}|w(n)|^2,
\]

so any bounded weight defines a bounded operator.

Take

\[
\ell(n)=\log(n+e),
\qquad
w(n)=\frac{\ell(T(n))}{\ell(n)}.
\tag{1}
\]

This is a positive bounded weight satisfying

\[
w(n)\longrightarrow1
\qquad (n\to\infty).
\tag{2}
\]

It is also a multiplicative coboundary:

\[
\prod_{j=0}^{k-1}w(T^j n)
=
\frac{\ell(T^k n)}{\ell(n)}.
\tag{3}
\]

This is the natural balance between exact cycle normalization and asymptotic unit weights.


\bigskip



\subsubsection*{2. Cycles give exact unit-modulus eigenvalues}


Suppose

\[
n_0\longmapsto n_1\longmapsto\cdots
\longmapsto n_{p-1}\longmapsto n_0
\]

is a cycle. By telescoping,

\[
\prod_{j=0}^{p-1}w(n_j)=1.
\tag{4}
\]

The span of \(e_{n_0},\dots,e_{n_{p-1}}\) is invariant under \(S_w\), and the restriction is a cyclic weighted permutation. Its characteristic equation is

\[
\lambda^p=\prod_{j=0}^{p-1}w(n_j)=1.
\]

Consequently every \(p\)-th root of unity is an eigenvalue of that restriction. In particular, any nontrivial Collatz cycle creates unit-modulus point spectrum.

So the cycle-detection half of the proposal works exactly.


\bigskip



\subsubsection*{3. A divergent orbit gives unit-circle Weyl sequences}


Suppose \(n_j=T^j(n_0)\) is nonperiodic and does not reach the trivial cycle. Then the \(n_j\) are distinct. Since they are distinct positive integers,

\[
n_j\longrightarrow\infty.
\]

By (2),

\[
w(n_j)\longrightarrow1.
\]

Fix \(|\lambda|=1\). For integers \(r,L\ge1\), define

\[
v_{r,L}
=
\frac1{\sqrt L}
\sum_{j=0}^{L-1}\lambda^{-j}e_{n_{r+j}}.
\]

A direct calculation gives

\[
\begin{aligned}
\|(S_w-\lambda I)v_{r,L}\|^2
={}&
\frac1L
+\frac{|w(n_{r+L-1})|^2}{L}\\
&+
\frac1L\sum_{j=0}^{L-2}|w(n_{r+j})-1|^2 .
\end{aligned}
\tag{5}
\]

Choose \(r=r(L)\) so far along the orbit that

\[
|w(n_j)-1|\le \frac1L
\qquad(j\ge r).
\]

Then (5) tends to zero as \(L\to\infty\). Thus

\[
\mathbb T\subseteq
\sigma_{\mathrm{ap}}
\left(S_w\big|_{\overline{\operatorname{span}}\{e_{n_j}\}}\right).
\tag{6}
\]

So the divergent-orbit detection half also works.


\bigskip



\subsubsection*{4. The exact exceptional-sector criterion}


Let

\[
B=\{n:T^k(n)\in\{1,2\}\text{ for some }k\ge0\},
\qquad
E=\mathbb N\setminus B.
\]

The set \(E\) is completely invariant:

\[
n\in E \iff T(n)\in E.
\]

Hence \(H_E=\ell ^2(E)\) is reducing for \(S_w\).

If \(E\neq\varnothing\), then an orbit in \(E\) either repeats, producing a nontrivial cycle and a unit eigenvalue, or is injective and divergent, producing (6). Therefore

\[
E\neq\varnothing
\quad\Longrightarrow\quad
r(S_w|_{H_E})\ge1.
\tag{7}
\]

Conversely, if Collatz holds, then \(E=\varnothing\) and the exceptional restriction is the zero operator. We therefore have the exact equivalence

\[
\boxed{
\text{Collatz holds}
\iff
r(S_w|_{\ell ^2(E)})<1.
}
\tag{8}
\]

The problem is that \(E\) is already the unknown exceptional set. Proving that this reducing subspace is zero is precisely the conjecture. Thus (8) is an exact reformulation but not yet a proof.

One might try to avoid this by removing only the known cycle \(\{1,2\}\). That is where a genuine obstruction appears.


\bigskip



\subsection*{5. No-go theorem: terminating trajectories force unit-circle approximate spectrum}


Let \(P\) be the projection onto

\[
\ell ^2(\mathbb N\setminus\{1,2\}),
\qquad
A=PS_wP.
\]

Even if every positive integer terminates, \(A\) has unit-circle approximate spectrum.


\subsubsection*{5.1 Long high-altitude even runs}


The orbit beginning at \(2^{2L}\) has the segment

\[
2^{2L}\to2^{2L-1}\to\cdots\to2^L.
\]

This is a terminating orbit, and all \(L+1\) displayed values tend uniformly to infinity. Hence all corresponding weights in (1) tend uniformly to \(1\). Applying the construction in (5) gives, for every \(|\lambda|=1\),

\[
\|(A-\lambda I)v_L\|\longrightarrow0.
\]

Thus powers of two alone already prohibit a strict global gap.

A designer might try to repair this by assigning a genuinely contracting weight to every even edge. The odd branch produces the same obstruction.


\subsubsection*{5.2 Arbitrarily long terminating all-odd runs}


For every \(L\ge1\), define

\[
q_L=\frac{2^{3^L}+1}{3^{L+1}},
\qquad
u_L=2q_L,
\qquad
N_L=2^Lu_L-1.
\tag{9}
\]

The lifting-the-exponent identity gives

\[
v_3(2^{3^L}+1)
=
v_3(2+1)+v_3(3^L)
=
L+1,
\]

so \(q_L\) is an odd positive integer and \(u_L\) is even.

For \(0\le j\le L\), put

\[
n_j=3^j2^{L-j}u_L-1.
\tag{10}
\]

Every \(n_j\) is odd, and for \(j<L\),

\[
T(n_j)
=
\frac{3n_j+1}{2}
=
3^{j+1}2^{L-j-1}u_L-1
=
n_{j+1}.
\]

The final odd step gives

\[
\begin{aligned}
T(n_L)
&=\frac{3^{L+1}u_L-2}{2}\\
&=3^{L+1}q_L-1\\
&=2^{3^L}.
\end{aligned}
\tag{11}
\]

After that, the orbit reaches \(1\) through \(3^L\) halvings.

Hence \(N_L\) has \(L+1\) consecutive odd shortcut steps, all at values tending uniformly to infinity, and then lands exactly on a power of two.

For example,

\[
3\to5\to8,
\]

and

\[
151\to227\to341\to512
\]

are the first two members.

Let \(M=L+1\) and define

\[
v_L=\frac1{\sqrt M}
\sum_{j=0}^{M-1}\lambda^{-j}e_{n_j}.
\]

Then

\[
\|(A-\lambda I)v_L\|^2
=
\frac{
1+|w(n_L)|^2+
\sum_{j=0}^{L-1}|w(n_j)-1|^2
}{L+1}.
\tag{12}
\]

Since \(\min_j n_j=N_L\to\infty\),

\[
\max_{0\le j\le L}|w(n_j)-1|\longrightarrow0.
\]

Consequently,

\[
\|(A-\lambda I)v_L\|^2
=
\frac{2+o(1)}{L+1}+o(1)
\longrightarrow0.
\tag{13}
\]

Therefore

\[
\boxed{\mathbb T\subseteq\sigma_{\mathrm{ap}}(A)}
\tag{14}
\]

using vectors supported entirely on \textbf{provably terminating orbits}.

The supports can be chosen disjoint, so these are escaping singular sequences rather than an artifact of one fixed finite-dimensional block.

This invalidates the desired implication

\[
\text{unit-circle approximate spectrum}
\Longrightarrow
\text{counterexample orbit}.
\]


\bigskip



\subsubsection*{6. Every finite parity word occurs on a terminating orbit}


The preceding examples are part of a general phenomenon.


\paragraph*{Finite-prefix extension lemma}


Let

\[
\varepsilon=(\varepsilon_0,\dots,\varepsilon_{k-1})
\in\{0,1\}^k
\]

be any shortcut-Collatz parity word. There exists a positive integer \(n\) whose orbit has this parity prefix and subsequently reaches a power of two.


\paragraph*{Proof}


The parity word determines one residue class

\[
n=a+2^kt.
\]

If \(s=\sum_j\varepsilon_j\), affine iteration gives

\[
T^k(a+2^kt)=b+3^st
\tag{15}
\]

for an integer \(b\).

If \(3\nmid b\), use that \(2\) generates the unit group modulo \(3^s\). Choose \(r\) such that

\[
2^r\equiv b\pmod{3^s}
\]

and sufficiently large that \(2^r\ge b\). Then

\[
t=\frac{2^r-b}{3^s}
\]

is a nonnegative integer, and the endpoint in (15) is \(2^r\).

If \(3\mid b\), first choose the parity of \(t\) so that \(b+3^st\) is odd and append one odd step. Its image is congruent to \(2\pmod3\), and the preceding argument applies with modulus \(3^{s+1}\). \ensuremath{\blacksquare}

Thus every finite parity pattern\textemdash{} including every finite prefix of a hypothetical divergent trajectory\textemdash{}has a terminating realization.

This has a decisive spectral consequence:


\begin{quote}
A bounded finite-memory or \(2\)-adically local weighting cannot distinguish an infinite positive-integer orbit from longer and longer terminating shadows merely through finite-block approximate eigenvectors.

\end{quote}


The missing property is \textbf{nestedness for one fixed positive seed}. Ordinary operator norm and ordinary approximate spectrum take a supremum over a different seed at every block length, so they discard exactly that information.


\bigskip



\subsection*{7. A stronger exact detector using an actual eigenvalue}


There is a simple weight for which every orbit-component produces an actual \(1\)-eigenvector.

Fix \(s>\tfrac12\), and define

\[
(K_sx)_n
=
\left(\frac{T(n)}{n}\right)^s x_{T(n)}.
\tag{16}
\]

This is bounded on \(\ell ^2(\mathbb N)\), since

\[
\begin{aligned}
\|K_sx\|^2
&=
\sum_m |x_m|^2
\sum_{T(n)=m}\left(\frac{m}{n}\right)^{2s},\\
\sum_{T(n)=m}\left(\frac{m}{n}\right)^{2s}
&\le
2^{-2s}+2^{2s}.
\end{aligned}
\]

Let \(C\) be any orbit-equivalence component of the Collatz graph, and set

\[
g_C(n)=n^{-s}\mathbf1_C(n).
\tag{17}
\]

Because \(s>\tfrac12\), \(g_C\in\ell ^2(\mathbb N)\). Moreover,

\[
(K_sg_C)_n
=
\left(\frac{T(n)}n\right)^s
T(n)^{-s}\mathbf1_C(T(n))
=
n^{-s}\mathbf1_C(n),
\]

so

\[
K_sg_C=g_C.
\tag{18}
\]

Conversely, if \(K_sx=x\), define

\[
y_n=n^sx_n.
\]

Equation (16) becomes

\[
y_n=y_{T(n)}.
\tag{19}
\]

Therefore \(y\) is constant on every orbit component. Hence

\[
\ker(I-K_s)
=
\left\{
\sum_C c_Cg_C:
\sum_C |c_C|^2
\sum_{n\in C}n^{-2s}<\infty
\right\}.
\tag{20}
\]

In particular, there is exactly one linearly independent \(1\)-eigenvector per Collatz component.

Since

\[
g(n)=n^{-s}
\]

is the sum of the component vectors, we obtain:

\[
\boxed{
\text{Collatz holds}
\iff
\ker(I-K_s)=\mathbb C\, (n^{-s})_{n\ge1}.
}
\tag{21}
\]

Equivalently:

\[
\boxed{
\text{Collatz holds}
\iff
1\text{ is a simple eigenvalue of }K_s.
}
\tag{22}
\]

This detects both possible failure modes:

\begin{itemize}
\item a nontrivial cyclic component contributes another \(g_C\);
\item a divergent component contributes another \(g_C\).
\end{itemize}

The original forward shift adjoint relation is

\[
K_s=S_s^*,
\qquad
S_se_n=\left(\frac{T(n)}n\right)^s e_{T(n)}.
\]

Thus this is genuinely a weighted-Collatz-shift statement.

But (19) exposes the remaining issue exactly: proving simplicity means proving that every component constant is the same, which means proving that the graph has only the component containing \(1\). No spectral estimate has yet supplied that arithmetic connection.


\bigskip



\subsection*{8. An anchored shift that removes transient spectral pollution}


There is also a useful diagnostic that does not refer to the unknown set \(E\).

On

\[
\mathcal H=\ell ^2(\mathbb N\times\mathbb N_0),
\]

let

\[
q_{n,j}
=
\mathbf1_{\{T^i(n)\notin\{1,2\}\text{ for }0\le i\le j\}}.
\]

For \(\epsilon>0\), define

\[
U_\epsilon e_{n,j}
=
e^{-\epsilon n}q_{n,j+1}e_{n,j+1}.
\tag{23}
\]

Every coefficient is computable by a finite orbit calculation.

Let \(\tau(n)\) be the first hitting time of \(\{1,2\}\), possibly infinite, and define

\[
m_k=\min\{n:\tau(n)>k\}.
\]

Then \(U_\epsilon\) is a weighted partial shift with

\[
\|U_\epsilon^k\|
=
e^{-\epsilon k m_k}.
\tag{24}
\]

If Collatz holds, then \(m_k\to\infty\), so

\[
r(U_\epsilon)=0.
\tag{25}
\]

If Collatz fails and \(m_*=\min E\), then \(m_k=m_*\) for all sufficiently large \(k\), giving

\[
r(U_\epsilon)=e^{-\epsilon m_*}>0.
\tag{26}
\]

Therefore

\[
\boxed{
\text{Collatz holds}
\iff
r(U_\epsilon)=0
\iff
U_\epsilon\text{ is compact}.
}
\tag{27}
\]

If a least counterexample \(m_*\) exists, its value is encoded spectrally by

\[
m_*=-\frac1\epsilon\log r(U_\epsilon).
\tag{28}
\]

This construction suppresses approximate modes whose starting seeds escape to infinity. Its cost is that a counterexample produces radius \(e^{-\epsilon m_*}\), not radius \(1\); as \(\epsilon\downarrow0\), this radius tends to \(1\).

Again, (27) is an exact detector rather than a proof: establishing \(r(U_\epsilon)=0\) is equivalent to establishing \(m_k\to\infty\).


\bigskip



\subsection*{9. What a successful spectral argument would still need}


The straightforward gap theorem is blocked for two separate reasons.

First, asymptotically unit weights produce Weyl sequences on high terminating paths. Second, every finite parity prefix has a terminating extension, so finite-block spectral data do not preserve a common starting seed.

A successful replacement would need a compactness theorem of the following form:


\begin{quote}
Every normalized approximate unit eigenvector in the nontrivial sector has a subsequence whose mass remains anchored to one fixed positive starting seed, rather than escaping through seeds \(n_L\to\infty\) or converging only \(2\)-adically.

\end{quote}


The explicit families above show that this is false in ordinary \(\ell ^2(\mathbb N)\). Obtaining it would require an additional nonlocal coordinate\textemdash{}most plausibly an anchored seed coordinate together with an archimedean/3-adic carry constraint\textemdash{}and a proof of nonescape of mass. That is substantially stronger than a routine weighted-shift estimate.

So the outcome of approach 66 is:

\[
\boxed{
\begin{array}{l}
\text{Cycle detector: successful.}\\
\text{Divergent-orbit Weyl detector: successful.}\\
\text{Global unit-circle gap: false for natural/local weights.}\\
\text{Exact repaired criterion: eigenvalue }1\text{ of }K_s
\text{ must be simple.}
\end{array}}
\]

The missing assertion

\[
\ker(I-K_s)=\mathbb C(n^{-s})
\]

is a clean spectral formulation, but I do not have a valid proof of it. Claiming that simplicity from positivity, Perron\textendash{}Frobenius theory, or a spectral-radius bound would silently assume the irreducibility\textemdash{}or equivalently the Collatz connectivity\textemdash{}that must be proved.
\subsection*{References}
\begin{enumerate}[leftmargin=2.2em]
\item \hypertarget{ref:R71:1}{}\hyperlink{cit:R71:1}{\texttt{turn864506view0}}
\item \hypertarget{ref:R71:2}{}\hyperlink{cit:R71:2}{\texttt{turn864506view2}}
\end{enumerate}

\clearpage
\section{Round 72: approach 67 --- Fredholm determinant and trace cancellation}
\subsection*{Prompt}
\begin{quote}
Complete a full proof of the Collatz conjecture using the following approach:

\begin{enumerate}
\item Fredholm determinant and trace cancellation
\end{enumerate}

\textbf{Speculative \textperiodcentered{} P for cycles}

Construct a trace-class regularization whose k-th trace counts admissible closed parity words with arithmetic weights. Prove exact cancellation of all nontrivial traces and identify the resulting Fredholm determinant.

Take as long as you need.
\end{quote}
\subsection*{Response}

\subsubsection*{Outcome}


I did \textbf{not} obtain a valid full proof of the Collatz conjecture. Current mathematical sources in 2026 still describe the problem as open, so I will not present the equivalences below as a solution. [\hypertarget{cit:R72:1}{}\hyperlink{ref:R72:1}{1}, \hypertarget{cit:R72:2}{}\hyperlink{ref:R72:2}{2}]

What can be completed rigorously is stronger than a merely formal proposal:

\begin{enumerate}
\item There is a natural trace-class Abel regularization of the Collatz adjacency operator.
\item Its traces count closed integer trajectories with strictly positive arithmetic weights.
\item Its Fredholm determinant has an exact Euler product, one factor for each primitive positive Collatz cycle.
\item The desired trace-cancellation identity is \textbf{equivalent} to the assertion that the trivial cycle is the only positive cycle.
\item In a faithful positive-weight construction, the nontrivial traces cannot cancel: they are sums of positive terms.
\item Every trace-class local regularization is spectrally blind to divergent acyclic trajectories.
\item A rank-one boundary closure extends the determinant method to an exact determinant formulation of the \textbf{full} Collatz conjecture, but proving the resulting boundary determinant identity is precisely the missing nonescape theorem.
\end{enumerate}

Thus approach 67 successfully produces the determinant, but it does not supply the hoped-for cancellation mechanism.


\bigskip



\subsection*{1. The trace-class Collatz operator}


Use the shortcut Collatz map

\[
T(n)=
\begin{cases}
n/2,&n\equiv0\pmod 2,\\[2mm]
(3n+1)/2,&n\equiv1\pmod 2.
\end{cases}
\]

Its trivial positive cycle is

\[
\mathcal C_0=(1,2).
\]

Let \(\varepsilon(n)\in\{0,1\}\) denote the parity of \(n\). Work on

\[
\mathcal H=\ell^2(\mathbb N)
\]

with standard orthonormal basis \((e_n)_{n\ge1}\).

For parameters

\[
0<\rho<1,\qquad u>0,
\]

define

\[
\boxed{\quad
\mathcal L_{\rho,u}e_n
   =\rho^n u^{\varepsilon(n)}e_{T(n)}.
\quad}
\]

The parameter \(\rho\) is an Abel regularization in the size of the integer, while \(u\) records odd steps.

An unweighted operator of the form \(e_n\mapsto e_{f(n)}\) has already appeared in Mori's operator-algebraic formulation; Mori proves that suitable irreducibility conditions are equivalent to the full Collatz conjecture. The point of inserting \(\rho^n\) here is specifically to make the operator trace class and hence give it a Fredholm determinant. [\hypertarget{cit:R72:3}{}\hyperlink{ref:R72:3}{3}, \hypertarget{cit:R72:4}{}\hyperlink{ref:R72:4}{4}]


\subsubsection*{Proposition 1.1 \textemdash{} Trace-class regularity}


For every \(0<\rho<1\) and \(u\in\mathbb C\), the operator \(\mathcal L_{\rho,u}\) is trace class.


\paragraph*{Proof}


We have the nuclear expansion

\[
\mathcal L_{\rho,u}
 =
 \sum_{n=1}^{\infty}
 \rho^n u^{\varepsilon(n)}
 \lvert e_{T(n)}\rangle\langle e_n\rvert .
\]

Each summand has trace norm

\[
\rho^n |u|^{\varepsilon(n)}.
\]

Consequently,

\[
\|\mathcal L_{\rho,u}\|_1
 \le
 \sum_{n=1}^{\infty}\rho^n |u|^{\varepsilon(n)}
 \le
 \max(1,|u|)
 \frac{\rho}{1-\rho}
 <\infty.
\]

So \(\mathcal L_{\rho,u}\) is trace class. \ensuremath{\blacksquare}


\bigskip



\subsection*{2. Exact trace formula}


Define the orbit statistics

\[
H_k(n)=\sum_{j=0}^{k-1}T^j(n),
\qquad
S_k(n)=\sum_{j=0}^{k-1}\varepsilon(T^j(n)).
\]

Here \(H_k(n)\) is the sum of the integer values traversed, and \(S_k(n)\) is the number of odd steps.


\subsubsection*{Proposition 2.1 \textemdash{} Powers of the regularized shift}


For every \(k\ge1\),

\[
\mathcal L_{\rho,u}^{\,k}e_n
 =
 \rho^{H_k(n)}
 u^{S_k(n)}
 e_{T^k(n)}.
\]


\paragraph*{Proof}


This follows immediately by induction, since the weight accumulated at the \(j\)-th step is

\[
\rho^{T^j(n)}u^{\varepsilon(T^j(n))}.
\]

Multiplying the \(k\) successive weights gives the stated formula. \ensuremath{\blacksquare}

Taking diagonal matrix coefficients gives the promised trace formula.


\subsubsection*{Theorem 2.2 \textemdash{} Closed-orbit trace formula}


For every \(k\ge1\),

\[
\boxed{\quad
\operatorname{Tr}\mathcal L_{\rho,u}^{\,k}
 =
 \sum_{\substack{n\ge1\\T^k(n)=n}}
 \rho^{H_k(n)}u^{S_k(n)}.
\quad}
\]


\paragraph*{Proof}


Since \(\mathcal L_{\rho,u}^{\,k}\) is trace class,

\[
\operatorname{Tr}\mathcal L_{\rho,u}^{\,k}
 =
 \sum_{n\ge1}
 \left\langle
 \mathcal L_{\rho,u}^{\,k}e_n,e_n
 \right\rangle .
\]

By Proposition 2.1, the diagonal coefficient is zero unless \(T^k(n)=n\), in which case it equals

\[
\rho^{H_k(n)}u^{S_k(n)}.
\]

The diagonal sum is absolutely convergent because the operator is trace class. \ensuremath{\blacksquare}

This is a genuine trace formula: transient paths do not appear, while every periodic point contributes a positive arithmetic weight.


\bigskip



\subsection*{3. Grouping traces by primitive cycles}


Let \(\mathfrak C\) denote the set of primitive positive cycles of \(T\). For a primitive cycle \(\mathcal C\), define

\[
\ell(\mathcal C)=|\mathcal C|,
\qquad
H(\mathcal C)=\sum_{n\in\mathcal C}n,
\qquad
s(\mathcal C)=\#\{n\in\mathcal C:n\text{ odd}\},
\]

and

\[
W_{\rho,u}(\mathcal C)
 =
 \rho^{H(\mathcal C)}u^{s(\mathcal C)}.
\]

If \(k=m\ell(\mathcal C)\), each of the \(\ell(\mathcal C)\) possible starting points on the cycle contributes

\[
W_{\rho,u}(\mathcal C)^m.
\]

Therefore

\[
\boxed{\quad
\operatorname{Tr}\mathcal L_{\rho,u}^{\,k}
 =
 \sum_{\substack{\mathcal C\in\mathfrak C\\
                   \ell(\mathcal C)\mid k}}
 \ell(\mathcal C)\,
 W_{\rho,u}(\mathcal C)^{k/\ell(\mathcal C)}.
\quad}
\]

For the trivial cycle,

\[
\ell(\mathcal C_0)=2,\qquad
H(\mathcal C_0)=1+2=3,\qquad
s(\mathcal C_0)=1,
\]

so

\[
W_{\rho,u}(\mathcal C_0)=\rho^3u.
\]

Its contribution to the \(k\)-th trace is therefore

\[
\mathbf 1_{2\mid k}\,2(\rho^3u)^{k/2}.
\]


\bigskip



\subsection*{4. The Fredholm determinant}


Since \(\mathcal L_{\rho,u}\) is trace class, its Fredholm determinant is

\[
D_{\rho,u}(z)
 =
 \det(I-z\mathcal L_{\rho,u}).
\]

Near \(z=0\),

\[
\log D_{\rho,u}(z)
 =
 -\sum_{k=1}^{\infty}
 \frac{z^k}{k}
 \operatorname{Tr}\mathcal L_{\rho,u}^{\,k}.
\]

Substituting the cycle expansion gives

\[
\begin{aligned}
\log D_{\rho,u}(z)
&=
-\sum_{\mathcal C\in\mathfrak C}
 \sum_{m=1}^{\infty}
 \frac{
   z^{m\ell(\mathcal C)}
   W_{\rho,u}(\mathcal C)^m
 }{m}  \\
&=
\sum_{\mathcal C\in\mathfrak C}
 \log\left(
 1-z^{\ell(\mathcal C)}
 W_{\rho,u}(\mathcal C)
 \right).
\end{aligned}
\]

Hence:


\subsubsection*{Theorem 4.1 \textemdash{} Exact Fredholm Euler product}


\[
\boxed{\quad
D_{\rho,u}(z)
 =
 \prod_{\mathcal C\in\mathfrak C}
 \left(
 1-z^{\ell(\mathcal C)}
   \rho^{H(\mathcal C)}
   u^{s(\mathcal C)}
 \right).
\quad}
\]

The product converges normally on compact subsets of the \(z\)-plane.


\paragraph*{Convergence argument}


Fix \(|z|\le R\) and put \(M=\max(1,|u|)\). A cycle is supported on a finite set of distinct positive integers, and

\[
|z|^{\ell(\mathcal C)}
|\rho|^{H(\mathcal C)}
|u|^{s(\mathcal C)}
\le
\prod_{n\in\mathcal C}
RM\rho^n.
\]

Thus

\[
\sum_{\mathcal C}
R^{\ell(\mathcal C)}
\rho^{H(\mathcal C)}
M^{\ell(\mathcal C)}
\le
\prod_{n=1}^{\infty}
\left(1+RM\rho^n\right)-1
<\infty.
\]

So the Euler product is normally convergent.

Separating the trivial cycle gives

\[
\boxed{\quad
D_{\rho,u}(z)
 =
 \left(1-z^2\rho^3u\right)
 D_{\rho,u}^{\mathrm{rel}}(z),
\quad}
\]

where

\[
D_{\rho,u}^{\mathrm{rel}}(z)
 =
 \prod_{\mathcal C\ne\mathcal C_0}
 \left(
 1-z^{\ell(\mathcal C)}
   \rho^{H(\mathcal C)}
   u^{s(\mathcal C)}
 \right).
\]

This is the exact infinite-dimensional counterpart of the finite determinant factorization obtained in approach 63.


\bigskip



\subsection*{5. Exact equivalence with positive-cycle exclusion}


Define the relative traces

\[
R_k(\rho,u)
 =
 \operatorname{Tr}\mathcal L_{\rho,u}^{\,k}
 -
 \mathbf1_{2\mid k}\,2(\rho^3u)^{k/2}.
\]

Then

\[
\boxed{\quad
R_k(\rho,u)
 =
 \sum_{\substack{\mathcal C\ne\mathcal C_0\\
                  \ell(\mathcal C)\mid k}}
 \ell(\mathcal C)
 \left(
 \rho^{H(\mathcal C)}
 u^{s(\mathcal C)}
 \right)^{k/\ell(\mathcal C)}.
\quad}
\]

For \(0<\rho<1\) and \(u>0\), every term on the right is nonnegative.


\subsubsection*{Theorem 5.1 \textemdash{} Fredholm criterion for cycles}


For any one fixed pair \(0<\rho<1\), \(u>0\), the following are equivalent:

\[
\begin{aligned}
&\text{The trivial cycle }(1,2)\text{ is the only positive cycle};\\
&D_{\rho,u}(z)=1-z^2\rho^3u;\\
&D_{\rho,u}^{\mathrm{rel}}(z)\equiv1;\\
&R_k(\rho,u)=0\quad\text{for every }k\ge1.
\end{aligned}
\]


\paragraph*{Proof}


If there are no nontrivial cycles, the Euler product contains only the trivial factor.

Conversely, suppose \(\mathcal C\ne\mathcal C_0\) is a primitive cycle of length \(\ell\). Then

\[
R_\ell(\rho,u)
\ge
\ell\,\rho^{H(\mathcal C)}u^{s(\mathcal C)}
>0.
\]

Thus the relative traces cannot all vanish, and the relative determinant cannot be identically one. \ensuremath{\blacksquare}

This completes the determinant construction exactly.

It does \textbf{not}, however, prove the needed identity.


\bigskip



\subsection*{6. Why the proposed trace cancellation does not occur}


The intended step was to ``prove exact cancellation of all nontrivial traces.'' In the faithful regularization above, there are no opposing terms to cancel:

\[
R_k(\rho,u)\ge0.
\]

A hypothetical cycle contributes positively to every trace whose order is a multiple of its length.

Moreover, within one primitive cycle \(\mathcal C\), every cyclic starting point contributes the same product of local weights. Its total contribution to the trace at its primitive length is

\[
\ell(\mathcal C)
\prod_{n\in\mathcal C}
\rho^n u^{\varepsilon(n)},
\]

not an alternating sum. Hence cyclic rotation cannot create cancellation.

Changing to signed or complex weights does not solve the logical problem. It can make the product around a genuine cycle zero or arrange cancellation between different cycles, but then vanishing traces no longer imply that cycles are absent.

This can be made completely formal by assigning an independent variable \(x_n\) to every vertex. The cycle factor becomes

\[
1-z^{\ell(\mathcal C)}
\prod_{n\in\mathcal C}x_n.
\]

Different cycles carry different vertex monomials. Consequently, no universal algebraic identity can remove a genuine cycle factor. Any specialization that does remove it has made the detector blind to that cycle.

Therefore:

\[
\boxed{
\text{For a faithful cycle detector, exact trace cancellation is
precisely the cycle-exclusion theorem itself.}
}
\]

It is not a consequence of trace class, cyclicity of trace, or Fredholm determinant algebra.


\bigskip



\subsection*{7. The parity-word arithmetic form of the same obstruction}


Let

\[
\boldsymbol\epsilon
 =
(\epsilon_0,\ldots,\epsilon_{k-1})
\in\{0,1\}^k,
\qquad
s=\sum_{j=0}^{k-1}\epsilon_j.
\]

The corresponding affine branch composition satisfies

\[
2^k T_{\boldsymbol\epsilon}^{\,k}(x)
 =
 3^s x+C(\boldsymbol\epsilon),
\]

where

\[
\boxed{\quad
C(\boldsymbol\epsilon)
 =
 \sum_{j=0}^{k-1}
 \epsilon_j\,2^j\,
 3^{\epsilon_{j+1}+\cdots+\epsilon_{k-1}}.
\quad}
\]

Thus the unique rational fixed candidate is

\[
x_{\boldsymbol\epsilon}
 =
 \frac{C(\boldsymbol\epsilon)}
      {2^k-3^s}.
\]

This is the classical parity-vector fixed-point formula: each finite periodic parity vector determines a unique rational periodic candidate. [\hypertarget{cit:R72:5}{}\hyperlink{ref:R72:5}{5}]

Put

\[
d_{k,s}=2^k-3^s.
\]

For \(s\ge1\), the word gives a positive integer periodic point exactly when

\[
d_{k,s}>0
\quad\text{and}\quad
d_{k,s}\mid C(\boldsymbol\epsilon).
\]

The parity consistency is automatic. Indeed, if

\[
2^k n=3^s n+C(\boldsymbol\epsilon),
\]

then modulo \(2\),

\[
n\equiv\epsilon_0\pmod2.
\]

After applying the first prescribed branch and dividing by \(2\), the same argument applies to the suffix word. Induction recovers every prescribed parity.

Let \(N_k\) be the number of positive fixed points of \(T^k\). The root-of-unity divisor filter gives

\[
\boxed{
\begin{aligned}
N_k
={}&
\sum_{\substack{1\le s\le k\\3^s<2^k}}
\frac{1}{d_{k,s}}
\sum_{r=0}^{d_{k,s}-1}
\sum_{\substack{
\boldsymbol\epsilon\in\{0,1\}^k\\
|\boldsymbol\epsilon|=s}}
\exp\left(
\frac{2\pi i r C(\boldsymbol\epsilon)}
     {d_{k,s}}
\right).
\end{aligned}
}
\]

The desired cycle identity is

\[
\boxed{\quad
N_k=1+(-1)^k.
\quad}
\]

Indeed, the right side equals \(0\) for odd \(k\) and \(2\) for even \(k\), accounting for the two points \(1,2\) of the trivial two-cycle.

But the Fourier cancellation occurs only word by word:

\[
\frac1d\sum_{r=0}^{d-1}e^{2\pi i rC/d}
=
\begin{cases}
1,&d\mid C,\\
0,&d\nmid C.
\end{cases}
\]

For an inadmissible word, the character sum cancels. For a genuine cycle word, every phase is \(1\), so it survives with weight \(+1\).

Thus the arithmetic trace-cancellation proposal becomes the exact divisibility assertion

\[
2^k-3^s\nmid C(\boldsymbol\epsilon)
\]

for every nonalternating positive-drift word. That is exactly the unresolved positive-cycle problem in parity-word form.


\bigskip



\subsection*{8. A trace-class regularization cannot detect divergent orbits}


There is a second, more structural obstruction.


\subsubsection*{Theorem 8.1 \textemdash{} Escape is invisible to local trace-class shifts}


Let \(F:\mathbb N\to\mathbb N\), and suppose

\[
Le_n=a_ne_{F(n)}
\]

is compact; in particular, suppose it is trace class. If

\[
n_0,n_1,n_2,\ldots,
\qquad n_{j+1}=F(n_j),
\]

is an infinite orbit with no repetitions, then the restriction of \(L\) to its orbit is quasinilpotent.


\paragraph*{Proof}


Because \(L\) is compact and \(e_n\rightharpoonup0\),

\[
\|Le_n\|=|a_n|\longrightarrow0.
\]

An infinite nonrepeating orbit in \(\mathbb N\) must satisfy \(n_j\to\infty\): otherwise infinitely many terms would lie in a finite set, forcing a repetition.

Hence

\[
a_{n_j}\longrightarrow0.
\]

On the orbit subspace,

\[
L^m e_{n_j}
 =
 \left(
 \prod_{r=0}^{m-1}a_{n_{j+r}}
 \right)e_{n_{j+m}}.
\]

Given \(\eta>0\), all sufficiently late orbit weights have modulus below \(\eta\). Only finitely many earlier weights can occur in a path without repetition. Therefore

\[
\limsup_{m\to\infty}\|L^m\|^{1/m}\le\eta.
\]

Since \(\eta\) is arbitrary, the spectral radius is zero. \ensuremath{\blacksquare}

For \(\mathcal L_{\rho,u}\), this is especially transparent: along an orbit tending to infinity,

\[
\rho^{n_j}u^{\varepsilon(n_j)}\longrightarrow0.
\]

Therefore a divergent acyclic Collatz orbit contributes:

\begin{itemize}
\item no diagonal term to any trace;
\item no Euler factor to the Fredholm determinant;
\item no nonzero eigenvalue.
\end{itemize}

So even the desired identity

\[
D_{\rho,u}(z)=1-z^2\rho^3u
\]

would prove only that there are no nontrivial positive cycles. It would not rule out divergence.

Tao's almost-all theorem gives extremely strong distributional control, but it does not supply the pointwise nonescape statement required here. [\hypertarget{cit:R72:6}{}\hyperlink{ref:R72:6}{6}]


\bigskip



\subsection*{9. Extending the determinant to detect escape}


There is an exact way to add the missing boundary information.

Fix \(N\ge3\), and let

\[
V_N=\{3,4,\ldots,N\}.
\]

Define the killed finite Collatz operator \(K_N\) on \(\ell^2(V_N)\) by

\[
K_Ne_j=
\begin{cases}
e_{T(j)},&T(j)\in V_N,\\
0,&T(j)\notin V_N.
\end{cases}
\]

Define the upward-boundary functional

\[
\beta_N(e_j)
 =
 \mathbf1_{\{T(j)>N\}}.
\]

For \(n\in V_N\), let

\[
\boxed{\quad
E_{n,N}(z)
 =
 z\,\beta_N(I-zK_N)^{-1}e_n.
\quad}
\]

Interpreting the inverse formally,

\[
E_{n,N}(z)
 =
 \sum_{r\ge1}
 z^r\,\beta_NK_N^{r-1}e_n.
\]

Because the dynamics is deterministic:

\[
E_{n,N}(z)=
\begin{cases}
z^r,&\text{if the orbit first exits above }N
       \text{ at time }r\\
   &\text{before reaching }\{1,2\},\\[1mm]
0,&\text{if it reaches }\{1,2\}
    \text{ before such an exit}.
\end{cases}
\]

If the orbit enters a nontrivial cycle inside \(V_N\), then \(E_{n,N}=0\) as well; this is why cycle exclusion must be imposed separately.


\subsubsection*{Rank-one determinant identity}


Adjoin an artificial edge sending every upward exit back to the starting point \(n\):

\[
\widetilde K_{n,N}
 =
 K_N+e_n\beta_N.
\]

The matrix determinant lemma gives

\[
\begin{aligned}
\det(I-z\widetilde K_{n,N})
&=
\det(I-zK_N)
\left(
1-z\beta_N(I-zK_N)^{-1}e_n
\right).
\end{aligned}
\]

Therefore

\[
\boxed{\quad
E_{n,N}(z)
 =
 1-
 \frac{
 \det(I-z\widetilde K_{n,N})
 }{
 \det(I-zK_N)
 }.
\quad}
\]

If there are no nontrivial cycles, then \(K_N\) is nilpotent and

\[
\det(I-zK_N)=1.
\]

If the orbit exits at time \(r\), the artificial boundary edge creates exactly one finite cycle of length \(r\), so

\[
\det(I-z\widetilde K_{n,N})=1-z^r.
\]

If the orbit reaches the trivial basin, no artificial cycle is created and the determinant remains \(1\).

This converts escape into another exact determinant calculation.


\bigskip



\subsection*{10. A full determinant-equivalent formulation of Collatz}


The following is a complete equivalence, though not a proof.


\subsubsection*{Theorem 10.1 \textemdash{} Fredholm-plus-boundary formulation}


The Collatz conjecture is equivalent to the conjunction of:


\paragraph*{Cycle determinant identity}


For one fixed pair \(0<\rho<1\), \(u>0\),

\[
\boxed{\quad
\det(I-z\mathcal L_{\rho,u})
 =
 1-z^2\rho^3u.
\quad}
\]


\paragraph*{Boundary determinant identity}


For every \(n\ge3\), there exists \(N\ge n\) such that

\[
\boxed{\quad
\det(I-z\widetilde K_{n,N})
 =
 \det(I-zK_N)
\quad}
\]

as polynomials in \(z\).


\paragraph*{Proof}


Assume Collatz. There are no positive cycles other than \((1,2)\), so the first identity follows from the Euler product. Given \(n\), choose \(N\) at least as large as the maximum value on its finite trajectory before it reaches \(1\). There is no upward exit, so \(E_{n,N}=0\), giving the second identity.

Conversely, the first identity excludes every nontrivial positive cycle. For a given \(n\), choose \(N\) from the second condition. Then \(E_{n,N}=0\). If the orbit did not reach \(\{1,2\}\), it would remain forever in the finite set \(V_N\), because it does not exit above \(N\). It would therefore repeat a value and enter a nontrivial cycle, contradicting the first identity. Hence it reaches \(\{1,2\}\), and therefore reaches \(1\). \ensuremath{\blacksquare}

So the full conjecture can indeed be written entirely in determinant language. The unresolved statements are now exactly visible:

\[
\begin{array}{c}
D_{\rho,u}^{\mathrm{rel}}\equiv1
\qquad\text{and}\qquad
E_{n,N}\equiv0\text{ for some }N\text{ for every }n.
\end{array}
\]


\bigskip



\subsection*{11. Why ordinary analytic limits still miss escape}


Suppose cycle exclusion has been proved and an orbit diverges. Let \(r_N\) be its first exit time above \(N\). Then

\[
E_{n,N}(z)=z^{r_N}.
\]

As \(N\to\infty\),

\[
r_N\to\infty.
\]

For every \(|z|<1\),

\[
z^{r_N}\longrightarrow0.
\]

But a convergent orbit also gives eventual value \(E_{n,N}(z)=0\). Thus, inside the open unit disk, the analytic limits of convergent and divergent orbits are identical.

The distinction exists only at the boundary:

\[
E_{n,N}(1)=
\begin{cases}
1,&\text{upward exit},\\
0,&\text{arrival in the trivial basin}.
\end{cases}
\]

This proves a useful no-go statement:

\[
\boxed{
\text{Local-uniform Fredholm convergence for }|z|<1
\text{ cannot by itself prove nonescape.}
}
\]

One must establish exact eventual vanishing, or obtain genuinely new control at \(z=1\).

There is a possible synthesis with the natural-boundary approach. For a divergent orbit, closing successive record exits produces artificial factors

\[
1-z^{r_1},\quad1-z^{r_2},\quad\ldots .
\]

An infinite product over record times would have roots of unity of unbounded orders accumulating around the unit circle. A proof that an independently constructed boundary determinant continued nonvanishingly across an arc would rule out such an orbit. At present, the required continuation theorem is missing.


\bigskip



\subsection*{Final assessment}


Approach 67 yields the exact and natural Fredholm determinant

\[
\boxed{
\det(I-z\mathcal L_{\rho,u})
 =
 \prod_{\mathcal C}
 \left(
 1-z^{\ell(\mathcal C)}
 \rho^{H(\mathcal C)}
 u^{s(\mathcal C)}
 \right).
}
\]

Its trivial Euler factor is

\[
1-z^2\rho^3u.
\]

The hoped-for identity

\[
\det(I-z\mathcal L_{\rho,u})
 =
 1-z^2\rho^3u
\]

is exactly equivalent to positive-cycle exclusion. In the faithful positive-weight model, the relative traces are

\[
R_k(\rho,u)
 =
 \sum_{\mathcal C\ne\mathcal C_0}
 \ell(\mathcal C)
 \left(
 \rho^{H(\mathcal C)}
 u^{s(\mathcal C)}
 \right)^{k/\ell(\mathcal C)}
 \ge0,
\]

so there is no hidden sign cancellation to exploit.

Moreover, trace-class regularization necessarily sends the weights on a divergent orbit to zero, making every divergent acyclic component quasinilpotent and invisible to the Fredholm determinant.

The rigorous conclusion is therefore:

\[
\boxed{
\begin{array}{l}
\text{The Fredholm construction is complete.}\\
\text{The proposed trace cancellation is equivalent to cycle exclusion.}\\
\text{The determinant alone cannot establish nonescape.}\\
\text{No full proof of Collatz results from approach 67 as presently stated.}
\end{array}
}
\]

The genuinely strengthened version is the Fredholm-plus-boundary formulation in Theorem 10.1; completing it requires either a new arithmetic proof of the positive trace identities or a boundary theorem forcing exact disappearance of every artificial escape factor.
\subsection*{References}
\begin{enumerate}[leftmargin=2.2em]
\item \hypertarget{ref:R72:1}{}\hyperlink{cit:R72:1}{\texttt{turn465927search1}}
\item \hypertarget{ref:R72:2}{}\hyperlink{cit:R72:2}{\texttt{turn733098view0}}
\item \hypertarget{ref:R72:3}{}\hyperlink{cit:R72:3}{\texttt{turn340692view0}}
\item \hypertarget{ref:R72:4}{}\hyperlink{cit:R72:4}{\texttt{turn340692view1}}
\item \hypertarget{ref:R72:5}{}\hyperlink{cit:R72:5}{\texttt{turn529457view0}}
\item \hypertarget{ref:R72:6}{}\hyperlink{cit:R72:6}{\texttt{turn130826search19}}
\end{enumerate}

\clearpage
\section{Round 73: approach 68 --- Numen-integrality exclusion}
\subsection*{Prompt}
\begin{quote}
Complete a full proof of the Collatz conjecture using the following approach:

\begin{enumerate}
\item Numen-integrality exclusion
\end{enumerate}

\textbf{Established/Hybrid \textperiodcentered{} B}

Recent "numen" formalisms assign multi-place values to symbolic Collatz sequences. Prove that every nontrivial periodic or nondecaying symbolic word has a nonintegral value at at least one completion. A symbolic sequence passing all local integrality tests would be a counterexample candidate.

Take as long as you need.
\end{quote}
\subsection*{Response}

\subsection*{Numen-integrality exclusion: exact outcome}


The proposed argument does \textbf{not} give a valid full proof of the Collatz conjecture. Formalizing it leads to one strong and useful theorem, but it also exposes two decisive obstructions:

\begin{enumerate}
\item For nonperiodic words, the natural numen is essentially a \textbf{backward-orbit invariant}, not a forward-orbit invariant. Positive integer numen values can have many irrational symbolic codes, including codes attached to integers already in the \(1\leftrightarrow2\) basin.
\item For periodic words, ``integral at every finite completion'' is \textbf{exactly equivalent} to being an integer Collatz cycle. Thus the desired local-integrality exclusion is not an easier lemma implying cycle exclusion; it is another formulation of cycle exclusion.
\end{enumerate}

The recent Hydra/numen preprint defines numens from finite affine compositions and then extends them separately to completions in which the corresponding limits converge. It explicitly notes that a fully rigorous adelic formulation remains future work. [\hypertarget{cit:R73:1}{}\hyperlink{ref:R73:1}{1}, \hypertarget{cit:R73:2}{}\hyperlink{ref:R73:2}{2}, \hypertarget{cit:R73:3}{}\hyperlink{ref:R73:3}{3}]

What follows is a complete proof of the strongest correct theorem obtainable from this approach, followed by explicit counterexamples to the missing correspondence.


\bigskip



\subsubsection*{1. The numen attached to a binary word}


Use the shortened Collatz map

\[
T(n)=
\begin{cases}
n/2,&n\equiv0\pmod2,\\[1mm]
(3n+1)/2,&n\equiv1\pmod2.
\end{cases}
\]

Write its two affine branches as

\[
H_0(x)=\frac{x}{2},
\qquad
H_1(x)=\frac{3x+1}{2}.
\]

Let

\[
\varepsilon=(\varepsilon_0,\varepsilon_1,\ldots)\in\{0,1\}^{\mathbb N},
\qquad
s_j=\sum_{i=0}^{j-1}\varepsilon_i,
\qquad
s_0=0.
\]

For a prefix of length \(N\), define

\[
M_N(\varepsilon)=\frac{3^{s_N}}{2^N}
\]

and

\[
\chi_N(\varepsilon)
=
\sum_{j=0}^{N-1}
\varepsilon_j\,\frac{3^{s_j}}{2^{j+1}}.
\]

A direct induction gives the fundamental affine identity

\[
H_{\varepsilon_0}\circ H_{\varepsilon_1}\circ\cdots
\circ H_{\varepsilon_{N-1}}(x)
=
\chi_N(\varepsilon)+M_N(\varepsilon)x.
\tag{1}
\]

The order here is important: the least significant symbolic digit describes the \textbf{last} branch of an ordinary forward trajectory. Thus the infinite numen code naturally parametrizes backward paths.

Suppose the \(1\)'s of \(\varepsilon\) occur at positions

\[
d_1<d_2<d_3<\cdots.
\]

Then

\[
\chi_N(\varepsilon)\longrightarrow
\chi_3(\varepsilon)
=
\sum_{r=1}^{\infty}
\frac{3^{r-1}}{2^{d_r+1}}
\quad\text{in }\mathbb Q_3.
\tag{2}
\]

Indeed, the \(r\)-th nonzero summand has \(3\)-adic valuation \(r-1\). Consequently,

\[
\boxed{\chi_3(\varepsilon)\in\mathbb Z_3}
\tag{3}
\]

for every infinite binary word; if there are only finitely many \(1\)'s, the sum is finite and is still in \(\mathbb Z_3\).

Thus the natural \(3\)-adic integrality test is not restrictive: \textbf{every word passes it}.


\bigskip



\subsubsection*{2. Exact integer-fiber theorem}


The following is the correct dynamical meaning of an ordinary positive integer numen value.


\paragraph*{Theorem 1 \textemdash{} Integer numen fibers are infinite backward orbits}


Let \(\varepsilon\) contain infinitely many \(1\)'s, and let \(n\in\mathbb N\). Then

\[
\chi_3(\varepsilon)=n
\tag{4}
\]

if and only if the recursion

\[
x_0=n,
\qquad
x_{j+1}=
\begin{cases}
2x_j,&\varepsilon_j=0,\\[2mm]
\dfrac{2x_j-1}{3},&\varepsilon_j=1
\end{cases}
\tag{5}
\]

produces positive integers for every \(j\).

In that case,

\[
T(x_{j+1})=x_j
\tag{6}
\]

for every \(j\). Hence \(\varepsilon\) is precisely an infinite backtracing parity vector starting at \(n\).


\paragraph*{Proof}


Let \(\sigma\) denote the left shift, and put

\[
x_j=\chi_3(\sigma^j\varepsilon).
\]

Splitting off the first numen digit gives the functional equations

\[
x_j=
\begin{cases}
x_{j+1}/2,&\varepsilon_j=0,\\[1mm]
(3x_{j+1}+1)/2,&\varepsilon_j=1.
\end{cases}
\tag{7}
\]

Assume first that \(x_0=n\in\mathbb N\).

If \(\varepsilon_j=0\), then \(x_{j+1}=2x_j\), so \(x_{j+1}\) is a positive integer.

If \(\varepsilon_j=1\), then

\[
x_{j+1}=\frac{2x_j-1}{3}.
\]

But \(x_{j+1}\in\mathbb Z_3\), because every shifted numen lies in \(\mathbb Z_3\). Therefore \(2x_j-1\) is divisible by \(3\) as an ordinary integer. It follows that \(x_{j+1}\) is an ordinary positive integer. Induction proves (5), and the parity of the two displayed preimages proves (6).

Conversely, suppose (5) always produces positive integers. Iterating (7), or applying (1), gives

\[
n=\chi_N(\varepsilon)+M_N(\varepsilon)x_N.
\tag{8}
\]

Because \(\varepsilon\) has infinitely many \(1\)'s,

\[
s_N\longrightarrow\infty.
\]

Since \(x_N\) is an integer,

\[
v_3\!\left(M_Nx_N\right)
=
s_N+v_3(x_N)
\geq s_N
\longrightarrow\infty.
\]

Hence \(M_Nx_N\to0\) in \(\mathbb Q_3\). Passing to the \(3\)-adic limit in (8) gives

\[
\chi_3(\varepsilon)=n.
\]

This proves the theorem. \(\square\)

This interpretation agrees with the established theory of infinite backtracing parity vectors in the \(3x+1\) graph. In particular, previous work distinguishes finite, immediately periodic, and irrational backtracing codes and explicitly constructs irrational examples. [\hypertarget{cit:R73:4}{}\hyperlink{ref:R73:4}{4}, \hypertarget{cit:R73:5}{}\hyperlink{ref:R73:5}{5}, \hypertarget{cit:R73:6}{}\hyperlink{ref:R73:6}{6}]


\bigskip



\subsubsection*{3. The genuine nondecay-exclusion theorem}


We now obtain the strongest successful part of Approach 68.


\paragraph*{Theorem 2 \textemdash{} Positive integer fibers have strictly decaying multiplier}


If

\[
\chi_3(\varepsilon)=n\in\mathbb N,
\]

then

\[
\boxed{\frac{3^{s_N}}{2^N}\longrightarrow0.}
\tag{9}
\]

Equivalently,

\[
N\log 2-s_N\log 3\longrightarrow+\infty.
\tag{10}
\]


\paragraph*{Proof}


By Theorem 1, there are positive integers \(x_N\) satisfying (8):

\[
n=\chi_N+M_Nx_N.
\]

Therefore

\[
0\leq\chi_N<n
\tag{11}
\]

for every \(N\).

At a position \(d_r\) occupied by a \(1\),

\[
\chi_{d_r+1}-\chi_{d_r}
=
\frac{3^{r-1}}{2^{d_r+1}}
=
\frac{M_{d_r}}{2}.
\tag{12}
\]

The left side consists of terms of a bounded increasing real series, by (11). Hence

\[
\frac{M_{d_r}}{2}\longrightarrow0,
\qquad\text{so}\qquad
M_{d_r}\longrightarrow0.
\tag{13}
\]

Between consecutive \(1\)'s, the multiplier is first multiplied by \(3/2\) and then repeatedly divided by \(2\). More explicitly, if

\[
d_r<N\leq d_{r+1},
\]

then

\[
M_N
=
\frac32\,M_{d_r}\,2^{-(N-d_r-1)}
\leq
\frac32 M_{d_r}.
\tag{14}
\]

Together with (13), this proves \(M_N\to0\). \(\square\)

Thus we obtain the unconditional exclusion

\[
\boxed{
\limsup_{N\to\infty}\frac{3^{s_N}}{2^N}>0
\quad\Longrightarrow\quad
\chi_3(\varepsilon)\notin\mathbb N.
}
\tag{15}
\]

This is a real theorem. It rules out every multiplier-nondecaying word as an ordinary positive integer numen fiber.

It does \textbf{not}, however, say that such a word is nonintegral in \(\mathbb Q_3\): equation (3) says the opposite. Its numen is still a perfectly integral \(3\)-adic integer; it is merely not a positive \textbf{diagonal} integer.


\bigskip



\subsubsection*{4. Irrational integer-valued numens are abundant}


A central claim in the recent numen correspondence is that an irrational symbolic code with an integral numen value should give a divergent trajectory. That implication is false for the Collatz numen.

In fact, every positive integer not divisible by \(3\) has continuum many irrational numen codes.


\paragraph*{Theorem 3 \textemdash{} Continuum-sized irrational fibers}


For every \(n\in\mathbb N\) with \(3\nmid n\), there are continuum many irrational

\[
z\in\mathbb Z_2\setminus\mathbb Q
\]

such that

\[
\chi_3(z)=n.
\tag{16}
\]


\paragraph*{Proof}


Start with \(x_0=n\). Suppose \(3\nmid x_r\).

Because \(2\) generates the unit group modulo \(9\), there is an

\[
a_r\in\{0,1,\ldots,5\}
\]

such that

\[
2^{a_r}x_r\equiv2\pmod9.
\]

Choose

\[
t_r=a_r+6r+6b_r,
\qquad b_r\in\{0,1\}.
\tag{17}
\]

Then also

\[
2^{t_r}x_r\equiv2\pmod9.
\]

Append to the symbolic word a block \(0^{t_r}1\), and define

\[
x_{r+1}
=
\frac{2^{t_r+1}x_r-1}{3}.
\tag{18}
\]

The numerator in (18) is \(3\pmod9\), so \(x_{r+1}\) is a positive integer satisfying

\[
x_{r+1}\equiv1\pmod3.
\]

The construction may therefore be repeated indefinitely.

The concatenated word

\[
\varepsilon
=
0^{t_1}1\,0^{t_2}1\,0^{t_3}1\cdots
\tag{19}
\]

is an infinite backtracing word from \(n\). By Theorem 1,

\[
\chi_3(\varepsilon)=n.
\]

Because \(t_r\to\infty\), the binary word has infinitely many \(1\)'s and unbounded gaps between them. It cannot be eventually periodic, so the associated \(2\)-adic integer is irrational.

Finally, each independent binary choice \(b_r\in\{0,1\}\) gives a different word, producing continuum many such codes. \(\square\)


\paragraph*{Explicit example over the trivial basin}


Take \(n=1\). Choose

\[
t_1=7,
\qquad
x_1=\frac{2^8-1}{3}=85,
\]

and then

\[
t_2=17,
\qquad
x_2=\frac{2^{18}\cdot85-1}{3}=7\,427\,413.
\]

The code begins

\[
0^7\,1\,0^{17}\,1\cdots.
\]

After the first two \(1\)'s, identity (8) is

\[
1
=
\frac1{2^8}
+
\frac3{2^{26}}
+
\frac{3^2}{2^{26}}\cdot7\,427\,413.
\tag{20}
\]

At the \(r\)-th completed block, the residual term has \(3\)-adic valuation \(r\), so it tends to zero in \(\mathbb Q_3\). Hence

\[
\chi_3(\varepsilon)=1.
\]

Yet \(1\) is periodic under Collatz, not divergent.

Moreover, for the choices above the same formal series converges in \(\mathbb R\) to

\[
\frac1{2^8}
+
\frac3{2^{26}}
+
\frac9{2^{46}}+\cdots
<\frac1{128},
\tag{21}
\]

not to \(1\). Thus even when a word has both a real and a \(3\)-adic numen sum, those sums need not be the same global number.

There is also an independently known example: the greedy backtracing vector \(V_4\) is irrational, while its numen value is \(4\); nevertheless

\[
4\mapsto2\mapsto1\mapsto2
\]

is preperiodic rather than divergent. [\hyperlink{ref:R73:6}{6}]

The 2026 preprint's Correspondence Principle II states that an irrational code with an integral numen value gives a divergent point. Its proof sketch excludes preperiodicity by invoking Part I, but Part I establishes only that a periodic point has \textbf{some} rational numen code. It does not imply that every code in the same numen fiber is rational. Theorem 3 supplies explicit noninjective fibers, so that inference cannot be used. [\hypertarget{cit:R73:7}{}\hyperlink{ref:R73:7}{7}, \hypertarget{cit:R73:8}{}\hyperlink{ref:R73:8}{8}]


\bigskip



\subsubsection*{5. Periodic words and local integrality}


Now let

\[
w=(\varepsilon_0,\ldots,\varepsilon_{k-1})
\]

be a finite block containing \(s\geq1\) ones, and consider the purely periodic word \(w^\infty\).

Define

\[
C(w)
=
\sum_{j=0}^{k-1}
\varepsilon_j\,3^{s_j}2^{k-j-1},
\qquad
D(w)=2^k-3^s.
\tag{22}
\]

Then

\[
\chi_k(w)=\frac{C(w)}{2^k},
\qquad
M_k(w)=\frac{3^s}{2^k}.
\]

In \(\mathbb Q_3\), the repeated-block series is geometric, because

\[
\left|M_k(w)\right|_3=3^{-s}<1.
\]

Therefore

\[
\chi_3(w^\infty)
=
\frac{\chi_k(w)}{1-M_k(w)}
=
\boxed{\frac{C(w)}{2^k-3^s}}.
\tag{23}
\]

The same rational expression gives the natural algebraic continuation at every completion.

Observe that

\[
\gcd(D(w),6)=1:
\tag{24}
\]

\(D(w)\) is odd, and modulo \(3\) it equals \(2^k\).


\paragraph*{Theorem 4 \textemdash{} Periodic local integrality is exactly cycle integrality}


For a periodic word \(w^\infty\) containing a \(1\), the following are equivalent:

\begin{enumerate}
\item \(C(w)/D(w)\in\mathbb Z_p\) for every finite prime \(p\);
\item \(D(w)\mid C(w)\);
\item \(C(w)/D(w)\) is an ordinary integer;
\item the branches encoded by \(w\) form an integer Collatz cycle.
\end{enumerate}

If additionally \(D(w)>0\), the cycle is positive.


\paragraph*{Proof}


For any rational number \(a/b\) in lowest terms,

\[
a/b\in\mathbb Z_p\quad\text{for every prime }p
\]

if and only if \(b=1\). This proves the equivalence of 1\textendash{}3.

Assume now that

\[
x_0=\frac{C(w)}{D(w)}
\]

is an integer. Let

\[
x_j=\chi_3(\sigma^j w^\infty).
\]

Each \(x_j\) is rational, because the corresponding code is periodic, and each lies in \(\mathbb Z_3\). The functional equations give

\[
x_{j+1}=
\begin{cases}
2x_j,&\varepsilon_j=0,\\[2mm]
\dfrac{2x_j-1}{3},&\varepsilon_j=1.
\end{cases}
\tag{25}
\]

Starting from the integer \(x_0\), induction shows that every \(x_j\) is an ordinary integer: in the second case, membership of \(x_{j+1}\) in \(\mathbb Z_3\) forces divisibility by \(3\).

Since \(w^\infty\) is \(k\)-periodic,

\[
x_{j+k}=x_j.
\]

Furthermore,

\[
T(x_{j+1})=x_j.
\]

Thus the \(x_j\)'s form an integer Collatz cycle. If \(D(w)>0\), then \(C(w)>0\) implies \(x_0>0\), and every member of the cycle is positive.

Conversely, reversing the parity block of any integer Collatz cycle gives a periodic word whose value (23) is the corresponding cycle element, hence an integer. \(\square\)

Consequently, the desired assertion

\[
\text{“Every nontrivial positive periodic word is nonintegral somewhere”}
\tag{26}
\]

is precisely the assertion

\[
\text{“There is no nontrivial positive Collatz cycle.”}
\tag{27}
\]

No simplification has occurred. The unknown divisibility

\[
2^k-3^s\nmid C(w)
\tag{28}
\]

is the usual cycle-exclusion problem in numen notation.

Because of (24), neither the \(2\)-adic nor the \(3\)-adic completion can detect such a denominator. Any detecting prime must satisfy

\[
p\geq5,
\qquad
v_p(2^k-3^s)>v_p(C(w)).
\tag{29}
\]

Showing that such a prime always exists is exactly what must still be proved.

The unrestricted version is false even more directly. For example,

\[
w=1
\]

has

\[
\frac{C(w)}{D(w)}=\frac1{2-3}=-1,
\]

and

\[
w=011
\]

has

\[
\frac{C(w)}{D(w)}=\frac5{8-9}=-5,
\]

corresponding to the negative cycle

\[
-5\mapsto-7\mapsto-10\mapsto-5.
\]

Thus positivity is essential.


\bigskip



\subsubsection*{6. Why there is no simultaneous local test for aperiodic words}


For an infinite word with \(1\)'s at \(d_r\), its nonzero numen terms are

\[
a_r=\frac{3^{r-1}}{2^{d_r+1}}.
\]

Their behavior at the finite completions is:

\[
|a_r|_3=3^{-(r-1)}\longrightarrow0;
\]

\[
|a_r|_2=2^{d_r+1}\longrightarrow\infty;
\]

and, for every prime \(p\geq5\),

\[
|a_r|_p=1.
\]

Therefore the raw numen series

\[
\sum_r a_r
\]

converges in \(\mathbb Q_3\), but does not converge in \(\mathbb Q_2\) or in any \(\mathbb Q_p\) with \(p\geq5\).

If the word is multiplier-nondecaying,

\[
\limsup M_N>0,
\]

its real terms also fail to tend to zero along a subsequence, so the real numen series diverges as well.

Thus a nondecaying aperiodic word does not possess a family of values

\[
(X_\infty,X_2,X_3,X_5,\ldots)
\]

whose local integrality can be checked. It normally has only the \(3\)-adic value, and that value is automatically in \(\mathbb Z_3\).

One could try to introduce separate regularizations at the other places, but they would not automatically form the embeddings of one common rational or algebraic number. This is exactly the compatibility issue that a genuinely adelic numen would have to solve; the current formalism acknowledges that such a fully rigorous adelic construction has not yet been supplied. [\hyperlink{ref:R73:3}{3}]

Even membership in every finite local integer ring would not by itself be enough. The profinite ring

\[
\widehat{\mathbb Z}=\prod_p\mathbb Z_p
\]

contains many non-diagonal elements, and the ordinary integers are dense in it. For example,

\[
N!\longrightarrow0
\quad\text{in every }\mathbb Z_p,
\]

while

\[
N!\longrightarrow+\infty
\quad\text{in }\mathbb R.
\]

Finite-place integrality therefore cannot supply the missing archimedean nonescape bound.


\bigskip



\subsubsection*{7. Final status of Approach 68}


The approach yields the rigorous theorem

\[
\boxed{
\chi_3(\varepsilon)\in\mathbb N
\quad\Longrightarrow\quad
\varepsilon\text{ is a positive backward Collatz path and }
\frac{3^{s_N}}{2^N}\to0.
}
\tag{30}
\]

This completely excludes multiplier-nondecaying words from the positive integer fibers of the \(3\)-adic numen.

But it does not complete Collatz, because:

\begin{itemize}
\item an irrational positive integer numen value does \textbf{not} imply a divergent forward trajectory;
\item periodic all-place integrality is exactly equivalent to the existence of an integer cycle;
\item aperiodic words do not naturally define one common value at all completions;
\item local integrality gives a profinite integer, not necessarily a diagonal ordinary integer;
\item there is no proved bridge taking a divergent forward orbit to a nondecaying word whose numen is a fixed positive ordinary integer.
\end{itemize}

A numen proof would therefore require an additional \textbf{adelic diagonalization and height theorem}: every exceptional forward orbit would have to produce a compatible multi-place numen value coming from one positive rational integer, while retaining nondecaying multiplier. Theorem 2 would then give a contradiction. The explicit irrational fibers above show that this diagonalization cannot follow merely from symbolic irrationality or local integrality.

Accordingly, Approach 68 provides a genuine backward-drift theorem and a precise local-global reformulation of cycles, but the proposed central correspondence is false and the remaining periodic condition is equivalent to the conjecture it was meant to prove.
\subsection*{References}
\begin{enumerate}[leftmargin=2.2em]
\item \hypertarget{ref:R73:1}{}\hyperlink{cit:R73:1}{\texttt{turn959590view0}}
\item \hypertarget{ref:R73:2}{}\hyperlink{cit:R73:2}{\texttt{turn959590view1}}
\item \hypertarget{ref:R73:3}{}\hyperlink{cit:R73:3}{\texttt{turn247946view0}}
\item \hypertarget{ref:R73:4}{}\hyperlink{cit:R73:4}{\texttt{turn191681view0}}
\item \hypertarget{ref:R73:5}{}\hyperlink{cit:R73:5}{\texttt{turn191681view1}}
\item \hypertarget{ref:R73:6}{}\hyperlink{cit:R73:6}{\texttt{turn696963view0}}
\item \hypertarget{ref:R73:7}{}\hyperlink{cit:R73:7}{\texttt{turn423697view0}}
\item \hypertarget{ref:R73:8}{}\hyperlink{cit:R73:8}{\texttt{turn423697view1}}
\end{enumerate}

\end{document}
